% 本文件由 LaTeX 排版 Agent 根据有效 translation.json 自动生成。
% author=Basil the Great; work_id=3202; title=书信集（圣巴西略）

\chapter{书信集（\Person{圣巴西略}）}

% structure: p0001 source=h1 target=omitted reason=page-title-already-rendered-by-parent

\Emph{\Person{圣巴西略}}\par

书信 1\newline{}书信 2\newline{}书信 3\newline{}书信 4\newline{}书信 5\newline{}书信 6\newline{}书信 7\newline{}书信 8\newline{}书信 9\newline{}书信 10\newline{}书信 11\newline{}书信 12\newline{}书信 13\newline{}书信 14\newline{}书信 15\newline{}书信 16\newline{}书信 17\newline{}书信 18\newline{}书信 19\newline{}书信 20\newline{}书信 21\newline{}书信 22\newline{}书信 23\newline{}书信 24\newline{}书信 25\newline{}书信 26\newline{}书信 27\newline{}书信 28\newline{}书信 29\newline{}书信 30\newline{}书信 31\newline{}书信 32\newline{}书信 33\newline{}书信 34\newline{}书信 35\newline{}书信 36\newline{}书信 37\newline{}书信 38\newline{}书信 39\newline{}书信 40\newline{}书信 41\newline{}书信 42\newline{}书信 43\newline{}书信 44\newline{}书信 45\newline{}书信 46\newline{}书信 47\newline{}书信 48\newline{}书信 49\newline{}书信 50\newline{}书信 51\newline{}书信 52\newline{}书信 53\newline{}书信 54\newline{}书信 55\newline{}书信 56\newline{}书信 57\newline{}书信 58\newline{}书信 59\newline{}书信 60\newline{}书信 61\newline{}书信 62\newline{}书信 63\newline{}书信 64\newline{}书信 65\newline{}书信 66\newline{}书信 67\newline{}书信 68\newline{}书信 69\newline{}书信 70\newline{}书信 71\newline{}书信 72\newline{}书信 73\newline{}书信 74\newline{}书信 75\newline{}书信 76\newline{}书信 77\newline{}书信 78\newline{}书信 79\newline{}书信 80\newline{}书信 81\newline{}书信 82\newline{}书信 83\newline{}书信 84\newline{}书信 85\newline{}书信 86\newline{}书信 87\newline{}书信 88\newline{}书信 89\newline{}书信 90\newline{}书信 91\newline{}书信 92\newline{}书信 93\newline{}书信 94\newline{}书信 95\newline{}书信 96\newline{}书信 97\newline{}书信 98\newline{}书信 99\newline{}书信 100\newline{}书信 101\newline{}书信 102\newline{}书信 103\newline{}书信 104\newline{}书信 105\newline{}书信 106\newline{}书信 107\newline{}书信 108\newline{}书信 109\newline{}书信 110\newline{}书信 111\newline{}书信 112\newline{}书信 113\newline{}书信 114\newline{}书信 115\newline{}书信 116\newline{}书信 117\newline{}书信 118\newline{}书信 119\newline{}书信 120\newline{}书信 121\newline{}书信 122\newline{}书信 123\newline{}书信 124\newline{}书信 125\newline{}书信 126\newline{}书信 127\newline{}书信 128\newline{}书信 129\newline{}书信 130\newline{}书信 131\newline{}书信 132\newline{}书信 133\newline{}书信 134\newline{}书信 135\newline{}书信 136\newline{}书信 137\newline{}书信 138\newline{}书信 139\newline{}书信 140\newline{}书信 141\newline{}书信 142\newline{}书信 143\newline{}书信 144\newline{}书信 145\newline{}书信 146\newline{}书信 147\newline{}书信 148\newline{}书信 149\newline{}书信 150\newline{}书信 151\newline{}书信 152\newline{}书信 153\newline{}书信 154\newline{}书信 155\newline{}书信 156\newline{}书信 157\newline{}书信 158\newline{}书信 159\newline{}书信 160\newline{}书信 161\newline{}书信 162\newline{}书信 163\newline{}书信 164\newline{}书信 165\newline{}书信 167\newline{}书信 168\newline{}书信 169\newline{}书信 170\newline{}书信 171\newline{}书信 172\newline{}书信 173\newline{}书信 174\newline{}书信 175\newline{}书信 176\newline{}书信 177\newline{}书信 178\newline{}书信 179\newline{}书信 180\newline{}书信 181\newline{}书信 182\newline{}书信 183\newline{}书信 184\newline{}书信 185\newline{}书信 186\newline{}书信 187\newline{}书信 188\newline{}书信 189\newline{}书信 190\newline{}书信 191\newline{}书信 192\newline{}书信 193\newline{}书信 194\newline{}书信 195\newline{}书信 196\newline{}书信 197\newline{}书信 198\newline{}书信 199\newline{}书信 200\newline{}书信 201\newline{}书信 202\newline{}书信 203\newline{}书信 204\newline{}书信 205\newline{}书信 206\newline{}书信 207\newline{}书信 208\newline{}书信 209\newline{}书信 200\newline{}书信 211\newline{}书信 212\newline{}书信 213\newline{}书信 214\newline{}书信 215\newline{}书信 216\newline{}书信 217\newline{}书信 218\newline{}书信 219\newline{}书信 220\newline{}书信 221\newline{}书信 222\newline{}书信 223\newline{}书信 224\newline{}书信 225\newline{}书信 226\newline{}书信 227\newline{}书信 228\newline{}书信 229\newline{}书信 230\newline{}书信 231\newline{}书信 232\newline{}书信 233\newline{}书信 234\newline{}书信 235\newline{}书信 236\newline{}书信 237\newline{}书信 238\newline{}书信 239\newline{}书信 240\newline{}书信 241\newline{}书信 242\newline{}书信 243\newline{}书信 244\newline{}书信 245\newline{}书信 246\newline{}书信 247\newline{}书信 248\newline{}书信 249\newline{}书信 250\newline{}书信 251\newline{}书信 252\newline{}书信 253\newline{}书信 254\newline{}书信 255\newline{}书信 256\newline{}书信 257\newline{}书信 258\newline{}书信 259\newline{}书信 260\newline{}书信 261\newline{}书信 262\newline{}书信 263\newline{}书信 264\newline{}书信 265\newline{}书信 266\newline{}书信 267\newline{}书信 268\newline{}书信 269\newline{}书信 270\newline{}书信 271\newline{}书信 272\newline{}书信 273\newline{}书信 274\newline{}书信 275\newline{}书信 276\newline{}书信 277\newline{}书信 278\newline{}书信 279\newline{}书信 280\newline{}书信 281\newline{}书信 282\newline{}书信 283\newline{}书信 284\newline{}书信 285\newline{}书信 286\newline{}书信 287\newline{}书信 288\newline{}书信 289\newline{}书信 290\newline{}书信 291\newline{}书信 292\newline{}书信 293\newline{}书信 294\newline{}书信 295\newline{}书信 299\newline{}书信 303\newline{}书信 306\newline{}书信 334\newline{}书信 335\newline{}书信 336\newline{}书信 337\newline{}书信 338\newline{}书信 339\newline{}书信 340\newline{}书信 341\newline{}书信 342\newline{}书信 343\newline{}书信 344\newline{}书信 345\newline{}书信 346\newline{}书信 347\newline{}书信 348\newline{}书信 349\newline{}书信 350\newline{}书信 351\newline{}书信 352\newline{}书信 353\newline{}书信 354\newline{}书信 355\newline{}书信 356\newline{}书信 357\newline{}书信 358\newline{}书信 359\newline{}书信 360\newline{}书信 366\par

\SourceRecord{Translated by Blomfield Jackson.；Nicene and Post-Nicene Fathers, Second Series；Vol. 8.；Edited by Philip Schaff and Henry Wace.；Buffalo, NY: Christian Literature Publishing Co.,；1895.；<http://www.newadvent.org/fathers/3202.htm>.}

% structure: p0001 source=h1 target=section reason=page-title

\section{第一封信}

\Emph{\Person{圣巴西略}\Place{凯撒勒雅}}\par

\Emph{致哲学家 \Person{尤斯塔修斯}}\par

我因所谓命运的嘲弄而深感苦恼，它似乎总在阻挠我与你会面。但你的来信令我无比欣慰和舒心，因为我此前一直在心中思忖：是否真如众人所言，有一种必然性或命运主宰着我们生活中大小诸事，而我们人类对任何事都毫无掌控？抑或，至少所有人生都是由某种运气所驱使？你若得知我为何产生这些想法，定会十分宽恕我。\par

听闻你的哲学后，我便对雅典的老师们心生轻蔑，离开了那里。赫勒斯滂上的那座城，我比任何经过塞壬歌声的尤利西斯都更无动于衷地经过。\par

我钦佩亚细亚，却急忙赶向其中至善的都会。及至抵达家乡，却未见到你——我热切所求的奖赏——于是开始了许多各种各样的意外阻碍。首先，我因病而错失了你；接着，当你启程往东方去时，我无法与你同行；然后，经历了无尽的劳顿，我到了叙利亚，但那位哲学家已动身去了埃及，我未能遇上。于是我必须启程前往埃及，路途遥远而疲惫，即使在那里我也未能达成目的。但我如此热切地渴望，要么我必须启程往波斯去，随你直至蛮荒最远之地（你已到了那里；我中了何等执拗之魔！），要么就此定居亚历山大。最后我选择了后者。我真以为，除非我像某种驯兽一样，追随一根伸出的树枝直至精疲力竭，你才会被驱赶着一直越过印度尼萨或任何更遥远的地区，在那里流浪漂泊。何必多说？\par

返家后，我因缠绵病痛而无法与你见面。若它们不见好转，我甚至在冬天也无法与你相会。这一切，正如你自己所言，岂非命运使然？这岂非必然？我的境遇岂不几乎胜过诗人笔下坦塔罗斯的故事？但我如前所述，收到你的信后感觉好转，如今心意已不再如故。当天主赐下美善时，我认为我们应当感谢祂，而不是在祂掌控分配时向祂生气。因此，若祂应允我与你相聚，我将认为这是最好、最可喜之事；若祂延迟，我将温顺地忍受这份失落。因为祂管理我们的生命，总比我们自己选择的更好。\par

\SourceRecord{Translated by Blomfield Jackson.；Nicene and Post-Nicene Fathers, Second Series；Vol. 8.；Edited by Philip Schaff and Henry Wace.；Buffalo, NY: Christian Literature Publishing Co.,；1895.；<http://www.newadvent.org/fathers/3202001.htm>.}

% structure: p0001 source=h1 target=section reason=page-title

\section{书信第二}

\Emph{\Person{圣巴西略}\Place{凯撒勒雅}}\par

\Emph{\Person{巴西略}致\Person{格列高利}}.\par

1. [我认出你的来信，正如人从显而易见的相似性认出朋友的子女一样。你说在让我听到任何关于我如何生活之前，描述我居住的地方，不足以说服你分享我的生活，这很像你；这配得上像你这样的灵魂，它轻视今生的一切，视为与应许给我们的未来福乐相比毫无价值。我自己在这偏远之地日夜所做的事，我羞于写下。我放弃了城里的生活，因为它必定导致无数祸患；但我尚未能摆脱我自己。我就像初次航海的旅客，晕船受苦，抱怨船太大而颠簸，当他们换到小艇或救生艇上时，到处都晕船受苦。无论他们去哪里，恶心和痛苦都跟着他们。我的处境与此类似。我随身携带自己的烦恼，所以到处我都处于类似的不适中。结果我并未从独处中得到太多好处。我应该做的事，那能使我紧跟随那引领我们救恩者的足迹的事——因为祂说：\BibleQuote{若有人愿意跟随我，该弃绝自己，背起自己的十字架来跟随我} \BibleRef{玛 16:24} ——就是这样.]\par

2. 我们必须追求宁静的心灵。正如眼睛若不安地上下左右游荡，不固定注视，就无法看清放在眼前的对象，同样，被千百种世俗忧虑所分心的心灵，也无法清楚地领会真理。尚未被婚姻束缚的人，受着疯狂的欲望、叛逆的冲动和无望的依恋所困扰；已找到伴侣的人，则被自己的烦忧所包围；若无子女，便渴望子女；若有子女，则忧虑其教育，关注妻子，照顾家庭，监督仆人，商业上的不幸，与邻舍的争吵，诉讼，商人的风险，农夫的辛劳。每一天都以其特有的方式使灵魂暗淡；一夜又一夜接过白天的忧虑，并以相应的幻象欺骗心灵。逃避这一切的一个方法就是与整个世界分离；也就是说，不是身体的分离，而是灵魂对身体的同情的断绝，并这样生活：没有城市、家庭、财产、社交、所有物、生计、事务、交往、人类学问，使心灵随时接受神圣教义的每一印记。心灵的准备就是摒弃邪恶交往的偏见。如同在写字前先把蜡板弄平滑。\par

现在，独处对此目的极为有用，因为它平息我们的激情，并让原则有机会将其从灵魂中切除。[正如动物在安抚时更容易被控制，贪欲和愤怒、恐惧和忧伤，这些灵魂的死敌，在通过静止平静下来，且没有持续刺激时，更能被理性所控制。] 那么，让我们有一个像我们这样的地方，远离与人交往，以便我们操练的连续性不被外界打断。虔诚的操练以神圣的思想滋养灵魂。还有什么状态比在地上模仿天使的合唱团更为有福？以祈祷开始一天，以赞美诗和歌曲尊崇我们的造物主？当白昼明亮时，我们将祈祷伴随始终，投入我们的劳动，并用赞美诗如盐般调味我们的工作？舒缓的赞美诗使心灵进入愉快平静的状态。因此，安静，如我所说，是我们圣化的第一步；舌头洁净了世界的闲言碎语；眼睛不被美丽的颜色或匀称的形体所刺激；耳朵不因悦耳的歌曲或特别的害处——轻浮之人和说笑者的谈话——而放松心态的基调。这样，心灵从外在的散漫中得救，不通过感官投向世界，便回到自身，从而上升到对天主的默观。[当那美丽照耀它时，它甚至忘记了自己的本性；它不再被食物的念头或衣着的焦虑所拖累；它从世俗的忧虑中休假，并将所有精力投入获取永恒的美好事物，只求如何使自制、勇敢、正义、智慧以及所有其他美德（这些美德分属这些类别，使善人妥善履行一切生活职责）在其中兴盛。]\par

3. 研读默感的圣经是寻求我们本分的主要途径，因为我们在其中既找到关于行为的训导，也找到有福之人的生平，它们被书写下来，犹如虔敬生活的活生生的肖像，供我们效法他们的善行。因此，无论谁在何处感到自己有所欠缺，若致力于此效法，便如同从某间药房找到对症的良药。喜爱贞洁的人思量若瑟的历史，从他学习贞洁的行为，发现他不仅具有克制快乐的自制力，而且在习性上具有德性的心态。忍耐则由约伯所教导：[他不仅在生活环境开始逆转、一瞬间从富贵跌入贫穷、从佳儿之父变为无子时，仍保持同一，灵魂的性情始终未被压碎，甚至对那些前来安慰他、践踏他、加重他痛苦的朋友，也不动怒。] 或者，若有人想学习如何既温良又宽宏大量，对罪刚毅，对人温良，他将发现达味在战事中英勇，在报复仇敌时温良而平静。梅瑟也是如此，对得罪天主的人勇毅地挺身而出，却以温良的灵魂忍受他人对他的诽谤。[因此，一般而言，如同画家在临摹其他画作时，不断注视范本，尽力将它的线条转移到自己的作品中；同样，凡渴望在一切美德上达至成全的人，也当将目光转向圣人们的生平，如同活生生的会移动的雕像，藉效法而拥有他们的美德。\par

4. 再者，阅读之后的祈祷，发现灵魂更为清新，并被对天主的爱更强烈地激动。那在灵魂中铭刻对天主的清晰概念的祈祷是好的；而藉记忆使天主安住在自身内，便是天主的居住。如此，我们成为天主的圣殿，当我们记忆的连续性不被尘世的忧虑切断时；当心神不被突如其来的感觉所困扰时；当敬拜者逃离一切事物，退避到天主那里，将所有引向自我放纵的情感都拉开，并度过时间在导向德性的追求中。\par

5. 这也是一个非常重要需要注意的点——懂得如何交谈：询问时不过于急切；回答时不炫耀；不打断有益的说话者，也不野心勃勃地想要插入自己的话；在说话和聆听上有节制；不以接受为耻，也不吝于提供信息；不把他人的知识当作自己的，如同堕落的女子把别人的孩子当作自己的，而是坦诚地归功于真正的父母。声音的中音最好，既不太低而听不见，也不因音调过高而失礼。人应当先思考自己要说什么，然后再发言；被招呼时要谦恭；社交时要和蔼可亲；不追求以滑稽取悦人，而是以良善的劝诫培育温柔。严厉总当撇弃，即使在责备时也是如此。[你越显出谦逊和谦卑，就越可能被需要你治疗的患者所接纳。然而，有许多场合我们应当善用先知所用的那种斥责方式：他并非亲自对达味犯罪后宣判定罪，而是通过虚构一个人物，使罪人成为自己罪的审判者，这样，在宣判自己之后，他就不能挑剔那揭发他的先见者。\par

6. 从谦卑和顺服的精神，产生忧愁和低垂的眼神，外表不修边幅，头发粗糙，衣服肮脏；以致哀悼者刻意呈现的样子，成为我们自然的状态。袍子应用腰带束在身体上，腰带不要像女人那样高过腰侧，也不要松弛，使袍子像懒散者那样松散。步履不应迟缓，这显示出缺乏活力的性格；另一方面，也不应推搡和傲慢，好像我们的冲动是鲁莽和狂野的。衣服的唯一目的，是在冬夏都足够遮盖。至于颜色，避免明亮；材料方面，避免柔软和精致。追求衣服的鲜艳颜色，就像妇女涂脂抹粉，用不同于自己本色的颜色染脸颊和头发。袍子应足够厚实，不需要其他帮助来保暖。鞋子应廉价但耐用。总之，关于衣着要注意的是必需品。饮食也是如此；因为健康的人面包就够了，水能解渴；可以加一些蔬菜，有助于增强体力以履行其职能。不应带着野蛮贪食的样子进食，而应在所有涉及我们快乐的事上保持节制、安静和自制；并且，始终不让心灵忘记思念天主，而要使我们食物的性质，以及进食身体的构造，成为光荣祂的基础和手段，我们想到各种适合身体需要的食物，都归功于宇宙的伟大管家的供应。饭前应念谢恩经，为感谢天主现在赐予的恩赐，也为祂将来保留的恩赐。饭后也要念谢恩经，感谢已赐的恩赐，并祈求应许的恩赐。每天进食应有一个固定的时间，总是在固定的规律中，这样一天二十四小时中只有这一小时花在身体上。其余时间，苦修者应当用于心智的锻炼。睡眠要轻而容易中断，正如清淡饮食后自然发生的那样；应当有意用关于伟大主题的思想来打断睡眠。被沉重的昏睡压倒，四肢无力，以至于轻易为狂野的幻想打开通道，就是陷入每日的死亡。对虔诚的运动员来说，午夜如同黎明；夜晚的寂静给他们的灵魂留下闲暇；没有有害的声音或景象侵入他们的心；心灵独自与自己和天主相处，通过回忆自己的罪过而修正自己，给予自己帮助避恶的箴言，并恳求天主的帮助以完成其所渴望的。\par

\SourceRecord{Translated by Blomfield Jackson.；Nicene and Post-Nicene Fathers, Second Series；Vol. 8.；Edited by Philip Schaff and Henry Wace.；Buffalo, NY: Christian Literature Publishing Co.,；1895.；<http://www.newadvent.org/fathers/3202002.htm>.}

% structure: p0001 source=h1 target=section reason=page-title

\section{第三书}

\Emph{\Person{圣巴西略·凯撒勒雅}}\par

\Emph{致坎迪迪亚努斯}。\par

1. 当我将你的信拿在手中时，我经历了一段值得讲述的经验。我以面对宣布国家大事的文书应有的敬畏看着它；在拆蜡封时，我感到一种恐惧，甚于有罪的斯巴达人见到拉克尼亚的\OriginalTerm{斯库塔莱密信}{scytale}时的恐惧。\par

然而，当我打开信并读完后，我忍不住笑了，一半是因为高兴发现信中没有任何令人惊恐的事，一半是因为我将你的处境比作德摩斯梯尼的处境。你记得，德摩斯梯尼在为一个小型的合唱舞队和乐师班提供费用时，要求不再被称为德摩斯梯尼，而被称为\Quote{\OriginalTerm{合唱队领队}{choragus}}。你无论是否扮演\Quote{\OriginalTerm{合唱队领队}{choragus}}的角色，始终如一。你确实是比德摩斯梯尼所供应必需品的人数多出无数倍的士兵们的\Quote{\OriginalTerm{合唱队领队}{choragus}}；然而，你在写信给我时并不摆架子，而是保持旧有的风格。你并未放弃对文学的研究，而是如柏拉图所说，在事务的风暴和喧嚣中，你仿佛站在一道坚固的墙下，使你的心不受任何干扰；不仅如此，你尽己所能，甚至不让别人受到干扰。这便是你的生活；对所有有眼可见的人，它是伟大而奇妙的；然而，对任何以你整个生活目的来判断的人来说，这并不奇怪。\par

现在让我说说我自己的事，确实非同寻常，但也在意料之中。\par

2. 一位与我们一起住在安内西的女仆，在我的仆人死后，不声称与他有任何违约行为，不接近我，不提出任何抱怨，不要求我自愿付给她任何款项，不威胁若得不到钱便使用暴力，突然，与几位像她自己一样的狂人袭击了我的家，野蛮地殴打负责看守的妇女，砸开大门，夺走了一些物品，并把剩下的东西许诺给任何想要的人，然后席卷一空。我不愿被视为\OriginalTerm{至极}{ne plus ultra}的无助，成为任何喜欢攻击我的人施暴的合适对象。因此，我现在请求你，请展示你一贯对我事务的善意关怀。只有一个条件可以确保我的安宁——即我确信有你精力充沛的支持。依我看来，若那人被地方长官逮捕并在监狱中关押一小段时间，便已足够惩罚。我不仅因所受的待遇而愤慨，而且希望未来有保障。\par

\SourceRecord{Translated by Blomfield Jackson.；Nicene and Post-Nicene Fathers, Second Series；Vol. 8.；Edited by Philip Schaff and Henry Wace.；Buffalo, NY: Christian Literature Publishing Co.,；1895.；<http://www.newadvent.org/fathers/3202003.htm>.}

% structure: p0001 source=h1 target=section reason=page-title

\section{第四封信}

\Emph{圣巴西略·凯撒勒雅}\par

\Emph{致\Person{奥林匹奥斯}}\par

亲爱的先生，您将我那亲爱的朋友、哲学之乳母——\Person{贫穷}——逐出我们的退隐之所，这是什么意思？倘若她能言语，我想您恐怕要在一桩非法驱逐案中成为被告。她可能会抗辩道：\Quote{我选择与这位\Person{巴西略}生活，他是\Person{芝诺}的仰慕者；\Person{芝诺}在沉船中失去一切时，曾以极大的刚毅喊道：\InnerQuote{干得好，\Person{命运女神}！你正使我回归那旧斗篷。}他也是\Person{克莱安特斯}的伟大仰慕者；\Person{克莱安特斯}靠从井中打水为生，不仅足以糊口，还能支付老师的学费。他还是\Person{第欧根尼}的狂热崇拜者；\Person{第欧根尼}以仅需必要之物而自豪，并在从某个少年那里学会弯腰用手心喝水之后，扔掉了他的碗。}我亲爱的伙伴\Person{贫穷}，因您的礼物而被逐出家园，她或许会用这样的话语责备您。她或许还会加上一句威胁：\Quote{倘若我再在此处捉着您，我将让您见识之前不过是西西里或意大利式的奢侈：我定要用自己的库存精确地回报您。}\par

不过，就此打住。我很高兴您已经开始一个疗程的药物治疗，并祈祷您能从中获益。一个适合无痛苦活动的身体状况，很配这样虔诚的灵魂。\par

\SourceRecord{Translated by Blomfield Jackson.；Nicene and Post-Nicene Fathers, Second Series；Vol. 8.；Edited by Philip Schaff and Henry Wace.；Buffalo, NY: Christian Literature Publishing Co.,；1895.；<http://www.newadvent.org/fathers/3202004.htm>.}

% structure: p0001 source=h1 target=section reason=page-title

\section{第五封信}

\Emph{圣巴西略（凯撒勒雅）}\par

\Emph{致内克塔里乌。}\par

1. 我听闻了您那难以忍受的损失，深感悲伤。三四天过去了，我仍心存疑虑，因为报信者无法向我清楚叙述这悲伤事件的细节。我外界的传言难以置信，因我祈祷它们不实，直到从主教处收到一封信，完全证实了这不幸的消息。我无需向您描述我是如何叹息哭泣。谁能如此铁石心肠、真正无情，竟然对已发生之事无动于衷，或仅以适度的悲伤来应对？他已离去：尊贵家族的后嗣，家族的支柱，父亲的希望，虔诚父母的子息，受无数祈祷养育，正值壮年，却从父亲手中被夺走。这些事足以使金刚石般的心破碎而感受痛苦。因此，我自然对这灾祸深感触动；我自始至终与您亲密相连，将您的喜乐与哀伤视为己有。但昨日，您似乎还少有烦恼，生命之流顺利前行。转瞬之间，因恶魔的恶意，家中所有的幸福、生命所有的光明，都遭摧毁，我们的生命成了哀伤的故事。若我们愿哀悼哭泣所发生之事，一生也不够；即使全人类与我们同哀，也无法使他们的哀悼与我们的损失相称。是的，即使所有溪流都化为泪水，也不足以充分哀哭我们的悲痛。\par

2. 但我们岂不是要发挥天主藏在我们心中的恩赐——即那清醒的理性？它在幸福的日子里常为我们的灵魂划定界限，当患难来临时，则使我们想起我们不过是人，并提示我们——这确实是我们所看见、所听见的——生命充满类似的灾祸，人类受苦的榜样不在少数。尤其重要的是，它将告诉我们，天主的命令是：我们这些信靠基督的人，不应为长眠的人悲伤，因为我们盼望复活；并且，为了作为巨大忍耐的奖赏，生命赛程的主宰为我们存留了伟大的荣冠。只要我们允许自己更智慧的思想以这种和谐的旋律向我们发言，我们或许就能找到一些减轻苦难的方法。那么，我恳求您，振作起来；打击是沉重的，但要站稳；不要被悲伤的重量压垮；不要灰心。请完全确信：虽然天主所命定之事的原因超乎我们，但那位智慧且爱我们的天主为我们安排的一切，无论多么难以忍受，都应当接受。祂自己知道如何为每个人安排最佳之事，以及为何祂为我们设定的生命期限不等。存在某种人无法理解的理由，为何有些人更早远离我们，而有些人则留得更久，以承受这痛苦人生的重担。因此，我们应当永远钦崇祂的仁慈，而不抱怨，并记得那位伟大的斗士约伯的著名话语——当他看到十个孩子在同一桌前转瞬之间被压死时，他说：\BibleQuote{主赐予，主也收回。}\BibleRef{约 1:21} 主看如何好，就如何成就。让我们采用这些奇妙的话语。在公义审判者手中，行同样善事的人将得到同样的赏报。我们并未失去这孩子；我们只是把他归还给了借主。他的生命并未消灭；而是改变得更好。我们所爱的人并非隐藏在地下；而是被接入天堂。让我们稍等片刻，我们将再次与他在一起。我们分离的时间不长，因为在今世，我们都如旅人，奔向同一个庇护所。一人已到达安息处，另一人到来，又一人匆忙赶路，但同一位终点等待着他们。他在路上超越了我們，但我们都将走同一条路，同一家旅店等待我们众人。惟愿天主赐予我们，因良善而能肖似他的纯洁，好使我们因生活的纯真而获得那赐予在基督内为孩童之人的安息。\par

\SourceRecord{Translated by Blomfield Jackson.；Nicene and Post-Nicene Fathers, Second Series；Vol. 8.；Edited by Philip Schaff and Henry Wace.；Buffalo, NY: Christian Literature Publishing Co.,；1895.；<http://www.newadvent.org/fathers/3202005.htm>.}

% structure: p0001 source=h1 target=section reason=page-title

\section{\WorkTitle{第六封书信}}

\Emph{\Person{圣巴西略·凯撒勒雅}}\par

\Emph{致纳克塔留斯之妻}。\par

1. 我犹豫是否要向您的尊贵致辞，因为我想，正如眼睛发炎时，即使最温和的药物也会引起疼痛，同样，当灵魂被深重的悲伤压伤时，在痛苦时刻所说的话语，即使它们带来许多安慰，也似乎有些不合时宜。但我想，我是在对一位基督徒妇女说话，她早已领受了虔敬的教训，对于人生变故并不陌生，于是我认为不应忽略我所承担的责任。我知道母亲的心是怎样，当我记起您待众人是何等良善温和时，我能估算出此刻您可能的痛苦程度。您失去了一位儿子，他在世时，所有母亲都称他有福，并祈祷自己的孩子能像他一样；他去世时，她们哀悼，仿佛各人都将自己的孩子埋葬了。他的死是两省的打击，既打击了我的省，也打击了奇里基雅。随着他的倒下，一个伟大而显赫的家族也倒下了，如同支柱被抽走而轰然坠地。唉！邪恶魔鬼的接触竟能造成如此巨大的祸害！大地啊，你被迫承受了何等灾祸！若太阳有感觉，人们会想它或许也对这悲惨景象感到战栗。谁能说出那无助心灵所想的一切呢？\par

2. 但我们的生命并非没有\OriginalTerm{天主眷顾}{Providence}。我们在福音中这样学到，因为若不是我们圣父的意愿，一只麻雀也不会掉在地上。\BibleRef{玛 10:29}所发生的一切，都是因着我们造物主的旨意而发生的。谁能抗拒天主的旨意呢？让我们接受所遭遇的事吧；因为如果我们对此不满，我们既不能改变过去，又会使自己毁灭。我们不要指责天主公义的审判。我们都太过无知，无法攻击祂不可言喻的判决。现在主正在考验您对祂的爱。如今您有机会，借着忍耐，分享殉道者的命运。玛加伯的母亲看到七个儿子的死亡，没有叹息，甚至没有流下一滴不相称的眼泪。她感谢天主，因为他们借着火、刀剑和残酷的打击，从肉体的束缚中得了解放。她赢得了天主的称赞和世人的美名。损失是巨大的，我自己也能这样说；但主为忍耐者预备的赏报也是巨大的。当您最初成为母亲，看到您的孩子，感谢天主时，您始终知道，您自己是一个凡人，所生下的也是一个凡人。凡人之死有什么可惊奇的呢？但我们因他未到天年而死感到悲伤。我们确定这不是他的时候吗？我们不知道如何为自己灵魂的利益挑选和选择，也不知道如何限定人的生命界限。环顾您所居住的整个世界；记住您所见的一切都是可朽的，都受败坏。仰望天空；它也要被解散；看太阳，连太阳也不会永远存留。所有星辰、所有陆地和海中的活物、地上一切美好的事物，甚至大地本身，都受朽坏；再过片刻，一切都要归于无有。让这些思考在您的困苦中给您一些安慰。不要单凭您的损失来衡量它；如果您这样做，它会显得难以忍受；但如果您把人类所有的事务都考虑在内，您会发现可以从其中得到一些安慰。最重要的是，我要强烈劝诫您一件事：爱惜您的丈夫。成为他人的安慰。不要让您因忧愁而耗损自己，加重他的痛苦。我知道单凭言语不能带来安慰。现在所需要的是祈祷；我确实祈求主自己用祂不可言喻的力量触动您的心，通过良善的思想使光照耀您的灵魂，使您在自己内有安慰的源泉。\par

\SourceRecord{Translated by Blomfield Jackson.；Nicene and Post-Nicene Fathers, Second Series；Vol. 8.；Edited by Philip Schaff and Henry Wace.；Buffalo, NY: Christian Literature Publishing Co.,；1895.；<http://www.newadvent.org/fathers/3202006.htm>.}

% structure: p0001 source=h1 target=section reason=page-title

\section{第七封信}

\Emph{\Person{凯撒勒雅的圣巴西略}}\par

\Emph{致我的友人额我略}。\par

当我写信给你时，我完全清楚，没有任何神学术语足以表达说话者的思想或提问者的需求，因为语言天生就过于软弱，无法服务于思想的对象。如果我们的思想是软弱的，而我们的舌头比思想更软弱，那么我所说的话又能期待什么呢？无非是被人指责表达贫乏而已。尽管如此，我不能对你的问题置之不理。如果我不愿意向那些爱主的人提供一个关乎天主的答案，那似乎是一种背叛。然而，无论所说的话看起来令人满意，还是需要进一步更仔细的补充，都需要一个恰当的时机来进行修正。目前，我恳求你，如同我以前恳求你一样，全然投身于真理的辩护，并善用天主赐予你的理智能力来确立善的事物。对此感到满足，不要向我要求更多。我实际上远不如人们所认为的那样有能力，而且更可能因我的软弱而损害我的话，而不是因我的辩护而增强真理。\par

\SourceRecord{Translated by Blomfield Jackson.；Nicene and Post-Nicene Fathers, Second Series；Vol. 8.；Edited by Philip Schaff and Henry Wace.；Buffalo, NY: Christian Literature Publishing Co.,；1895.；<http://www.newadvent.org/fathers/3202007.htm>.}

% structure: p0001 source=h1 target=section reason=page-title

\section{第八封书信}

\Person{凯撒勒亚的圣巴西略}\par

\Emph{致凯撒勒亚人}。\Emph{为他退隐的辩护，及论信德}。\par

1. 我常常惊奇你们对我有如此的感情，并且像我这样一个微不足道、似乎并不值得爱的小人物，如何能赢得你们如此的忠诚。你们提醒我友谊和祖国的要求，并迫切地催促我回到你们那里，好像我是从父亲的内心和家中逃出来的人。我承认我是个逃亡者，不愿否认；既然你们已经想念我，我就告诉你们原因。我像一个被突然响声震惊的人那样惊愕。我没有压抑我的思绪，而是在逃亡中沉思它们，如今我已离开你们相当长的时间。然后我开始渴望神圣教义以及与之相关的哲学。我说，我如何才能克服居住在我们中间的邪恶？谁是我的\Person{拉班}，使我摆脱\Person{厄撒乌}，并引导我走向至高的哲学？赖天主的助佑，我已尽我所能达到我的目标；我找到了一个被选的器皿，一口深井；我是指\Person{额我略}，基督的口。因此，我恳求你们，给我一点时间。我不是在拥抱城市生活。我很清楚恶者如何借此为人类设计欺骗，但我确实认为与圣人们的交往最为有益。因为在对神圣教义的更为持续的思想交流中，我正养成持久的默观习惯。这就是我目前的情况。\par

2. 诸位敬虔挚爱的朋友，我恳求你们，要提防培肋舍特人的牧人；不要让他们不知不觉堵塞你们的意志，不要让他们玷污你们对信德认知的纯洁。他们的目的始终是，不是用圣经教导纯朴的灵魂，而是用外教哲学破坏真理的和谐。难道他不是公开的培肋舍特人，将\Quote{\OriginalTerm{无生}{unbegotten}}和\Quote{\OriginalTerm{受生}{begotten}}这两个词引入我们的信仰，并断言曾经有一个时候，永生者不存在；那位在本性和永恒上是父的，变成了父；圣神不是永恒的吗？他迷惑我们先祖的羊群，使她们不喝\BibleQuote{那涌到永生的水泉的井} \BibleRef{若 4:14}，反而把先知的话招到自己身上：\BibleQuote{他们离弃了我这活水的泉源，却为自己挖掘了不能蓄水的漏水水池} \BibleRef{耶 2:13}；而他们本该承认父是天主，子是天主，圣神是天主，正如他们从神圣的话语以及从那些以最高意义理解话语的人所学的。对于那些指责我们是三神论者的人，应当回答：我们承认一位天主，不是在数目上，而是在本性上。因为凡在数目上被称为一的东西，并非绝对是一，也不是本性上简单的；但天主普遍被承认为简单而非复合的。因此天主不是在数目上是一。我这样说。我们说世界在数目上是一，但在本性上不是一，也不是简单的；因为我们将世界分解为它的构成元素：火、水、空气和土。同样，人被称为在数目上的一。我们常说一个人，但人由身体和灵魂组成，不是简单的。同样，我们说一个天使在数目上的一，但在本性上不是一，也不是简单的，因为我们设想天使的\OriginalTerm{位格}{hypostasis}是本质加上圣化。因此，如果凡在数目上是一的东西在本性上不是一，而凡在本性上是一且简单的东西在数目上不是一；如果我们称天主在本性上是一，那么我们怎么能被指控数目呢？我们完全将数目排除在那有福的和属灵的本性之外。数目与量有关，量与肉体本性相连，因为数目属于肉体本性。我们相信我们的主是肉体的创造者。因此，每一个数目都指示那些接受了物质和有限本性的事物。反之，\OriginalTerm{太一}{Monad}和\OriginalTerm{统一}{Unity}则指示简单且不可理解的本性。因此，凡承认天主子或圣神是数目或是受造物的人，就不自觉地引入了物质和有限的本性。我所说的有限，不仅指在空间上受限制，而且指一种本性，能由那将要把它从无中引出到有的那一位在预知中理解，并且能被科学所理解。因此，每一种本性有限、圣洁是获得的圣物，都不免受恶的影响。但是圣子和圣神是圣化的根源，每个有理性的受造物都按其德性藉此被圣化。\par

3. 我们依照真道，论及圣子时，既不说祂相似圣父，也不说祂不相似圣父。这两个词同样不可能，因为相似和不相似是关于性质的谓语，而天主性不受性质限制。相反，我们承认本性同一，接受同体，并摒弃圣父有组合成分，祂在本体上是天主，生了圣子，圣子在本体上也是天主。由此证明同体。因为在本质或本体上的天主，与在本质或本体上的天主是同本质或同体的。但当人也被称为\BibleQuote{神}，如经上说：\BibleQuote{我曾说你们是神}，以及被称为\BibleQuote{邪魔}，如经上说：\BibleQuote{外邦的神都是邪魔}；前者是出于恩宠的名号，后者则是虚假的。唯独天主在本体上、在本质上是天主。当我说\BibleQuote{唯独}时，我是指天主神圣而未被造的本质和本体。因为\BibleQuote{唯独}一词可用于任何个体，也可泛指人性。用于个体，如保禄，他唯独被提到第三层天，并\BibleQuote{听见了不可言传的话，那是人不能说的}（\BibleRef{格后 12:4}）；用于人性，如达味说：\BibleQuote{至于世人，他的岁月如同青草}；这不是指某个具体的人，而是泛指人性，因为每个人都是短暂而会死的。同样，我们理解这些话是论及本性：\BibleQuote{唯独祂永不死}（\BibleRef{弟前 6:16}），\BibleQuote{独一全智的天主}（\BibleRef{罗 16:27}），\BibleQuote{除了天主一个外，没有谁是良善的}（\BibleRef{路 18:19}），这里的\BibleQuote{一个}意思与\BibleQuote{唯独}相同。同样，\BibleQuote{惟有祂独自铺张诸天}（\BibleRef{约 9:8}），以及\BibleQuote{你要朝拜上主你的天主，唯独事奉祂}，\BibleQuote{除我以外没有别的神}。在圣经中，\BibleQuote{一个}和\BibleQuote{唯独}论及天主时，不是为了区别于圣子和圣神，而是为了排除那些虚假被称为神的。例如：\BibleQuote{上主独自引导他们，没有外邦神与他们同在}，\BibleQuote{那时以色列子民除掉巴耳和阿舍辣，唯独事奉上主}（\BibleRef{撒上 7:4}）。同样，圣保禄说：\BibleQuote{因为虽然有称为神的，或在天上，或在地上，就如那许多\BibleQuote{神}，许多\BibleQuote{主}，然而我们只有一个天主，就是圣父，万物都出于祂；也只有一个主耶稣基督，万物是藉着祂而有的}（\BibleRef{格前 8:5}）。这里我们要问，为什么当祂说\BibleQuote{一个天主}时还不满足？因为我们说过，\BibleQuote{一个}和\BibleQuote{唯独}应用于天主时，是指明本性。为什么他又加上\BibleQuote{圣父}一词，并提及基督？保禄——蒙选的器皿——我想，并不认为仅仅宣讲圣子是神、圣神是神就足够了，他已经用\BibleQuote{一个天主}这词表达了这一点，但并没有通过进一步添加\BibleQuote{圣父}来表达万物所出的祂；也没有通过提及\BibleQuote{主}来指明万物藉以存在的圣言；更进一步，加上\BibleQuote{耶稣基督}这个词，宣布了降生，陈述了苦难，宣告了复活。因为\BibleQuote{耶稣基督}这个词向我们提示了所有这些观念。为此，我们的主在受难之前不愿被称为\BibleQuote{耶稣基督}，并吩咐祂的门徒\BibleQuote{不要告诉任何人祂是耶稣、基督}（\BibleRef{玛 16:19}）。因为祂的旨意是在完成救恩计划、从死者中复活、升天之后，才交付给他们宣讲祂是耶稣、基督。这就是\BibleQuote{使他们认识你，唯一的真天主，和你所派遣的耶稣基督}（\BibleRef{若 17:3}）以及\BibleQuote{你们信天主，也当信我}（\BibleRef{若 14:1}）这些话的含义。圣神处处确保我们对祂的理解，以免我们在一方面前进时却在另一方面跌倒，注重神学而轻视救恩计划，从而因疏忽而陷入不敬虔。\par

4. 现在让我们考查，并尽己所能解释圣经话语的含义，对手们抓住这些话，曲解为自己的意思，用来攻击我们，以毁灭独生子的荣耀。首先，取\BibleQuote{我因父而生活}这句话，因为这是那些不敬虔地使用它的人射向天上的箭矢之一。我认为这些话不是指永恒的生命；因为任何因他物而生活的，都不能自存，正如被另一个温暖的事物不能是温暖本身；但那位是我们的基督和天主说：\BibleQuote{我是生命}（\BibleRef{若 11:25}）。我理解\BibleQuote{因父而生活}指的是这肉身的、现世的生活。祂自愿来度人类的生活。祂没有说\BibleQuote{我曾经因父而生活}，而是说\BibleQuote{我因父而生活}，清楚指明现在的时间，而基督，在祂内有天主的话语，能够称祂所度过的生活为生命；而祂的这个意思我们将从下文得知。祂说：\BibleQuote{吃我的人}，祂说，\BibleQuote{也要因我而生活}；因为我们吃祂的肉，喝祂的血，藉着祂的降生成人和祂可见的生活，成为祂的圣言和祂的智慧的分享者。因为祂把祂在我们当中的全部奥秘旅程称为血肉，并提出包含实践科学、物理学和神学的教导，藉此我们的灵魂得到滋养，同时被训练去默观实际的实在。这也许是祂所说的话的意图。\par

5. 再者，\BibleQuote{我的父比我大。}\BibleRef{若 14:28}忘恩负义之徒，恶者的孽种，也利用这节经文。我相信，即使从此节经文也可看出圣子与圣父的同体。因为我知道，在具有同一性质的事物之间才能恰当地进行比较。我们说天使大于天使，人比人更义，鸟比鸟更快。那么，如果比较是在同一种类的事物之间进行的，并且圣父被说成比圣子大，那么圣子就与圣父是同一实体。但这句话背后还有另一层意义。那\BibleQuote{是圣言，且成了血肉}\BibleRef{若 1:14}的祂，在荣耀中比天使卑微，在形象上不如人，如此承认祂的父比祂大，这有什么可奇的呢？因为\BibleQuote{祢使他稍微逊于天使}，又说\BibleQuote{那暂时成为比天使小的}\BibleRef{希 2:9}，以及\BibleQuote{我们看见祂，祂没有俊美，也没有容貌，祂的容貌比众人更残缺。}这一切祂都是为着祂对自己作品的丰富慈爱而忍受的，为要拯救失落的羊，并在拯救后领它回家，将那个从\Place{耶路撒冷}下\Place{耶里哥}而落到强盗手中的人，安然无恙地带回他自己的地方。异端者岂会因祂的愚昧无知而将马槽掷于祂的脸上，而祂正是藉着理性的圣言喂养了他？岂会因木匠的儿子没有床铺就抱怨祂贫穷？这就是为什么圣子小于圣父；祂为你们而死，为使你们脱离死亡，并使你们分享天上的生命。这正如同有人责怪医生俯就疾病，呼吸病气，以便医治病人。\par

6. 祂不知道审判的时辰和日子，这是为了你们。然而，真智慧无所不知，因为\BibleQuote{万物是藉着祂造成的}\BibleRef{若 1:3}；就连人也不会不知道自己所造的东西。但这是祂因你们的软弱而作的安排，免得罪人因所定时期过短而陷入绝望，不给他们留下悔改的机会；另一方面，也免得那些长期与敌人势力争战的人因时间拖延而擅离岗位。祂藉着自己\OriginalTerm{假定的无知}{assumed ignorance}为这两类人作了安排：为前者缩短时间，使他们进行光荣的战斗；为后者提供悔改的机会，因了他们的罪。在福音书中，祂将自己列于无知者之中，如我所说，是为了大多数人的软弱。在\WorkTitle{宗徒大事录}中，祂似乎是对完美的人说话，说道：\BibleQuote{父以自己的权柄所定的时候和日期，不是你们应当知道的。}\BibleRef{宗 1:7}在此，祂含蓄地把自己排除在外。以上只是初步攻击的粗略陈述。现在让我们从更高的观点探究这段经文的含义。让我叩响知识之门，也许我能唤醒那家的主人，祂把神粮赐给求祂的人，因为我们急于款待的是朋友和兄弟。\par

7. 我们救主的圣洁门徒，在超越了人类思想的极限，并藉着圣言得到净化之后，探询那结局，渴望认识我们的主宣称连祂的天使和祂自己都不知道的终极福乐。祂称对天主计划的一切精确领悟为一天，而将对惟一无二与合一的默观——祂将此知归给圣父独自享有——称为一个时辰。因此，我理解：天主被说成知道自己存在的事，而不知道不存在的事；天主，祂按其本性就是公义和智慧，被说成知道公义和智慧，但对不义和邪恶一无所知；因为创造我们的天主不是不义和邪恶。那么，如果天主被说成知道自己存在的事，而不知道不存在的事；如果我们的主，按降生成人的目的和较粗浅的道理，并非最终渴望的对象；那么我们的救主就不知道那结局和终极福乐。但祂说天使也不知道\BibleRef{玛 13:32}——也就是说，连在他们内的默观和他们的职务方式也不是最终渴望的对象。因为他们的知识，与面对面的知识相比，也是粗浅的。祂说，只有圣父知道，因为祂自己就是结局和终极福乐；因为当我们不再在镜子里、不直接地认识天主，而是作为惟一无二的一位接近祂时，我们就将认识那终极结局。因为一切属物质的知识都被称为基督的王国；而属非物质的知识，可以说是实际神性的知识，则是天主圣父的知识。但我们的主，按圣言的目的，祂自己也就是结局和终极福乐；因为祂在\WorkTitle{福音}中说了什么？\Quote{在末日我要使他复活。}\BibleRef{若 6:40}祂称从物质知识到非物质默观的过渡为复活，将那种知识之后别无其他的知识称为末日；因为当我们的理智默观圣言的惟一无二与合一时，它就被提升、被唤醒到福乐的高峰。但既然我们的理智变得粗浊、束缚于尘世，并与泥土混杂，不能专注地凝视纯粹的默观，而只能通过与其身体相关的装饰被引导。它思量造物主的工程，并暂时凭其效果判断它们，以便逐渐成长，终有一天变得足够强壮，以致能接近那实际无遮蔽的天主性。我想，这就是\Quote{我的父比我大}\BibleRef{若 14:28}这句话的含义，也是\Quote{不是我可以给的，而是为我父所预备的}这句话的意思。基督\Quote{将王国交给天主圣父}\BibleRef{格前 15:24}这句话也是这个意思；因为按较粗浅的道理——如我所说，这是相对于我们而非圣子本身而言——祂不是结局，而是初果。正是根据这一观点，当祂的门徒在\WorkTitle{宗徒大事录}中再次问祂：\Quote{你什么时候复兴以色列国？}时，祂回答说：\Quote{父以自己的权柄所定的时候和日期，不是你们可以知道的。}\BibleRef{宗 1:6}也就是说，这样的国度的认识不是为那束缚在血肉之中的人。圣父将这默观放在自己的权柄中，\Quote{权柄}意指那些有权能的人，\Quote{自己的}意指那些不为下面事物的无知所束缚的人。我恳求你，不要想着感官的时间或日期，而是以心灵感知的太阳所认识的知识的某种区分。因为我们主的祈祷必须实现。是耶稣祈祷说：\Quote{父啊！愿他们在我们内合而为一，如同我与你是一体。}因为当天主——祂是惟一无二——在每人内时，祂使众人合一；数目消失在合一的内在居住中。\par

这是我第二次尝试阐释这经文。若有人能给出更好的解释，并能本着真宗教修正我所说的，就让他说，让他修正，主必因我而赏报他。我心中没有嫉妒。我研究这些段落，不是为了争辩和虚荣，而是为了帮助我的弟兄们，免得那些心硬如石、未受割礼的人——他们的武器是愚妄的智慧——似乎欺骗了那盛有天主宝藏的瓦器。\par

8. 再者，正如通过智者撒罗满在\WorkTitle{箴言}中所说的：\Quote{祂被造；}且被称为\Quote{道路的起始}——这道路是福音之路，引领我们走向天国。祂在本体和本体上不是受造物，而是按\OriginalTerm{救恩计划}{œconomy}被造成了一条\Quote{道路}。\Quote{被造成}与\Quote{被创造}意义相同。正如祂被造成道路，祂也被造成门、牧人、天使、羔羊，以及大司祭和宗徒\BibleRef{希 3:1}——这些称谓是在其他意义上使用的。再者，论到被称为不受臣服的天主，以及祂为我们成了罪，异端者又能说什么呢？因为经上记载：\Quote{那时，万物都屈服于祂，圣子自己也要屈服于那使万物屈服于祂的。}先生，你难道不怕天主被称为不受臣服的？因为祂把你的屈服当作自己的；并且因你抗拒良善，祂称自己不受臣服。同样，祂也曾说自己被迫害——\Quote{扫禄，扫禄，你为什么迫害我？}\BibleRef{宗 9:4}那时扫禄正赶往大马士革，想要囚禁门徒。祂又称自己赤身露体，当祂任何一位弟兄赤身时。\Quote{我赤身露体，你们给我穿了衣裳。}\BibleRef{玛 25:36}同样，当别人在监里时，祂说自己被囚禁，因为祂亲自承担了我们的罪愆，背负了我们的疾病。我们的软弱之一就是不受臣服，祂背负了这一点。因此，凡我们所遭遇的一切伤害，祂都当作自己的，在祂与我们共融中承担了我们的痛苦。\par

9. 但那些反抗天主、使其听者败坏的家伙，又抓住了另一段经文：即\Quote{圣子凭自己不能做什么}\BibleRef{若 5:19}。在我看来，这句话也显然宣称圣子与圣父本性相同。因为如果每个有理性的受造物凭自己都能有所作为，且每个人趋向恶或善的倾向都在自己的权下，而圣子却\Quote{圣子凭自己不能做什么}，那么圣子就不是受造物。如果祂不是受造物，那么祂就与圣父同一本体和实体。再者，没有受造物能做它喜欢的事。但圣子在天地间实现了祂的意愿。因此圣子不是受造物。再者，一切受造物要么由对立面构成，要么能接纳对立面。但圣子是真正的正义，且是非物质的。因此圣子不是受造物，并且如果祂不是受造物，祂就与圣父同一本体和实体。\par

10. 对于眼前的这些经文，凭我的能力做了足够的考察。现在让我们把讨论转向那些攻击圣神的人，并摧毁他们理智中一切高举起来反对天主知识的傲慢之事\BibleRef{格后 11:5}。你说圣神是一个受造物。而一切受造物都是造物主的仆人，因为\Quote{万物都是祢的仆人}。如果祂是仆人，那么祂的圣洁是获得的；而一切圣洁是获得的之事物，都能接纳邪恶；但圣神本质上是圣洁的，被称为\Quote{圣洁之源}\BibleRef{罗 1:4}。因此圣神不是受造物。如果祂不是受造物，祂就与圣父同一本体和实体。告诉我，你怎能将仆人之名加给那藉着圣洗使你脱离奴役的那一位？经上说：\Quote{生命之圣神的神律使我脱离了罪的法律}\BibleRef{罗 8:2}。但你永远不敢称祂的本性为可变的，只要你顾及与祂相对立的仇敌势力的本性——那势力如同闪电从天坠下，从真生命中陨落，因为它的圣洁是获得的，它的邪恶谋划导致了它的改变。当它从合一中坠落，丢弃了天使的尊严时，便按它的品性被称为\Quote{魔鬼}，它从前福乐的状态熄灭，而这敌对的势力被燃起。\par

再者，若有人称圣神为受造物，他就是在描述祂的本性有限。那么以下两段经文怎能成立？\Quote{上主的神充满了世界}\BibleRef{智 1:7}；以及\Quote{我往何处去躲避祢的神？}但看样子，他并不承认祂本性是单纯的，因为他描述祂在数目上是一。正如我已经说过的，凡数目上为一的事物都不是单纯的。如果圣神不是单纯的，祂就由本体和圣化构成，因此是合成的。但谁会疯狂到形容圣神是合成的、不单纯的，且与圣父和圣子\OriginalTerm{同一实体}{consubstantial}？\par

11. 如果我们应当进一步推进论证，将审视转向更高深的话题，让我们特别从以下观点默观圣神的属神本性。在圣经中，我们读到三次创造。第一次是从无到有的演化。第二次是从坏到好的改变。第三次是死者的复活。在这三次中，你将发现圣神与圣父和圣子合作。诸天的存在产生，达味说什么？\Quote{因上主的话诸天造成，因祂口中的气万象成就。}再者，人藉着圣洗被创造，因为\Quote{谁若在基督内，他就是一个新受造物}\BibleRef{格后 5:17}。为什么救主对门徒们说：\Quote{你们要去使万民成为门徒，因父及子及圣神之名给他们授洗}？在这里你也看见圣神与圣父和圣子同在。对于死者的复活，当我们将朽坏并归于尘土时，你又会说什么？我们是尘土，也将归于尘土。祂将派遣圣神，创造我们，并更新地面。因为圣保禄所谓复活，达味称之为更新。让我们再听听那被提到第三层天上的人。他说什么？\Quote{你们是圣神的圣殿，祂住在你们内}\BibleRef{格前 6:19}。如今，每一座圣殿都是天主的圣殿；如果我们是一座圣神的圣殿，那么圣神就是天主。它也被称为撒罗满的圣殿，但这是就他是其建造者而言。如果我们在这一意义上是一座圣神的圣殿，那么圣神就是天主，因为\Quote{建造万有的就是天主}\BibleRef{希 3:4}。如果我们是一座受敬拜并住在我们内者的圣殿，就让我们承认祂是天主，因为\Quote{你要朝拜主你的天主，唯独事奉祂}。假若有人反对\Quote{神}这个称呼，就让他们学习这个词的含义。神之所以被称为\OriginalTerm{神}{\GreekText{Θεὸς}}，要么是因为祂放置（\OriginalTerm{放置}{\GreekText{τεθεικέναι}}）万物，要么是因为祂观看（\OriginalTerm{观看}{\GreekText{Θεᾶσθαι}}）万物。如果祂被称为\OriginalTerm{神}{\GreekText{Θεὸς}}是因为祂\Quote{放置}或\Quote{观看}万物，并且圣神知道天主的一切，正如在我们内的圣神知道我们的一切，那么圣神就是天主。再者，如果圣神的剑是天主的话语\BibleRef{弗 6:17}，那么圣神就是天主，因为剑属于祂，其话语也归于祂。祂被称为圣父的右手吗？因为\Quote{上主的右手大显威能}，并且\Quote{上主，祢的右手击碎了敌人}\BibleRef{出 15:6}。但圣神是天主的手指，正如经上说：\Quote{我若靠天主的指头驱魔}\BibleRef{路 11:20}，在另一部福音中的版本是\Quote{我若靠天主的圣神驱魔}\BibleRef{玛 12:28}。所以圣神与圣父和圣子\OriginalTerm{同一实体}{consubstantial}。\par

12. 关于可钦崇的、至圣的天主圣三，以上所述已足够目前之用。如今不可能再进一步探讨此事。你们要从像我这样卑微的人身上取些种子，自己培育成熟的穗子，因为你们知道，在这种情况下我们期待利息。但我信赖天主，因你们纯洁的生活，必结出三十倍、六十倍、一百倍的果实。因为经上记载：\BibleQuote{心里洁净的人是有福的，因为他们要看见天主。}\BibleRef{玛 5:8} 并且，我的弟兄们，不要对天国抱有别的概念，只把它视为对实在的静观本身。神圣的圣经称之为真福。因为\Quote{天国就在你们内。}\par

内在的人不外乎静观。因此，天国必须是静观。如今我们仿佛在镜中看见它们的影子；将来，脱离这尘世的身体，穿上不朽坏的、不死的外衣，我们将看见它们的原型——也就是说，如果我们正确驾驶了自己生命的航程，并持守了正确的信仰，否则无人能见主。因为经上记载：\BibleQuote{智慧不进入恶意的灵魂，也不住在一个屈服于罪的身体内。}\BibleRef{智 1:4} 但愿没有人提出反对，说我无视眼前的事物，却向他们空谈无体无形之存有。在我看来，感官被允许自由活动于其适当质料的同时，却唯独将心智排除于其特有作用之外，这完全是荒谬的。正如感官触及可感之物，心智也领悟可领悟之物。还应当说，天主我们的创造主并没有将自然官能归入可教授的事物之列。没有人教视觉辨识颜色与形状，也没有人教听觉辨识声音与言语，嗅觉辨识香臭，味觉辨识味道与滋味，触觉辨识软硬、冷热。同样，没有人教导心智去触及可领悟之物；正如感官若有某种疾病或损伤，只需适当治疗，便轻易履行其功能；同样，心智被囚于肉体内，充满由此而来的思虑，就需要信德和正直的品行，这使\Quote{它的脚快如母鹿的蹄，立于高处。} 同样的劝告由智慧的撒罗满给予我们，他一方面以勤劳的蚂蚁为例，推荐其积极的生活；另一方面以智慧蜜蜂筑巢的工作为例，暗示一种自然的静观，其中也包含至圣天主圣三的道理——只要从受造之物的美丽中相应默观造物主即可。\par

但我要感谢圣父、圣子及圣神，就此结束这封信，因为正如谚语所说：\OriginalTerm{凡事皆适中为佳}{\GreekText{πᾶν} \GreekText{μέτρον} \GreekText{ἄριστον}}。\par

\SourceRecord{Translated by Blomfield Jackson.；Nicene and Post-Nicene Fathers, Second Series；Vol. 8.；Edited by Philip Schaff and Henry Wace.；Buffalo, NY: Christian Literature Publishing Co.,；1895.；<http://www.newadvent.org/fathers/3202008.htm>.}

% structure: p0001 source=h1 target=section reason=page-title

\section{第九封信}

\Emph{\Person{凯撒勒雅的圣巴西略}}\par

\Emph{致哲学家\Person{玛克西穆}}\par

1. 言语实乃心灵之像：我从你的书信中得以认识你，正如谚语所说，我们可以\Quote{从爪子认出狮子。}\par

我很欣喜地发现你强烈的倾向指向那首要且至善之物——即爱天主与爱邻人。从你待我的仁慈中，我找到了后者的证据；从你渴求知识的热情中，我找到了前者的证据。基督的每一位门徒都深知，一切尽包含在这二者之中。\par

2. 你向我要狄约尼削的著作；它们确实到了我手，数量很多；但我身边没有这些书，所以未能寄出。然而，我的看法如下。我并非欣赏他所写的一切；有些事我完全不赞同。因为或许，那如今我们听到如此之多的不敬之论——我指的是\OriginalTerm{阿诺米派}{Anomœan}——据我所知，正是他首先为人们播下了种子。我不将他的这种行为归因于任何心智的败坏，而只归因于他抵抗\Person{撒伯里乌}的强烈意愿。我常把他比作一个试图矫直一棵弯曲树苗的樵夫，他用力过猛地向相反方向拉，以至超出了中点，反而将树苗拉向了另一侧。这正是我们在\Person{狄约尼削}身上所见的。当他激烈反对那个利比亚人的不敬时，不知不觉被自己的热忱卷入相反的错误。他本可以指出圣父与圣子并非本质相同，以此对抗那亵渎者，这原本完全足够。但是，为了赢得一个明确而充分的胜利，他不满足于确立位格的不同，还坚持要主张本质的不同、能力的减弱和光荣的可变性。于是他以一种祸害换取另一种，偏离了正统教义的道路。在他的著作中，他表现出混杂的不一致：有时他对\OriginalTerm{同质（同本体）}{homoousion}不忠，因为他的对手滥用它来摧毁位格；而在另一处，他在给同名的\Person{狄约尼削}的\WorkTitle{护教书}中承认了它。此外，他还说了关于圣神非常不当的话，将祂与天主性、与崇拜的对象分开，并将祂置于与被造和附属本性同等的较低地位。此人便是如此。\par

3. 若我不得不说出自己的看法，那就是如此。\Quote{本质相似}这个短语，如果加上\Quote{毫无差异}来理解，我接受它传达与\OriginalTerm{同质（同本体）}{homoousion}相同的含义，符合同质的正确意义。\Place{尼西亚}的教父们本着这种思想，称独生子为\Quote{光之光}、\Quote{真天主之真天主}等等，然后一致加上了同质。没有人可能持有光相对于光、真理相对于真理有可变性的观念，也不可能认为独生子的本质相对于圣父的本质有可变性。因此，如果这个短语在这个意义上被接受，我对此没有异议。但如果有人像在\Place{君士坦丁堡}所做的那样，从\Quote{相似}一词中砍掉\Quote{毫无差异}这个限定，那么我对此短语持怀疑态度，认为它有损独生子的尊位。我们常常习惯于在模糊的相似、与原物相去甚远的情况下持有\Quote{相似}的观念。我自己赞成同质，因为它更不容易被不当解释。但是，我亲爱的先生，你为何不来看我，好让我们当面谈论这些崇高的主题，而非将它们托付给无生命的书信——尤其是我已决定不发表我的看法？并且请不要对我套用\Person{狄奥根尼}对\Person{亚立山大}说的话：\Quote{你到我这里与我到你那里一样远。}我因健康不佳几乎不得不像植物一样停留在一处；此外，我认为\Quote{默默无闻地生活}是主要的善之一。我听说你身体健康；你使自己成了世界公民；你来看我时，不妨视作回家。对于一个行动者而言，你拥有群众和城市来展示你的善行，是十分恰当的。对我来说，安静是沉思和心智操练的最佳帮助，借此我紧贴天主。我在退隐中丰富地培育这种安静，依靠其赐予者天主的帮助。然而，若你不得不讨好伟人，而轻视我这样伏于低处的人，那你就写信吧，以此使我的生活更快乐。\par

\SourceRecord{Translated by Blomfield Jackson.；Nicene and Post-Nicene Fathers, Second Series；Vol. 8.；Edited by Philip Schaff and Henry Wace.；Buffalo, NY: Christian Literature Publishing Co.,；1895.；<http://www.newadvent.org/fathers/3202009.htm>.}

% structure: p0001 source=h1 target=section reason=page-title

\section{第十封信}

\Emph{\Person{圣巴西略·凯撒勒雅}}\par

\Emph{致一位寡妇。}\par

捕捉鸽子的技艺如下。当从事此艺的人捉到一只鸽子后，他们便驯养它，使它与自己一同进食。然后他们用甘饴之油涂抹它的翅膀，再放它飞回外面的鸽群。于是那甘饴之油的气味使自由的鸽群成为驯鸽主人的财产，因为其他鸽子都被芳香吸引，落在屋内。但我为何这样开始写信呢？因为我已经接纳了你的儿子\Person{狄约尼削}（他曾名\Person{狄奥墨得斯}），并用天主的甘饴之油涂抹了他灵魂的翅膀，将他派到你那里，为的是你能与他一同飞翔，并飞向他在我屋檐下所筑的巢。若我得以见到此事，并见到你——我可敬的朋友——被提升到我们崇高的生活，我便需要许多配得上天主的人，以向祂献上应得的一切荣耀。\par

\SourceRecord{Translated by Blomfield Jackson.；Nicene and Post-Nicene Fathers, Second Series；Vol. 8.；Edited by Philip Schaff and Henry Wace.；Buffalo, NY: Christian Literature Publishing Co.,；1895.；<http://www.newadvent.org/fathers/3202010.htm>.}

% structure: p0001 source=h1 target=section reason=page-title

\section{第十一封信}

\Emph{\Person{圣巴西略（凯撒勒雅）}}\par

\Emph{无地址。致某些朋友}。\par

因天主的恩宠，我与我们的儿子们共度了圣日，并因他们对天主极大的爱，向主举行了真正圆满的庆节。之后，我送他们平安地前往阁下处，并向我们仁爱的天主祈祷，求祂赐予他们一位平安天使，予以帮助和陪伴，并恩准他们见到您健康安宁，为的是无论您的命运落在何处，我终我一生，每当我听到您的消息，想到您事奉和感谢主时，都感到欣慰。如果天主容许您迅速摆脱这些忧虑，我恳求您不要让任何事情妨碍您来我处同住。我想您会发现没有人像我这样爱您，或如此珍视您的友谊。因此，只要圣者命定这种分离，您务必不要错过任何一个以书信安慰我的机会。\par

\SourceRecord{Translated by Blomfield Jackson.；Nicene and Post-Nicene Fathers, Second Series；Vol. 8.；Edited by Philip Schaff and Henry Wace.；Buffalo, NY: Christian Literature Publishing Co.,；1895.；<http://www.newadvent.org/fathers/3202011.htm>.}

% structure: p0001 source=h1 target=section reason=page-title

\section{第十二封书信}

\Emph{\Person{凯撒勒雅的圣巴西略}}\par

\Emph{致\Person{奥林比乌斯}}。\par

从前你曾给我写寥寥数语，如今连寥寥数语也没有了。你的简短很快会变成沉默。回到你从前的样子吧，不要让我因你的惜字如金而责备你。但即使是一封短信，我也乐意视作你深切爱情的标志。只要写信给我。\par

\SourceRecord{Translated by Blomfield Jackson.；Nicene and Post-Nicene Fathers, Second Series；Vol. 8.；Edited by Philip Schaff and Henry Wace.；Buffalo, NY: Christian Literature Publishing Co.,；1895.；<http://www.newadvent.org/fathers/3202012.htm>.}

% structure: p0001 source=h1 target=section reason=page-title

\section{第十三封信}

\Emph{\Person{凯撒勒雅的圣巴西略}}\par

\Emph{致\Person{奥林匹乌斯}}。\par

正如季节的所有果实都在它们适当的时间来临：春天有花朵，夏天有谷物，秋天有苹果，而冬天的果实就是交谈。\par

\SourceRecord{Translated by Blomfield Jackson.；Nicene and Post-Nicene Fathers, Second Series；Vol. 8.；Edited by Philip Schaff and Henry Wace.；Buffalo, NY: Christian Literature Publishing Co.,；1895.；<http://www.newadvent.org/fathers/3202013.htm>.}

% structure: p0001 source=h1 target=section reason=page-title

\section{\\WorkTitle{第十四封信}}

\\Emph{\\Person{圣巴西略·凯撒利亚}}\par

\\Emph{致其友\\Person{额我略}}.\par

我的兄弟\\Person{额我略}写信告诉我，他长久以来一直渴望与我在一起，并补充说你也有同样的心意；然而，我未能等待，一方面因为我难以相信，考虑到我曾多次失望，另一方面因为我发现自己被各种事务牵扯。我必须立即前往\\Place{本都}，在那里，若蒙天主允许，我也许能结束漂泊。在艰难地放弃了我曾对你抱有的虚幻希望——或者更确切地说，是梦幻（因为有人说希望是醒着的梦）——之后，我出发前往\\Place{本都}，寻找一个居住的地方。在那里，天主为我开启了一处完全符合我心意的所在，以至于我亲眼看到了我常在闲想中想象的情景：有一座高山，覆盖着茂密的森林，北面被清凉而清澈的溪流滋润。山脚下有一片平原，因不断从山上流下的水而肥沃，周边自然生长着茂密的树木，几乎形成一道篱笆；甚至超过了\\Place{卡吕普索岛}，\\Person{荷马}似乎认为那是世上最美的地方。这地方确实像一座岛，四面环绕：两侧被深谷切断；前面是一条刚刚从悬崖落下的河流，水流湍急，如同一道不可逾越的城墙；后面的山脉延伸出来，与深谷形成新月形，封锁了山脚的小路。只有一条通道，我掌握了它。在我住所的背后，又有一道峡谷，向上延伸到一道山脊，俯瞰着平原的广阔和界定平原的溪流，这条溪流在我看来，其美丽不亚于从\\Place{安菲波利斯}远眺的\\Place{斯特里蒙河}。因为后者缓缓流淌，几乎涨成一个湖，过于平静不像一条河；而前者是我所知最湍急的河流，且因上游的岩石而稍显浑浊；从那里奔腾而下，在深潭中回旋，为我或任何人形成了一幅极其宜人的景象；并且为乡民提供了无穷的资源，因为其深处有数不清的鱼。至于大地的气息或河流的微风，又何须多言？另一个人也许会赞叹繁花和鸣禽；但我无暇顾及这些思绪。然而，这地方最值得称赞的是，它适宜产出各种物产，并且培育出在我看来最甜美的产物——宁静；确实，它不仅远离城市的喧嚣，甚至很少有旅人经过，除非偶尔有猎人。这里猎物丰富，还有其他东西，但我很高兴地说，没有你们那里的熊和狼，而是鹿、野山羊、野兔之类。你不觉得我差点犯了一个愚蠢的错误吗？我曾急切地要用这个地方换取你的\\Place{提贝里纳}，那个全地最糟糕的坑穴。\par

那么，若我现在执着于此，请原谅我；因为我想，即使是\\Person{阿尔克迈翁}本人，在找到\\Place{埃基纳泽斯群岛}之后，也无法忍受继续漂泊了。\par

\SourceRecord{Translated by Blomfield Jackson.；Nicene and Post-Nicene Fathers, Second Series；Vol. 8.；Edited by Philip Schaff and Henry Wace.；Buffalo, NY: Christian Literature Publishing Co.,；1895.；<http://www.newadvent.org/fathers/3202014.htm>.}

% structure: p0001 source=h1 target=section reason=page-title

\section{\WorkTitle{第十五封信}}

\Emph{\Person{凯撒勒雅的圣巴西略}}\par

\Emph{致帝国司库\Person{阿尔卡狄乌斯}}.\par

我所属都会的市民们赐予我的恩惠大于他们所领受的，因为他们给了我一个机会向阁下写信。他们凭这封信从我这里赢得的善意，在我写信之前就已经因您惯常且与生俱来的对所有人的礼貌而得到了保证。但我认为这是一个极大的益处，能有幸向阁下写信，祈求圣洁的天主，使我得以继续喜悦，分享您慷慨施与者的欢乐，而您越来越中悦祂，同时您崇高地位的荣耀继续增长。我祈求，在适当的时候，我能满心欢喜地再次迎接那些将我这封信呈递给您的人，并派遣他们离开，称赞您对他们的体恤待遇，正如许多人所做的那样，并且我相信，他们对我的推荐在接近阁下崇高的尊荣时不会没有用处。\par

\SourceRecord{Translated by Blomfield Jackson.；Nicene and Post-Nicene Fathers, Second Series；Vol. 8.；Edited by Philip Schaff and Henry Wace.；Buffalo, NY: Christian Literature Publishing Co.,；1895.；<http://www.newadvent.org/fathers/3202015.htm>.}

% structure: p0001 source=h1 target=section reason=page-title

\section{第十六封信}

\Signature{圣巴西略·凯撒里亚}\par

\Emph{驳斥异端者欧诺弥。}\par

凡主张有可能发现实际存在之事的人，无疑是通过某种有序方法，借助对实际存在之事的认识来提升自己的理智。他先通过领会微小而易于理解的对象来训练自己，然后将自己的领会能力运用于超越一切理智的事物上。他以自己确实达到了对实际存在的理解为荣；那么，让他向我们解释一下最微小的可见事物的本性吧；让他告诉我们关于蚂蚁的一切。它的生命依赖于呼吸和气息吗？它有骨骼吗？它的身体是由筋腱和韧带连接的吗？它的筋腱是否被肌肉和腺体包裹？它的骨髓是否从额头到尾巴沿着背椎延伸？它是否通过包裹的神经膜推动其运动的肢体？它有肝脏吗，肝脏附近有胆囊吗？它有肾脏、心脏、动脉、静脉、膜、软骨吗？它有毛还是无毛？它有分趾的蹄还是脚分成几部分？它能活多久？它如何繁殖？它的妊娠期是多长？为什么蚂蚁既不是全部爬行也不是全部飞行，而有些属于爬行动物，有些则在空中穿行？这位以自己知道真实存在为荣的人，应该在此期间告诉我们关于蚂蚁的本性。然后，让他给我们一个类似的生理学描述，关于那超越一切人类理解力的权能。但是，如果你的知识尚未能领会微不足道的蚂蚁的本性，你怎能夸口自己能构想不可理解的天主的权能呢？\par

\SourceRecord{Translated by Blomfield Jackson.；Nicene and Post-Nicene Fathers, Second Series；Vol. 8.；Edited by Philip Schaff and Henry Wace.；Buffalo, NY: Christian Literature Publishing Co.,；1895.；<http://www.newadvent.org/fathers/3202016.htm>.}

% structure: p0001 source=h1 target=section reason=page-title

\section{第十七封信}

\Emph{\Person{凯撒勒雅的圣巴西略}}\par

\Emph{致\Person{奥力振}}。\par

聆听您令人愉悦，阅读您亦然；我认为您的著作给了我更大的快乐。一切感谢归于我们仁慈的天主，祂没有容许真理因国家主要权势者的背叛而受苦，反而藉着您使真宗教教义的辩护得以圆满而令人满意。像毒芹、附子和其它有毒的草药一样，它们开花片刻之后，便会迅速凋谢。但主为报答您为祂的圣名所说的一切而赐予您的赏报，却永远常新。因此，我祈求天主赐给您家中一切幸福，并使祂的祝福降于您的儿子们。我非常高兴看到并拥抱那些高尚的男孩，他们是您卓越善良的生动形象，我为他们的祈祷祈求一切父亲所能祈求的。\par

\SourceRecord{Translated by Blomfield Jackson.；Nicene and Post-Nicene Fathers, Second Series；Vol. 8.；Edited by Philip Schaff and Henry Wace.；Buffalo, NY: Christian Literature Publishing Co.,；1895.；<http://www.newadvent.org/fathers/3202017.htm>.}

% structure: p0001 source=h1 target=section reason=page-title

\section{第十八封信}

\Emph{\Person{圣巴西略·凯撒勒雅}}\par

\Emph{致\Person{玛加略}与\Person{若望}}。\par

田间的劳作对农夫而言并不新奇；水手若在海上遇到风暴，也不会惊讶；夏日炎炎中的汗水是雇工常见的经历；同样，对那些选择过圣洁生活的人而言，今世的磨难也不能预料不到。他们各自行当的劳苦是众所周知且与生俱来的，并非因其本身而被选择，而是为了期待中美好事物的享受。在这些情况下，每一件事中作为患难慰藉的，正是那真正构成一切人类生活的纽带与联系——希望。如今，为地上的果实或为属世之事而劳作的人，有些只是想象中享受他们所期待的，结果却大失所望；即使在另一些人中，结果符合了期望，也很快需要另一个希望，因为第一个希望转瞬即逝，消逝不见。唯独为圣洁与真理而劳作的人，他们的希望不会被任何欺骗所摧毁；任何结果都不能毁灭他们的劳苦，因为等待他们的天国是坚固而确定的。因此，只要真理之言在我们一边，绝不要因谎言的诽谤而丝毫沮丧；不要让帝王的恐吓吓倒你；不要因亲友的嘲笑与讥讽而悲伤，也不要因那些假装关心你、并以其最具诱惑力的诱饵——假装给予善意劝告——而谴责你的人而难过。让健全的理智与这一切争战，并呼求我们的主耶稣基督——真宗教的导师——的庇护与援助，为祂受苦是甘甜的，\Quote{死亡便是获益。} \BibleRef{斐 1:21}\par

\SourceRecord{Translated by Blomfield Jackson.；Nicene and Post-Nicene Fathers, Second Series；Vol. 8.；Edited by Philip Schaff and Henry Wace.；Buffalo, NY: Christian Literature Publishing Co.,；1895.；<http://www.newadvent.org/fathers/3202018.htm>.}

% structure: p0001 source=h1 target=section reason=page-title

\section{第十九封信}

\Emph{圣\Person{巴西略}（凯撒勒雅的）}\par

\Emph{致我友\Person{额我略}。}\par

前天我收到了你的信。证明这是你的信，与其说是笔迹，不如说是独特的风格。寥寥数语表达了丰富的含义。我没有当场回复，因为我当时不在家，而送信人在把包裹交给我的一个朋友后就离开了。现在，我可以通过\Person{伯多禄}写信给你，同时既回复你的问候，又给你一个回信的机会。写一封像你写给我的那样简洁的信件当然不费什么功夫。\par

\SourceRecord{Translated by Blomfield Jackson.；Nicene and Post-Nicene Fathers, Second Series；Vol. 8.；Edited by Philip Schaff and Henry Wace.；Buffalo, NY: Christian Literature Publishing Co.,；1895.；<http://www.newadvent.org/fathers/3202019.htm>.}

% structure: p0001 source=h1 target=section reason=page-title

\section{\WorkTitle{第二十書信}}

\Person{凱撒勒雅的聖巴西略}\par

\Emph{致智者萊昂提烏斯}.\par

我也並非時常寫信給你，但並不比你寫給我的少，儘管有許多人從你那裡來到此處。你若託付他們每一位接連不斷地帶信，便沒有什麼能阻止我似乎實際與你相處，並如我們在一起般享受其樂，因為來者絡繹不絕。但你為何不寫信呢？對一位智者而言，寫信並非難事。不，即使你的手疲倦了，也無需親自書寫；別人會代勞。只需你的口舌即可。即使它不對我說話，也必定會對你的某位同伴說話。若無人與你同在，它會自言自語。當然，一位智者與雅典人的口舌，正如春天喚醒夜鶯歌唱時一般，很少沉默。至於我，目前所從事的眾多事務或許可以為我缺乏書信提供些許藉口。或許，由於我的風格因經常與日常談話接觸而受損，使我有些猶豫於向像你這樣的智者致函，因為除非你聽到合乎你智慧的話語，否則你必定會惱怒且毫不留情。相反，你理當利用一切機會使你的聲音傳揚於外，因為在我所知的所有希臘人中，你是最好的演說家；我認為我認識你們中最著名的人；因此，你的沉默實在沒有藉口。但關於這點，就此為止。\par

我已將我反對歐諾米烏斯的著作寄給了你。它們應被稱為兒戲，還是稍微嚴肅些的東西，由你判斷。就你自己而言，我認為你不再需要它們；但我希望，若你遇到任何乖僻之人，它們不會是無用的武器。我這樣說，與其說是因為我對自己論著的力量有信心，不如說是我深知你是那種能舉一反三的人。若你有任何覺得比應有的軟弱之處，請不要猶豫，告訴我錯誤。朋友與奉承者之間的主要區別在於：奉承者說話為了取悅，朋友則不會省略即使令人生厭的話。\par

\SourceRecord{Translated by Blomfield Jackson.；Nicene and Post-Nicene Fathers, Second Series；Vol. 8.；Edited by Philip Schaff and Henry Wace.；Buffalo, NY: Christian Literature Publishing Co.,；1895.；<http://www.newadvent.org/fathers/3202020.htm>.}

% structure: p0001 source=h1 target=section reason=page-title

\section{\WorkTitle{第二十一封信}}

\Signature{圣巴西略·凯撒勒雅}\par

\Emph{致智者莱昂提乌}\par

那位优秀的犹利安似乎从普遍的状况中为自己的私事捞到了一些好处。如今一切都充斥着要求征收的税款，他也被猛烈地催讨和控告。只是问题不在于欠缴的税款，而在于信件。但他如何成了欠款者，我不知道。他向来付出一封信，也收到一封信——正如他收到这一封一样。但或许你偏爱那著名的\Quote{四倍之多}。因为连毕达哥拉斯学派都那么喜爱他们的四元组，正如这些现代的税吏喜爱他们的\Quote{四倍之多}。然而或许更公平的做法恰恰相反，像你这样一位词藻如此丰富的智者，本该向我抵押\Quote{四倍之多}。但千万别以为我写这些都是出于不悦。我能从你那里得到哪怕一顿训斥，也已是喜出望外。据说善与美的人做任何事都带着善与美的附加。甚至他们的悲伤与愤怒也是得体的。无论如何，人宁愿看到朋友对自己发怒，也不愿看到别人对自己奉承。因此，不要停止像上次那样提出指控！这指控本身意味着一封信；对我而言，没有什么比这更珍贵、更令人愉悦的了。\par

\SourceRecord{Translated by Blomfield Jackson.；Nicene and Post-Nicene Fathers, Second Series；Vol. 8.；Edited by Philip Schaff and Henry Wace.；Buffalo, NY: Christian Literature Publishing Co.,；1895.；<http://www.newadvent.org/fathers/3202021.htm>.}

% structure: p0001 source=h1 target=section reason=page-title

\section{第二十二封信}

\Emph{\Person{圣巴西略（凯撒勒雅）}}\par

\Emph{无收信人。论独修者生活的成全。}\par

1. 受默感的圣经为所有渴望悦乐天主的人制定了诸多规范。但就目前而言，我认为只需就你们中间近来引起争议的问题作一简要提醒，这是我通过研读受默感的圣经本身所学到的。我愿为勤勉的学生留下易于理解的详细证据，以便他们也能教导他人。基督徒应当有与天上召叫相称的心志，其生活与言行应与基督的福音相称。基督徒不应心持二意，也不应被任何事物牵引而忘却天主及其旨意与审判。基督徒应在一切事上超越律法之下的义，既不发誓也不说谎。他不可毁谤；\BibleRef{铎 3:2}不可使用暴力；\BibleRef{弟前 2:13}不可争斗；\BibleRef{弟后 2:24}不可报复；\BibleRef{罗 12:19}不可以恶报恶；\BibleRef{罗 12:17}不可发怒。\BibleRef{玛 5:22}基督徒应当忍耐，\BibleRef{雅 5:8}无论受何苦，都要按主的诫命，适时指正作恶者，\BibleRef{铎 2:15}不是出于自我辩解，而是为了弟兄的悔改。\BibleRef{玛 15:18}基督徒不应在弟兄背后说任何毁谤他的话，因为即使所言属实，这也是诽谤。他应当远离那恶言攻击他的弟兄；不应耽于戏谑；\BibleRef{弗 5:4}不应嬉笑，甚至不应容忍引人发笑之事。他不可闲谈，说些对听者无益的话，也不可说些超出天主允许我们必要及合宜使用的话；\BibleRef{弗 5:4}好使工作者尽可能在沉默中尽力工作；受托谨慎分施圣言以建立信德的人，应向工作者提出善言，免得使天主的圣神忧伤。任何人进来，在那些受委托管理全体纪律的人认定此事悦乐天主、有益公共福祉之前，不得擅自与任何弟兄交谈或说话。基督徒不应受酒所奴役；\BibleRef{伯前 4:3}也不应贪求肉食，\BibleRef{罗 14:21}一般而言，不应贪图饮食之乐，\BibleRef{弟后 3:4}\BibleQuote{因为凡在竞赛中争取桂冠的，在一切事上都有节制。}\BibleRef{格前 9:25}基督徒应将所给予他使用的一切视为属于主，而非据为己有或积蓄起来；应谨慎留意一切属于主的事物，若有被丢弃或忽视的东西，不可忽略。基督徒不应自以为主，而应如同被天主赐予，作同心合意弟兄的奴仆；但\Quote{各人在自己的行列中。}\par

2. 基督徒应在必需品匮乏时、在劳苦工作中从不抱怨，因为此事责任在于有权限的人。绝不可有喧哗、暴躁或激动表现以表怒意，也不可因分心而丧失对天主亲临的确信。\par

声音应有节制；没有人可以粗暴或轻蔑地回答问题或行事，凡事要温和，尊重众人。不可有狡诈的眼色，也不可有使弟兄忧伤或表示轻蔑的行为与姿态。\BibleRef{罗 14:10}衣着与鞋履的炫耀应当避免，这是虚荣的浪费。身体所需之物应使用廉价物品；不可花费超乎必要或仅为奢侈之物；这是滥用我们的财产。基督徒不应追求尊荣或争先。\BibleRef{谷 9:37}每个人都应将他人置于自己之前。\BibleRef{斐 2:3}基督徒不应悖逆。\BibleRef{铎 1:10}能工作的人不应吃闲饭，\BibleRef{得后 3:10}即使忙于为基督的光荣而行的善工，也当勉力去做力所能及的工作。\BibleRef{得前 4:11}每位基督徒，在上级的许可下，凡事都当合理而确信地去做，甚至饮食，也要为光荣天主而行。\BibleRef{格前 10:31}基督徒未经负责安排此类事务之人的许可，不应从一项工作换到另一项工作；除非有必须的急事突然召唤某人去援助无助者。每个人都应坚守自己的岗位，不可越界介入未被吩咐的事，除非主管当局认为某人需要帮助。没有人应当从一个工场游荡到另一个工场。凡事不可出于竞争或争执。\par

3. 基督徒不应嫉妒他人的名声，也不应因任何人的过犯而喜悦；\BibleRef{格前 13:6} 他应在基督的爱里，为弟兄的过犯忧伤痛苦，并为弟兄的善行而喜乐。\BibleRef{格前 12:26} 他不应对罪人冷漠或缄默。\BibleRef{弟前 5:20} 指出他人错误的人，应以温柔的心去做，\BibleRef{弟后 4:2} 怀着天主的敬畏，并以归化罪人为目的。\BibleRef{弟后 4:2} 被证明有错或受责备的人，应欣然接受，认识到自己因被纠正而得益。当有人被指控时，另一个人不应在他本人或任何人面前反驳指控者；但若在任何时候指控在某人看来似乎毫无根据，他应私下与指控者讨论，要么提出证据，要么使其信服。每个人都应尽可能与对他有怨言的人和好。没有人应向悔改的罪人怀恨在心，而应真心宽恕他。\BibleRef{格后 2:7} 说自己已悔罪的人，不仅应为自己的罪感到痛悔，还应结出与悔改相称的果实。\BibleRef{路 3:8} 首次犯错受纠正并获宽恕者，若再犯罪，是为自己预备比先前更重的义怒审判。\BibleRef{希 10:26} 经过第一次和第二次劝诫后仍坚持错误者，\BibleRef{铎 3:10} 应被带到掌权者面前，或许经众人责备后他会羞愧。即使这样仍不能改正，就应将他从其他人中除去，如同使人绊倒的人，视他为外邦人和税吏，\BibleRef{玛 18:17} 以保护顺服的人，如所说：\Quote{当恶人跌倒时，义人战栗。}他应被当作从身体上割下的肢体而哀伤。太阳不应落在弟兄的愤怒上，\BibleRef{弗 4:26} 免得黑夜隔绝弟兄与弟兄，并在审判之日使控告成立。基督徒不应等候改正自己的机会，因为明天不可靠；许多人计谋多端，却未活到明天。他不应因暴饮暴食而受迷惑，因为夜里会生梦境。他不应为过度劳作而分心，也不应超出足够的界限，如宗徒所说：\Quote{只要有衣有食，我们就当知足。} \BibleRef{弟前 6:8} 不必要的富足看似贪婪，贪婪被谴责为偶像崇拜。\BibleRef{哥 3:5} 基督徒不应贪爱钱财，也不应为无益的目的积蓄财宝。来到天主面前的人，应在一切事上拥抱贫穷，并紧守对天主的敬畏，如所说：\Quote{求你将我的肉体钉在你的敬畏中，因为我害怕你的判断。} 愿主赐予你完全确信我所言，结出与圣神相称的果实，为光荣天主，赖天主的恩惠，和我们的主耶稣基督的合作。\par

\SourceRecord{Translated by Blomfield Jackson.；Nicene and Post-Nicene Fathers, Second Series；Vol. 8.；Edited by Philip Schaff and Henry Wace.；Buffalo, NY: Christian Literature Publishing Co.,；1895.；<http://www.newadvent.org/fathers/3202022.htm>.}

% structure: p0001 source=h1 target=section reason=page-title

\section{第二十三封书信}

\Emph{\Person{凯撒勒雅的圣巴西略}}\par

\Emph{致一位隐修士}.\par

有一个人，如他所说，因谴责今世之虚荣，并觉察其喜乐止于此世，只因它们只提供\OriginalTerm{永火}{eternal fire}的材料并迅速消逝，便来我处，渴望脱离这邪恶悲惨的生活，弃绝肉体的享乐，并在将来踏上一条通往主之居所的道路。现若他对其真福之志向坚定不渝，且灵魂中怀有光荣可赞的热忱，尽心、尽力、尽意爱主他的天主，则阁下有必要向他指明那窄路和狭道的艰辛与困苦，并以那尚不可见、却为所有堪当主的人所应许的美好事物之望德坚定他。为此我写信恳求您在基督内无与伦比的成全，若可能，塑造其品格，并在我缺席的情况下，按天主所喜悦的方式促成其弃绝，并确保他接受依照圣教父们所规定并以文字颁布的基本训导。且务要使他面前摆上一切苦修纪律所必需的事，使他在自愿接受为信仰所做的劳苦、顺服主轻松的轭、并效法那为我们成了贫穷（\BibleRef{格后 8:9}）且取了肉身的祂而度一种效法祂的生活之后，能被引入那生命，从而无阻地奔向至高召唤的奖赏，并获主的嘉许。他渴望在此领受天主之爱的华冠，我却暂缓他，因我愿与阁下您一同为他从事这类争斗而傅油，并从您中间指派一人为他所选作其训练者，高尚地训练他，以恒常而蒙福的关怀使他成为经过考验的角力者，击伤并推翻今世黑暗之王及邪恶的属灵势力，对此，如蒙福宗徒所说，\BibleQuote{我们的角力。}（\BibleRef{弗 6:12}）便是。凡我愿与您一同做的，愿您在基督内的爱代我而行。\par

\SourceRecord{Translated by Blomfield Jackson.；Nicene and Post-Nicene Fathers, Second Series；Vol. 8.；Edited by Philip Schaff and Henry Wace.；Buffalo, NY: Christian Literature Publishing Co.,；1895.；<http://www.newadvent.org/fathers/3202023.htm>.}

% structure: p0001 source=h1 target=section reason=page-title

\section{第二十四封信}

\Emph{圣巴西略·凯撒勒雅}\par

\Emph{致\Person{阿塔纳西}，\Place{安基拉}主教\Person{阿塔纳西}之父}\par

即使并非不可能，最难实现的事情之一就是超然于毁谤之上，这一点我本人深信不疑，我想阁下也是如此。然而，不因自己的行为给社会的吹毛求疵者或伺机抓我们把柄的恶意中伤者留下把柄，不仅可能，而且是一切按照真正宗教准则智慧地安排生活者特有的品质。请不要以为我如此单纯轻信，以致不经适当调查就接受任何人的贬损之言。我牢记圣神的告诫：\Quote{不可接受虚假的报告。}但你们这些博学之士自己说：\Quote{可见之事表明不可见之事。}因此我恳请——（请勿见怪，我似乎像在教训人；因为\BibleQuote{天主拣选了世上软弱的}和\BibleQuote{世上卑贱的，}\BibleRef{格前 1:27}常常藉此拯救那些得救的人）；我所说和所敦促的是：在言语和行为上严加注意合宜，并按照宗徒的训诫，\BibleQuote{凡事都不叫人有妨碍。}\BibleRef{格后 6:3}一个在获取知识上辛勤劳作、治理过城市和国家、并热切效法祖先崇高品格的人，其生活本身应当成为崇高品格的典范。你现在不应仅仅在言语上向你的子女表现你的慈爱，正如你从成为父亲以来长久表现的那样；你不仅应表现出禽兽也有的那种自然亲情，如你自己所说的，经验也证实了这一点。你应当使你的爱更进一步，成为一种更加个人化、更加自愿的爱，因为你看到你的子女配得上父亲的祈祷。在这一点上，我不需要被说服。事实的证据已经足够。不过，为了真理，我要说一件事：不是我们的兄弟\Person{弟茂德}，那位\OriginalTerm{乡村主教}{Chorepiscopus}将外面传闻的消息告诉了我。因为他无论是口传还是书信，从未传达过任何毁谤之言，无论大小。我确实听到了一些事，但并非弟茂德控告你。然而，无论我听到什么，我至少会效法\Person{亚历山大}的榜样，为被告保留一只耳朵。\par

\SourceRecord{Translated by Blomfield Jackson.；Nicene and Post-Nicene Fathers, Second Series；Vol. 8.；Edited by Philip Schaff and Henry Wace.；Buffalo, NY: Christian Literature Publishing Co.,；1895.；<http://www.newadvent.org/fathers/3202024.htm>.}

% structure: p0001 source=h1 target=section reason=page-title

\section{第二十五封书信}

\Emph{圣巴西略·凯撒勒雅}\par

\Emph{致\Place{安齐拉}主教\Person{亚大纳削}。}\par

1. 我从那些来自安齐拉的人那里得到消息，他们人数众多，不可胜数，但所说的话却都一致：您，我非常亲爱的人（我该如何说才能不冒犯您呢？），您提及我时，并非以令人愉快的言辞，也非以您的品格所应许的态度。然而，我早已懂得人性的软弱及其易于从一个极端转向另一个极端的倾向；所以，请您确信，关于人性的一切都不能使我惊讶，也没有任何变化是出乎意料的。因此，我的处境变坏，昔日的尊重变成了责备和侮辱，我并不介意。但有一件事确实令我既惊讶又骇异，那就是您竟然对我有这种看法，以至于对我感到愤怒和义愤，而且，如果从您听众的报告中可信，您已经走到极端，发出了威胁。对于这些威胁，我不否认，我确实笑了。真的，如果被这些吓唬人的东西吓倒，我不过是个孩子。然而，令我感到惊恐和痛苦的是，您——我曾相信，您是为教会的安慰而被保存的，是在许多人背离时的真理支柱，是古老而真实之爱的种子——竟如此顺应时势，以至于被您偶然遇到的第一个人的诽谤所影响，甚于您长期以来对我的了解，并无端地被引诱去怀疑荒谬之事。\par

2. 但是，如我所说，目前我暂不深究此事。亲爱的先生，难道在朋友之间，通过一封短信讨论您想提出来的问题，会是一件过于困难的事吗？或者，如果您不愿将此类事情托付于笔墨，是否可以让我到您那里去呢？但如果您忍不住要发泄，而您那无法控制的愤怒又不容片刻耽搁，至少您本可以利用您身边那些天生适于处理机密事务的人，作为与您沟通的媒介。然而现在，在所有因这样或那样的原因接近您的人中，谁的耳朵没有灌满这样的说法：我撰写并创作了某些\Quote{害物}？因为那些逐字引用您的话的人说，这正是您使用的词语。我越是集中思想考虑此事，就越感到困惑。我产生了这样的想法：难道某位异端者触怒了您正统的信仰，并恶意地将我的名字安在他自己的著作上，从而驱使您说出那个词？因为您，这位为真理进行过伟大而著名的斗争的您，是绝不能容忍对我所写的反对那些胆敢说天主圣子在本性上不像天主圣父，或亵渎地将圣神描述为受造之物之人的作品施加如此暴行的。您自己可以解除我的困惑，只要您坦率地告诉我，是什么激怒了您，使您对我如此不满。\par

\SourceRecord{Translated by Blomfield Jackson.；Nicene and Post-Nicene Fathers, Second Series；Vol. 8.；Edited by Philip Schaff and Henry Wace.；Buffalo, NY: Christian Literature Publishing Co.,；1895.；<http://www.newadvent.org/fathers/3202025.htm>.}

% structure: p0001 source=h1 target=section reason=page-title

\section{第二十六封信}

\Emph{\Person{圣巴西略}（\Place{凯撒勒雅}）}\par

\Emph{致\Person{凯撒略}，\Person{格列高利}的兄弟。}\par

感谢天主，因祂在你身上显示了奇妙的权能，将你从如此可怕的死亡中保存下来，使你和你的祖国，以及我们这些朋友仍能相聚。我们理当不要忘恩负义，也不可辜负如此宏大的仁慈，而应尽己所能，讲述天主的奇妙化工，以行动庆祝我们所经验的恩惠，而非仅以言语报谢。我们应当真正成为我因你身上所行的奇迹而相信你现在已是的人。我们更劝诫你事奉天主，不断加深敬畏，迈向成全，好使我们成为明智的管家，管理我们的生命，天主的良善为我们保留这生命。因为如果对我们所有人来说，都有命令要\BibleQuote{将自己奉献给天主，如同从死者中复活的人一样}，\BibleRef{罗 7:13}，那么对那些已从死亡之门被提起的人，这命令岂不更加强烈？我相信，若我们常能保持我们在危难时刻的心思，这事就能达到最佳效果。因为，我想，我们生命的虚浮呈现在我们面前，我们感到所有属于人的东西，既暴露于变迁之中，就没有什么可靠，没有什么坚固。我们感到，正如所料，对过去的悔改；我们为未来立下承诺，倘若得救，必事奉天主，并谨慎省察自己。若死亡的迫近给了我任何反省的缘由，我想你也一定被相同或相近的思想所触动。因此，我们有义务偿还一项必须偿还的债务，既因天主赐予我们的恩典而喜乐，同时又为将来焦虑。我斗胆向你提出这些建议。现在取决于你，要像我们面对面交谈时你常做的那样，善意而和善地接受我的话。\par

\SourceRecord{Translated by Blomfield Jackson.；Nicene and Post-Nicene Fathers, Second Series；Vol. 8.；Edited by Philip Schaff and Henry Wace.；Buffalo, NY: Christian Literature Publishing Co.,；1895.；<http://www.newadvent.org/fathers/3202026.htm>.}

% structure: p0001 source=h1 target=section reason=page-title

\section{第二十七书}

\Emph{圣巴西略（凯撒勒雅的）}\par

\Emph{致\Person{欧瑟比}，\Place{萨莫萨塔}主教}。\par

当我藉天主的恩宠并赖您的祈祷的助佑，似乎已从病中稍有康复，并恢复气力时，冬天却来临了，将我禁锢在家中，迫使我留在原地。诚然，它的严酷程度远不及往年，但即便如此，也足以使我在整个冬季不仅不能旅行，甚至不敢把头探出门外。但对我而言，得以——即便仅凭书信——与尊驾相通，并在期待回音中安享宁静，绝非等闲之事。然而，若季节许可，且若许我活得更久，并若匮乏不阻止我踏上旅程，或许藉您的祈祷的助佑，我能实现这恳切的心愿，能在您自己的炉边找到您，并在充裕的闲暇中，饱飨您浩博的智慧宝藏。\par

\SourceRecord{Translated by Blomfield Jackson.；Nicene and Post-Nicene Fathers, Second Series；Vol. 8.；Edited by Philip Schaff and Henry Wace.；Buffalo, NY: Christian Literature Publishing Co.,；1895.；<http://www.newadvent.org/fathers/3202027.htm>.}

% structure: p0001 source=h1 target=section reason=page-title

\section{第二十八封书信}

\Person{凯撒勒雅的圣巴西略}\par

\Emph{致新凯撒勒雅教会：慰问书}\par

1. 那降临在你们身上的事强烈地推动我去探望你们，带着双重的目的：一是与你们这些亲近我、为我所爱的人们一同向那位已蒙福的亡者致以全部敬意；二是通过亲眼目睹你们的哀伤，更紧密地与你们共担患难，从而能够与你们商议当行之事。然而，有许多障碍阻止我亲自到你们当中，因此我只能以书信与你们相通。这位已故者令人钦佩的品质——我们正是因这些品质而估量我们损失之重大——实在多得无法在信中一一列举；此外，现在也不是我们因悲伤而精神萎靡时讨论他众多善行的时候。因为，对于他所做的一切，我们何曾能忘记？我们有什么认为值得缄默不语？要一下子全部道出是不可能的；若只说一部分，我担心会违背真理。一位在一切人力所能及的善事上超越了他所有同代的人逝去了：他是国家的支柱，教会的光荣，真理的柱石和支撑，基督信仰的稳固，朋友的保护者，反对者的顽强敌人，祖先原则的守护者，革新的仇敌；他在自己身上展示了教会古老的风范，并使他治下的教会状况符合那古老的法度，如同符合一个神圣的模型，以致所有与他共同生活的人都仿佛与那些两百多年前如同光耀在世界中的人们同处一个社会。所以，你们的主教并未提出任何出于自己的、新奇的东西；而是如同梅瑟的祝福所说，他懂得从自己心中隐秘而良善的府库里拿出\BibleQuote{陈年旧藏，并因着新物而显为陈旧}\BibleRef{肋 26:10}。因此，在主教同侪的集会上，他并非按年龄排列，而是因着他智慧的古老，所有人一致同意他居于众人之上。无人观察你们的状况后需远寻便可看出这种训练方式的好处。因为，就我所知，在这样的大风大浪之中，唯独你们，或至少你们与极少数人，在他的良好引领下得以生活而不被波涛动摇。你们从未被异端那冲击的狂风击中——这风使不稳定的灵魂遭遇翻船和淹没；为使你们永远不受其害，我祈求那统管万有的主，祂曾赐予那建立你们教会的第一位奠基者——祂的仆人\Person{额我略}——长久的安宁。\par

现在不要失去那安宁；不要因过度的哀悼和完全沉溺于悲伤，而把行动的机会交到那些图谋你们祸患的人手中。如果你们必须哀悼（我不允许这样，免得你们在这方面如同\BibleQuote{那些没有希望的人}\BibleRef{得前 4:13}），那么，若你们觉得合适，可以像一些涉水的歌队那样，选出一个领唱者，随他一同唱起眼泪的哀歌。\par

2. 然而，如果你们所哀悼的那位并未达到极高的年迈，那么，就他对你们教会的治理而言，他确实被赋予了不短的生命期限。他拥有足够强的体力，使他能在困苦中展现精神的力量。也许你们中有人以为时间会增加同情、增添情感，而不会导致厌倦；因此，你们领受善待的时间越长，就越敏感于失去它。你们或许认为，对于一个义人，好人甚至尊敬他的影子。愿你们中许多人都能这样觉得！我绝无贬低我们朋友之意！但我劝告你们要以男子汉的坚忍承受痛苦。我自己绝非对那些哀哭他们损失的人们所说的一切无动于衷。那曾如汹涌河流灌满我们耳朵的言语如今沉寂了：一颗从未被探测过的心灵深处，按人的说法，已如虚幻的梦境消逝了。谁的目光有他那样锐利，能预见未来？谁有他那样的坚定和心智力量，能闪电般投入行动之中？哦，\Place{新凯撒勒雅}，你已遭受诸多苦难，从未被如此致命的损失击打过！如今你的容颜失色了，你的美丽枯萎了；你的教会沉默无语；你的集会充满了哀容；你的神圣会议渴求带领者；你的圣言等待阐释者；你的孩童失去了父亲，你的长者失去了兄弟，你的显贵失去了居首者，你的民众失去了保护者，你的穷人失去了支持者。每个人用最贴近自己心坎的称呼呼唤他，扬起那在各自悲伤看来最恰当、最合适的哀号。但我的话语被这含泪的喜乐带到了何处？我们不会警醒吗？我们不会聚会吗？我们不会仰望我们共同的主吗？祂允许祂的每位圣人在自己的世代中服务，并在祂自己指定的时刻召回他们归于自己。现在，要适时记起那位在向你们宣讲时经常说的话：\BibleQuote{你们要提防狗，提防作恶的工人}\BibleRef{斐 3:2}。狗很多。为何我说狗？不如说凶恶的狼，它们披着羊皮隐藏诡计，在全地撕裂基督的羊群。你们必须提防它们，在某个警醒的主教保护之下。这样的主教是你们应当祈求的，从你们心中除去一切竞争和野心；这样的主教是主将指示给你们的。那一位主，从\Person{额我略}——你们教会的伟大捍卫者——的时代直到这位已蒙福的亡者，为你们设立了一位又一位，并时常彼此配合，如同宝石紧密相连，为你们的教会赐下了光荣的装饰。所以，你们无须对将来的主教绝望。主认识属于祂的人。祂或许会将那些我们或许未曾料想的人带到我们中间。\par

3. 我原已打算早早结束，但心中的痛苦不容我这样做。如今我以诸位教父、以真信德、并以我们蒙福的友人恳求你们：振奋你们的心灵，各人将眼前所行之事视为自己的切身要务，各人思量自己将首先承受结局的后果——无论结局如何——免得你们遭遇那屡见不鲜的境况：人人都把众人的共同利益留给邻人；于是，当各人在心中轻看所发生之事时，你们全体便因疏忽而不知不觉地为自己招致本可避免的灾祸。我恳请你们以善意接纳我所说的话，无论这是被视为邻人的同情，抑或主内同道的情谊，或者——更接近实情——是遵守爱德法律、又对沉默的危险心存畏惧之人的心声。我深信，直到主的日子，你们是我的夸耀，正如我是你们的夸耀；并且，赐予你们的牧者如何，将决定我是因爱德的纽带与你们更紧密相连，还是与你们完全隔绝。愿天主禁止后者！因天主的恩宠，必不会如此；此刻我本不愿说一句不恩宠的话。但我愿你们知道：尽管在为教会和平的努力中，那位蒙福之人并非常常在我身边——因为他亲口承认，因某些偏见——然而，我从未在意见上与他有分歧，并且在我与异端者斗争时，我始终呼求他的帮助。为此，我呼求天主及所有最了解我的人作证。\par

\SourceRecord{Translated by Blomfield Jackson.；Nicene and Post-Nicene Fathers, Second Series；Vol. 8.；Edited by Philip Schaff and Henry Wace.；Buffalo, NY: Christian Literature Publishing Co.,；1895.；<http://www.newadvent.org/fathers/3202028.htm>.}

% structure: p0001 source=h1 target=section reason=page-title

\section{第二十九封}

\Person{凯撒勒雅的圣巴西略}\par

\Emph{致安基拉教会。慰唁书。}\par

你们遭遇的灾祸消息极为沉痛，我的惊愕长久使我缄默无言。我犹如一人被巨雷震聋了耳朵。随后，我这哑然无语的状态略有好转。如今，我为这事哀悼，无人能免于哀悼，并在我的哀叹中，将这封信寄给你们。我写此信，并非为了安慰你们——因为谁能找到言语来治愈如此巨大的灾祸呢？——而是为了借此方式，尽我所能，向你们表达我内心的痛苦。此刻我需要\WorkTitle{耶肋米亚的哀歌}，或是任何其他曾动情哀悼过大灾难的圣人的哀歌。一位真正作为教会柱石和支柱的人倒下了；更确切地说，他已被从我们这里带走，进入了有福的生命，而危险不小：恐怕许多人在这支柱被抽离后也随之倒下，恐怕一些人的不健全被暴露。一张善言正义、以恩宠之语造就弟兄团体的口被封住了。那真正在天主内运行的思想的谋划消失了。啊！我必须自责，我多少次曾对他感到愤怒，因为他全心想离世与基督同在，而不愿为我们留在肉身内！今后我能将教会的忧患托付给谁？谁能分担我的烦恼？谁能分享我的喜乐？啊，何等可怕而悲伤的孤独！我岂不像旷野中的鹈鹕？然而，事实上，教会成员被他以领导联结如一人，结合成紧密的情感与共融，并将在和平的纽带中，为着精神的共融，得到保存并将永远保存。天主赐予我们恩惠，使那有福灵魂在教会中所做的一切崇高事业，坚固而不可动摇。但斗争并非轻微：唯恐因着主教选举再次兴起纷争和分裂，你们的一切工作被某些争吵所颠覆。\par

\SourceRecord{Translated by Blomfield Jackson.；Nicene and Post-Nicene Fathers, Second Series；Vol. 8.；Edited by Philip Schaff and Henry Wace.；Buffalo, NY: Christian Literature Publishing Co.,；1895.；<http://www.newadvent.org/fathers/3202029.htm>.}

% structure: p0001 source=h1 target=section reason=page-title

\section{\WorkTitle{书信三十}}

\Emph{\Person{圣巴西略}}\par

\Emph{致萨莫萨塔的厄乌塞伯。}\par

若我尽述迄今使我留在家中、无法启程拜见尊驾的所有原因，尽管我渴望前往，那将是一个无尽的故事。疾病接踵而至、严冬气候、工作压力，这一切您已尽知。如今，因我的罪过，我失去了母亲，这是我生活中唯一的慰藉。请勿见笑，我虽年迈，仍哀叹自己成了孤儿。请原谅我无法忍受与一个灵魂分离，在我面前未来的一切中，我视其无可比拟。于是，我的哀怨再次回到我身；我再次卧床，在虚弱中辗转反侧，几乎每时每刻都在期待生命的终结；而教会的状况与我的身体大致相似，没有良好的希望照耀，且每况愈下。与此同时，\Place{新凯撒勒雅}和\Place{安齐拉}已决定为死者选立继任者，目前尚属平安。那些阴谋反对我的人尚未被允许做出任何符合其苦毒和愤怒的事。对此我们毫不隐瞒，归功于您为教会所做的祈祷。那么，请不要厌倦为教会祈祷并恳求天主。请向那些有幸服事您的人致以所有问候。\par

\SourceRecord{Translated by Blomfield Jackson.；Nicene and Post-Nicene Fathers, Second Series；Vol. 8.；Edited by Philip Schaff and Henry Wace.；Buffalo, NY: Christian Literature Publishing Co.,；1895.；<http://www.newadvent.org/fathers/3202030.htm>.}

% structure: p0001 source=h1 target=section reason=page-title

\section{\WorkTitle{第三十一封书信}}

\Emph{\Person{圣巴西略（凯撒勒雅）}}\par

\Emph{致\Person{厄瑟比乌斯}，\Place{萨莫萨塔}主教。}\par

死亡仍与我们同在，因此我不得不留在此地，部分是由于分施职责，部分是由于对受苦者的同情。因此，即使现在，我也未能陪同我们可敬的弟兄\Person{希帕提乌斯}——我称他为弟兄，不仅是出于习惯用语，而是因为亲属关系，因为我们同出一源。你知道他病得很重。想到他一切安慰的希望都被切断，这令我痛苦，因为那些拥有医治恩赐的人未被允许在他的病例中使用他们惯常的疗法。因此，他再次恳求你祈祷的帮助。请接受我的恳求，请你给予他惯常的保护，既为你的缘故——因为你总是善待病人，也为我——我在为他求情。如果可能，请召唤那些非常神圣的弟兄到你身边，使他在你的眼前得到治疗。如果不可能，请发善心以书信送他前行，并把他推荐给更远的朋友们。\par

\SourceRecord{Translated by Blomfield Jackson.；Nicene and Post-Nicene Fathers, Second Series；Vol. 8.；Edited by Philip Schaff and Henry Wace.；Buffalo, NY: Christian Literature Publishing Co.,；1895.；<http://www.newadvent.org/fathers/3202031.htm>.}

% structure: p0001 source=h1 target=section reason=page-title

\section{\WorkTitle{书函三十二}}

\Signature{圣巴西略（凯撒勒雅的）}\par

\Emph{致\Person{索弗若纽}督军。}\par

我们的天主——亲爱的弟兄、主教\Person{额我略}，分担了这个时代的患难，因为他也像其他人一样，被接二连三的暴行所困扰，如同一个遭受意外打击的人。因为那些不敬畏天主的人，可能是被巨大的苦难所逼迫，正在辱骂他，理由是他们曾借钱给\Person{凯撒里乌斯}。确实，任何损失的问题并不严重，因为他早已学会轻视财富。问题在于，那些如此慷慨地分发了\Person{凯撒里乌斯}所有值钱财产的人，实际上得到的很少，因为他的财产掌握在奴隶和比奴隶好不了多少的人手中，并没有给执行人留下多少。这一点点财产他们以为没有抵押给任何人，就立刻花在了穷人身上，这不仅出于他们自己的偏好，也出于死者的嘱咐。因为\Person{凯撒里乌斯}临终时据说曾说：\Quote{我希望我的财产归于穷人。}于是他们遵从\Person{凯撒里乌斯}的意愿，适当地分配了这些财产。如今，\Person{额我略}身为基督徒，却贫穷，却陷入了商贾的喧闹之中。因此我想到将此事禀报阁下，以便您就\Person{额我略}之事向\OriginalTerm{财政大臣}{Comes Thesaurorum}表明您的意见，从而尊荣您多年所认识的人，荣耀主，祂视对祂仆人所作的就是对祂自己作的，也尊荣我，我特别与您相连。我希望，以您的大智，能想出办法使这些无法无天的人和难以忍受的烦扰得到解脱。\par

2. 没有人会如此不了解\Person{额我略}，以至于对他因贪财而提供不准确情况的说法有任何可耻的怀疑。证明他慷慨的证据并不难找。\Person{凯撒里乌斯}所剩的财产，他乐意地放弃给国库，以便财产可以保存在那里，国库官员可以回答那些攻击它并要求证据的人；因为我们不擅长此类事务。阁下可以知道，只要可能，没有人得不到他想要的东西就离开，每个人都毫无困难地得到了他所要求的。结果当然是很多人后悔起初没有要求更多；这又产生了更多的反对者，因为有了先前成功申请者的例子，一个虚假的索赔人接踵而至。我恳请阁下反对这一切，像一道介入的溪流一样，打断这些烦恼的连续性。您最知道如何帮助解决问题，无需等待我指教。我对此生的事务没有经验，无法看出摆脱困境的出路。以您的大智慧，发现一些帮助的方法。做我们的顾问。做我们的捍卫者。\par

\SourceRecord{Translated by Blomfield Jackson.；Nicene and Post-Nicene Fathers, Second Series；Vol. 8.；Edited by Philip Schaff and Henry Wace.；Buffalo, NY: Christian Literature Publishing Co.,；1895.；<http://www.newadvent.org/fathers/3202032.htm>.}

% structure: p0001 source=h1 target=section reason=page-title

\section{第三十三信}

\Emph{\Person{圣大巴西略}}\par

\Emph{\Person{阿布尔吉乌斯}}.\par

有谁像您一样懂得尊重旧谊、崇尚德性、并同情病者呢？如今，我的天主所爱的兄弟、主教\Person{额我略}卷入了一些无论如何都令人不快、且完全不合其性向的事务。因此，我认为最好是投靠您的护佑，力求从您那里获得解决我们困境的方案。一个天生和心性都不适于此类事务的人，竟被迫承担如此责任；穷人被索要钱财；一个早已决心度隐修生活的人被拖到公众面前——这实在是难以容忍的局面。至于您认为向\OriginalTerm{财政大臣}{Comes Thesaurorum}或其他人士进言是否有用，则取决于您的明智判断。\par

\SourceRecord{Translated by Blomfield Jackson.；Nicene and Post-Nicene Fathers, Second Series；Vol. 8.；Edited by Philip Schaff and Henry Wace.；Buffalo, NY: Christian Literature Publishing Co.,；1895.；<http://www.newadvent.org/fathers/3202033.htm>.}

% structure: p0001 source=h1 target=section reason=page-title

\section{第三十四书}

\Emph{\Person{凯撒勒亚的圣巴西略}}\par

\Emph{致\Place{萨莫萨塔}主教\Person{欧瑟伯}书。}\par

在此关头，我怎能缄默不语？若不能缄默，我又如何找到合乎时事之辞，使我的声音不像单纯的叹息，反倒像一场清晰表明灾祸之大的哀哭？唉！\Place{塔尔索}已亡。这是一项难以承受的痛苦，但它并非单独降临。更令人痛心的是，像这样一个与\Place{基里基雅}、\Place{卡帕多细雅}和\Place{亚述}相连的城邑，竟因两三个人的疯狂而轻易被弃，而你们却一直犹豫不决，商议对策，面面相觑。早该像医生那样行事（我久病在身，这类比喻并不缺乏）。当病人的疼痛加剧时，医生便使其失去知觉；同样，我们应祈祷我们的灵魂对痛苦变得麻木，免得陷入无法忍受的剧痛。在此困境中，我尚不缺少一种慰藉之法。我仰望你的仁爱；我试图借着对你的思念和追忆来减轻我的痛苦。当眼睛凝视过明亮的物体后，转而观看蓝绿色之物便可缓解；同样，对你的仁爱与关怀的追忆对我的灵魂起相同的作用；它是一帖温和的疗法，除去我的苦楚。当我想到你个人已尽人力之所能时，我更深切地感受到这一点。你若公正地判断，便已充分向我们证明，这场灾难决非由于你个人所致。你因热忱维护正义而从天主手中获得的赏报并不微小。愿主将你赐予我，也赐予祂的教会，以改善生活、引领灵魂，并愿祂再次恩准我得见你的殊荣。\par

\SourceRecord{Translated by Blomfield Jackson.；Nicene and Post-Nicene Fathers, Second Series；Vol. 8.；Edited by Philip Schaff and Henry Wace.；Buffalo, NY: Christian Literature Publishing Co.,；1895.；<http://www.newadvent.org/fathers/3202034.htm>.}

% structure: p0001 source=h1 target=section reason=page-title

\section{第三十五封}

\Emph{圣巴西略·凯撒勒雅}\par

\Emph{无收件人地址。}\par

我曾写信给你提到许多如同我自己一样的人；现在我要提及更多。穷人永不会断绝，我也永不会说\Quote{不}。没有任何人比我更亲密，更能在他力所能及的地方为我施恩，除了尊敬的弟兄\Person{莱昂提乌斯}。所以，请如此对待他的家，如同你发现我并非处于现今因天主助佑我所过的贫穷之中，而是拥有财富和地产。毫无疑问，你不会使我贫穷，反而会看顾我所拥有的，甚至增加我的产业。这就是我请求你在\Person{莱昂提乌斯}的家中行事的方式。你将从我这里得到你惯常的奖赏；因你劳苦所显出的善良和真诚，以及提前回应穷人的恳求，我向圣洁的天主献上祈祷。\par

\SourceRecord{Translated by Blomfield Jackson.；Nicene and Post-Nicene Fathers, Second Series；Vol. 8.；Edited by Philip Schaff and Henry Wace.；Buffalo, NY: Christian Literature Publishing Co.,；1895.；<http://www.newadvent.org/fathers/3202035.htm>.}

% structure: p0001 source=h1 target=section reason=page-title

\section{第三十六书}

\Emph{\Place{凯撒勒雅}的\Person{圣巴西略}}\par

\Emph{无收信人地址}。\par

我认为，贵下早已知道，本地的司铎是我的乳兄弟。我还能多说什么，以劝您仁慈地以友善的眼光看待他，并在他的事务中予以帮助？如果您爱我，正如我知道您爱我一样，我相信您会尽力帮助任何一个我看作是另一个我的人。那么，我要求什么？就是他不失去他原有的评级。实在，他在照料我的需要上不辞辛苦，因为我，如您所知，自己一无所有，而是依靠亲友的财物。那么，请您看顾我兄弟的家，如同看顾我的家一样，或者更确切地说，如同看顾您自己的家。天主必将因您对他的仁慈，而不断帮助您本人、您的家及您的家人。切望特别留心，不要让他在评级均等中受到任何损害。\par

\SourceRecord{Translated by Blomfield Jackson.；Nicene and Post-Nicene Fathers, Second Series；Vol. 8.；Edited by Philip Schaff and Henry Wace.；Buffalo, NY: Christian Literature Publishing Co.,；1895.；<http://www.newadvent.org/fathers/3202036.htm>.}

% structure: p0001 source=h1 target=section reason=page-title

\section{第三十七封信}

\Emph{凯撒勒雅的圣巴西略}\par

\Emph{无收信人姓名。}\par

我怀疑书信泛滥。我违背本意，因无法拒绝请愿者的缠扰，被迫开口。我写信是因为我想不出别的解脱方法，只能同意那些不断向我求取书信之人的恳求。我实在担心，既然许多人带着信离开，其中一封会被认为是那位兄弟。我承认，在我的故乡有许多朋友和亲属，并因主所赐予我的职位，我被置于\Emph{\OriginalTerm{代行父母职责}{in loco parentis}}的地位。其中就有这位我的乳兄弟，我乳母的儿子，我祈求我长大的那所房子能维持原有的税额，以免大人您在我们中间的居留——这对我们大家如此有益——反成为他的祸因。如今我也依赖同一所房子的供给，因为我身无长物，依靠爱我的人。所以我恳求您宽免我乳养之家的负担，如同您继续维持对我的供养。愿天主回报您，赐予您永远的安息。不过有一件事，且千真万确，我认为大人应该知道：那些奴隶中的大部分，从一开始就是我父母作为我养育费的等价物赐予他的。同时，这不应被视为绝对的赠与；他仅有终身使用权，所以若因他们而对他发生什么严重的事，他可随意将他们送回给我，于是我便以另一种方式为赋税和收税人负责。\par

\SourceRecord{Translated by Blomfield Jackson.；Nicene and Post-Nicene Fathers, Second Series；Vol. 8.；Edited by Philip Schaff and Henry Wace.；Buffalo, NY: Christian Literature Publishing Co.,；1895.；<http://www.newadvent.org/fathers/3202037.htm>.}

% structure: p0001 source=h1 target=section reason=page-title

\section{第三十八书}

\Emph{\Person{圣巴西略·凯撒勒雅}}\par

\Emph{致其弟额我略，论} \OriginalTerm{本质}{\GreekText{οὐσία}} \Emph{与} \OriginalTerm{位格}{\GreekText{ὑπόστασις}} \Emph{之区别。}\par

1. 许多人在研习神圣教义时，未能区分本质或实体中共同之处与位格的意义，因而得出相同概念，以为谈论\OriginalTerm{本质}{\GreekText{οὐσία}}或位格并无差别。结果，有些未加探究便接受这些论题之人，喜欢说\Quote{一个位格}，正如他们说一个\Quote{本质}或\Quote{实体}；而另一些接受三个位格的人，则以为根据这一宣认，必须藉数目类比也主张三个本质或实体。为避免你陷入同样错误，我特为你撰写此简短备忘录。简言之，这些词的意义如下：\par

2. 在所有名词中，有些用于复数且数目不同的主体，其意义较为普遍，例如\Emph{人}。当我们这样说时，使用该名词指示共同的本性，并不将意义局限于任何一个以此名称被认识的特定之人。例如，\Person{伯多禄}并不比\Person{安德肋}、\Person{若望}或\Person{雅各伯}更是\Emph{人}。因此，这谓词既是共同的，且延及所有归在同一名称下的个体，就需要某种区别记号，使我们能理解的不是一般意义上的人，而是特定的\Person{伯多禄}或\Person{若望}。\par

另一些名词的指称则较为有限；藉助这种限制，呈现在我们面前的不是共同的本性，而是任何事物的限制，就其特殊性而言，与同类的其他事物毫无共同之处；例如\Person{保禄}或\Person{弟茂德}。总而言之，这类名词不会扩展到本性中共同之物；而是从一般观念中分离出某些限定的概念，并通过它们的名称来表达。假设两人或多人放在一起，例如\Person{保禄}、\Person{息耳瓦诺}和\Person{弟茂德}，并对人的本质或实体进行探究；没有人会为\Person{保禄}给出一个本质或实体的定义，而为\Person{息耳瓦诺}给出第二个，为\Person{弟茂德}给出第三个；相反，用以阐述\Person{保禄}本质或实体的同一词语也适用于其他二人。那些由同一本质或实体的定义所描述的人，具有相同的本质或实体。当探究者了解了共同之物，转而注意那些使一人区别于另一人的区别特性时，用来认识每个人的定义将不再与另一个人的定义在所有细节上一致，尽管在某些点上可能发现一致。\par

3. 我的主张如下：凡是特殊而独特地谈论的对象，由位格之名指示。假设我们说\Quote{一个人}。该词不确定的意义在耳中留下某种模糊的印象。本性质被指出，但那个实际存在且由名称特别而独特地指示的对象并不明确。假设我们说\Quote{\Person{保禄}}。我们便通过名称所指之物，显明了实际存在的本性。\par

这就是位格，或者说\Quote{\Emph{理解}}；它不是本质或实体的不确定概念（因为所表示的是一般性的，所以没有\Quote{\Emph{立足点}}），而是通过所表达的特性给那一般而无限制的概念以\Emph{立足点}和限定。在圣经中，这样的区分是常见的，在\BibleBook{约伯}{传}的历史以及许多其他段落中皆是如此。当\Person{约伯}打算叙述他的生平事迹时，他先提到共同点，说\Quote{一个人}；然后立即通过加上\Quote{某一个}加以具体化。关于本质的描述，因与其著作的范围无关，他保持沉默；但通过特殊的身份标记，提及地点、性格特点以及那些能个体化并从共同一般观念中分离出来的外在条件，他具体指明那\Quote{某一个人}，以便从姓名、地点、心理品质和外部环境各方面，使所述之人的生平描述完全清晰。如果他当时是在阐述本质，那么在解释本性时便不会提及这些事项。同样，对于\Person{叔亚人彼耳达得}、\Person{纳阿玛人左法尔}以及那里所提到的每一个人，其描述也是如此。\BibleRef{约 2:11} 那么，请将你在人间事物中认识的本质与位格方面的区别标准，同样运用于神圣教义，必不会出错。你对圣父存在方式所想到的一切，也当同样思想圣子，同样也思想圣神。因为试图以任何孤立的概念来限制超乎一切概念的理智是徒劳的。因为圣父、圣子、圣神方面，关于非受造和不可理解的描述是同一的。没有一位比另一位更不可理解和非受造。既然必须藉着区别标记，在圣三中保持区别而不混淆，我们就不应考虑那些共同观察到的特性（如非受造、超乎一切理解或任何此类性质）来评估区别；我们应只关注如何使每个特定概念与共同概念清晰而明确地区分开来。\par

4. 依我之见，引领我们探索的正确方式如下。我们说，凡因天主的眷顾临于我们的一切美善，都是那在我们内运行万事的恩宠的运作，正如宗徒所说：\BibleQuote{但这都是同一圣神所运行，随他的心愿个别分配给各人。}\BibleRef{格前 12:11} 如果我们问，如此降临于圣人们的美善供应是否只源于圣神，我们另一方面受圣经引导而相信，那藉圣神在我们内所成就的美善供应，其创始者和原因是独生天主；因为圣经教导我们：\BibleQuote{万物是藉着他造的}\BibleRef{若 1:3}，以及\BibleQuote{万有赖他而存在}\BibleRef{哥 1:17}。当我们被提升到这个概念时，再次受天主启示的引导，我们被教导：万物由那能力从无到有被带出，但并非排除原始而藉那能力。另一方面，有一种能力无需受生、无需起源而自存，它是万有之原因的原因。因为那藉以创造万有、且圣神与祂不可分割地共想的圣子，出自圣父。因为任何人若未预先受圣神光照，就不可能想及圣子。所以，既然圣神——一切受造美善的供应源——与圣子相连，并与祂不可分割地被领悟，且其存在依附于圣父作为原因，祂也由圣父所发；祂有其特殊位格本性的这一特征：即在圣子之后并与圣子一同被认识，且祂的自立体出自圣父。圣子宣告圣神藉祂自己并与祂自己一起由圣父发出，从非受生之光单独且以独生方式照耀出来；就特殊特征而言，祂与圣父或圣神没有任何共同之处。惟独祂由所述标记而被认识。但那超越万有的天主，独自有其位格的一个特殊标志：即祂是父，且祂的位格不由任何原因而来；藉此标志，祂被特异地认识。因此，在本体的共融中，我们认为在圣三内所察觉的那些指示标记没有相互接近或互通，这些标记显示了信仰所传授的诸位格的固有特性，每一位格由其自身标记而独特地被认识。因此，按照所述指示标记，我们发现了位格的分离；然而就无限、不可理解、非受造、不受限制及类似属性而言，在赋予生命的本性中并无变化；我指的是圣父、圣子、圣神的本性，但在祂们中看到一种不可分解、持续不断的共融。藉同样的考虑，一位善于反思的学生若能察觉圣三中任何一位格的伟大，他将毫无差别地继续前行。当以心神注视圣父、圣子、圣神中的荣耀时，始终未发现任何虚空间隔可在祂们之间穿行，因为祂们之间没有任何植入物；在神性之外也没有任何如此自存之物，能以任何外来之物的介入使那本性自自分隔。也没有任何缺少自存的虚空间隔，能在神性本质的相互和谐中造成断裂，或以空无的插入破坏连续性。凡知觉圣父、且凭自身知觉祂的，同时也在心智上知觉圣子；凡接受圣子的，并不将祂与圣神分开，而是在秩序上的承继、在本性上的结合中，在自身内将三者混合在一起的信仰表达出来。凡独提圣神的，在其宣认中也包含了祂是那一位的神。既然圣神是基督的，也是天主的，如\Person{保禄}所说，那么，正如抓住链子一端的人将另一端拉向自己，凡\Quote{吸引圣神}的人，如先知所说，也藉着祂同时将圣子和圣父吸引到自己这里。如果有人切实接受了圣子，他将从两面抓住祂：圣子一面将自己的父拉向他，另一面将自己的圣神拉向他。因为那永恒存在于圣父内的，绝不可能与圣父分离；那藉圣神成就万事的，也绝不可能与自己的圣神分离。同样，凡接受圣父的，实际上同时接受了圣子和圣神；因为绝不可能持有分离或分割的观念，以至于认为圣子与圣父分离，或圣神与圣子分离。但在祂们内所察觉的共融与区别，在某种意义上是不可言说、不可思议的：本性的连续性绝不为位格的区别所撕裂，恰当区别的特征也不在本质的共通中被混淆。所以，不要惊奇于我说同一事物既是结合又是分开，并如谜语般隐晦地思想一种新奇而陌生的结合的分隔与分隔的结合。实际上，甚至在感官可察觉的对象中，任何以坦诚无争之心探讨此主题的人，都可能发现事物的类似情况。\par

5. 然而，请将我的话当作对真理的一个标记和反射，而非真理本身。因为标记中所见与采用标记所指向的对象之间，不可能完全一致。那么，我为何说在可感对象中能找到分离与结合的类比？你曾在春天见过云中虹的亮光；我是指那通常被称为虹（Iris）的弓，据熟悉此道者说，它是在一定湿气与空气混合，风力将已变为云雾的蒸气中浓密潮湿的部分挤压成雨时形成的。据说虹的形成如下：当太阳光线斜穿云层中浓密阴暗的部分，将其圆盘直射在某朵云上时，光辉便从潮湿发光处反射回来，结果便像光自身弯曲和返回。因为火焰般的闪光具有这样的性质：若落在任何光滑表面，它们会折射回自身；而太阳的形状，借助光线在潮湿光滑的空气部分上形成，是圆的。因此，必然的结果是，云旁空气凭借辐射光辉，按照太阳圆盘的形状被标出。这光辉既是连续的又是可分的。它颜色众多，形态众多；它不知不觉地浸染在其染料的绚丽明亮色调中；它不知不觉地使我们视觉中众多颜色的组合变得抽象，以至于在蓝色与火焰色之间、火焰色与红色之间、红色与琥珀色之间，都无法辨别任何空间——这空间在自身内混合或分离颜色的差异。因为所有的光线同时可见，都远亮，且它们既无相互组合的迹象，又无法检验，以致不可能发现火焰色或翠绿色光部分的界限，以及每一色在显现其光辉之前源于何处。正如在标记中我们清楚分辨颜色的差异，却无法凭感官把握它们之间的间隔；同样，我请你如此推论，以便推理有关神圣教义：\OriginalTerm{位格}{hypostases}的特性，就像虹中所见的颜色，在我们相信存在于天主圣三中的每一位上闪耀其光辉；但在固有本性上，无法设想一位与另一位之间存在任何差异，因为独特的特性在\OriginalTerm{本质}{essence}的共融中照于每位。就连在我们的例子中，发出多彩光辉并被太阳光线折射的那个本质，是一个本质；多形的是现象的色彩。我的论证就这样教导我们，甚至借助可见受造物，每当遇到难以解决的问题，以及一想到接受所提出的东西就头脑发晕时，也不要在教义要点上感到苦恼。在可见物方面，经验似乎优于因果理论；同样，在超越一切知识的事物中，论证的理解逊于那同时教导我们\OriginalTerm{位格}{hypostasis}的区别与本质的结合的信德。既然我们的讨论已包括天主圣三中共同的和区分的，那么共同的是指本质；\OriginalTerm{位格}{hypostasis}则是若干区分的标记。\par

6. 然而，或许会有人以为，这里对\OriginalTerm{位格}{hypostasis}的说明与宗徒论及主时的话意义不符，他说：\BibleQuote{他是天主光荣的光辉，是他位格的真像，}\BibleRef{希 1:3}因为我们若教导\OriginalTerm{位格}{hypostasis}是若干特性的汇合，并且既然承认，就如圣父方面可默观到某种专属的、特殊的、唯独凭之认出的东西，同样也相信关于\OriginalTerm{独生子}{Only-begotten}的事；那么，圣经何以在此处将\OriginalTerm{位格}{hypostasis}之名只归于圣父，而将圣子描述为那\OriginalTerm{位格}{hypostasis}的形象，且不是凭他自己的专有标识，而是凭圣父的标识？因为若\OriginalTerm{位格}{hypostasis}是若干存在的标记，而圣父的特性仅限于\OriginalTerm{无始}{unbegotten}的存在，圣子是按圣父的特性塑造的，那么\Quote{无始}一词便不能再独有地用于圣父，因为\OriginalTerm{独生子}{Only-begotten}的存在是由圣父的特性标记所标示的。\par

7. 然而，我的意见是，在此段落中，宗徒的论证指向一个不同的目的；正是着眼于这一点，他使用了\Quote{荣耀的光辉}和\Quote{位格的真像}。任何人若仔细留意这一点，便会发现这与我所言毫无冲突，而论证是在一种特殊且独特的意义上进行的。因为宗徒论证的目的并非藉由明显的标记来区分\OriginalTerm{位格}{hypostasis}；而是领悟圣子与圣父之间本然的、不可分离的密切关系。他没有说\Quote{他是圣父的荣耀}（尽管事实上他就是）；他将这一点视为当然而略过，然后努力教导我们，不可认为在圣父那里有一种荣耀形式，在圣子那里有另一种；他将独生子的荣耀定义为圣父荣耀的光辉，并且通过光的比喻，使人将圣子与圣父视为不可分离的结合。因为正如火焰发出光辉，光辉并非在火焰之后，而是在同一瞬间火焰照耀、光明闪耀；同样，宗徒意指圣子从圣父获得存在，且独生子并非因任何空间的间隔而与圣父的存在分离，而是受生者始终与原因一同被思考。完全同样，他如同对前述原因的意义进行解释，并以物质例证引导我们理解不可见者，他也提到了\Quote{位格的真像}。因为身体完全在于形相，但身体的定义与形相的定义不同，没有人希望给出其中一个的定义会发现与另一个一致；然而，即使理论上你将形相与身体分开，本性却不允许区分，二者不可分离地被理解；同样，宗徒认为，即使信仰的道理将位格的区别表现为无混且分明，他也必须通过他的语言来呈现独生子与圣父之间的连续以及某种具体的关系。他这样说，并非似乎独生子没有位格性的存在，而是因为合一不允许圣子与圣父之间有任何间隔，结果，凡用灵魂之眼专注凝视独生子真像的人，也会感知到圣父的位格。然而，在它们之中所观察到的固有性质既不会改变，也不会混杂，以至于我们不应将受生的起源归于圣父，或将不受生的起源归于圣子；而是如此：即使我们能够克服将两者分开的不可能性，使一方能被单独理解，因为单单名字就隐含了圣父，任何人若不理解圣父，就不可能甚至称呼圣子。\par

8. 因此，正如主在福音中所说，\BibleRef{若 14:9}看见圣子的，也看见圣父；为此他说，独生子是他父位格的真像。为使此事更明了，我还要引用宗徒另外的段落，其中他称圣子为\Quote{不可见天主的肖像}，\BibleRef{哥 1:15}以及\Quote{他良善的肖像}。并非因为肖像与其原型在不可分性和良善的定义上有所不同，而是为了表明，即使肖像不同，它也与原型相同。因为肖像的概念若未完全保持清晰不变的相似性，就会丧失。因此，感知肖像之美的人，也得以感知原型。所以，凡是对圣子的形相有某种心灵领悟的人，就印下了圣父位格的真像，在圣子中看见圣父，不是在反映中看见圣父的非受生之存在（因为那样就会完全同一，毫无区别），而是凝视在受生者中的非受生之美。正如一个人在光滑的镜子中看到形相的反射，如同对所表现之面容的清晰认识；同样，认识圣子的人，透过对圣子的认识，在心中接受圣父位格的真像。因为凡属于圣父的，都在圣子中看见；凡属于圣子的，都属于圣父；因为整个圣子在圣父内，而圣父全部在圣子内。因此，圣子的位格成为认识圣父的形相与面貌，而圣父的位格在圣子的形相中被认识，而在其中所观察到的固有性质，则保留为对位格的清晰区分。\par

\SourceRecord{Translated by Blomfield Jackson.；Nicene and Post-Nicene Fathers, Second Series；Vol. 8.；Edited by Philip Schaff and Henry Wace.；Buffalo, NY: Christian Literature Publishing Co.,；1895.；<http://www.newadvent.org/fathers/3202038.htm>.}

% structure: p0001 source=h1 target=section reason=page-title

\section{第三十九书}

\Signature{凯撒勒雅的圣巴西略}\par

\Person{儒利安}致\Person{巴西略}。\par

谚语说：\Quote{你不是在宣告战争，}并且，容我加上，出自喜剧：\Quote{金言使者啊。}那么，请用行动证明这话，并速来见我。你将作为朋友来访朋友。对将事务视为次要之人，显眼且不懈的投入似乎是重担；然而勤奋者谦逊，如我所信，明智且随时准备应对任何紧急状况。我允许自己放松，以便一个不疏忽任何事的人也能得享安息。我们的生活方式并无宫廷虚伪——我想你对此已有一些经验——按照那种虚伪，恭维意味着比我们对最坏敌人更致命的仇恨；但以适当的自由，我们在应受责备时责备并斥责，却以最亲密朋友的爱去爱。因此，我可以说，以全然的真诚，既能勤于放松，又能工作时不觉疲惫，并安睡无忧；因为我醒时不仅为自己，也恰当地为每一个人警醒。恐怕这有些愚蠢和琐碎，因为我觉得有些懒惰（我像\Person{阿斯提达马斯}一样自我夸赞），但我写信是要向你证明，能见到你这位智者，对我而言，好处多于麻烦。因此，如我所说，勿失时机：火速赶来。在你如愿与我共度一段长时间后，你可继续你的旅程，无论去向何方，我都致以最良好的祝愿。\par

\SourceRecord{Translated by Blomfield Jackson.；Nicene and Post-Nicene Fathers, Second Series；Vol. 8.；Edited by Philip Schaff and Henry Wace.；Buffalo, NY: Christian Literature Publishing Co.,；1895.；<http://www.newadvent.org/fathers/3202039.htm>.}

% structure: p0001 source=h1 target=section reason=page-title

\section{第四十封书信}

\Emph{\Person{圣巴西略·凯撒勒雅}}\par

\Emph{儒里安致巴西略。}\par

迄今为止，我自幼便自然具有的温和与仁慈一直显现出来，我已使所有居于太阳之下的人都归服于我。因为，看！每一个野蛮部落，直到大洋之滨，都来将他们的礼物置于我的脚下。同样，居住在多瑙河彼岸的萨加达雷斯人，他们身上有奇异的刺青，几乎不像人类，如此野蛮和怪异，现在也匍匐在我脚下，并宣誓服从我的主权所加给他们的一切命令。我还有进一步的目标。我必须尽快出兵波斯，击败那个大流士的后裔沙普尔，使他成为进贡者。我也打算蹂躏印度人和撒拉逊人的国土，直到他们都承认我的优越并成为我的进贡者。然而，你却自称拥有超越这一切的智慧；你自称披着虔诚的外衣，但你的衣服实际上是厚颜无耻，到处诽谤我，说我不配皇帝的尊严。难道你不知道我是杰出的君士坦提乌斯的孙子吗？我知道你如此，但我并未改变从前我们两人年轻时所持有的对你的情感，以及你对我的情感。然而，出于我的仁慈之意，我命令你，当我经过凯撒勒雅时，从你那里送给我一千磅黄金；因为我仍在行军，正在以最快的速度赶去波斯战役。如果你拒绝，我已准备好摧毁凯撒勒雅，推翻那些长久以来装饰它的建筑；在它们的原址建立庙宇和雕像；并以此迫使所有人服从罗马皇帝，而不自高自大。因此，我责令你务必经过一位可靠的使者，在适当点数和称量之后，用你的戒指封好，将规定的黄金如数送到我手中。这样，我或许会因你的错误而向你显示仁慈，只要你承认，尽管为时已晚，任何借口都无济于事。我已经学会认识并谴责我从前所读到的一切。\par

\SourceRecord{Translated by Blomfield Jackson.；Nicene and Post-Nicene Fathers, Second Series；Vol. 8.；Edited by Philip Schaff and Henry Wace.；Buffalo, NY: Christian Literature Publishing Co.,；1895.；<http://www.newadvent.org/fathers/3202040.htm>.}

% structure: p0001 source=h1 target=section reason=page-title

\section{第四十一书}

\Person{圣巴西略} \Place{凯撒勒雅}\par

\Emph{巴西略致尤利安书}。\par

一、尊驾目前的壮举微不足道，您对我的猛烈攻击，或不如说是对您自己的攻击，也微不足道。当我想到您身披紫袍，戴着冠冕在您那蒙羞的头上——由于真宗教的缺席，这王冠与其说荣耀您的帝国，不如说使其蒙羞——我便颤抖不已。而您本人，虽已变得如此崇高伟大，但如今卑鄙且憎恨荣誉的恶魔将您引至此境，您不仅开始超乎一切人性之上，甚至起来反对天主，侮辱祂的教会——万有的母亲与养育者——竟派遣最微不足道的我，命令我向您进献一千磅黄金。我并非因黄金的重量而如此震惊，尽管它确实沉重；但它使我为您如此迅速的堕落流下苦泪。我想到您我一同学习最优良、最神圣之书的功课。我们各自通过了神圣而受天主启示的圣经。那时，于您一无所隐。如今，您已迷失了正当情感，被傲慢围困。您的殿下并非昨日才初次发现我并不活在极度富足之中。今日您却向我索要一千磅黄金。但愿殿下屈尊宽免我。我的财产之微，实在连今日想吃个够也不够。在我屋檐下，厨艺已绝迹。我仆人的刀从不沾血。最重要的菜肴——我们丰盛之所在——是野菜叶搭配极粗的面包和酸酒，这样我们的感官不致因贪食而迟钝，也不至于纵情过度。\par

二、您出色的护民官劳苏斯，您命令的忠实执行者，也向我报告说，有一妇人因儿子被毒杀而向殿下求诉；您断定毒杀者不得容留；若有，则当处死，或只留下那些与野兽搏斗的人。您如此判断正确，但在我看来却奇怪，因为您试图用微小的补救治疗大创伤的痛苦，这荒谬至极。在侮辱天主之后，您再关心寡妇孤儿也无用。前者是疯狂危险的；后者是仁慈善良之人的作为。像我这样一个平民对皇帝说话是件严重的事；您对天主说话将更为严重。没有人在天主与人之间做中介。您所读的，您并不明白。若您明白了，就不会定罪。\par

\SourceRecord{Translated by Blomfield Jackson.；Nicene and Post-Nicene Fathers, Second Series；Vol. 8.；Edited by Philip Schaff and Henry Wace.；Buffalo, NY: Christian Literature Publishing Co.,；1895.；<http://www.newadvent.org/fathers/3202041.htm>.}

% structure: p0001 source=h1 target=section reason=page-title

\section{第四十二书}

\Person{圣巴西略}\par

\Emph{致其门徒奇洛}。\par

我真正的弟兄，若你欣然接受我对你当行之路的劝告，尤其是在你向我请示的几点上，你将把救恩归功于我。许多人曾有勇气度隐修生活，但能坚持到底者，或许寥寥无几。终点并非必然包含在初衷中，然而辛苦的酬报在于终点。因此，那些未能坚持到底、只是开始时过隐修生活的人，得不到任何益处。反倒使他们的宣认显得可笑，并被外界指责为缺乏勇气和意志不坚。此外，主论及这样的人说，想要建造房屋的，\Quote{岂不先坐下计算费用，是否够完成？} 恐怕安了地基，却不能完成，\Quote{旁观者就讥笑他说，} 这个人安了地基，\Quote{却不能完成。} 因此，让这开端意味着你诚心在德行上进步。高贵的战士保禄，不愿我们安于过去行善的日子，而要每日取得更大进步，说：\Quote{忘记背后，努力面前的，向着标竿直跑，要得天主召我来得的奖赏。}\BibleRef{斐 3:13} 整个人生确实如此，不满足于既往，与其说靠过去滋养，不如说是靠未来。因为一个人昨日果腹，今日若饥饿未能从食物中获得适当满足，于他有何益处？同样，灵魂不能从昨日的德行获益，除非有今日的正直行为相随。因为经上说：\Quote{我按我所见判断你。}\par

2. 如此，义人的劳苦便是枉然，罪人的道路也无可指责，若发生转变，前者从善变恶，后者从恶变善。我们从厄则克耳听见他仿若以主的名教导说：\Quote{若义人转离他的义而行恶，我不再纪念他先前所行的义；他必因他的罪而死。} 对罪人也是如此：若他转离他的恶，行正直的事，他必存活。天主的仆人梅瑟的一切辛劳何在？因一刻的悖逆而被禁止进入应许之地。革哈齐与厄里叟的同伴关系后来如何？因贪婪而自染麻风病。撒罗满的广博智慧和他先前对天主的尊敬有何益处？后来因疯狂爱恋女子而陷入偶像崇拜。就连有福的达味也非无可指责，当他的心念偏离，对乌黎雅的妻子犯了罪。一个例子足以警告度敬虔生活的人：犹达斯从善堕落到恶，他作基督门徒如此之久，却为微利出卖了师傅，自己得了绳索。弟兄，要晓得，开始得好的人并非成全者，而是在天主眼中终结得好的人，才蒙悦纳。因此，不要容你的眼睛睡觉，也不容你的眼皮打盹，好使你得以逃脱，\Quote{如麋鹿脱离网罗，如飞鸟脱离捕鸟人的圈套。} 看，你正行经网罗之中；你正走在高墙之顶，跌落对跌落者极其危险；因此不要立即尝试极端的操练；尤其要谨慎自信，免得因缺乏训练而从操练的高处跌落。最好循序渐进。逐步从生活的享乐中退出，渐渐破坏你所有的旧习惯，免得同时激发所有情欲而招来众多诱惑。当你制伏了一个情欲，再开始向另一个宣战，如此你终将胜过一切。就名称而言，放纵只有一个，但其实际运作各不相同。首先，弟兄，以坚忍迎接每一个诱惑。信德之人通过多种诱惑被考验：世俗损失、控告、谎言、反对、诽谤、迫害！这些和类似的是对信德之人的考验。此外，要安静，言语不轻率，不争辩，不争论，不贪图虚浮的荣耀，不急于获取知识而非给予，不多言，总愿学习多于施教。不要为世俗生活忧虑；从中你得不到任何益处。经上说：\Quote{愿我的口不说人的行为。} 喜欢谈论罪人行为的人，很快激起自我放纵的欲望；最好专注于善人的生平，这样你将为自己获得益处。不要着急从村庄到村庄、从房屋到房屋四处游走；反倒要避开它们，如同灵魂的陷阱。倘若有人出于真正的怜悯，多次恳请你进入他家，要告诉他效法那位百夫长的信德：当时耶稣急忙要去他那里治病，他却恳求耶稣不要这样做，说：\Quote{主啊，我不配你到我的屋顶下来，只要你说一句话，我的仆人就会痊愈。}\BibleRef{玛 8:8} 耶稣对他说：\Quote{去吧，照你所信的，给你成就。}\BibleRef{玛 8:13} 他的仆人就在那一刻痊愈了。弟兄，要晓得，是求告者的信德，而非基督的临在，使病人得救。现在也同样，无论你在何处祈祷，若病人相信你的祈祷会帮助他，一切将如愿成就。\par

3. 你爱亲属不可胜过爱主。祂说：\Quote{谁爱}父亲、母亲或兄弟\Quote{超过我，不配是我的。}主的诫命有何含义？\Quote{谁不背起自己的十字架跟随我，也不能作我的门徒。}如果你同基督一起、按肉身已死于你的亲属，为什么还想同他们再活？若你为亲属的缘故，重建那为基督的缘故所拆毁的，你就使自己成了犯法者。那么，不要为亲属的缘故放弃你的地方：你若放弃你的地方，或许你也会放弃你的生活方式。不要爱人群，不要爱乡村，也不要爱城市；要爱旷野，常独处自守，心无游荡，视祈祷与赞颂为你一生的事业。切勿忽视阅读，尤其是新约圣经，因为阅读旧约时常会导致祸害；并非所写的有害，而是受伤者的心灵软弱。一切面包都有营养，但对病人却可能有害。同样，全部圣经都是天主所默感、有益的，其中没有不洁之物：惟有那以为不洁的人，在他才是不洁。\Quote{凡事要考验；好的，要持守；各样恶事要远避。}\Quote{凡事都可行，但不全有益。}\BibleRef{格前 6:12}在你所接触的众人中，凡事总不要使人绊倒，要喜乐，\Quote{有手足之情，}\BibleRef{伯前 3:8}和蔼可亲，谦卑自下，不可因丰盛的肴馔而失却款待的准绳，常以手中之物为足。不要从任何人那里收取超出独修生活日常所需之物。尤其要远避金钱，视之为灵魂之敌、罪恶之父、魔鬼的工具。不要假借周济穷人的名义，使自己受贪婪的指控；若有人为你带来给穷人的钱，而你知道有谁需要，就当劝那主人亲自将钱送去给他急需的弟兄，免得你因接受金钱而玷污了良心。\par

4. 要远避享乐；追求节制；锻炼身体适应辛劳；使灵魂习惯试探。视灵魂与身体的分离为脱离一切邪恶的解脱，期待所有圣人共同分享的永恒美善的享受。要时刻如同持天平对抗魔鬼的每一个建议，投入一个神圣的念头，随着天平的倾向，你就随它而去。尤其是当恶念升起，说：\Quote{你在此地度过一生有何益处？你退隐人群得到了什么？难道你不知道，那些由天主祝圣为天主教会主教的人们，常与同伴交往，不知疲倦地参加灵修聚会，与会者从中获益甚多？在那里可以享受难解之言的讲解、宗徒教训的阐述、福音思想的诠释、神学的课程以及灵性弟兄的交往，他们仅凭面容的显现就对所遇之人大有裨益。而你，却决定与这一切美善隔绝，像野兽一样坐在这荒野之中。你环顾四周，是一望无际的旷野，几乎没有一个同类；缺乏一切训导，与弟兄疏远，你的灵魂在遵行天主的诫命上懈怠无力。}当恶念用所有这些巧妙的借口攻击你，想要毁灭你时，你要用敬虔的反思和你自己的实际经验来对抗它，说：你告诉我世上的事物是好的；我来这里的原因，是因为我判断自己不配拥有世上的美善。世上的美善中掺杂着邪恶，而邪恶明显占上风。有一次我参加灵修聚会，好不容易找到一个在我看来敬畏天主的弟兄，但他却是魔鬼的猎物，我从他口中听到有趣的轶事和编造出来欺骗遇见者的故事。之后我又遇见许多强盗、掠夺者、暴君。我看见可耻的醉汉；我看见被压迫者的血；我看见女人的美貌，折磨着我的贞洁。我逃避了实际的淫乱，却以心中的思念玷污了我的童贞。我听到许多对灵魂有益的讲道，但在任何一位教师身上，我都无法发现他的生活与他的言词相称。此后，我又听到大量的戏剧，以放荡的歌曲引人入胜。然后我听到竖琴的悦耳演奏，杂技演员的掌声，小丑的谈话，各样的玩笑和蠢行，以及人群的喧嚣。我看见被抢者的眼泪，受暴政者的痛苦，受折磨者的哀号。我定睛一看，啊，那里没有灵修聚会，只有一片大海，被风吹动翻腾，试图用它的波浪一下子淹没每一个人。告诉我，恶念！告诉我，短暂享乐与虚荣的恶魔！当我无力救助任何受这些苦害的人时，当我既不被允许保护无助者，也不被允许纠正堕落者时，或许我还可能毁掉自己——我看到、听到这一切，有什么益处呢？因为正如一股风沙尘土就能吹走极少量的淡水，同样，我们在今世认为所做的善行，也被众多的邪恶所淹没。今世为人演出的戏剧，借着喜乐和欢娱，如同木桩一般打入他们的心中，使他们敬拜的光辉暗淡下来。而受同僚冤屈者的哀号与悲叹，却被引入来显示穷人的忍耐。\par

5. 那么，除了丧失我的灵魂，我还能得到什么好处呢？为此，我像飞鸟一样逃往山中。\Quote{我好像飞鸟，从捕鸟人的网罗里逃脱。} \newline{} 邪恶的念头啊，我正生活在主曾生活过的旷野里。这里有玛默勒的橡树；这里有雅各伯所看见的通往天堂的梯子和天使的堡垒；这里有洁净了的百姓接受法律、从而进入预许之地并看见天主的旷野。这里有厄里亚曾旅居并悦乐天主的加尔默耳山。这里有厄斯德拉退隐的平原，他按天主的命令说出了所有受天主默感的书。这里有真福若翰吃蝗虫并向人宣讲悔改的旷野。这里有橄榄山，基督曾到此祈祷，并教导我们祈祷。这里是喜爱旷野的基督，因为祂说：\Quote{哪里有两个或三个人，因我的名字聚在一起，我就在他们中间。} \Quote{导入生命的门是窄的，路是狭的。} \BibleRef{玛 7:14} 这里有\Quote{在旷野、山岭、山洞和地穴中漂流无定}\BibleRef{希 11:38} 的导师和先知。这里有宗徒、传福音者和远离城市的独修生活。我全心拥抱这一切，为赢得应许给基督的殉道者及祂所有其他圣人的一切，从而我能真正说：\Quote{因了你口中的言语，我遵行了艰难的道路。} 我听说，天主的朋友亚巴郎听从神圣的声音，去了旷野；依撒格服从权威；圣祖雅各伯离开家园；贞洁的若瑟被出卖；三个青年学会守斋，并与烈火搏斗；达尼尔两次被投入狮子圈；耶肋米亚放胆直言，被扔进污泥坑；依撒意亚看见不可言喻的事，被锯子锯开；以色列人被掳走；责斥奸淫的若翰被斩首；基督的殉道者被杀害。但我何必多说？在这里，我们的救主亲自为我们被钉在十字架上，为借祂的死亡赐给我们生命，并训练和吸引我们所有人忍耐。我向祂、也向圣父和圣神奋力迈进。我努力被证明是真诚的，自认不配享有这世界的财物。然而，不是因为我而有了世界，而是因为世界而有我。你要在心里思想这一切，热忱地追随它们；要战斗，正如你所受的命令，为真理奋斗至死。因为基督\Quote{服从}至\Quote{死。} \BibleRef{斐 2:8} 宗徒说：\Quote{你们要小心，免得你们中间有人起背信的恶心……背离了生活的天主；反之，只要还有\Quote{今天}之称，你们要天天互相劝勉……（并彼此建树\BibleRef{得前 5:11}）。} \BibleRef{希 3:12} \Quote{今天}就是指我们整个生命的时间。这样生活，兄弟，你将拯救自己，你将使我喜乐，你将光荣天主，从永远到永远。阿们。\par

\SourceRecord{Translated by Blomfield Jackson.；Nicene and Post-Nicene Fathers, Second Series；Vol. 8.；Edited by Philip Schaff and Henry Wace.；Buffalo, NY: Christian Literature Publishing Co.,；1895.；<http://www.newadvent.org/fathers/3202042.htm>.}

% structure: p0001 source=h1 target=section reason=page-title

\section{第四十三封书信}

\Person{圣巴西略} \Place{凯撒勒雅}\par

\WorkTitle{劝青年书}\par

有信德的独修者，实践真宗教的人，请学习福音的生活方式：克制身体、温良的精神、心灵的纯洁、消灭骄傲。若被强迫服役，为主还要加增礼物；被抢夺，绝不诉讼；被憎恨，却要爱；被迫害，却要忍耐；被毁谤，却要恳求。死于罪恶，被钉死于天主。将一切挂虑托付于主，好使你被发现于千万天使、首生者的集会、先知们的宝座、圣祖们的权杖、殉道者的冠冕、义人的赞美之中。切望在基督耶稣我们的主内，与那些义人一同被数算。愿光荣归于祂，直到永远。阿们。\par

\SourceRecord{Translated by Blomfield Jackson.；Nicene and Post-Nicene Fathers, Second Series；Vol. 8.；Edited by Philip Schaff and Henry Wace.；Buffalo, NY: Christian Literature Publishing Co.,；1895.；<http://www.newadvent.org/fathers/3202043.htm>.}

% structure: p0001 source=h1 target=section reason=page-title

\section{书信 44}

\Emph{\Person{凯撒勒雅的圣巴西略}}\par

\Emph{致一位失足的隐修士}。\par

1. 我不愿你欢喜，因为恶人没有欢喜。即使现在我也无法相信；我的心灵无法想象如此重大的不义，即你所犯的罪行；如果，也就是说，事实真的是人们普遍理解的那样。我茫然不解，如此深奥的智慧如何能被消除；如此严谨的纪律如何能被废弃；何处的如此深沉盲昧能围绕你；你如何以极其轻率的态度造成如此灵魂的毁灭。如果这是真的，你已将自己的灵魂交付深渊，并松懈了所有听见你不敬之事之人的热忱。你藐视了信德；你错过了光荣的战斗。我为你悲伤。哪一位神职人员听闻后不哀叹？哪一位教会人士不捶胸？哪一位平信徒不沮丧？哪一位苦修者不悲伤？或许，连太阳也因你的堕落而昏暗，天上的权势因你的毁灭而震动。甚至无知的石头也为你的疯狂流泪；甚至你的敌人也因你极大的不义而哭泣。哦，心硬！哦，残忍！你不敬畏天主；你不尊重世人；你毫不关心你的朋友；你一朝使一切沉没；你一朝失去了一切。我再次为你悲伤，不幸的人。你曾向众人宣讲天国的权能，而你却从中堕落了。你曾使众人因你的教导而战栗，而你眼前却没有敬畏天主。你曾宣讲纯洁，而你却被发现污秽了。你曾以你的贫穷自夸，而你却被证实贪得无厌；你曾展示并解释天主的惩罚，而你自己却将惩罚引到了自己头上。我该如何哀悼你，如何为你悲伤？早晨升起的晨星路济弗尔是如何坠落、跌碎在地呢？凡听见的人，耳朵都要发鸣。比黄金更明亮的纳齐尔人，如何变得比沥青更黯淡？\Place{熙雍}光荣的儿子如何成了无用的器皿！至于他，对圣经的记忆曾在众人口中，如今他的记忆随着声音消失了。那机敏聪慧的人迅速消逝了。那多才多智的人做了多重的恶事。所有曾因你的教导受益的人，都因你的堕落受了伤害。所有曾来聆听你言谈的人，都因你的堕落掩耳不听。我，忧伤沮丧，各方面软弱，以灰为食，以麻布裹伤，就这样数说你的赞美；或者更确切地说，无人安慰，无人医治，我正撰写墓志铭。因为安慰从我眼前隐藏了。我没有药膏，没有油，没有绷带可敷。我的创伤疼痛，我如何能痊愈？\par

2. 如果你对救恩还有任何希望；如果你对天主有丝毫思念，或对将来的福乐有任何渴望；如果你对为不悔改者保留的惩罚有任何恐惧，就请立刻醒来，举目向天，清醒过来，停止你的恶行，摆脱笼罩你的昏沉，对抗击败你的仇敌。努力从地上起来。记住那将要跟随并拯救你的善牧。即使只剩下两条腿或一个耳垂，也要从伤你的野兽那里跳回。记住天主的仁慈，以及祂如何用油和酒医治。不要对救恩绝望。回想你记忆中所记载的圣经：跌倒的必起来，转离的必回转；受伤的必得医治，被野兽掠取的必逃脱；承认自己罪过的必不被拒绝。主不愿罪人死亡，而愿他悔改而存活。不要像邪恶深坑中的恶人一样轻蔑。有忍耐的时候，有长期受苦的时候，有医治的时候，有改正的时候。你跌倒了？起来。你犯了罪？停止。不要站在罪人的道路上，而要跳开。当你悔改并叹息时，你必得救。劳苦带来健康，汗水带来救恩。当心，免得因你希望持守某些义务，而破坏你在许多证人面前向天主宣认的义务。我恳求你，不要因任何尘世的顾虑而犹豫前来见我。当我找回我的死者时，我要哀悼，我要照料他，我要哭泣，\BibleQuote{因为我人民的女儿被掠夺。}\BibleRef{依 22:4}众人都准备好欢迎你，众人都会分担你的努力。不要退缩。回忆昔日时光。有救恩；有改正。鼓起勇气；不要绝望。这不是一条毫不怜悯地定罪至死的法律，而是慈悲，免除惩罚并等候改善。门尚未关闭；新郎仍在倾听；罪不是主人。再努力一次，不要犹豫，可怜你自己和我们众人，靠我们的主耶稣基督，愿光荣和权能归于祂，从今世，直到永远。阿们。\par

\SourceRecord{Translated by Blomfield Jackson.；Nicene and Post-Nicene Fathers, Second Series；Vol. 8.；Edited by Philip Schaff and Henry Wace.；Buffalo, NY: Christian Literature Publishing Co.,；1895.；<http://www.newadvent.org/fathers/3202044.htm>.}

% structure: p0001 source=h1 target=section reason=page-title

\section{\WorkTitle{书信第四十五}}

\Person{圣巴西略}\par

\Emph{致一位失足的隐修士}\par

1. 我内心深感双重惊骇，而你是缘由。要么我成了某种不友善的成见的牺牲品，被迫提出不友爱的指控；要么，尽管我极愿同情你，温和地对待你的困扰，却被迫采取一种不同且不友善的态度。因此，当我提笔书写时，我已通过反思鼓起不情愿的手；但我因你而忧伤，面色低垂，无力改变。我为你感到如此羞耻，以致我的嘴唇转为哀悼，我的口立刻垂下。哀哉！我该写什么？我困惑中该想什么？\par

若我回想你从前空虚的生活方式，当你财富滚滚、拥有许多微小的世俗声誉时，我战栗；那时你身后跟着一群谄媚者，短暂享受奢华，伴随明显的危险和不义之财；一方面，对官员的恐惧驱散了你对救恩的关心，另一方面公共事务的动荡扰乱了你的家，持续的烦恼使你的心转向那位能帮助你的天主。然后，你逐渐开始寻求救主，祂为你带来有益的恐惧，祂拯救并保护你，尽管你在安稳中嘲笑祂。那时你开始训练自己转向一种有价值的生活，将你所有危险的财产视为粪土，抛弃家庭和妻子的陪伴。你像陌生人和流浪者一样四处漂泊，穿越城镇和乡村，去了耶路撒冷。在那里，我自己曾与你同住，并为你苦行修炼的劳苦而称你有福，当你连续几周禁食，在天主面前沉思，避开同伴的交往，如同一个溃逃的逃亡者。然后你为自己安排安静独处的生活，拒绝社会的一切烦扰。你用粗糙的苦衣刺痛身体；你在腰间束紧硬皮带；你勇敢地给骨头施加压力；你使两侧从前面到后面松弛，因禁食而中空；你不愿穿柔软的绷带，将胃像葫芦一样收缩，使其紧贴肾脏周围。你清空了肉中的一切脂肪；你干涸了腹部以下的所有通道；你因缺乏食物而将腹部折叠起来；你的肋骨像房子的洞穴，遮蔽了腰部周围的部分，全身收缩，在漫长的夜间将忏悔倾泻给天主，使胡须被泪水的溪流浸湿。为何细说？还记得你亲吻了多少圣人的口，拥抱了多少身体，多少人握住你纯洁的手，多少天主的仆人如同敬拜般跑来抱住你的膝？\par

2. 这一切的结局是什么？我的耳朵被一项奸淫的指控所伤，它飞得比箭还快，以更尖锐的刺痛刺穿我的心。什么巫师的狡猾诡计将你驱入如此致命的陷阱？什么多网的魔鬼之网缠住了你，瘫痪了你德行的一切力量？你劳苦的故事成了什么？还是我们必须不信它们？当我们看到明显的事实时，怎能避免相信早已隐藏的事？我们该说什么：你以巨大的誓言束缚了那些逃向天主求助的灵魂，而比是或非更多的东西被仔细地归于魔鬼？你为致命的伪证作了担保；通过轻视苦修者的品格，你甚至把责备加于宗徒和主自己身上。你羞辱了纯洁的夸耀。你玷污了贞洁的许诺；我们成了俘虏的悲剧，我们的故事在犹太人和外邦人面前被演成一出戏。你分裂了独修士的精神，使更严格操练的人陷入恐惧和怯懦，他们仍对魔鬼的力量感到惊讶，又引诱粗心的人模仿你的不节制。尽你所能，你摧毁了基督的夸耀，祂说：\Quote{你们放心吧，我已战胜了世界，}\BibleRef{若 16:33} 和世界的王。你为你的国家调制了一碗恶名。的确，你证实了那句谚语：\Quote{如鹿被击穿肝脏。}\par

但现在呢？力量之塔并未倒塌，我的弟兄。纠正的药方并未被嘲笑；避难之城并未关闭。不要停留在邪恶的深处。不要把自己交给灵魂的凶手。主知道如何扶起那些被击倒的人。不要试图远远逃离，但要赶快到我这里来。重新开始你年轻时的劳作，通过新的善行行为，摧毁那匍匐在地面上的卑劣放纵。仰望那已如此接近我们生命的终结。看哪，现在犹太人和外邦人的子孙如何被驱使敬拜天主，不要完全否认世界的救主。绝不要让那最可怕的判决临到你：\Quote{离开我，我从来不认识你们。}\BibleRef{路 13:27}\par

\SourceRecord{Translated by Blomfield Jackson.；Nicene and Post-Nicene Fathers, Second Series；Vol. 8.；Edited by Philip Schaff and Henry Wace.；Buffalo, NY: Christian Literature Publishing Co.,；1895.；<http://www.newadvent.org/fathers/3202045.htm>.}

% structure: p0001 source=h1 target=section reason=page-title

\section{第四十六书}

\Emph{\Person{凯撒勒雅·圣巴西略}}\par

\Emph{致一位失足的贞女}\par

1. 现在正是引用先知之言的时候，要说：\BibleQuote{但愿我的头是水泉，我的眼是泪泉，好能昼夜为我人民女儿的伤亡哭泣。}\BibleRef{耶 9:1} 虽然她们深陷沉默，被自己的不幸击昏，因致命的打击失去了一切感觉，但我绝不可让这样的堕落无人哀悼。如果对耶肋米亚来说，那些在战争中身体受伤的人堪受无数的哀哭，那么对这样的灵魂灾难又该说什么呢？经上说：\BibleQuote{我的被杀者，不是被刀剑所杀，也不是战死。}\BibleRef{依 22:2} 但我哀悼的是真正死亡的刺伤，是罪的沉重和恶者燃烧的箭，它们凶残地焚烧了灵魂和身体。诚然，天主的律法看到地上如此重大的玷污，必会大声叹息。它们早已颁布禁令：\BibleQuote{不可贪恋你近人的妻子}；\BibleRef{申 5:21} 并通过神圣的福音说：\BibleQuote{凡注视妇女而有贪念的，已经心里与她犯了奸淫。}\BibleRef{玛 5:28} 如今它们看到主的净配——以基督为元首的——竟胆敢犯奸淫。同样，众圣者的团体也会叹息。热诚的丕乃哈斯，因为他再也不能执枪在手，为身体的凌辱复仇；洗者若翰，因为他不能离开上天之境——如同在世时离开旷野——赶去谴责罪恶，即使为这事必须死，他宁可失去头颅也不失去言论的自由。但是，也许像有福的亚伯尔，他虽死仍向我们说话，如今更比若翰当年论及黑落狄雅时大声疾呼：\BibleQuote{你不可占有她。}\BibleRef{玛 14:4} 因为即使若翰的身体顺从自然法则，接受了天主的判决，他的舌头沉默了，然而\BibleQuote{天主的言决不受束缚。}\BibleRef{弟后 2:9} 若翰看到同仆的婚约被藐视，敢于斥责至死；他看到如此凌辱在主的洞房上，又会作何感想呢？\par

2. 你抛弃了那神圣结合的轭；你逃离了真正君王的无玷洞房；你可耻地堕入这耻辱而不敬神的败坏之中；如今你既无法逃避这痛苦的指控，也无办法或计谋隐藏你的烦恼，你便冲入傲慢之中。恶人一旦落入罪恶的深坑，就开始轻视一切，而你正在否认你与真正新郎的实际盟约；你说你不是贞女，没有作过许诺，尽管你承担并公开宣认了许多童贞的誓愿。要记住你在天主、天使和人面前见证的那美好宣认。要记住那神圣的交往，贞女的神圣团体，主的集会，圣者的教会。要记住你的祖母，她在基督里年迈，却仍保持青春、德行刚健；你的母亲在主里与她竞赛，试图以奇异而非常规的辛劳打破平常的生活；要记住你的姐妹，她模仿她们的行为，甚至努力超过她们，在童贞的恩宠中超越先辈的善行，并以言以行热切地挑战你——她认为是她的姐妹——作同样的努力，同时热切祈祷你的童贞得以保存。要回想这一切，以及你与她们一同对天主的神圣事奉，虽在肉身内却是属灵的生命，虽在地上却是天上的交谈。要回想平静的白天，明亮的夜晚，属灵的诗歌，甜美的圣咏旋律，圣善的祈祷，纯洁无瑕的床榻，贞女的队伍，适度的餐饮。你庄重的容貌，优雅的举止，贞女相宜的朴素衣饰，良好的羞赧之红晕，因节制和守夜而来的苍白、明亮、胜过一切脸色的光华，如今都怎样了？你曾多少次含着泪祈祷，求能保持童贞毫无玷污！你曾多少次写信给圣者，恳求他们为你献上祈祷，不是为让你出嫁，更不是为陷入这可耻的败坏，而是为叫你不致离开主耶稣！你曾从新郎那里领受了多少恩赐？何必列举你为祂的缘故从祂的人那里所得的尊荣？何必述说你与贞女的共处，与她们同行，因童贞而受到她们的赞美，贞女们的称誉，写给贞女的信函？然而，如今有一阵来自空中之灵的微风——\BibleQuote{就是如今在悖逆之子身上运行的}——\BibleRef{弗 2:2}，你便弃绝了这一切；你以那珍贵、值得一切代价争取的宝藏，换取了短暂的享乐，这享乐一时满足口腹之欲，将来你必会觉得它比苦胆更苦。\par

3. 有谁不为此哀哭，并说：\BibleQuote{忠信之城怎会变成了娼妓？} \BibleRef{依 1:21} 主自己岂不会对那些现在依\Person{耶肋米亚}精神行事的人说：\Quote{你看见以色列的贞女对我做了什么事？} 我曾以忠诚、纯洁、正义、审断、怜悯和仁慈聘娶了她，正如我藉\Person{欧瑟亚}先知向她应许的。但她却爱上了外人，我——她的丈夫——尚在人间，她就被称为淫妇，毫不畏惧地归属另一个丈夫。那么，新娘的领路人——神圣而蒙福的\Person{保禄}，无论是古代的那一位，还是今日后来那一位（在你藉他的中介和训导离开父家、与主结合时），在忧伤中岂不会说：\BibleQuote{我所恐惧的临到了我，我所害怕的向我袭来。} \BibleRef{约 3:25} \BibleQuote{我曾把你们许配给一个丈夫，如同把贞洁的童女献给基督。} \BibleRef{格后 11:2} 我确实常常害怕\BibleQuote{恐怕你们的心意被败坏，就像蛇用狡诈诱惑了\Person{厄娃}一样；} \BibleRef{格后 11:3} 因此我采用无数对抗的咒语来控制你感官的骚动，并用无数保障来守护主的新娘。所以我不断阐述未嫁女子的生活，描绘\BibleQuote{未嫁的女子独自关心主的事，为叫她在身心上成为圣洁。} \BibleRef{格前 7:34} 我常描述童贞的高贵尊荣，并把你们当作天主的宫殿，仿佛给你的热心添上翅膀，竭力将你提升到\Person{耶稣基督}面前。但出于对邪恶的恐惧，我藉\BibleQuote{若有人毁坏天主的宫殿，天主必要毁坏他} \BibleRef{格前 3:17} 这句话来防止你跌倒。因此，我藉祈祷试图使你更稳固，\BibleQuote{好使你的身体、灵魂和心神得以保存，无可指摘，直到我们的主\Person{耶稣基督}来临。} \BibleRef{得前 5:23} 然而，我为你所做的一切辛劳都白费了。那些甜蜜的劳苦对我而言结局是苦涩的。如今我不得不在本应欢喜的事上叹息。你被蛇欺骗得比\Person{厄娃}更惨；不仅你的心意，连你的身体也受了玷污。甚至那最后可怕的事也发生了，我羞于说出，但又不能不说，因为它如燃烧的火焰在我骨中，我耗尽而无从忍受。你拿了基督的肢体，作了娼妓的肢体。\BibleRef{格前 6:15} 这是一件无与伦比的邪恶。这在生命中是新的暴行。\BibleQuote{请你们过海到基廷群岛观察，打发人到\Place{刻达尔}去留意查考，看看有没有这样的事。岂有一个民族更换了自己的神明？它们根本不是神！} \BibleRef{耶 2:10} 但这位贞女却更换了自己的荣耀，她的荣耀就在她的羞耻中。诸天为此惊骇，大地为此战栗，主说，因为这位贞女做了两件恶事：她离弃了我——圣洁灵魂的真实而圣洁的新郎——并投靠了一个不虔不义、毁灭身心的家伙。她背弃了拯救她的天主，将自己的肢体献为不洁和不义的奴仆。她忘记了我，去追随了她的情人，从那里她得不到任何好处。\par

4. 若有人把磨石套在他的脖子上，投在海里，也比让他得罪主的贞女更好。哪个奴隶曾狂妄至如此疯狂，竟扑到主人的床上？哪个强盗曾愚昧至如此地步，竟伸手取天主的供物——不是死物的器皿，而是活的身体，内里珍藏着一个按天主的肖像所造的灵魂？\par

有谁曾胆大包天，在城中心和正午时分，在皇家的雕像上刻上污秽的猪形？若有人藐视人的婚姻，在两三个证人面前被处死，不得怜悯。那么，你们想，那践踏了天主子、玷污了祂聘娶的新娘、侮辱了童贞之神的人，该受怎样更重的惩罚呢？但女人，他辩称，是自愿的，我没有对她施暴。所以，那位埃及的荡妇曾对俊美的\Person{若瑟}燃起情欲，但这位贞洁青年的德性并未被那恶妇的疯狂所克服，甚至当她向他伸手时，他也没有被迫行恶。但他仍辩称，对她来说，这并非新鲜事；她已不是处女；如果我不愿意，她也会被另一个人糟蹋。是的，经上记载：人子固然被预定要被出卖，但那出卖人子的人有祸了！绊倒人的事是免不了的，但那绊倒人的有祸了！\par

5. 在如此情形下，\Quote{他们跌倒，不再起来吗？他转面离去，不再回来吗？} \BibleRef{耶 8:4} 为何童女可耻地转面离去，尽管她已听见她的新郎基督通过耶肋米亚的口说：\Quote{我说：她行了这一切事（行了这一切淫乱，按七十贤士译本）之后，你当归向我，她却不回来？} \BibleRef{耶 3:7} \Quote{基肋阿得没有乳香吗？那里没有医生吗？为何我人民的女儿不得痊愈？} \BibleRef{耶 8:22} 你实在可以在圣经中找到许多治恶的良药，许多拯救人脱离毁灭、引向健康的药物：死亡与复活的奥迹、可怕审判与永罚的判决、悔改与罪赦的道理，以及无数悔改的比喻——那块银币、那只羊、那个与娼妓耗尽家财的儿子，他迷失了又被寻回，他死了又复活了。我们不要滥用这些良药；让我们藉这些方法医治我们的灵魂。要思念你的末日，因为你必不像其他妇女一般永远活着。那痛苦、喘息、临终时刻、即将来临的天主的审判、匆匆赶来的天使、因罪恶意识而恐惧战栗、痛苦不堪的灵魂，它可怜地转向今生的事物，又不得不面对将来在别处度过的漫长生命。请你想象，随着你脑海中的画面，一切人生的结局——当天主子在祂的荣耀中与祂的天使降临时，\Quote{因为祂要来，并不缄默；} 当祂来审判生者死者，按各人的行为施行报应时；当那可怕的号角以巨大的声音唤醒沉睡万代的人，行善的复活得生，作恶的复活定罪时。要记得达尼尔神的异象，看他如何将审判呈现在我们面前：\Quote{我观看，直到宝座安置好，万古常存者坐上去，祂的衣服洁白如雪，头发如纯净的羊毛；……祂的宝座是火焰，轮子是烈火。一道火河涌出，从祂面前流出来；有千千万万事奉祂，有亿亿万万侍立在祂面前；审判者已坐定，案卷已展开，} \BibleRef{达 7:9} 在众天使和众人面前清楚揭露一切善恶之事，公开与隐秘的行为、言语、思想，同时呈现。那么，那些生活邪恶的人将是怎样呢？那时，那灵魂在所有这些观众眼前，突然以其全部羞耻显露出来，它将躲到哪里去呢？它将用怎样的身体承受那永不止息、不可忍受的剧痛——在永不熄灭的火中，在永远不死、永不消灭的虫子里，在黑暗可怖的哈德斯深渊里；痛苦的哀号、大声的哭泣、切齿与无尽的苦楚？死亡之后，无法摆脱这一切灾祸；没有任何方法或途径能从痛苦刑罚中解脱。\par

6. 我们现在还能逃脱。趁着还能，让我们从跌倒中起来：只要我们离弃邪恶，就绝不要对自己绝望。耶稣基督来到世上，是为拯救罪人。\Quote{请来，让我们俯伏敬拜，在祂面前哭泣。} 邀请我们悔改的圣言高呼：\Quote{凡劳苦和负重担的，你们都到我这里来，我必要使你们安息。} \BibleRef{玛 11:28} 所以，如果我们愿意，就有救恩之路。\Quote{死亡曾逞凶吞噬，但上主再次擦去悔改者脸上的所有眼泪。} 上主在祂所有的话语中都是信实的。祂说：\Quote{你们的罪虽似朱红，将变成雪一样洁白；虽红如丹砂，将白如羊毛。} \BibleRef{依 1:18} 祂并不说谎。灵魂的伟大良医，祂随时准备好解救，不仅是你，也是所有被罪恶奴役的人，祂愿意医治你的疾病。正是祂说出的这些话，正是祂甘甜而拯救的嘴唇说：\Quote{健康的人不需要医生，有病的人才需要……我来不是为召义人，而是为召罪人悔改。} \BibleRef{玛 9:12} 当祂这样说时，你还有什么借口？任何人还有什么借口？上主愿意洁净你，使你摆脱疾病的困扰，并在黑暗之后向你显示光明。那好牧人，祂撇下那些没有走迷的羊，正在寻找你。如果你将自己交给祂，祂不会拒绝。祂出于慈爱，甚至不耻于将你扛在肩上，为寻回迷失的羊而欢喜。圣父站着等待你从流浪中归来。你只要回来，还远远的时候，祂就会跑过来，扑在你颈上，在你藉悔改得以洁净之后，用慈爱的拥抱环绕你。祂会为上等袍子，给那脱去旧人及其一切行为的灵魂穿上；祂会把戒指戴在洗去死亡之血的手上，给那从邪恶之路转向福音平安之路的脚穿上鞋。祂会向自己的子民，包括天使和人，宣报喜乐欢庆的日子，并四方传扬你的救恩。因为祂说：\Quote{我实在告诉你们：在天主面前，为一个悔改的罪人，有喜乐。} 如果那些自以为站得稳的人因你迅速被接纳而责备，好父亲会亲自为你用这些话回答：\Quote{我们应当欢宴喜乐，因为} 我的这个女儿 \Quote{死了又复活了，迷失了又被寻回了。} \BibleRef{路 15:32}\par

\SourceRecord{Translated by Blomfield Jackson.；Nicene and Post-Nicene Fathers, Second Series；Vol. 8.；Edited by Philip Schaff and Henry Wace.；Buffalo, NY: Christian Literature Publishing Co.,；1895.；<http://www.newadvent.org/fathers/3202046.htm>.}

% structure: p0001 source=h1 target=section reason=page-title

\section{第四十七封书信}

\Person{圣巴西略·凯撒勒雅}\par

\Emph{致\Person{额我略}}。\par

谁能赐给我像鸽子一样的翅膀？或者我衰迈之年怎能如此更新，使我能够旅行到你的怀抱，满足我深切的渴望见到你，告诉你我灵魂的一切烦忧，并从你那里得到一些安慰来减轻我的痛苦呢？因为当有福的主教\Person{欧瑟伯}安眠之后，我们并非没有极大的忧虑，恐怕反对我们大都会教会的阴谋者，企图以他们异端的莠子充满其中，会抓住眼前的机会，用他们邪恶的教导将许多辛劳播撒在人心中的真信德连根拔除，并破坏其合一。这已是他们许多教会行动的结果。然而，当我收到司铎们的来信，劝勉我不要在这样的危急时刻忽略他们的需要时，我环顾四周，想起了你的仁爱精神、你的正确信德以及你为天主的教会所怀的不懈热忱。因此，我已派遣可爱的执事\Person{欧斯大削}去邀请你的尊驾，并恳求你为教会再增添这一劳苦。我也恳求你以见你一面来慰藉我的晚年；并为真教会维护其著名的正统信仰，倘若我配得上与你联合，就请与我一同在这善工上努力，按照主的旨意给她一位牧人，能够正确引导祂的子民。我眼前有一个连你自己也不陌生的人。只要我们配得获得他，我确信我们将获得对天主的坦然无惧，并对那呼求我们援助的人民施以极大的恩惠。现在再次，甚至多次，我呼求你，放下一切犹豫，来与我会面，并在冬天的困难来临之前动身。\par

\SourceRecord{Translated by Blomfield Jackson.；Nicene and Post-Nicene Fathers, Second Series；Vol. 8.；Edited by Philip Schaff and Henry Wace.；Buffalo, NY: Christian Literature Publishing Co.,；1895.；<http://www.newadvent.org/fathers/3202047.htm>.}

% structure: p0001 source=h1 target=section reason=page-title

\section{第四十八封信}

\Emph{\Person{圣巴西略（凯撒勒雅）}}\par

\Emph{致\Place{萨摩撒塔}主教\Person{尤西比乌斯}}。\par

我颇费周折才找到一位信使将信函呈递给阁下，因为我们的人如此畏惧冬季，甚至几乎不敢将头探出屋外。我们遭遇了一场极重的降雪，以至于连人带屋都被掩埋其中，整整两个月来，我们生活在洞穴和窑洞里。您了解卡帕多齐亚人的性格，要让我们动弹一下是多么困难。那么，请宽恕我未能及早写信，将安提约基雅的最新消息禀报尊座。如今再向您讲述这一切，而您很可能早已得知，便显得乏味而无趣。然而，即使讲述您已知之事，我也不以为烦劳，我已将那位读经员带来的信函寄给了您。关于此事，我不再多言。君士坦丁堡现已有一段时间由\Person{德莫菲卢斯}掌管，此信的携带者会亲自告诉您，而且您的圣座无疑已收到报告。从所有来自该城的人那里，关于他一致报告说有一种对正统和纯正信仰的仿冒，以至于甚至城中分裂的各派都已达成一致，一些邻近的主教也接受了和解。我们这里的人并不比我所预料的更好。您一走，他们就来，说了许多令人痛心的话，做了许多令人痛心的事，最后又回去了，从而使他们与我的分离更加扩大了。将来是否会有更好的事发生，他们是否会放弃恶行，除了天主，无人知晓。以上就是我们目前的状况。教会其余部分，因天主的恩宠，依然稳固，并祈祷春天我们能再次与您同在，因您的好建议而更新。我的健康状况一如既往，没有好转。\par

\SourceRecord{Translated by Blomfield Jackson.；Nicene and Post-Nicene Fathers, Second Series；Vol. 8.；Edited by Philip Schaff and Henry Wace.；Buffalo, NY: Christian Literature Publishing Co.,；1895.；<http://www.newadvent.org/fathers/3202048.htm>.}

% structure: p0001 source=h1 target=section reason=page-title

\section{第四十九封书信}

\Emph{圣巴西略（凯撒勒雅）}\par

\Emph{致主教\Person{阿尔卡狄乌斯}}。\par

我最虔诚的兄弟，当我读到你的信函时，我感谢了神圣的天主。我祈求我可能不辜负你对我的期望，并因你因我们的主耶稣基督之名给予我的尊荣而获得圆满的赏报。我非常喜悦地听说你从事了一件极其适合基督徒的事，为光荣天主建造了一座房屋，并且实际热切地爱慕——如经上所载——\BibleQuote{主殿的华丽}\BibleRef{咏 27:4}，并为你自己预备了那座天上的居所，那是主在他的安息中为爱主的人所预备的。若我能找到任何殉道者的遗骸，我祈求能参与你热诚的努力。若\BibleQuote{义人要永远被记念}\BibleRef{咏 111:6}，我无疑将有分于圣者要赐给你的美誉。\par

\SourceRecord{Translated by Blomfield Jackson.；Nicene and Post-Nicene Fathers, Second Series；Vol. 8.；Edited by Philip Schaff and Henry Wace.；Buffalo, NY: Christian Literature Publishing Co.,；1895.；<http://www.newadvent.org/fathers/3202049.htm>.}

% structure: p0001 source=h1 target=section reason=page-title

\section{书信五十}

\Emph{\Person{圣巴西略·凯撒利亚}}\par

\Emph{致\Person{依诺增爵}主教。}\par

确实，谁能比你更适合鼓励胆怯者、唤醒沉睡者呢？我敬爱的主，你已在此事上显出你的全面卓越，因为你甘愿降卑到你卑微的下属中间，如同那位说过\BibleQuote{我在你们中间，不像同席者而是如同服事者。}\BibleRef{路 22:27}因你已屈尊将属灵的喜乐服事给我们，用你尊贵的书信更新我们的灵魂，仿佛将你伟大的手臂环绕在孩童的幼稚之上。因此，我们恳求你善良的灵魂祈祷，使我们堪当从像你这样伟大的人那里获得援助，并能有口才与智慧，得以与所有像你一样被圣神引导的人一同应和。我听说你是他的朋友和真诚的朝拜者，我深为感谢你对天主坚定不渝的爱。我祈求我的份能与真诚的朝拜者一同找到，我们确信你的阁下必列于其中，连同那位以奇妙工作充满普世的伟大而真实的主教。\par

\SourceRecord{Translated by Blomfield Jackson.；Nicene and Post-Nicene Fathers, Second Series；Vol. 8.；Edited by Philip Schaff and Henry Wace.；Buffalo, NY: Christian Literature Publishing Co.,；1895.；<http://www.newadvent.org/fathers/3202050.htm>.}

% structure: p0001 source=h1 target=section reason=page-title

\section{第五十一书}

\Emph{\Person{凯撒勒雅的圣巴西略}}\par

\Emph{致\Person{波西波留}主教}。\par

你认为我听到那些不惧怕审判者——祂\Quote{将消灭说谎者}——的人对我堆砌的诽谤时，我的心是何等痛苦？我几乎整夜未眠，思念你充满爱的话语；悲伤紧紧抓住了我的心扉。诚如\Person{撒罗满}所言，诽谤使人卑微。没有人会如此无感觉，以至内心不被触动，并俯伏于地，若他遇到惯于说谎的嘴唇。但我们必要忍受一切，承担一切，将我们的辩护交托给主。祂不会轻视我们；因为\BibleQuote{欺压穷人的，就是辱骂造他的主}。\BibleRef{箴 14:31} 那些拼凑出这出亵渎新悲剧的人，似乎已对主失去了一切信心，主曾宣告，在审判之日，我们连一句闲话也要交账。\BibleRef{玛 12:36} 至于我，告诉我，我绝罚了可敬的\Person{迪亚尼乌斯}吗？因为这就是他们对我所说的话。在哪里？何时？在谁面前？什么借口？仅仅口头，还是书写？跟随他人，还是我自己作为此事的发起者与创作者？哀哉！这些人厚颜无耻，什么话都说得出口！哀哉！他们藐视天主的审判！除非他们真的在虚构中又加上一条，说我有一次神志不清，不知道自己在说什么。因为只要我神志清醒，我就知道我从未做过任何此类事，也从未有过丝毫愿望这样做。我所确实意识到的乃是：自孩提时代起，我便在爱中受教养，当我注视他时，思念他的仪容，观其何等可敬、何等威严、何等尊荣。及至年长，我因他灵魂的美善而认识他，乐于与他为伴，逐渐学会感知他品格的纯朴、高贵与宽宏，以及一切最卓越的特质：他灵魂的温柔，兼具宽宏与谦和，举止合宜，节制脾气，和蔼可亲与庄重威严相融合。由此，我将他列于品行最为卓越之人中。\par

然而，在他生命的末期（我不隐瞒真相），我与本国许多敬畏主的人一起，为他忧心如焚，因为他签署了\Person{乔治}从\Place{君士坦丁堡}带来的信经。后来，他满有仁慈与温柔，并愿以慈父之心使众人满意，当他已染上致死之病时，他派人来召我，并呼主为证，说他出于心地纯朴同意了从\Place{君士坦丁堡}送来的文件，但绝无拒绝圣教父们在\Place{尼西亚}所颁布的信经之意，他的心意自始至终从未改变。他还祈祷，但愿他不被排除在那三百一十八位可祝福的主教——他们曾向世界宣告了敬虔法令——的行列之外。因这令人满意的声明，我消除了所有焦虑与疑虑，如你所知，与他共融，不再悲伤。我与\Person{迪亚尼乌斯}的关系就是如此。若有人声称他知晓任何我针对\Person{迪亚尼乌斯}的卑劣诽谤，就让他不要像奴隶一样在角落里窃窃私语；让他大胆出来，在光天化日之下定我的罪。\par

\SourceRecord{Translated by Blomfield Jackson.；Nicene and Post-Nicene Fathers, Second Series；Vol. 8.；Edited by Philip Schaff and Henry Wace.；Buffalo, NY: Christian Literature Publishing Co.,；1895.；<http://www.newadvent.org/fathers/3202051.htm>.}

% structure: p0001 source=h1 target=section reason=page-title

\section{书信五十二}

\Emph{\Person{凯撒勒雅的巴西略}}\par

\Emph{致女隐修院院长们。}\par

1. 一个痛苦的传闻传入我耳中，令我非常忧伤；但我的弟兄，天主所爱的\Person{波私波里}主教，却带来了关于你们的更令人满意的报告，同样令我喜悦。他因天主的恩宠断言，所有关于你们的谣言都是那些对你们实情不确切了解的人捏造的。他还补充说，他在你们中间发现了一些关于我的不虔敬的诽谤，此类诽谤出自那些不指望在祂公义报应的日子向审判者连一句闲话都要交账的人。因此，我感谢天主，既因为我从对你们有害的看法中得了痊愈——这看法来自人们的诽谤——又因为我听到你们在听到我弟兄的保证后，放弃了那些关于我的毫无根据的想法。他在一切出于自己所说的话中，也完全表达了我自己的感受。因为我们两人在信仰上同心，都是继承同一批教父，他们曾在\Place{尼西亚}颁布了关于信仰的伟大法令。其中一些部分被普遍接受而无争议，但\OriginalTerm{同质}{homoousion}在某些地方未被妥善接受，至今仍被一些人拒绝。对于这些反对者，我们完全可以责备他们，然而另一方面，又认为他们值得宽恕。拒绝追随教父，不认为他们的声明比自己的意见更有权威，这种行为是应受责备的，因为它充满了自以为是。另一方面，他们怀疑一个被敌对派别曲解的用语，这一事实确实在某种程度上减轻了他们的责任。此外，事实上，为审理\Person{萨摩萨塔的保禄}案件而召开的各次主教会议，确实批评了这个用语，认为它是不恰当的。\par

因为他们主张，\OriginalTerm{同质}{homoousion}既提出了本质的概念，也提出了从本质所出之物的概念，以至于本质在被分割时，将\OriginalTerm{同本质}{co-essential}的称谓赋予了它所分成的各部分。就青铜和用它制成的钱币而言，这种解释有一定道理；但在天主圣父和天主圣子方面，并不存在先于或甚至作为两者基础的本体；仅仅想到或说出这样的事情，就是至不虔敬的极端。在无始者之先，还能想到什么？由于这种亵渎，对圣父和圣子的信仰就被摧毁了，因为由同一个本体构成的事物，彼此之间的关系是兄弟关系。\par

2. 因为在那个时候，有些人断言圣子是从无中产生的，所以采用了\OriginalTerm{同质}{homoousion}一词，以根除这种不虔敬。因为圣子与圣父的结合是没有时间、没有间隔的。前面的话已表明这就是其本意。因为在说了圣子是光之光，由圣父的本体所生，而非受造之后，他们接着加上了\OriginalTerm{同质}{homoousion}，从而表明，任何人赋予圣父多少光，也必赋予圣子多少光。因为真光相对于真光，按照光的实际概念，不会有任何变化。既然圣父是无始之光，圣子是受生之光，但两者都是光，是光；他们正确地说了\Quote{同一本体}，以表明本性的同等尊贵。具有兄弟关系的事物，并非如某些人所认为的那样是同质的；但当原因和从原因中由其本性而存在之物具有同一本性时，它们才被称为\Quote{同质}。\par

3. 这个术语也纠正了\Person{撒贝略}的错误，因为它除去了位格同一性的概念，并完善地引入了\OriginalTerm{位格}{hypostases}的概念。因为没有任何东西是与自身同质的，而是一物与另一物同质。因此，这个词具有卓越而正统的用法，既界定了\OriginalTerm{位格}{hypostases}各自的特征，又表明了本性的不变性。当我们被告知圣子是由圣父的本体所生、而非受造时，我们不要陷入关系的物质性理解。因为本体并没有从圣父分离而赐予圣子；本体也不是通过流溢，也不是像植物结出果实那样而产生的。神圣生育的方式是无法言说、非人类思想所能理解的。用可朽的时间之物来比较永恒之物，并想象神如同物质之物一样生育，这确实是贫乏、属血肉的智力的特征；我们更应该通过相反的论证来提升到真理，并说：既然可朽者是这样，那不朽者就不是这样。因此，我们既不可否认神圣生育，也不可用物质性的感官玷污我们的理智。\par

4. 圣神也与圣父和圣子同列，因为祂超越受造界，且按主在福音中所教导我们的：\Quote{你们要去因父及子及圣神之名给他们授洗。}\BibleRef{玛 28:19} 反之，那将圣神置于圣子之前，或声称祂比圣父更古老的人，是抗拒天主的诫命，与纯正的信德无分，因为他没有保持他所领受的赞美颂的形式，反而为了取悦人而采用某种新奇的发明。经上记载：\BibleQuote{圣神是出于天主的，}\BibleRef{格前 2:12} 如果祂是出于天主的，祂怎能比祂所出自的那位更古老呢？既然有一位 \OriginalTerm{无始者}{Unbegotten}，还说有什么比那无始者更优越，那岂不是愚蠢至极？祂甚至也不先在，因为圣子与圣父之间没有任何间隔。然而，如果祂不是出于天主，而是通过基督，那么祂根本就不存在。由此可见，关于顺序的这种新发明实际上摧毁了圣神本身的存在，并且否定了整个信德。将祂降低到受造物的地位，或将祂置于圣子或圣父之下，无论是在时间上还是在位阶上，同样都是不虔敬的。这些就是我所听说你在询问的问题。如果主恩准我们见面，我或许会就这些主题多说一些，而且我自己也可能就我正在探讨的问题从你那里得到满意的答复。\par

\SourceRecord{Translated by Blomfield Jackson.；Nicene and Post-Nicene Fathers, Second Series；Vol. 8.；Edited by Philip Schaff and Henry Wace.；Buffalo, NY: Christian Literature Publishing Co.,；1895.；<http://www.newadvent.org/fathers/3202052.htm>.}

% structure: p0001 source=h1 target=section reason=page-title

\section{\WorkTitle{书信五十三}}

\Person{圣巴西略，凯撒勒雅}\par

\Emph{致乡村主教}\par

1. 我所写之事极为重大，单就此事已引起众人猜疑与议论，为此我心深感痛楚。但迄今为止，我仍觉难以置信。因此，我愿所写之事为有罪者成为良药，为无辜者成为警诫，为漠不关心者——我相信你们中无人属此类——成为见证。我所说的是何事？有报告称，你们中有人从晋秩候选人手中收受金钱，并以宗教为借口加以掩饰。这实是更糟。若有人假借善行作恶，他应受双倍惩罚：因为他不仅行不善之事，而且，可谓使善成了罪的同谋。若此事属实，愿它不再发生。愿更好的局面开始。对受贿者，必须如宗徒对那愿以金钱购买圣神恩赐之人所说：\BibleQuote{你的银子与你一同丧亡！}（\BibleRef{宗 8:20}）无知而想购买天主的恩赐，其罪轻于出售。出售确是事实；若你出售白白领受的恩赐，你将失去这恩惠，如同自己被卖给撒殚一般。你将小贩的交易带入了属灵之事，带入了我们受托保管基督体血的教会。这些事绝不可行。我还要指出其中一种巧妙的诡计：有人以为因不是在晋秩之前而是在之后收钱，所以不算犯罪；但无论何时，收受就是收受。\par

2. 因此，我劝告你们：舍弃这种收益，或更确切地说，舍弃这通往地狱之路。不要因以这类贿赂玷污双手，而使自己不配举行神圣奥迹。但请宽恕我：我起初以不信为由，如今却如确信般发出威胁。若在我这封信之后，仍有人行此类事，他必远离此处的祭台，去寻找一个他能买卖天主恩赐的地方。我们与天主的众教会并无这种习俗。再说一句，我便结束：这些事源于贪财。贪财是万恶之根，且被称为偶像崇拜。因此，不要为了一点金钱而将偶像置于基督之上。不要效法犹达斯，再次为贿赂出卖那为我们被钉十字架者。因为所有如此获利之人的田地和双手，都将被称为血田。\par

\SourceRecord{Translated by Blomfield Jackson.；Nicene and Post-Nicene Fathers, Second Series；Vol. 8.；Edited by Philip Schaff and Henry Wace.；Buffalo, NY: Christian Literature Publishing Co.,；1895.；<http://www.newadvent.org/fathers/3202053.htm>.}

% structure: p0001 source=h1 target=section reason=page-title

\section{第五十四封信}

\Signature{凯撒勒雅的圣巴西略}\par

致\OriginalTerm{乡村主教}{Chorepiscopi}\par

我非常痛心，因为教父的教规已被废弃，教会的严格纪律从你们中间消失。我担心，随着这种漠不关心滋长，教会事务会逐渐陷入混乱。按照天主的教会中遵守的古代惯例，教会中接纳神职人员时，须经过仔细审查；他们的一生受到查考，查明他们不是毁谤者，也不是酒徒，不轻易争吵，使他们的青春受约束，以便能持守\BibleQuote{没有圣德，谁也不能见主}\BibleRef{希 12:14}。这审查由与他们一起生活的长老和执事进行。然后他们将这些人带到乡村主教那里；乡村主教在得到见证人确证属实，并告知主教后，如此就接纳神职人员进入司祭职。然而现在，你们完全忽略了我；你们甚至没有屈尊征询我的意见，并将全部权力转归自己。此外，你们漠不关心，任由长老和执事将不配的人引入教会，只要是他们所选的人，不作任何生活和品格的预审，单凭徇私，基于亲属关系或其他关系。结果是，每个村庄都有许多神职人员，却没有一个配得上祭台服务的人。你们自己从选立合适候选人的困难中提供了证明。因此，当我看出这恶已逐渐发展到无法挽回的地步，尤其此刻因惧怕征兵而有大量人员自荐供职时，我不得不求助于恢复教父的教规。于是我书面命令你们，将各村庄的神职人员名册送交给我，写明每人由谁引入，以及其生活方式。你们手中存有名册，以便你们的版本可与我的文档对照，无人能随意加入自己的名字。因此，若有人在初次任命后由长老引入，就应予以拒绝，并使他们列于平信徒之中。然后你们必须重新开始对他们的审查，若证明合格，就由你们决定接纳。将不配的人逐出教会，如此洁净教会。今后，通过审查测试合格的人，然后接纳他们；但在报告我之前，不得将他们计入人数。否则，你们要清楚明白：凡未经我授权而被接纳供职的人，仍将是平信徒。\par

\SourceRecord{Translated by Blomfield Jackson.；Nicene and Post-Nicene Fathers, Second Series；Vol. 8.；Edited by Philip Schaff and Henry Wace.；Buffalo, NY: Christian Literature Publishing Co.,；1895.；<http://www.newadvent.org/fathers/3202054.htm>.}

% structure: p0001 source=h1 target=section reason=page-title

\section{第五十五封书信}

\Emph{\Person{圣巴西略}（\Place{恺撒勒雅}）}\par

\Emph{致\Person{帕雷哥里乌斯}司铎}\par

我已耐心细读你的信函，令我惊讶的是，你明明能立即采取行动，向我提供简短而容易的辩解，却偏要坚持我抱怨之事，并试图通过长篇大论来医治那不可医治的。我不是第一个，帕雷哥里乌斯，也不是唯一一个定下此规条的人，即女人不可与男子同住。请读一读我们圣教父们在尼西亚会议所颁布的教规，它明确禁止有暗中同住的女人。独身生活之所以可敬，正在于它断绝与所有女性的交往。因此，如果有人外表上声称独身，实际上却效仿那些娶妻的人，那他显然只是在名义上追求童贞的名声，并未远离不当的放纵。你本应更乐意顺从我的要求，因为你自己声称已摆脱一切肉体欲望。我并不认为一个七十岁的老人与女人同住是出于这种欲望，我也并非因任何罪行已犯而作出决定。但我们从宗徒那里学到：\BibleQuote{不可将绊脚石或障碍物放在弟兄的路上}；\BibleRef{罗 14:13} 而且我知道，某些人做得非常妥当的事，自然而然地会成为别人犯罪的诱因。因此，我遵照圣教父们的训令，下令你必须与那女人分开。那么，你为何要责备乡村主教呢？重提旧怨有何益处？你为何怪我轻信诽谤？为何不归咎于你自己，不肯断绝与那女人的关系？把她从你家中赶出去，安置在一所隐修院里。让她与贞女同住，而你自己由男人伺候，免得天主的名因你而受到亵渎。在你做完这些之前，你在信中使用的无数论据对你毫无用处。你将徒然死去，并需向天主交账你的无用。如果你坚持不改正自己的行径而紧握圣职地位，你将在全体人民面前受诅咒，而所有接待你的人，也将在整个教会中受绝罚。\par

\SourceRecord{Translated by Blomfield Jackson.；Nicene and Post-Nicene Fathers, Second Series；Vol. 8.；Edited by Philip Schaff and Henry Wace.；Buffalo, NY: Christian Literature Publishing Co.,；1895.；<http://www.newadvent.org/fathers/3202055.htm>.}

% structure: p0001 source=h1 target=section reason=page-title

\section{第五十六封信}

\Person{圣巴西略}（\Place{凯撒勒雅}）\par

\Emph{致\Person{培尔加米乌斯}}。\par

我天生容易遗忘，近来又多事，以致这一天然弱点加剧。因此，我深信您曾写信给我，尽管我不记得从阁下处收到过任何信；因为我相信您不会说不实之事。但若说没有回信，错不在我，而在于那位未曾索取回信的人。不过现在，您收到了这封信，它包含了我对过去的辩护，并为第二次问候提供了基础。所以，当您写信给我时，不要以为您是在开启另一段通信。您只是在此履行您应尽的本分。因为事实上，尽管我这封信是对您前一封信的回复，但由于它比您那封信厚两倍以上，它将实现双重目的。您瞧，我的慵懒把我驱向了怎样的诡辩。但是，亲爱的先生，请不要用几句话就提出严重的指控，实则是最严重的指控。对朋友的遗忘以及因高位而产生的忽视，是包括各种错误的过失。我们若未能按主的诫命去爱，那么我们就失去了烙在我们身上的标记。我们若被空虚的骄傲和傲慢填满，我们就陷入了魔鬼无可避免的定罪。因此，如果您说这些话是因为您对我有这样的看法，那么请祈祷我能逃离您在我身上发现的邪恶；然而，如果您舌头说出这些话，是由于某种轻率的俗套，那么我将自我安慰，并请您不吝提出一些证据来支持您的指控。请确信，我目前的焦虑给了我谦卑的机会。当我停止认识自己时，我将开始忘记您。那么，绝不要以为一个人很忙碌，他就是有缺陷的人。\par

\SourceRecord{Translated by Blomfield Jackson.；Nicene and Post-Nicene Fathers, Second Series；Vol. 8.；Edited by Philip Schaff and Henry Wace.；Buffalo, NY: Christian Literature Publishing Co.,；1895.；<http://www.newadvent.org/fathers/3202056.htm>.}

% structure: p0001 source=h1 target=section reason=page-title

\section{第57封信}

\Emph{\Place{凯撒勒雅}的\Person{圣巴西略}}\par

\Emph{致\Place{安提约基雅}主教\Person{默勒丢}}。\par

如果您的圣座知道每次您写信给我时给我带来的快乐有多大，我知道您绝不会放过任何给我寄信的机会；不仅如此，您还会每次写多封信，因为知道我们慈爱的主为安慰受苦者所预备的赏报。这里的一切仍处于非常痛苦的境况，唯有对您圣座的思念将我从自己的烦恼中唤醒；通过与您那封充满智慧和恩宠的信函的交流，这个思念对我来说更加清晰。因此，当我将您的信拿在手中时，首先我会看它的篇幅，我因它如此之大而更加喜爱；然后，当我阅读时，我为其中发现的每一个字而欢喜；当我接近结尾时，我开始感到悲伤；因为您所写的每一个字都如此美好。一颗善心所溢出的就是善。然而，若蒙允准，因应您的祈祷，在我今生仍活于世上时，能与你面对面相见，并享受你生动声音的有益教导，或任何帮助我现世或来生生活的援助，我会将此视为最好的祝福，是天主仁慈的前奏。我本应早已坚持这一意向，若不是被真诚而仁爱的弟兄们阻止。我已告诉我兄弟\Person{特奥弗拉斯图斯}向您详细报告事务，关于这些事我并未将意向付诸笔端。\par

\SourceRecord{Translated by Blomfield Jackson.；Nicene and Post-Nicene Fathers, Second Series；Vol. 8.；Edited by Philip Schaff and Henry Wace.；Buffalo, NY: Christian Literature Publishing Co.,；1895.；<http://www.newadvent.org/fathers/3202057.htm>.}

% structure: p0001 source=h1 target=section reason=page-title

\section{第五十八封书信}

\Emph{\Person{圣巴西略·凯撒勒雅}}\par

\Emph{致我弟\Person{额我略}}\par

我如何能与你写信争辩？我怎能充分抓住你，带着你所有的纯朴？告诉我，谁曾第三次落入同一张网？谁曾第三次陷入同一个陷阱？甚至一头牲畜也很难做到。你伪造了一封信，带给我，仿佛它来自我可敬的主教叔父，试图欺骗我，我不知道为什么。我把它当作主教写并由你送交的信函收下了。我为什么不应该呢？我很高兴；我把它给许多朋友看；我感谢天主。伪造被发现，因主教亲自否认。我深感羞愧；因被狡猾欺骗和谎言的耻辱所覆盖，我祈祷大地为我裂开。然后他们给了我第二封信，据称是主教本人经由你的仆人阿斯德里乌送来的。连这第二封也并非主教真正送来的，正如我可敬的弟兄安提默告诉我的。现在阿达曼提乌带着第三封到来。我应如何接收由你或你的人带来的一封信？我本可祈祷有一颗铁石心肠，既不记得过去，也不感受现在；像牲畜一样，低着头承受每一次打击。但我应如何想呢，既然在第一次和第二次经历之后，没有确凿证据我什么都不能相信？因此我写信指责你的纯朴，我清楚地看到这既不是一般基督徒所应有的，也不适合当前的紧急情况；我写这信，为叫你在将来至少照顾好自己并放过我。我必须极其坦诚地对你说，我告诉你，你是一个不配服务于如此大事的人。然而，无论信是谁写的，我都已作了合宜的回复。那么，无论你是自己在试探我，还是你送来的信确实是从主教们那里收到的，你已得到了我的答复。在这样的时刻，你本应记着你是我的弟兄，尚未忘记天性的联系，不要把我当作敌人看待，因为我所进入的生活正在耗尽我的力量，且远远超过我的能力，甚至伤害我的灵魂。然而尽管如此，既然你已决定向我宣战，你本应到我这里来，分担我的困苦。因为经上说：\BibleQuote{弟兄们和援助是在患难的时候。}\BibleRef{训 40:24}如果可敬的主教们确实愿意见我，请告知我一个地点和时间，并差他们自己的人来邀请我。我不拒绝见我的叔父，但除非邀请以正当合适的方式送达，我不会去。\par

\SourceRecord{Translated by Blomfield Jackson.；Nicene and Post-Nicene Fathers, Second Series；Vol. 8.；Edited by Philip Schaff and Henry Wace.；Buffalo, NY: Christian Literature Publishing Co.,；1895.；<http://www.newadvent.org/fathers/3202058.htm>.}

% structure: p0001 source=h1 target=section reason=page-title

\section{第五十九封书信}

\Emph{圣\Person{巴西略}（\Place{凯撒勒雅}）}\par

\Emph{致其叔父\Person{额我略}}。\par

1. 我沉默了许久。难道我要永远沉默吗？我是否还要继续对自己施加沉默这最严厉的惩罚，既不写信，也不接受他人的任何陈述？因着坚守这一严厉的决心直到如今，我能将先知的话用在自己身上：\Quote{我如临产妇人，隐忍等待。}然而，我始终渴望能亲自或通过书信与你交流，却因自己的罪过而始终未能如愿。因为我不能设想所发生之事有任何其他原因，除了我所确信的这一个：因着与你的爱隔绝，我是在补赎旧罪；——倘若我在你的事上使用\InnerQuote{隔绝}一词没有错的话，对任何人，尤其对我，你向来待我如父。如今，我的罪如同浓云笼罩着我，使我忘却了这一切。当我想到所发生之事仅给我带来悲伤，我怎能不将其归咎于自己的邪恶？但如果事情要追溯到罪过，那就让这成为我痛苦的终结；如果这其中有任何惩戒的意图，那么你的目的已完全达成，因为惩罚已持续多时。因此，我不再叹息，而是率先打破沉默，恳请你记念我与你，你一生中对我所怀有的强烈关爱，远超亲情所要求。如今，我恳求你，为了我而善待这座城市。不要因我的缘故而疏远它。\par

2. 因此，若在基督内有任何安慰，圣神有任何相通，有任何怜悯与慈悲，就请成全我的祈求。止住我的沮丧。让未来有一个更光明的开端。你自己在通往一切至善的道路上作他人的领袖，不要跟随任何人步入歧途。温柔与和平之于你的灵魂，如同任何特征之于任何人的身体一样显著。像你这样的一位，理应吸引众人归向自己，并使所有亲近你的人都被你良善的本性所浸润，如同没药的馨香。因为尽管目前存在一定程度的反对，但不久之后，人们必会认识到和平的福祉。然而，只要纷争所生的诽谤还有立足之地，猜疑就必定日益恶化。其他人不理会我固然不妥，但在尊驾身上尤其不妥。若我错了，我因受责备而变得更好；但若我们从未见面，这便不可能。但若我没有做错，又为何被人厌恶？我以此为自己辩护。\par

3. 至于众教会可能为自己所说的话，也许我最好保持沉默：他们承受着我们分歧的后果，这对他们并无益处。我说话并非为了沉溺于悲伤，而是为了止息悲伤。我相信，你的智慧不会让任何事情逃过你的眼睛。你本人将比我所想象的更能辨明并向他人讲述远为重要的事。你比我更早看到教会的祸患，你比我更忧伤，因为你早已从主那里学会了不可轻看最小的一个。如今，祸患不只限于一两个人，而是整个城市和人民都分担着我的灾祸。何需述说在我疆界之外，关于我的传闻会是怎样？像你这样宽厚的人，最好将好斗之心留给别人；不，如果可能，更要从他们心中根除它，而你则以忍耐战胜可悲之事。任何愤怒的人都能自卫，但超越愤怒本身，只属于你，以及任何像你这样善良的人，如果有的话。有一件事我不愿说：那对我怀恨的人，是让他的怒气落在无辜者身上。所以，请安慰我的灵魂：或来我这里，或写信，或邀请我去你那里，或通过其他方式。我祈祷，愿你的虔诚在教会中显明，愿你借着你的显现和你恩惠的话语，同时医治我和众人。若此事可能，那是最好的；你若决定采取其他途径，我也乐于接受。只求你应允我的恳求：明确告知我你的智慧所决定的事。\par

\SourceRecord{Translated by Blomfield Jackson.；Nicene and Post-Nicene Fathers, Second Series；Vol. 8.；Edited by Philip Schaff and Henry Wace.；Buffalo, NY: Christian Literature Publishing Co.,；1895.；<http://www.newadvent.org/fathers/3202059.htm>.}

% structure: p0001 source=h1 target=section reason=page-title

\section{\WorkTitle{第六十封书信}}

\Emph{\Person{圣巴西略}（\Place{凯撒勒雅}）}\par

\Emph{致其叔父\Person{额我略}}.\par

从前，我因见到我的兄弟而喜乐。为何不呢？既然他是我的兄弟，且是这般的一位兄弟？如今，我以同样的心情接待他前来探望我，我的亲情丝毫未减。愿天主决不让我如此感觉，以致忘记天性的联系，与亲近我的人为敌。我发现他的临在对我身体的疾病和灵魂的其他困扰是一种安慰；尤其是他带来的从阁下那里来的信，令我格外欣喜。我长久以来一直盼望这封信的到来，唯一的原因就是：我不愿我的生命中再添上一段亲族间争吵的悲伤插曲，这必然使仇敌喜悦、朋友忧愁，并使天主不悦，因为祂已经将完美的爱定为祂门徒的特征。因此，我确实有义务回复，恳切请求您为我祈祷，并像对待您的亲人一样，在一切事上照顾我。由于我缺乏信息，无法清楚理解所发生之事的意义，我断定应该接受您好心给予我的叙述的真实性。其余的事，如我们会面、合适的时间和便利的地点，就由您凭智慧来决定。如果阁下真的不嫌弃屈尊就我，与我交谈，无论您希望会见在他人面前还是在私下，我绝不反对，因为我已下定决心在爱中顺从您，并毫无例外地执行阁下为光荣天主所命令我的一切。\par

我没有迫使我可敬的兄弟必须亲口向您报告任何事情，因为先前场合他所说的话并未被事实证实。\par

\SourceRecord{Translated by Blomfield Jackson.；Nicene and Post-Nicene Fathers, Second Series；Vol. 8.；Edited by Philip Schaff and Henry Wace.；Buffalo, NY: Christian Literature Publishing Co.,；1895.；<http://www.newadvent.org/fathers/3202060.htm>.}

% structure: p0001 source=h1 target=section reason=page-title

\section{第六十一封书信}

\Emph{\Person{圣巴西略·凯撒勒雅}}\par

\Emph{致\Person{阿塔纳削}，亚历山大主教}。\par

我读过尊驾的信函，您在信中表达了对不幸的利比亚总督的痛心。我深感悲哀，自己的国家竟滋生并养育了如此恶习。我也悲哀，邻国利比亚因我们的罪恶而受苦，且被交付于一个其生命同时被残酷与罪行所标记之人的不人道。然而，这正合乎训道者的智慧：\Quote{祸哉，你这土地，当你的王是孩童；} \BibleRef{训 10:16}（一个更深重的麻烦）并且他的\Quote{王子}不在夜间\Quote{吃}，而是在正午纵饮，追逐他人之妻，比野兽更少理智。此人必定期待公义审判者的鞭笞，以他先前加诸圣人身上的同样报应回报他。根据尊驾的信函，我的教会已接到通知，他应被所有人视为可憎，断绝火、水和住所——如果对于如此着魔之人，全体一致的谴责还有作用的话。他的恶名已足够，阁下信函已广为传阅，我也必将其展示给他的所有亲友。确然，即使报应未如对待法老般立即临到他，日后也必将带给他沉重严厉的惩罚。\par

\SourceRecord{Translated by Blomfield Jackson.；Nicene and Post-Nicene Fathers, Second Series；Vol. 8.；Edited by Philip Schaff and Henry Wace.；Buffalo, NY: Christian Literature Publishing Co.,；1895.；<http://www.newadvent.org/fathers/3202061.htm>.}

% structure: p0001 source=h1 target=section reason=page-title

\section{\WorkTitle{书信第六十二}}

\Emph{\Person{圣巴西略·凯撒勒雅}}\par

\Emph{致帕尔纳索斯教会}。\par

遵循一项多年以来的古老习俗，并同时向你们呈现在天主内的爱德，那是圣神的果实，我如今，我敬虔的朋友们，致这封信给你们。我立即与你们一同感受你们在临到你们的事件中的悲伤，以及你们手中所关顾的事。关于这一切困扰，我只能说，我们被给予机会注视宗徒的训诫，不要忧伤，\BibleQuote{像那些没有希望的人一样。}\BibleRef{得前 4:13}我并不是说我们应当对所遭受的损失无动于衷，而是说我们不应屈服于忧伤，同时认为牧者在他结局中是有福的。他在高龄中逝世，并在主赐给他的极大尊荣中找到了安息。\par

至于未来，我有以下建议给你们。你们现在必须放下一切哀悼；你们必须醒悟；你们必须起来管理教会的必要事务；为使圣洁的天主垂顾祂自己的小小羊群，并按照祂自己的旨意赐给你们一位牧者，他能明智地牧养你们。\par

\SourceRecord{Translated by Blomfield Jackson.；Nicene and Post-Nicene Fathers, Second Series；Vol. 8.；Edited by Philip Schaff and Henry Wace.；Buffalo, NY: Christian Literature Publishing Co.,；1895.；<http://www.newadvent.org/fathers/3202062.htm>.}

% structure: p0001 source=h1 target=section reason=page-title

\section{第六十三书}

\Emph{圣巴西略·凯撒勒雅}\par

\Emph{致新凯撒勒雅总督}。\par

\Quote{智者，即使他居住在远方，即使我从未见过他，我也视之为友。}悲剧家欧里庇得斯如此说。因此，如果我虽从未有幸亲自与阁下会面，却以熟友自居，请勿将此视为空洞的恭维。普遍的声音，热烈宣扬您的普遍仁爱，在此正是友谊的促进者。事实上，自从我遇见十分可敬的\Person{厄尔丕迪乌斯}，我便同样认识了您，并完全被您吸引，如同与您长久相处并亲身经验了您的卓越品质。因为他不停地向我讲述您，逐一提及您的宽宏大量、崇高情操、温和举止、办事才干、智慧、愉快调和的尊严，以及口才。他在与我长谈中列举的其他所有细节，我无法在信中一一写出，否则会超出合理的篇幅。我怎能不爱这样的人？我怎能如此克制，不将我内心的感受高声宣告呢？那么，最卓越的先生，请接受我向您致意，因为这是出于真实而无伪的友谊。我厌恶一切谄媚的恭维。请将我列入您的朋友名单，并常常写信给我，使您在我面前出现，并在我缺席时安慰我。\par

\SourceRecord{Translated by Blomfield Jackson.；Nicene and Post-Nicene Fathers, Second Series；Vol. 8.；Edited by Philip Schaff and Henry Wace.；Buffalo, NY: Christian Literature Publishing Co.,；1895.；<http://www.newadvent.org/fathers/3202063.htm>.}

% structure: p0001 source=h1 target=section reason=page-title

\section{第六十四书}

\Signature{凯撒勒雅的圣巴西略}\par

\Emph{致\Person{赫西基乌斯}}。\par

从一开始，我就在许多方面与阁下的卓越之处有共通之处：您对文学的热爱，所有体验过的人都到处传扬，以及我与可敬的\Person{泰伦提乌斯}之间的旧谊。但自那最卓越之人——他对我来说是友谊所能要求的一切——我贤明的兄弟\Person{埃尔皮迪乌斯}遇见我，并向我讲述了您所有的美德（还有谁比他更能同时察觉一个人的美德并加以描述呢？）——他激起了我如此强烈的见您之愿，以至于我祈祷您有朝一日能来我久居之地探望我，使我不仅能听闻您的美德，更能亲身实际体验。\par

\SourceRecord{Translated by Blomfield Jackson.；Nicene and Post-Nicene Fathers, Second Series；Vol. 8.；Edited by Philip Schaff and Henry Wace.；Buffalo, NY: Christian Literature Publishing Co.,；1895.；<http://www.newadvent.org/fathers/3202064.htm>.}

% structure: p0001 source=h1 target=section reason=page-title

\section{第六十五封信}

\Emph{\Person{圣巴西略·凯撒勒雅}}\par

\Emph{致阿塔比乌}.\par

如果我继续坚持因年长而享有的特权，等待你主动通信，而你，我的朋友，却想更顽固地坚持你那不行动的恶谋，我们的沉默何时终了？但是，凡涉及友谊之处，在我看来，失败即是胜利，因此我完全准备让你优先，并退出这场坚持己见的竞赛。我首先提笔写信，因为我知道\BibleQuote{爱德包容一切……忍耐一切……不求己益}，因此\BibleQuote{永不止息}。在爱中服从近人者绝不会受辱。我恳求你，无论如何，今后彰显圣神最初和最大的果实——爱；除去你以沉默向我显示的愤怒之人的愠怒，恢复你心中的喜乐、与你同心合意的兄弟之间的和平，以及对主的教会持续安全的热情和焦虑。如果我不像正统教义的反对者倾覆和彻底毁灭教会那样，为教会付出同样艰辛的努力，你可以确信，没有什么能阻止真理被敌人扫除和毁灭，而我也会因此受到谴责，因为我没有以全部热情和急切，在相互同心和虔诚一致中，为教会的合一而表现出一切可能的焦虑。因此我劝你，从你的头脑中驱除你不需要与任何人共融的想法。将自己与兄弟的联系切断，不是行于爱中的人的特征，也不是满全基督诫命的行为。同时，我确实希望你在所有善意中考虑到，现在环绕我们四周的战争灾难有一天可能临到我们自己家门口，如果我们像其他所有人一样遭受暴行，我们将找不到任何人甚至同情我们，因为在我们繁荣的时刻，我们拒绝向受冤屈者施与我们的同情。\par

\SourceRecord{Translated by Blomfield Jackson.；Nicene and Post-Nicene Fathers, Second Series；Vol. 8.；Edited by Philip Schaff and Henry Wace.；Buffalo, NY: Christian Literature Publishing Co.,；1895.；<http://www.newadvent.org/fathers/3202065.htm>.}

% structure: p0001 source=h1 target=section reason=page-title

\section{\WorkTitle{第六十六封信}}

\Signature{凯撒勒雅的圣巴西略}\par

\Emph{致\Place{亚历山大}的主教\Person{阿塔纳削}}.\par

我敢肯定，没有人比阁下更为当前的状况——更确切地说，为教会的恶劣状况——感到痛苦；因为你将现在与过去相比，并认识到已发生了多大的变化。你很清楚，如果不阻止我们目睹的迅速恶化，很快就会无法阻止教会的彻底转变。如果教会的衰败在我看来如此可悲，那么——我常在孤独的默想中思量——一个亲身体验过主教会昔日的安宁和他们关于信德的一心一意的人，又会有何感受呢？既然阁下对此痛苦感受最深，我认为你的智慧更应该被强烈地推动去关心教会的事务。就我个人而言，以我有限的智力对时事的判断，我早已意识到东方教会得救的唯一途径在于获得西方主教们的同情。因为只要那些主教们愿意为旅居在我们这片区域的基督徒表现出他们在西方对一两个被定罪违反正统信仰的人所展示的同样热忱，我们的共同利益很可能会获得不小的益处，我们的君主们就会尊重人民的权威，而各地的平信徒也会毫不犹豫地跟随他们。但是，要实现这些目标，谁比阁下更有能力、更有智慧、更有远见呢？谁更敏锐地看到需要采取的必要行动？谁在推行有益政策方面更有实际经验？谁更深刻地感受到弟兄们的苦难？在整个西方，有谁比你的可敬白发更受尊荣？最可尊敬的父啊，请留下一些配得上你生命和品格的纪念。藉此一举，为你为真宗教所付出的无数努力加冕。请从你管辖下的圣教会派遣精通纯正道理的人到西方的主教们那里。向他们讲述我们遭受的困扰。建议某种解除之道。成为教会的撒慕尔。与受围困的人同忧。为和平祈祷。向主恳求恩惠，求祂给教会赐下和平的纪念。我知道信件在如此重要的事上多么无力；但你自己不需要别人的劝勉，正如最高贵的运动员不需要孩童的欢呼一样。并不是我在教导一个无知的人；我只是给一个已经振作的人新的推动力。至于东方其余的事务，或许你还需要更多帮助，我们必须等待西方人。但显然，安提约基雅教会的纪律取决于尊座能否控制一些人，使另一些人缄默，并藉和睦恢复教会的力量。没有人比你更清楚，你应当像所有明智的医生一样，从最要害的部位开始治疗；对于普世教会来说，有什么部位比安提约基雅更要害呢？只要安提约基雅恢复和谐，就没有什么能阻止她像健康的头一样，为整个身体提供健全。诚然，那座城的病症——它不但被异端者分裂，且被那些自称彼此同心的人撕碎——需要你的智慧和福音的同情。使那些分裂的部分再次合一，带来一个身体的和谐，唯独属于那一位——祂凭其不可言喻的德能，甚至使枯骨重新长出筋和肉。但主常藉那些堪当祂的人行其大能之事。我们再次相信，在这事上，如此重要的事务的职分也适宜于阁下，以致你将平息人民的暴风，消除党派优越，使众人因爱彼此服从，并将古时的力量归还给教会。\par

\SourceRecord{Translated by Blomfield Jackson.；Nicene and Post-Nicene Fathers, Second Series；Vol. 8.；Edited by Philip Schaff and Henry Wace.；Buffalo, NY: Christian Literature Publishing Co.,；1895.；<http://www.newadvent.org/fathers/3202066.htm>.}

% structure: p0001 source=h1 target=section reason=page-title

\section{第六十七书}

\\Emph{\\Person{凯撒勒雅的圣巴西略}}\par

\\Emph{致亚历山大的主教\\Person{亚大纳削}}.\par

在我先前的信中，我以为向尊下指出这一点就够了：安提约基雅圣教会中所有信仰纯正的信徒，都应被引导至和睦与合一。我的目的是要表明，那些现已分裂为几个派别的部分，应当在蒙主所爱的主教\\Person{默肋提}之下联合起来。如今，同一蒙爱的执事\\Person{多罗德}请求就这些事作更明确的陈述，因此我不得不指出，整个东方教会的祈祷，以及像我这样完全与他合一之人的热切愿望，就是看到他掌管主的教会。他是一位信仰无瑕可指的人；他的生活方式无比卓越；他可以说是整个教会身体的元首，其余的人都不过是脱节的肢体。因此，基于一切理由，其他人与他联合是必要且有益的，正如较小的溪流与较大的河流汇合一般。然而，对于其余的人，需要一定程度的处理，以适合他们的地位，并可能安抚民众。这符合您本人的智慧，以及您著名的果断和精力。然而，您的聪明才智丝毫没有忽略这样的事实：同一做法已经受到与您一致的西方人的赞同——这是我从有福的\\Person{西尔瓦诺}带给我的书信中得知的。\par

\SourceRecord{Translated by Blomfield Jackson.；Nicene and Post-Nicene Fathers, Second Series；Vol. 8.；Edited by Philip Schaff and Henry Wace.；Buffalo, NY: Christian Literature Publishing Co.,；1895.；<http://www.newadvent.org/fathers/3202067.htm>.}

% structure: p0001 source=h1 target=section reason=page-title

\section{第六十八封信}

\Emph{\Person{圣巴西略}}\par

\Emph{致\Person{梅肋提乌斯}，\Place{安提约基雅}主教。}\par

我原想将可敬的执事兄弟\Person{多罗德}留在身边，直到谈判结束，以便通过他向阁下禀报一切细节。但日复一日，拖延渐长；如今，当我想到一个在我困境中尽可能可行的方案时，便派他前往圣德座前，亲自向您报告所有情况，并呈上我的备忘录，以便若我所想之事对您有益，阁下能促其实现。简言之，大家认为最好让我们的这位兄弟\Person{多罗德}前往\Place{罗马}，促使一些\Place{意大利}人起航海路来访，以避免所有可能设置障碍的人。我这样做的理由是，我发现那些在皇帝面前极有权势的人，既不愿也不能向他提出任何关于被流放者的建议，他们只庆幸教会没有遭遇更糟的事。若我的计划在您的明智看来也可取，请您惠予起草信函并指示备忘录中他应详述的要点，以及他应联系何人。为使您的信函具有分量和权威，请附上所有与您意见相同的人，即使他们不在当地。此处一切未定；\Person{欧伊丕乌斯}已抵达，但迄今未有任何表示。然而，他和来自\Place{亚美尼亚四城}及\Place{基里基雅}的同道们正威胁要举行一次骚乱的集会。\par

\SourceRecord{Translated by Blomfield Jackson.；Nicene and Post-Nicene Fathers, Second Series；Vol. 8.；Edited by Philip Schaff and Henry Wace.；Buffalo, NY: Christian Literature Publishing Co.,；1895.；<http://www.newadvent.org/fathers/3202068.htm>.}

% structure: p0001 source=h1 target=section reason=page-title

\section{Letter 69}

\Emph{\Person{圣巴西略（凯撒勒雅）}}\par

\Emph{致\Person{亚大纳削}，\Place{亚历山大里亚}主教。}\par

1. 随着时光流逝，我长久以来对您圣德所持的看法不断得到证实；更确切地说，这一看法因每日发生的事件而愈加坚定。大多数人确实满足于各守本分，但您并非如此；您对所有教会的关切，丝毫不亚于您对那由我们共同的主特别托付给您的教会所怀的关心；因为您从不间断地说话、劝勉、写信、派遣使者，这些使者随时就每一紧急情况提出最佳建议。如今，从您神圣的圣职人员行列中，您派遣了可敬的弟兄伯多禄前来，我满心欢喜地接待了他。我也赞同他此行的良好目标，即按照阁下您的命令，在他遇到对立之处实现和解，以合一取代分裂。为了对此事正在采取的行动有所贡献，我认为最合适的开端莫过于向阁下您求助，如同向一切的首脑和元首，并将您视为这项事业的顾问和指挥官。因此，我决定派遣我们的弟兄、执事多罗特乌斯前往尊前，他属于那在极其可敬的默肋提乌斯主教治下的教会；他既是正统信仰的积极支持者，又热切渴望看到教会的和平。我希望结果将是，他遵循您的建议（您因年龄和处理事务的经验，并因您比别人更多地拥有圣神的助佑，所以较不易失败），从而尝试实现我们的目标。我相信您会欢迎他，并会以友善的目光看待他。您将以祈祷的帮助坚强他；您将给他一封信作为路上的供应；您将赐给他一些您身边的正直忠信之人作为同伴；这样您将加速他踏上前方的道路。在我看来，最好也寄一封信给罗马主教，恳请他审视我们的状况；并且，由于以全体主教会议决议从西方派遣代表有困难，建议他在这件事上行使他个人的权威，挑选合适的人来承担旅途的劳苦——这些人也应以性格的温和与坚定来纠正我们这里的不驯者；能恰如其分、合宜地说话，并完全了解里米尼之后为推翻那里的强暴措施所做的一切。我建议，在无人知晓的情况下，他们应尽可能不引人注意地由海路前来，以避开和平的敌人。\par

2. 此处有些人也极力主张（连我自己也显然觉得十分必要），他们应驱除马尔塞禄的异端，因为它有害、有损，并违反纯正的信仰。因为直到现在，在他们所写的所有信函中，他们始终彻底诅咒声名狼藉的亚略，并把他从教会中逐出。但他们并未归咎于马尔塞禄；马尔塞禄提出了一种与亚略截然相反的异端，不敬地攻击独生天主性的存在本身，并错误地理解了\Quote{圣言}这个名称。他承认独生子被称为\Quote{圣言}，是因为在需要和适当的时候出现，但他声称他又回到了他出来的地方，在他出来之前不存在，回去之后也没有位格。我手中那些含有他不义著作的书卷，就是我所说的证据。然而，他们从未公开谴责过他，并因此而有罪，因为他们起初不认识真理，就接纳他进入教会的共融。当前的事态尤其需要提请他注意，以便寻求机会的人无法得逞，因为正直的人与您的圣德联合，所有在真信仰中跛足的人都会公开显明；这样我们就能知道谁在我们一边，而不致如在夜战中，无法分清敌友。只是我恳求您，将我提到的执事用最早可能的邮船遣回，以便我们祈求的至少一些目标能在来年实现。然而，有一件事，即使我不提，您也完全明白，并且我相信会留意：当他们来到时，若天主愿意，他们不可在教会中引发分裂；即使他们发现有些人因个人原因彼此不和，他们也应竭尽全力联合所有意见相同的人。因为我们应当优先考虑和平的利益，并首先要关注安提约基雅的教会，以免它健全的部分因个人分歧而患病。但您自己将对这些事给予更全面的关注，一旦蒙天主降福，您发现每个人都把确保教会和平的责任托付给您。\par

\SourceRecord{Translated by Blomfield Jackson.；Nicene and Post-Nicene Fathers, Second Series；Vol. 8.；Edited by Philip Schaff and Henry Wace.；Buffalo, NY: Christian Literature Publishing Co.,；1895.；<http://www.newadvent.org/fathers/3202069.htm>.}

% structure: p0001 source=h1 target=section reason=page-title

\section{\WorkTitle{Letter 70}}

\Emph{\Person{凯撒勒雅的圣巴西略}}\par

\Emph{无收信人地址。}\par

更新古老爱德的律法，并使那因岁月而枯萎的、基督属天且救恩的恩赐——我指的是教父们的和平——重新恢复活力，这工作对我而言确实是必要且有益的，但也是令人愉快的，我相信它在你那爱基督的禀性中也是如此。因为还有什么比看到那被距离如此遥远分隔的众人，因着爱的结合而在基督身体内成为一和谐肢体更为可喜的呢？整个东方（我用这名指代从\Place{伊利里亚}到\Place{埃及}的所有地区），可敬的教父，正被一场可怕的风暴和狂涛所搅动。那由真理的敌人\Person{亚略}所播下的古老异端，如今已大胆而无耻地重现了。如同一毒根，它正结出致命的果实，并且得势。其原因在于，在每个地区，正统教义的捍卫者都因诽谤和迫害而被逐出他们的教会，事务的管理被交给了那些俘虏较为单纯弟兄灵魂的人。我曾将你仁慈的到访视为我们困难唯一的可能解决之道。过去我常因你非凡的爱德而得到安慰；虽曾一度因那令我心喜的报告——我们将蒙你到访——而心中欢愉，但我失望了，因此不得不以书信恳请你被触动来帮助我们，并派遣一些与我们同心的人，要么劝解分歧者，使天主的教会重归友好合一，要么至少让你更清楚地看到谁应对我们这不安的境况负责，以便你将来明白与谁共融是适宜的。我这样做绝非提出什么新的请求，而只是要求那些在我们之前、蒙福且为天主所爱的人，尤其是你自己，所惯常做的事。因为我清楚记得从我们教父被问及的回答以及仍保存在我们中的文献中得知，那位在你教区中因信仰纯正及其他一切美德而著称的荣耀且蒙福的主教\Person{狄约尼削}，曾以书信访问我的\Place{凯撒勒雅}教会，并以书信劝告我们的教父，并派遣人赎回我们被俘的弟兄。但现在我们的境况更加痛苦和阴郁，需要更谨慎的处理。我们所哀悼的不仅仅是地上建筑的倒塌，而是教会的被俘；我们眼前所见的不仅是身体的奴役，而是灵魂被掳——日日被异端捍卫者所行。若你至今仍不动心援助我们，不久所有人将落入异端统治之下，你将找不到任何可以伸手相助的人了。\par

\SourceRecord{Translated by Blomfield Jackson.；Nicene and Post-Nicene Fathers, Second Series；Vol. 8.；Edited by Philip Schaff and Henry Wace.；Buffalo, NY: Christian Literature Publishing Co.,；1895.；<http://www.newadvent.org/fathers/3202070.htm>.}

% structure: p0001 source=h1 target=section reason=page-title

\section{\WorkTitle{第七十一书}}

\Emph{\Place{凯撒勒雅}的\Person{圣巴西略}}\par

\Emph{\Person{巴西略}致\Person{额我略}}.\par

1. 我已收到圣座经由最可敬的\Person{赫肋纽斯}兄弟送来的信函，凡你所暗示之事，他已原原本本地告诉我。对于此事我作何感受，你绝无可疑。然而，我已决定我对你的情谊应超越一切痛苦，无论这痛苦如何，我都已照我所当行的方式接受了它。我祈求神圣的天主，使我在余下的时日或时刻中，对待你的态度能与往昔一样谨慎无误；在此期间，我的良心告诉我，我从未在大小事上亏缺于你。[但那个自夸如今才开始窥视基督徒生活的人，以为与我有所来往便可获得些许声誉，竟捏造他未曾听闻之事，讲述他从未经历过之事，这丝毫不令人惊讶。令人惊讶而非同寻常的是，他竟能使我在\Place{纳西盎}的弟兄中的几位挚友听信于他；不仅听信，而且似乎接受了他说的话。从许多方面看，诽谤者竟是这般品性，而我竟成为受害者，这或许令人惊讶；但这些动荡的时代已教导我们耐心忍受一切。因我的罪，比这更甚的轻慢久已成为我的家常便饭。我从未向此人的弟兄们显明我对天主的看法，如今我也不作回答。那些不为长久经验所说服的人，恐怕也不会为短短一封信所说服。若前者已足够，就让诽谤者的指控如同虚妄之言。但若我任由放肆的口舌和未受教的心灵随意议论他们所欲议论之人，而自己却始终竖耳恭听，那么我将不是唯一听到他人谈论的人；他们也将不得不听到我的回应。]\par

2. 我深知导致这一切的原因，并已想尽种种话题试图阻止；但如今我已厌倦这个话题，不再多言，我指的是我们那少许的交往。因为若我们彼此坚守旧日承诺，并适当顾及教会对我们的要求，我们本该一年中有大部分时间共处；这样就不会给这些诽谤者留下可乘之机。恳请你切勿理会他们；让我说服你前来帮我从事劳作，尤其是在我正与那个攻击我的人的争战中。你的出现本身就会使他住手；一旦你向这些扰乱我们家园之人表明，你将因天主的降福，亲自领导我们一方，你便会瓦解他们的结党，封闭每一个不义之口，那说那不义之言反对天主的口。这样，事实将表明，谁是你追随者中的良善者，谁是蹒跚而行、怯懦地背叛真理之言者。然而，若教会遭背叛，那么我将无需费心用言语为自己辩护，因为那些尚未学会衡量自己的人，对我的看法自是如此。或许，藉天主的恩宠，不日之内，我就能以实际行为驳斥他们的诽谤；因为我似乎很快将为真理遭受比平时更大的苦楚；我所能期望的最好结果便是流放；若此希望落空，毕竟基督的审判席已近在咫尺。[你若为教会的缘故要求会面，我随时准备前往你所邀请的任何地方。但若仅是为驳斥这些诽谤，我实无时间予以回应。]\par

\SourceRecord{Translated by Blomfield Jackson.；Nicene and Post-Nicene Fathers, Second Series；Vol. 8.；Edited by Philip Schaff and Henry Wace.；Buffalo, NY: Christian Literature Publishing Co.,；1895.；<http://www.newadvent.org/fathers/3202071.htm>.}

% structure: p0001 source=h1 target=section reason=page-title

\section{第七十二封信}

\Emph{\Person{圣巴西略·凯撒勒雅}}\par

\Emph{致\Person{赫西基乌斯}}。\par

我知道你对我怀有感情，并对一切善事充满热忱。我极切望安抚我心爱的儿子\Person{卡利斯忒涅斯}，心想若能得你相助，或许更容易达成目的。\Person{卡利斯忒涅斯}对\Person{欧斯托基乌斯}的举止甚感不快，且他确有充分的理由如此。他指责\Person{欧斯托基乌斯}的家人对他傲慢无礼、施以暴力。我正恳请他息怒，满足于已使那些无礼之徒及其主人受到惊吓，并予以宽恕，了结这场纷争。如此将产生双重结果：他将赢得人的尊敬，并在天主面前得蒙称赞，只要他能将宽容与威胁相结合。若你与他有友谊和亲密关系，请向他提出此请求；若你知道城中有任何能影响他的人，请让他们与你一同行动，并告诉他们这将使我特别欣慰。一旦他的使命完成，即速将那位执事遣回。那些人既已逃到我这里避难，若我无法对他们有所帮助，我将感到羞愧。\par

\SourceRecord{Translated by Blomfield Jackson.；Nicene and Post-Nicene Fathers, Second Series；Vol. 8.；Edited by Philip Schaff and Henry Wace.；Buffalo, NY: Christian Literature Publishing Co.,；1895.；<http://www.newadvent.org/fathers/3202072.htm>.}

% structure: p0001 source=h1 target=section reason=page-title

\section{\WorkTitle{第七十三封书信}}

\Emph{\Person{圣巴西略·凯撒勒雅}}\par

\Emph{致\Person{卡利斯忒涅斯}}。\par

一、当我读过你的来信后，我首先感谢天主：因为一个渴望尊敬我的人向我致意，我确实非常珍视与高尚之人的交往；其次，我因被记挂而感到欣喜。因为来信是记挂的记号；当我收到你的信并了解其内容后，我惊讶地发现，正如大家所同意的那样，你对我表达了如同儿子对父亲般的尊敬。一个人在愤怒和恼怒之中，急于惩罚那些惹恼他的人，却放下了一大半的激烈情绪，并授权我来决定此事，这使我在灵性上如同得到儿子一般喜乐。作为回报，我除了为你祈求一切祝福之外，还有什么可做的呢？愿你成为朋友的喜悦、敌人的畏惧、众人尊敬的对象，以至于任何亏欠你责任的人，当他们得知你的温良时，只会自责不该错待像你这样品性的人！\par

二、我很想知道你的善心有何目的，即命令将仆人们送到他们犯下无礼行为的地点。如果你亲自前来，亲自执行因过失应得的惩罚，那些奴隶将会在那里。如果你已下定决心，还有其他可能吗？只是我不知道，如果我没能免除这些年轻人的惩罚，我还会得到什么进一步的恩惠。但如果事务在路上耽搁了你，谁在那里接应他们？谁代替你惩罚他们？但如果你已决定亲自与他们见面，且这已确定，就命令他们在\Place{萨西玛}停下，并在那里展示你的温良和宽宏大量。在将攻击你的人置于你的权力之下，并以此向他们表明你的尊严不可轻慢之后，就让他们毫发无损地离开，正如我在上一封信中恳求你的那样。这样你将给予我恩惠，并从天主那里得到善行的回报。\par

三、我这样说，不是因为事情本应就此了结，而是对你激动情绪的让步，并担心你的一些愤怒可能仍然未消。当一个人的双眼发炎时，最柔和的敷料也显得疼痛，我担心我所说的反而会激怒你，而非安抚你。真正最合宜、为你带来极大荣誉、也为我及我的朋友们和同时代人带来不小荣光的事，就是你将惩罚留给我。尽管你发誓要按照法律规定的处决他们，但我的责备作为惩罚并不逊色，神圣法律也不比世上现行的法律更不重要。但他们有可能通过在这里接受我们法律的惩戒——这也正是你自身的救恩所在——来解除你的誓言，并承受与他们的过失相称的惩罚。\par

但我的信又一次写得太长了。我极其恳切地想要说服你，无法忍受遗漏任何我想到的恳求，并且我十分担心，如果我未能说出所有能表达我意思的话，我的恳求会落空。现在，教会真实而尊贵的儿子，请证实我对你所怀的希望；请证实众人一致对你和善与温良的一致见证。命令那位士兵立即离开我；他现在已经尽可能地令人厌烦和粗鲁；他显然更愿意不给你带来任何烦恼，而不是使这里的所有人成为他的密友。\par

\SourceRecord{Translated by Blomfield Jackson.；Nicene and Post-Nicene Fathers, Second Series；Vol. 8.；Edited by Philip Schaff and Henry Wace.；Buffalo, NY: Christian Literature Publishing Co.,；1895.；<http://www.newadvent.org/fathers/3202073.htm>.}

% structure: p0001 source=h1 target=section reason=page-title

\section{书信第七十四}

\Person{圣巴西略·凯撒勒雅}\par

\Person{致马丁尼亚努斯}。\par

你可曾想过，我多么珍视我们再次相见的喜悦？能与你共度更长的时光，享受你所有的美德，那是多么令人愉快的事！如果见识众多人的城邦、了解众多人的习俗能证明一个人的教养，那么我相信，在你的陪伴下，这种教养能迅速获得。因为，逐一认识许多人，与认识一位集所有经验于一身的人，有何区别？我要说，后者具有莫大的优越性：它使我们毫不费力就获得对善与美的事物的认识，并将美德的教育置于我们伸手可及之处，纯洁无瑕，不受任何邪恶的沾染。谈及高尚的事迹、值得传颂的言语、超凡入圣之人的制度？这一切都珍藏在你心灵的宝库中。因此，我不会祈祷只像阿尔基努斯听奥德修斯那样倾听你一年，而是愿用我的一生来倾听；为此，我愿祈祷生命长久，即使我的境况并不轻松。那么，我现在为何写信而不亲自来见你？因为我的祖国在苦难中不可抗拒地召唤我回到她身边。朋友，你知道她所受的痛苦。她被真正的狂女——魔鬼——撕成碎片，如同彭透斯一般。他们正在分割她，一而再地分割，如同无知的庸医，使伤口更糟。她遭受这种分割，我唯有如同照料病人一样照料她。凯撒勒雅的居民紧急来信向我求助，我必须前往，并非以为我能有所帮助，而是为了避免疏忽的责备。你知道人在困境中多么容易抱有希望，也多么容易吹毛求疵，总是将他们的烦恼归咎于未做的事。\par

正因如此，我本应来见你，向你倾诉我的想法，或者更确切地说，恳求你考虑一些与你智慧相称的有力措施；不要对我那屈膝求救的祖国视而不见，而要亲自前往宫廷，以你特有的胆量，让他们明白他们不能认为自己拥有两个行省而非一个。他们并没有从世界的其他地方引入第二个行省，而是如同一个马或牛的主人，将其砍成两段，然后以为有了两个，而非一个——实际上既没有变成两个，反而毁掉了原有的一个。告诉皇帝和他的大臣们，他们并非以此方式扩展帝国，因为权力不在于数量，而在于状态。我相信，现在人们忽视了事态的发展，有些人可能是因为不了解真相，有些人是因为不愿说不中听的话，有些人则是因为事不关己。最有益且符合你高尚原则的做法是，如有可能，你亲自面见皇帝。如果这既因季节又因你的年龄（如你所说，懒惰与年龄相伴）而难以实现，那么你至少可以毫无困难地写信。你若是如此以书信援助我们的祖国，首先你将因知道自己已竭尽全力而心满意足；其次，即使只是表面上的同情，也会给病人带来莫大的安慰。但愿你能亲自来到我们中间，亲眼目睹我们的悲惨境遇！那么，或许在眼前事实的激励下，你能说出与你慷慨大度和凯撒勒雅的苦难相称的话语。但请不要怀疑我所说的。确实，我们需要一个西蒙尼德斯或其他类似的诗人，来根据实际经历哀叹我们的苦难。但为何提到西蒙尼德斯？我更应提及埃斯库罗斯，或任何其他以类似言辞描述巨大灾难并放声哀号的人。\par

如今，我们再也没有聚会，没有辩论，没有智者在广场的集会，没有一切曾使我们的城市闻名的事物。如今，在我们的广场上，出现一位博学多识或口若悬河的人，比之古时的雅典出现一个手带罪孽印记或双手污秽的人更为罕见。取而代之的是马萨革泰人和西徐亚人带来的粗野习气。只听得见讨价还价者、交易失败者和受鞭笞者的喧嚣声。两旁的柱廊回荡着凄厉的回声，仿佛它们正发出自然而恰当的声音，为所发生的一切呻吟。我们的痛苦使我们无暇顾及紧锁的体育馆和没有火把的夜晚。危险不小，因为我们的官员被撤走，一切可能像支柱倒塌一样轰然坍塌。有什么言辞能充分描述我们的灾难？有些人已流亡他乡，相当一部分元老院成员——且并非最不重要的——选择了永久流放，而不愿前往波丹杜斯。当我提到波丹杜斯时，请想象斯巴达的塞达斯，或你曾见过的任何天然的深渊，那里散发着有毒的蒸气，有人不由自主地称之为卡戎泉。请设想波丹杜斯的邪恶与此地相当。因此，三分之一的居民离乡背井，处于流亡之中，妻儿老小和家园尽弃；三分之一像俘虏一样被带走，城中大多数优秀人物，成为朋友们眼中凄惨的景象，实现了敌人的祈祷——如果有人曾祈求如此可怕的诅咒降在我们头上的话。还有三分之一留下：他们既无法忍受老友的遗弃，又无法自给自足，不得不痛恨自己的生命本身。\par

这就是我恳求你，用你那独特的口才和你的生活方式所赋予你的公义的直言不讳，到处宣扬的事情。要明确陈述一点：除非当权者很快改变他们的计谋，否则他们将发现没有人可以对其行使仁慈了。你要么对国家有所助益，至少也要像\Person{梭伦}那样行事：当他在\Place{雅典卫城}被占领时无法保卫被遗弃的同胞，便穿上盔甲，坐在城门前，以此表明他不赞同所发生的事。我深信，即使目前可能有人不赞同你的建议，但不久之后，当他们看到事件与你的预言相符时，他们会因你的仁爱和远见而极大地称赞你。\par

\SourceRecord{Translated by Blomfield Jackson.；Nicene and Post-Nicene Fathers, Second Series；Vol. 8.；Edited by Philip Schaff and Henry Wace.；Buffalo, NY: Christian Literature Publishing Co.,；1895.；<http://www.newadvent.org/fathers/3202074.htm>.}

% structure: p0001 source=h1 target=section reason=page-title

\section{书信七十五}

\Emph{\Person{圣巴西略·凯撒勒雅}}\par

\Emph{致\Person{阿布尔吉乌斯}。}\par

你有许多品质使你超乎常人，但没有比你对你祖国的热忱更显着的特点。因此，你虽已升到如此高度，在整个世界闻名遐迩，却仍对生养你的土地给予正义的回报。然而，你的母城，那生育你、哺育你的城，如今竟落入了古代传说般难以置信的境况。造访\Place{凯撒勒雅}的人，即便是最熟悉她的人，也认不出她现在的样子；她竟如此突然地沦为彻底的荒废，她的许多行政长官先前已被调离，如今几乎所有官员都被转移到了\Place{波丹杜斯}。余下的人，如同从肢体上被割下的残肢，自己也陷入了完全的绝望，并造成了如此普遍的沮丧，以至于城中人口已很稀少；这地方看起来像一片沙漠，对所有爱她的人是一个悲惨的景象，对长期图谋我们倾覆的人则是欢喜和鼓励的缘由。那么，谁会伸手援助我们？谁会怜悯我们的信德而洒泪？你曾同情过一座同样处于困境的异乡城市，难道你的仁厚之德不会为生养你的母城感到伤痛吗？你若有什么影响力，请在当前的急需中显示出来。你当然从天主那里得到极大的帮助，祂从未舍弃你，并赐予你许多祂慈爱的明证。只要你愿意为我们出力，并使用你所拥有的一切影响力来支援你的同胞。\par

\SourceRecord{Translated by Blomfield Jackson.；Nicene and Post-Nicene Fathers, Second Series；Vol. 8.；Edited by Philip Schaff and Henry Wace.；Buffalo, NY: Christian Literature Publishing Co.,；1895.；<http://www.newadvent.org/fathers/3202075.htm>.}

% structure: p0001 source=h1 target=section reason=page-title

\section{第七十六封书信}

\Emph{凯撒勒雅圣巴西略}\par

\Emph{致大臣索弗罗尼乌斯}。\par

降临在我们本城的灾祸之巨，似乎迫使我亲自前往朝廷，向阁下以及所有处理政务的要员陈述凯撒勒雅所处的悲惨境地。然而，我因体弱多病并需照料各教会而滞留于此。因此，我急忙以书信向主上禀告我们的困苦，并使您知晓：从未有船只因狂风怒海而如此突然沉没，从未有城市因地震而倾覆或因洪水而湮灭，如此迅速地从眼前消失——就像我们这座城市，被这新宪政吞噬，已彻底毁灭。我们的不幸已成传说，我们的制度已成往事；所有高级文官因市政官的遭遇而绝望，弃城离乡，漂泊于田野。因此，政务运作中断，这座曾因博学之士和富裕城镇中众多人才而自豪的城市，如今呈现出极不体面的景象。我们在患难中只剩下唯一的安慰，就是向阁下倾诉我们的不幸，并恳求您，若可能，向匍匐在您面前的凯撒勒雅伸出援助之手。至于您如何能够帮助我们，我自己无法说明；但我确信，以您的智慧，发现方法并不困难，而借着天主赐予您的权柄，使用这些方法也并不难。\par

\SourceRecord{Translated by Blomfield Jackson.；Nicene and Post-Nicene Fathers, Second Series；Vol. 8.；Edited by Philip Schaff and Henry Wace.；Buffalo, NY: Christian Literature Publishing Co.,；1895.；<http://www.newadvent.org/fathers/3202076.htm>.}

% structure: p0001 source=h1 target=section reason=page-title

\section{第七十七封书信}

\Person{凯撒勒雅的圣巴西略}\par

\Emph{无题：关于\Person{特拉西乌斯}}.\par

我们从伟大的\Person{特拉西乌斯}的治理中确实获得了一件好事，那就是你经常来访我们。现在，哎呀！我们失去了我们的总督，连这个好事也被剥夺了。但既然天主曾赐予我们的恩赐保持不变，而且，虽然我们身体分离，却藉记忆固定在我们各人的灵魂中，让我们不断写信，互相传达我们的需要。我们在目前这样做也很恰当，此时风暴暂时休战。我相信你不会离开那令人钦佩的\Person{特拉西乌斯}，因为我认为分担他的巨大忧虑是非常合宜的，而且我高兴你有机会既看到朋友，也被他们看到。关于许多事我有很多话要说，但推迟到我们见面时，因为，我认为，将如此重要的事托付给书信几乎是不安全的。\par

\SourceRecord{Translated by Blomfield Jackson.；Nicene and Post-Nicene Fathers, Second Series；Vol. 8.；Edited by Philip Schaff and Henry Wace.；Buffalo, NY: Christian Literature Publishing Co.,；1895.；<http://www.newadvent.org/fathers/3202077.htm>.}

% structure: p0001 source=h1 target=section reason=page-title

\section{第七十八封书信}

\Emph{\Person{凯撒勒雅的圣巴西略}}\par

\Emph{无题署，为\Person{厄尔彼底乌斯}求情。}\par

我没有忽视你对我们可敬的朋友\Person{厄尔彼底乌斯}所表现出的关怀；你以你惯常的智慧给了\OriginalTerm{总督}{prefect}一个展示他仁慈的机会。我现在写信请求你的是，使这一恩惠圆满，并建议\OriginalTerm{总督}{prefect}通过一项特别命令，将那位对公共事务充满一切可能的关心的人安置在我们的城中。因此，你将有许多令人钦佩的理由向\OriginalTerm{总督}{prefect}说明，命令\Person{厄尔彼底乌斯}留在\Place{凯撒勒雅}。无论如何，你不需要我来教导你，因为你自己非常清楚，现状如何，以及\Person{厄尔彼底乌斯}在行政管理中的能力。\par

\SourceRecord{Translated by Blomfield Jackson.；Nicene and Post-Nicene Fathers, Second Series；Vol. 8.；Edited by Philip Schaff and Henry Wace.；Buffalo, NY: Christian Literature Publishing Co.,；1895.；<http://www.newadvent.org/fathers/3202078.htm>.}

% structure: p0001 source=h1 target=section reason=page-title

\section{第七十九书}

\Emph{凯撒勒雅的圣巴西略}\par

\Emph{致\Person{塞巴斯特亚主教欧斯塔修}。}\par

甚至在我收到您的信函之前，我就知道您愿意为每个人承受怎样的困难，特别是为我这个卑微的人，因为在这场斗争中我暴露无遗。所以当我从可敬的\Person{厄肋乌西尼}手中收到您的信函，并亲眼看到他本人时，我赞美天主，因祂赐予我这样一位斗士和同伴，在藉着圣神的助佑为真宗教而奋斗的过程中。请尊贵的您知晓，至今我承受了一些来自高级官员的攻击，而且并非轻微；总督和高级内侍都同情地为我对手求情。但是，到目前为止，我靠着天主的仁慈，没有被动摇，这仁慈藉着圣神赐予我助佑，并透过祂坚固我的软弱。\par

\SourceRecord{Translated by Blomfield Jackson.；Nicene and Post-Nicene Fathers, Second Series；Vol. 8.；Edited by Philip Schaff and Henry Wace.；Buffalo, NY: Christian Literature Publishing Co.,；1895.；<http://www.newadvent.org/fathers/3202079.htm>.}

% structure: p0001 source=h1 target=section reason=page-title

\section{第八十封信}

\Emph{\Place{凯撒勒雅}的\Person{圣巴西略}}\par

\Emph{致\Place{亚历山大里亚}主教\Person{亚大纳削}}.\par

众教会之疾病越严重，我们就越发转向您的尊贵，相信您的护佑是我们患难中唯一的慰藉。因您祈祷的能力，以及您在紧急情况下所提供的最佳建议的知识，凡对您的尊贵稍有认识——无论传闻或亲历——都相信您能拯救我们脱离这场可怕的风暴。因此，我恳求您不要停止为我们的灵魂祈祷，并不断以书信激励我们。如果您知道这些书信对我们何等有益，您就不会错过任何一个写信的机会。只要我能藉着您的祈祷，被认为配得见您、享受您的美德，并在我的生命历程中增添与您真正伟大的灵魂的会面，那么我就会确信，我已从天主的仁慈中获得了与我一生所有苦难相等的慰藉。\par

\SourceRecord{Translated by Blomfield Jackson.；Nicene and Post-Nicene Fathers, Second Series；Vol. 8.；Edited by Philip Schaff and Henry Wace.；Buffalo, NY: Christian Literature Publishing Co.,；1895.；<http://www.newadvent.org/fathers/3202080.htm>.}

% structure: p0001 source=h1 target=section reason=page-title

\section{\WorkTitle{第八十一封书信}}

\Signature{凯撒勒雅的圣巴西略}\par

\Emph{致依诺森主教}。\par

我很高兴收到你深情寄来的信；但同样令我悲伤的是，你加给了我一个超出我能力的责任重负。我如此远离，怎能承担如此重大的职责？只要教会拥有你，它就如靠在它自己的支柱上。若主愿意在你生命的事上有所安排，我能从我们这里派谁去接管对弟兄们的照料，且能与你本人一样受尊敬呢？你在信中所表达的是一个非常明智且恰当的愿望：当你还在世时，就能见到在你之后被预定来引导主所选羊群的继承人（如同真福梅瑟，他既愿望又见到了）。由于这个地方重大而著名，你的工作有伟大而广泛的名声，且时代艰难，因连续的风暴和浪潮攻击教会，需要一个并非无足轻重的引导者，我顾及自己的灵魂，认为不可轻率处理此事，尤其当我牢记你写信的措辞时。因为你说，你指责我忽视教会，要在主面前反对我。那么，为了不与你争论，反而让你在我于基督面前所作的辩护中站在我这边，我在环视城中司铎集会之后，选定了那位极可敬重的器皿，即真福赫摩革内的后裔，他在大公会议上写下了伟大而不可战胜的信经。他是教会的一位司铎，多年资历，性格坚定，精通教会法规，信德准确，至今过着节制和苦修的生活，虽然他严峻生活的严厉已制服了肉身；他是一个贫穷的人，在此世没有资源，以至于甚至没有基本的食粮，而是靠双手劳动与同住的弟兄们维持生计。我打算派他去。那么，如果这就是你想要的那种人，而不是一个仅适合被派遣和履行此世普通职责的年轻人，请尽快写信给我，我就可以给你派去这人，他是天主所特选的，适应眼前的工作，受到所有遇见他的人的尊敬，并以柔和教导所有与他意见不同的人。我本可以立刻派他去，但由于你自己先于我要了一个品格高尚、为我所爱但远不及我所指明的那位的人，我想让你知道我对这件事的想法。因此，如果这就是你想要的那种人，或者派一位弟兄在斋期去接他，或者如果你无人能来我这里，请写信告诉我。\par

\SourceRecord{Translated by Blomfield Jackson.；Nicene and Post-Nicene Fathers, Second Series；Vol. 8.；Edited by Philip Schaff and Henry Wace.；Buffalo, NY: Christian Literature Publishing Co.,；1895.；<http://www.newadvent.org/fathers/3202081.htm>.}

% structure: p0001 source=h1 target=section reason=page-title

\section{第八十二封书信}

\Emph{圣巴西略·凯撒勒雅}\par

\Emph{致\Person{圣亚大纳削}，\Place{亚历山大城}主教。}\par

当我将目光转向世界，察觉一切追求美善的努力都如同戴着脚镣行走的人一样受阻时，我便陷入失望。但随后我将目光转向尊驾，想起我们的主已任命您为教会疾病的医师，于是我恢复精神，从绝望的消沉中振作，期盼更好的事。如您的智慧深知，整个教会已濒临崩溃。您用心灵之眼审视各方，如同从高塔上瞭望，正如海上许多船只一起航行，因波涛汹涌而互相碰撞，船只失事有的源于外界狂风怒海，有的源于水手们的混乱、互相阻碍。描绘这幅景象便足够了，无需多言。您的智慧无需更多，而当前的局势也不容我自由言说。在这样的风暴中，哪里能找到有能力的舵手？谁配唤醒主斥责风和海？除了那位自幼为真道打美好仗的人，还能有谁？如今，我们中所有健全的人确实正朝着与同道共融合一的方向迈进，因此我们满怀信心地恳求您寄给我们一封信，指示我们当行的事。他们希望通过这种方式有一个促进合一的交流的开端。或许，当您回忆往事时会对他们有所怀疑，因此，最蒙天主所爱的父，请这样做：将致主教们的信函，要么通过您在\Place{亚历山大城}所信任的人，要么通过我们的兄弟、执事\Person{多罗特乌斯}之手寄给我；我收到这些信函后，将等到获得主教们的答复才转交；否则，让我\Quote{永远承担责任。}\BibleRef{创 43:9} 这事本当更令那位最初对他父亲说出此话的人感到畏惧，而我现在对属灵父亲说同样的话，更应如此。然而，如果您完全放弃这个希望，至少让我在行事上不受责备，因为我受托传递信息和调解，完全是出于真诚和纯朴，出于对和平的渴望，以及所有在主内同道之间的彼此往来。\par

\SourceRecord{Translated by Blomfield Jackson.；Nicene and Post-Nicene Fathers, Second Series；Vol. 8.；Edited by Philip Schaff and Henry Wace.；Buffalo, NY: Christian Literature Publishing Co.,；1895.；<http://www.newadvent.org/fathers/3202082.htm>.}

% structure: p0001 source=h1 target=section reason=page-title

\section{第八十三封书信}

\Emph{\Person{凯撒利亚的圣巴西略}}\par

\Emph{致一位地方官}\par

我与阁下相识相交时日虽短，但通过传闻，我与许多有地位有影响的人物有所交往，从这些传闻中我对您已有不少并非微不足道的了解。至于我是否因传闻而在您心目中有什么分量，您最清楚。无论如何，我在您心目中的声誉正如我所说。然而，既然天主召唤您从事一项使您能够施恩的职业，并且在其行使中，您有能力使我这已夷为平地的城市得以重建，我想，我有责任提醒阁下：要指望天主必将给予的酬报，您应显出这样的品格，以赢得不朽的记忆，并因使受难者的痛苦变得难以承受而承受永恒的安息。我在\Place{查马内内}有一份产业，恳请您像对待自己的一样关照它的利益。请勿惊讶我称朋友的产业为自己的，因为在其他美德中，我也学会了友谊之道，并记得那位智者的话：\Quote{朋友是另一个自己。}因此，我如同对自己的产业一样，将这份属于朋友的产业托付给阁下。恳请您考虑这户人家的不幸，给予他们过去的安慰，并使这地方未来对他们更加舒适；因为由于强加于其上的重重捐税，它已无人居住。我将尽力与阁下会面，商谈细节。\par

\SourceRecord{Translated by Blomfield Jackson.；Nicene and Post-Nicene Fathers, Second Series；Vol. 8.；Edited by Philip Schaff and Henry Wace.；Buffalo, NY: Christian Literature Publishing Co.,；1895.；<http://www.newadvent.org/fathers/3202083.htm>.}

% structure: p0001 source=h1 target=section reason=page-title

\section{\WorkTitle{Letter 84}}

\Emph{\Person{凯撒勒雅的圣巴西略}}\par

\Emph{致总督}。\par

1. 你几乎不会相信我将要写的内容，但为了真理的缘故，必须写下来。我非常渴望能尽可能频繁地与阁下通信，但当我得到这次写信的机会时，并没有立刻抓住这个幸运的时机。我犹豫并退缩了。令人惊奇的是，当我得到我一直在祈求的东西时，我却没有接受。其原因在于，我实在每次写信给你时都感到羞愧，因为这并非纯粹出于友谊，而是带有求取某物的目的。但随后我思忖（当你考虑这一点时，我希望你不会认为我与你通信更多是为了交易而非友谊），接近一位官员和一位普通人的方式必须有所不同。我们不会像对任何无名之辈那样对医生说话；也不会像对私人那样对官员说话。我们试图从对方的技艺或地位中获取一些益处。在阳光下行走，你的影子就会跟随你，无论你愿意与否。同样，与大人交往必然会带来好处，即对受苦之人的援助。我写信的首要目的在能够问候阁下时已经实现。真的，即使我完全没有其他写信的理由，这也必须被视为一个极好的主题。那么，亲爱的先生，请接受我的问候；愿你在担任一任又一任官职时，被全世界保护，并藉你的权威时而帮助这人，时而帮助那人。这样的问候是我惯常的；所有在你任期内多少经历过你的仁慈的人都应向你致以这样的问候。\par

2. 现在，在作了这一祈祷之后，请听我为那位可怜的老人所作的恳求。皇帝的谕令已豁免他担任任何公职；但实际上，我可以说，年老先于皇帝给了他退伍。你自己已满足了上级赐予他的恩惠，既出于对自然衰弱之躯的尊重，我想，也出于对公共利益的考虑，以免国家因一个年迈而变得愚钝的人蒙受损害。但是，亲爱的先生，你如何无意中将他拖入了公共生活，竟下令让他的孙子，一个尚未满四岁的孩子，列入元老院名册？你做的正是同一件事：通过其后代将老人再次拖入公务。但现在，我恳求你，怜悯这两个年龄，并让他们因各自的可怜之处而获得自由。一个从未见过父母，从未认识他们，自幼便失去双亲，靠陌生人的帮助才进入人生；另一个得以存活如此之久，以至遭受了各种灾难。他看到了儿子的早逝；他看到了后继无人的家宅；现在，除非你以你的仁慈想出某种相称的补救办法，他将看到丧子之痛的唯一安慰被用作无尽烦恼的缘由。因为，我想，那小子永远不会担任元老，收税或发饷；但老人的白发必将再次蒙羞。请依照法律且合乎自然地赐予恩惠；下令允许这孩子等到成年，而老人则在床上安详地等待死亡。别人若愿意，尽可以推托公务繁忙、迫不得已。但是，即使你公务繁忙，你也不会轻视受难者，轻慢法律，或拒绝听从朋友的祈祷。\par

\SourceRecord{Translated by Blomfield Jackson.；Nicene and Post-Nicene Fathers, Second Series；Vol. 8.；Edited by Philip Schaff and Henry Wace.；Buffalo, NY: Christian Literature Publishing Co.,；1895.；<http://www.newadvent.org/fathers/3202084.htm>.}

% structure: p0001 source=h1 target=section reason=page-title

\section{第八十五封书信}

\Signature{\Person{圣巴西略}（\Place{凯撒勒雅}）}\par

\Emph{不应当起誓。}\par

我一向的习惯是在每次会议上抗议，并在私下交谈中敦促，不该由收税人将关于税赋的誓言加诸农夫。我同样以书面方式，在天主和人面前，就相同事项作证：你们应停止加给人的灵魂以死亡，并想出一种其他的征收方法，同时让人保持灵魂不受伤害。我这样写信给你们，并非好像你们需要口头的劝诫（因为你们自己有现成的动机敬畏主），而是为使你们所有的属下从你们学习，不激怒圣者，也不因错误的熟悉而让一项被禁止的罪变得无关紧要。起誓对他们没有任何益处，以使他们缴付所征收的，而且他们灵魂遭受无可否认的伤害。因为当人习惯了发假誓，他们就不再给自己压力去偿还，反而认为在誓言中发现了诡计的手段和拖延的机会。因此，如果主严厉惩罚发假誓的人，当债务人因惩罚而灭亡时，被传唤时就无人应答。另一方面，如果主以忍耐容忍，那么，正如我之前所说，那些试探主耐心的人就会藐视祂的良善。不要让他们徒然违反法律；不要让他们招惹天主的义怒反对自己。我已经说了我该说的话。悖逆的人将看见。\par

\SourceRecord{Translated by Blomfield Jackson.；Nicene and Post-Nicene Fathers, Second Series；Vol. 8.；Edited by Philip Schaff and Henry Wace.；Buffalo, NY: Christian Literature Publishing Co.,；1895.；<http://www.newadvent.org/fathers/3202085.htm>.}

% structure: p0001 source=h1 target=section reason=page-title

\section{第八十六书}

\Person{圣巴西略·凯撒利亚}\par

致总督\par

我知道，阁下的首要目标是在各方面支持正义；其次，造福朋友，并为那些投靠您保护的人尽力。在我们面前的这件事上，这两点都得到了结合。我们所辩护的事业是正义的；它对我——您朋友中的一员——来说是可贵的；它也是对那些在苦难中祈求您坚毅援助的人所应得的。我那非常亲爱的兄弟\Person{多罗修斯}仅有的维持生计所需的谷物，被\Place{贝里西}的一些负责管理事务的当局人员夺走了。他们或是出于自愿，或是受他人唆使而施行了这种暴力。无论哪种情况，这都是可起诉的罪行。因为，一个人的邪恶源于自身，与仅作为他人邪恶的施行者相比，何者罪行更轻？对受害者而言损失是一样的。因此，我恳求您，使\Person{多罗修斯}能从他遭抢的人那里收回谷物，不容许他们把暴行的罪责推到别人肩上。若您允准我的请求，我将根据为获得食物所需的重要性，来衡量您赐予的恩惠的价值。\par

\SourceRecord{Translated by Blomfield Jackson.；Nicene and Post-Nicene Fathers, Second Series；Vol. 8.；Edited by Philip Schaff and Henry Wace.；Buffalo, NY: Christian Literature Publishing Co.,；1895.；<http://www.newadvent.org/fathers/3202086.htm>.}

% structure: p0001 source=h1 target=section reason=page-title

\section{第八十七书}

\Emph{\Person{圣凯撒勒雅巴西略}}\par

\Emph{无收信人，同一主题}。\par

我甚为惊讶，既有您可申诉，竟对那位司铎犯下如此严重之过，以致他被剥夺了唯一的生计。此事最严重之处在于，犯事者将其行为的罪责转嫁给您；而您始终有责任不仅不容许此类行为发生，更要运用全部权力防止它们发生在任何人身上，尤其是那些与我同心、行走在同一真宗教之路上的司铎们。若您有意使我满意，请务必立即纠正这些事。因为，在天主助佑下，您能成就此事，甚至更大的事，凡您所愿者。我已写信给我本省的总督，告知若他们不肯自愿行事，可凭法庭的压力强制他们如此行。\par

\SourceRecord{Translated by Blomfield Jackson.；Nicene and Post-Nicene Fathers, Second Series；Vol. 8.；Edited by Philip Schaff and Henry Wace.；Buffalo, NY: Christian Literature Publishing Co.,；1895.；<http://www.newadvent.org/fathers/3202087.htm>.}

% structure: p0001 source=h1 target=section reason=page-title

\section{第八十八封书信}

\Emph{\Person{凯撒勒雅的圣巴西略}}\par

\Emph{无称呼，论征税事宜。}\par

您的阁下比任何人都更清楚凑齐捐献的金银的困难。对于我们的贫困，没有比您本人更好的见证了，因为您以极大的仁慈同情我们，并且至今，凡在您能力范围内，您一直容忍我们，从未因上级权威引起的恐慌而偏离您天生的宽厚。现在，总额中尚有一些欠缺，必须从我们向全城推荐的捐献中收取。我所请求的是，您能赐予我们一些宽限，以便向乡间居民发出提醒，而我们的多数官员都在乡下。如果可能，按照我们仍欠缺的磅数来缩减缴纳，我很乐意您这样安排，所缺金额稍后补足。然而，如果必须一次缴清全部金额，那么我重申我的第一个请求：请允许我们更长的恩惠期限。\par

\SourceRecord{Translated by Blomfield Jackson.；Nicene and Post-Nicene Fathers, Second Series；Vol. 8.；Edited by Philip Schaff and Henry Wace.；Buffalo, NY: Christian Literature Publishing Co.,；1895.；<http://www.newadvent.org/fathers/3202088.htm>.}

% structure: p0001 source=h1 target=section reason=page-title

\section{第八十九封信}

\Person{凯撒勒雅的圣巴西略}\par

致\Person{安提约基雅主教梅肋提乌}。\par

1. 仁慈的天主赐我向您的尊威致候的机会，使我深切的渴望得以缓解。祂亲自作证，我多么热切渴望见到您的面容，并享受那使灵魂复苏的良善教导。如今，藉由我敬爱而卓越的兄弟、执事\Person{多罗特乌}出发之际，我首先恳求您为我祈祷，使我不致成为百姓的绊脚石，也不致妨碍您恳求主的祈祷。其次，我建议您通过上述兄弟妥善安排一切；若认为应当寄一封信给西方人士，因为理应通过我们自己的信使将信息以书面形式传达，那么请您口述这封信。我已会晤了他们派来的执事\Person{萨比诺}，并写信给\Place{伊利里亚}、\Place{意大利}和\Place{高卢}的主教们，以及几位私下写信给我的人。因为应当派遣某人代表主教会议的共同利益，携带一封我恳请您写就的第二封信。\par

2. 关于可敬的\Person{阿塔纳修}主教，您已明悉我将提及之事：除非他通过某种方式从您那里领受共融——当时您推迟了此事——否则我的信件无法推进任何事，也无法实现任何可期望的目标。据说他非常渴望与我联合，并愿意尽其所能贡献一切，但遗憾的是他未被授予共融而离开，且承诺仍未兑现。\par

东方所发生的事不可能不传到您的尊威耳中，但上述兄弟将口头传达更准确的信息。请务必在复活节后立即派遣他，因他正等待\Place{萨莫萨塔}的答复。请善意看待他的热忱，以您的祈祷坚固他，并派遣他执行此任务。\par

\SourceRecord{Translated by Blomfield Jackson.；Nicene and Post-Nicene Fathers, Second Series；Vol. 8.；Edited by Philip Schaff and Henry Wace.；Buffalo, NY: Christian Literature Publishing Co.,；1895.；<http://www.newadvent.org/fathers/3202089.htm>.}

% structure: p0001 source=h1 target=section reason=page-title

\section{第九十封书信}

\Signature{圣巴西略}\par

\Emph{致西方诸位圣善的弟兄、主教们。}\par

良善的天主时常在痛苦中掺和安慰，如今在悲痛之中，也借着我们所敬重的教父、主教亚大纳削从你们那里收到并转寄给我的信件，赐予我些许慰藉。因为那些信件包含着纯正信德的明证，以及你们无可侵犯的同心同德的凭据，表明牧者们追随先祖的足迹，以知识牧养主的子民。这一切使我的心灵如此欢畅，驱散了沮丧，在我现在所处的令人忧虑的境遇中，为我的灵魂带来一丝笑意。主又通过我儿、可敬的执事撒彼诺将安慰延伸给我，他向我精确叙述了你们的情况，使我的灵魂得到鼓舞；他本人也将根据自己的亲身经历，将我们的情况清楚地告知你们，以便你们首先以热切不断的祈祷向天主援助我的困苦，其次，请赐予你们力所能及的安慰给我们受难的教会。因为在这里，最可敬的弟兄们，一切都已衰微；教会面对敌人持续的攻击已经退却，如同海中的小舟在波涛的连续打击下颠簸；除非天主的仁慈迅速施以援手。因此，正如我们把你们的彼此同情和团结视为我们自身的重要福祉，我们同样恳求你们怜悯我们的分裂；不要因为我们相隔遥远就将我们与你们分离，而要允许我们因同一圣神的共融而进入同一身体的同心同意。\par

我们的苦难已是众所周知，即使我们闭口不言，因为如今它们的声音已传遍天下。教父的教义被藐视，宗徒传统被废弃；革新者的计谋在教会中流行；如今人们更擅长狡猾体系的编造者而非神学家；此世的智慧赢得了最高奖项，拒绝了十字架的荣耀。牧人被驱逐，取而代之的是凶残的豺狼冲散基督的羊群。祈祷的殿宇无人聚集；旷野中充满哀哭的人群。年长者哀叹今昔之比。年轻人更值得怜悯，因为他们不认识自己被剥夺了什么。这一切足以触动那些学会了基督之爱的人的怜悯之心；然而，与实际情况相比，言语远远不足。因此，若有什么爱德的安慰，圣神的共融，慈悲的心肠，请被激动来帮助我们。为真道而热忱，将我们从这场风暴中拯救出来。愿那著名的教父教义永远在我们中间被勇敢地宣讲，它摧毁了亚略臭名昭著的异端，并建立了教会于纯正的教义中，即宣认圣子与圣父同一实体，圣神同等尊荣、同受钦崇，好使由于你们的祈祷和合作，主将赐予我们同样为真理的勇敢和因宣认神圣而赐救恩的天主圣三而得荣耀，就如祂已赐予你们的一样。但前述的执事将详细告知你们一切。我们欢迎你们宗徒般的正统热诚，并同意所有由诸位尊座合乎教规所行的一切。\par

\SourceRecord{Translated by Blomfield Jackson.；Nicene and Post-Nicene Fathers, Second Series；Vol. 8.；Edited by Philip Schaff and Henry Wace.；Buffalo, NY: Christian Literature Publishing Co.,；1895.；<http://www.newadvent.org/fathers/3202090.htm>.}

% structure: p0001 source=h1 target=section reason=page-title

\section{第九十一封书函}

\Signature{圣巴西略·凯撒勒雅}\par

\Emph{致\Person{瓦莱里亚努斯}，\Place{伊利里库姆}主教}。\par

感谢归于主，祂允许我通过你贞洁的品行看到初期爱德的果实。虽然你身体远离，但你藉着来信与我结合，以属灵而神圣的渴望拥抱了我，在我灵魂中注入了不可言说的感情。如今我明白了这句箴言的力量：\BibleQuote{凉水给口渴的灵魂，来自远方的佳音也是如此。} \BibleRef{箴 5:25}尊敬的长兄，我实在渴求爱德。原因并不难找，因为邪恶增多，许多人的爱已冷淡。因此，你的信对我弥足珍贵，我藉着可敬的兄弟\Person{萨比努斯}回信。通过他，我向你表明自己，并恳请你在为我们祈祷上保持警醒，好使天主有朝一日赐予这里的教会平静与安宁，并斥责这风和海，使我们摆脱现在每时每刻都预期被淹没的风暴与动荡。但在这诸多患难中，天主赐给了我们一件大恩：听闻你们彼此完全一致合一，真正信仰的教义在你们中间毫无阻碍地被宣讲。因为终有一天——除非这个世界的时期已经结束，并且人间的日子尚存——藉着你们，信仰必将在东方得以更新，你们也必在适当的时候回报她所赐予你们的祝福。我们这里保存着祖先纯正信仰的健康部分遭受重创，魔鬼以其诡诈通过许多各样狡猾的攻击粉碎了它。但是，藉着你们爱主之人的祈祷，愿邪恶诡诈的\OriginalTerm{亚略}{Arian}异端之错误被熄灭；愿在\Place{尼西亚}聚会的祖先们的美好教导发光照耀，使荣耀归于受祝福的\OriginalTerm{天主圣三}{Trinity}，并按照救恩的\OriginalTerm{圣洗}{Baptism}的格式。\par

\SourceRecord{Translated by Blomfield Jackson.；Nicene and Post-Nicene Fathers, Second Series；Vol. 8.；Edited by Philip Schaff and Henry Wace.；Buffalo, NY: Christian Literature Publishing Co.,；1895.；<http://www.newadvent.org/fathers/3202091.htm>.}

% structure: p0001 source=h1 target=section reason=page-title

\section{第九十二封书信}

\Signature{圣巴西略·凯撒利亚}\par

致意大利人和高卢人。\par

1. 致我们在意大利和高卢供职的、极其虔诚和圣洁的弟兄们，与我们同心合意的主教们：我们，即\Person{梅利提乌斯}、\Person{优西比乌}、\Person{巴西略}、\Person{巴苏斯}、\Person{额我略}、\Person{佩拉吉乌斯}、\Person{保禄}、\Person{安提穆斯}、\Person{狄奥多图斯}、\Person{比图斯}、\Person{阿布拉米乌斯}、\Person{约比努斯}、\Person{则诺}、\Person{狄奥多勒图斯}、\Person{马尔基亚努斯}、\Person{巴拉库斯}、\Person{阿布拉米乌斯}、\Person{利巴尼乌斯}、\Person{塔拉西乌斯}、\Person{若瑟}、\Person{波厄图斯}、\Person{亚特里乌斯}、\Person{狄奥多图斯}、\Person{欧斯塔提乌斯}、\Person{巴尔苏马斯}、\Person{若望}、\Person{科斯洛埃斯}、\Person{约萨克斯}、\Person{纳尔塞斯}、\Person{马里斯}、\Person{额我略}和\Person{达夫努斯}，在主内问候你们。痛苦的灵魂在内心深处发出一声又一声叹息时，会找到一些安慰；甚至流下一滴眼泪也能减轻痛苦的重压。但叹息和泪水给我们的安慰，还不如向你们的爱德倾诉我们的烦恼的机会多。此外，我们还因一个更好的希望而振奋：或许，如果我们向你们宣布我们的烦恼，可能会促使你们给予我们那种我们长久以来盼望你们给予东方教会的援助，但我们至今尚未收到；天主以其智慧安排一切，祂必是按照其公义的隐秘判断，命定我们在这试探中经受更长时间的考验。我们处境的声誉已传遍地极，你们并非不知道；你们也并非不同情与你们自己同心合意的弟兄，因为你们是那位教导我们爱邻舍成全法律的宗徒的门徒。但是，正如我们所说的，天主的公义判断——它规定我们必须承受因我们的罪而带来的苦难——阻止了你们。但当你们从我们可敬的弟兄、执事\Person{萨比努斯}那里得知一切，特别是那些迄今尚未传到你们耳中的事时，他将能亲自叙述我们信中所遗漏的部分，我们恳求你们被激发起对真理的热忱和对我们的同情。我们恳求你们存怜悯的心肠，抛开一切犹豫，承担起爱德的劳苦，不计路途遥远、你们自己的事务或任何其他世间的挂虑。\par

2. 不仅仅是一个教会处在危险中，也不仅仅是两三个教会陷入这场可怕的风暴。这个异端的祸害几乎从\Place{伊利里库姆}的边界蔓延到\Place{特拜德}。它的坏种子首先由臭名昭著的\Person{亚略}播下；然后，在他和我们之间的时代里，通过许多人的努力，它深深扎根。现在，它们结出了致命的果实。真宗教的教义被推翻了。教会的法律陷入混乱。那些不敬畏天主的人的野心冲上高位，显赫的职位现在公开被认为是邪僻的奖赏。结果是，一个人越亵渎，人民就越认为他适合做主教。圣职的尊严已成往事。完全缺乏有知识的人牧养主的羊群。野心勃勃的人不断将供应穷人的财物用于自己的享乐和分送礼物。没有精确的法规知识。犯罪完全不受惩罚；因为那些靠人的恩宠获得职位的人，必须通过不断宽容犯罪者来回报恩惠。公义的审判已成往事；每个人都按自己心里的欲望行事。邪恶没有界限；人不知约束。当权者不敢说话，因为那些靠人的利益获得权力的人，是那些使他们升迁者的奴隶。现在，对正统的辩护在某些地方被视为相互攻击的机会；人们隐藏私下的恶意，假装他们的敌意全然是为了真理。另一些人害怕因可耻的罪行被定罪，就煽动人民陷入骨肉相争的争吵，以便在普遍的苦难中自己的行为不被注意。因此，战争没有休战的可能，因为作恶的人害怕和平，因为和平可能会揭开他们隐秘的丑行。与此同时，不信的人嘲笑；信心软弱的人动摇；信仰不确定；灵魂在无知中沉浸，因为圣言的掺杂者模仿真理。真信者的口被堵住，而每一根亵渎的舌头自由摆动；圣物被践踏；较好的平信徒躲避教堂，视之为邪僻学府；他们在旷野中举手叹息流泪，向天上的主呼求。你们一定也听说过我们大多数城市里发生的事：我们的人民带着妻子儿女，甚至老人，涌出城墙，在露天祈祷，耐心忍受天气的一切不便，等待主的帮助。\par

3. 有哪一声哀叹能匹配这些苦难？又有多少眼泪的泉源足以为此流泪？因此，当一些人似乎仍然站立，旧日状态的痕迹尚存时，在教会完全船毁之前，请尽快到我们这里来，现在就到我们这里来，真正的弟兄们，我们恳求你们；我们双膝跪下恳求你们，伸出援手。愿你们兄弟般的慈肠被感动；愿同情的泪水流淌；不要无动于衷地看着半个帝国被错谬吞噬；不要让信仰之光在它最初照耀的地方熄灭。\par

至于你当以何种行动有助于事态，以及你应如何向受困者表示同情，无需我们告知；圣神自会启示你。但无疑地，若要拯救幸存者，便需迅速行动，并需相当数量的弟兄到来，使来访者得以凑足主教会议的法定人数，以便在推行改革时具有分量——不仅因其所代表的尊严，亦因其自身的人数：如此，他们将恢复我们教父在\Place{尼西亚}所撰写的信经，禁绝异端，并使所有同心者达成一致，向各教会宣告和平。因为这一切中最可悲的是，健全的部分自相分裂，我们遭受的困扰恰似当年\Person{韦斯巴芗}围攻\Place{耶路撒冷}时所发生的那样。当时的犹太人既受外敌围攻，又被内部的叛乱消耗。在我们这里也一样，除了异端者的公开攻击，教会因那些本应视为正统者之间的战火而陷入极度无助。因此，我们确实渴望你的帮助，以便将来所有宣认宗徒信德的人，可以终止他们可悲地制造的分裂，并从此归服于教会的权威；这样，基督的身体再次得以完整，所有肢体恢复合一。如此，我们不仅能赞美他人的福祉（这是我们现在唯一能做的），还能看到我们自己的教会再次恢复其昔日正统的荣耀。因为，主赐予你的恩惠确实最值得庆贺——你辨别假冒与真实纯洁的能力，以及你毫不掩饰地宣讲教父的信德。我们接受了那信德；我们知道那信德印有宗徒的标记；我们赞同那信德，也赞同主教会议信函中一切合乎教规和法律的宣告。\par

\SourceRecord{Translated by Blomfield Jackson.；Nicene and Post-Nicene Fathers, Second Series；Vol. 8.；Edited by Philip Schaff and Henry Wace.；Buffalo, NY: Christian Literature Publishing Co.,；1895.；<http://www.newadvent.org/fathers/3202092.htm>.}

% structure: p0001 source=h1 target=section reason=page-title

\section{书信九十三}

\Emph{\Person{凯撒勒雅的圣巴西略}}\par

\Emph{致贵妇凯撒利亚：论领圣体}\par

每天领圣体，领受基督的圣体圣血，是良好而有益的。因为祂明说：\BibleQuote{谁吃我的肉，并喝我的血，必得永生。}\BibleRef{若 6:54} 谁又会怀疑，常分享生命，就等于拥有丰富的生命呢？我本人每周领圣体四次：在主日、星期三、星期五、安息日，以及在其他日子如果有纪念某位圣人时。无需指出，在迫害时期，任何人在没有司铎或执事在场的情况下，被迫自己亲手领圣体，并不算严重过失，因为长久以来的习俗本身已认可这种做法。在旷野中，所有隐修士，在没有司铎的地方，就自己领圣体，并在家里保存圣体。在\Place{亚历山大里亚}和\Place{埃及}，大部分平信徒都各自在家里保存圣体，并在愿意时领受。因为一旦司铎完成了祭献并分送了圣体，领受者每次完整地参与领受时，必须相信他从分送者那里正确地取用并领受了。即使在教堂里，当司铎分送圣体时，领受者完全掌握它，并用自己的手举到唇边。无论是从司铎一次领受一块或多块，其效力相同。\par

\SourceRecord{Translated by Blomfield Jackson.；Nicene and Post-Nicene Fathers, Second Series；Vol. 8.；Edited by Philip Schaff and Henry Wace.；Buffalo, NY: Christian Literature Publishing Co.,；1895.；<http://www.newadvent.org/fathers/3202093.htm>.}

% structure: p0001 source=h1 target=section reason=page-title

\section{书信第九十四号}

\Emph{\Person{凯撒勒雅的圣巴西略}}\par

\Emph{致行省总督埃利亚斯}。\par

我也一直非常渴望见到阁下，以免因未能如此而比我的控告者更不利；但身体的疾病阻止了我，比平时更加严重，所以我被迫写信给您。不久前，最卓越的先生，当我得以与阁下会面时，我急于与您的智慧商讨关于我的一切事务；我也急于为您代表教会讲话，以免为将来的诽谤留下任何把柄。但我克制了自己，认为给一个肩负如此众多责任的人添上他自身必要事务之外的麻烦，完全是多余且令人厌烦的。同时（因为必须说出真相），我确实不愿因我们相互指责而被迫伤害您的灵魂，而您本应因对天主的纯全虔诚而获得完美的虔诚赏报。因为，实在说，如果我吸引您注意我，我将使您对公务的闲暇所剩无几；这就像一个在波涛汹涌中操纵新船的人，本应减轻货物、尽力使船轻便，却额外添加行李一样。正因如此，我想，我们伟大的皇帝在看出我何等忙碌之后，让我独自管理教会。现在，我想请那些围攻您公正之耳的人问问，我治理政府究竟遭受了什么损害？通过我对教会的管理，任何公共利益，无论大小，遭受了什么贬值？不过，也许有人会辩称，我因为天主建造一座辉煌的教堂，并在其周围建造一所慷慨分配给主教的住宅，以及下面分配给教会职员的其它房屋，而损害了政府。这两者都可供各位行政长官及其随从使用。但建造一个招待客旅的地方，包括旅行者和因病需要治疗的人，从而为他们提供所需的安慰、医生、交通工具和随从，我们又伤害了谁呢？所有这些人都必须学习生活所必需的职业，这些职业被认为对体面的事业至关重要；他们还必须有适合其工作的建筑，所有这些都使该地增光，并且由于他们的声誉归功于我们的总督，也给他带来荣耀。当然，这并非因此使您不情愿地接受统治我们的责任，因为您凭借高尚的品格，足以恢复我们的废墟，使荒芜的地区有人居住，将荒野变成城镇。在履行这些职责时，是骚扰和烦恼，还是尊重和敬重一位同伴更好呢？最卓越的先生，请不要认为我说的只是空话。同时，我们已经开始准备材料。在我们的统治者面前，我们的辩护就到此为止。至于要如何回答控告者的指控，对于一个基督徒和一位关心我意见的朋友，我现在不能再多说；这个主题对于一封信来说太长，而且也不便安全地写下来。但为了避免在我们有机会见面之前，您因一些人的诽谤而放弃对我的任何善意，请像亚历山大那样做。故事是，正如您记得的，当他的一个朋友被诽谤时，他让一只耳朵向诽谤者敞开，而小心地用另一只手捂住另一只耳朵，目的是表明负责判断的人不应轻易完全听信首先占据他注意的人，而应保持一半的听力为缺席者辩护。\par

\SourceRecord{Translated by Blomfield Jackson.；Nicene and Post-Nicene Fathers, Second Series；Vol. 8.；Edited by Philip Schaff and Henry Wace.；Buffalo, NY: Christian Literature Publishing Co.,；1895.；<http://www.newadvent.org/fathers/3202094.htm>.}

% structure: p0001 source=h1 target=section reason=page-title

\section{第95封信}

\Person{圣巴西略·凯撒勒雅}\par

\Emph{致\Place{萨莫萨塔}主教\Person{尤西比乌斯}}。\par

我此前曾致书尊驾，谈及我们彼此相会及其他事宜，但遗憾的是，我的信未能送达阁下。因为那已故的执事\Person{德奥佛拉斯图斯}受托保管书信后，我因不得已的旅行外出，他未将信转交给尊驾，因为他患病而亡。因此，我的写信便耽搁如此之久，以致现在时间已极其紧迫，我不指望这封信会有多大用处。虔诚的主教\Person{梅利提乌斯}和\Person{德奥多图斯}曾极力敦促我去拜访他们，表示会面将是友爱的证明，并希望补救目前令人担忧的困境。他们指定了将近六月中旬作为我们相会的时间，地点则是\Place{法尔加穆斯}，一个以殉道者荣耀和每年在此召开大型会议而闻名的地方。我一返回，便听闻那位已故执事的死讯，我的信函也闲置在家中，我深感不能懈怠，因为距离预定时间仍剩三十三天，于是我匆忙将信送交非常尊敬的\Person{尤斯塔修斯}——我的同工，以便由他转交给尊驾，并迅速得到回音。因此，如果可能且您乐意前来，我也将前来。若不然，但愿天主允许，我将偿还去年所欠的会面之债：除非因我的罪过而再次出现某种阻碍，那样的话，我只能将与众主教的会面推迟到另一时间。\par

\SourceRecord{Translated by Blomfield Jackson.；Nicene and Post-Nicene Fathers, Second Series；Vol. 8.；Edited by Philip Schaff and Henry Wace.；Buffalo, NY: Christian Literature Publishing Co.,；1895.；<http://www.newadvent.org/fathers/3202095.htm>.}

% structure: p0001 source=h1 target=section reason=page-title

\section{第九十六封信}

\Emph{凯撒勒雅的圣巴西略}\par

\Emph{致索弗洛尼乌斯阁下}。\par

有谁像您这样热爱自己的城邦，以赤子之心尊敬生养自己的地方，为全城和城中的每一个人祈祷，而且不仅祈祷，还用您自己的方式证实您的祈祷？因为靠着天主的助佑，您能够做到这些，而且，善良的人啊，愿您长久能力如此。然而，在您主政的时期，我们的城邦只享受了短暂的繁荣之梦，因为它曾交到这样一位长官手中——根据我们最古老编年史的研究者所说，像他这样的人从未坐过总督席位。但现在，城邦突然失去了他的服务，因为有些人的邪恶，他们以他的宽厚和公正为攻击点，并在未告知阁下您的情况下，编造了针对他的诽谤。因此，我们普遍感到沮丧，失去了这样一位总督：他有着独特的能力来振兴我们颓丧的团体，他是正义的真正守护者，冤屈者可以接近他，违法者畏惧他，对待富人和穷人一视同仁，而最重要的是，他使基督徒的利益恢复了往日的尊荣。他是我所知所有人中最不受贿赂的，也从未对任何人行不义的偏袒——我认为这一点与他其他的德行相比还是小事。我如今作证确实太迟了，就像那些唱挽歌以自慰的人，得不到任何实际的救助。然而，让他的记忆留在您慷慨的心中，并让您因他作为故乡的恩人而感谢他，并非无用。如果那些因他不牺牲正义以满足他们利益而对他怀恨的人攻击他，您若能为他辩护和保护他，那就好了。这样，您将向所有人明确表示，您把他的利益视为自己的利益，而您与他如此亲近的充分理由就是他无可指摘的政绩和他任职期间其治理的卓越表现。因为别人在许多年里无法做到的事，他在短时间内就迅速完成了。如果您能向皇帝推荐他，并驱散针对他的诽谤指控，那将是对我极大的恩惠，也是在此情况下的慰藉。请相信我，我在这里不只是为自己说话，而是为整个团体说话，我们一致祈祷他能从您的阁下的帮助中获益。\par

\SourceRecord{Translated by Blomfield Jackson.；Nicene and Post-Nicene Fathers, Second Series；Vol. 8.；Edited by Philip Schaff and Henry Wace.；Buffalo, NY: Christian Literature Publishing Co.,；1895.；<http://www.newadvent.org/fathers/3202096.htm>.}

% structure: p0001 source=h1 target=section reason=page-title

\section{\WorkTitle{第97封书信}}

\Person{圣巴西略（凯撒勒雅的圣巴西略）}\par

\Place{提雅纳元老院}\par

那揭示隐藏之事、显明人心思虑的主，甚至连卑微之人都赐予了认识看似难解之诡计的知识。\newline{} 没有任何事逃过我的注意，也没有任何单独行动不为人知。然而，除了天主的平安和与之相关的一切，我什么也看不见，什么也听不见。\newline{} 他人或许伟大、有力、自信，但我什么也不是，毫无价值，因此我绝不敢自以为无需支持就能处理事务。我完全明白，我比一只手需要另一只手更需要每一位弟兄的帮助。\newline{} 的确，主藉着我们身体的构造，教导了我们共融的必要性。当我观看我的肢体，发现没有一个是自足的，我怎能认为自己足以履行生活的职责呢？一只脚若无另一只脚的支持，不能安稳行走；一只眼若无另一只眼的配合并与之共同注视对象，也不能看得清楚。听觉因双耳接收声音而更准确，握力因手指的协作而更牢固。\newline{} 总之，在自然和意志所行的一切事上，我看见没有一件是离开同力协作而完成的。甚至连祈祷，若未被联合的祈祷，就失去其自然的力量；主告诉我们，当两三个人同心合意呼求祂时，祂就在他们中间。\newline{} 主亲自担负了救恩计划，为藉着祂十字架的血，使天上地上的事物得以和好。\newline{} 因此，为了这一切理由，我祈祷在我余下的日子能安于和平；我祈求能在和平中安眠。为了和平的缘故，我不逃避任何劳苦，不退缩任何谦卑的行动或言语；我不计算旅途的远近，愿承受任何不便，只求能获致和平的赏报。\newline{} 若有人在这方向上跟随我，那很好，我的祈祷得以垂允；但若结果不同，我也不会退缩。在报应之日，各人将按自己的作为领取果实。\par

\SourceRecord{Translated by Blomfield Jackson.；Nicene and Post-Nicene Fathers, Second Series；Vol. 8.；Edited by Philip Schaff and Henry Wace.；Buffalo, NY: Christian Literature Publishing Co.,；1895.；<http://www.newadvent.org/fathers/3202097.htm>.}

% structure: p0001 source=h1 target=section reason=page-title

\section{第九十八封信}

\Emph{\Person{凯撒勒雅的圣巴西略}}\par

\Emph{致\Person{萨摩沙塔主教优西比乌}}。\par

1. 收到您的圣座的信函后，您在信中说您不会来，我非常渴望动身前往\Place{尼科波利斯}，但我的意愿渐渐减弱，并想起了我的一切软弱。我也想到那些邀请我的人的举止缺乏庄重。他们通过我们可敬的兄弟，\Place{纳齐盎司}的税务监督\Person{赫楞尼乌斯}，随便邀请了我，却从未费心派信使提醒我或派人护送我。由于我的罪，我对他们来说是个可疑对象，我避免因我的在场而玷污他们聚会的明亮。与您的阁下在一起时，我并不回避全力以赴应对重大考验；但离开了您，我甚至觉得自己几乎无力面对日常的烦恼。既然我与他们的会面原是为了教会事务，我让节庆的时光过去，将会议推迟到一个安息且无干扰的时期，并决定前往\Place{尼科波利斯}，与虔诚的\Person{梅利蒂乌}主教讨论各教会的需要，以防他不愿前往\Place{萨摩沙塔}。如果他同意，我将赶紧去见他，只要这事由你们两位——他回复我（我已写信）和您——向我明确说明。\par

2. 我们本应会见\Place{第二卡帕多奇亚}的主教们，但他们一旦被划分到另一个管辖区，就突然觉得自己成了外人，与我疏远了。他们忽视我，好像从未受我管辖，与我毫无关系。我也期待着与可敬的\Person{尤斯塔提乌斯}主教的第二次会面，实际上也发生了。因为许多人指责他损害信仰，我见了他，发现藉天主的恩宠，他衷心追随一切正统信仰。由于那些本应传递我的信的人的错误，那位主教的信未能呈递给您的阁下，而且我事务繁多，竟忘记了。\par

我也希望我们的兄弟\Person{额我略}能治理一个与他能力相称的教会；那将会是普世教会聚集在一处。但既然不可能，就让他做主教吧，不是从他的教座获得尊荣，而是他自己赋予教座尊荣。因为一个真正伟大的人物，不仅足以成就大事，还能凭自己的影响使小事变为伟大。\par

但是，面对\Person{帕尔马提乌斯}，在弟兄们多次劝诫之后，他仍然帮助\Person{马克西穆斯}进行迫害，我们该怎么做？即使现在，他们也不犹豫给他写信。他们自己因身体软弱和各自的事务而不能前来。请相信我，非常虔诚的父，我们自己的事务非常需要您的在场，您必须再次行动您可敬的老年，以便您能支持现在摇摇欲坠、有倾倒危险的\Place{卡帕多奇亚}。\par

\SourceRecord{Translated by Blomfield Jackson.；Nicene and Post-Nicene Fathers, Second Series；Vol. 8.；Edited by Philip Schaff and Henry Wace.；Buffalo, NY: Christian Literature Publishing Co.,；1895.；<http://www.newadvent.org/fathers/3202098.htm>.}

% structure: p0001 source=h1 target=section reason=page-title

\section{第九十九封信}

\Emph{\Person{圣巴西略·凯撒勒雅}}\par

\Emph{致\Person{特伦提乌斯}伯爵}。\par

我一直渴望并确实尽力遵守皇帝的命令和大人友善的来信，即使只是部分遵守。我深信您的每一句话、每一个想法都充满善意和正直的情感。但我未能通过实际行动表现出我的欣然赞同。最真实的原因是我的罪恶，它们总在我面前升起，总是阻碍我的脚步。其次，还有那位被指派与我合作的主教的疏远，原因我不知道；但是我可敬的弟兄\Person{狄奥多图斯}主教，他起初答应与我合作，并诚挚地从\Place{格塔萨}邀请我到\Place{尼科波利}。然而，当他在镇上看到我时，竟如此震惊于我的罪恶，如此害怕我的罪，以至于他无法忍受带我去参加晨祷或晚祷。他这样做，就我的功过而言，完全是公正的，也非常符合我的生活方式，但不适宜促进教会的利益。他声称的原因是，我接纳了非常可敬的弟兄\Person{欧斯塔修斯}进入共融。我做的事情如下。当应我们弟兄\Person{狄奥多图斯}的邀请参加一个会议时，出于爱，我希望服从召唤，以免使聚会无果而虚，我急切地想与上述弟兄\Person{欧斯塔修斯}进行沟通。我向他提出了关于信德的指控，这些指控是我们的弟兄\Person{狄奥多图斯}对他提出的，我问他，如果他遵循正统信德，就向我表明，这样我就能与他共融；如果他另有想法，他必须清楚地知道我会与他分离。我们就此话题进行了许多交谈，那一天都花在审查上；到了晚上，我们没有任何明确结论就分开了。第二天，我们上午又进行了一次会议，讨论了同样的问题，加入了\Place{塞巴斯特亚}的司铎我们的弟兄\Person{珀梅尼乌斯}，他激烈地与我争论。我逐点澄清了他似乎指控我的问题，并使他们同意我的提议。结果是，藉着主的恩宠，我们发现即使在最微小的细节上也相互一致。所以大约第九时辰，感谢天主恩准我们思想一致、言语一致之后，我们起身去祈祷。除此之外，我应该从他那里获得一些书面声明，以便他的同意能让反对者知晓，他的意见的证明对其他人也足够。但我自己急于追求极大的精确性，想会见我的弟兄\Person{狄奥多图斯}，从他那里获得一份关于信德的书面陈述，并把它提交给\Person{欧斯塔修斯}；这样两个目标可以同时实现：\Person{欧斯塔修斯}宣认正统信德，以及\Person{狄奥多图斯}和他朋友们完全满意，他们在接受自己的提议后将没有反对的理由。但\Person{狄奥多图斯}，在不知道我们为何会面以及我们交流的结果之前，就决定不允许我们参加聚会。所以我们在半路上折返，失望于我们为教会和平所做的努力被阻挠。\par

3. 此后，当我被迫踏上前往\Place{亚美尼亚}的旅程时，我知道此人（指\Person{狄奥多图斯}）的品性，并且既为了在有资格的证人面前为自己所发生的事辩护，也为了满足他，我旅行到了\Place{格塔萨}，在非常虔诚的\Person{梅利提乌斯}主教的领地内，上述\Person{狄奥多图斯}与我同行。在那里，当他指控我与\Person{欧斯塔修斯}共融时，我告诉他，我们交流的结果是我发现\Person{欧斯塔修斯}在一切事情上都与我一致。然后他坚持说，\Person{欧斯塔修斯}离开我后否认了这一点，并对他的门徒断言他从未与我达成关于信德的任何协议。因此，我驳斥了这一说法；看，最卓越的人，我的回答是否非常公正和完全？我说，我确信，根据\Person{欧斯塔修斯}的品性来判断，他不可能如此轻易地从一种方向转向另一种方向，现在宣认现在否认他所说的话；一个即使在任何小事上也避免撒谎如同可怕之罪的人，不太可能在如此重大且众所周知的事情上选择违背真理。但如果你们中间所报告的是真的，他必须面对一份包含正统信德完整阐述的书面声明；然后，如果我发现在书写中他愿意赞同，我将继续与他共融；但是，如果我发现在考验面前他退缩，我将断绝与他的一切来往。\Person{梅利提乌斯}主教同意这些论证，在场的\Person{狄奥多鲁斯}弟兄司铎，然后是\Person{狄奥多图斯}可敬的弟兄都表示同意，并邀请我到\Place{尼科波利}，既为访问那里的教会，也为陪他一直到\Place{萨塔拉}。但他把我留在\Place{格塔萨}，当我到达\Place{尼科波利}时，他忘记了他从我所听到的一切以及与我达成的协议，便打发我离开，遭受了我所提及的侮辱和羞辱。\par

4. 那么，尊敬的大人，我怎能执行任何加给我的命令，并为\Place{亚美尼亚}提供主教呢？当我的责任分担者对我如此态度时——正是我原指望藉着他的帮助能够找到合适人选的那个，因为在他的教区里有备受尊敬和博学的男士，善于言辞，并了解该民族的其他特点——我如何能行事？我知道他们的名字，但我将避免提及它们，以免将来为\Place{亚美尼亚}利益服务时出现任何障碍。\par

现在，我在这样健康状况下到达了\Place{萨塔拉}之后，似乎靠着天主的恩宠处理了其余的事。我在\Place{亚美尼亚}的主教们之间缔造了和平，并向他们发表了一场合适的讲话，敦促他们抛弃惯常的冷漠，重新燃起在主的事业上的古老热忱。此外，我向他们传交了规则，阐明他们应当如何关注在\Place{亚美尼亚}普遍实行的不义行为。我还接受了\Place{萨塔拉}教会的一项决定，请求通过我给他们一位主教。我也仔细查问了针对我们\Place{亚美尼亚}主教\Person{西里尔}弟兄所传播的诽谤，靠着天主的恩宠，我发现它们是由他的仇敌的谎言中伤所引发的。他们向我承认了这一点。我似乎在一定程度上使人民与他和解，以致他们不再回避与他的共融。这些成就或许微不足道，价值不大，但由于魔鬼的诡计所引起的互相不和，我不可能取得更多成果。就连这些我也不该说，以免显得在宣扬自己的耻辱。但由于我无法以其他方式在阁下面前辩护我的案件，我不得不告诉你全部真相。\par

\SourceRecord{Translated by Blomfield Jackson.；Nicene and Post-Nicene Fathers, Second Series；Vol. 8.；Edited by Philip Schaff and Henry Wace.；Buffalo, NY: Christian Literature Publishing Co.,；1895.；<http://www.newadvent.org/fathers/3202099.htm>.}

% structure: p0001 source=h1 target=section reason=page-title

\section{第一百封信}

\Emph{\Place{凯撒勒雅}的\Person{圣巴西略}}\par

\Emph{致\Place{萨莫萨塔}主教\Person{欧瑟伯}}\par

当我看到您充满友爱的信函时，正值我身处\Place{亚美尼亚}边境地区，这封信如同在远处高举的火炬，为海上的水手指引方向，尤其当海面被风吹动之时。您阁下的信函本身令人愉悦，充满安慰；但它天然的优美因抵达的时刻而大大增强，这一时刻对我来说如此痛苦，以至于我一旦决心忘记其烦恼，便几乎不知如何描述。不过，我的执事会向您全面说明。我的体力完全衰竭，以至于即使最轻微的活动也会带来痛苦。尽管如此，我确实祈祷，借助您的祈祷，我的渴望能得以实现；尽管我的行程因此事耗费了很长时间，导致我自己的教会事务被忽视，从而给我带来了巨大困难。但是，如果在我有生之年，天主恩赐我在我的教会中见到您阁下，那么我真切地希望，即使未来，我也并未完全被排除在天主的恩赐之外。如果可能，我恳求我们之间的这次会面可以在我们每年为纪念真福殉道者\Person{奥普西基}而举行的主教会议上进行，该会议即将于九月七日召开。我因各种忧虑所困，需要您的帮助和同情，无论是关于主教的任命，还是关于\Person{额我略}·\Place{尼撒}的单纯所带给我的麻烦，他正在\Place{安基拉}召集一次主教会议，竭尽全力与我对抗。\par

\SourceRecord{Translated by Blomfield Jackson.；Nicene and Post-Nicene Fathers, Second Series；Vol. 8.；Edited by Philip Schaff and Henry Wace.；Buffalo, NY: Christian Literature Publishing Co.,；1895.；<http://www.newadvent.org/fathers/3202100.htm>.}

% structure: p0001 source=h1 target=section reason=page-title

\section{第一○一号书信}

\Emph{凯撒勒雅的圣巴西略}\par

\Emph{慰藉书}.\par

这是我写给你的第一封信，我本希望它的主题更加光明。若是如此，事情就会如我所愿，因为我希望所有立志生活在真宗教中的人都幸福地度过一生。但是，主以祂不可言喻的智慧安排我们的道路，祂安排这一切为我们的灵魂有益，一方面使你的生活悲伤，另一方面激起像我这样与你因敬主之爱相连的人的同情。因此，当我从我的弟兄们那里得知你所遭遇的事时，我觉得不能不尽我所能给你一些安慰。如果真有可能，我会尽力前往你现在居住的地方。但是，我的健康不佳和目前占据我的事务，使我已经进行的这次旅行对我的教会的利益造成了损害。因此，我决定写信给阁下，提醒你这些苦难并非统治我们的主无故赐给天主的仆人，而是作为我们对于神圣造物主之爱是否真诚的考验。正如运动员通过竞技场上的拼搏赢得桂冠，基督徒也通过试探的考验得以成全，只要我们学会以应有的忍耐和感恩之心接受主所赐予的一切。一切都是因主的爱而安排的。我们不可将所遭遇的任何事视为痛苦，即使目前它影响我们的软弱。也许我们不知道主为何将每一件发生在我们身上的事作为祝福赐给我们，但我们应该相信，所有发生在我们身上的事都是为我们的好处，或是为了我们忍耐的赏报，或是为了我们已领受的灵魂，以免它因在今生停留太久而充满这世界上的邪恶。如果基督徒的望德仅限于今生，那么过早与身体分离可能被视为悲惨的命运；但是，对于爱天主的人而言，灵魂脱离这些身体束缚是我们真正生命的开始，我们为何要像没有望德的人那样悲伤呢？\BibleRef{得前 4:12}。愿你因此得到安慰，不要陷在苦难中，而要显出你胜过它们，并能超越它们。\par

\SourceRecord{Translated by Blomfield Jackson.；Nicene and Post-Nicene Fathers, Second Series；Vol. 8.；Edited by Philip Schaff and Henry Wace.；Buffalo, NY: Christian Literature Publishing Co.,；1895.；<http://www.newadvent.org/fathers/3202101.htm>.}

% structure: p0001 source=h1 target=section reason=page-title

\section{第一〇二封信}

\Emph{\Person{圣巴西略}}\par

\Emph{致\Place{萨塔拉}市民}\par

被你们的迫切请求以及你们全体人民的恳求所感动，我承担了你们教会的职责，并在主前许诺，在我能力范围内，决不缺少你们任何事物。因此，我被迫，正如经上所载，\BibleQuote{如同触摸我眼中的瞳仁}。这样，我对你们的高度敬意使我既不顾念亲情，也不顾念自幼与所述之人的亲密——这份亲密本应比你们的请求对我有更强的要求。我忘记了所有使他亲近可爱的私人考虑，不计较我的人民因失去他的治理而发出的叹息，也不计较他所有亲属的眼泪；更没有挂虑他年迈母亲的痛苦——她只靠他的帮助维持生活。所有这些重大而众多的考虑，我都放在一边，只着眼于一个目标：为你们的教会赐予这样一位贤者治理的祝福，并帮助她——她现在因长期没有首领而痛苦，需要强大有力的支持才能重新站立。以上是我自己方面的。现在，另一方面，我请求你们不要辜负我所怀的希望以及我对他所作的承诺——我已将他送到亲密的朋友那里。我请求你们每一个人都要在爱戴和情谊上超越其余的人。我恳请你们展现这种值得称许的竞争，并以你们对他极大的关怀安慰他的心，使他能忘记自己的家、忘记自己的亲属、忘记如此依赖他治理的人民，如同从母亲怀中断奶的婴孩。\par

我已先派遣\Person{尼基雅}向各位阁下说明一切，好使你们定一个日子欢庆，并感谢主，祂已使你们祈祷得到应允。\par

\SourceRecord{Translated by Blomfield Jackson.；Nicene and Post-Nicene Fathers, Second Series；Vol. 8.；Edited by Philip Schaff and Henry Wace.；Buffalo, NY: Christian Literature Publishing Co.,；1895.；<http://www.newadvent.org/fathers/3202102.htm>.}

% structure: p0001 source=h1 target=section reason=page-title

\section{第一〇三封信}

\Emph{凯撒勒雅的圣巴西略}\par

\Emph{致萨塔拉的民众}.\par

主应允了祂子民的祈祷，藉着我这卑微的器皿，赐给了他们一位名副其实的牧者；他并非像许多人那样以言语牟利，而能够充分满足你们——你们热爱正统教义，并奉主名接受了合乎主诫命的生活——主以祂自己的属神恩宠充满了他。\par

\SourceRecord{Translated by Blomfield Jackson.；Nicene and Post-Nicene Fathers, Second Series；Vol. 8.；Edited by Philip Schaff and Henry Wace.；Buffalo, NY: Christian Literature Publishing Co.,；1895.；<http://www.newadvent.org/fathers/3202103.htm>.}

% structure: p0001 source=h1 target=section reason=page-title

\section{第一〇四书}

\Emph{\Person{圣巴西略}}\par

\Emph{致\Person{莫德斯图}总督。}\par

仅向如此伟大之人致函，即使别无他故，亦当视为莫大荣耀。因与高位者通信，凡获准者均得荣光。然而，今我因全乡境况之极度忧患，不得不向尊驾陈情。恳求您仁慈相待，依您自身品性，向我现已屈膝之乡土伸出援手。我恳求之直接缘由如下：按旧有户籍，天主的圣职人员，即\OriginalTerm{司铎}{presbyters}与\OriginalTerm{执事}{deacons}，皆获豁免。然近期登记官未得阁下授权，竟将彼等列入册籍，惟少数因年迈获免。故我请求您留此善举之纪念，使您美名永存万世：即依旧律，圣职员免于赋役。我并非请求仅对现今在册者个别施恩，因恩惠将传于其继任者，而继任者未必常堪\OriginalTerm{圣职}{sacred ministry}。我建议对圣职员作普遍让步，按公开登记册之形式，使各地教会管理者得以豁免其属下职员。此恩惠必将为尊驾善行带来不朽荣耀，且促使多人祈求皇室。亦实有益于政府，若我们并非普惠全体圣职员，而仅对每遇困窘者施以免役之惠。此乃凡愿知之者皆可明晓，乃我们得自由时所行之道。\par

\SourceRecord{Translated by Blomfield Jackson.；Nicene and Post-Nicene Fathers, Second Series；Vol. 8.；Edited by Philip Schaff and Henry Wace.；Buffalo, NY: Christian Literature Publishing Co.,；1895.；<http://www.newadvent.org/fathers/3202104.htm>.}

% structure: p0001 source=h1 target=section reason=page-title

\section{第一〇五封信}

\Emph{\Person{圣巴西略}（\Place{凯撒勒雅}）}\par

\Emph{致\Person{特伦提乌斯}伯爵的女儿们，女执事们。}\par

来到\Place{萨莫萨塔}时，我曾指望能见到两位阁下，及至希望落空，我实难承受。我说：何时能再与你们为邻？何时你们愿意来我处？这一切，仍须听从主的圣意。至于当前，当我得知吾子\Person{索弗若尼乌斯}要起程前往你们那里，便欣然托他带去此信，以传达我的问候，并告诉你们：赖天主的恩宠，我不断忆念你们，并因你们而感谢主，因为你们是好树上的好枝条，在善工上硕果累累，实在如同荆棘中的百合花。你们身处那些败坏真道之人的可怕悖逆包围中，却不屈从他们的诡计；你们没有放弃宗徒信仰的宣告，没有转而追随当今成功的新奇事物。这岂非深切感谢天主的理由？这岂不应为你们带来极大的声誉？你们已宣认对圣父、圣子和圣神的信仰。切莫放弃这一寄托——圣父是万有之源；圣子是\OriginalTerm{独生子}{Only-begotten}，由祂所生，是\OriginalTerm{真天主}{very God}，\OriginalTerm{完美中之完美}{Perfect of Perfect}，是\OriginalTerm{活的肖像}{living image}，在自身内显示整个圣父；圣神\OriginalTerm{由天主而自有}{having His subsistence of God}，是圣德之源，赋予生命的能力，使人成全的恩宠，藉着祂人得被收养，必朽的成为不朽的，在一切事上——在光荣与永恒、在权能与国度、在主权与神性上——与圣父和圣子相联；这由救恩的圣洗圣事的传统所见证。\par

但所有主张圣子或圣神是受造物，或完全将圣神贬至服役与奴仆地位的人，都远离真理。远离他们的共融。回避他们的教导。他们危害灵魂。若主赐我们相见，我将进一步与你们论述信仰，好使你们藉着圣经的证明，同时看清真理的能力和异端的腐朽。\par

\SourceRecord{Translated by Blomfield Jackson.；Nicene and Post-Nicene Fathers, Second Series；Vol. 8.；Edited by Philip Schaff and Henry Wace.；Buffalo, NY: Christian Literature Publishing Co.,；1895.；<http://www.newadvent.org/fathers/3202105.htm>.}

% structure: p0001 source=h1 target=section reason=page-title

\section{第一〇六书}

\Emph{\Person{圣巴西略}（\Place{凯撒勒雅}）}\par

\Emph{致一位士兵。}\par

我有许多理由感谢天主在我的旅程中所赐予的恩惠，但我认为没有比认识你的卓越更蒙福的事，这是由我们善主的仁慈所赐予的。我认识到一个人，他证明即使在军旅生活中也能保持对天主的完美爱德，并且我们必须分辨基督徒不是由他穿着的样式，而是由他灵魂的倾向。与你相见是我极大的喜乐；如今每当我记起你，我便感到非常欢喜。要刚强；要坚强；努力滋养并增长对天主的爱德，好使祂赐给你更大的福佑。我不需要更多的证据证明你记得我；从你所做的事上我已得到明证。\par

\SourceRecord{Translated by Blomfield Jackson.；Nicene and Post-Nicene Fathers, Second Series；Vol. 8.；Edited by Philip Schaff and Henry Wace.；Buffalo, NY: Christian Literature Publishing Co.,；1895.；<http://www.newadvent.org/fathers/3202106.htm>.}

% structure: p0001 source=h1 target=section reason=page-title

\section{\WorkTitle{第一〇七书}}

\Emph{\Person{圣巴西略} \Place{凯撒勒雅}}\par

\Emph{给寡妇\Person{尤利塔}}。\par

读到夫人您的来信，得知您仍陷于同样的困境，我深感忧虑。对于如此反复无常、出尔反尔、从不坚守承诺之人，能做什么呢？如果他在我本人及前任长官面前许下诺言后，如今却仿佛什么都未曾说过，试图缩短恩宠的期限，那么在我看来，他确实已丧失了一切羞耻心。尽管如此，我还是写信责斥了他，并提醒他所许的承诺。我也给赫拉狄乌斯写了信——他是长官的家臣——以便通过他传递关于您事务的消息。我本人犹豫未敢对如此位高权重的官员贸然行事，因为我从未就私事给他写过信，也担心他做出不利的裁决，如您所知，大人物们往往对此类事情容易恼怒。然而，若此事能有所进展，必是通过赫拉狄乌斯，他是个杰出之人，对我怀有好意，敬畏天主，并能自如地觐见长官。圣者能救您脱离一切苦难，只要我们真诚地将全部希望寄托于祂。\par

\SourceRecord{Translated by Blomfield Jackson.；Nicene and Post-Nicene Fathers, Second Series；Vol. 8.；Edited by Philip Schaff and Henry Wace.；Buffalo, NY: Christian Literature Publishing Co.,；1895.；<http://www.newadvent.org/fathers/3202107.htm>.}

% structure: p0001 source=h1 target=section reason=page-title

\section{第一〇八书}

\Emph{\Person{圣巴西略·凯撒利亚}}\par

\Emph{致\Person{尤利塔}继承人的监护人。}\par

我十分惊讶地听说，在您作出慷慨承诺之后——这些承诺原本符合您宽厚的品格——如今您却忘记了它们，并对我们的姐妹施加暴力和严厉的压力。在这种情况下，我实在不知道该作何感想。我从许多体验过您的慷慨并为之作证的人那里知道，您的慷慨是多么伟大；我也记得您在我和前任总督面前所作的承诺。您说，您在书面中指定了较短的时间，但您愿意给予更长的恩宠期限，因为您希望满足实际需要，并向那位寡妇施恩——她现在被迫立即从自己的财产中支付如此大的一笔款项。我无法想象这种变化的原因。但无论如何，我恳求您记住自己宽厚的品格，并仰望赏报善行的主。我恳求您给予起初所承诺的宽限时间，以便她们能够卖掉财产并还清债务。我清楚地记得，您承诺过，如果您收到商定的款项，您将把所有原定的文件——无论是在长官面前签署的还是私人文件——都归还给寡妇。因此，我恳求您，为了给我面子，也为您自己从主那里获得大祝福。记住您自己的承诺，承认您是人，自己也必须仰望那个您需要天主帮助的时候。不要因您现在的严苛而将自己隔绝于那帮助之外；相反，要通过对受苦者显示一切仁慈和宽容，为您自己吸引天主的怜悯。\par

\SourceRecord{Translated by Blomfield Jackson.；Nicene and Post-Nicene Fathers, Second Series；Vol. 8.；Edited by Philip Schaff and Henry Wace.；Buffalo, NY: Christian Literature Publishing Co.,；1895.；<http://www.newadvent.org/fathers/3202108.htm>.}

% structure: p0001 source=h1 target=section reason=page-title

\section{\WorkTitle{书信第一〇九}}

\Emph{\Person{圣巴西略}}\par

\Emph{致\Person{赫拉迪乌斯伯爵}}\par

因君之高位，吾常不敢以琐事相扰，恐有滥用友誼之嫌；然而事势所迫，不得不陈情。吾之亲属，今为寡妇，哀痛之中，须照料其孤儿之产业。见其不胜重负，吾心甚怜，故急求援手，望君俯允扶持其所遣之人，俾其偿清当面承诺之本金后，不再受逼迫。彼曾同意，还本则免息。然而今管理其继承人产业者，欲并征本息。主，如你所知，亲自将照顾孤儿寡妇视为己事，君亦当尽力而为，冀获天主之报偿。窃思吾等可敬之仁慈总督，闻本金已清，必怜此遭难不幸之家，今已力竭，不能再御所加之损害。恕我冒昧，恳请君按基督所赐之权能；君是良善忠实之人，善用君之才能，助成此事。\par

\SourceRecord{Translated by Blomfield Jackson.；Nicene and Post-Nicene Fathers, Second Series；Vol. 8.；Edited by Philip Schaff and Henry Wace.；Buffalo, NY: Christian Literature Publishing Co.,；1895.；<http://www.newadvent.org/fathers/3202109.htm>.}

% structure: p0001 source=h1 target=section reason=page-title

\section{第一一〇书}

\Emph{\Person{圣巴西略·凯撒勒雅}}\par

\Emph{致总督\Person{莫德斯妥}}。\par

承蒙大人慈祥纡尊降贵前来见我，赐我莫大殊荣，并允我极大自由；我同样，甚至更甚，祈求我们良善的主在整个您的一生中赐予大人这一切。我久已想写信给您，并承您赐予荣誉，但出于对您伟大尊位的敬畏，我始终克制自己，小心谨慎，以免似乎滥用您所赐予的自由。然而现在，我不得不鼓起勇气，这不仅是由于我已获得您无与伦比的卓越大人的许可可以写信，而且也是因为受难者的迫切需要。如果即使是小人物们的祈祷也能对伟人有所裨益，那么，最卓越的先生，请您好意垂怜，赐予如今处境悲惨的农村居民以救济，并下令使由产铁的\Place{托鲁斯}地区的居民所缴纳的铁税变得可以承受。请准许此事，以免他们被彻底压垮，而无法持续为国效力。我相信您那令人钦佩的仁爱必将成全此事。\par

\SourceRecord{Translated by Blomfield Jackson.；Nicene and Post-Nicene Fathers, Second Series；Vol. 8.；Edited by Philip Schaff and Henry Wace.；Buffalo, NY: Christian Literature Publishing Co.,；1895.；<http://www.newadvent.org/fathers/3202110.htm>.}

% structure: p0001 source=h1 target=section reason=page-title

\section{第111封信}

\Emph{\Person{圣巴西略}}\par

\Emph{致总督\Person{莫德斯图}}。\par

在通常情况下，我不敢冒昧打搅阁下，因为我知道如何衡量自己的分量，也懂得尊重高位。但现在，我见到一位朋友因被传唤到您面前而陷入困境，便斗胆送上了这封信。我希望他以此信作为某种和好的象征，能得到仁慈的体恤。确实，尽管我微不足道，但谦逊本身或许能赢得最仁慈的总督的谅解，并为我求得宽恕。这样，若我的朋友并无过错，他仅凭真理的力量便可获救；若他犯了错，也可因我的恳求而得到宽恕。\par

我们在此地的处境，没有人比您更清楚，因为您洞察每个人的弱点，并以您卓越的先见之明治下一切。\par

\SourceRecord{Translated by Blomfield Jackson.；Nicene and Post-Nicene Fathers, Second Series；Vol. 8.；Edited by Philip Schaff and Henry Wace.；Buffalo, NY: Christian Literature Publishing Co.,；1895.；<http://www.newadvent.org/fathers/3202111.htm>.}

% structure: p0001 source=h1 target=section reason=page-title

\section{第一一二号书信}

\Emph{\Place{凯撒勒雅}的\Person{圣巴西略}}\par

\Emph{致\Person{安多尼库斯}将军}。\par

1. 若我的健康状况允许我不费力地出门旅行，并能忍受冬天的严寒，我就不写信，而亲自前往拜谒阁下，这有两个原因。第一，偿还旧债，因为我知道我曾承诺到\Place{塞巴斯蒂亚}去，并乐于拜见阁下；我确实去了，但因到达时间稍晚于阁下而未能与您会面；第二，亲自担任使节，因为迄今为止我一直不敢派人送信，觉得我微不足道，不配获得如此恩惠；同时我认为，仅仅写信，没有人能像亲自面谈那样有效地说服担任公职或私人身份的人，为任何人求情；在面谈中可以澄清指控中的某些要点，为其他事项恳求，并请求宽恕；这些目的都难以通过一封信轻易达成。如今，面对这一切，我只能把一件事物摆出来，就是您最优秀的本人；因为只要向您说明我对事情的想法，您自然会补足一切欠缺，所以我才斗胆写了这封信。\par

2. 但您看到，由于我的犹豫，因为我把说明求情理由的事推迟了，我写得迂回曲折。这个\Person{多米提安}从一开始就是我和我父母的密友，待我如兄弟。我为何不说实话呢？当我得知他陷入目前困境的原因时，我说他罪有应得。因为我希望，任何人无论犯了多大过或小过，都不能逃脱惩罚。但是当我看到他过着不安全、不体面的生活，并感到他唯一的希望取决于您的决定时，我认为他受到惩罚已经够了；所以我恳求您以宽宏和人道的态度看待他的案件。拥有对手的生杀大权，对一个有魄力、有权势的人来说是合适且正当的；但对失势者仁慈温和，则是灵魂伟大和严酷性上超越常人者的标志。因此，如果您愿意，您完全有能力在惩罚他和拯救他时展示您的宽宏大量。让\Person{多米提安}对他所猜疑之事和他自知应受惩罚之事的恐惧，作为他惩戒的全部。我恳求您不要再给他增加惩罚，因为请考虑这一点：古代有许多人，没有记载传下，他们曾将得罪自己的人置于权力之下。但那些在哲学上超越同侪的人，并没有坚持愤怒，他们的记忆一直流传下来，永垂不朽。愿这荣耀加于历史上对您的评价。请赐予我们这些渴望赞美您的人，能够超越古代传唱的仁慈实例。据说，\Person{克罗伊斯}就是这样对杀他儿子的人息怒的，因为那人自首接受惩罚；而伟大的\Person{居鲁士}在战胜后也对这个\Person{克罗伊斯}友好。我们将把您与这些人并列，并尽我们的全力宣扬您的荣耀，除非我们被视为不配为如此伟大人物的传报者。\par

3. 我还应该提出另一个理由：我们惩罚犯错者，并非因为过去的事已无法挽回（因为有什么方法能够弥补过去呢？），而是为了他们将来能改过，或者成为他人循规蹈矩的榜样。现在，没有人能说本案缺少这两点；因为\Person{多米提安}将终生记住所发生的事；而且我认为，其余所有人，以他为戒，都吓得半死。在这种情况下，我们对他惩罚的任何增加，都只会显得是发泄我们自己的愤怒。我认为这远非您的情况。事实上，若不是我看到给予者比接受者得到更大的祝福，我是绝对不会说这样的话的。您的宽宏大量也不会只有少数人知道。整个\Place{卡帕多西亚}都在观望要看事情如何发展，我祈祷他们能把这算作您善行中的一件。我迟迟不敢结束这封信，生怕有任何遗漏对我不利。但我要补充一点：\Person{多米提安}收到了许多为他求情的人的信，但他认为我的信最重要，因为他从某处（我不知道从何处）得知我对阁下有影响力。不要让他寄托在我身上的希望破灭；不要让我在我的人中失去信用；请接受恳求，杰出的先生，赐予我所求之物。您像任何哲学家一样清晰地洞察人生，您知道给所有帮助他人者积存了何等美好的宝藏。\par

\SourceRecord{Translated by Blomfield Jackson.；Nicene and Post-Nicene Fathers, Second Series；Vol. 8.；Edited by Philip Schaff and Henry Wace.；Buffalo, NY: Christian Literature Publishing Co.,；1895.；<http://www.newadvent.org/fathers/3202112.htm>.}

% structure: p0001 source=h1 target=section reason=page-title

\section{第一一三书}

\Emph{\Emph{\Person{圣巴西略}\Place{凯撒勒雅}}}\par

\Emph{致\Place{塔尔索}的司铎}.\par

与这位弟兄会面后，我衷心感谢天主，因祂藉此探望，在诸多苦难中安慰了我，并透过他清楚向我显示了你们的爱德。我似乎从一个人的性情中看到了你们所有人对真理的热忱。他将告诉你们我们彼此交谈的内容。你们应当直接从我这学到的是以下这些。\par

我们生活在似乎教会倾覆迫在眉睫的日子里，这是我一直以来所认识到的。教会没有建树，错误没有纠正，弱小者没有同情，健全的弟兄没有辩护；既找不到医治已缠累我们的疾病的补救，也无法预防我们所预期的。总之，教会的状况（如果我可以用一个虽然看似过于简陋的明喻）就像一件旧外衣，总是被撕裂，再也不能恢复原来的强度。因此，在这样的时刻，需要巨大的努力和勤奋，使教会在某些方面得到裨益。迄今分离的部分若能联合起来是一件好事。如果我们在不损害灵魂的情况下愿意迁就较弱者，联合就能实现；既然许多张口反对\OriginalTerm{圣神}{Holy Ghost}，许多舌头磨快对祂的亵渎，我们恳求你们，尽你们所能，使亵渎者的人数减少，并接纳所有不宣称圣神是\OriginalTerm{受造物}{creature}的人进入共融，使亵渎者被孤立，要么羞愧而回归真理，要么若仍坚持错误，则因人数稀少而不再重要。因此，我们不要再寻求更多，只向所有愿意与我们联合的弟兄提出\OriginalTerm{尼西亚信经}{Nicene Creed}。如果他们同意，我们再要求不应称圣神为受造物，也不接纳任何这样说的人进入共融。我认为我们不应坚持超出此范围的事。因为我相信，通过更长时间的交流和相互经验，若无争斗，若还需要添加任何解释，那 \BibleQuote{为爱祂的人促成万有的主}\BibleRef{罗 8:28} 必会赐予。\par

\SourceRecord{Translated by Blomfield Jackson.；Nicene and Post-Nicene Fathers, Second Series；Vol. 8.；Edited by Philip Schaff and Henry Wace.；Buffalo, NY: Christian Literature Publishing Co.,；1895.；<http://www.newadvent.org/fathers/3202113.htm>.}

% structure: p0001 source=h1 target=section reason=page-title

\section{\WorkTitle{第一一四封信}}

\Signature{\Person{凯撒勒雅的圣巴西略}}\par

\Emph{致 \Person{基里雅科}，于 \Place{塔尔索}。}\par

我几乎不必告知和平之子，和平的祝福是何等伟大。但现在，这一祝福——伟大、奇妙，值得所有爱主的人竭力追求——因邪恶增多，多数人的爱德冷淡，而面临沦为虚名的危险。因此，我认为所有真正侍奉主的人所应追求的惟一伟大目标，就是使那些\BibleQuote{多次并以多种方式}彼此分离的教会重新合一。我自己尝试实现这一点，不能被合理地指责为多管闲事，因为没有什么比作和平缔造者更能体现基督徒的特质，为此，我们的主应许赐予和平缔造者极高的赏报。\par

因此，当我与弟兄们会面，得知他们何等深厚的兄弟之爱、对您的敬意，以及更甚的对基督的爱，还有他们在一切有关信仰之事上的精确和坚定，此外，他们热切追求两个目标：既不与您的爱分离，也不放弃他们纯正的信仰，我便赞赏他们良好的心意；我现在写信给尊驾，恳求您以全部的爱留住他们，使他们与您真正合一，并与您一同为教会忧虑。我已向他们为您的正统性作了担保，并担保您也藉天主的恩宠被登记在册，为真理竭力奋战，无论您可能为真道忍受什么。我本人的意见是，以下条件不会与您自己的感受相悖，且足以使上述弟兄们满足：即，你们应宣认由我们的教父在\Place{尼西亚}聚会时所颁布的信仰，不可省略其中任何一条命题，而要牢记那聚集在一起毫无纷争的三百一十八位教长并非没有圣神的运作，并且不得在该信经中添加圣神是\OriginalTerm{受造物}{creature}的说法，也不得与如此说的人保持共融，为使天主的教会得以纯洁，没有任何毒麦的邪恶混杂。如果您的善意给予他们这充分的保证，他们准备向您给予适当的服从。我自己也为弟兄们承诺，他们不会反对，而会完全顺从，只要阁下乐意应允他们所求的这一件事。\par

\SourceRecord{Translated by Blomfield Jackson.；Nicene and Post-Nicene Fathers, Second Series；Vol. 8.；Edited by Philip Schaff and Henry Wace.；Buffalo, NY: Christian Literature Publishing Co.,；1895.；<http://www.newadvent.org/fathers/3202114.htm>.}

% structure: p0001 source=h1 target=section reason=page-title

\section{第一一五函}

\Emph{圣巴西略·凯撒勒雅}\par

\Emph{致异端者\Person{辛普利基雅}}\par

我们常常不明智地憎恨上司，疼爱下属。因此我本人缄默不言，对我所受侮辱的羞辱保持沉默。我等待至高的审判者，祂最终知道如何惩罚一切邪恶，即使人倾出金子如沙；让他践踏正义吧，他只是伤害自己的灵魂。天主常要求祭献，我想并非因为祂需要它，而是因为祂接纳虔诚正直的心作为珍贵的祭献。但当一个人因他的过犯践踏自己时，天主就视他的祈祷为不洁。那么，请你思念最后的一天，请不要试图教导我。我知道的比你多，而且没有如此被内心的荆棘堵塞。我不介意十倍的邪恶伴有少许善行。你激起蜥蜴和蟾蜍反对我，确实是春天的野兽，但无论如何是不洁的。但天上的鸟会来吞吃它们。我所要交账的，不是照你的想法，而是照天主认为合适来审判。如果需要证人，站到审判者面前的不会是奴隶；也不是一群可耻可憎的宦官；既不是女人也不是男人，好色、嫉妒、受贿、暴躁、柔弱、肚腹的奴隶、疯狂求金、无情、抱怨晚餐、反复无常、吝啬、贪婪、不知足、野蛮、嫉妒。还需要我说什么？他们一出生就注定被刀割。他们的脚歪曲，心怎能正直？他们因刀而贞洁，这算不得什么功劳。他们无目的地好色，出于本性的卑劣。这些不是站到审判中的证人，而是义人的眼目和成全者的目光，是所有那时将用眼看见他们如今以理智所见之人。\par

\SourceRecord{Translated by Blomfield Jackson.；Nicene and Post-Nicene Fathers, Second Series；Vol. 8.；Edited by Philip Schaff and Henry Wace.；Buffalo, NY: Christian Literature Publishing Co.,；1895.；<http://www.newadvent.org/fathers/3202115.htm>.}

% structure: p0001 source=h1 target=section reason=page-title

\section{\WorkTitle{第116号书信}}

\Emph{\Person{圣巴西略}（\Place{凯撒勒雅}）}\par

\Emph{致\Person{菲尔米尼乌斯}。}\par

你很少写信，信也很短，也许是因为你懒得写信，或是为了避免过度的饱足，又或者是为了训练自己简短的言语。我确实从不满足，无论你的沟通多么丰富，都小于我的渴望，因为我希望知晓关于你的每一个细节。你健康如何？苦修纪律如何？你是否坚持你最初的目标？还是根据情况改变了主意？如果你保持不变，我就不需要大量的信。我会非常满足于\Quote{我很好，也希望你很好。}但我听到令我羞愧的事：你离开你有福的祖先的行列，投靠你父系的祖父，渴望成为一个\Person{布雷塔尼乌斯}，而不是\Person{菲尔米尼乌斯}。我非常渴望听到这件事，并了解是什么原因导致你过这种生活。你自己保持沉默；我想你是为你的意图感到羞愧，因此我必须恳求你不要怀有任何与羞耻相关的计划。如果任何这样的想法进入你的头脑，把它抛开，重新振作，永远告别军旅、武器和营地的辛劳。回家去，像你祖先以前所想的那样，认为在城中获得高位就足以获得安逸的生活和一切可能的荣誉。我确信，考虑到你的天赋和很少的竞争者，你将毫不费力地实现这一点。如果这不是你最初的意思，或者形成后又放弃了，就立刻告诉我。如果相反——愿天主不许——你仍持同样的想法，就让麻烦自行到来。我不想收到信。\par

\SourceRecord{Translated by Blomfield Jackson.；Nicene and Post-Nicene Fathers, Second Series；Vol. 8.；Edited by Philip Schaff and Henry Wace.；Buffalo, NY: Christian Literature Publishing Co.,；1895.；<http://www.newadvent.org/fathers/3202116.htm>.}

% structure: p0001 source=h1 target=section reason=page-title

\section{\WorkTitle{第一一七封信}}

\Emph{\Person{凯撒勒雅的圣巴西略}}\par

\Emph{无收信人。}\par

因为许多理由，我知道自己是尊驾的债务人；如今我身处焦虑之中，这迫使我不得不寻求此类服务，尽管我的顾问们不过是偶然相遇的人，不似您那样与我有许多不同的联系。无需回顾过去。我可以说是自己困难的根源，因为我决心离开那唯一导向救恩的良好纪律。结果是，在这次困境中，我很快就陷入了诱惑。所发生的事似乎值得一提，以免我再次陷入类似的苦恼。至于未来，我愿向尊驾充分保证，靠天主的恩宠，一切都会顺利，因为这项程序是合法的，并且没有任何困难，因为我在宫廷的许多朋友都愿意帮助我。因此，我将草拟一份请愿书，类似提交给代理官的那份格式；如果没有延误，我将很快获得解脱，并一定会将正式文件寄给您，使您感到欣慰。我相信，在此事上，我自己的信念比皇帝的命令更有力量。如果我在最高生活中显示出这种坚定而稳固的信念，靠着天主的助佑，我贞洁的持守将是不可侵犯和确定的。我很高兴看到您托付给我的兄弟，并把他视为我的亲密朋友。我相信他能证明自己堪当天主和您的美言。\par

\SourceRecord{Translated by Blomfield Jackson.；Nicene and Post-Nicene Fathers, Second Series；Vol. 8.；Edited by Philip Schaff and Henry Wace.；Buffalo, NY: Christian Literature Publishing Co.,；1895.；<http://www.newadvent.org/fathers/3202117.htm>.}

% structure: p0001 source=h1 target=section reason=page-title

\section{第一百一十八封信}

\Emph{圣巴西略}\par

\Emph{致佩拉主教约维努斯。}\par

你欠我一份人情。因为我曾借给你一份善意，现在理当连本带利收回——这种利息，我们的主并不拒绝。那么，我的朋友，请通过来访来偿还我吧。以上是本钱；利息呢？利息就在于你这位来访者，你比我优胜的程度，如同父亲胜于子女。\par

\SourceRecord{Translated by Blomfield Jackson.；Nicene and Post-Nicene Fathers, Second Series；Vol. 8.；Edited by Philip Schaff and Henry Wace.；Buffalo, NY: Christian Literature Publishing Co.,；1895.；<http://www.newadvent.org/fathers/3202118.htm>.}

% structure: p0001 source=h1 target=section reason=page-title

\section{第一一九封书信}

\Signature{圣巴西略}\par

\Emph{致塞巴斯特亚主教尤斯塔修斯。}\par

我藉着极为尊贵可敬的兄弟\Person{伯多禄}致信于您，恳求您现在并时常为我祈祷，使我得以远离危险当避的道路，终有一日堪当承受基督之名。我虽不言，您二人必谈论我的事，他将向您详述所发生的一切。然而，您未加审察便轻信那卑劣的猜疑，这猜疑必来自那些曾侮辱我——违反正直人的敬畏天主与顾及人的看法——的人。我羞于告诉您我从那位杰出的\Person{巴西略}那里所受的对待，我曾接受他作为您尊者为我生命提供的保护。但当您听到我们的兄弟所言，便会知道一切细节。我如此说并非为报复他，因为我祈祷主不要归罪于他，而是为使我与您之间的爱心保持坚定，因我担心它会被那些人在为自己的堕落辩护时必然捏造的恶毒诽谤所动摇。无论他们提出什么指控，我希望您的明智会向他们质询这些问题：他们是否曾正式控告我？他们是否曾寻求纠正他们所指摘的错误？他们是否曾明确陈述对我的不满？事实上，他们卑怯的逃避已表明，在他们欢愉的面容和虚假的情感表达下，内心深处一直隐藏着极深的诡诈与怨恨。对于这一切，无论我是否叙述，您的明智都完全清楚他们带给我何等悲伤，以及带给那些始终对此悲惨城市中的虔诚生活表示憎恶的人何等嘲笑——他们断言德行的伪装不过是骗取信任的伎俩、欺骗的假象。因此如今在此地，没有任何生活方式像苦修的宣认那样被怀疑为恶习。您的明智将考虑这一切的最佳疗方。\par

至于\Person{索弗罗尼}对我捏造的指控，它们远非祝福的前奏，而是分裂与分离的开端，甚至可能使我的爱心变得冷淡。我恳求您以慈悲的良善阻止他有害的努力，愿您的爱心努力加固那正在破裂的纽带，而非加入那些渴望分裂的人而加深分离。\par

\SourceRecord{Translated by Blomfield Jackson.；Nicene and Post-Nicene Fathers, Second Series；Vol. 8.；Edited by Philip Schaff and Henry Wace.；Buffalo, NY: Christian Literature Publishing Co.,；1895.；<http://www.newadvent.org/fathers/3202119.htm>.}

% structure: p0001 source=h1 target=section reason=page-title

\section{书信第一二〇号}

\Emph{凯撒勒雅的圣巴西略}\par

\Emph{致安提约基雅主教梅肋提乌}\par

我收到了极敬天主的欧瑟比主教的一封信，他在信中嘱咐就某些教会事务再写一封信给西方人。他表达了意愿，希望此信由我起草，并由所有共融者签署。我无法就他的这些意愿写信，因此将他的备忘录寄给您的圣德，以便您在阅读并留意我们同僚司铎、至亲的弟兄桑克提西慕斯所提供的信息后，您本人能仁爱地就这些要点起草一封信，按您认为最合适的方式。我们准备同意此信，并尽快将其送交与共融者，以便所有人都签署后，由那位即将启程拜访西方主教的信使携带。请下令尽快将您圣德的决议告知我，以免我不知晓您的意向。\par

至于目前在安提约基雅正在策划或已经策划的针对我的阴谋，同一位弟兄将向您的圣德传达消息，除非事件的报告先于他而澄清了情况。我们有理由希望那些威胁即将结束。\par

我希望您敬爱知道，我们的弟兄安提慕斯未经投票，又代替我们可敬的弟兄济利禄，任命了与教宗同住的福斯图斯为主教。如此，他使亚美尼亚充满了分裂。我认为应当将此告知您的敬爱，以免他们诬陷我，让我对这些混乱的举动负责。您自然认为应将此事告知其他人。我认为这种不合规矩的行为会使许多人忧心。\par

\SourceRecord{Translated by Blomfield Jackson.；Nicene and Post-Nicene Fathers, Second Series；Vol. 8.；Edited by Philip Schaff and Henry Wace.；Buffalo, NY: Christian Literature Publishing Co.,；1895.；<http://www.newadvent.org/fathers/3202120.htm>.}

% structure: p0001 source=h1 target=section reason=page-title

\section{\WorkTitle{书信一二一}}

\Emph{\Place{凯撒勒雅}的圣\Person{巴西略}}\par

\Emph{致\Place{尼科波利斯}主教\Person{德奥多图斯}}\par

冬季严寒而漫长，以至于我甚至难以获得信函的慰藉。为此，我很少致函尊驾，也少闻尊驾音讯。但现在，我亲爱的兄弟、共司铎\Person{桑克蒂西穆斯}启程前往贵城。藉他，我向阁下致敬，并恳请您为我祈祷，且垂听\Person{桑克蒂西穆斯}之言，好从他得知教会所处的境况，并对向您提出的诸事予以最大关注。您须知，\Person{福斯图斯}带着教宗给我的信函前来，请求祝圣为主教。然而，当我向他索取来自您本人及其他主教的推荐信时，他轻视了我，径自前往\Person{安提穆斯}处。他回来了，被\Person{安提穆斯}祝圣，而此事未向我作任何通报。\par

\SourceRecord{Translated by Blomfield Jackson.；Nicene and Post-Nicene Fathers, Second Series；Vol. 8.；Edited by Philip Schaff and Henry Wace.；Buffalo, NY: Christian Literature Publishing Co.,；1895.；<http://www.newadvent.org/fathers/3202121.htm>.}

% structure: p0001 source=h1 target=section reason=page-title

\section{第一百二十二封信}

\Emph{\Place{凯撒利亚}的圣巴西略}\par

\Emph{致\Person{颇厄墨尼乌斯}}，\Emph{\Place{萨塔拉}主教}\par

当亚美尼亚人经由你处返回时，你无疑曾向他们索要信件，并得知了我为何没有将信交给他们。如果他们像爱慕真理者那样说话，你当时就原谅了我；倘若他们有所隐瞒（我并不这样认为），无论如何，请听我说明。\par

最显赫的\Person{安提摩斯}，早已与我修和，却趁满足自己虚荣并使我烦恼之机，凭其权威亲手祝圣了\Person{福斯图斯}，未等你们任何选举，并讥讽我在此类事上的拘泥小节。既然如此，他既混淆了古老秩序，又轻视了你们——我正在等待你们的选举——并且以我看来，他的行为令天主不悦，因此我对他们感到痛心，没有给任何亚美尼亚人带信，甚至也没有给您大人的。我甚至连共融也不接纳\Person{福斯图斯}，以此明确表明：除非他给我带来您的信，我将永远与他疏远，并影响与我同心的人同样待他。如果这些事还有补救，请务必亲自写信给我，为他作证，若您见他的生活良好；并劝勉其余的人。反之，如果伤害不可救药，请让我确切了解，好使我不再把他们放在心上；尽管事实上，他们已证明，今后同意将他们的共融转向\Person{安提摩斯}，以示藐视我和我的教会，好像我的友谊不再值得拥有。\par

\SourceRecord{Translated by Blomfield Jackson.；Nicene and Post-Nicene Fathers, Second Series；Vol. 8.；Edited by Philip Schaff and Henry Wace.；Buffalo, NY: Christian Literature Publishing Co.,；1895.；<http://www.newadvent.org/fathers/3202122.htm>.}

% structure: p0001 source=h1 target=section reason=page-title

\section{第一二三封信}

\Emph{\Person{圣巴西略（凯撒勒雅）}}\par

\Emph{致隐修士\Person{乌尔比奇乌斯}}。\par

你本应来见我（祝福已经临近），以你的指尖凉爽我，因我在诱惑中如火焚。但后来如何？我的罪从中作梗，阻碍了你的启程，以致我身患无药可医之疾。正如波浪在我们周围，一个沉下，另一个涌起，又有一个阴森可怖地显现，我的烦恼也是如此：有些已经平息，有些正与我同在，有些还在前方。通常来说，这些烦恼的唯一良药便是向危机让步，远离迫害我的人。然而，请来我这里，安慰我，劝诫我，甚至与我同行；无论如何，单是见到你，就会使我好转。最重要的是，祈祷，再祈祷，愿我的理智不被烦恼的波浪所淹没；愿我自始至终持守一颗悦乐天主的心，不与被列为邪恶的仆人，他们只在主人赐予好处时感恩，却在主人以逆境责罚他们时拒绝顺服；但愿我从试炼中获益，在我最需要天主时，最信赖祂。\par

\SourceRecord{Translated by Blomfield Jackson.；Nicene and Post-Nicene Fathers, Second Series；Vol. 8.；Edited by Philip Schaff and Henry Wace.；Buffalo, NY: Christian Literature Publishing Co.,；1895.；<http://www.newadvent.org/fathers/3202123.htm>.}

% structure: p0001 source=h1 target=section reason=page-title

\section{第一二四号书信}

\Signature{圣\Person{巴西略}·\Place{凯撒勒雅}}\par

致\Person{德奥多禄}。\par

常有人说，为爱欲之情所奴役的人，当因某种不可避免的必要而与所渴望的对象分离时，若偶能瞻仰所爱之人的肖像，便可藉满足视觉来抑止其情欲的猛烈。此事是否真实，我不敢断言；但我的朋友，就你而言，我所遭遇的，也与此无甚差别。我对于你虔诚无伪的灵魂，若可如此说，感到一种近乎爱慕之情；然而我欲望的满足（如同一切其他福佑）却因我罪的阻碍而变得困难。然而，在我至为可敬的弟兄们面前，我似乎看到了一个非常像你的肖像。倘若命运使我在远离他们时遇见你，我便会以为在他们身上看到了你。因为你们各自的爱是如此深厚，以至于你们二人之间显然在争夺优越。为此我感谢了天主。若我还有更长的余生，我祈求因着你们，我的生命得以甜美；正如现在我视生命为可厌的、应逃避的可怜之物，因为我与最心爱之人的陪伴分离了。因为依我判断，当人与真正爱我们的人分离时，没有什么是能令人愉悦的。\par

\SourceRecord{Translated by Blomfield Jackson.；Nicene and Post-Nicene Fathers, Second Series；Vol. 8.；Edited by Philip Schaff and Henry Wace.；Buffalo, NY: Christian Literature Publishing Co.,；1895.；<http://www.newadvent.org/fathers/3202124.htm>.}

% structure: p0001 source=h1 target=section reason=page-title

\section{第一二五函}

\Emph{圣巴西略}\par

\Emph{这份信德抄录由圣巴西略口授，并由塞巴斯特主教欧斯塔修签署。}\par

1. 无论是那些先入为主地持有异端信经、如今希望改宗正统信仰的人，还是那些初次渴望接受真理教导的人，都必须学习在尼西亚召开的大公会议上由可敬教父们制定的信经。同样的教导对于所有被怀疑对纯正教义怀有敌意、并借助巧妙而貌似合理的借口隐藏其邪恶观点的人，也极为有益。因为对他们而言，这信经就是所需的一切。他们要么被治愈隐藏的不健全身心，要么继续隐藏，从而自己承担因不诚实而应受的判决，并在审判之日为我们提供辩护，那时主将揭开黑暗隐藏之事，并\BibleQuote{显明人心中的计谋}（\BibleRef{格前 1:5}）。因此，不仅应当接受他们宣认相信我们教父在尼西亚所陈述的言辞，也应按那些言辞所表达的纯正意义来接受。因为有人甚至在这信经中歪曲真理之言，按自己的观念曲解其中词语的含义。例如，\Person{玛尔塞鲁}对我们的主耶稣基督的位格发表不虔的见解，将他描述为\OriginalTerm{圣言}{Logos}而非其他，竟敢声称从那信经中为他的原则找到借口，给\OriginalTerm{同质}{Homoousion}一词附上不当的意义。此外，利比亚的\Person{撒伯里乌}的一些不虔追随者，认为\OriginalTerm{位格}{hypostasis}和\OriginalTerm{本体}{substance}是同一的，也从同一来源为建立他们的亵渎找到根据，因为信经中写着\Quote{若有人说圣子具有不同的本体或位格，公教会与宗徒教会弃绝他}。但他们并未在那处陈述位格和本体是同一的。如果这两个词表达同一含义，何需两者？相反，显然有些人否认圣子与圣父同体，另一些人声称他不在本体之内，而是具有其他位格，因此他们谴责这两种意见都不属于教会所持的观点。当他们陈述自己的观点时，他们宣称圣子出于圣父的本体，但没有加上\Quote{出于位格}的字样。前一句用于谴责错误观点；后一句明确陈述了救恩的教义。因此，我们有义务宣认圣子与圣父同体，正如所记载的；但同时，圣父存在于他自己的位格中，圣子在他的位格中，圣神在他的位格中，正如他们自己清楚传授的教义。他们确实在\OriginalTerm{光之光}{Light of Light}这些词中明确而充分地宣告：生光者和受生光者是有区别的，然而都是光；所以，本体的定义是同一且相同的。现在我将附上在尼西亚制定的实际信经。\par

2. \Quote{我们信唯一的天主，全能的圣父，一切有形无形之物的创造者}\par

\SourceRecord{Translated by Blomfield Jackson.；Nicene and Post-Nicene Fathers, Second Series；Vol. 8.；Edited by Philip Schaff and Henry Wace.；Buffalo, NY: Christian Literature Publishing Co.,；1895.；<http://www.newadvent.org/fathers/3202125.htm>.}

% structure: p0001 source=h1 target=section reason=page-title

\section{书信一二六}

\Emph{\Person{圣巴西略（凯撒勒雅）}}\par

\Emph{致\Person{阿塔比乌}。}\par

我抵达\Place{尼科波利斯}，怀抱着双重希望：解决已发生的纷扰，并尽可能矫正那些违反教会法律、混乱无序的措施，却因未能与你见面而深感失望。我听闻你已仓促离开，甚至就是从你自己正在主持的会议中退出的。因此，我不得不诉诸笔墨，以此信命你到我面前来，好让你亲自缓解我心中的痛苦（这痛苦几乎致命），因为我听说你竟敢在教堂中央做出我至今闻所未闻的行为。这一切，虽然痛苦严重，但尚可忍受，因为它是发生在一个已将受苦应有的惩罚交托于天主、并全心致力于和平、防止因自己的过失而带给天主子民任何伤害的人身上。然而，由于几位可敬的、完全可信的弟兄告诉我，你在信仰中引入了某些新异之说，并发表了反对健全道理的言论，在此情况下我更加不安，忧虑至极，唯恐除了那些背叛福音真理的人给教会造成的无数创伤之外，还会有一场新的灾祸兴起，即\Person{撒伯略}——这位教会的敌人——的古代异端死灰复燃；因为弟兄们报告说，你的言论与此相似。因此，我写信吩咐你，不要害怕行程短途来见我，并在这事上给我完全的保证，以缓解我的痛苦，并安慰天主的各教会，因为它们因你的行为和你所传的言语，已感到严重甚至无法忍受的痛苦。\par

\SourceRecord{Translated by Blomfield Jackson.；Nicene and Post-Nicene Fathers, Second Series；Vol. 8.；Edited by Philip Schaff and Henry Wace.；Buffalo, NY: Christian Literature Publishing Co.,；1895.；<http://www.newadvent.org/fathers/3202126.htm>.}

% structure: p0001 source=h1 target=section reason=page-title

\section{\WorkTitle{第一二七号书信}}

\Emph{\Person{凯撒勒雅的圣巴西略}}\par

\Emph{致\Place{萨莫萨塔}主教\Person{欧瑟比}}。\par

我们的仁慈天主，祂使安慰与患难相称，并安慰卑微的人，免得他们不知不觉在过度的忧伤中被淹没，藉着适时地带来蒙主钟爱的主教\Person{约比诺}，赐下与我在\Place{尼科波利斯}所遭受的患难相等的安慰。他必须亲自告诉您他的来访是多么及时。我不愿写长信，所以保持沉默。我更倾向于沉默，免得我似乎给那些已经转变并开始对我表示尊敬的人贴上标记，提及他们的跌倒。\par

愿天主恩赐您能来到我的家，以便我能拥抱您的尊驾并详细告诉您一切。因为我们时常发现，讲述痛苦的实际经历能带来一些安慰。然而，无论如何，请称颂那位极其敬虔的主教所做的一切，完全符合他对我的爱，尤其是在严格遵守教会法规方面。此外，感谢天主，您的学生在各处都展现出您尊驾的品格。\par

\SourceRecord{Translated by Blomfield Jackson.；Nicene and Post-Nicene Fathers, Second Series；Vol. 8.；Edited by Philip Schaff and Henry Wace.；Buffalo, NY: Christian Literature Publishing Co.,；1895.；<http://www.newadvent.org/fathers/3202127.htm>.}

% structure: p0001 source=h1 target=section reason=page-title

\section{第128封信}

\Emph{\Place{凯撒勒雅}的\Person{圣巴西略}}\par

\Emph{致\Place{萨莫萨塔}主教\Person{欧瑟伯}。}\par

1. 迄今为止，我一直未能以实际行动充分证明我热切渴望使主的各教会得享平安。但我在心里肯定，我有如此强烈的渴望，以至于我乐意甚至舍弃自己的生命，只要由此能扑灭由恶者点燃的仇恨火焰。若我不是为了这寻求平安的渴望而同意来到\Place{科洛尼亚}，愿我的生命不蒙平安之福。我所寻求的平安是真正的平安，是主亲自留给我们的；我所请求并得以确认的保证，属于那除了真平安之外一无所求的人，虽然有些人曲解真理为另一种意义。任凭他们随意使用舌头，但可以肯定，他们终有一天会为自己的话后悔。\par

2. 现在我恳求你的圣德记住最初的提案，不要被那些答非所问的回答所误导，也不要给予那些毫无辩论能力却非常巧妙地根据自己的想法曲解真理的人以实际重视。我提出了一些非常简单、清晰且易于记住的提案：我们是否拒绝与那些不肯接受\OriginalTerm{尼西亚信经}{Nicene Creed}的人共融？我们是否拒绝与那些胆敢宣称圣神是受造物的人有分？然而，他并没有逐字回答我的问题，而是编造了你寄给我的声明——这并非出于愚钝（如可能想象的那样），也不是因为他无法看到后果。他以为，通过驳斥我的提案，他将向人们暴露自己的真面目；而如果他同意，他将偏离那条迄今为止在他看来比任何其他立场都更可取的\OriginalTerm{中间道路}{via media}。愿他不要试图欺骗我，也不要与其他人一起欺骗你的智慧。愿他就我的问题给出一个简洁的回答，即他是否接受还是拒绝与信仰的敌人共融。如果你能使他这样做，并给我一个我所祈求的明确答复，我承认自己在之前所有事情上都错了；我将一切责任归咎于自己；然后向我求一个谦卑的证明。但如果没有这样的事发生，请原谅我，最敬爱天主的父，我无法以虚伪接近天主的祭台。若不是因为这敬畏，我为何要与\Person{欧依丕乌斯}分离，他是如此博学，如此年迈，又被那么多情感纽带与我相连？然而，如果在这种情况下我做得对，我相信，通过那些和蔼可亲的人的中介，与那些持有与\Person{欧依丕乌斯}相同观点的人联合，将是荒谬的。\par

3. 我不是认为我们有绝对的责任与那些不接受信仰的人断绝关系，而是要按照爱的古老法律照顾他们，并一致写信给他们，以怜悯给予他们一切劝诫，并向他们提出父辈的信仰，邀请他们合一。如果我们成功，就应该与他们共融；如果我们失败，就必须满足于彼此，并摒弃这种不确定的精神，恢复那些从起初就接受圣言的人所遵循的福音和简朴的交往。据说，\BibleQuote{他们，}据说，\BibleQuote{都是一心一意。}\BibleRef{宗 4:32} 如果他们听从你，这是最好的；如果不听从，就要认清战争的真正制造者，今后不要再给我写任何关于和解的信。\par

\SourceRecord{Translated by Blomfield Jackson.；Nicene and Post-Nicene Fathers, Second Series；Vol. 8.；Edited by Philip Schaff and Henry Wace.；Buffalo, NY: Christian Literature Publishing Co.,；1895.；<http://www.newadvent.org/fathers/3202128.htm>.}

% structure: p0001 source=h1 target=section reason=page-title

\section{第一二九封信}

\Emph{\Person{圣巴西略·凯撒勒雅}}\par

\Emph{致\Place{安提约基雅}主教\Person{默肋提乌斯}}\par

1. 我知道最近对多嘴的\Person{阿波利纳里}提出的指控，在阁下听来必定感到奇怪。我自己直到现在才知道他受到了指控；然而，当前，\Place{塞巴斯特人}在某处搜罗之后，将这些事情拿了出来，并且他们正流传着一份文件，他们特别试图据此指控我持有相同的主张。文件中包含以下语句：\Quote{因此，处处都必须将第一次同一性理解为与第二次同一性相连，更确切地说，与第三次同一性合一，并宣称第二次与第三次是相同的。因为圣父首先是什么，圣子其次就是什么，圣神第三就是什么。反之，圣神首先是什么，圣子其次就是什么，只要圣神是主；圣父第三就是什么，只要圣神是天主。并且，为最有力地表达那不可言说的，圣父是按父义而言的圣子，圣子是按子义而言的圣父，对于圣神也是如此，只要天主圣三是独一天主。}这就是正在流传的言论。我绝不相信这是由那些公布它的人杜撰的，尽管在他们对我的诽谤之后，我认为没有什么超出他们的胆量。因为写信给他们同党中的某些人时，他们对我提出了虚假的指控，然后加上了我引用的那些话，称其为异端者的作品，却对文件的作者只字不提，以便让普通人认为它是出自我的手笔。然而，在我看来，他们的智慧不足以将这些短语组合在一起。为此，为了驳斥针对我日益增长的亵渎，并向全世界表明我与那些说这些话的人毫无共同之处，我不得不提到\Person{阿波利纳里}，因为他近似于\Person{撒伯里乌}的邪恶。关于这个话题，我不再多说。\par

2. 我从宫廷收到消息，在皇帝最初的冲动（这是由我的诽谤者促成的）之后，通过了第二道法令，即我不得被交给我的控告者，也不得任由他们处置，如同最初命令的那样；但在此期间有些拖延。如果这一法令得以保持，或者决定了任何更温和的措施，我会通知你。如果前者占据上风，那也不会是你不知道的。\par

3. 我们的兄弟\Person{桑克提西穆斯}一定在你那里已经很久了，你已了解他此行的目的。因此，如果致西方人的信对你来说似乎包含了任何必要的东西，请费心将它写出来并转交给我，以便我能让与我们意见相同的人签名，并保持签署预备好，且写在另一张纸上，我们可以将其附在我们的弟兄和同司铎所携带的信上。由于我在备忘录中没有发现任何结论性的东西，我对于该向西方人写些什么感到困难。必要之点已被预先处理，写多余的东西是无用的，而且在这类问题上表示情绪岂不荒谬？然而，有一个主题在我看来至今尚未触及，并提供了一个写作的理由，那就是劝告他们不要不加区别地接受来自东方的共融；而是在选择一边之后，根据同证者的见证接纳其余的人，并且不要同意每个以正统为借口撰写\OriginalTerm{信经}{creed}的人。如果他们这样做，他们将发现与互相争斗的人共融，这些人常常提出相同的公式，却彼此激烈争战，如同那些最疏远的人。因此，为使\OriginalTerm{异端}{heresy}不至更加蔓延，而那些彼此争执的人相互反对各自的公式，应当劝告他们区分偶然到来者带给他们的共融行为，以及根据教会规则适当起草的共融行为。\par

\SourceRecord{Translated by Blomfield Jackson.；Nicene and Post-Nicene Fathers, Second Series；Vol. 8.；Edited by Philip Schaff and Henry Wace.；Buffalo, NY: Christian Literature Publishing Co.,；1895.；<http://www.newadvent.org/fathers/3202129.htm>.}

% structure: p0001 source=h1 target=section reason=page-title

\section{第一三〇書}

\Emph{凱撒勒雅的聖巴西略}\par

\Emph{致尼科波利的德奧多特主教。}\par

1. 可敬愛的主內弟兄，您非常正確且合理地責備我，自從我離開尊駕，將那些關於信德的建議傳遞給尤斯塔修斯以來，關於他的事情，我對您可說是大大小小一概未曾提及。這種疏忽確實不是因為我輕視他對待我的方式，而僅僅是因為這件事如今已傳遍眾人之耳，無人需要我的教導來了解他的意圖。因為他非常留意，彷彿唯恐對他的意見見證太少，便將他針對我所寫的書信寄到天涯海角。因此，他斷絕了與我的共融。他未曾同意在約定地點與我會面，也沒有像他所許諾的那樣帶領他的門徒。相反地，他在公開的會議中，與西利西亞的德奧菲魯斯一道，公開且毫不掩飾地誹謗我，說我在眾人的靈魂中播下了與他自己教導相悖的學說。這已足以破壞我們之間的一切聯合。後來他去了西利西亞，遇見一位名叫格拉西烏斯的人，便向他展示了一篇信經，而這信經只有亞略派或亞略的徹底門徒才能簽署。那時，我確實更加堅定了與他分離的立場。我感覺，埃塞俄比亞人不會改變他的膚色，豹子也不會改變牠的斑點，一個在乖謬學說中長大的人，也絕不可能洗去他異端的污點。\par

2. 除此之外，他還厚顏無恥地撰文反對我，或者說，寫了長篇大論，充滿各種辱罵和誹謗。對於這些，我至今未作任何回應，這是因為我們受到宗徒的教導，不要為自己報仇，而要給憤怒留有餘地。\BibleRef{羅 12:19} 再者，想到他一直以來接近我時所懷的深刻虛偽，我甚至驚愕得說不出話來。但是，即便這一切從未發生，面對這種新的狂妄之舉，誰能不感到恐懼和厭惡呢？如今，據我聽聞——若這消息屬實而非誹謗——他竟膽敢重新祝聖某些人；這種行徑，迄今為止連任何異端者都沒有膽敢嘗試。那麼，我如何能平靜地忍受這樣的對待？我如何能將此人的錯誤視為可治癒的？因此，請警惕，不要被謊言引入歧途；不要被那些傾向於從壞處看一切的人們的猜疑所動搖，就好像我輕視這些事一樣。因為，我極親愛且可敬的朋友，請確信，我從未像現在這樣，聽到教會律法的混亂而感到如此悲傷。唯願主恩賜我不要在憤怒中採取任何行動，而是持守愛德——它不作非禮之事，不自我膨脹。請看，沒有愛德的人是如何超越一切人的界限，行為不合體統，做出了整個過去歷史上未有先例的膽大妄為之事。\par

\SourceRecord{Translated by Blomfield Jackson.；Nicene and Post-Nicene Fathers, Second Series；Vol. 8.；Edited by Philip Schaff and Henry Wace.；Buffalo, NY: Christian Literature Publishing Co.,；1895.；<http://www.newadvent.org/fathers/3202130.htm>.}

% structure: p0001 source=h1 target=section reason=page-title

\section{\newline{}\WorkTitle{书信一三一}}

\newline{}\Emph{\Person{圣凯撒勒雅的巴西略}}\par

\newline{}\Emph{致\Person{奥林比乌}}\par

\newline{}1. 真料想不到的消息令人双耳发鸣。我正是如此。那些针对我的作品正在四处流传，它们落入了我早已相当习以为常的耳中，因为我先前曾收到那封信——这信配得上我的罪过，但我从未料到它竟会由那些寄出它的人所写。然而，随后所发生的事，在我看来，竟是如此出奇的残忍，以致抹去了之前的一切。当我读到那封致可敬的弟兄\Person{达齐纳}的信时，我怎能不几乎失去理智？那封信充满了无礼的侮辱、诽谤以及对我的攻击，仿佛我已被证实对教会怀有许多有害的阴谋。此外，他们立即拿出了证据，证明那些针对我的诽谤是真实的，这些证据来自一份我不知道作者是谁的文件。我承认，我认出其中一部分是由劳狄刻雅的\Person{阿波里纳}所写。这些部分我甚至有意从未读过，只是从他人的传闻中听说过。我还发现其中包含了其他部分，这些部分我从未读过，也从未从任何其他人那里听说过；对此，天上有一位忠实的见证者。那么，那些回避谎言、知道爱是法律的满全、声称要承担软弱者重担的人，怎能同意将这些诽谤加诸于我，并依据他人的著作来定我的罪呢？我常常自问这个问题，但我想不出原因，除非如我一开始所说，我这所有的痛苦是我因罪过应受惩罚的一部分。\par

\newline{}2. 首先，我心灵悲伤，因为真理在人类之子中减少了；其次，我为自己担忧，唯恐在我其他罪过之上，我还变成一个厌世者，不再相信任何人心中存有真理和荣誉，如果我最为信任的那些人竟被证明对我及对真理持有如此态度的话。那么，我的弟兄，以及每一位真理的朋友，请确信：那篇作品不是我的；我不赞同它，因为它不是按我的观点写成的。即使我在许多年前曾写信给\Person{阿波里纳}或其他任何人，我也不应受到责备。如果某个团体的任何成员被切断而陷入异端（你们完全知道我说的是谁，尽管我没有指名道姓），我本人并不挑剔，因为每个人都会因自己的罪恶而死。\par

\newline{}这就是我回复那份寄给我的文件的答复，好使你们知道真理，并让所有不愿存着不义而压制真理的人都明了。如果需要我在每一条上更详细地为自己辩护，在天主的助佑下，我必这样做。我，\Person{奥林比乌}弟兄，既不主张三个天主，也不与\Person{阿波里纳}共融。\par

\SourceRecord{Translated by Blomfield Jackson.；Nicene and Post-Nicene Fathers, Second Series；Vol. 8.；Edited by Philip Schaff and Henry Wace.；Buffalo, NY: Christian Literature Publishing Co.,；1895.；<http://www.newadvent.org/fathers/3202131.htm>.}

% structure: p0001 source=h1 target=section reason=page-title

\section{第一三二封信}

\Emph{\Person{巴西略} \Place{凯撒勒雅}}\par

\Emph{致\Person{阿布拉米乌斯}，\Place{巴特奈}主教}\par

自从秋天以来，我对尊驾的行踪一直一无所知；因为我不断听到不确定的传闻，有人说您在\Place{萨摩撒塔}停留，有人说您在乡下，而另一些人则声称曾在\Place{巴特奈}见过您。这就是我不常写信的原因。现在，听说您下榻于\Place{安提约基雅}，在可敬的\Person{撒杜尔尼诺}伯爵府上，我很高兴把这封信交给我们所敬爱和尊敬的弟兄\Person{桑克提西摩斯}，我们的同作司铎，由他问候您，并劝勉您，无论身在何处，首先要纪念天主，其次要纪念我——您从一开始就决心爱戴并以之作为最亲密朋友之一的人。\par

\SourceRecord{Translated by Blomfield Jackson.；Nicene and Post-Nicene Fathers, Second Series；Vol. 8.；Edited by Philip Schaff and Henry Wace.；Buffalo, NY: Christian Literature Publishing Co.,；1895.；<http://www.newadvent.org/fathers/3202132.htm>.}

% structure: p0001 source=h1 target=section reason=page-title

\section{第一三三号书信}

\Person{圣巴西略·凯撒勒雅}\par

致\Place{亚历山大}主教\Person{伯多禄}。\par

肉眼所见带来肉体的友谊，长久的相伴使之加固，但真诚的关切是圣神的恩赐，祂将远隔重洋者联合，使朋友们彼此相识，不凭肉体特质，而藉灵魂的特质。主的恩宠赐予我这恩惠，允许我用灵魂之眼看见你，以真诚的挚爱拥抱你，仿佛被拉近到你身边，在信德的共融中与你紧密合一。我相信，你作为如此伟大人物的门徒，且长久与他为伴，必会以同样的精神行走，遵循同样真实宗教的教义。因此，我致书阁下，恳求你在其他方面继承那位伟人之外，也在对我的爱上继承他：你会经常写信告诉我你的消息，并以同样的关爱和热忱照顾普世的弟兄团体，正如那位最蒙福的人始终向一切真诚爱慕天主者所显示的那样。\par

\SourceRecord{Translated by Blomfield Jackson.；Nicene and Post-Nicene Fathers, Second Series；Vol. 8.；Edited by Philip Schaff and Henry Wace.；Buffalo, NY: Christian Literature Publishing Co.,；1895.；<http://www.newadvent.org/fathers/3202133.htm>.}

% structure: p0001 source=h1 target=section reason=page-title

\section{第一三四封书信}

\Emph{\Person{圣巴西略·凯撒勒雅}}\par

\Emph{致司铎\Person{佩奥尼乌斯}}。\par

从你信函的内容，你可猜出你给我的信带给我多少喜悦。你信中所表露的，清楚地显示出那颗纯洁的心，正是这心孕育了那些言辞。一涓细流，让人知晓其源头；同样，言谈的方式也揭示了它的来源。我不得不承认，一件非常奇特而又不大可能的事发生在我身上。因为我一直深深地渴望收到阁下的信函，可当我拿着你的信读完后，我并未因你所写的内容而高兴，反而因计算你在长期缄默中给我造成的损失而烦恼。既然你已开始写信，请不要停止。你将给我比给吝啬鬼送去大量金钱更大的快乐。我身边没有执笔之人，既无缮写员，亦无速记者。在我偶然雇佣的那些人中，有的已回归他们原有的生活方式，有的因长期疾病而不适于工作。\par

\SourceRecord{Translated by Blomfield Jackson.；Nicene and Post-Nicene Fathers, Second Series；Vol. 8.；Edited by Philip Schaff and Henry Wace.；Buffalo, NY: Christian Literature Publishing Co.,；1895.；<http://www.newadvent.org/fathers/3202134.htm>.}

% structure: p0001 source=h1 target=section reason=page-title

\section{第一三五号书信}

\Emph{\Person{圣巴西略}}\par

\Emph{致安提约基雅司铎狄奥多若}。\par

1. 我已读毕贤明所赠之书。对第二卷，我甚为喜悦，不仅因其篇幅简短——正如一位身体欠佳、偏好懒散的读者所愿——更因其思想丰富，且编排得当，使反对者的异议与相应的答覆清晰可辨。其质朴自然的文风，在我看来，恰合一位基督徒著述者的身份，他写作不为沽名钓誉，而为大众福祉。前一部著作虽实质相同，但修辞华丽，辞藻丰富，比喻繁多，对话精巧，似乎需要读者花费大量时间与心力方能领会其意并铭记于心。其中穿插的对反对者的抨击与对本方的支持，虽看似给论文增添了辩证的魅力，却因打断与延误而破坏了思想的连贯性，削弱了论证的力量。我深知，你的才智完全明白，撰写对话的异教哲学家，如亚里士多德与泰奥弗拉斯托斯，皆直指要点，因为他们自知不具柏拉图的文采。反之，柏拉图凭借其卓越的写作能力，同时攻击各种意见，并偶尔嘲弄其笔下人物，时而抨击特拉西马库斯的鲁莽与轻率，希庇阿斯的轻浮与浅薄，以及普罗泰戈拉的傲慢与浮夸。然而，当他在对话中引入无显著特征的人物时，他利用对话者来阐明观点，却不让任何属于人物特性的东西渗入论证。\WorkTitle{法律篇}便是一例。\par

2. 我们写作，既非出于虚妄的野心，而是有意将健全教义的忠告遗赠给弟兄们，那么，若引入某个举世皆知其言行冒昧的人物，以便按该人物的特质在论证中穿插某些点，这固然是好的——除非我们根本无权放下手头的工作去指责他人。但若对话主题宽泛而普遍，对人身的离题话便会打断其连贯性，且无益于任何善果。我写这些，是为了证明你并未将你的著作送给一个阿谀奉承者，而是与一位真正的弟兄分享了你的辛劳。我此番话并非为了纠正已完成的著作，而是为了预防将来；因为一个惯于写作、且勤于写作的人，必定不会犹豫继续写作；更何况，给他提供题材的人并未减少。于我而言，阅读你的著作已足够。我几乎要说出这样的话：正如我远离健康或从工作中得不到丝毫闲暇一样，我也同样远离写作的能力。然而，我现在已将两卷书中较大且较早的一卷退还，并尽可能读完了它。第二卷我保留了下来，希望抄录一份，但至今未能找到一位速记员。卡帕多细亚人令人羡慕的状况竟已沦落到如此贫困的地步！\par

\SourceRecord{Translated by Blomfield Jackson.；Nicene and Post-Nicene Fathers, Second Series；Vol. 8.；Edited by Philip Schaff and Henry Wace.；Buffalo, NY: Christian Literature Publishing Co.,；1895.；<http://www.newadvent.org/fathers/3202135.htm>.}

% structure: p0001 source=h1 target=section reason=page-title

\section{第一三六封信}

\Emph{\Person{凯撒勒雅的圣巴西略}}\par

\Emph{致\Place{萨莫萨塔}主教\Person{欧瑟伯}}。\par

善良的\Person{伊撒克}发现我处于何种状态，他本人会向您解释得最好；尽管他的言辞不足以悲剧性地描述我的痛苦，因为我病得如此之重。然而，任何稍微了解我的人都能猜想得到。因为，若在我被视为健康之时，我甚至比那些已被放弃的人更虚弱，您可想象我病重时会是什么样子。然而，既然疾病是我的自然状态，那么（让热病开个玩笑吧）在这种习性的变化中，我的健康反而变得特别旺盛。但这是主的鞭笞，它根据我的罪债不断增加我的痛苦；因此我遭受接二连三的病痛，以至于现在连孩子也能看出我这副躯壳必然要崩溃，除非天主的仁慈由于祂的宽容，赐予我悔改的时间，并且如今，像往常一样，将我从人力无法医治的灾祸中解救出来。愿这一切按祂所喜悦的和对我的益处而发生。\par

我几乎不必告诉您教会的处境是多么可悲和无望。如今，为了自身的安全，我们忽略了邻人的安危，甚至似乎看不到普遍的灾难牵连着个人的毁灭。尤其不必向您——像您这样的人——说这些，您从远处预见了未来，曾预言并宣告，并且率先觉醒，也唤醒他人，写信、亲自来访，没有留下任何未做的事、任何未说的话。我每一件事都记得，然而我们并无起色。的确，若不是我的罪过从中作梗（首先，我亲爱的弟兄、可敬的执事\Person{欧斯塔修斯}病重，整整拖延了我两个月，我日复一日地盼望他恢复健康；然后我周围的人全都病了；\Person{伊撒克}兄弟会告诉您其余的事；最后，我自己也被这疾病侵袭），我早就该去拜见您的尊驾了，并非想借此改善普遍的局势，而是为了从与您的交往中获得益处。我曾下定决心要避开教会的炮火，因为我毫无准备去面对敌人的攻击。愿天主的全能之手保全您，为我们众人做信德的高贵护卫者、教会的警觉斗士；并在我死前恩准我见到您，以安慰我的灵魂。\par

\SourceRecord{Translated by Blomfield Jackson.；Nicene and Post-Nicene Fathers, Second Series；Vol. 8.；Edited by Philip Schaff and Henry Wace.；Buffalo, NY: Christian Literature Publishing Co.,；1895.；<http://www.newadvent.org/fathers/3202136.htm>.}

% structure: p0001 source=h1 target=section reason=page-title

\section{第一三七号书信}

\Emph{凯撒勒雅的圣巴西略}\par

\Emph{致\Person{安提帕特}，论其就任\Place{卡帕多西亚}总督}\par

我现在真正感到因病所受的损失；因为，当这样一位人物就任我国总督时，我因必须调养身体而被迫缺席。整整一个月，我都在接受天然温泉的治疗，希望从中获得一些益处。但我似乎在孤寂中徒劳地自寻烦恼，或者确实应当成为世人嘲笑的对象，因为我忽略了那句谚语：\Quote{温暖对死者无益。}即便如此处境，我仍急切地想放下一切事务，亲自前往阁下那里，以便享受您所有高贵品质所带来的好处，并藉您的善良，妥善处理我在此处的所有家事。我们可敬的\Person{帕拉迪亚}母亲的房屋是我的，因为我不仅是她的近亲，还因她的品格而视她为母亲。如今，关于她的房屋引起了一些纷扰，我请求阁下暂时推迟调查，等待我到来；这绝不是为了使正义不得伸张——因为即使死一万次，我也不愿向一位热爱法律与正义的法官请求那样的恩惠——而是为了让您从我口中得知一些我不便书面陈述的事情。您若这样做，将绝不有违真理，而我们也不会遭受任何损害。因此，我恳请您将有关当事人安全拘押在军队监管下，并且不要拒绝赐予我这无害的恩惠。\par

\SourceRecord{Translated by Blomfield Jackson.；Nicene and Post-Nicene Fathers, Second Series；Vol. 8.；Edited by Philip Schaff and Henry Wace.；Buffalo, NY: Christian Literature Publishing Co.,；1895.；<http://www.newadvent.org/fathers/3202137.htm>.}

% structure: p0001 source=h1 target=section reason=page-title

\section{\WorkTitle{第一三八封信}}

\Emph{\Person{巴西略·凯撒勒雅}}\par

\Emph{致\Person{欧瑟伯}，\Place{萨摩撒塔}主教}\par

1. 当你虔诚的信件送达时，你想我当时的心境如何？想到它字里行间所表达的情感，我恨不得立即飞往\Place{叙利亚}；但想到那束缚我的身体病痛，我自觉不仅无力飞行，连在床上翻身都难。今天，我们亲爱的可敬弟兄兼执事\Person{厄耳彼狄}抵达，是我生病的第五十天。高烧让我极度虚弱。由于缺乏可吞噬之物，它像一盏油尽灯枯的灯，在这干枯的肉体上徘徊，由此引发了消耗性的顽疾。接着，我的旧疾——肝脏问题——又发作，令我无法进食，妨碍睡眠，将我置于生死边缘，仅赐予足以感受其折磨的生命。因此，我不得不求助于温泉，并借助医师的帮助。\par

但这一切都敌不过这祸患。或许别人能忍受，但既然它不期而至，无人能坚强到承受它。我长期受其困扰，但从未像现在这样，因它阻止我与你相见、享受你真挚友谊而痛苦。我知道自己被剥夺了何等喜乐，尽管去年我曾指尖触到你教会那甘甜的蜂蜜。\par

2. 出于许多紧迫的理由，我觉得必须与你的尊驾会面，既为与你商议诸多事宜，也为向你学习许多。在这里，连真诚的爱也难以找到。即使能找到一位真正的朋友，也没有人能像你那样，以在教会事务中多年操劳所获得的智慧和经验，在当前紧急情况下给我建言。其余的我不能写。然而，以下内容我可以放心地说。安提阿人彭佩雅诺之子、司铎\Person{厄瓦格略}，不久前随福者\Person{欧瑟伯}前往西方，现已从\Place{罗马}返回。他要求我按照西方人指定的精确措辞写一封信。他把我自己的信带回来给我，并报告说当地更精确的权威人士对此不满意。他还要求尽快派遣一位有名望的人组成代表团，以便他们有合理的借口来拜访我。我在\Place{塞巴斯特}的同情者已揭开\Person{欧达提}信仰不正的隐痛，并要求我给予教会关怀。\par

\Place{依科尼雍}是\Place{彼西底}的一座城市，古时仅次于最大城市，现为部分地区的首府，由不同部分组成，获准统治一个独立的行省。\Place{依科尼雍}也召唤我去访她，并为她选立一位主教，因为\Person{福斯提诺}已去世。我是否应避免边界外的祝圣；我应如何回答塞巴斯特人；我应对\Person{厄瓦格略}的提议采取何种态度——所有这些我都渴望通过与你面谈得到答案，因为现在我在软弱中与一切隔绝。如果你能找到近期往这边来的人，请赐予我你对所有这些问题的答复。若不能，请祈祷主所喜悦的事进入我心中。在你的主教会议中，也请提及我，并请你自己为我祈祷，并让你的子民与你一同祈祷，使我在此世寄居的剩余日子或时辰里，能得以以悦乐主的方式继续我的服务。\par

\SourceRecord{Translated by Blomfield Jackson.；Nicene and Post-Nicene Fathers, Second Series；Vol. 8.；Edited by Philip Schaff and Henry Wace.；Buffalo, NY: Christian Literature Publishing Co.,；1895.；<http://www.newadvent.org/fathers/3202138.htm>.}

% structure: p0001 source=h1 target=section reason=page-title

\section{第一三九函}

\Emph{\Person{圣巴西略}（\Place{凯撒勒雅}）}\par

\Emph{致\Place{亚历山大里亚}人}。\par

1. 我已经听说在亚历山大里亚和埃及其他地区的迫害，正如所料，我深感震动。我观察到魔鬼作战方式的巧妙。当他看到教会在敌人的迫害下反而增长，更加兴旺时，他就改变了策略。他不再进行公开的战争，而是为我们设下隐秘的陷阱，用他们（迫害者）所承载的名字来掩盖敌意，以至于我们既像我们的父辈一样受苦，同时又似乎不是为基督受苦，因为我们的迫害者也带着基督徒的名字。怀着这些思绪，我们久久静坐，对发生的事感到茫然，因为说实话，在听到你们迫害者的无耻和非人的异端时，我们耳朵都在发痒。他们不尊重年龄，不尊重对社会服务，也不尊重人民的爱戴。他们施加酷刑、侮辱和流放；他们掠夺所能找到的一切财产；他们既不在乎人的谴责，也不在乎公义审判者手中那可怕的未来报应。这一切使我惊愕，几乎使我失去理智。我的思考中还加上了这个念头：难道主已经完全抛弃了他的教会吗？末时到了吗？那\Quote{背叛}就这样临到我们，使那\BibleQuote{无法无天者，即那敌对者，且高举自己在一切称为神或受人敬拜者以上的，那丧亡之子，得以显露出来}\BibleRef{得后 2:4}吗？但如果这考验是暂时的，那么你们这些基督的高贵战士，要忍受。如果这世界将被交付于完全的、最终的毁灭，让我们不要灰心，而要等待从天上来的启示，和我们的伟大天主和救主耶稣基督的显现。如果所有受造物都要消解，这个世界的形态也要改变，为什么我们作为受造物的一部分，会为承受普遍的灾祸而惊讶，并被交付给我们公义的天主按我们力量的程度所加给我们的苦难，他\BibleQuote{不让我们受超过我们能力的试探，但在试探中给我们开一条出路，使我们能忍受得住}\BibleRef{格前 10:13}呢？弟兄们，殉道者的冠冕在等着你们。宣认者的队伍已准备好伸手向你们，欢迎你们进入他们的行列。请记住，古代圣人中，没有一个人是靠奢华生活、受人奉承而赢得忍耐的冠冕的；他们都是经过大苦难的火炼而受到考验的。\Quote{因为有些人受了戏弄和鞭打的酷刑，有些人被锯断，被刀杀。}这些是圣人的荣耀。能够为基督受苦的人是有福的；受苦更大的人更有福，因为\BibleQuote{现在的苦楚微不足道，无法与将来要显示给我们的光荣相比}\BibleRef{罗 8:18}。\par

2. 如果我能够前往你们那里，我最高兴的莫过于与你们相会，好让我看到并拥抱基督的运动员，分享你们的祈祷和神恩。但如今我的身体因长病而衰弱，几乎离不开床，又有许多像饿狼一样的人埋伏着，等待时机撕裂基督的羊群。因此我不得不以书信来探访你们；我首先极力劝你们为我祈祷，使我在余下的日子或时辰里，能够按天主国的福音侍奉主。其次，我请求你们原谅我的缺席和迟未写信。我好不容易才找到一个能完成我心意的人。我指的是我的儿子、隐修士\Person{欧仁}，我藉着他恳求你们为我并为整个教会祈祷，并回信告诉我你们的消息，好使我听到后更加愉快。\par

\SourceRecord{Translated by Blomfield Jackson.；Nicene and Post-Nicene Fathers, Second Series；Vol. 8.；Edited by Philip Schaff and Henry Wace.；Buffalo, NY: Christian Literature Publishing Co.,；1895.；<http://www.newadvent.org/fathers/3202139.htm>.}

% structure: p0001 source=h1 target=section reason=page-title

\section{\WorkTitle{第一四〇号信函}}

\Signature{圣巴西略（凯撒勒雅的）}\par

\Emph{致安提约基雅教会}。\par

1. \Quote{但愿我有鸽子般的翅膀，好能飞去}到你们那里，满足我渴望与你们相见的愿望。但是现在，我所欠缺的不仅是翅膀，而是整个身体，因为长期患病，如今已因不断的苦难而衰朽。因为任何人，无论心如铁石，全无同情与仁爱，当他听到从四面八方传来，犹如悲伤的合唱队齐奏哀歌的叹息时，怎能不心中忧伤，俯首至地，被这些无法弥补的烦恼消磨殆尽？但是圣洁的天主有能力为无法弥补之事提供医治，并赐予你们久劳之后的喘息。我希望你们感受到这份安慰，并在喜乐中盼望安慰，忍受当前苦难的痛苦。我们是在为自己的罪孽偿付代价吗？那么，我们的灾祸将使我们今后免受天主的忿怒。我们是否因这些试探而为真理争战？那么，公义的赏报者不会让我们受超过所能承受的试探，而是以我们先前的奋斗为回报，赐予我们忍耐并盼望祂的冠冕。因此，让我们不要因惧怕为真理打那美好的仗而退缩，也不要因绝望而抛弃我们已经完成的劳苦。因为灵魂的力量不表现于一次勇敢的行为，也不表现于短时间的努力；但那试验我们心灵的主，愿意我们在长久持续的考验之后获得正义的冠冕。只要我们的精神保持不折，我们在基督内的信德坚定不移，我们的护卫者就要显现；祂必来临，决不迟延。期待患难一个接一个，希望上加希望；还有一点点时候，还有一点点时候。圣神知道如何以将来的应许安慰祂的儿女。患难之后是希望，我们所盼望的并不遥远，因为即使有人列举人的一生，与我们所盼望的无穷岁月相比，也不过是短暂的间隙。\par

2. 我不接受别人为我撰写的新信经，也不敢凭自己的理智提出结论，以免使真宗教的言语沦为凡人之言；而是将圣教父们所教导的，向所有询问我的人宣布。在我的教会中，尼西亚圣教父在会议中撰写的信经被使用。我相信在你们那里也念诵同一信经；但我不拒绝在信中抄写其确切条文，以免被指责为懒惰。信经如下：这就是我们的信仰。但关于圣神，未作任何定义，因为那时圣神对抗者尚未出现。因此，没有提到需要诅咒那些说圣神是被造和受造本性的人。因为在天主圣三中，没有什么是受造的。\par

\SourceRecord{Translated by Blomfield Jackson.；Nicene and Post-Nicene Fathers, Second Series；Vol. 8.；Edited by Philip Schaff and Henry Wace.；Buffalo, NY: Christian Literature Publishing Co.,；1895.；<http://www.newadvent.org/fathers/3202140.htm>.}

% structure: p0001 source=h1 target=section reason=page-title

\section{第一四一章}

\Emph{\Person{凯撒勒雅的圣巴西略}}\par

\Emph{致\Person{欧瑟伯}，\Place{萨摩萨塔}主教。}\par

1. 我现在收到了两封来自您神圣而卓越的智慧的书信。信中明确告诉我，您圣座辖下的平信徒如何期待我，以及我因未能出席神圣主教会议而使他们失望。另一封信，从笔迹推断是较早的，虽然送达较晚，却给了我忠告，既令您尊贵，又对我必要：不要忽视天主教会的事务，也不要让事务的引导逐渐落入我们的对手手中，从而他们的利益必兴盛，我们的则受损。我想我两封都回复了。但是，由于我不确定那些受托送信的人是否保存了我的回信，我将再次为自己辩护。关于我的缺席，我可以提出一个无可挑剔的理由，我认为此消息必已传到您的圣座，即我一直被疾病所困，那病几乎使我到了死亡的门槛。即使现在写作，疾病的余孽仍在我身上。它们对另一个人来说可能是无法忍受的。\par

2. 关于并非因我的疏忽而使教会利益被出卖给对手这一事实，我希望阁下知晓：与我共融的主教们，或因缺乏热忱，或因怀疑我并不对我坦诚，又或因魔鬼常阻碍善工，他们不愿与我合作。从前，我们中的多数人确实彼此合一，包括杰出的博斯波里乌斯。但实际上，在最关键的事情上，他们并未给我帮助。这一切的结果是，我的康复在很大程度上因我的痛苦而受阻，我所感到的悲伤使我的最坏症状复发。然而，当教规——正如您自己所知——不允许此类事情由一人决定时，我独自一人且无人帮助又能做什么呢？然而，我尝试过什么补救措施呢？我未能提醒他们什么决定呢？有的通过信件，有的亲自。当他们听到我死亡的传闻时，甚至来到城里；因天主的旨意，他们发现我还活着，我便向他们作了合于时宜的讲话。在我面前，他们尊重我，承诺一切合宜之事，但一回去，他们就回到自己的意见。在这一切中，我像其他人一样受苦，因为主显然已离弃了我们，由于罪恶增多，我们的爱德冷淡了。愿您伟大而有力的代祷在天主前为我足够。也许我们会成为有用的人，或者即使我们达不到目标，也能逃避谴责。\par

\SourceRecord{Translated by Blomfield Jackson.；Nicene and Post-Nicene Fathers, Second Series；Vol. 8.；Edited by Philip Schaff and Henry Wace.；Buffalo, NY: Christian Literature Publishing Co.,；1895.；<http://www.newadvent.org/fathers/3202141.htm>.}

% structure: p0001 source=h1 target=section reason=page-title

\section{\WorkTitle{第一四二封书信}}

\Person{凯撒勒雅的圣巴西略}\par

\Emph{致总督会计}。\par

我召集了我所有担任\OriginalTerm{乡村主教}{\GreekText{χωρεπίσκοποι}}的弟兄，在真福殉道者欧普西奇的教务会议上，将他们引见给阁下。因您缺席，只得藉信函将他们呈于台前。故此，请明鉴此弟兄敬畏上主，堪受您的信赖。关于他为穷人向您的善意陈情之事，请视其为可靠之人，并尽己所能援助受苦者。我相信您会欣然照顾他辖区内的穷人医院，并全然豁免其税赋。您的同僚已恩准将穷人的微薄产业免于课税。\par

\SourceRecord{Translated by Blomfield Jackson.；Nicene and Post-Nicene Fathers, Second Series；Vol. 8.；Edited by Philip Schaff and Henry Wace.；Buffalo, NY: Christian Literature Publishing Co.,；1895.；<http://www.newadvent.org/fathers/3202142.htm>.}

% structure: p0001 source=h1 target=section reason=page-title

\section{\WorkTitle{第一四三封信}}

\Emph{\Person{圣巴西略·凯撒利亚}}\par

\Emph{致另一位会计师}。\par

我若能遇见尊驾，本应亲自向您呈上我所忧虑之事，并为受苦者辩护，但我因病和事务繁忙而受阻。因此，我派遣这位\OriginalTerm{村司铎}{\GreekText{χωρεπίσκοπος}}、我的兄弟代我前来，恳求您给予他援助，利用他并听取他的意见，因为他的诚实和明智使他适合提出此类建议。若您惠然视察他管理的穷人医院（我相信您必不会过而不入，因您在这方面经验丰富；我听说您用主所祝福的资财支持了\Place{阿马西亚}的一家医院），我相信，在看到之后，您会给予他所求的一切。您的同僚已答应为这些医院提供一些帮助。我告诉您这些，不是要您效仿他，因为您很可能在善工上成为他人的表率，而是要您知道，在这件事上其他人已对我表示关心。\par

\SourceRecord{Translated by Blomfield Jackson.；Nicene and Post-Nicene Fathers, Second Series；Vol. 8.；Edited by Philip Schaff and Henry Wace.；Buffalo, NY: Christian Literature Publishing Co.,；1895.；<http://www.newadvent.org/fathers/3202143.htm>.}

% structure: p0001 source=h1 target=section reason=page-title

\section{第一四四封信}

\Signature{凯撒勒雅的圣巴西略}\par

\Emph{致长官的官员}\par

你在城内遇见来人时便认识他。然而我写信将他托付给你，希望他在许多你所关心的事上能对你有帮助，因为他能给出虔诚而明智的建议。现在正是实现你私下对我说过的话之时；我的意思是，当这位兄弟向你说明穷人的状况时。\par

\SourceRecord{Translated by Blomfield Jackson.；Nicene and Post-Nicene Fathers, Second Series；Vol. 8.；Edited by Philip Schaff and Henry Wace.；Buffalo, NY: Christian Literature Publishing Co.,；1895.；<http://www.newadvent.org/fathers/3202144.htm>.}

% structure: p0001 source=h1 target=section reason=page-title

\section{\WorkTitle{第一四五封書信}}

\Person{聖巴西略·凱撒勒雅}\par

\Emph{致\Person{歐瑟伯}，\Place{薩摩薩塔}主教}\par

我知道你爲了天主的眾教會所受的無數勞苦；我知道你事務繁忙，你盡責並非視之爲次要之事，而是按天主的旨意行事。我知道那個人，他彷彿在緊緊圍攻你，使你像老鷹遮蔽下的鳥兒一樣，不敢遠離你的庇護。這一切我都知道。但是渴望是強烈的，既盼望不可能的事，又嘗試辦不到的事。不如說：對天主的\OriginalTerm{望德}{hope}是萬物中最強的。因爲我並非出於不合理的願望，而是出於\OriginalTerm{信德}{faith}的力量，期待甚至從最大的困難中也能找到出路，並期待你找到克服一切障礙的方法，來看你最心愛的\OriginalTerm{教會}{Church}，也讓她看見你。她所珍視的萬善之中，最渴望的就是見到你的面容，聽到你的聲音。所以，務必不要使她的希望落空。去年我從\Place{敘利亞}回來時，報告了你給我的承諾，你無法想象我讓她心中充滿了何等的希望。我的朋友，不要把你的來訪推遲到另一個時候。即使有一天你可能見到她，但你可能見不到我和她了，因爲疾病正催促我離開這痛苦的人生。\par

\SourceRecord{Translated by Blomfield Jackson.；Nicene and Post-Nicene Fathers, Second Series；Vol. 8.；Edited by Philip Schaff and Henry Wace.；Buffalo, NY: Christian Literature Publishing Co.,；1895.；<http://www.newadvent.org/fathers/3202145.htm>.}

% structure: p0001 source=h1 target=section reason=page-title

\section{第一四六封信}

\Emph{\Person{凯撒勒雅的圣巴西略}}\par

\Emph{致\Person{安提约古}}。\par

我不能指责你疏忽或不在意，因为当有写信的机会时，你什么也没说。因为我视你用你那可敬的手向我送来的问候，胜过许多书信。作为回报，我向你致意，并恳切地请求你留意你灵魂的救恩，以理性约束一切肉体的私欲，并且恒常将天主的思念建立在你的灵魂内，如同在一个至圣的圣殿中。在每一个行动和每一句话语中，都要将基督的审判摆在眼前，这样每一个个别的行为，由于被引用到那精确而可畏的审查中，就会在赏报的日子为你带来光荣，那时你将得到一切受造物的称赞。如果那位伟人能够来拜访我，我将很高兴在这里与你一同见到他。\par

\SourceRecord{Translated by Blomfield Jackson.；Nicene and Post-Nicene Fathers, Second Series；Vol. 8.；Edited by Philip Schaff and Henry Wace.；Buffalo, NY: Christian Literature Publishing Co.,；1895.；<http://www.newadvent.org/fathers/3202146.htm>.}

% structure: p0001 source=h1 target=section reason=page-title

\section{第一四七封信}

\Emph{\Place{凯撒勒雅}圣\Person{巴西略}}\par

\Emph{致\Person{阿布尔吉乌斯}}\par

直到如今，我读\Person{荷马}诗歌的第二部分（其中叙述\Person{尤利西斯}的历险）时，一向以为\Person{荷马}是虚构的。然而，降临在最杰出的\Person{马克西穆斯}身上的灾难，使我过去认为荒诞不经、难以置信的事变得极为可能。\Person{马克西穆斯}曾统治一个并非微不足道的民族，正如\Person{尤利西斯}是\Person{刻法勒尼亚人}的首领一样。\Person{尤利西斯}拥有巨大财富，却空手而归，一无所有。灾难已将\Person{马克西穆斯}逼到如此窘境，他可能不得不穿着借来的破衣回家。或许他遭受这一切，是因为他激怒了一些\Person{莱斯特律戈涅斯人}，又遇到了某个以女子之形隐藏着狗的凶残与狂暴的\Person{斯库拉}。自那以后，他几乎无法从这无法脱身的漩涡中游出。他通过我恳求你，本着仁爱，为他不应遭受的不幸而悲伤，不要对他的需要保持沉默，而是将其告知当权者。他期望因此能得到一些帮助，以对抗针对他的诽谤；如若不然，至少也要公开那位对他表现出如此敌对狂热的敌人的意图。当一个人受了冤屈，能使敌人的邪恶公之于众，对他来说便是极大的安慰。\par

\SourceRecord{Translated by Blomfield Jackson.；Nicene and Post-Nicene Fathers, Second Series；Vol. 8.；Edited by Philip Schaff and Henry Wace.；Buffalo, NY: Christian Literature Publishing Co.,；1895.；<http://www.newadvent.org/fathers/3202147.htm>.}

% structure: p0001 source=h1 target=section reason=page-title

\section{\WorkTitle{第一四八封信}}

\Emph{\Person{圣巴西略}\Place{凯撒勒雅}}\par

\Emph{致\Person{图拉真}}。\par

即便是哀叹自身灾祸的能力，也能给受苦者带来极大的安慰；尤其当他们遇到那些因崇高品格而能同情其悲伤的人时，更是如此。因此，我尊贵的弟兄\Person{玛克西穆斯}，在担任我国总督之后，遭受了前无古人的苦楚；他被剥夺了所有祖传和自积的财产，身体因在世界各地流浪而备受各种折磨，甚至无法保全其公民身份——为了保持这一身份，自由人通常不辞辛劳——他向我诉说了所遭遇的一切，并恳求我向您简要描述他所陷入的苦难的\WorkTitle{伊利亚特}。而我，既然无法以其他方式减轻他的困苦，便欣然应允，将我所听到的一部分简要地向阁下您叙述。他对我而言，似乎因要直述自己的灾祸而羞愧。如果所发生的事表明施加不义者是恶人，那么至少也证明受难者值得极大的怜悯；因为陷入由天主眷顾所施加的患难这一事实，似乎在一定程度上表明此人已被指定为受苦者。但如果您能仁慈地垂顾他，并将那所有领受者都无法穷尽的丰盛恩惠——我指的是您的仁慈——也施予他，那对他来说就是足够的安慰了。我们都确信，在法庭上您的庇护将是迈向胜利的一大步。那恳求我的信函以求帮助的人，是所有人中最正直的。但愿我能看到他，与众人一起，竭尽全力高声歌颂您阁下的赞美。\par

\SourceRecord{Translated by Blomfield Jackson.；Nicene and Post-Nicene Fathers, Second Series；Vol. 8.；Edited by Philip Schaff and Henry Wace.；Buffalo, NY: Christian Literature Publishing Co.,；1895.；<http://www.newadvent.org/fathers/3202148.htm>.}

% structure: p0001 source=h1 target=section reason=page-title

\section{\WorkTitle{第一四九封信}}

\Signature{\Person{圣巴西略}·\Place{凯撒勒雅}}\par

\Emph{致\Person{图拉真}}\par

您亲眼目睹了\Person{玛克西穆}的悲惨境况——他曾是德高望重之人，如今却最堪怜悯，曾是我故国的总督。但愿他从未如此！我想，许多人很可能会回避省督的职位，如果他们的尊荣竟以这样的结局告终。那么，对于一个凭借其敏捷才智、能从少许情况推知其余的人来说，我几乎无需详述我亲眼所见、亲耳所闻的一切。然而，或许我提及以下情况并不显得多余：在您到来之前，虽然有许多可怕的事曾胆大妄为地加于他，但后来所发生的事，竟使先前的行为被视为仁慈；当权者此后对他采取的行动，其暴行、伤害乃至人身虐待，竟至于此。如今他由一名护卫押送到此，若非您愿意伸出援助之手保护受难者，他的恶行便将满盈。在敦促您的善心行善时，我感到自己在做一件多余的事。然而，既然我渴望为\Person{玛克西穆}效劳，我恳求阁下为我的缘故，在您对正义的热诚之上再加一点援手，好让他清楚看到，我为他所作的干预对他是有益的。\par

\SourceRecord{Translated by Blomfield Jackson.；Nicene and Post-Nicene Fathers, Second Series；Vol. 8.；Edited by Philip Schaff and Henry Wace.；Buffalo, NY: Christian Literature Publishing Co.,；1895.；<http://www.newadvent.org/fathers/3202149.htm>.}

% structure: p0001 source=h1 target=section reason=page-title

\section{第一百五十封信}

\Signature{\Person{圣巴西略（凯撒勒雅的）}}\par

\Emph{致盎菲罗削，以赫辣克利达斯之名。}\par

1. 我记得我们往日的交谈，对你我所说的话都没有遗忘。如今世俗生活已无法束缚我。因为尽管我内心如常，尚未脱去旧人，但在外表上，藉着远远远离世俗生活，我似乎已经开始踏上基督徒生活的道路。我独坐一旁，如同即将航行深海的人，眺望眼前的一切。的确，水手需要顺风才能航行顺利；而我却需要一位引导者，牵着我的手，安全地带领我穿越人生的苦浪。我感到首先需要一道缰绳来约束我的青年气盛，然后需要刺棒驱使我走向虔敬的道路。这两者似乎都由理性提供，它时而管束我灵魂的桀骜不驯，时而激励我的怠惰。我又需要其他的疗方，好洗去习惯的污秽。你知道，我长久习惯于论坛，言语铺张，不能防范恶者注入我心中的思想。我也是荣誉的奴仆，难以轻易放弃自视过高。针对这一切，我感到需要一位伟大的导师。此外，我还断定，净化灵魂的眼睛，使其从一切无知的黑暗（如同从某种盲目的体液）中解脱后，能凝视天主荣耀的美，这并非小事，也不仅是一时之益。所有这些，我知道你的智慧都明了；我知道你希望我能得到某人给予这样的帮助，如果天主曾恩赐我与你相遇，我相信我会更清楚自己应留心什么。因为如今，由于极大的无知，我甚至难以判断自己缺乏什么。然而我并不后悔最初的冲动；我的灵魂并未从虔敬生活的目标退缩，正如你为我所担心，你高尚而合宜地尽己所能，免得我像我听说的那个女人一样，转身变成一根盐柱。然而，我仍受外在权威的约束；因为行政长官像追捕逃兵一样追捕我。但影响我最大的还是自己的心，它向我证实了我向你述说的一切。\par

2. 既然你提到我们的盟约，并宣布要起诉，你使我在沮丧中发笑，因为你仍是一位辩护律师，不放弃你的精明。我认为，除非我像无知的人一样完全误解真理，否则通往主的道路只有一条，所有走向祂的人都是一同旅行，并按照同一生活的盟约行走。如果是这样，无论我走到哪里，我怎能与你分离？我怎能停止与你共同生活，并一同事奉我们俩都投奔的天主？我们的身体可能因距离而分离，但天主的眼目无疑仍注视着我们两人；如果像我这样的生活确实配得上神圣眼目的眷顾的话；因为我在圣咏某处读到：主的眼目注视义人。我确实祈祷，与你以及所有与你同心的人，我能在身体上联合，并在日日夜夜与你及任何其他真实敬拜天主的人一同向我们在天之父屈膝；因为我知道在祈祷中团契能带来巨大益处。如果，每当我偶然躺在不同的角落呻吟时，我总是被指控说谎，我无法与你的论点争辩，并且已经判定自己为说谎者，如果由于我的疏忽，我说了什么让我受此指控的话。\par

3. 我最近在凯撒勒雅，为了解那里的情况。我不愿留在城内，便去了附近的医院，以便在那里得到我所需的信息。按他的习惯，那位十分虔诚的主教访问了那里，我向他请教了你敦促我的问题。我不可能记住他回答的一切；它远超过一封信的篇幅。然而，简而言之，他关于贫穷的论述如下：规则应该是每人只应限制自己的财产为一件衣服。为此他引用了若翰洗者的话：\BibleQuote{有两件衣裳的，就当分给那没有的；} \BibleRef{路 3:11} 另一个证据是主禁止门徒有两件衣裳。 \BibleRef{玛 10:10} 他又补充说：\BibleQuote{你若愿意是成全的，去！变卖你所有的，施舍给穷人。} \BibleRef{玛 19:21} 他还说珍珠的比喻也与此相关，因为那商人寻得贵重珍珠后，就去变卖一切所有，买了它；他还补充说，任何人甚至不应允许自己分配自己的财产，而应将其交给负责管理穷人事务的人手中；他从宗徒大事录中证明了这一点，因为那些人变卖财产，拿来放在宗徒们脚前，由他们照各人所需分给每人。 \BibleRef{宗 4:35} 因为他说，需要有经验来区分真实需要和贪婪乞讨。因为凡施与受苦者的，就是施与主，且必从主那里得到赏报；但施与每个流浪者的，是扔给狗，这狗虽因纠缠而令人讨厌，却不因缺乏而值得怜悯。\par

4. 此外，他首先简短地谈论了一些日常生活的职责，这正符合主题的重要性。关于这一点，我希望你亲耳听他讲述，因为由我来削弱他教导的力量是不合适的。我祈祷我们能一起拜访他，这样你既能准确地记住他所说的话，又能凭自己的才智补充任何遗漏。在我听到的许多话中，我记得有一句是：教导基督徒生活的方式，与其说在于言语，不如说在于每日的榜样。我知道，若不是因为你照料年迈父亲的职责所阻，没有什么比与这位主教见面更让你珍视，你也不会建议我离开他去荒野中漫游。洞穴和岩石随时为我们准备着，但从同伴那里得到的帮助却并非总是唾手可得。因此，如果你愿意接受我的建议，请让你父亲明白，允许你暂时离开他，去与这样一位人物会面是值得的——这位人物无论从自身经验还是从个人智慧来说，都所知甚多，并能将其传授给所有接近他的人。\par

\SourceRecord{Translated by Blomfield Jackson.；Nicene and Post-Nicene Fathers, Second Series；Vol. 8.；Edited by Philip Schaff and Henry Wace.；Buffalo, NY: Christian Literature Publishing Co.,；1895.；<http://www.newadvent.org/fathers/3202150.htm>.}

% structure: p0001 source=h1 target=section reason=page-title

\section{第一百五十一书}

\Signature{圣巴西略·凯撒勒雅}\par

\Emph{致医生厄斯塔丢}\par

如果我的书信有任何益处，请速写信给我，并催我写信。当我们阅读爱主的智士的书信时，我们无疑会更愉快。应当由你来判断，你是否觉得我写的东西有任何值得注意之处。要不是事务繁多，我不会剥夺自己频繁通信的乐趣。恳请你，负担较轻，用书信安慰我。俗话说，水井越用越清。你从你的职业中得到的劝诫显然不切题，因为下刀的并不是我，而是那些日子已尽的人，他们正在自残。斯多葛派的格言说：\Quote{既然事情不按我们所愿发生，我们就喜欢已发生的事；}但我无法使自己的心顺应正在发生的事。有些人因为不得已而做自己不喜欢的事，我并无异议。你们医生并不因为喜欢才给病人灼烙或使其受其他痛苦，而是出于病情的需要常采用这种疗法。水手们并不情愿把货物抛到海里，但为了免于船难，他们忍受损失，宁过贫苦的生活也不愿死。我确实以忧伤和许多叹息看待那些疏远者的分离。但我还是忍受了。对于真理的热爱者来说，没有任何事物能置于天主和对他所怀的希望之前。\par

\SourceRecord{Translated by Blomfield Jackson.；Nicene and Post-Nicene Fathers, Second Series；Vol. 8.；Edited by Philip Schaff and Henry Wace.；Buffalo, NY: Christian Literature Publishing Co.,；1895.；<http://www.newadvent.org/fathers/3202151.htm>.}

% structure: p0001 source=h1 target=section reason=page-title

\section{第一五二号书信}

\Emph{圣巴西略（凯撒勒雅）}\par

\Emph{致\Person{维克多}，统领。}\par

若我对其他任何人疏于写信，或许有理由被指责为疏忽或遗忘。但忘记您是不可能的，因为您的名字常挂在众人嘴边。然而，我不能对一位或许在帝国中比任何人都更杰出的人粗心大意。我沉默的原因显而易见：我害怕打扰如此伟大的人物。但是，如果您在一切其他美德之外，还加上不但接收我所寄出的信函，而且主动询问缺失的信件，那么，看！我现在正怀着喜悦的心情给您写信，并且将来会继续写信，祈求天主报答您对我的尊重。因为教会，您已先于我的恳求，做了我本应请求的一切。您行事并非为取悦人，而是为取悦天主，祂曾荣耀了您；祂在此生赐给您一些美物，在来世也将赐给您其他美物，因为您以真理行走在祂的道路上，并且从始至终，您的内心一直固守于正确的信德。\par

\SourceRecord{Translated by Blomfield Jackson.；Nicene and Post-Nicene Fathers, Second Series；Vol. 8.；Edited by Philip Schaff and Henry Wace.；Buffalo, NY: Christian Literature Publishing Co.,；1895.；<http://www.newadvent.org/fathers/3202152.htm>.}

% structure: p0001 source=h1 target=section reason=page-title

\section{\WorkTitle{第一五三封书信}}

\Emph{\Person{圣巴西略（凯撒勒雅）}}\par

\Emph{致前任执政官维克托}。\par

每次我读到阁下的信件，我便感谢天主，因为您继续记念我，并且没有被任何诽谤所动摇，减少您曾经愿意因您明智的判断或友善的交往而对我抱持的爱。我于是向圣洁的天主祈祷，愿您对我保持这份心意，并愿我配得上您所赐予的荣耀。\par

\SourceRecord{Translated by Blomfield Jackson.；Nicene and Post-Nicene Fathers, Second Series；Vol. 8.；Edited by Philip Schaff and Henry Wace.；Buffalo, NY: Christian Literature Publishing Co.,；1895.；<http://www.newadvent.org/fathers/3202153.htm>.}

% structure: p0001 source=h1 target=section reason=page-title

\section{第一五四封信}

\Person{凯撒勒雅的圣巴西略}\par

\Emph{致\Person{阿斯科利乌斯}，\Place{得撒洛尼}主教}。\par

您按着属神之爱的法律，首先写信给我，并以您的善表激发我采取同样的行动，实在是做得好。世俗的友谊确实需要实际的见面和往来，以便由此产生亲密。然而，凡知道以精神相爱的人，不需要肉体来促进感情，而是因着信德的共融被引向精神的相通。为此，感谢主，祂安慰了我的心，向我表明爱德并没有在所有人身上冷淡，而是世上仍有显出基督门徒证据的人。您那里的情况，似乎有点像夜晚的星辰，有的在天空的这一部分发光，有的在那一部分；它们的明亮是迷人的，而正因为出其不意，更加迷人。您正是如此，教会的明灯，最多只有少数，且在这昏暗的景况中屈指可数，好似在没有月光的夜里发光；您因德行而受欢迎，又因难得一见而更加令人渴望。您的来信使您的性情对我完全显明了。虽然就字数而言是简短的，但就其情感的纯正而言，完全足以向我证明您的心思和意向。您对真福的\Person{圣亚大纳削}事业的热情，清楚地证明您在最重要的事情上是健全的。为了回报我从您的信中获得的喜乐，我非常感谢我可敬的儿子\Person{欧菲米乌}，我祈求\OriginalTerm{圣者}{Holy One}赐予他一切帮助，并恳请您与我一同祈祷，以便我们很快能接他和他非常可敬的妻子、我在主里的女儿回来。至于您本人，我恳求您不要使我们的喜乐停留在开端，而是要在一切可能的机会上写信，并通过不断的通信增进您对我的好感。我恳求您，告诉我关于您的教会的情况，以及它们在合一方面的状况。请为我们这里祈祷，愿我们的主斥责风和海，使我们也能享受伟大的平静。\par

\SourceRecord{Translated by Blomfield Jackson.；Nicene and Post-Nicene Fathers, Second Series；Vol. 8.；Edited by Philip Schaff and Henry Wace.；Buffalo, NY: Christian Literature Publishing Co.,；1895.；<http://www.newadvent.org/fathers/3202154.htm>.}

% structure: p0001 source=h1 target=section reason=page-title

\section{\WorkTitle{第一五五号书信}}

\Emph{\Person{圣巴西略（凯撒勒雅）}}\par

\Emph{无收信人地址}。\Emph{关于一位训练者}。\par

我不知道该如何为自己辩护，以回应阁下惠赐的第一封也是唯一一封信中所包含的所有抱怨。并非因为我理亏，而是因为在众多指控中，难以选出最关键的，并确定我应从何处开始补救。或许，若我按照来信的顺序，便依次触及每一个。直至今日，我对那些前往\Place{斯基泰亚}的人一无所知；也没有人告诉我任何从您府上来的，以便我能通过他们向您问候，尽管我渴望抓住每一个问候阁下的机会。在祈祷中忘记您是不可能的，除非我先忘记天主召叫我的工作，因为借着天主的恩宠，您是忠实的，您记得教会的一切祈祷：我们也会为旅途中的弟兄祈祷，在圣教会中为从军者奉献祈祷，为那些为主之名说话的人祈祷，也为那些显示圣神果实的人祈祷。在大部分或所有这些中，我认为阁下也被包括在内。就我个人而言，我怎能忘记您呢？我有如此多的理由唤起回忆：这样一位姐妹，这样一些侄儿，这样一些亲属，如此善良，如此爱我，还有房屋、家人和朋友？所有这一切，即使违背我的意愿，也迫使我记起您的善意。然而，关于此事，我们的兄弟没有带来不愉快的消息，我也没有作出任何可能伤害他的决定。那么，请免除那位\OriginalTerm{乡村主教}{chorepiscopus}和我的一切责备，并为那些造谣中伤的人感到悲伤。如果我们博学的朋友想起诉我，有法庭和法律在。在此，我恳求您不要责备我。在您所做的一切善事中，您为自己积攒财宝；您在报应之日为自己预备了同样的安慰，就像您为主之名受迫害的人所提供的那样。若您送殉道者的遗骸回家，您会做得很好；正如您信中所写，那里的迫害甚至现在仍在为主产生殉道者。\par

\SourceRecord{Translated by Blomfield Jackson.；Nicene and Post-Nicene Fathers, Second Series；Vol. 8.；Edited by Philip Schaff and Henry Wace.；Buffalo, NY: Christian Literature Publishing Co.,；1895.；<http://www.newadvent.org/fathers/3202155.htm>.}

% structure: p0001 source=h1 target=section reason=page-title

\section{第一五六封信}

\Emph{凯撒勒雅的圣巴西略}\par

\Emph{致司铎埃瓦格里乌斯}。\par

1. 我向你保证，面对你来信的篇幅，我不仅没有不耐烦，反而因为阅读时的喜悦而觉得它甚至太短。因为有什么比和平的思想更令人愉快？有什么比采取行动实现和平更适于神圣的职务，更中悦主呢？愿你获得和平制造者的赏报，因为如此蒙福的职务是你良好愿望和热切努力的目标。同时，我所尊敬的朋友，请相信我，在热切希望和祈祷看到那一天——所有同心同德的人共聚一堂——方面，我不亚于任何人。事实上，对教会的分裂与不和感到愉悦，却不认为最大的益处在于将基督身体的肢体紧密连结，那将是荒谬的。但是，唉！我的能力正如我的愿望一样真实。没有人比你自己更清楚，时间才是医治时间所造成病患的唯一良药。此外，需要一种强烈且有力的治疗才能触及病症的根源。你会明白这个暗示，不过我没有理由不直说。\par

2. 自大，一旦通过习惯在心灵中扎根，就无法通过一个人、一封信或短时间消除。除非有一位所有各方都信任的仲裁者，否则猜疑和冲突永远不会完全停止。的确，如果天主的恩宠的影响倾注在我身上，并且我在言语、行为和神恩上得着能力来影响这些对立派别，那么这项大胆的尝试可能会被要求于我；不过，也许即使那样，你也不会建议我只靠自己、而没有主教（主要承担教会责任的人）的合作来尝试调整局面。但他不能来这里，而我在冬季也不容易进行长途旅行，或者更确切地说，我无论如何也不能，因为\Place{亚美尼亚}山脉很快将变得难以通行，即使是年轻力壮的人也是如此，更不用说我一直以来的身体疾病了。我不反对写信告诉他这一切；但我并不指望写信会带来什么结果，因为我了解他谨慎的性格，而且毕竟书面的文字很难说服人心。有那么多事情需要敦促、承受、回应和反对，而一封信是没有灵魂的，实际上不过是废纸。然而，正如我说过的，我会写信。只是，最虔诚亲爱的兄弟，请相信在这件事上我没有私人感情。感谢天主，我对任何人都没有这样的感情。我没有忙于调查针对这个或那个人所提出的或真或假的抱怨；因此，我的意见值得你关注，因为它来自一个真正不能出于偏袒或偏见而行事的人。我只希望，藉着主的善意，一切都能按照教会应有的秩序进行。\par

3. 我从我亲爱的儿子\Person{多罗泰}（我们事奉中的同伴）那里得知你不愿与他共融，这使我感到烦恼。就我所记得的，这不是你曾与我交谈的那种方式。至于我派人去\Place{西方}，这完全不可能。我没有合适的人选。事实上，当我环顾四周，我似乎没有人在我这一边。我只能祈祷我能在那七千位没有向巴耳屈膝的人之列。我知道目前迫害我们所有人的仇敌正在索要我的性命；但这不会丝毫减少我对天主的教会所应尽的热情。\par

\SourceRecord{Translated by Blomfield Jackson.；Nicene and Post-Nicene Fathers, Second Series；Vol. 8.；Edited by Philip Schaff and Henry Wace.；Buffalo, NY: Christian Literature Publishing Co.,；1895.；<http://www.newadvent.org/fathers/3202156.htm>.}

% structure: p0001 source=h1 target=section reason=page-title

\section{第一五七号书信}

\Emph{\Place{凯撒勒雅}的\Person{圣巴西略}}\par

\Emph{致\Person{阿米奥古斯}}\par

您可以想象，夏季未能与您会面，我是多么失望；虽然往年的会面已足以使我满足，但即使在梦中看见所爱之人，也能给爱他们的人带来些许安慰。但您甚至不写信，您是如此懒散，我认为您的缺席只能归因于您不愿为亲情而旅行。关于这一点，我不再多言。请为我祈祷，求主不要舍弃我，而是如同祂已使我脱离过去的诱惑一样，也使我脱离那将临的诱惑，以使我所信赖的那位的圣名得光荣。\par

\SourceRecord{Translated by Blomfield Jackson.；Nicene and Post-Nicene Fathers, Second Series；Vol. 8.；Edited by Philip Schaff and Henry Wace.；Buffalo, NY: Christian Literature Publishing Co.,；1895.；<http://www.newadvent.org/fathers/3202157.htm>.}

% structure: p0001 source=h1 target=section reason=page-title

\section{第一五八书}

\Emph{\Person{圣巴西略（凯撒勒雅的）}}\par

\Emph{致\Person{安提约古}}。\par

我的罪阻止了我实现长久以来想与你会面的愿望。请允许我以书信为我的缺席致歉，并恳求你切莫忘记在祈祷中纪念我，好使我若尚存于世，得以享受你的陪伴；若不然，也因你的祈祷，使我满怀善望离开此世。我将我们的兄弟——那骆驼夫——托付给你。\par

\SourceRecord{Translated by Blomfield Jackson.；Nicene and Post-Nicene Fathers, Second Series；Vol. 8.；Edited by Philip Schaff and Henry Wace.；Buffalo, NY: Christian Literature Publishing Co.,；1895.；<http://www.newadvent.org/fathers/3202158.htm>.}

% structure: p0001 source=h1 target=section reason=page-title

\section{\WorkTitle{第159封信}}

\Emph{\Place{凯撒勒雅}的\Person{圣巴西略}}\par

\Emph{致\Person{欧帕特里乌}及其女。}\par

1. 你们不难想象，尊贵之人的来信给了我多大的快乐，仅从其内容便可想而知。的确，对于一个常常祈祷能与敬畏主的人相通并分享其福乐的人来说，还有什么能比这样一封询问关于天主知识的信更令人欣慰呢？因为，如果对我而言，\BibleQuote{生活就是基督}（\BibleRef{斐 1:21}），那么我的言语本当关乎基督，我的一切思想行为本当遵循祂的诫命，我的灵魂也当效法祂。因此，我欢喜被问及这样的事，并祝贺提问者。简言之，在我眼中，在尼西亚聚会的教父们的信仰超越一切后来的发明。在此信仰中，子被宣认为与父\OriginalTerm{同质}{con-substantial}，且在本质上与生养祂的那一位相同，因为祂被宣认为光之光、天主之天主、善之善，等等。那些圣人们宣告了同样的道理，而我也如此宣告，并祈祷能步武他们的芳踪。\par

2. 然而，由于那些总是企图引入新奇事物的人现在提出的问题（即关于圣神的问题），在古时因为未被否认而未被详尽阐述，至今仍未得到充分解释，我将依据圣经的意义对此加以说明。我们怎样受洗，就怎样宣信。我们怎样宣信，就怎样赞美。因此，既然圣洗是救主以父及子及圣神之名赐予我们的，我们就按照圣洗宣写信经，并按照信经献上赞美。我们将圣神与父及子一同荣耀，确信祂并不与天主性分离；因为本性上外在的事物并不分享同样的尊荣。凡称圣神为受造物的人，我们怜悯他们，因为他们借此言论陷入了不可赦免的亵渎圣神之罪。对于稍微受过圣经训练的人，我无需费言证明受造物与天主性有别。受造物是奴隶，而圣神使人自由。受造物需要生命，而圣神是生命的赐予者（\BibleRef{若 6:63}）。受造物需要教导，而圣神教导（\BibleRef{若 14:26}）。受造物被圣化，而圣神圣化（\BibleRef{罗 15:16}）。不论你们提名天使、总领天使或天上万军，他们都通过圣神得到圣化，但圣神自己本质上是圣的，并非因恩赐而得，而是本性所有；因此祂获得了\Quote{圣}之名。因此，我们不允许将本性为圣的那一位（如父本性为圣，子本性为圣）与神圣而可赞美的天主圣三分离割裂，也不接受那些鲁莽地将祂算作受造物一部分的人。我虔诚的朋友们，这简短的摘要对你们足矣。从小小的种子，借着圣神的合作，你们将收获更丰盛的虔敬。\BibleQuote{教导智慧人，他就越有智慧}（\BibleRef{箴 9:9}）。更充分的论证我将推迟到我们见面时。届时，我将能够回答异议，提供更充分的圣经证据，并确认所有正确的信仰准则。目前请原谅我的简短。我本不该写信，但若完全拒绝你们的请求，在我看来比部分允准对你们伤害更大。\par

\SourceRecord{Translated by Blomfield Jackson.；Nicene and Post-Nicene Fathers, Second Series；Vol. 8.；Edited by Philip Schaff and Henry Wace.；Buffalo, NY: Christian Literature Publishing Co.,；1895.；<http://www.newadvent.org/fathers/3202159.htm>.}

% structure: p0001 source=h1 target=section reason=page-title

\section{第一六〇封书信}

\Emph{\Place{凯撒勒雅}的\Person{圣巴西略}}\par

\Emph{致\Person{狄奥多鲁斯}}。\par

1. 我已收到那封以狄奥多鲁斯之名递到我这儿的信，但其内容却更值得归于任何人而非狄奥多鲁斯。似乎有一位善用心计之人冒用了你的名字，意图在听众面前博得赞誉。有人问他，与亡妻的姐妹缔结婚姻是否合法；他非但不对此问题感到战栗，反而无动于衷地听着，并且相当大胆而勇敢地支持了这不体面的欲望。若我手边有此人的信，我必将其寄给你，你便能为自己和真理辩护。但将信给我看的人又将其取走，并把它当作一件胜利的纪念品四处炫耀，反对我——我起初就禁止此事，他却宣称已得到书面许可。如今我写信给你，为要加倍攻击那伪造的文件，使它没有任何力量能伤害它的读者。\par

2. 首先我必须强调，在这种事上最重要的，是我们自己的习惯——它具有法律效力，因为这些规则是由圣人们传授给我们的。习惯如下：若有人因不洁而陷于与两姐妹的非法结合，这不应被视为婚姻，他们也绝不可被接纳进入教会，直到他们彼此分离。因此，即使不能再说什么，这习惯也足以维护正义。但既然那封信的作者试图通过错误论证将这种恶行引入我们的实践，我便不得不借助理性的帮助，尽管在显而易见的事上，每个人的直觉感受比论证更有力。\par

3. 他说，在\BibleBook{肋未}{纪}上写着：\Quote{你不可娶一个妇人来加在她的姊妹身上，\Emph{在她活着的时候}揭露她的下体。}\BibleRef{肋 18:18} 由此他论证说，显然妻子死后娶其姐妹是合法的。对此我的首要回答是：凡律法所说的，都是对律法之下的人说的；否则我们便须服从割损、安息日、禁戒肉类。因为我们绝不可在发现合乎自己喜好之事时，就使自己受律法的奴役；尔后若律法中某事显得严苛，又转而依靠在基督内的自由。有人问，是否写着可以在妻子死后娶其姐妹。让我们说一个稳妥且真实的话：并未写着。但通过推理得出沉默之事，是立法者的职责，而非引用律法条文者的职责。实际上，按此条件，任何人也可以自由地在妻子活着时娶其姐妹。同样的诡辩在此也适用。他说，写着：\Quote{你不可娶一个妇人来加在姊妹身上，}所以，若不致加害，便不禁止娶她。那想纵欲之人会声称，姐妹之间的关系使她们不致彼此加害。去掉了禁止与两人同住的原因，有什么能阻止一人娶两姐妹呢？我们将会说，这并未写着。前者也没有明确写出。推理的推演在这两种情况下都允许自由。但解决困难的办法，或许可以稍稍回溯律法条文的背景。看来立法者并未包括每一种罪恶，而是特别禁止了埃及人的行为——以色列人从那里出来——以及迦南人的行为——以色列人正往那里去。话是这样说的：\Quote{你们不可照埃及地的行为居住在那里，也不可照迦南地的行为，就是我领你们去的地方的行为，也不可随从他们的习俗。}\BibleRef{肋 18:3} 很可能这类罪恶在当时的外邦人中并未流行。在此情况下，立法者大概无需防范它；不成文的习惯足以定罪这罪行。那么，为什么在禁止较大的罪恶时，却对较小的保持沉默呢？因为许多放纵肉体以致与姐妹同居之人，认为祖宗的榜样是有害的。我应当怎样做？引用圣经，还是阐述圣经未明言之事？在这些律法中，并未写着父亲和儿子不应拥有同一妾室，但在先知书中，这被视为最应受谴责的事，经上说：\Quote{人与他的父亲\InnerQuote{同去亲近同一少女}。}\BibleRef{亚 2:7} 还有许多其他形式的不洁情欲，由魔鬼的学校发明，而神圣的圣经对它们保持沉默，不愿以污秽之物的名号玷污自己的尊严，只是笼统地谴责它们的不洁。正如宗徒保禄所说：\Quote{至于邪淫和一切不洁……在你们中间连提也不可提，如圣人所相称的，}\BibleRef{弗 5:3} 这样，他将男女间不可言说的行为都归于不洁之名。由此可见，沉默绝不给予纵欲者自由。\par

4. 然而，我坚持认为这一点并非被忽略，而是立法者明确禁止了。\BibleQuote{你们中任何人不可接近任何自己的近亲，暴露他们的下体}\BibleRef{肋 18:6} 也包括了这种亲属关系，因为还有什么比自己的妻子，即自己的骨肉，更亲近呢？\BibleQuote{他们不再是两个，而是一体了}\BibleRef{玛 19:6}。因此，通过妻子，妻子的姐妹就成了丈夫的亲属。因为正如他不可娶妻子的母亲，也不可娶妻子的女儿，因为他不可娶自己的母亲，也不可娶自己的女儿；同样，他不可娶妻子的姐妹，因为他不可娶自己的姐妹。另外，妻子与丈夫的亲属结合也是不合法的，因为亲属关系的权利在双方都有效。至于我，我要向每一个考虑婚姻的人作证：\BibleQuote{这世界的局面正在逝去}\BibleRef{格前 7:31}，时日无多了：\BibleQuote{从今以后，那些有妻子的，要像没有一样}\BibleRef{格前 7:29}。如果有人不适当地引用\BibleQuote{你们要生育繁殖，充满大地}\BibleRef{创 1:28}这条命令，我就嘲笑他，因为他没有辨别时代的征兆。第二次婚姻是治疗淫乱的补救措施，而不是好色的手段。经上说：\BibleQuote{如果他们不能节制，就让他们结婚}\BibleRef{格前 7:9}；但如果他们结婚，就不可违反法律。\par

5. 但是那些被可耻的私欲弄瞎了灵魂的人，甚至不顾从古时就区分家族名称的自然。因为那些缔结这些婚姻的人将如何称呼他们的儿子？他们称他们为兄弟还是堂兄弟？由于混乱，两个名称都会适用。人啊，不要让姨母变成小儿的继母；不要用无法平息的嫉妒武装她，而她本应以母爱养育他们。只有继母才将仇恨延续到死亡之后；其他敌人与死者和解；继母却在死后开始仇恨。我所说的话总结如下：如果有人想缔结合法的婚姻，整个世界向他敞开；如果他只是被私欲所驱使，就让他更受限制，\BibleQuote{好使他学习怎样用圣洁和尊贵持守自己的器皿，而不是放纵邪淫的私欲}。我本想说更多，但我信件的篇幅不允许了。我祈祷我的劝诫能战胜私欲，或者至少这种污秽不会出现在我自己的教区。它在何处被冒险施行，就让它留在那里吧。\par

\SourceRecord{Translated by Blomfield Jackson.；Nicene and Post-Nicene Fathers, Second Series；Vol. 8.；Edited by Philip Schaff and Henry Wace.；Buffalo, NY: Christian Literature Publishing Co.,；1895.；<http://www.newadvent.org/fathers/3202160.htm>.}

% structure: p0001 source=h1 target=section reason=page-title

\section{第一六一封信}

\Person{凯撒勒雅的圣巴西略}\par

致\Person{盎波罗削}，于其被祝圣为主教之际。\par

1. 愿天主受赞美，祂从世世代代拣选祂所喜悦的人，分别出蒙选之器，并用他们为圣人服务。虽然你曾试图逃避——如你所承认的，不是逃避我，而是逃避你预料要通过我而来的召叫——祂却以恩宠的稳妥之网捕获了你，将你带到了\Place{丕息狄雅}的中心，为的是为主捕人，将魔鬼的猎物从深渊带到光明中。你也可以像蒙福的\Person{达味}那样说：\Quote{我往何处，才能脱离你的圣神？我往何处，才能逃避你的面容？}这就是我们慈爱的主的奇妙作为。\Quote{驴丢了}，好使以色列有了一位君王。\Person{达味}既是以色列人，便被赐予以色列；然而，那养育了你、并使你达到如此崇高德行的土地，不再拥有你，却看见邻邦因自己的装饰而变得华美。但所有在基督内的信徒都是一个子民；所有基督的子民，尽管祂从许多地区受到欢呼，却是一个教会；因此，我们的家乡为天主的安排而欢喜快乐，不认为自己缺少了一个人，反而认为藉着一个人，她拥有了整个教会。唯愿主恩准我既能亲眼见到你，又在我与你分离的日子里，听到你在福音上的进步和你各教会的良好秩序。\par

2. 那么，你要刚强，要勇敢，在至高者托付给你手中的人民面前行走。如同一位有技巧的舵手，你的心灵要超越那被异端之风掀起的每一个波浪；要保护船只不被虚假教义的咸苦波涛所淹没；并等待主要平息风浪的时刻，一旦有相称的声音唤醒祂去斥责风和海。你若想来看望如今因久病而趋近不可避免终点的我，不要等待时机或我的话语。你知道，在一位父亲的心中，任何时刻都适合拥抱他所深爱的儿子，而亲情胜过言语。不要为超过你力量的责任而悲伤。如果你注定要独自承担重担，那不仅是沉重的，更是无法忍受的。但若主与你分担重担，\Quote{将你们的一切挂虑都托付给主}，祂自会行动。只愿你时刻警惕，免得被邪恶的习俗所迷惑。反而要藉着天主赐给你的智慧，将以往的一切恶习转变为善。因为基督派遣你不是去跟随别人，而是你自己去引导所有正在得救的人。我恳求你为我祈祷，如果我仍在此世，愿我能见到你和你的教会。然而，如果命定我如今就要离去，但愿我在主内见到你们众人，你的教会如同一棵葡萄树，因善工而开花结果；而你自己，如同一精明的农夫和良善的仆人，按时给同伴分粮，并领取那明智忠信管家的赏报。所有与我同在的人都问候你的尊威。愿你在主内坚强而喜乐。愿你在圣神和智慧的恩宠中荣耀地保持完好。\par

\SourceRecord{Translated by Blomfield Jackson.；Nicene and Post-Nicene Fathers, Second Series；Vol. 8.；Edited by Philip Schaff and Henry Wace.；Buffalo, NY: Christian Literature Publishing Co.,；1895.；<http://www.newadvent.org/fathers/3202161.htm>.}

% structure: p0001 source=h1 target=section reason=page-title

\section{第一六二函}

\Emph{\Person{凯撒勒雅的圣巴西略}}\par

\Emph{致\Person{萨莫萨塔主教欧瑟伯}}\par

同一原因似乎既使我迟疑写信，又证明我不得不写。当我想到我欠的拜访，并估计与你会面的收获时，我不禁轻视书信，认为与实在相比，它们连影子都不是。另一方面，当我想到我唯一的安慰——我既被剥夺了一切最好和最重要的事物——就是问候这样一个人，并照例恳求他不要在他的祈祷中忘记我，我便想到书信并非毫无价值。我自己并不愿放弃拜访的一切希望，也不对见到你感到绝望。我若不能显得对你的祈祷有如此大的信心，以致期望从老人变成青年（倘若需要如此），而不仅仅是从现在的病弱憔悴变得稍强健一些，那我将感到羞愧。用言语表达我未能与您同在的原因并不容易，因为我不仅受实际疾病的阻碍，而且甚至没有足够的口才随时向您解释如此多样而复杂的病情。我只能说，从复活节至今，发烧、腹泻和肠道紊乱如波浪般淹没我，不让我抬头。\Person{巴拉古斯}弟兄或许能告诉你我症状的性质，即便不能如严重性所当，至少足够清晰地让你明白我耽搁的原因。你若衷心地加入我的祈祷，我毫不怀疑我的痛苦将轻易逝去。\par

\SourceRecord{Translated by Blomfield Jackson.；Nicene and Post-Nicene Fathers, Second Series；Vol. 8.；Edited by Philip Schaff and Henry Wace.；Buffalo, NY: Christian Literature Publishing Co.,；1895.；<http://www.newadvent.org/fathers/3202162.htm>.}

% structure: p0001 source=h1 target=section reason=page-title

\section{第一六三封信}

\Emph{凯撒勒雅的圣巴西略}\par

\Emph{致约维努斯伯爵}.\par

从你的信中可以看出你的灵魂，因为实际上没有任何画家能如此精确地捕捉外在的相似，如同说出的话语能够映照灵魂的秘密。当我读你的信时，你的言辞准确地刻画了你的坚毅、你真正的尊严、你始终如一的真诚；尽管不能见到你，但在所有这些事情上，它极大地安慰了我。所以，要永远抓住每一个给我写信的机会，使我能享受与你远距离交谈的乐趣；因为我现在因自己健康状况的痛苦而不能与你当面相见。这有多严重，你将从天主所爱的主教\Person{安菲洛基乌斯}那里得知，他因常与我在一起而能向你报告，并且完全有能力告诉你他所见到的。但我之所以希望你知道我的痛苦，唯一的原因就是你将来会原谅我，并免除我缺乏活力的责任，如果我没能来看你；虽然实际上，我的损失并不太需要我的辩护，而更需要你的安慰。如果我能到你那里去，我本会更愿意看到你的阁下，胜过别人认为值得努力的所有目标。\par

\SourceRecord{Translated by Blomfield Jackson.；Nicene and Post-Nicene Fathers, Second Series；Vol. 8.；Edited by Philip Schaff and Henry Wace.；Buffalo, NY: Christian Literature Publishing Co.,；1895.；<http://www.newadvent.org/fathers/3202163.htm>.}

% structure: p0001 source=h1 target=section reason=page-title

\section{书信一六四}

\Emph{\Person{圣巴西略·凯撒勒雅}}\par

\Emph{致\Person{阿斯科利乌}}\par

1. 我实在难以言表，读到您圣德的来信时我是多么喜悦。我的言辞太过贫弱，无法表达我所有的感受；然而，您应当能从您所写之优美中揣测到我的心情。因为您的信函中难道不包含什么呢？它包含对天主的爱；对殉道者的奇妙描述，如此清晰地呈现他们英勇奋斗的样式，使我仿佛亲眼目睹；对我本人的爱与善意；以及无比优美的言辞。因此，当我将它拿在手里，反复阅读，并看出它充满圣神的恩宠时，我以为自己回到了那美好的古老时代，那时天主的教会欣欣向荣，扎根于信德，在爱德中合一，所有肢体如同在一个身体内和谐共处。那时迫害者是明显的，被迫害者也是明显的。那时人民因受攻击而更加增多。那时殉道者的鲜血灌溉了教会，滋养了更多真宗教的斗士，每一代人都以先辈的热忱投入战斗。那时我们基督徒彼此和睦相处，拥有主留给我们的平安，然而我们如此残酷地将它从我们中间驱走，如今连一丝痕迹也不剩。然而，当一封来自远方的信函，闪耀着爱德之美，以及一位来自\Place{多瑙河}以外蛮荒之地的殉道者亲自来到我这里，以自身宣讲那里所持守的信仰之精确性时，我的灵魂确实回到了那古老的福乐中。谁能述说这一切给我灵魂带来的喜悦？有什么言论能胜任描述我内心深处的一切感受？然而，当我看到这位斗士时，我祝福他的训练者：他在公义的审判者面前，在坚固了许多人为真宗教而战之后，也必将领受正义的冠冕。\par

2. 您使我记起真福的\Person{欧提基乌}，并因我的故乡播下了真宗教的种子而尊荣它，这既因您对往事的提醒而使我喜悦，又因您对现今的深信而使我忧愁。现在我们中没有一个人能接近\Person{欧提基乌}的良善；我们不但没有以圣神的柔和力量及祂恩宠的运作来驯化蛮族，反而因我们的大罪将心地温和的人变为蛮族，因为我认为我们自己和我们的罪导致了异端的影响如此广泛传播。或许世上没有哪个部分逃脱了异端的烈火。您告诉我斗士们的奋斗，身体为真理而受鞭打，凶猛的怒火被无畏的心所轻视，迫害者的各种折磨，以及摔跤手们贯穿这一切的坚毅，还有殉道者致死的刀斧和水刑。而我们的情况又如何？爱德变冷了；教父的教导正在荒废；信仰处处遭遇海难；信徒们的口舌沉默；人民被赶出祈祷之所，在露天向他们在天上的主举起双手。我们的苦难深重，殉道无处可见，因为虐待我们的人与我们同称一名。因此，请您亲自向主祈祷，并请所有基督的崇高斗士与您一同为教会祈祷，好使若这世界尚存些许时日，万物尚未被推向毁灭，愿天主与祂自己的教会和好，并恢复它们古老的平安。\par

\SourceRecord{Translated by Blomfield Jackson.；Nicene and Post-Nicene Fathers, Second Series；Vol. 8.；Edited by Philip Schaff and Henry Wace.；Buffalo, NY: Christian Literature Publishing Co.,；1895.；<http://www.newadvent.org/fathers/3202164.htm>.}

% structure: p0001 source=h1 target=section reason=page-title

\section{\WorkTitle{第一六五号书信}}

\Person{圣巴西略．凯撒勒雅}\par

致德撒洛尼主教\Person{阿斯科略}\par

天主满全了我昔日的祈祷，屈尊让我收到您真正圣德的来信。我最渴望的是与您相见、被您看见，并在实际的交往中享受您所领受的一切圣神的恩宠。但因我们相隔遥远，且各自事务繁忙，这无法实现。因此，我其次祈求我的灵魂能因您在基督内的爱的频繁来信而得滋养。如今，当我手中握着您的信函时，这已蒙恩赐予。我因享受您的通讯而双倍喜悦。我感到仿佛真的能在您的话语中如镜般看见您的灵魂；我且欣喜至极，不仅因为您如一切证言所述，更因您的高贵品质是我故乡的光荣。您像一棵从荣耀的根上发出的旺盛枝条，以属灵的果实充满了我们疆界之外的土地。因此，我们的故乡为她自己的苗裔而欢庆。当您为信德而战时，她听说教父们的美好遗产在您内得以保存，她便光荣了天主。如今您在做什么？您因送一位刚刚在你们边界的蛮族土地上打了一场好仗的殉道者给生养您的土地而荣耀了它，如同一位感恩的园丁将他初熟的果实送给赐予他种子的人一般。这礼物确实堪当基督的斗士，一位刚戴上正义冠冕的真理殉道者；我们欣然迎接他，光荣那如今在普世实现了祂的基督的福音的天主。请允许我在您内心爱我的祈祷中纪念我，并为我灵魂的缘故恳切祈求主，使我有一天也能堪当开始按照祂为赐予我们救恩而颁布的诫命之道事奉天主。\par

\SourceRecord{Translated by Blomfield Jackson.；Nicene and Post-Nicene Fathers, Second Series；Vol. 8.；Edited by Philip Schaff and Henry Wace.；Buffalo, NY: Christian Literature Publishing Co.,；1895.；<http://www.newadvent.org/fathers/3202165.htm>.}

% structure: p0001 source=h1 target=section reason=page-title

\section{第一六七号书信}

\Emph{圣巴西略·凯撒勒雅}\par

\Emph{致 \Person{欧瑟伯}，\Place{萨莫撒塔}主教}。\par

您记得我并写信给我，我已感到喜悦；而更为重要的是，您在信中赐予我您的祝福。倘若我堪当承受您的劳苦及其在基督事业中的奋斗，我便应获准来到您身边，拥抱您，并以您为忍耐的楷模。但既然我不配如此，且被诸多痛苦和繁忙所困，我便行其次：我向阁下致敬，并恳求您不要厌倦对我的怀念。因为收到您的书信的荣幸与喜乐，不仅对我有益，还使我在世人面前有夸耀和骄傲的资本：我竟能受到一位德行如此高超、与天主如此亲密共融的尊者的尊崇。这位尊者既能以其教导，也能以其榜样，使他人与他一同融于此共融之中。\par

\SourceRecord{Translated by Blomfield Jackson.；Nicene and Post-Nicene Fathers, Second Series；Vol. 8.；Edited by Philip Schaff and Henry Wace.；Buffalo, NY: Christian Literature Publishing Co.,；1895.；<http://www.newadvent.org/fathers/3202167.htm>.}

% structure: p0001 source=h1 target=section reason=page-title

\section{第一六八封书信}

\Emph{凯撒勒雅的圣巴西略}\par

\Emph{致安提约古}。\par

我为那失去如此牧者引导的教会哀恸。但我更有理由祝贺你，因你在这时刻配享与那位为真理打这美好仗者交往的特恩，并且我确信，你这崇高支持并激励其热忱的人，必被主看为配得与他同样的命运。能在这不间断的安宁中，与这位学问渊博、阅历丰富的人交往，是何等的福分！如今，至少我确信你必须知道他有多智慧。往昔，他的心思必然分散于许多挂虑，而你也过于忙碌，无法全心留意那从他纯洁心灵涌出的灵性泉源。愿天主赐你能安慰他，而你自身永不需要他人的慰藉。我确信你内心的倾向，既从我与你的短暂交往中得知，也因你那卓越教师的高超训导——与他一日的相处，便足以为救恩之旅预备充分。\par

\SourceRecord{Translated by Blomfield Jackson.；Nicene and Post-Nicene Fathers, Second Series；Vol. 8.；Edited by Philip Schaff and Henry Wace.；Buffalo, NY: Christian Literature Publishing Co.,；1895.；<http://www.newadvent.org/fathers/3202168.htm>.}

% structure: p0001 source=h1 target=section reason=page-title

\section{第一六九封书信}

\Emph{\Person{凯撒勒雅的圣巴西略}}\par

\Emph{\Person{巴西略}致\Person{额我略}}.\par

您承担了一项仁慈而富有爱德的工作，把傲慢的\Person{格利塞里乌斯}（目前我必须这样写）的被掳的队伍聚集起来，并尽可能您所能，掩盖了我们共同的耻辱。您的阁下理应充分了解关于他的事实，以消除这一耻辱。\par

这位严肃而可敬的\Person{格利塞里乌斯}，被我祝圣为\Place{维内萨}教会的执事，以协助司铎并照顾教会的工作，尽管这个人在其他方面难以驾驭，但在手工劳动上他天生灵巧。他一就职就忽略自己的工作，仿佛完全没有事可做。但他凭借自己的私人权力和权威，聚集了一些可怜的贞女，其中有些是自愿来找他的（您知道年轻人多么容易倾向于这类事），另一些则是不情愿地被迫接受他作为她们团体的领袖。然后他自称宗主教，突然开始扮演大人物的角色。他达到这一点并非基于任何合理或虔诚的理由；他唯一的目标是获得谋生手段，就像有些人从事这个行业，有些人从事那个行业一样。他几乎颠覆了整个教会，鄙视自己的司铎——一位品性和年龄都可敬的老人，鄙视他的乡村主教，也鄙视我，完全不当一回事，不断用混乱和扰动充满城镇和全体圣职人员。现在，当我和他的乡村主教温和地责备他，以免他轻视我们（因为他试图煽动年轻人同样不服从）时，他正在策划最胆大妄为和不人道的行为。他抢劫了尽可能多的贞女后，夜间逃跑了。我相信这一切在您看来一定非常可悲。也想想时间：当时那里正在举行庆节，自然聚集了大量人群。然而，他那边却带出了他自己的队伍，他们跟随年轻男子并围着他们跳舞，使所有正派人士极为痛心，而轻浮的饶舌者则放声大笑。即使这样还不够，尽管丑闻巨大。我听说，连贞女的父母，因无法忍受失去女儿，希望把离散的队伍带回家，自然而然地叹息流泪跪在女儿脚下，却被这位优秀青年和他的匪徒队伍侮辱和凌辱。我相信您的阁下会认为这一切不可容忍。它的嘲笑同样落到我们所有人身上。首先，命令他带着贞女们回来。如果他带着您的一封信回来，或许能得到一些怜悯。如果您不采取这一措施，至少把贞女们送回她们的母亲——教会。如果这也不能做到，至少不要对愿意回来的人施加任何暴力，而是让她们回到我这里。否则，我呼求天主和人作证，这一切都是恶行，违反了教会法律。最好的办法是让\Person{格利塞里乌斯}带着信回来，以适当而正确的心态；否则，就撤销他的牧职。\par

\SourceRecord{Translated by Blomfield Jackson.；Nicene and Post-Nicene Fathers, Second Series；Vol. 8.；Edited by Philip Schaff and Henry Wace.；Buffalo, NY: Christian Literature Publishing Co.,；1895.；<http://www.newadvent.org/fathers/3202169.htm>.}

% structure: p0001 source=h1 target=section reason=page-title

\section{第一七〇书}

\Emph{\Person{凯撒勒雅的圣巴西略}}\par

\Emph{致\Person{格莱契乌斯}}.\par

你的疯狂愚妄要到何时尽头？你还要多久才停止图谋危害自己？你还要多久继续激怒我，并使全体独修者的共同规诫蒙羞？回转吧！信赖天主，也信赖我——我效法天主的慈爱。我虽如父亲般责备了你，也要如父亲般宽恕你。这就是你将从我这里得到的对待，因为有许多人在为你求情，而首当其冲的便是你自己的司铎，我对他那满头白发和慈悲心肠深怀敬意。你若继续远避我，就必完全失去你的品位。你也要远离天主，因为你的歌声和装扮正在引导那些年轻女子，并非走向天主，而是走向深渊。\par

\SourceRecord{Translated by Blomfield Jackson.；Nicene and Post-Nicene Fathers, Second Series；Vol. 8.；Edited by Philip Schaff and Henry Wace.；Buffalo, NY: Christian Literature Publishing Co.,；1895.；<http://www.newadvent.org/fathers/3202170.htm>.}

% structure: p0001 source=h1 target=section reason=page-title

\section{第一七一号书信}

\Emph{圣巴西略·凯撒勒雅}\par

\Emph{致\Person{额我略}}.\par

我不久前写信给你，论及\Person{格利色里乌斯}和那些贞女。至今他们仍未返回，依旧犹豫不决，如何为何我不得而知。我若将此事归咎于你，以为你如此行事是为使我蒙羞，或因你对我有所不满，或为取悦他人，我将感到遗憾。那么就让他们来，无所畏惧。请你担保他们这样做。因为我被切除肢体感到痛苦，虽然他们被切除是合宜的。若他们固执己见，担子将落在他人身上。我对此撒手不管。\par

\SourceRecord{Translated by Blomfield Jackson.；Nicene and Post-Nicene Fathers, Second Series；Vol. 8.；Edited by Philip Schaff and Henry Wace.；Buffalo, NY: Christian Literature Publishing Co.,；1895.；<http://www.newadvent.org/fathers/3202171.htm>.}

% structure: p0001 source=h1 target=section reason=page-title

\section{第一七二封书函}

\Emph{\Person{圣巴西略} \Place{凯撒勒雅}}\par

致\Person{索弗若尼}主教\par

无需多言，你的来信令我何等欣喜。你自己便能推知我收到信时的感受。你在信中向我展示了圣神首生的果实，即爱德。在现今的时势下，还有什么比这更珍贵的呢？因为不义增多，许多人的爱德已冷淡了。\BibleRef{玛 24:12} 如今，少有比与弟兄的精神交往、和平的言语以及我在你身上所找到的精神共融更稀罕的事了。为此我感谢主，恳求祂使我得以分享在你内圆满的喜乐。如果你的信尚且如此，那么与你本人相见又当如何？若你远在异地已如此打动我，当你面对面相见时，你又会如何对待我呢？你要确信，若非我被无数事务缠身，且被种种无法避免的忧虑所束缚，我早已赶去拜会你的尊驾了。尽管我旧有的疾患对我行动大为妨碍，但鉴于我所期待的益处，我本不会让此成为障碍。能得与一位怀抱相同见解且尊崇祖先信德的人会面——正如我们可敬的弟兄和同司铎所说的你所行的那样——实则是回归教会古时的福乐，那时因不当争辩而受苦的人极少，众人都生活在和平中，是遵守诫命的\Quote{工人}，不\Quote{需要羞愧}，\BibleRef{弟后 2:15} 以单纯明净的宣认服事主，并以纯正不可侵犯的态度保存其在对圣父、圣子及圣神的信仰。\par

\SourceRecord{Translated by Blomfield Jackson.；Nicene and Post-Nicene Fathers, Second Series；Vol. 8.；Edited by Philip Schaff and Henry Wace.；Buffalo, NY: Christian Literature Publishing Co.,；1895.；<http://www.newadvent.org/fathers/3202172.htm>.}

% structure: p0001 source=h1 target=section reason=page-title

\section{第一七三封书信}

\Emph{\Person{圣巴西略} \Place{凯撒勒亚}}\par

\Emph{致\Person{德敖多拉}修女院院长}。\par

我本该更勤于写信给你，但因为我相信，我的朋友，我的信件并非总能到达你本人手中。我担心，因我所依赖之人的顽劣，尤其在此整个世界处于混乱之际，许多其他人会截获这些信件。因此我等待被责备，并热切期待有人来索要我的信，这样我才能证明它们已送达。然而，无论我写或不写，有一件事我从不懈怠，那就是在心中铭记对阁下的记忆，并祈求主恩赐你完成你所选择的善生旅程。因为事实上，对于一个许下誓愿的人来说，履行誓愿所包含的一切并非易事。任何人都可以接受福音生活，但据我所知，只有极少数人能在最细微的细节上完全履行其职责，且不忽略其中包含的任何事项。只有极少数人能够按照福音的要求，一贯克制舌头、引导眼睛；以取悦天主的目标用手劳作；移动双足，使用每个肢体，如同造物主从起初所规定的那样。衣着得体，与人交往时谨慎，饮食有节，避免在获取必需品时过度；这些事若只是这样提及，似乎微不足道，但据我经验，一贯遵守它们需要不小的挣扎。再者，如此完美的谦卑——甚至不记得家族的高贵，也不因我们可能拥有的任何身体或心灵的天然优点而自高，也不容许他人对我们的看法成为骄傲与高抬的理由——这一切都属于福音生活。还有持久的自制，勤于祈祷，手足之爱的同情，对穷人的慷慨，卑逊的性情，痛悔的心，健全的信德，在抑郁中的平静，同时我们永不忘记那可怕且无可避免的审判座。我们都在匆匆奔赴那审判，但记得它并忧虑其后之事的人却极少。\par

\SourceRecord{Translated by Blomfield Jackson.；Nicene and Post-Nicene Fathers, Second Series；Vol. 8.；Edited by Philip Schaff and Henry Wace.；Buffalo, NY: Christian Literature Publishing Co.,；1895.；<http://www.newadvent.org/fathers/3202173.htm>.}

% structure: p0001 source=h1 target=section reason=page-title

\section{Letter 174}

\Emph{凯撒勒雅的圣巴西略}\par

\Emph{致一位寡妇}\par

我一直极愿常写信给尊驾，但我屡次克制自己，生怕因那班对我怀有恶意之人，给您带来任何试探。据我所知，他们的仇恨甚至到了这种地步：若有人偶然收到我一封信，他们便要大惊小怪。如今您已亲自写信给我——您这样做实在很好，将您心中所想的一切必要信息都告诉我——我便被激励而回信给您。那么，让我弥补过去所缺失的，同时也答复尊驾所写之事。真有福的灵魂，她昼夜没有别的挂虑，只想着当那伟大的日子来临，所有受造物都将站在审判者面前，并为自己的行为交账时，她也能轻松摆脱生命的清算。\par

因为那常将那日子和那时刻摆在眼前，并不断思想在无可推诿的审判台前如何辩护的人，将完全不犯罪，或至少不犯重罪，因为我们开始犯罪，是由于我们内心缺乏对天主的敬畏。当人清楚地意识到所威胁他们的事时，他们内在的敬畏决不允许他们陷入轻率的行动或思想。因此，要记念天主。将祂的敬畏存记在心中，并动员众人与您一同祈祷，因为那些能以恳切之心感动天主之人的帮助是伟大的。切勿停止这样做。即使我们在此肉身生命中生活，祈祷也将成为我们的大能助佑；当我们离开此世时，它也将成为我们前往来世旅途的充足预备。\par

忧虑是好的；但另一方面，沮丧、气馁和绝望于我们的救恩，却对灵魂有害。因此，要信赖天主的良善，期待祂的援助，知道倘若我们正确而真诚地转向祂，祂不仅永远不会永远抛弃我们，而且甚至在我们正在说出祈祷的言语时，祂会对我们说：\Quote{看！我与你同在。}\par

\SourceRecord{Translated by Blomfield Jackson.；Nicene and Post-Nicene Fathers, Second Series；Vol. 8.；Edited by Philip Schaff and Henry Wace.；Buffalo, NY: Christian Literature Publishing Co.,；1895.；<http://www.newadvent.org/fathers/3202174.htm>.}

% structure: p0001 source=h1 target=section reason=page-title

\section{书信第一七五封}

\Person{圣巴西略·凯撒勒雅}\par

\Emph{致伯爵\Person{玛格尼阿努斯}}.\par

近日阁下致函于我，除其他事宜外，明确嘱咐我论及信德之事。我欣赏尔于此事的热情，并向天主祈祷，愿你常怀对善事的选择，且在知识和善工上不断进步，得以完善。但我不愿留下一部论信德的著作，或写下不同信经，因此谢绝寄送你所索要的文件。在我看来，你身边充斥着那些人等的喧嚷——无所事事之辈，他们编造某些言辞来诽谤我，以为通过可耻地对我撒谎便可改善自身的处境。过去他们的所作所为已昭然若揭，而未来的经验将以更为鲜明的色彩将其暴露。然而，我奉劝所有信赖基督的人，切勿忙于反对古代信德，而应如我们相信的那样受洗，如我们受洗的那样奉献光荣颂。我们只需宣认那从圣经所领受的名字，并避免关于它们的一切创新。我们的救恩不在于发明称谓的方式，而在于对天主性的正确宣认——我们所信仰的天主性。\par

\SourceRecord{Translated by Blomfield Jackson.；Nicene and Post-Nicene Fathers, Second Series；Vol. 8.；Edited by Philip Schaff and Henry Wace.；Buffalo, NY: Christian Literature Publishing Co.,；1895.；<http://www.newadvent.org/fathers/3202175.htm>.}

% structure: p0001 source=h1 target=section reason=page-title

\section{第一七六封信}

\Emph{\Person{凯撒勒雅的圣巴西略}}\par

\Emph{致\Person{依科尼雍}主教\Person{安斐罗希}}。\par

但愿当你收到这封信时，天主赐你健康、安闲，并如你所愿。因为那时，我邀请你莅临本城，为增添我们教会每年为纪念殉道者所举行庆典的荣光，便不会徒然。我至尊敬且亲爱的朋友，请你确信，我们这里的子民虽曾经验过许多人，却无人像你那样热切渴望你的临在；你与他们短暂的交往留下了如此亲切的印象。因此，为使主受光荣，子民欢欣，殉道者得尊崇，并使我晚年能从我真正的儿子那里得到应有的关怀，请勿推辞，尽快来我处。我也恳请你提前来到集会之日，以便我们安闲交谈，藉互相分享神恩而彼此安慰。日期是九月五日。那么，请你提前三天来到，好使你也以你的临在尊荣医院教会。愿主藉恩宠保祐你在主内身心安康，精神愉快，为我及天主的教会祈祷。\par

\SourceRecord{Translated by Blomfield Jackson.；Nicene and Post-Nicene Fathers, Second Series；Vol. 8.；Edited by Philip Schaff and Henry Wace.；Buffalo, NY: Christian Literature Publishing Co.,；1895.；<http://www.newadvent.org/fathers/3202176.htm>.}

% structure: p0001 source=h1 target=section reason=page-title

\section{第一七七书}

\Emph{\Person{凯撒勒雅的圣巴西略}}\par

\Emph{致\Person{索弗若尼乌斯}阁下}.\par

要一一列举那些因我的缘故蒙受阁下恩惠的人，并非易事；因我所知，借由您的慷慨援助，我已使许多人受益，这是主在如此严峻的时期赐予我的恩赐。其中最值得称道的，便是此信引荐给您的可敬弟兄\Person{优西比乌}；他遭到荒唐的诽谤，而惟有您——以您的正直——才能消除这诽谤。因此，我恳请您，既出于正义，亦出于仁爱，赐予我惯常的恩惠：将\Person{优西比乌}的事当作您自己的事，予以支持、捍卫，同时亦捍卫真理。他拥有正义在侧，这并非小事；若他不在当前危机中倒下，必将毫不费力地、明白无误且无可辩驳地证明这一点。\par

\SourceRecord{Translated by Blomfield Jackson.；Nicene and Post-Nicene Fathers, Second Series；Vol. 8.；Edited by Philip Schaff and Henry Wace.；Buffalo, NY: Christian Literature Publishing Co.,；1895.；<http://www.newadvent.org/fathers/3202177.htm>.}

% structure: p0001 source=h1 target=section reason=page-title

\section{第一七八号书信}

\Emph{圣巴西略（凯撒勒雅）}\par

\Emph{致阿布尔吉乌}。\par

我知道我已多次将许多人推荐给阁下，并在危急时刻对困境中的朋友大有帮助。但我认为，我从未曾向您推荐过一位比我更敬重的人，或是一位参与更重要斗争的人，就是我现在把此信交到您手中的、我十分亲爱的儿子\Person{欧瑟伯}。若有机会，他将亲自告知阁下他所陷入的困境。至少，我应当说明以下几点：他不应被误解，也不应因许多人被证明行为不端而受到普遍怀疑。他应当得到公正的审判，他的生活也必须接受调查。这样，对他的指控的不实之处将得以显现，而他在享受您的正义保护之后，必将永远宣扬他对您仁慈的感激。\par

\SourceRecord{Translated by Blomfield Jackson.；Nicene and Post-Nicene Fathers, Second Series；Vol. 8.；Edited by Philip Schaff and Henry Wace.；Buffalo, NY: Christian Literature Publishing Co.,；1895.；<http://www.newadvent.org/fathers/3202178.htm>.}

% structure: p0001 source=h1 target=section reason=page-title

\section{第一百七十九号书信}

\Person{圣巴西略} \Place{凯撒勒雅}\par

致\Person{阿林陶}。\par

您天性崇高，平易近人，使我视您为自由与人的朋友。因此，我毫不犹豫地代一位因祖先世系而显赫、但因自身天赋的善良气质而更值得尊敬和荣誉的人向您进言。我恳请您，因我的请求，在他面临一项事实上可笑、但因指控严重而难以应对的法律诉讼时，给予支持。若您肯为他说句好话，对他的成功将至关重要。首先，您将帮助正义；其次，您将以此向我——您的朋友——表示您一贯的尊重和善意。\par

\SourceRecord{Translated by Blomfield Jackson.；Nicene and Post-Nicene Fathers, Second Series；Vol. 8.；Edited by Philip Schaff and Henry Wace.；Buffalo, NY: Christian Literature Publishing Co.,；1895.；<http://www.newadvent.org/fathers/3202179.htm>.}

% structure: p0001 source=h1 target=section reason=page-title

\section{书信一八〇}

\Person{圣巴西略（凯撒勒雅）}\par

\Quote{致索夫罗尼乌斯大师，为尤纳提乌斯陈情。}\par

我遇见一位配受尊敬的之人身陷极大的困境，心中甚是忧伤。作为人，我怎能不同情一位品格高尚、却遭际与其功德不相称之苦的人呢？当我思量如何方能对他有所裨益时，我确曾找到一个助其脱离困境的方法，那就是使他与大人您相识。如今，也请您将曾施予许多人的同样善举，也推广到他身上。您可以从他呈递给皇帝们的请愿书中知晓此事的全部情况。我恳请您亲自过目此件文书，并全力帮助他。您将帮助一位基督徒、一位士绅，一位其深厚学识理应赢得尊重的人。若我再补充一句，即帮助他便是对我施予极大的恩惠——尽管我的小事确实微不足道，但因您常如此善待，使之变得重要，故您对我的恩赐将非同小可。\par

\SourceRecord{Translated by Blomfield Jackson.；Nicene and Post-Nicene Fathers, Second Series；Vol. 8.；Edited by Philip Schaff and Henry Wace.；Buffalo, NY: Christian Literature Publishing Co.,；1895.；<http://www.newadvent.org/fathers/3202180.htm>.}

% structure: p0001 source=h1 target=section reason=page-title

\section{\WorkTitle{信函第一八一}}

\Emph{\Person{凯撒勒雅的圣巴西略}}\par

\Emph{致\Place{梅利泰内}的主教\Person{奥特雷乌斯}}\par

我知道，阁下对于非常敬爱天主的主教\Person{优西比乌}的离世感到的悲痛，不亚于我。我们两人都需要安慰。让我们彼此安慰吧。请你写信告诉我你在\Place{萨摩萨塔}听到的消息，我也会把我从\Place{色雷斯}可能获得的任何情况告知你。\par

知晓民众的坚贞，对我而言是当前苦难不小的慰藉。同样，你若得到我们共同的神父的消息，想必也是如此。当然，我现在无法通过书信告诉你这些，但我向你推荐一位完全知情的人，他会向你报告他离开他时的情况，以及他如何承受他的磨难。那么，请为他和我祈祷，愿主赐他早日脱离苦难。\par

\SourceRecord{Translated by Blomfield Jackson.；Nicene and Post-Nicene Fathers, Second Series；Vol. 8.；Edited by Philip Schaff and Henry Wace.；Buffalo, NY: Christian Literature Publishing Co.,；1895.；<http://www.newadvent.org/fathers/3202181.htm>.}

% structure: p0001 source=h1 target=section reason=page-title

\section{\WorkTitle{第一八二号书信}}

\Emph{\Person{圣巴西略}}\par

\Emph{致\Place{萨莫萨塔}的司铎}。\par

虽然我为教会的荒凉感到悲伤，但我仍要恭贺你们这么快就被带到了你们艰苦斗争的极限。愿天主赐予你们耐心度过这一关，好使你们因忠信的管理，以及你们为基督的事业所表现出的高贵坚毅，而获得巨大的赏报。\par

\SourceRecord{Translated by Blomfield Jackson.；Nicene and Post-Nicene Fathers, Second Series；Vol. 8.；Edited by Philip Schaff and Henry Wace.；Buffalo, NY: Christian Literature Publishing Co.,；1895.；<http://www.newadvent.org/fathers/3202182.htm>.}

% structure: p0001 source=h1 target=section reason=page-title

\section{\WorkTitle{第一八三篇书信}}

\Person{圣巴西略}\par

\Emph{致\Place{萨摩萨塔}元老院。}\par

鉴于我目睹诱惑现已遍布普世，叙利亚的大城市也与你们同样遭受了相同的磨难，（尽管，诚然，没有任何地方的元老院像你们那样因善行而受称许和著称，你们因公义的热忱而闻名，）我几乎要感谢降临于你们的患难。\par

因为若非这场苦难发生，你们在考验中的证明就永远不会被人知晓。对一切为任何善行而热切努力的人而言，他们因对天主的望德而忍受的苦难，有如黄金的熔炉。\par

那么，最杰出的诸位先生，请奋起，使你们即将承担的辛劳不辜负你们已忍受的那些，使你们在一个坚固的基础上被看到加上一个更相称的完成。请奋起，使你们能围在教会的牧者周围，当主恩准他出现在自己的宝座上时，依次告诉你们每一个人：一些为天主的教会而行的善行。在主的大日，每个人将按自己辛劳的比例，从慷慨的主那里获得自己的酬报。请记念我并尽可能经常写信给我，回复我便是行公义，并且也会给我极大的喜乐，因为用文字传递一个声音的明显标记，这声音是我乐意听到的。\par

\SourceRecord{Translated by Blomfield Jackson.；Nicene and Post-Nicene Fathers, Second Series；Vol. 8.；Edited by Philip Schaff and Henry Wace.；Buffalo, NY: Christian Literature Publishing Co.,；1895.；<http://www.newadvent.org/fathers/3202183.htm>.}

% structure: p0001 source=h1 target=section reason=page-title

\section{\WorkTitle{第一八四封信}}

\Person{凯撒勒雅的圣巴西略}\par

致\Place{希默里亚}主教\Person{欧斯塔修斯}。\par

我知道，孤寂令人非常凄苦，且带来许多劳苦，因为它使我们失去了那些管理我们的人。因此我推断，你之所以不写信给我，是因为你对你所遭遇的事感到沮丧，同时你现在正忙于巡视基督的羊群，因为羊群四面受敌。然而，每一种忧伤都能在与同情者的交流中找到安慰。那么，我恳求你，尽可能多地写信给我。你同我说话，会使你自己得到振作，而我也因收到你的信得到安慰。我将在我的工作允许下，尽力同样待你。你自己要祈祷，并恳切地求全体弟兄向主恳求，求祂赐我们有一天能从目前的困苦中解脱。\par

\SourceRecord{Translated by Blomfield Jackson.；Nicene and Post-Nicene Fathers, Second Series；Vol. 8.；Edited by Philip Schaff and Henry Wace.；Buffalo, NY: Christian Literature Publishing Co.,；1895.；<http://www.newadvent.org/fathers/3202184.htm>.}

% structure: p0001 source=h1 target=section reason=page-title

\section{\WorkTitle{书信185}}

\Emph{\Place{凯撒勒雅}的\Person{圣巴西略}}\par

\Emph{致\Place{贝雷亚}主教\Person{泰奥多图斯}。}\par

虽然你不写信给我，但我知道你心中记念我；我推断这一点，并非因为我值得任何好的纪念，而是因为你的灵魂充满丰富的爱。然而，尽你所能，利用一切机会写信给我，使我既欢喜听到你的消息，又有机会向你传达我的消息。这是居住遥远的人唯一的交流方式。不要让我们彼此剥夺这种机会，只要我们的劳作允许。但我祈求天主，使我们能亲自见面，使我们的爱增长，并对我们的主因祂更大的恩赐倍增感恩。\par

\SourceRecord{Translated by Blomfield Jackson.；Nicene and Post-Nicene Fathers, Second Series；Vol. 8.；Edited by Philip Schaff and Henry Wace.；Buffalo, NY: Christian Literature Publishing Co.,；1895.；<http://www.newadvent.org/fathers/3202185.htm>.}

% structure: p0001 source=h1 target=section reason=page-title

\section{\WorkTitle{书信一八六}}

\Emph{\Person{圣巴西略·凯撒勒雅}}\par

\Emph{致\Person{安提帕特}总督}。\par

哲学是极好的事物，仅凭这一点就够了：它甚至能以微小的代价治愈其门徒；因为在哲学中，同一事物既是美味又是健康的食物。我听说您藉着腌白菜恢复了衰退的胃口。过去我厌恶它，既因为那句谚语，也因为它让我想起与之相伴的贫穷。如今，我却被逼改变想法。我嘲笑那句谚语，当我看到白菜竟是如此\Quote{养育男子汉的好慈母}，并且使我们总督恢复了青春活力。将来我必定认为没有任何东西比得上白菜，甚至\Person{荷马}的莲花，甚至那神食（无论它是什么），那养育了\Person{奥林波斯诸神}的神食，都比不上。\par

\SourceRecord{Translated by Blomfield Jackson.；Nicene and Post-Nicene Fathers, Second Series；Vol. 8.；Edited by Philip Schaff and Henry Wace.；Buffalo, NY: Christian Literature Publishing Co.,；1895.；<http://www.newadvent.org/fathers/3202186.htm>.}

% structure: p0001 source=h1 target=section reason=page-title

\section{第一八七书}

\Emph{\Place{凯撒勒雅}的\Person{圣巴西略}}\par

\Emph{\Person{安提帕特}致\Person{巴西略}}。\par

\Quote{两次白菜便是死，}刻薄的谚语说。而我，虽常呼唤它，却将死一次。是的：即使我从未呼唤过它！若你无论如何都要死，就不要害怕品尝一道美味佳肴，尽管被那谚语不公地诋毁！\par

\SourceRecord{Translated by Blomfield Jackson.；Nicene and Post-Nicene Fathers, Second Series；Vol. 8.；Edited by Philip Schaff and Henry Wace.；Buffalo, NY: Christian Literature Publishing Co.,；1895.；<http://www.newadvent.org/fathers/3202187.htm>.}

% structure: p0001 source=h1 target=section reason=page-title

\section{第一八八封信}

\Emph{\Person{圣巴西略（凯撒勒雅）}}\par

（第一法典书函）\par

\Emph{致\Person{安菲洛基}，论法典}。\par

经上记载：\Quote{连愚人}，据说，\Quote{当他提问时}，就被视为有智慧。但智者提问，连愚人也变得智慧。感恩天主，每当我收到你勤勉的来信时，正是如此。因为只是通过提问，我们就比从前更有学识、更有智慧，只因我们学到许多先前不知道的事；而急于回答它们，又反过来教导了我们。的确，在目前，尽管我此前从未留意过你所提出的问题，但我已被迫进行精确的探究，并在心中重温我从长老们那里听到的一切，以及所有与他们的教诲相符的道理。\par

一、关于你询问的\OriginalTerm{卡塔利派}{Cathari}，先前已有所陈述，你也适当地提醒我，应当遵循各地既有的惯例，因为当时就此问题作出决定的人，对卡塔利派的圣洗持有不同意见。但我认为\OriginalTerm{佩普齐派}{Pepuzeni}的圣洗绝无权威；我深感惊奇，精通法典的\Person{迪约尼修}竟未觉察此事。古代权威决定，只接受与信仰毫无偏离的圣洗。因此，他们区分了异端、分裂和非法集会等名称。所谓\Emph{异端}，是指那些在有关真正信仰的事上完全断绝关系、疏离的人；所谓\Emph{分裂}，是指那些因某些教会原因和可相互解决的问题而分离的人；所谓\Emph{非法集会}，是指由不守规矩的司铎或主教，或无知的平民主持的聚会。例如，若有人被定罪，被禁止履行圣职，却拒绝服从法典，擅自篡取主教和圣职的权利，而有人离开天主教会加入他，这就是非法集会。与教会成员在悔改问题上意见不同，是\Emph{分裂}。\Emph{异端}的例子有\OriginalTerm{摩尼派}{Manichæans}、\OriginalTerm{瓦伦廷派}{Valentinians}、\OriginalTerm{马西昂派}{Marcionites}以及这些佩普齐派；因为他们的分歧直接涉及对天主的真正信仰。因此，古代权威认为，应完全拒绝异端者的圣洗，但接受分裂者的圣洗，理由是后者仍属于教会。\par

至于那些参加非法集会的人，他们的决定是，经过适当的悔改和责备，使他们改善后，再重新接纳他们加入教会；因此，在许多情况下，当有圣职的人与不守规矩的人一同反叛时，在他们悔改后，恢复同等的品级。如今佩普齐派显然是异端，因为他们非法并可耻地将\OriginalTerm{蒙丹}{Montanus}和\OriginalTerm{普里西拉}{Priscilla}称为\OriginalTerm{护慰者}{Paraclete}，从而亵渎了圣神。因此，他们因将神性归于人类而应受谴责；因将圣神与人相比而侮辱圣神。他们因此也当受永罚，因为亵渎圣神是不得赦免的。那么，接受那些因圣父、圣子及蒙丹或普里西拉之名施洗的人的圣洗，有何理由呢？因为凡不是因我们所领受的名受圣洗的，根本未受圣洗。所以，尽管此事逃过了伟大的\Person{迪约尼修}的注意，我们决不可效法他的错误。这种立场的谬误一目了然，凡是稍有推理能力的人都看得清楚。\par

卡塔利派是分裂者；但古代权威，我是指\Person{西彼廉}和我们的\Person{斐米利安}，认为应以同一谴责拒绝所有这些派别——卡塔利派、\OriginalTerm{禁戒派}{Encratites}和\OriginalTerm{水礼派}{Hydroparastatæ}，因为分裂的起源始于分党，而那些从教会背道的人，身上已不再有圣神的恩宠，因为连续性一断，恩宠就不再赐予。最初的分离者从教父们领受了圣职，并藉按手持有神恩。但那些被分裂出去的人已成为平信徒，既然他们无法将自己已失落的圣神恩宠赐予他人，便无权付洗或授圣职。因此，凡不时受他们圣洗的人，被命令如同由平信徒付洗一般，到教会来，藉教会真正的圣洗得以洁净。然而，既然亚细亚的一些人认为，为了管理多数人的缘故，应接受他们的圣洗，那就接受吧。但我们必须认识到禁戒派的恶行；他们为了阻止自己被教会重新接纳，试图用一种独特的圣洗预先阻止日后的复和，甚至以此违背他们自己的特殊惯例。因此，我的意见是，既然对此没有明确规定，我们有责任拒绝他们的圣洗；凡从他们那里领受圣洗的人，来到教会时，我们应为他付洗。然而，如果此举可能有损普世纪律，我们就必须回归惯例，并遵循那些规定了我们当行之事的教父。因为我有些担心，我们若因决定严厉，阻止他们付洗，反而会成为得救之人的障碍。如果他们接受我们的圣洗，不要因此困扰。我们绝无义务回报他们同样的恩惠，只需严格遵从法典。无论如何，应当规定，凡从他们的圣洗来到我们这里的人，应在信友面前接受傅油，只有这样才可领受奥迹。我知道我已从禁戒派的团体中接纳了\Person{伊佐伊}和\Person{撒杜尼努}担任主教职。因此，我不能将那些与他们联合的人从教会中分离，因为我通过接受那些主教，公布了某种共融的法典。\par

二、故意毁灭自己胎儿的女子，犯了杀人罪。在我们看来，无需精细探究胚胎是否成形。在此案中，不仅是要为即将出生的生命伸张正义，也是为那攻击自己的女子辩护；因为大多数这样尝试的女子都会死亡。毁灭胚胎是一项附加的罪行，是第二次杀人，至少若我们视其为故意为之，则如此。不过，对这些女子的惩罚不应是终身的，而是为期十年。而且她们的处分不单取决于时间流逝，也取决于她们悔改的诚意。\par

三、执事在被任命为执事后犯奸淫者，应被免职。但当他被撤职并列为平信徒后，不应被排除于共融之外。因为有一条古代教规规定：凡从自己的品级堕落者，仅受此等惩罚。\par

我想，古代权威于此遵循了那条古老规则：\Quote{你不可为一事复仇两次}。还有一个进一步的理由：平信徒被逐出信友席位后，有时可恢复其原先的等级；但执事一旦承受了免职的永久惩罚，便不再恢复执事圣秩，他们只受此一罚。以上所说的是依据明文法规。但通常更真实的疗法是远离罪恶。因此，若有人为了肉欲的放纵而拒绝恩宠，后来藉着克己和自制来折磨肉体、制服它，从而放弃了曾征服他的享乐，此人将给我充分证明他得以痊愈。因此，我们应当既知道明确的规定，也知道习俗；在那些无法接受最高处理方式的案件中，应遵循传统指引。\par

四、对于\OriginalTerm{三婚}{trigamy}和\OriginalTerm{多婚}{polygamy}，他们按比例规定了与\OriginalTerm{二婚}{digamy}相同的准则：即二婚为一年（有些权威规定两年）；三婚者被隔离三至四年；但这已不再被称为婚姻，而是多婚，更确切地说，是有限度的奸淫。为此，主对曾有五个丈夫的撒玛黎雅妇人说：\BibleQuote{你现在所有的，不是你的丈夫。}他不认为那些超过第二次婚姻界限的人配得丈夫或妻子的称号。在三婚的事例上，我们接受了五年的隔离，这不是依据\OriginalTerm{教规}{canon}，而是遵循我们前任的训令。这类犯过者不应完全被剥夺教会的特权；他们应被认为在二至三年后堪当聆听，之后获准站在他们的位置；但必须被阻止领受美好的恩赐，只在显示出某种悔改的果实后才恢复共融的位置。\par

五、临终悔改的异端者应被接纳；但当然不是不加区别地接纳，而是经过考验，表明他们真实悔改，并结出果实证明他们热切渴望得救。\par

六、\OriginalTerm{正式圣职人员}{canonical persons}的奸淫不应被视为婚姻，他们的结合应被完全解除，因为这样既有利于教会的安全，也能防止异端者向我们发起攻击的借口，好像我们藉着允许犯罪的自由来吸引人加入我们。\par

七、与男人行淫者、与兽行淫者，以及杀人者、\OriginalTerm{巫师}{wizard}、奸淫者和拜偶像者，应受同样的惩罚。你在其他案例中所持的规则，也应在此类中遵守。不过，对于那些因无知而犯不洁之罪已三十年并悔改的人，我们无疑应当接纳。此案中有宽恕的理由：因无知、自愿告解以及时间漫长。也许他们曾被交与撒殚整整一个人世代的时期，好叫他们学会不再行为不端；因此，命你们立即接纳他们，特别是若他们流泪感动你的仁慈，并表现出堪受怜悯的生活方式。\par

八、盛怒之下拿起斧头攻击自己妻子的男子是杀人犯。但凭你的智慧，你理应恰当地提醒我在这点上更详细地说明，因为\OriginalTerm{故意}{intent}与非故意之间必须作出广泛区分。纯粹\OriginalTerm{非故意}{unintentional}的行为，且与行为人的目的相距甚远，例如一个人向狗或树扔石头，却打中了人。目的是赶走野兽或震落果实。偶然路过的人恰好被击中，这行为是非故意的。同样，任何人用皮带或软棍打别人以示惩戒，而被挨打者死亡，这行为也是非故意的。在此情况下必须考虑，目的不是杀人，而是改善犯过者。此外，在非故意行为中，还须算上打斗中的人，在抵挡敌人用棍棒或手的攻击时，毫不留情地击中某要害，使之受伤，虽不完全致死。然而，这已非常接近故意；因为使用这类武器自卫的人，或毫不留情地打击的人，显然因被激情控制而没有饶恕对手。同样，任何人使用重棍，或大得令人无法站立的石头，也被算作非故意，因为他没有做他所想的事：他在盛怒下打出如此一击，致受害人死亡，然而他心中所想只是揍他一顿，并非要杀死他。但是，若有人使用了剑或类似物品，则无可推诿；若他扔斧头，更是如此。因为他不是用手击打，以使打击的力度可以自己控制，而是投掷，因此由于铁器的重量和刃口以及投掷的力量，伤口必定致命。\par

另一方面，在战争或抢劫的攻击中所做的事显然是故意的，毫无疑问。强盗为了贪婪和逃避定罪而杀人。在战争中致死的士兵，其明显目的不是战斗，也不是惩罚，而是杀死对手。如果有人为了其他原因调制了某种\OriginalTerm{魔法药水}{magic philtre}，并导致死亡，我视此为故意。妇女常常试图通过\OriginalTerm{咒语}{incantations}和\OriginalTerm{魔法结}{magic knots}来吸引男人爱她们，并给他们服用使智力迟钝的药物。这样的妇女，当她们导致死亡时，尽管结果可能并非她们所愿，然而由于她们的行为涉及魔法且被禁止，应被视为故意杀人。此外，服用药物导致堕胎的妇女，以及服用毒药摧毁未出生婴儿的妇女，都是杀人犯。就此话题。\par

九. 主的话说，除了因淫乱的缘故，不可离开婚姻结合（\BibleRef{玛 5:32}），按论证，这话同样适用于男女。然而习俗并非如此。但是，关于妇女，可以找到非常严格的表述；例如，宗徒的话：\BibleQuote{与娼妓联合的，成为一体}（\BibleRef{格前 6:16}），以及耶肋米亚的话：如果妻子\BibleQuote{成为别人的妻子，他还能回到她那里吗？那地岂不是大大玷污了吗？}（\BibleRef{耶 3:1}），又说：\Quote{拥有奸妇的人是愚昧和邪恶的。}然而，习俗规定，犯奸淫和行淫乱的男子应由其妻子保留。因此，我不知道与那被休的男子一起生活的女子是否能恰当地被称为奸妇；此案的指控针对的是休夫的女子，取决于她离开婚姻的原因。如果她因被打而拒绝顺从，她最好忍受而不是与丈夫分离；如果她因经济损失而反对，即使这样她也没有充分的理由。如果她的理由是丈夫生活在淫乱中，我们在教会的习俗中没有找到这一点；但妻子不可离开不信的丈夫，而应留下，因为结局不确定。\BibleQuote{因为妻子，你哪里知道是否能救你的丈夫？}（\BibleRef{格前 7:16}）因此，妻子如果离开丈夫而跟从别人，就是奸妇。但被抛弃的男子是可原谅的，与这样的男子一起生活的女子不被定罪。但如果抛弃妻子的男子跟从别人，他本人就是奸夫，因为他使她犯了奸淫；与他一起生活的女子就是奸妇，因为她使别人的丈夫归向了她。\par

十. 那些发誓不接受圣秩、以誓言拒绝圣职的人，不可强迫他们犯假誓，尽管似乎有一条教规对此类人作出让步。然而我凭经验发现，发假誓的人从未有好结果。但必须考虑誓言的格式、条款、发誓时的心态，以及条款中最细微的附加内容，因为如果在任何地方都找不到解救的理由，这样的人必须被开除。然而，\Person{塞维鲁}的情况，我指的是由他按立的司铎，在我看来确实允许这种解救，如果您允许的话。请指示将置于\Place{梅斯蒂亚}下的区域（此人被派往该地）划归\Place{瓦索达}管辖。这样，他将因不离开原地而不违背誓言，而\Person{隆基努斯}有\Person{基里亚库斯}在身边，就不会使教会无人照管，他自己也不会犯疏忽职守之过。此外，我也不会因向\Person{基里亚库斯}作出让步而被视为违反任何教规，他曾发誓留在\Place{明达纳}，却接受了调任。他的返回将符合他的誓言，他服从安排不会被视为假誓，因为在他的誓言中没有附加他不可离开\Place{明达纳}（即使很短时间），而是将来要留在那里。我将宽恕以遗忘为借口的\Person{塞维鲁}，只告诉他，那位洞察隐秘的主不会忽视这样一个蹂躏祂教会的人：此人首先不按教规任命，然后违反福音加给人誓言，接着在调任之事上让人发假誓，最后假装遗忘而说谎。我不是内心的审判者；我只凭所闻判断；让我们将惩罚留给主，而我们自己宽恕这常见的遗忘之人为过失，并毫无质疑地接纳那人。\par

十一. 犯下\OriginalTerm{非故意杀人}{unintentional homicide}罪的人，在十一年中已做了足够的补赎。我们无疑会遵守梅瑟对受伤之人的规定：人若被打后躺下，又扶着拐杖起来行走，就不算杀人（\BibleRef{出 21:19}）。但如果被打后不再起来，然而由于没有杀人的意图，那打人的人是杀人者，但是非故意杀人。\par

十二. 教规绝对排除再婚者担任圣职。\par

十三. 战争中的杀人，我们的教父不视为杀人；我猜想，是因为他们愿意对为贞洁和真宗教而战斗的人作出让步。然而，或许最好劝告那些手不洁的人，只三年内不领圣体。\par

十四. 放高利贷的人，如果同意将不义之财用于穷人，并且今后摆脱贪婪的祸害，可以接纳进入牧职。\par

XV. 我惊讶于你们要求 \WorkTitle{圣经} 的精确性，并主张那传达原文含义的译文措辞有些牵强，未能准确转达希伯来词的意思。然而，我不可轻率放过这由探究之心提出的问题。在创造世界时，空中的飞鸟和海中的鱼类有相同的起源；因为两者都出自水。原因是两者具有相同的特性：后者在水中游，前者在空中飞。因此它们被一同提及。这种表达方式并非无区别地使用，而是对一切水中生物都恰当地运用。空中的飞鸟和海中的鱼类都受制于人；不仅是它们，还有一切在海路中穿梭的生物。因为并非所有水中生物都是鱼，例如海怪、鲸、鲨、海豚、海豹，甚至海马、海狗、锯鳐、剑鱼和海牛；此外，你若喜欢，还有海荨麻、蛤蜊及一切硬壳生物，其中没有一种是鱼，它们都在海路中穿梭；因此有三种：空中的飞鸟、海中的鱼类，以及一切不同于鱼、在海路中穿梭的水中生物。\par

XVI. \Person{纳阿曼}并非在主面前为大，而是在他主人面前为大；即他是叙利亚王的主要将领之一。\BibleRef{列下 5:1} 请仔细阅读你的 \WorkTitle{圣经}，你将在那里找到你问题的答案。\par

\SourceRecord{Translated by Blomfield Jackson.；Nicene and Post-Nicene Fathers, Second Series；Vol. 8.；Edited by Philip Schaff and Henry Wace.；Buffalo, NY: Christian Literature Publishing Co.,；1895.；<http://www.newadvent.org/fathers/3202188.htm>.}

% structure: p0001 source=h1 target=section reason=page-title

\section{第一八九封书信}

\\Person{凯撒勒雅的圣巴西略}\par

\\Emph{致医生\\Person{欧斯塔基}}\par

人道是所有从事医学之人的日常要务。依我之见，将医术置于人生追求之首，是合理而正确的决定。无论如何，若生命不能健康地活着，则痛苦且不值得活；而健康又依赖于你们的技艺。在您身上，医学仿佛有两只右手；您并未将技艺的应用局限于人的身体，而是也关注治愈他们灵魂的疾病，从而扩展了公认的仁爱界限。我这样写，不仅是依据公论，也是基于我多次亲身体验，尤其是在当前，当我们的敌人那难以言表的邪恶如毒流般泛滥我们生命时，您极巧妙地驱散了它，并以抚慰的话语浇灌，平息了我心中的炎症。鉴于敌人对我接连不断、花样百出的攻击，我以为应当保持沉默，承受他们的连续打击而不回击，也不试图反驳那些以谎言武装的敌人——那可怕的武器常常刺穿真理本身的心脏。您做得好，敦促我不要放弃为真理辩护，反而要定罪我们的诽谤者，免得谎言得逞，许多人受害。\par

我的对手们对我表现出出乎意料的仇恨，似乎是重演伊索寓言中的旧事。他让狼指控羔羊某些罪状，仿佛真的羞于在没有合理借口的情况下杀死一个无辜的动物；当羔羊轻易驳斥了诽谤时，狼依然继续攻击，虽然在公义上失败，却在撕咬中获胜。那些似乎将仇恨我视为美德的人也是如此。他们也许因无端恨我而羞愧，于是捏造借口和指控，不坚持任何一项控诉，却以时而这个、时而那个、时而另一件事作为他们憎恨的理由。他们的恶意没有一次是一致的；当一项指控受阻，他们就抓住另一项；这项失败，又求助于第三项；即使所有指控都被推翻，他们也不放弃敌意。他们说我在宣讲三个神，向民众的耳朵灌输这一指控，并貌似有理地、顽固地诽谤。然而，真理在我一方；无论在公共场合对众人，还是在私下对我遇到的每个人，我都证明我诅咒任何主张三个神的人，甚至不认为他是基督徒。他们一听到这话，撒伯里乌就顺手被用来指责我，并传出风声说我的教导沾染了他的错误。我再次以惯用的真理武器捍卫自己，证明我对撒伯里乌主义的恐惧如同对犹太教一样。\par

那么，经过所有这些努力，他们疲倦了吗？他们停止了吗？一点也没有。他们指控我创新，并将此指控建立在我承认三个位格的基础上，指责我肯定一个美善、一个能力、一个天主性。在这一点上，他们没有偏离真理，因为我确实这样肯定。他们抱怨说，他们的习惯不接受这一点，圣经也不符合。我的回答是什么？我认为他们中间流行的习惯不应被视为正统的法则和标准。如果要以习惯来证明什么是正确的，那么我同样可以提出我这里存在的习惯。如果他们拒绝，我们显然不必跟随他们。因此，让天主启示的圣经在我们之间作裁决；哪一方找到与天主之言和谐的教义，真理的票就投给那一方。那么，指控是什么？对我的指控同时提出了两点。一方面，我被指控分离位格；另一方面，我从未使用任何与天主性相关的名词的复数形式，而总是使用单数，如我所说的，一个美善、一个能力、一个天主性等等。关于分离位格，那些在神圣本性中主张本质区别的人不应反对或抗拒。因为主张三个本质却反对三个位格是不合理的。因此，剩下的唯一指控就是我在谈论天主性时使用单数形式。\par

我对付第二项指控有些困难。谁若谴责那些主张天主性为一的人，就必然要么同意那些主张多神的人，要么同意那些主张无神的人。没有第三种可能。受默感的圣经教导不允许我们说多个天主性，但凡提到天主性，都是用单数；例如：\\BibleQuote{因为天主性的一切丰满，有形有体地居住在他内。}\\BibleRef{哥 2:9}又如：\\BibleQuote{其实，自从天主创世以来，他那看不见的美善，即他永远的大能和他为天主的天主性，都可凭他所造的万物，辨认洞察出来。}\\BibleRef{罗 1:20}因此，如果使天主性增多是多神错误受害者的特殊标志，而完全否认天主性则是陷入无神论，那么我承认一个天主性，这指控有何意义？但他们更清楚地暴露了他们的目的：即在父方面承认他是天主，同样同意子也应以天主性受到尊崇，但拒绝将圣神包含在与父和子同一的天主性观念中。他们允许天主性的能力从父延伸到子，却将圣神的本性与神圣荣耀分开。针对这一观点，我必须尽我所能，对自己的立场做一简短的辩护。\par

5. 那么，我的论点是什么？主在将得救的信德传授给那些在祂的教导中成为门徒的人时，连同圣父和圣子也一起包含了圣神。我认为，一旦被包含，就处处时时都被包含。这里不存在在某些方面一起排序、在其他方面却分离的问题。在使我们的本性从腐朽的生命转化为不朽的生命力的德能中，圣神的德能是与圣父和圣子一同被包含的；并且在许多其他事例中，例如在善、圣、永恒、智慧、正义、至高、有效等概念中，以及一般地在所有具有更高含义的术语中，祂都是不可分割地结合的。因此，我认为，那在如此多崇高和神圣的意义上与圣父和圣子结合的圣神，是永远不会分离的。事实上，我不知道在关于天主性的术语中有任何更好或更差的等级，也无法想象，允许圣神参与那些较低尊荣的属性，而认为祂不配享有更高尊荣，会是虔诚和正确的。因为所有关于天主的概念和术语彼此之间具有同等的尊荣，因为它们在其所指的主题意义上并无不同。我们的思想并非因善的归属而引向一个主体，因智慧、大能、正义的归属而引向另一个主体；无论你提到哪个属性，所指的事物是同一个。而且，如果你称\Quote{天主}，你指的是你通过其他术语所理解的同一存在。因此，既然所有用于天主性的术语在其所描述的对象方面彼此具有同等效力，一个强调这一点，另一个强调那一点，但都使我们的理智专注于同一对象的默观；那么，有什么理由允许圣神在所有其他术语上与圣父和圣子共融，却唯独在天主性上孤立祂？我们无法逃避这一立场：要么在此处允许共融，要么处处都拒绝。如果祂在其他每方面都配得上，那么祂在此方面当然也不会不配。如果，正如我们的对手所争论的，祂过于渺小，不能被允许与圣父和圣子在天主性中共融，那么祂也不配分享任何一件神圣属性；因为当这些术语被仔细考虑并相互比较时，借助每个术语所特有的意义，人们会发现它们所涉及的丝毫不亚于\Quote{天主}这个称号。我所说的一个证明在于这样一个事实：甚至许多低等对象也被赋予这个名称。不仅如此，圣经甚至毫不避讳地将这个术语用于性质完全相反的事物，例如将\Emph{神}这个称号用于偶像。如经上所载：\Quote{愿那些没有创造天地的神，被除掉，被丢到地底下；}又说：\Quote{外邦人的神是偶像。}而那位女巫，当她为撒乌耳召来所需的亡魂时，据说看见了神。\BibleRef{撒上 28:13} 巴郎，一个占卜者和先见者，手持神谕，如圣经所说，当他以他神圣的巧计获得了魔鬼的教训时，圣经描述他与天主商议。从圣经中许多类似的事例可以证明，\Quote{天主}这个名称并不比其他用于神圣事物的词汇更有优先权，因为，如前所述，我们发现它甚至不加区别地被用于性质完全相反的事物。另一方面，圣经教导我们，圣、不朽、正义、善这些名称从未被不加区别地用于不值得的对象。因此，如果他们不否认圣神在那些真正宗教用法中专门用于天主性本身的名称上与圣子和圣父相关联，那么就没有合理的理由拒绝在那一个词上允许同样的关联，而这个词，如我已指出的，甚至被公认是魔鬼和偶像的\OriginalTerm{同音异义词}{homonym}。\par

6. 但他们争辩说，这个称号表明了它所应用的事物的本性；圣神的本性不是与圣父和圣子共享的本性；因此，圣神不应被允许共同使用这个名称。因此，他们必须指出是通过什么途径察觉出本性的这种差异。如果确实有可能直接默观天主性本身；如果可以通过可见事物发现什么是属于它的，什么是外在于它的；那么我们当然不需要用言语或其他标记来引导我们理解探究的对象。但天主性过于崇高，无法像探究对象那样被感知，而对于超出我们知识的事物，我们只能根据或然证据进行推理。因此，我们在探究天主性时必然以其活动为指导。假设我们观察到圣父、圣子、圣神的活动彼此不同，那么我们就会从活动的多样性推断出活动主体的本性也是不同的。因为不可能在本质上不同的东西在活动的形式上彼此关联；火不会使水结冰；冰不会使水变暖；本性的差异意味着源于它们的活动的差异。因此，假设我们察觉圣父、圣子和圣神的活动是同一个，毫无差异或变化；那么从这种活动的同一性，我们必然推断出本性的合一。\par

7. 圣父、圣子和圣神一同圣化、赋生、启迪和安慰。任何人听到福音中救主向圣父论及祂的门徒说\Quote{在祢的名内圣化他们}之后，都不会将一种特殊而独特的圣化行动归于圣神的运作。同样地，所有其他行动都由圣父、圣子和圣神平等地在所有堪当领受的人身上实行：每一恩宠和德性、引导、生命、安慰、变化为不朽、通往自由的过渡，以及降临于人的一切其他美善。甚至那超越我们、关于受造界（无论理智或感官方面的）的安排，如果确实有可能从我们所知的事对超越我们的事形成任何推测的话，也不是离开圣神的行动和德能而构成的；每个个体都按各自的尊严和需要分沾祂的帮助。确实，对于超越我们本性之存在的秩序与治理，我们的认知是模糊的；然而，任何人从已知事物推论，会更合理地认为圣神的德能甚至在这些存在中运作，而非认为祂被排除在超世上物的治理之外。如此断言乃是提出一种赤裸裸且无根据的亵渎；它是用谬误支持荒谬。另一方面，同意甚至那超越我们的世界也是由圣神的德能以及圣父和圣子的德能所治理，则是提出一种基于人类生活中所见之明显证实的论点。圣父、圣子和圣神的运作的同一性清楚地证明了本性的不变性。因此，即使天主性这个名称确实表示本性，本质的共通性也证明这个称号非常适当地应用于圣神。\par

8. 然而，我无法理解我们的对手何以能以其全部才智，用\OriginalTerm{天主性}{Godhead}这个称号来证明本性，好像他们从未从圣经中听说过本性并非出于制度和任命。当神圣的声音说\BibleQuote{看，我已立你为法郎的天主}时，梅瑟被立为埃及人的天主。\BibleRef{出 7:1}因此，这个称号确实证明了某种监督或行动的权柄。另一方面，神圣本性在所有被造出的词语中始终是无法解释的，正如我常教导的。我们了解到它是慈善的、审判的、正义的、良善的等等，因而学到了行动的差异。然而，我们仍无法通过行动的概念来理解行动者的本性。请任何人解释这些名称中的每一个，以及它们所应用的实际本性，就会发现两个定义并不相同。哪里定义不相同，本性就不同。因此，必须区分本质（至今尚未发现任何解释性术语）与名称按其行动或尊严所应用的意义。从术语的共同性我们推断行动没有差异。但是，我们并没有得出本性变化的明确证据，因为如前所述，行动的同一性表明本性的共同性。因此，如果天主性是行动的名称，我们说天主性是唯一的，因为圣父、圣子和圣神的行动是唯一的；然而，如果如同普遍所认为的，天主性这个名称表示本性，那么既然我们发现本性没有变化，我们合理地定义天主圣三为同一的天主性。\par

\SourceRecord{Translated by Blomfield Jackson.；Nicene and Post-Nicene Fathers, Second Series；Vol. 8.；Edited by Philip Schaff and Henry Wace.；Buffalo, NY: Christian Literature Publishing Co.,；1895.；<http://www.newadvent.org/fathers/3202189.htm>.}

% structure: p0001 source=h1 target=section reason=page-title

\section{书信第一九〇号}

\Emph{圣巴西略}\par

\Emph{致依科尼雍主教\Person{安斐罗西}}。\par

1. 你对伊索利亚教会事务所表现出的关切，正是你那不断令我钦佩的热诚与合宜行为的自然流露。最不经意的观察者也必能立即看出，将忧虑与关注分散于多位主教在各方面都更为有利。此事未能逃过你的观察，你考察并告知我事态的进程，做得很好。然而，找到合适的人选并非易事。因此，当我们渴望获得人数带来的声誉，并藉由更多职员使天主的教会得到更有效的管理时，我们务必小心，以免因被召任职者的品格不称而无意中使圣言遭受轻视，并惯坏平信徒的冷淡。你自己深知，被治理者的行为通常与治理者如出一辙。因此，也许更好的办法是任命一位经过充分认可的人（尽管这也不容易）来监督整个城市，并委托他自行负责管理细节。只要他是一位天主的仆人，\BibleQuote{无需羞愧的工人，}\BibleRef{弟后 2:15}不是\BibleQuote{只求自己的事，}\BibleRef{斐 2:4}而是求多数人的事，\BibleQuote{使之得救。}\BibleRef{得前 2:16}若他感到责任过重，他会为自己联合其他收割的工人。只要我们能找到这样的人，我承认我认为一人可抵多人，并且以这种方式安排灵魂的牧养，对教会更有益处，对我们风险更小。然而，若此路困难，让我们首先尽力为那些自古即为主教驻地的小城镇或村庄任命监督者。然后我们将再次任命城市的主教。若不采取此步骤，被任命者可能会妨碍后续的管理，因他希望统治更大的教区并拒绝接受其他主教的祝圣，我们可能突然陷入内部纷争。若此路困难且时间不允许，请确保通过祝圣一些邻近主教来严格限定伊索利亚主教在其界限内。将来，待我们仔细审查了自己认为最合适的人选之后，再于合宜的季节为其余地区任命主教。\par

2. 我已按你的请求询问了\Person{乔治}。他的回答如你所报告。在这一切事上，我们必须保持安静，将家务交托于主。因为我信赖圣洁的天主，祂必藉我的协助，以某种方式赐予他脱离困境，并赐我无忧无虑地生活。若不能如此，请惠然告知我需要关注哪些部分，好让我开始向每位有权势的朋友请求这份恩惠，或分文不取，或以适度代价，按主所赐的顺利。\par

我已按你的请求写信给\Person{瓦肋里}弟兄。尼撒的事务正如阁下所留下的那样进展，并藉你的圣德而不断改善。那些当时与我分离的人中，有些已前往宫廷，有些留下等待消息。主既能挫败后者的期望，也能使前者的归来归于无用。\par

3. \Person{斐罗}依据某些犹太传统解释，玛纳的性质随食者的口味而变化：本身像蜜煮小米粒，有时当作面包，有时当作鸟兽之肉，有时当作蔬菜，各随所好；甚至当作鱼，以致每种味道都准确地在食者口中重现。\par

圣经承认有三人乘坐的战车，因为其他战车只有两人（驭手和武士），而法郎的战车有三人：两名武士和一名持缰者。\par

\Person{辛皮乌斯}写信给我，表达了尊敬与共融。我将回复的信件寄给阁下，若您完全赞同，请加上您自己的问候转交给他。愿您藉圣者的慈爱，为我和天主的教会得以保全，身体健康，在主内喜乐，并常为我祈祷。\par

\SourceRecord{Translated by Blomfield Jackson.；Nicene and Post-Nicene Fathers, Second Series；Vol. 8.；Edited by Philip Schaff and Henry Wace.；Buffalo, NY: Christian Literature Publishing Co.,；1895.；<http://www.newadvent.org/fathers/3202190.htm>.}

% structure: p0001 source=h1 target=section reason=page-title

\section{第一九一封书信}

\Emph{圣巴西略·\Place{凯撒勒亚}}\par

\Emph{致依科尼雍主教\Person{安斐洛希}}。\par

阅读尊驾的书信，我衷心感谢天主。因为我在您的措辞中发现了古时情谊的痕迹。您不像大多数人，没有坚持拒绝开始一封充满情谊的通信。您已学到了为谦卑所应许给圣人们的巨大赏报，所以您选择居于第二位，从而走在我前面。在基督徒中，这就是得胜的条件；是那甘居第二位的人赢得荣冠。但我不可在这场美德的竞赛中落后，因此我以此回信问候尊驾；并向您告知我的心意：既然在我们中间信仰一致已确立，再没有甚么能阻止我们成为一身一灵，正如我们蒙召于同一盼望。那么，凭您的爱德，请继续这良好的开端，召集志同道合的人站在您身边，并指定会面的时间和地点。这样，借天主的恩宠，透过相互的迁就，我们得以用古老的爱德治理教会：接纳从对面而来的弟兄们为我们的成员，如同派遣给我们亲属的人，又如同从我们自己亲属处接纳。这的确是教会昔日的荣耀。从各教会来的弟兄们，从世界一端旅行到另一端，携带着小小的凭证，发现所有人都是父兄。这是基督的教会之仇敌所剥夺我们的特恩，正如其余一切一样。我们各自被困在自己的城里，每个人都以不信任的目光看着邻人。还有甚么可说的呢？唯有我们的爱德已冷淡了（\BibleRef{玛 24:12}），而我们的主曾告诉我们，唯独借此，祂的门徒被辨认出来（\BibleRef{若 13:35}）。首先，若您愿意，请彼此认识，好让我知道应与谁达成一致。如此，我们可共同商定一个对双方都便利的地点，并在适宜旅行的季节，加速彼此会面；主将指引我们的道路。再会。愿你心安。请为我祈祷。愿我因圣者的恩宠而获得您？\par

\SourceRecord{Translated by Blomfield Jackson.；Nicene and Post-Nicene Fathers, Second Series；Vol. 8.；Edited by Philip Schaff and Henry Wace.；Buffalo, NY: Christian Literature Publishing Co.,；1895.；<http://www.newadvent.org/fathers/3202191.htm>.}

% structure: p0001 source=h1 target=section reason=page-title

\section{\WorkTitle{书函一九二}}

\Emph{\Person{圣巴西略·凯撒勒雅}}\par

\Emph{致尚书索夫罗纽}\par

以您行善的非凡热忱，您写信给我，说您自己欠我双倍的感谢：第一，因收到我的信；第二，因我为您效劳。那么，我岂不更该感谢您吗？既因读到您那令人至为愉悦的文字，又因我所寄望之事如此迅速地实现！那消息本身已极为可喜，但因您是我所欠恩惠的朋友，更使我倍加欣慰。愿天主恩准，不久我能见到您，亲自向您道谢，并享受与您交往的极大快乐。\par

\SourceRecord{Translated by Blomfield Jackson.；Nicene and Post-Nicene Fathers, Second Series；Vol. 8.；Edited by Philip Schaff and Henry Wace.；Buffalo, NY: Christian Literature Publishing Co.,；1895.；<http://www.newadvent.org/fathers/3202192.htm>.}

% structure: p0001 source=h1 target=section reason=page-title

\section{第一九三篇书信}

\Emph{\Person{圣巴西略}（\Place{凯撒勒雅}）}\par

\Emph{致医师\Person{梅肋提乌斯}}。\par

我无法像鹤那样逃避冬季的不适，尽管在预见未来方面，我与鹤一样聪明。但在生活的自由上，这些鸟几乎比我还超前，正如它们在飞翔能力上超过我一样。首先，我被一些世俗的事务所耽搁；然后，持续而猛烈的发热发作使我如此虚弱，以至于我甚至比以往更瘦——我比任何时候都瘦。除此以外，间日疟的发作已经持续了二十多次。现在，我似乎已经不发烧了，但我如此虚弱，连蜘蛛网都不如。因此，最短的旅程对我也太远了，每一阵风对我来说都比大海中的巨浪更危险。我别无选择，只能躲在我的小屋里等待春天，只要我能撑到那个时候，并且没有被我从未摆脱的内在疾病提前带走。如果主用祂大能的手拯救我，我将欣然前往你那偏远的地方，欣然拥抱如此亲爱的朋友。只求你为我祈祷，使我的生命能按最有益于我灵魂的福祉来安排。\par

\SourceRecord{Translated by Blomfield Jackson.；Nicene and Post-Nicene Fathers, Second Series；Vol. 8.；Edited by Philip Schaff and Henry Wace.；Buffalo, NY: Christian Literature Publishing Co.,；1895.；<http://www.newadvent.org/fathers/3202193.htm>.}

% structure: p0001 source=h1 target=section reason=page-title

\section{\WorkTitle{书信194}}

\Emph{\Person{凯撒勒雅的圣巴西略}}\par

\Emph{致\Person{佐依禄}}。\par

最卓越的先生，您在谦卑上先知先觉，这是怎么回事？您既有教养，又能写出您所寄来的那样一封书信，却仍在我面前请求宽恕，仿佛您从事了什么鲁莽越份之举似的。但玩笑归玩笑。请随时写信给我。难道我完全是文盲吗？阅读一位雄辩家的信函是令人愉快的。难道我从圣经学不到爱是何等美善的事吗？我视与一位仁爱的朋友交往为无价之宝。而且我确实希望您能告诉我我为您祈求的一切美好恩赐：最健康的身体，以及您全家的一切福乐。至于我自己的情况，我的状况并不比往常更可忍受。只告诉您这一点就够了，您就会明白我健康状况的恶劣。它确实已达到了如此极端的痛苦，以至于既难以形容，也难以体验——如果您的体验不及我的话。然而，这是良善的天主的工作，赐给我力量，使我得以忍耐承受为我的益处而由仁慈的主所加给我的一切试探。\par

\SourceRecord{Translated by Blomfield Jackson.；Nicene and Post-Nicene Fathers, Second Series；Vol. 8.；Edited by Philip Schaff and Henry Wace.；Buffalo, NY: Christian Literature Publishing Co.,；1895.；<http://www.newadvent.org/fathers/3202194.htm>.}

% structure: p0001 source=h1 target=section reason=page-title

\section{书信第一九五}

\Emph{\Person{圣巴西略（该撒利亚的）}}\par

\Emph{致\Person{厄弗若纽}，\Place{亚美尼亚的科洛尼亚}主教。}\par

\Place{科洛尼亚}，就是主置于你权下的那个地方，远离通常的路径。结果是，尽管我经常给\Place{小亚美尼亚}其余弟兄写信，我却犹豫给阁下写信，因为我不期望能找到任何人传递我的信。现在，既然我盼望要么你亲临，要么我的信会由我曾写信的一些主教转交给你，我这样就写信向你致意。我想告诉你，我似乎还活着，同时劝你为我祈祷，愿主减轻我的痛苦，并从我心中卸下像云一样压着的沉重痛楚。如果主只愿迅速恢复那些因忠于真信仰而被迫流散各地的虔诚主教们，我将得到这种缓解。\par

\SourceRecord{Translated by Blomfield Jackson.；Nicene and Post-Nicene Fathers, Second Series；Vol. 8.；Edited by Philip Schaff and Henry Wace.；Buffalo, NY: Christian Literature Publishing Co.,；1895.；<http://www.newadvent.org/fathers/3202195.htm>.}

% structure: p0001 source=h1 target=section reason=page-title

\section{书信 196}

\Emph{\Person{圣巴西略} \Place{凯撒勒雅}}\par

\Emph{致\Person{阿布尔吉乌}}.\par

传闻，这好消息的使者，不停地报告你如何如同星辰，时而在这里，时而在那里，出现在蛮族地区；时而为军队供给粮草，时而以华丽的仪仗出现在皇帝面前。我祈求天主，使你所行之事如其所应得地亨通，并使你获得卓越的成功。我祈求，在我还活着并呼吸这空气的日子里（因为我的生命如今不过是喘息），我们的祖国能时常得见你。\par

\SourceRecord{Translated by Blomfield Jackson.；Nicene and Post-Nicene Fathers, Second Series；Vol. 8.；Edited by Philip Schaff and Henry Wace.；Buffalo, NY: Christian Literature Publishing Co.,；1895.；<http://www.newadvent.org/fathers/3202196.htm>.}

% structure: p0001 source=h1 target=section reason=page-title

\section{书信第一九七}

\Emph{\Person{凯撒勒雅的圣巴西略}}\par

\Emph{致\Person{盎博罗削}，\Place{米兰}主教}。\par

1. 主的恩赐永远伟大而繁多；其伟大不可测量，其数目不可计算。对于那些不麻木于祂仁慈的人，这些恩赐中最大的之一就是我现在所利用的——虽然我们相隔遥远，却能通过书信彼此交谈。祂赐予我们两种相识的方式：一是亲身交往，二是书信往来。现在，我通过你的言语认识了你。我并非指我的记忆留下了你外在的面貌，而是你话语的丰富多样使我体会到了你内在之人的美，因为我们都 \Quote{口里说的，心里充满的。}\BibleRef{玛 12:34} 我光荣了天主，祂在每一世代拣选祂所喜悦的人；祂昔日从羊圈中拣选了一位王子作祂子民的领袖；祂藉着圣神赐给牧羊人亚毛斯能力，提升他做先知；如今，祂从帝国城市中选召了一个人来照顾基督的羊群，委托他治理整个民族，他在品格、家世、地位、口才以及这个世界所赞赏的一切上都出类拔萃。这个人抛弃了世上的一切利益，视之为损失，为要赢得基督，\BibleRef{斐 3:8} 他亲手执掌了那艘伟大的、因信德而闻名于天主面前的船——基督的教会。那么，来吧，天主的人！你所领受的基督福音，不是从人来的，也不是人所教导的，而是主亲自将你从地上的审判官席位转移到了宗徒的宝座上；打那美好的仗；医治百姓的软弱，若有因亚流狂妄之病而感染的人；重振教父们古老的足迹。你已为我奠定了情谊的基础；愿你通过频繁的问候努力在这基础上建造。这样，我们虽在地上居所相隔遥远，却能在精神上彼此亲近。\par

2. 因你在真福主教狄约尼削一事上的热忱与热诚，你证明了你对主的一切爱德、对前人的尊敬，以及对信德的热诚。因为我们对待忠信同仆的态度，是归于他们所侍奉的主。凡尊敬为信德争战之人者，证明他也有同样的热诚。一个行为就证明了许多美德。\par

我想告知你在基督内的爱德：由你的尊位委派在这件善工中为你效劳的那群非常热忱的弟兄，因他们态度的和蔼赢得了全体圣职人员的称赞；因为他们个人的谦逊与温和，展现了所有人的健全状态。此外，他们以全副热诚和勤勉，勇敢面对了严酷的季节；并以不懈的坚持，说服了那位真福者遗体的忠实守护者，将他们视为生命保障之物的保管权移交给他们。你应当明白，若不是这些弟兄的坚持促使他们遵从，他们是绝不会被任何人世权威或主权强迫的。无疑，对我们深爱的、可敬的司铎儿子德拉西乌斯的在场，对达成所愿的目标大有帮助。他自愿承担了旅途的全部辛劳；他缓和了当地信友的热忱；他以论点说服了反对者；在司铎、执事以及许多敬畏主的人面前，他以一切应有的敬意取起了圣髑，并协助弟兄们保存之。请你以与那些守护者交出并送来给你时的痛苦相等的喜悦来接收这些圣髑。任何人都不要争论，不要怀疑。这里你有那位不可征服的战士。这些骨骼曾与真福灵魂一同战斗，为主所认识。在主公正报应的日子里，祂要将这些骨骼与那灵魂一同加冕，正如经上所载：\BibleQuote{我们都要站在基督的审判台前，好使每人按他肉身所行的，或善或恶，领取相当的报应。}一口棺材安放了那受尊敬的尸体。没有别的尸体躺在他旁边。葬礼是隆重的；人们向他行了殉道者的敬礼。那些曾接待他为客人、并亲手将他安葬的基督徒，现在将他启了出来。他们如同失去了父亲和勇士的人一样哭泣。但他们把他送到你这里，因为他们把你的喜乐置于自己的安慰之上。给予的手是虔诚的；接受的手是审慎小心的。没有欺骗的余地，没有狡诈的余地。我对此作证。愿纯洁的真理为你们所接受。\par

\SourceRecord{Translated by Blomfield Jackson.；Nicene and Post-Nicene Fathers, Second Series；Vol. 8.；Edited by Philip Schaff and Henry Wace.；Buffalo, NY: Christian Literature Publishing Co.,；1895.；<http://www.newadvent.org/fathers/3202197.htm>.}

% structure: p0001 source=h1 target=section reason=page-title

\section{第一九八封书函}

\Emph{\Person{凯撒勒雅的圣巴西略}}\par

\Emph{致撒摩撒塔主教\Person{欧瑟伯}}。\par

在官员们带给我的那封信之后，我又收到了另一封后来寄出的信。我自己没有多寄信，因为我找不到往您那里去的旅人。但我寄出的信不止四封，其中也包括在您圣德的第一次书信之后从撒摩撒塔转交给我的那些信。我将这些信密封好，寄给我们可敬的兄弟、尼西亚的\OriginalTerm{税务官}{peraequator}\Person{莱昂提乌斯}，敦促他设法将信转交给我们可敬的兄弟\Person{索弗洛尼乌斯}的管家，由他安排送到您处。由于我的信要经多人之手，很可能因为某人太忙或太大意，导致您尊驾永远收不到。那么，我恳求您原谅我信写得少。以您平时的智慧，您已经恰当地责备了我，说我没有在必要时按应有的方式派出自己的信使；但您要知道，我们经历了一个如此严酷的冬天，所有道路都封锁到复活节，而且没有人愿意冒旅途的艰难。因为虽然我们的神职人员看起来人数众多，但他们缺乏旅行经验，从不经商，也宁愿不远行，大多数人从事坐着的手艺来维持每日生计。我现在派到您尊驾那里的这位兄弟，是我从乡间召来的，派他送信给您圣德，以便他能向您清楚报告我和我的情况，并且，靠天主的恩宠，也能给我带回关于您和您情况的清楚而迅速的消息。我们亲爱的兄弟、读经员\Person{欧瑟伯}早已渴望赶往您圣德那里，但我为了天气好转而把他留在这里。即使现在，我也颇感焦虑，怕他旅行经验不足会给他带来麻烦，引发疾病；因为他身体并不强壮。\par

2. 关于东方的革新，我无需在信中多言，因为兄弟们可以亲自向您准确报告。您要知道，我尊敬的朋友，当我写这些话时，我已病得失去了对生命的一切希望。我无法一一列举我所有的痛苦症状——我的虚弱、发烧的猛烈、以及整体糟糕的健康状况。只举一点：我如今在今生痛苦而悲惨的岁月已到尽头。\par

\SourceRecord{Translated by Blomfield Jackson.；Nicene and Post-Nicene Fathers, Second Series；Vol. 8.；Edited by Philip Schaff and Henry Wace.；Buffalo, NY: Christian Literature Publishing Co.,；1895.；<http://www.newadvent.org/fathers/3202198.htm>.}

% structure: p0001 source=h1 target=section reason=page-title

\section{\WorkTitle{书信第199号}}

\Emph{\Person{凯撒勒雅的圣巴西略}}\par

\OriginalTerm{Canonica Secunda}{Canonica Secunda}.\par

\Emph{致\Person{安菲洛基乌斯}，论规条}.\par

我前些时候写了答复尊座的问题，但没有寄出那封信，部分是因为我长期危险疾病没有时间，部分是因为没有可差遣的人。我身边只有少数有旅行经验且适合此类服务的人。当您了解到我延迟的原因时，请原谅我。我对您的学习热忱和谦逊感到惊讶。您被托付了教师的职务，却屈尊学习，而且向我这样自称无甚知识的人学习。然而，由于您出于对天主的敬畏，同意做别人可能犹豫的事，我就有义务甚至超越我的能力来辅助您的热心和正义热忱。\par

XVII. 您曾问我关于司铎\Person{比亚诺}的事——因为他发过誓，能否被录入圣职人员？我知道我已对安提约基雅的圣职人员下过关于所有与他一同发誓者的通令；即他们应回避公开集会，但可在私下进行司铎职务。此外，他因自己神圣的职责不在安提约基雅，而在\Place{依科尼雍}，故有更多执行司祭职务的自由；因为，正如您亲自写信给我，他选择居住在后地而非前地。因此，此人可以被接纳；但尊座必须要求他悔改，因为他在未信者面前轻率鲁莽地发誓，不能忍受那微小危险的困扰。\par

XVIII. 关于堕落的贞女，她们在主前宣示了贞洁生活后，因陷于肉体的情欲而使誓愿落空，我们的教父们温柔谦和地宽容跌倒者的软弱，规定她们可被接纳，但要经过一年，并将她们列入二次婚姻者之列。然而，随着天主恩宠的助佑，教会在前进中日益强大，贞女团体也越来越多，我的判断是：应当审慎考虑这一行为，正如从思考中所呈现的那样，以及从上下文中可发现的圣经的心意。寡妇地位低于贞女；因此寡妇的罪远不及贞女的罪。让我们看看保禄给\Person{弟茂德}写了什么。\BibleQuote{至于年轻的寡妇，你要拒绝：因为当她们放纵情欲违背基督时，就想再嫁；她们因背弃了起初的信德而该受惩罚。} \BibleRef{弟前 5:11} 因此，如果寡妇因蔑视对基督的信德而负重罪，那么，作为基督的新娘、奉献于主的特选器皿的贞女，我们该怎样想呢？即使是奴隶，私行婚配，给家中带来污秽，且因邪恶生活亏欠主人，也是重大的过错；然而，新娘成为奸妇，羞辱与新郎的结合，放纵于不洁的享乐，则更加令人震惊。寡妇，如同堕落的奴隶，确实被定罪；但贞女则犯奸淫罪。凡与别人妻子同居的，我们称之为奸夫，在他停止犯罪之前，我们不接纳他共融；同样，我们也将在占有贞女的人身上这样规定。然而，有一点必须事先确定，就是\Quote{贞女}这一名称是献给自愿奉献给主、弃绝婚姻、接受圣洁生活的女子。我们只承认从完全理性的年龄开始的宣愿。因为在这种情况下，接受单纯的儿童之言是不对的。但一个十六七岁的女孩，心智成熟，经过严格考察，并且坚定，持续请求被接纳，那么就可以被列入贞女之列，她的宣愿予以批准，其违犯严加惩罚。许多女孩在未成熟时由父母、兄弟和其他亲属带来，她们内心并无独身生活的冲动。亲友的目的是为了自己的利益。这样的女子不可轻易接纳，除非我们公开调查她们自己的意愿。\par

XIX. 我不承认男人的宣愿，除非是那些登记在隐修士团体中、似乎已秘密接受独身生活的人。然而，我认为在他们身上，应事先进行考察，并接受他们明确的宣愿，以便当他们可能回到肉体享乐的生活时，可受奸淫者应得的惩罚。\par

XX. 我认为不应谴责那些在异端中宣示贞女身份、而后又选择结婚的女子。\BibleQuote{凡法律所说的，是对那属于法律的人说的。} \BibleRef{罗 3:19} 那些尚未背负基督之轭的人，不承认主的法律。因此，她们当在教会中被接纳，如同因对基督的信德，在这些罪上也像在一切罪上一样获得赦免。一般规则是，所有在慕道阶段犯的罪不予追究。教会不接纳未经圣洗的人；而在这种情况下，非常必要遵守长子的权利。\par

XXI. 如果一個男人與妻子同住卻不滿足於婚姻而犯姦淫，我視他為姦淫者，並延長其懲罰期。然而，我們沒有教規將之定為通姦罪，如果所犯之罪是針對未婚女子。因為關於通姦者，經上說：\BibleQuote{被玷污的已被玷污}（\BibleRef{耶 3:1}），她不可回到丈夫那裡；並且\Quote{與通姦者同居的是愚昧不敬虔的人。} 但犯姦淫者不被剝奪與妻子共同生活的權利。因此，妻子應從丈夫犯姦淫後接納他，但丈夫應從家中驅逐被玷污的女人。這裡的論證不易，但習俗如此通行。\par

XXII. 關於男人強搶婦女後與之同居，如果她們原是許配給別人的，則必須先將婦女歸還給原訂婚者，然後才可接納這些男人，無論原訂婚者願意接受她們還是與她們分離。至於未訂婚的女子被搶，應先將女子歸還給她自己的親屬，並交由她的親屬（父母、兄弟或任何有監護權的人）決定。如果他們願意讓她與搶婚者同居，可以維持；但如果他們拒絕，不得使用暴力。對於用勾引（秘密或暴力）得到妻子的男人，必須認定他犯姦淫罪。姦淫者的懲罰定為四年：第一年，他們必須被驅逐出祈禱，在教堂門口哭泣；第二年，可以被接納聽講道；第三年，進行補贖；第四年，與民眾站在一起，但不得參與獻祭。最後，他們可以被准許領受善恩的共融。\par

XXIII. 關於男人娶兩姐妹，或女人嫁兩兄弟，我寫了一封短函，副本已寄給尊駕。娶自己兄弟妻子的人，必須離開她後才可被接納。\par

XXIV. 一位寡婦，若其姓名列於寡婦名單中，即由教會供養，若她再婚，宗徒命令不再供養她。（\BibleRef{弟前 5:11}）\par

對於鰥夫沒有特別規定。為再婚所定的懲罰足以適用。若一位六十歲的寡婦選擇再與男人同居，她將被認為不配領受善恩的共融，直到她不再被不潔的欲望所動。如果我們在她六十歲之前這樣做，責任在我們，不在那女人。\par

XXV. 一個男人若強行玷污女人並留她為妻，應受強姦之罰，但可合法娶她為妻。\par

XXVI. 姦淫不是婚姻，也不是婚姻的開始。因此，最好盡可能分開因姦淫結合的人。如果他們堅持同居，讓他們接受姦淫的懲罰。允許他們繼續同居，以免發生更壞的事。\par

XXVII. 關於司鐸因無知捲入非法婚姻，我作了適當規定：他可以保留席位，但必須停止其他職務。因為在這種情況下，寬恕足夠了。一個需要治療自己創傷的人卻祝福他人，是不合理的，因為祝福是傳遞聖潔。一個人因無知而犯罪，自己都沒有聖潔，如何能傳遞聖潔？讓他無論公開還是私下都不要祝福，也不要向他人分送基督的身體，也不要執行任何其他神聖職務，但應滿足於他的榮耀席位，以哭泣懇求主，使他無知所犯的罪得赦免。\par

XXVIII. 在我看來，任何人發誓戒食豬肉都是可笑的。請教導人們避免愚蠢的誓言和承諾。表明使用豬肉是完全無關緊要的。凡天主所造的物，若感謝著領受，就沒有一樣可棄絕的。（\BibleRef{弟前 4:4}）這誓言荒唐，禁食不必要。\par

XXIX. 特別值得注意的是，有權勢的人發誓要傷害屬下。補救方法有二：首先，教導他們不要隨意起誓；其次，不要堅持邪惡的決定。任何人在執行誓言傷害他人時被制止，應對他的魯莽誓言表示懺悔，決不可藉著虔敬的藉口來鞏固他的邪惡。黑落德履行誓言並無益處，為了免於偽誓，他成了先知的謀殺者。（\BibleRef{瑪 14:10}）起誓是絕對禁止的（\BibleRef{瑪 5:34}），所以傾向於惡的誓言理應被譴責。因此，起誓者必須改變心意，不要堅持確認他的不敬。再深入思考這事之荒謬。假設某人發誓要挖出弟弟的眼睛：他是否應該執行誓言呢？或者殺人？或者違反其他任何誡命？\Quote{我起誓，並且必履行}，不是為犯罪，而是為\Quote{遵守你的公義判決。} 我們有責任消除和摧毀罪惡，正如我們有責任以不可改變的決意堅守誡命。\par

XXX. 關於綁架者，我們沒有古代規則，但我表達了個人判斷。懲罰期為三年，罪犯及其同謀被排除在服務之外。沒有暴力的行為不受懲罰，只要不是事先有侵犯或搶劫。寡婦是獨立的，跟隨與否在於她自己的權力。因此，我們不必理會藉口。\par

XXXI. 一個婦女的丈夫離家失踪，她在未獲知其死亡證據前再嫁，犯通姦罪。犯致死罪的聖職人員，被剝奪聖秩，但不被排除於平信徒的共融之外。不可為同一過失懲罰兩次。\par

XXXIII. 對在路上生下孩子而置之不理的婦女，應以殺人罪起訴。\par

XXXIV. 犯奸淫的妇女，因热心而承认自己的过犯，或以任何方式被定罪者，我们的先父不允许她们被公开暴露，以免我们在定罪后导致她们死亡。但他们命令应将她们排除于共融之外，直到完成补赎的期限。\par

XXXV. 对于被妻子离弃的男子，必须考虑离弃的原因。如果她似乎无故离弃他，他应得宽恕，而妻子应受惩罚。宽恕将给予他，使他能与教会共融。\par

XXXVI. 军人的妻子在丈夫缺席期间再婚，将与那些在丈夫旅行时未等其归来的妻子适用同一原则。然而，她们的情况确实允许一些宽限，因为更有理由怀疑死亡。\par

XXXVII. 男子拐走他人妻子后与之结婚，将因第一个行为而承担奸淫的罪责，但对第二个行为则无罪。\par

XXXVIII. 女子违背父亲的意愿与人私奔，犯奸淫；但如果她们的父亲与她们和好，此行为似乎可得到补救。然而，她们不是立即恢复共融，而是受罚三年。\par

XXXIX. 与奸夫同居的女人，整个时期都是奸妇。\par

XL. 女子违背主人的意愿委身于男子，犯奸淫；但若后来她接受自由婚姻，她便结婚。前者是奸淫，后者是婚姻。不独立之人的契约无效。\par

XLI. 寡居且独立的女子，若无人阻止其同居，可与丈夫同住而无罪；因为宗徒说：\Quote{但如果她的丈夫死了，她就自由了，可以随意嫁人，只要在主内。} \BibleRef{格前 7:39}\par

XLII. 未经当权者许可而缔结的婚姻是奸淫。如果父亲和主人都已去世，订婚双方无罪；正如若当权者同意同居，则婚姻具有确定性。\par

XLIII. 击打邻人致死的人，是杀人犯，无论他是先出手还是自卫。\par

XLIV. 与异教徒犯奸淫的\OriginalTerm{女执事}{deaconess}，可以被接纳悔改，并在第七年被准许参与祭献；当然，如果她保持贞洁。异教徒在相信后转向偶像崇拜，又回到自己的呕吐物。然而，我们不放弃女执事的身体供肉体之用，因为她已被祝圣。\par

XLV. 若有人在接受了基督徒的名号后侮辱基督，他从这名号中得不到任何益处。\par

XLVI. 女子不情愿地嫁给一个当时被妻子离弃的男子，之后因前妻回到他身边而被休弃，她犯奸淫，但非自愿。因此，她将不被禁止再婚；但最好保持现状。\par

XLVII. \OriginalTerm{节制派}{Encratitæ}、\OriginalTerm{麻衣派}{Saccophori}和\OriginalTerm{弃绝派}{Apotactitæ}并不被视为与\OriginalTerm{诺瓦替安派}{Novatians}相同，因为对后者已有明文规定，尽管不同；而对前者则未提及。所有这些人都按同一原则我给他们重新授洗。如果你们那里禁止为他们重新授洗，为了某种安排，但让我的原则得胜。他们的异端，可以说是\OriginalTerm{马西昂派}{Marcionites}的一个分支，因为他们憎恶婚姻，拒绝饮酒，并称天主的受造物为不洁。因此，我们不接纳他们进入教会，除非他们接受我们的洗礼而受洗。不要让他们说他们已经因圣父、圣子及圣神之名受洗，因为他们效法\Person{马西昂}和其他异端，将天主视为邪恶的制造者。因此，如果此决定通过，更多的主教应聚集在同一地方，颁布这样的教规，以便行动安全，并对此类问题的回答赋予权威。\par

XLVIII. 被丈夫离弃的女子，依我判断，应保持现状。主说：\Quote{任何人休妻，除非为了奸淫的缘故，便是使她犯奸淫；} \BibleRef{玛 5:22} 因此，称她为奸妇，主就禁止她与另一男子发生关系。因为既然男子有罪，因为他导致了奸淫，而女子也被主称为奸妇，原因在于她与另一男子发生关系，那么她怎能无罪呢？\par

XLIX. 遭受强暴不应成为定罪的理由。因此，女奴若被自己的主人强迫，则无罪。\par

L. 关于三次婚姻没有法律：第三次婚姻不受法律约束。我们将此类事情视为教会的玷污。但我们不使其受公开谴责，因为总比无节制的奸淫好。\par

\SourceRecord{Translated by Blomfield Jackson.；Nicene and Post-Nicene Fathers, Second Series；Vol. 8.；Edited by Philip Schaff and Henry Wace.；Buffalo, NY: Christian Literature Publishing Co.,；1895.；<http://www.newadvent.org/fathers/3202199.htm>.}

% structure: p0001 source=h1 target=section reason=page-title

\section{\WorkTitle{第200封书信}}

\Emph{\Person{凯撒勒雅的圣巴西略}}\par

\Emph{致\Person{安斐洛基}，\Place{依科尼雍}主教}\par

我接二连三地患病，所有的工作——不仅是教会的事务，还有那些扰乱教会的人加给我的工作——使我整个冬天直至如今都不得闲。因此我根本无法派人到你那里去，也无法亲自拜访你。我猜想你的处境也类似；当然不是指疾病，愿主保佑你继续健康，以便遵行祂的诫命。但我知道，你对众教会的关怀也带给我与你同样的忧虑。我现在正打算派一个人去准确了解你的情况。然而，当我心爱的儿子\Person{梅肋提乌斯}——他正负责招募新兵——提醒我可通过他向你致意时，我欣然借此机会写信，并求助于送信人的热心服务。他本人就是一封书信，既是由于他亲切的品性，也因为他对我的一切了如指掌。因此，我恳请你的敬长特别为我祈祷，求主使我摆脱这劳苦的身体；赐祂的教会和平；赐你安息；并当你以宗徒的方式照你所开始的处理完\Place{吕考尼雅}的事务后，也能有机会访问此地。无论我仍寄居肉身，或已蒙召离世归主，我希望你会关心我们这一地区如同你自己的（它确实是你的），坚固一切软弱的，激发一切懈怠的，并借着住在你内的圣神的能力，使一切都转变为悦乐主的状况。我极其尊敬的儿子\Person{梅肋提乌斯}和\Person{梅利提乌斯}，你早已认识并知道他们忠诚于你，请继续善待他们并为他们祈祷。这足以保全他们。我恳求你以我的名义问候与你圣德同在的所有人，包括所有圣职人员和所有在你牧养下的平信徒，以及我的虔敬的弟兄和同工。请纪念真福殉道者\Person{欧普西基乌斯}，不必等我再次提及。不要费心在准确那天到来，而要提前，这样若我仍活在世上，就给我喜乐。在此期间，愿圣者的恩宠保佑你为我并为天主的教会而存留，健康富足于主内，并为我祈祷。\par

\SourceRecord{Translated by Blomfield Jackson.；Nicene and Post-Nicene Fathers, Second Series；Vol. 8.；Edited by Philip Schaff and Henry Wace.；Buffalo, NY: Christian Literature Publishing Co.,；1895.；<http://www.newadvent.org/fathers/3202200.htm>.}

% structure: p0001 source=h1 target=section reason=page-title

\section{书信第201封}

\Signature{\Person{凯撒勒雅的圣巴西略}}\par

\Emph{致\Place{依科尼雍}主教\Person{安斐洛纠}}\par

我渴望与你会面，原因甚多：既可在所处理之事上受益于你的忠告，又能在久别之后见到你，以稍慰你缺席之苦。然而，我俩皆因相同缘由受阻——你因降临于你的疾病，我因缠绕已久仍未离去的病痛——若你同意，让我们彼此宽恕，好使双方皆可免于责难。\par

\SourceRecord{Translated by Blomfield Jackson.；Nicene and Post-Nicene Fathers, Second Series；Vol. 8.；Edited by Philip Schaff and Henry Wace.；Buffalo, NY: Christian Literature Publishing Co.,；1895.；<http://www.newadvent.org/fathers/3202201.htm>.}

% structure: p0001 source=h1 target=section reason=page-title

\section{书信第202封}

\Emph{\Person{凯撒勒雅的圣巴西略}}\par

\Emph{致\Place{依科尼雍}主教\Person{安菲洛基乌斯}。}\par

若在其他情况下，我会认为与尊驾相会是一种特别的荣幸，但尤其如今，当我们聚首之事如此重要之时。然而，我仍残余的疾病足以阻止我哪怕移动极短的距离。我曾尝试驾车直到殉道者之处，却几乎复发了旧病。因此，务请宽恕我。若此事能推迟数日，我必赖天主的恩宠与您会合，分担您的忧虑。若事情紧迫，就靠天主的助佑完成必须做的事；但请把我看作与您同在，并参与您可敬的作为。愿您，藉圣者的恩宠，得蒙保存于天主的教会，在主内坚强喜乐，并为我祈祷。\par

\SourceRecord{Translated by Blomfield Jackson.；Nicene and Post-Nicene Fathers, Second Series；Vol. 8.；Edited by Philip Schaff and Henry Wace.；Buffalo, NY: Christian Literature Publishing Co.,；1895.；<http://www.newadvent.org/fathers/3202202.htm>.}

% structure: p0001 source=h1 target=section reason=page-title

\section{信函 203}

\Person{凯撒勒雅的圣巴西略}\par

\Emph{致沿海诸主教}\par

我一直强烈渴望与你们相见，但屡次有阻碍临到，使我无法实现我的目的。我或因疾病受阻——你们清楚知道，从我青年时代直到现在的老年，这病痛一直是我常伴的伙伴，与我一同成长，并因天主公义的审判（祂以智慧安排万事）而惩戒我；或受教会事务的阻碍，或与真理教义的反对者斗争而受阻。[直到今日，我仍生活在极大的痛苦和忧伤中，心中常存一种感觉：你们对我是欠缺的。因为当天主告诉我——祂正是为此而取了肉身的居留，以便借着职责的榜样来规范我们的生活，并亲自向我们宣讲天国的福音——当祂说：\Quote{如果你们彼此相爱，众人因此就可认出你们是我的门徒}，而主在即将完成肉身中的救恩计划时，将自己的平安作为临别礼物留给门徒，说：\Quote{我把平安留给你们，我将我的平安赐给你们}，我无法说服自己，若没有对他人的爱，以及尽我所能对众人的和睦，我可以被称为耶稣基督的忠仆。我长久等待你们仁爱的机会来访问我们。因为你们并非不知，我们如同海中突出的礁石，承受着异端浪潮的怒涛，它们在我们周围碎裂，并未淹没后方的区域。我说\Quote{我们}，为的是将之归于天主的恩宠，而非人的能力，祂藉着人的软弱彰显祂的大能，正如先知在主位中所说：\Quote{你们不惧怕我吗？我曾以沙为海的界限}。因为藉着最弱小、最卑微的沙，全能者限制了浩瀚的大海。既然我们的处境如此，你们的仁爱本应常派真正的弟兄来探望我们这些在风暴中劳苦的人，并更经常地寄来仁爱的书信，一部分是为了坚定我们的勇气，一部分是为了纠正我们的错误。因为我们承认，我们身为人类，生活在肉身中，难免犯下无数错误。]\par

2. 但迄今为止，极可敬的弟兄们，你们未曾给我应有的对待；这有两个原因：或者你们没有看出正确的做法，或者你们受了一些关于我的流言的影响，认为我不值得你们以友爱来探望。现在，因此，我亲自采取主动。我请求表明，我完全准备好在你们面前洗脱对我的指控，但条件是，控告我的人被允许与我在你们尊前当面对质。如果我被定罪，我不会否认我的错误。你们在定罪之后，将从主那里获得赦免，因你们远离了我这个罪人的共融。成功的控告者也将因我的隐秘邪恶被揭露而得到报偿。然而，如果你们在得到证据之前就定我的罪，我并不会因此受损，除了失去我最为珍视的财产——你们的友爱；而你们，作为回报，将因失去我而遭受同样的损失，并且似乎与福音的话相悖：\BibleQuote{我们的法律岂能不经审问就定人的罪吗？}\BibleRef{若 7:51} 此外，控告者如果不能证明他所说的话，将显示出他除了坏名声之外，什么也没有从恶言中得到。除了诽谤者所声称的、因他所犯的恶行而应得的名称之外，还能给他什么恰当的名称呢？因此，让控告者不是作为诽谤者，而是作为控告者出现；不，我不称他为控告者，我宁愿将他视为一个弟兄，以仁爱劝诫，并为了我的改正而提出证据。你们不要作流言的听者，而要作证据的检验者。我也不应因我的罪未显明而得不到医治。\par

[3. 不要让这个想法影响你们：\Quote{我们住在海上，我们不受普通人的苦难，我们不需要别人的帮助；那么，外来的共融对我们有什么好处？}因为那位用海将岛屿与大陆分开的主，也用爱把岛屿上的基督徒与大陆上的基督徒联系在一起。弟兄们，除了故意的疏远，没有什么能把我们彼此分开。我们同有一位主，一个信德，同一个望德。双手需要彼此，双脚互相支撑。眼睛因一致而获得清晰的视力。我们这边承认自己的软弱，并寻求你们的同情。因为我们确信，即使你们身体不在场，但借着祈祷的帮助，你们在这最危急的时刻将给我们带来极大的益处。你们若做出连不认识天主的外邦人都不认可的声明，那在人们面前既不体面，也不取悦于天主。甚至外邦人，据我们所闻，尽管他们居住的国度自给自足，但因未来的不确定性，仍重视彼此结盟，并寻求相互往来，认为这对他们有利。然而，我们——作为那些曾经立下法规，仅凭简短的书信便可将共融的凭证从地极传到地极，使所有人与所有人成为同胞和亲友的父辈的子孙——如今却将自己与整个世界隔绝，既不为我们的孤立感到羞愧，也不战栗于主那可怕的预言已应验在我们身上：\Quote{由于罪恶的增多，许多人的爱德必要冷淡。}]\par

4. 至可敬的弟兄们，不要，千万不要容忍此事。反之，要以和平的书信和爱德的问候，为往事安慰我。你们从前的疏忽在我心中造成创伤。请以温柔的触摸来抚慰其痛苦。或者你们愿意到我这里来，亲自调查你们所听到关于我软弱之事的真相，或者因更多谎言的添加，我的罪过被报告给你们更加严重，我也必须接受这一切。我准备张开双臂欢迎你们，并将自己交给最严格的检验，只愿爱德主持整个过程。或者你们愿意指定你们地区的某个地点，让我可以前去拜访你们，尽我当尽之责，并尽可能接受审查，以治愈过去并防止将来的诽谤，我也接受。虽然我的肉体软弱，但只要我还有一口气，我就有责任妥善履行每一项有助于基督众教会建立的责任。我恳求你们，不要轻视我的恳求。不要迫使我把我的苦楚向他人诉说。迄今为止，弟兄们，正如你们所深知，我一直把悲伤埋在心里，因为我羞于向远处与我们共融的人提及你们对我的疏远。我既怕使他们痛苦，又怕让恨我的人高兴。现在我独自写这封信；但我以\Place{卡帕多细雅}全体弟兄的名义寄出，他们嘱咐我不要随便派一位信使，而要派一位，万一我在急于避免冗长而遗漏重要之处时，能用天主所赐的智慧补足的人。我所指的就是我可爱的、可敬的同司铎\Person{伯多禄}。请以爱德接待他，并平安地打发他到我这里来，使他能成为向我报告好消息的信使。\par

\SourceRecord{Translated by Blomfield Jackson.；Nicene and Post-Nicene Fathers, Second Series；Vol. 8.；Edited by Philip Schaff and Henry Wace.；Buffalo, NY: Christian Literature Publishing Co.,；1895.；<http://www.newadvent.org/fathers/3202203.htm>.}

% structure: p0001 source=h1 target=section reason=page-title

\section{书信二〇四}

\Person{凯撒勒雅的圣巴西略}\par

\Signature{致新凯撒利亚人}\par

1. 我们双方长期保持沉默，可敬且亲爱的弟兄们，仿佛彼此心怀怨恨。然而，有谁对那些伤害过自己的人如此阴郁、不肯和解，以至于把始于厌恶的怨恨几乎延续整个人生呢？在我们这种情况中，据我所知，并无任何正当理由导致疏远，相反，从一开始就有许多强有力的理由要求最亲密的友谊与合一。其中最大也是首要的，便是我们主的诫命，祂明确地说：\Quote{如果你们彼此相爱，众人因此就可认出你们是我的门徒。}\BibleRef{若 13:35} 再者，宗徒明确地将爱德之益处摆在我们面前，告诉我们爱德是法律的满全；\BibleRef{罗 13:10} 又在另一处说爱德是一件美事，应置于一切伟大美好事物之上：\Quote{我若能说人间和天使的言语，若没有爱德，我就成了发声的锣、响亮的钹。我若有先知之恩，又明白一切奥秘和各种知识；我若有全备的信德，甚至能移山，若没有爱德，我什么也不算。我若把我所有的财产全施舍给人，并舍身投火被烧，若没有爱德，为我毫无益处。}\BibleRef{格前 13:1} 这并不是说上述各项若没有爱德就无法做到，而是圣者愿意，正如祂亲口所说，通过夸张的修辞手法将这诫命的卓越性彰显出来。\par

2. 其次，如果拥有相同的教师能极大地促进亲密关系，那么你们和我便拥有相同的神圣奥迹的教师和神父，他们从一开始就是你们教会的奠基人。我指的是伟大的额我略，以及所有相继登上你们主教宝座的人，如同相继升起的星辰，循着同样的轨迹，为一切渴望的人留下了天上政制的清晰印记。如果血缘关系不可轻视，反而大大有助于持久的合一与共融，那么这些权利对你们和我而言也是自然存在的。那么，可敬的城市（因为我通过你们向整个城市发言），为什么没有来自你们的公正文书，没有令人欣慰的声音，而你们的耳朵却向那些致力于诽谤的人敞开呢？因此，我越看到他们的目的，就越应该叹息。毫无疑问，谁是诽谤的始作俑者。他以许多恶行而闻名，但尤以此种邪恶而著称，故此罪过成了他的名字。但你们必须容忍我的直言。你们将双耳向我的诽谤者敞开。你们由衷欢迎所听到的一切，不加任何查问。你们中没有一人能分辨谎言与真理。当一个人孤军奋战时，谁曾因缺乏恶毒的指控而受苦？当受害者不在场时，谁曾被证实说谎？当诽谤者坚持某事如此，而受谤者既不在场也听不到指控时，有什么理由在听者听来不是貌似有理的呢？难道世间的普遍习俗没有教导你们，关于这些事，若想做一个公正不偏的听者，就不应完全被第一个发言者左右，而应等待被告的辩护，以便通过双方论点的比较来证明真理？\Quote{你们要作正义的判断。}\BibleRef{若 7:24} 这条诫命是最为必要的救恩诫命之一。\par

3. 我说这话时，并未忘记宗徒的话，他逃避人的法庭，将一生的辩护保留给无误的审判席，他说：\Quote{我被你们论断，或被人的法庭审判，为我都是极小的事。}\EditorialNote{格前 4} 你们的耳朵已被虚妄的诽谤所占据，这些诽谤触及我的品行，触及我对天主的信德。然而，我深知诽谤者同时伤害三个人：受害者、听者和自己。至于我自己的冤屈，我本应保持沉默，请相信，并非因为我不在乎你们的好评（否则我现在为何写信并竭力恳求不失去它？），而是因为我看到三个受害者中，我自己受伤最轻。的确，我将失去你们，但你们正失去真理，而这一切的始作俑者正在使你们与我分离，但他自己却使自己远离主，因为没有人能通过做禁止的事而亲近主。我恳求，与其说是为了我自己，不如说是为了你们，为使你们免于无法承受的冤屈。因为谁能承受比失去最珍贵的东西——真理——更大的灾难呢？\par

4. [弟兄们，我说什么？并非我无罪，并非我的生活没有无数过犯。我认识自己；确实，我不停止为我的罪流泪，但愿我能以此平息我的天主，并逃脱威胁我的惩罚。但我说：那审判我的人，若他能说自己眼目清明，就让他寻找我眼中的木屑吧。] 弟兄们，我承认，我需要健全健康之人的照顾，且需要很多。若他不能说自己是清明的——越清明越不会这样说——（因为完美的人不抬高自己；若他们抬高，必会陷入法利塞人骄傲的罪责，他自我称义，却定了税吏的罪）就让他与我同去医生那里；让他不要\Quote{在时候未到以前审判，等主来，祂要照出暗中的隐情，显明人心的意念。}\BibleRef{格前 4:5} 让他记着这话：\Quote{你们不要判断，你们就不受判断；}\BibleRef{玛 7:1} 还有\Quote{你们不要定罪，你们就不被定罪。}\BibleRef{路 6:37} [总之，弟兄们，若我的过犯尚有可治，为何这样的人不听从教会导师的话：\Quote{劝勉、斥责、告诫}？\BibleRef{弟后 4:2} 若我的不义已不可救药，为何他不当面抵抗我，公布我的过犯，使教会脱离我带来的祸害？] 不要容忍在齿缝间对我进行的诽谤。这是从磨坊女仆口中可能发出的辱骂；这是你在街头流浪汉身上可能看到的炫耀；他们的舌头磨尖了准备任何毁谤。但是[有主教们在；让他们被申诉。有天主各教区的神职；让最杰出者聚集。无论何人，让他自由发言，使我面对一项指控，而非诽谤。] 让我隐秘的恶行暴露无遗；让我不再被恨恶，而是被当作弟兄劝诫。我们有罪之人被那蒙福无罪者怜悯，比被他们怒气对待更为公正。\par

5. [若过错涉及信仰，请向我指出文献。同样，让一次公平公正的调查被指定。让控告书被宣读；让它经受检验，看是否源于控告者的无知，而非写作材料本身的错误。因为正确之事，在判断不精确的人看来常非如此。当天平两臂不等时，等重的砝码看起来也不等。] 因疾病而味觉被毁的人，有时会以为蜂蜜是酸的。有病的眼看不见许多存在之物，却注意到许多不存在之物。同样的事常发生在对言辞的力量的判断上，当批评者劣于作者时。批评者本当以与作者大致相同的训练和装备出发。一个对农业无知的人，完全不能批评农事；和谐与不和谐的区别只能由训练有素的音乐家来恰当地判断。但任何人，只要愿意，都可以自封为文学批评家，尽管他不能告诉我们他在哪里上学，或他的教育花了多少时间，且对文字一无所知。我清楚看见，即使对圣神的话语，术语的研究并非每个人都可以尝试，而只有那有辨别神恩的人，如宗徒在神恩的区别中所教导的——\Quote{这人蒙圣神赐他智慧的言语，那人也蒙同一位圣神赐他知识的言语，又有人蒙同一位圣神赐他信德，还有人蒙同一位圣神赐他治病的奇恩，又有人能行奇迹，有人能说预言，有人能辨别神恩。}\BibleRef{格前 12:8} 因此，若我的恩赐是属神的，那想要判断它们的人必须证明他自己拥有\Quote{辨别神恩}之恩赐。反之，若照他诽谤的主张，我的恩赐是属这世界的智慧，就让他显示他是这世界智慧的行家，我将服从他的裁决。且[不要让人以为我在找借口逃避指控。至亲的弟兄们，我将此交给你们，让你们亲自调查针对我的指控。难道你们如此迟钝，竟全然依赖辩护人发现真理？若所讨论之点使你们觉得本身就很清楚，就劝说那些戏谑者放弃争论。[若有什么你们不明白，可透过指定且公正待我的人向我提问；或以书面形式向我询问。并尽一切努力，使什么都不被遗漏。]]\par

6. 有什么比我的信德更清楚的证据呢？我就是由我的祖母——那位有福的妇人，来自你们中的著名的\Person{玛克利纳}——抚养长大的。她教导我那位有福的\Person{额我略}的言语；这些言语，就记忆所及保存到她那个时代，她珍藏在自己心中，同时在我年幼时，以虔敬的教义塑造和培养了我。当我获得思考能力时，我的理性因年龄成熟而发展，我走遍许多海陆，凡遇到走在所交付的虔敬规范中的人，那些我尊为教父，并使他们在我灵魂朝向上主的旅程中成为向导。直到今日，借着那以祂神圣的召叫召叫我认识祂自己的那位的恩宠，我知道没有任何与纯正教导相悖的道理渗入我心；我的灵魂也从未被\Person{亚略}那臭名昭著的亵渎所玷污。如果我曾接纳过任何来自那位导师、将他们的不健全深藏于内、或说出虔敬话语、或至少不反对我所说的话的人，我是基于这些条件接纳他们；我并没有完全依靠自己来判断他们，而是遵循教父们关于他们的决定。因为我收到了那位极其有福的教父、亚历山大里亚主教\Person{亚太纳削}的书信，这信我拿在手中，向任何询问的人展示。他在信中明确宣告：任何愿意从亚略异端归正并接受尼西亚信经的人，应当毫不犹豫、毫不困难地被接纳，并引用马其顿和亚洲主教的全体同意来支持他的观点。我，自认为有责任追随这样一位伟人和制定规则之人的崇高权威，并且我自己也渴望获得那应许给缔造和平之人的赏报，遂将所有接受那信经的人登记在共融者的名册上。\par

7. \Quote{公平的判断，不应基于一两个行为不正的人，而应基于全世界众多与我藉着主的恩宠相连的主教。请问询\Place{彼息狄雅}人、\Place{里考尼雅}人、\Place{依撒乌利亚}人、两省\Place{夫黎基雅}人、你们邻邦\Place{亚美尼亚}人、\Place{马其顿}人、\Place{亚该亚}人、\Place{依里利亚}人、\Place{高卢}人、\Place{西班牙}人、整个\Place{意大利}、\Place{西西里}人、\Place{非洲}人、\Place{埃及}健康的部分、\Place{叙利亚}余剩的部分——所有这些人都写信给我，也收到我的回信。} 从这些信函，无论是由他们寄出的，还是由我们发给他们的，你都可以知道我们同心合意，意见一致。 \Quote{凡回避与我共融的人，以你的精确性必能看出，他是自绝于整个教会。请环顾四周，弟兄们，你们与谁共融？若你们不接受从我而来的，还有谁承认你们？请不要逼我不得不对我如此珍爱的教会说出任何不愉快的话。} 有些事我现在隐藏在我心底，暗自叹息哀叹我们生活的邪恶时代，因为那些长期以兄弟之爱彼此联合的最大教会，如今无缘无故地彼此对立。请不要，喔请不要逼我向所有与我共融的人抱怨这些事。不要逼我说出那些我至今靠反省抑制、藏在心里的话。倒不如我被除掉，而教会合一，远胜于天主子民因我们幼稚的恶意遭受如此祸害。 \Quote{请问你们的祖先，他们必告诉你们，虽然我们的地区在位置上分开，但在心灵上却是一体，并由同一情感统治。民众的交往频繁，圣职人员的访问频繁；牧者们彼此也有如此深情，以致各自以对方为在主之事上的教师和向导。}\par

\SourceRecord{Translated by Blomfield Jackson.；Nicene and Post-Nicene Fathers, Second Series；Vol. 8.；Edited by Philip Schaff and Henry Wace.；Buffalo, NY: Christian Literature Publishing Co.,；1895.；<http://www.newadvent.org/fathers/3202204.htm>.}

% structure: p0001 source=h1 target=section reason=page-title

\section{书信 205}

\Emph{凯撒勒雅的圣巴西略}\par

\Emph{致主教厄尔丕迪乌斯。}\par

我再次派遣我极敬爱的司铎默肋提乌斯，向阁下传达我的问候。我原本决意让他休息，因为他为了基督的福音而制服自己的身体，自愿承受了软弱。但我觉得，由他这样的人代为问候是适宜的：他能凭自身弥补我书信的一切不足，无论对写信者还是收信者，都是一封活生生的书信。我也在实现他长久以来的强烈愿望，即见到阁下的尊贵，因为他早已体验到您的高尚品质。因此，我现在恳请他前往您处，藉着他，我偿还了欠您的拜访之债，并恳求您为我及天主的教会祈祷，愿主使我摆脱福音敌人的伤害，以平安和宁静度过此生。然而，若您明智地认为，我们有必要前往同一地点，与其他可敬的海滨地区主教弟兄们会面，那么请您亲自指定合适的地点与时间。请写信给我们弟兄们，以便各位在预定时间放下手头事务，能为建立天主的教会有所贡献，消除我们因彼此猜疑而承受的痛苦，并建立爱德——主曾亲自规定，若缺少爱德，遵守任何诫命都归于无效。\par

\SourceRecord{Translated by Blomfield Jackson.；Nicene and Post-Nicene Fathers, Second Series；Vol. 8.；Edited by Philip Schaff and Henry Wace.；Buffalo, NY: Christian Literature Publishing Co.,；1895.；<http://www.newadvent.org/fathers/3202205.htm>.}

% structure: p0001 source=h1 target=section reason=page-title

\section{书信二〇六}

\Emph{\Person{圣巴西略}}\par

\Emph{致主教\Person{艾尔彼丢}。慰问函。}\par

如今，我尤其感到自己身体的软弱，因它阻碍了我灵魂的益处。若事情如我所愿，我此刻便不会以书信或信使与你交谈，而会亲自前去偿还爱德之债，并当面享受属灵的益处。然而，我如今的处境是，即使能在自己教区内进行必要的堂区巡视，也已是万幸。愿主赐予你力量和心甘情愿，也赐予我，除了热切的渴望外，还有能力在\Place{科马纳}地区与你相聚时享受你的陪伴。我怕你家庭中的困扰可能成为某种障碍，因为我已得知你因失去幼子而悲痛。对祖父而言，他的去世无疑是令人伤痛的。然而，对于一个已达如此崇高德行的人，并且同时凭借对现世的经验和属灵的培育而深知人性者，亲人的离世不应是完全无法忍受的。主向我们要求的，并非他向每个人所要求的。普通大众凭习惯生活，但基督徒的生活准则乃是主的诫命，以及古圣先贤的榜样，他们的伟大心灵尤其体现在逆境之中。因此，为使你自己也为后世留下坚韧与对盼望之事的真诚信靠的榜样，请表明你并非被悲伤所胜，而是超脱于哀痛之上，在困苦中忍耐，在盼望中喜乐。切莫让这些事妨碍我们预期的会面。孩童因其年幼而可视为无罪，但你和我则肩负着按照主的命令事奉主，并随时准备管理教会事务的责任。为妥善履行此职责，主为忠信而又精明的管家预备了极大的赏报。\par

\SourceRecord{Translated by Blomfield Jackson.；Nicene and Post-Nicene Fathers, Second Series；Vol. 8.；Edited by Philip Schaff and Henry Wace.；Buffalo, NY: Christian Literature Publishing Co.,；1895.；<http://www.newadvent.org/fathers/3202206.htm>.}

% structure: p0001 source=h1 target=section reason=page-title

\section{书信二〇七}

\Emph{\Person{圣巴西略} \Place{凯撒勒雅}}\par

\Emph{致\Place{新凯撒勒雅}的神职人员。}\par

你们众人都一致憎恨我。你们全体跟随着那攻击我的首领。因此我原打算不向任何人说一句话。我决定不写任何友爱的信函，不发起任何通讯，只将自己的忧伤默然藏在心里。然而面对诽谤保持沉默是不对的；并非为了借反驳而为自己辩护，而是为了不让谎言继续传播，不让其受害者受到伤害。所以我认为有必要将这件事也呈在你们面前，并写信给你们；尽管我最近曾给全体司铎团写过一封公函，你们却不屑于给我回信。我的弟兄们，请不要满足那些以有害思想充塞你们心灵者的虚荣。不要同意袖手旁观，当你们知道天主子民正被不虔敬的教导颠覆时。除了利比亚人撒伯里乌和加拉提雅人玛尔凯禄，没有人胆敢教导并写出你们民众的首领试图在你们中间作为他们自己的私人发现而提出的东西。他们大肆谈论，却完全无力给他们的诡辩和谬误哪怕一丝真理的色彩。在他们的长篇攻击中，他们对我不惜任何邪恶，并固执地拒绝与我见面。为什么？岂非因为他们害怕因自己的邪恶意见而被定罪吗？是的；而且在攻击我时，他们已变得如此无耻，竟编造某些梦来诋毁我，同时妄称我的教导是有害的。让他们把秋月的所有幻象都揽到自己头上吧；他们不能将任何亵渎加于我身上，因为在每一个教会中都有许多人为真理作证。\par

当被问及这场狂暴无休止的战争的原因时，他们举出圣咏和一种与你们中盛行的习俗不同的音乐，以及诸如此类他们本应羞愧的借口。此外，我们被指控因为坚持那些弃绝世俗以及主比作荆棘、不让圣言结出果实的一切生活挂虑、修行真正宗教的人。这种人身上带着耶稣的死，背起了自己的十字架，跟随了天主。如果这些真是我的过错，如果我有选择这种苦修生活、认我为师的人与我在一起，我甘愿舍弃生命。我听说这种美德如今在埃及可以找到，或许在巴勒斯坦也有一些人的言谈遵循福音的训导。我还听说在美索不达米亚可以找到一些成全有福的人。与成全者相比，我们是孩子。但如果妇女也选择了按福音生活，宁愿童贞胜过婚姻，俘虏了肉体的欲望，并生活在被称为有福的哀恸中，那么无论她们身在何处，她们都是蒙福的。然而我们很少有这方面的例子，因为在我们这里，人们仍处于初级阶段，正逐渐被引向虔诚。如果有人对我们妇女的生活提出混乱的指控，我不为之辩护。但我只说一点：这些大胆的心，这些无缰的嘴，一直无畏地吐露着谎言之父撒殚至今未能说出的话。我希望你们知道，我们欣喜地拥有男女集会，他们的生活是在天上，他们已经把肉体连同肉体的私欲和情欲钉在十字架上；他们不为食物和衣服忧虑，而是宁静地守在主身边，日夜不断祈祷。他们的口不说人的作为，而是不断向天主唱赞美诗，亲手做工，以便能够分施给有需要的人。\par

至于关于唱圣咏的指控——我的毁谤者尤其以此惊吓较单纯的人——我的答复是这样。现今流行的习俗与天主所有教会的习俗一致。在我们中间，人们在夜间前往祈祷之所，在忧愁、苦难和不断的眼泪中向天主告解，最后从祈祷中起身，开始唱圣咏。然后分为两部分，彼此对唱，这样既巩固他们对福音的学习，同时也为他们自己产生一种警醒的心态和免于分心的心灵。然后他们又把曲调的前奏交给一个人，其余人接续；这样在种种圣咏中度过后半夜，当黎明渐临时，大家不时祈祷，然后一齐如同以同一声音、同一心灵向主高唱忏悔圣咏，各人为自己形成忏悔的表达。如果因为这些原因你们与我断绝关系，你们也将断绝与埃及人的关系；你们将断绝与利比亚人、底比斯人、巴勒斯坦人、阿拉伯人、腓尼基人、叙利亚人、幼发拉底河居民的关系；总而言之，与所有那些看重守夜、祈祷和共同唱圣咏的人的关系。\par

4. 但是，有人声称，这些做法在伟大的\Person{额我略}时代并未被遵守。我的反驳是：就连你们现在使用的连祷，在他时代也没有被使用。我这样说并非要指责你们；因为我祈祷你们每个人都生活在眼泪和不断的补赎中。至于我们，我们常为自己的罪过献上恳求，但我们平息我们天主的愤怒，不是像你们那样用人的言语，而是用圣神的话语。你们有什么证据证明这个习俗在伟大的\Person{额我略}时代没有被遵行呢？你们直到现在也没有遵守他的任何习俗。\Person{额我略}祈祷时不蒙头。他怎么能呢？他是宗徒的真门徒，宗徒说：\BibleQuote{凡男人祈祷或说预言，若蒙着头，就是羞辱自己的头。}\BibleRef{格前 11:4}又说：\BibleQuote{男人本不该蒙头，因为他是天主的肖像。}\BibleRef{格前 11:7}\Person{额我略}，那纯洁的灵魂，堪当与圣神共融，他远离誓愿，满足于是或否，遵从了主的诫命，主说：\BibleQuote{我告诉你们：总不可发誓。}\BibleRef{玛 5:34}\Person{额我略}不能忍受称呼他的兄弟为愚人，因为他敬畏主的威胁。激情、愤怒和苦涩从未从他口中发出。他痛恨毁谤，因为毁谤不能进入天国。嫉妒和傲慢被摒除在那无罪的灵魂之外。他在与兄弟和好之前，绝不会站在祭台前。谎言或任何诽谤他人的言语，他都深恶痛绝，因为他知道谎言来自魔鬼，主必要毁灭一切说谎者。如果你们没有这些，而且完全清白，那么你们确实是主的门徒的门徒；如果不是，那么要当心，恐怕你们在争论怎样唱圣咏时，却滤出蠓虫，而忽略了最大的诫命。\par

我被迫说出这些话，是由于我辩护的迫切性，好使你们学会先除去自己眼中的梁木，然后才去移除别人眼中的木屑。尽管如此，我仍然让步一切，尽管在天主面前没有不被审查的事。但愿重要的事占上风，不要让信仰中的创新被人听见。不要忽视\OriginalTerm{位格}{hypostases}。不要否认基督的名字。不要曲解\Person{额我略}的话。如果你们这样做，只要我还有气息、还有说话的能力，当我看见灵魂如此被毁灭时，我不能保持沉默。\par

\SourceRecord{Translated by Blomfield Jackson.；Nicene and Post-Nicene Fathers, Second Series；Vol. 8.；Edited by Philip Schaff and Henry Wace.；Buffalo, NY: Christian Literature Publishing Co.,；1895.；<http://www.newadvent.org/fathers/3202207.htm>.}

% structure: p0001 source=h1 target=section reason=page-title

\section{书信第二〇八}

\Emph{圣凯撒勒雅的巴西略}\par

\Emph{致\Person{尤兰修斯}}\par

长久以来你都保持沉默，尽管你拥有很强的口才，并精于谈话的艺术以及以雄辩展示自己。或许是\Place{新凯撒勒雅}导致你未写信给我。我想我必须将此视为一种善意，如果那里的人不记得我，因为据我所得知的消息，那些听闻者所报告的提及我的言辞并不友善。然而，你过去本是那种因我而被厌恶的人，而非那种因他人而厌恶我的人。我希望这一描述将继续适用于你，无论你在何处，你都会写信给我，并对我怀有善意，只要你多少关心公平与正当。公平地讲，那些率先表达感情的人理应得到同样的回报。\par

\SourceRecord{Translated by Blomfield Jackson.；Nicene and Post-Nicene Fathers, Second Series；Vol. 8.；Edited by Philip Schaff and Henry Wace.；Buffalo, NY: Christian Literature Publishing Co.,；1895.；<http://www.newadvent.org/fathers/3202208.htm>.}

% structure: p0001 source=h1 target=section reason=page-title

\section{第209号书信}

\Emph{\Place{凯撒勒雅}的圣\Person{巴西略}}\par

\Emph{无收信人抬头}。\par

分担我的忧苦，并为我作战，乃是你的本分，这正是你刚毅的明证。安排我们生命的天主，将更大的扬名机会赐给那些能承受大争战的人。你确实为朋友的缘故，像金子经火一般，以自己的生命为英勇的试验。我祈求天主，使他人得以改善，而你则保持现状；并且请你一如既往地继续责备我，指责我不常给你写信——这对我而言是一种过错，对你则造成极大的伤害。这种指责唯有出自朋友之口。请你坚持索取这类债务的偿还。我并非在回报情感的要求上过于不讲理。\par

\SourceRecord{Translated by Blomfield Jackson.；Nicene and Post-Nicene Fathers, Second Series；Vol. 8.；Edited by Philip Schaff and Henry Wace.；Buffalo, NY: Christian Literature Publishing Co.,；1895.；<http://www.newadvent.org/fathers/3202209.htm>.}

% structure: p0001 source=h1 target=section reason=page-title

\section{\WorkTitle{书信二一〇}}

\Emph{\Person{圣巴西略·凯撒勒雅}}\par

\Emph{致\Place{新凯撒勒雅}显要人士}。\par

我实在没有义务向你们公开我的想法，或陈述我目前在此逗留的原因；我不习惯自我炫耀，此事也不值得公开。我想，我并非在随从自己的意愿；我是在回应你们首领的挑战。我一直努力被人忽视，其热切程度远胜过追逐名声者追求出名。但有人告诉我，你们城中的每个人都听得津津有味，而某些专门造谣的谎话制造者正把我和我的所作所为讲给你们听。因此我认为，我不应该坐视你们受邪恶意图和污秽口舌的教导；我认为我有责任亲自告诉你们我的处境。我自幼就熟悉这个地方，因为我是由祖母在此抚养长大的；我常退隐此处，并在努力逃避公共事务的喧嚣时在此度过多年，因为经验告诉我，此地的宁静和独处有利于严肃的思考。此外，由于我的兄弟们现在住在这里，我欣然退居此处，从缠扰我的诸多辛劳中获得短暂喘息，不是为了以此为据点给他人带来麻烦，而是为了满足我自己的愿望。\par

2. 那么，求助于梦境、雇佣解梦者、让我成为公共宴会上的谈资，又有何必要呢？倘若诽谤是从别处针对我的，我本应请你们来作证以证明我的想法，而现在我请求你们每个人回想从前那段日子：那时你们城邀请我负责青年的教育，你们中显要人士的代表团曾来见我。后来，你们所有人聚集在我周围，你们有什么不愿给的？有什么不愿承诺的？然而，你们未能留住我。既然如此，那个当时不肯听从你们邀请的我，如今怎会不请自来地试图强加于你们呢？那个在你们赞美和钦佩我时躲避你们的我，如今在你们诽谤我的时候，又怎会打算讨好你们呢？绝非如此，诸位；我并非那么廉价。一个理智的人不会登上没有舵手的船，也不会靠近一所由那些坐在舵位上的人正在掀起风暴与风浪的教会。全城一片混乱是谁之过？有人无人追赶却逃跑，有人没有敌人逼近却溜走，所有的巫师和梦者都在挥舞他们的鬼怪？要不然是谁之过？难道每个孩童不知道那是群氓首领们的过错？他们对我怀恨的原因，由我说出是不合礼数的；但你们很容易理解。当苦毒与分裂达到了野蛮的极致，而对原因的说明却毫无根据且荒谬可笑时，那么心理疾病便显而易见：确实对他人的安宁构成威胁，但对患者本人来说却是一场巨大而个人的灾难。他们还有一个可爱之处。他们内心被撕裂、备受煎熬，却又因羞愧而无法坦言。他们的处境不仅可以从他们对我的行为，而且可以从他们其余的行为中看出。如果这尚未为人所知，那倒也不打紧。但他们回避与我交流的真正原因，你们中大多数人可能尚未察觉。请听；我这就告诉你们。\par

3. 在你们中间有一种运动，对信仰有害，对福音教义不忠，对真正伟大的圣额我略以及直到真福穆索纽斯的继任者的传统也不忠，而穆索纽斯的教导仍在你们耳中回响。因为那些人，出于对反驳的恐惧，正在捏造针对我的谎言，他们企图重兴撒伯流那早已开始、又被伟大的圣额我略的传统所扑灭的古老祸害。但是，请你们告别那些醉醺醺的脑袋，他们被狂饮升腾的烟雾弄得昏昏沉沉，而我这清醒的人，出于对天主的敬畏无法保持沉默，请听我告诉你们你们中间蔓延的灾祸。\OriginalTerm{撒伯流异端}{Sabellianism}不过是犹太主义伪装成基督教引入福音宣讲之中。因为如果有人称圣父、圣子和圣神是一个具有许多面孔的事物，并使三位的\OriginalTerm{位格}{hypostasis}成为一个，那岂不是否认\OriginalTerm{独生子}{Only-begotten}的永恒先在？他也否认了主在世上的\OriginalTerm{降生成人}{incarnation}、下降阴府、复活、审判；他也否认了圣神的独特运作。我听说，在你们中间，有人甚至敢做出比愚蠢的撒伯流更大胆的创新。据说，而且有亲耳听见的人作证，你们那些聪明人竟极端到说，没有独生子名字的传统，却有仇敌名字的传统；对此他们非常高兴和得意，仿佛这是他们自己的发现。因为经上说：\Quote{我因我父的名而来，你们却不接待我；若有人因自己的名而来，你们倒要接待他。}又因经上说：\Quote{你们要去使万民成为门徒，因父及子及圣神之名给他们授洗。}（见：\BibleRef{玛 28:19}）显然，他们辩称，名字是一个，因为经文不是说\Quote{因\Emph{诸}名}，而是说\Quote{因\Emph{那}名}。\par

4. 我这样写信给你们，实在感到羞愧，因为这些犯罪的是我的骨肉之亲；我也为自己的灵魂叹息，好比拳击手同时与两人搏斗，我只能借着证据击中并击倒左右两边的错误道理，才能使真理发挥应有的力量。一方是阿诺米尤斯派攻击我，另一方是撒伯里乌斯派攻击我。我恳求你们，不要理会这些可恶而无力的诡辩。要知道，那超乎万名之上的名字，就是基督被称为天主之子，正如伯多禄所说：\Quote{除他以外，天下没有赐下别的名字，使我们赖以得救的。}\BibleRef{宗 4:12}至于\Quote{我因我父的名而来}这句话，应当理解为他说这话，是表明父是他自己的根源和原因。如果又说\Quote{你们要去，因父及子及圣神之名给他们授洗}，我们不可以为这里只传给我们一个名字。因为正如那说\Quote{保禄、息耳瓦诺、弟茂德}的人提到了三个名字，并用\Quote{和}字将它们彼此连接，同样，那提到父、子、圣神之名的人，也提到了三个名字，并用连词将它们连接，教导我们每一个名字都应当理解其固有的含义；因为名字代表事物。凡是具有哪怕最小一点理智的人，都不会怀疑事物所具有的存在是独特而自足的。因为父、子、圣神是同一天主性，同一个本性；但这些是不同的名字，向我们展示含义的界限和精确性。因为除非每个名字的特质含义不被混淆，否则无法恰当地向父、子、圣神献上赞颂。\par

如果他们否认他们曾如此说、如此教导，那我的目的就达到了。然而我发现，这样的否认并不容易，因为我们有许多证人听到了这些话。但过去的事就让它过去吧；只要他们现在正确就行。如果他们仍坚持同样的旧错误，我就不得不将你们的灾祸通告其他教会，并请更多的主教给你们写信。在我努力粉碎这日益悄悄增长的不敬神的大山时，我或者会对我的目标有所成效；或者至少我现在的见证会在审判之日洗清我的罪责。\par

5. 他们已经把这些话写进了自己的著作中。他们先寄给了天主的仆人默肋提主教，在收到他适当的答复后，就像畸形怪物的母亲，因自己天生的畸形而羞愧，这些人就在适当的黑暗中生出并养育他们那可憎的产物。他们还通过书信企图试探我亲爱的安提默，提亚纳的主教，理由是额我略在信仰解说中说：父与子思想上为二，但位格上为一。这些自诩智力敏锐的人竟未能看出，这话并非针对教义主张，而是为了反驳埃利安。在这场辩论中，抄写员的错误不少，若蒙天主恩准，我将就实际使用的措辞加以说明。但额我略在说服外邦人时，认为自己无需对所用的措辞过于挑剔；他有时明智地根据被说服者的特点做出让步，以免错失良机。因此，你们会发现其中有许多现在为异端提供极大支持的语句，例如\Quote{受造物}和\Quote{受造之物}之类。但那些无知地批评这些著作的人，将许多针对与人结合而说的话，引到天主性的问题上；正如他们到处兜售的这段经文。因为必须清楚理解：正如不承认本质或实体的共同性就会陷入多神论，同样，拒绝承认位格的区别就会陷入犹太主义。因为我们必须将我们的思想，可以说，保持在某个潜在的主题上，通过形成对其区别特征的清晰印象，从而达至所欲的目标。因为如果我们不想想父的身份，也不记住这个特质所区分的那位，我们怎能理解天主父的概念？仅仅列举位格的区别是不够的；我们必须承认每个位格都具有真实位格中的自然存在。撒伯里乌斯甚至不反对形成没有位格的位格，他说同一的天主，在质料上是一，却随需要而变形，时而被称为父，时而被称为子，时而又被称为圣神。这个无名异端的发明者，正在复活那早已熄灭的旧错误；我说的就是那些摒弃位格、否认天主之子名号的人。他们必得停止对天主说出不义，否则他们将与那些否认基督的人一同哀哭。\par

6. 我深感有必要以这些话写信给你们，好使你们提防不良教导所引起的祸害。如果我们确可把有害的教导比作毒药，如同你们的解梦者所说，这些道理就是毒参和乌头，或任何其他致人于死命之物。正是这些东西毁灭灵魂，而非我的话语，如同这些尖声叫嚷的醉醺醺的渣滓——充满了他们境遇的幻想——所声称的那样。他们若有理智，就该知道：在纯洁、摆脱了一切污秽的灵魂中，先知恩赐明亮地闪耀。在污秽的镜子里，你无法看清映像是什么；同样，一个被此生的忧虑所占据、又被肉欲情欲所昏蒙的灵魂，也不能接受圣神的光辉。并不是每个梦都是预言，正如匝加利亚所说：\BibleQuote{主要使云彩明亮，赐给他们阵雨，……因为偶像说了虚妄，占卜者讲了假梦。}\BibleRef{匝 10:1}那些如依撒意亚所说，做梦并贪睡在床上的人，忘记了有一种错误的作用被赐给了\BibleQuote{悖逆之子}。\BibleRef{弗 2:2}并且有一个说谎的灵，在假预言中兴起，欺骗了阿哈布。\BibleRef{列上 22:22}知道这事，他们就不该如此自高，以致将预言的恩赐归于自己。他们被显明甚至远不及先见者巴郎的情形；因为巴郎受摩阿布王以重金厚礼邀请时，尚且不敢说出任何违背天主意愿的话，也不愿诅咒主未曾诅咒的以色列。\BibleRef{户 22:11}如果他们的睡梦幻想与主的诫命不符，那就让他们满足于福音好了。福音无需梦来提高其可信度。主已将祂的平安赐给我们，并给我们留下一条新诫命，就是彼此相爱；但梦却带来纷争、分裂和爱的毁灭。因此，不要给魔鬼趁他们睡觉时攻击他们灵魂的机会；也不要使他们的想象具有比救恩的训导更大的权威。\par

\SourceRecord{Translated by Blomfield Jackson.；Nicene and Post-Nicene Fathers, Second Series；Vol. 8.；Edited by Philip Schaff and Henry Wace.；Buffalo, NY: Christian Literature Publishing Co.,；1895.；<http://www.newadvent.org/fathers/3202210.htm>.}

% structure: p0001 source=h1 target=section reason=page-title

\section{第211封书信}

\Emph{\Person{圣巴西略·凯撒里亚}}\par

\Emph{致\Person{奥林匹乌斯}。}\par

确实，当我读到阁下之信时，我感到异乎寻常的喜悦与欢愉；当我遇到您可爱的儿子们时，我仿佛看到了您本人。他们见我处于极大的痛苦中，但他们的举止使我忘记了毒芹——那是您那些梦想家和兜售梦想者为了讨好雇用他们的人而四处散布以伤害我的。有些信我已寄出；其余的，若您愿意，随后送上。我只希望它们对收信人有所裨益。\par

\SourceRecord{Translated by Blomfield Jackson.；Nicene and Post-Nicene Fathers, Second Series；Vol. 8.；Edited by Philip Schaff and Henry Wace.；Buffalo, NY: Christian Literature Publishing Co.,；1895.；<http://www.newadvent.org/fathers/3202211.htm>.}

% structure: p0001 source=h1 target=section reason=page-title

\section{书信第二一二号}

\Emph{\Person{圣巴西略}}\par

\Emph{致\Person{希拉略}}。\par

1. 你可以想象我的感受，以及当我来到\Place{达济蒙}却发现你在我到达前几日离开时的心情。自幼年起，我就钦佩你，因此自我们学生时代以来，一直珍视与你的交往。然而，我这样做的另一个原因是，如今没有什么比一颗热爱真理、在实务中具有健全判断的灵魂更为宝贵。我认为，这一点能在你身上找到。我看到大多数人，如同在赛马场上，分成派别，有的支持一方，有的支持另一方，并随其党派呐喊。而你超越恐惧、谄媚以及一切卑劣的情感，因此自然而然地以不偏不倚的眼光注视真理。我也看到你对教会的事务深感兴趣，关于这些事，你在最后一封来信中说过曾给我寄过一封信。我想知道是谁负责传递那封较早的信，以便我知道是谁因丢失信件而损害了我。至今我尚未收到你关于此事的任何信件。\par

2. 那么，我多么希望能见到你，好能向你倾诉我所有的烦恼！当一个人痛苦时，如你所知，即使是描述它也是一种缓解。我多么乐意回答你的问题，不依靠无生命的信件，而是亲自向你逐一叙述。活泼言语的说服力更为有效，也不像信件那样容易受到攻击和歪曲。因为如今没有人留下任何未尝试的手段，连那些我最信任的人——当我看到他们在人群中时，我曾认为他们超乎常人——也从某人那里收到文件，并将它们当作我的文件发送出去，不管它们是什么，并以此在弟兄们面前诽谤我，好像现在虔诚的信徒所应憎恶的没有什么比我的名字更甚。从一开始，我的目标就是默默无闻地生活，无人能及任何考虑人性软弱的人所达到的程度；但现在，仿佛相反地，我的目的就是让自己举世闻名，我已被全世界谈论，我还可以说全海也如此。因为那些走到不敬虔最极端的人，并在教会中引入\OriginalTerm{不相似}{Unlikeness}的无神论观点的人，正在向我开战。那些持守他们所认为的\OriginalTerm{中庸之道}{via media}的人，虽然从相同原则出发，却不遵循其逻辑结论，因为他们与多数人的观点如此对立，也同等敌视我，尽力以他们的责难压倒我，并对我进行各种暗算。但是主已使他们的努力归于徒然。\par

这难道不是一种令人悲痛的情形吗？难道这不会使我的生活痛苦吗？无论如何，在我的烦恼中，我有一个安慰，就是我身体的软弱。我确信，这不会使我在这悲惨的生命中存留太久。这一点不再多言。我也劝你，在你身体的软弱中，要勇敢地、相称于召叫我们的天主生活。如果祂看到我们以感恩之心接受我们目前的处境，祂要么会像对待\Person{约伯}那样除去我们的烦恼，要么会在来世以忍耐的荣耀冠冕赏报我们。\par

\SourceRecord{Translated by Blomfield Jackson.；Nicene and Post-Nicene Fathers, Second Series；Vol. 8.；Edited by Philip Schaff and Henry Wace.；Buffalo, NY: Christian Literature Publishing Co.,；1895.；<http://www.newadvent.org/fathers/3202212.htm>.}

% structure: p0001 source=h1 target=section reason=page-title

\section{\WorkTitle{书信二一三}}

\Person{圣巴西略·凯撒勒雅}\par

\Emph{无收信人}。\par

1. 愿那在我苦难中速速救助我的主，以你们写信给我所给予的慰藉之帮助回报你们，以真实的、来自圣神的大喜乐酬谢你们对我卑微之人的安慰。因为我确实心中沮丧，当我看见在众多人群中那近乎愚蠢且不理智的麻木不仁，以及他们的领袖那根深蒂固、难以根除的不满足。但我看到了你的信；我看到了其中所包含的爱德宝藏；于是我明白，那安排我们一切生活的主，在我生命的苦涩中使我得享一些甜蜜的慰藉。因此我回礼问候你的圣德，并如我所惯常那样劝勉你，不要停止为我这不幸的生命祈祷，好使我永不沉溺于此世的虚幻，忘记那\BibleQuote{从尘埃中提拔穷人的}天主；好使我永不因骄傲而自高自大，陷入魔鬼的审判；好使我永不为主发现我疏忽管家的职分而沉睡；也永不误尽本分，伤害同仆的良心；且永不与醉酒的人为伍，在天主对恶仆的公正审判中承受所警告的惩罚。因此我恳求你，在你的一切祈祷中，祈求天主使我事事警醒；使我永不因我内心的秘密在吾主\Person{耶稣基督}显现的伟大日子被揭露，而成为基督之名的羞辱或耻辱。\par

2. 你要知道，我预计会因异端者的恶行，以和平的名义被传唤到朝廷。也请知道，收到此消息后，这位主教写信给我，要我赶往\Place{美索不达米亚}，并在召集了当地与我们同心同德、巩固教会状况的人之后，与他们一同前往皇帝那里。但也许我的健康状况不足以让我在冬天出行。事实上，迄今为止我并不认为此事紧迫，除非你建议。因此我将等候你的意见，以便我下定决心。那么，我恳请你，不要迟延，立即通过我们可信任的弟兄之一，告知我你这由天主引导的智慧阁下认为最佳的做法。\par

\SourceRecord{Translated by Blomfield Jackson.；Nicene and Post-Nicene Fathers, Second Series；Vol. 8.；Edited by Philip Schaff and Henry Wace.；Buffalo, NY: Christian Literature Publishing Co.,；1895.；<http://www.newadvent.org/fathers/3202213.htm>.}

% structure: p0001 source=h1 target=section reason=page-title

\section{\WorkTitle{第二一四封信}}

\Emph{\Person{凯撒勒雅的圣巴西略}}\par

\Emph{致\Person{特伦提乌斯}伯爵。}\par

当我听说大人您再次被迫参与公务时（实情必须说出来），我立即感到难过，因为想到这必定与您的意愿多么相悖：您在放弃了官场焦虑、让自己有空照顾灵魂之后，竟又被逼回到旧生涯。但后来我又想，或许主这样安排，是愿意您再次公开露面，以便赐予这唯一的缓解，以解目前困扰我们这一带教会的无数痛苦。此外，想到我将在离世之前再次见到大人，我也感到欣慰。\par

然而又有传闻传到我这，说您正在安提约基雅，与主要官员处理手头的事务。除此之外，我还听说属于保利诺一方的弟兄们正就与我们联合的问题与大人进行某些讨论；我所说的\Quote{我们}，是指那位天主的人、真福梅里提乌斯的支持者。我还听说，保利诺派带着一封西方人的信，将安提约基雅教会的主教职位指派给他们，但信中对天主真正教会的那位卓越主教梅里提乌斯持有错误印象。对此我并不惊讶。他们完全不了解此地发生的事；对方虽然可能自以为了解，却向西方人作了以派别置于真理之上的报告。他们要么不知真相，要么甚至故意隐瞒真福主教阿塔纳修写信给保利诺的原因，这也是可以预料的。但大人您当地有能够准确告诉您约维安在位时主教们之间发生了什么的人，我恳求您向他们获取信息。我不控告任何人；我祈求对众人怀有爱德，\InnerQuote{尤其是对信徒一家的人}（\BibleRef{迦 6:10}）；因此我为那些收到罗马来信的人感到高兴。尽管这是对他们的盛大见证，我只希望它是真实的，并经事实证实。但我绝不会因此而说服自己忽视梅里提乌斯，忘记他领导下的教会，或视那些引起分裂的问题为微小、对真宗教无关紧要。我决不会仅仅因为某人因收到人的来信而非常得意就让步。即使那信自天而降，但他若不同意健全的信仰之道，我就不能视他与圣人们共融。\par

我卓越的朋友，请好好想一想：那些篡改真理、将亚流派分裂引入教父健全信仰中的人，之所以拒绝接受教父虔诚的意见，无非是因为他们恶性地理解\OriginalTerm{同体}{homoousion}的意思，并诽谤整个信仰，声称我们主张圣子在\OriginalTerm{位格}{hypostasis}上是\OriginalTerm{同质}{consubstantial}的。如果我们因被那些提出这些言论（及其类似言论）的人——与其说是出于恶意，不如说是出于幼稚——所左右而给他们任何机会，那就没有什么能阻止他们为自己提供反驳我们的无可辩驳的论据，并确认那些人的异端，这些人所有关于教会的言论，其目的与其说是建立自己的立场，不如说是诽谤我的立场。还有什么诽谤比这更严重？还有什么比这更能搅扰大多数人的信仰，即我们中有些人竟被证明宣称圣父、圣子、圣神只有一个\OriginalTerm{位格}{hypostasis}？我们明确设定位格是有区别的；但这一说法已被\Person{撒伯里乌}预见到，他断言天主在\OriginalTerm{位格}{hypostasis}上是一，但圣经以不同的位格描述祂，按照每个具体情形的需要：有时以圣父的名字，当需要这一位格时；有时以圣子的名字，当下降到人类利益或任何\OriginalTerm{救恩工程}{Œconomy}的行动时；有时以圣神的位格，当情形要求这样措辞时。那么，如果我们中有人被证明宣称圣父、圣子、圣神在实体上是同一，而我们却主张三个完满的位格，我们怎能避免清楚地、无可辩驳地证明关于我们的断言的真实性？\par

4. 我认为，\OriginalTerm{位格}{hypostasis}与\OriginalTerm{本体}{ousia}的非同一性，甚至由我们西方的弟兄所暗示，他们因怀疑自己语言的不足，而将\OriginalTerm{本体}{ousia}一词保留为希腊文，以便任何可能的含义差异都能在清晰而不混淆的术语区分中得以保存。若你问我简要陈述自己的见解，我将陈述：\OriginalTerm{本体}{ousia}与位格的关系，正如普遍与特殊的关系。我们每一个人都因共有的\Emph{\OriginalTerm{本质}{essence}}（\Emph{\OriginalTerm{本体}{ousia}}）一词而共享存在，又因自身的属性而成为如此这般的人。同样，在所论之事中，\Emph{\OriginalTerm{本体}{ousia}}一词是共通的，如良善、天主性或其他类似属性；而位格则见于父职、子职或圣化能力的特殊属性中。若他们描述三位格为无位格，则此说法\Emph{本身}是荒谬的；但若他们承认三位格存在于真实的位格中，正如他们所承认的，就让他们如此计算，使\Emph{\OriginalTerm{同本体}{homoousion}}的原则在天主性的统一中得以保存，并且所宣讲的道理成为对真宗教的承认，即在所言及的每一位格的完全而圆满的位格中承认父、子、圣神。尽管如此，有一件事我愿向阁下强调：您及所有像您一样关心真理、并尊重那为真宗教事业而征战的人，应当等待教会的首要权威们率先采取行动来实现这一合一与和平；我将他们视为真理和教会的柱石和基础，并因他们被遣去受罚、被流放远离家乡而更加尊敬他们。我恳求您，使自己免于偏见，使我在您身上——天主赐予我作我万事中的杖和扶持的您——能得安息。\par

\SourceRecord{Translated by Blomfield Jackson.；Nicene and Post-Nicene Fathers, Second Series；Vol. 8.；Edited by Philip Schaff and Henry Wace.；Buffalo, NY: Christian Literature Publishing Co.,；1895.；<http://www.newadvent.org/fathers/3202214.htm>.}

% structure: p0001 source=h1 target=section reason=page-title

\section{\WorkTitle{书信第二一五封}}

\Person{圣巴西略·凯撒勒雅}\par

致司铎\Person{多罗特}。\par

我趁最早的机会写信给最令人钦佩的伯爵\Person{特伦提乌斯}，认为通过陌生人就当前事宜写信给他更好，且担心我们非常亲爱的兄弟\Person{阿加修}不会因任何延误而不便。因此我把信交给政府司库，他正乘坐帝国驿马旅行，并嘱咐他先将信呈给您。我无法理解为什么没有人告诉您，通往\Place{罗马}的道路在冬季完全无法通行，因为\Place{君士坦丁堡}与我们本地区之间的乡村满是敌人。如果必须走海路，季节将是有利的；如果确实我主内亲爱的兄弟\Person{额我略}同意此行并承担关于这些事务的使命。就我而言，我不知道谁能与他同行，而且我知道他在教会事务上相当缺乏经验。对于一个性格友善的人，他可能相处得很好并受到尊重，但是，通过一个骄傲自满、因此完全无法从较低立场聆听向他宣讲真理的人，与我兄弟这样完全不知道卑躬屈膝为何物的人之间的交流，能对事业产生什么好处呢？\par

\SourceRecord{Translated by Blomfield Jackson.；Nicene and Post-Nicene Fathers, Second Series；Vol. 8.；Edited by Philip Schaff and Henry Wace.；Buffalo, NY: Christian Literature Publishing Co.,；1895.；<http://www.newadvent.org/fathers/3202215.htm>.}

% structure: p0001 source=h1 target=section reason=page-title

\section{书信216}

\Emph{\Person{凯撒勒雅的圣巴西略}}\par

致\Person{安提约基雅的默肋提乌斯}主教。\par

许多其他旅程使我离家。我曾远至\Place{丕息狄雅}，与丕息狄雅的主教们一同解决有关\Place{依扫利雅}弟兄们的事务。然后我旅行到\Place{本都}，因为\Person{欧斯塔提乌斯}在\Place{达齐蒙}引起了不小的骚动，并导致那里从我们的教会中分离出相当一部分人。我甚至到了我兄弟\Person{伯多禄}的家，由于那里离\Place{新凯撒勒雅}不远，给新凯撒勒雅人带来了相当大的麻烦，也使我受到了许多无礼对待。有些人在无人追赶时却逃跑了。而我被当作是不请自来，仅仅是为了从那里的人那里得到恭维。我一回到家，因恶劣的天气和焦虑患了重病，立即收到一封来自东方的信，告诉我\Person{保利诺}收到了来自西方的一些信函，承认他有某种更高的权利；安提约基雅的叛乱者因此大为得意，接下来他们准备了一份信经，并提出以其条款作为与我们教会联合的条件。除此之外，还报告说他们已将那位最卓越的人\Person{泰伦提乌斯}诱入他们的派别。我立即尽可能有力地写信给他，劝他暂停；并试图指出他们的不诚实。\par

\SourceRecord{Translated by Blomfield Jackson.；Nicene and Post-Nicene Fathers, Second Series；Vol. 8.；Edited by Philip Schaff and Henry Wace.；Buffalo, NY: Christian Literature Publishing Co.,；1895.；<http://www.newadvent.org/fathers/3202216.htm>.}

% structure: p0001 source=h1 target=section reason=page-title

\section{\WorkTitle{书信二一七}}

\Signature{圣巴西略（凯撒勒雅）}\par

\Emph{致安菲洛奇书——《圣教法典》}\par

在我远行归来之后（因我已前往\Place{本都}办理教会事务并探访亲属），身体虚弱多病，精神颇受打击，我拿起尊座的信函。一收到那对我而言最为甜美之声的迹象和那亲爱之手的手迹，我立刻忘记了所有烦恼。若我因收到您的信函而如此愉悦，您应能猜到我何等珍视您的亲自临在。愿圣者赐我此恩，待您方便并亲自邀请我之时。若您前来\Place{优菲米亚}的住所，我将乐于与您相聚，摆脱我在此地的烦恼，奔向您真诚的情谊。或许还有其他原因，因至受天主钟爱的\Person{格列高利}主教的突然离去，我可能被迫远赴\Place{纳齐盎}。此事如何发生、为何发生，我至今未获消息。关于我曾向阁下提及的那人，原以为此刻已准备就绪，须知他患了慢性病，且现因旧疾及新患眼部疾病而受苦，完全无法工作。我身边别无他人。因此，尽管事情已由他们交由我处理，更好是由他们中推举一人。事实上，不得不认为那些话只是必要的形式，他们真正所愿是他们最初请求的——被选为领导者的人应出自他们中间。若最近受洗者中有任何人，无论\Person{马其顿尼}是否赞同，就让他被委任。您将教导其职责，主在万事中与您合作，也为此工作赐予您恩宠。\par

第五十一。关于圣职人员，圣教法典毫不区分地规定：凡失足者，一律处以撤职，不论他们属圣秩品，或留在不经按手授予的职务上。\par

第五十二。凡生育婴儿并将其弃于路旁的女人，如果她能救活却疏忽，或想借此掩盖自己的罪过，或出于野蛮而非人道的动机，则应按谋杀案审判。反之，若她无力抚养，婴儿因暴露和缺乏生活必需品而死亡，则母亲应得宽恕。\par

第五十三。寡妇身份的奴隶，若在强奸的借口下再婚，则不犯严重跌倒。她不应因此被指控。更需考察的是目的而非借口。但显然，她应受重婚的处罚。\par

第五十四。我知道我已就非故意杀人案中应遵守的区分，尽我所知致信尊座。关于这一点，我不能再多说。加重或减轻惩罚的严厉程度，取决于尊座的明智，视个别案例需要而定。\par

第五十五。袭击强盗者，若为教外人，则被禁止领受善物的共融。若为圣职人员，则被剥夺圣秩。因为经上说：\Quote{凡持剑的，必死于剑下。}\BibleRef{玛 26:52}\par

第五十六。故意杀人者，若事后悔改，将被停止参与圣事二十年。二十年分配如下：前四年哭泣，站在祈祷堂门外，恳求进堂的信友为他祈祷，并承认自己的罪过。四年后，他被允许列于听道者，并随他们出去五年。七年中随同跪伏者祈祷。四年中仅与信友并立，不参与奉献。期满后，方可参与圣事。\par

第五十七。非故意杀人者，将被停止参与圣事十年。十年分配如下：前二年哭泣，三年中列于听道者，四年为跪伏者，一年仅与信友并立。然后方可参与圣事。\par

第五十八。犯奸淫者，将被停止参与圣事十五年。四年哭泣，五年听道，四年跪伏，二年并立而不领共融。\par

第五十九。行淫者，七年内不得参与圣事：哭泣二年，听道二年，跪伏二年，并立一年；第八年方可领受共融。\par

第六十。凡宣发誓愿守贞而背誓的女人，应按照奸淫案的期限完成其节欲。对宣示独身生活而后失足的男子，也遵守同样规则。\par

第六十一。偷窃者，若自动悔改并自首，只应禁止领受圣事一年；若被定罪，则两年。期限分为跪伏和并立。然后应被视为堪当领受共融。\par

第六十二。凡对男性行为不检者，应受与奸淫者同样的惩戒。\par

第六十三。凡承认与兽类行淫者，应受同样期限的补赎。\par

第六十四。发虚誓者，应被停止领受圣事十年：哭泣二年，听道三年，跪伏四年，并立一年；然后应被视为堪当领受共融。\par

第六十五。凡承认行巫术或邪术者，应作谋杀期限的补赎，并应与自认此罪者同样处理。\par

第六十六。盗墓者，应被停止领受圣事十年：哭泣二年，听道三年，跪伏四年，并立一年；然后方可被接纳。\par

第六十七。与姊妹乱伦者，应受与谋杀同样的补赎。\par

第六十八。亲属间在禁止通婚等次内结合者，若发现已犯奸淫行为，应受奸淫的惩罚。\par

LXIX. 凡在婚前与未婚妻行房的读经员，应停职一年后准予读经，但不得晋升。若未订婚而秘密行房，则应撤去其职务。辅礼员亦同。\par

LXX. 凡在口舌上被玷污的执事，若已告明此罪，应从其职务上撤除。但应准许他与执事们一同领受圣事。司铎亦然。若有人被查出犯有更严重的罪，无论其品级如何，应予以撤职。\par

LXXI. 凡知晓上述任何一项罪行，且未经告明而被定罪者，应与实际犯罪者受同等期限的惩罚。\par

LXXII. 凡将自己托付于占卜者或任何此类人物者，应与杀人者受同等时期的惩戒。\par

LXXIII. 凡否认基督、并得罪救恩奥迹者，应终身哭泣，并须持守补赎，唯在临终时，因信赖天主的仁慈，堪当领受圣事。\par

LXXIV. 然而，若先前犯罪者因补赎而改过，则蒙天主仁慈授予释放与束缚权柄者，若因看见罪人补赎的极大诚意而减轻刑罚期限，不应受谴责，因为圣经中的历史教导我们：凡以更大热忱行补赎者，必迅速获得天主的仁慈。\par

LXXV. 凡与自己的姐妹（无论是同父或同母所生）行淫者，不得进入祈祷之所，直至他放弃那邪恶不法之行。当他认识到这可怕之罪后，应让他站立在祈祷之所门外哭泣三年，恳求进入祈祷的人们，各人以热忱为他向主仁慈祈祷。此后，再接受他仅能听道三年；当聆听圣经和教训时，应将他逐出，不准他参与祈祷。之后，若他流泪祈求，并怀着痛悔之心和极大的自卑跪伏在主前，再准他跪伏三年。如此，当他堪当显示悔改的果实后，在第十年准他与信友一同祈祷，但不献祭；与信友一同祈祷站立两年后，然后，才可堪当领受善事的共融。\par

LXXVI. 同样规则适用于娶自己儿媳者。\par

LXXVII. 凡遗弃合法结合之妻子者，按主的判决应受奸淫之罚。但我们的教父们已立为教规：此类罪人应哭泣一年，听道两年，跪伏三年，第七年与信友同站；如此，若他们以眼泪悔改，则堪当领受献祭。\par

LXXVIII. 同样规则适用于娶两姐妹者，即便是在不同时期。\par

LXXIX. 凡对继母行淫者，与对姐妹行淫者受同样教规处罚。\par

LXXX. 关于多妻制，教父们保持沉默，因其兽性且完全不人道。此罪在我看来比奸淫更重。因此，此类罪人应受教规处罚，即哭泣一年，跪伏三年，然后接纳。\par

LXXXI. 在蛮族入侵期间，许多人曾发异教誓言，尝过在魔法庙中非法献祭之物，从而背弃了对天主的信仰。这些人应按照我们教父所订的教规予以处理。那些遭受严酷折磨、因无法承受痛苦而被逼否认信仰者，可被排除三年，听道两年，跪伏三年，然后接纳进入共融。那些未受大苦而放弃对天主的信仰，触摸魔鬼的祭桌并发出异教誓言者，应被排除三年，听道两年。当他们跪伏祈祷三年，并与信友一同祈求站立另外三年后，然后接纳他们进入善事的共融。\par

LXXXII. 关于发假誓者，若因暴力胁迫而违背誓言，所受刑罚较轻，因此可在六年后被接纳。若未受胁迫而背信，应让他们哭泣两年，听道三年，跪伏祈祷五年，期间两年被接纳进入祈祷的共融，但不献祭，如此最后，在显示适当悔改的证据后，他们将恢复与基督身体的共融。\par

LXXXIII. 凡请教占卜者、追随异教习俗、或带人进屋以寻求疗法和施行净化者，应受六年教规处罚。哭泣一年，听道一年，跪伏三年，与信友同站一年后，如此方可被接纳。\par

LXXXIV. 我写下这一切，是为了检验悔改的果实。我并非绝对根据时间来决定这类事情，而是注意补赎的方式。如果人们处于一种难以摆脱自己道路的状态，宁愿服侍肉体的享乐而不愿侍奉主，并且拒绝接受福音的生活，那么我和他们之间就没有共同的基础。在悖逆和违抗的人民中，我被告知要听到这些话：\Quote{拯救你自己的灵魂。}那么，我们不可赞同与这样的罪人一同灭亡。让我们畏惧可怕的审判。让我们把主报复的可怕日子放在眼前。我们不可同意在别人的罪中灭亡，因为如果主的恐怖没有教导我们，如果如此巨大的灾祸没有让我们意识到，是因我们的不义，主抛弃了我们，把我们交在蛮族手中，百姓在敌人面前被俘虏，并被分散，因为携带基督之名的人竟敢做出这样的行为；如果他们不知道也不明白，正是因这些原因，天主的愤怒临到我们身上，那么我和他们还有什么共同的论据呢？\par

但我们应当日夜向他们作证，无论是在公开场合还是私下。我们不可同意被他们拖入他们的邪恶中。让我们首先祈祷，使我们能对他们行善，并从恶者的陷阱中拯救他们。如果我们不能做到这一点，至少让我们尽力拯救我们自己的灵魂免于永罚。\par

\SourceRecord{Translated by Blomfield Jackson.；Nicene and Post-Nicene Fathers, Second Series；Vol. 8.；Edited by Philip Schaff and Henry Wace.；Buffalo, NY: Christian Literature Publishing Co.,；1895.；<http://www.newadvent.org/fathers/3202217.htm>.}

% structure: p0001 source=h1 target=section reason=page-title

\section{第二百一十八封信}

\Emph{凯撒勒雅的圣巴西略}\par

致\Place{依科尼雍}主教\Person{安斐罗西}。\par

\Person{厄里亚奴斯}弟兄已亲自完成了前来的事务，并不需要我的帮助。然而，我对他有双重感谢，既因他为我带来阁下的信函，也因他给了我给你写信的机会。因此，藉着他，我问候阁下真诚无伪的爱，并恳求你如今更迫切地为我祈祷，因为我极其需要你祈祷的帮助。我的健康因前往\Place{本都}的行程而严重受损，病情难以忍受。有一件事我长久以来渴望告知你。我并非说我因其他原因而遗忘，但现在我想提醒你派人去\Place{利基亚}，查访谁持有正信，因为或许他们不应被忽视，如果从我那里一位热心旅人带来的报告属实，即他们已完全疏远亚细亚人的观点，并希望与我们共融。若有人前往，让他分别在\Place{科律达拉}询问已故隐修士\Person{亚历山大}主教，在\Place{利米拉}询问\Person{迪迪默}，在\Place{米拉}询问司铎\Person{塔提安}、\Person{颇肋摩}和\Person{玛卡略}，在\Place{帕塔辣}询问\Person{欧德默}主教，在\Place{特尔美索}询问\Person{希拉略}主教，在\Place{菲洛}询问\Person{拉里亚奴}主教。我已得知他们以及更多其他人在信德上是纯正的，并且我感谢天主，即使是在亚细亚地区，竟然有人免于异端祸害。那么，如有可能，在此期间我们亲自查问他们。待获得信息后，我打算写一封信，并渴望邀请其中一人来与我见面。愿天主保佑那令我极其珍爱的\Place{依科尼雍}教会一切顺利。藉着阁下，我问候所有可敬的圣职人员和所有与阁下共事的人。\par

\SourceRecord{Translated by Blomfield Jackson.；Nicene and Post-Nicene Fathers, Second Series；Vol. 8.；Edited by Philip Schaff and Henry Wace.；Buffalo, NY: Christian Literature Publishing Co.,；1895.；<http://www.newadvent.org/fathers/3202218.htm>.}

% structure: p0001 source=h1 target=section reason=page-title

\section{Letter 219}

\Emph{圣巴西略，凯撒勒雅}\par

致萨莫萨塔的圣职人员。\par

主吩咐\Quote{一切事物皆以衡量和重量,}\BibleRef{智 11:20}，并赐予我们不超过承受能力的试探，藉苦难考验所有为真道争战的人，不让他们受试探超过所能忍受的。祂将泪水大量赐给那些应当藉着感情表明他们是否保持对祂感恩的人。尤其在你身上，祂的慈爱彰显在祂的处置中，不容许敌人加诸于你的逼迫大到使你们中有人动摇或偏离对基督的信德。祂使你们与微小且易抵挡的对手相配，并藉着战胜他们为你们的坚忍预备了奖赏。但我们生命的共同仇敌，以其诡计对抗天主的良善，因见你们如坚固的城墙轻视外来的攻击，便如我所闻，设计使你们中间兴起彼此间的冒犯和争执。这些起初微不足道且易治愈，然而随着时间推移，因争论而加剧，往往导致不可补救的祸患。因此，我藉此信劝诫你们。若可能，我愿亲自前来当面恳求你们。但现况阻止了此事，故以书信代替恳求，向你们呈上此信，望你们尊重我的请托，止息彼此间的竞争，并早日将你们中间一切冒犯之事已息的好消息告知于我。\par

2. 我非常渴望你们知道：凡谦卑服从邻人、承受对自己的指责而不觉羞愧，即使它们不实，以将和平的宏大祝福带给天主的教会，此人在天主面前是伟大的。\par

我希望你们中间兴起一种友善的竞争，看谁先配称为天主的儿子，因做和平者而获得这地位。你们极其蒙天主所爱的主教也已致信你们，指示当行之道。他将按其职责再次写信。但我既已被准许靠近你们，也不能忽视你们的处境。所以，当极虔诚的兄弟、副执事狄奥多若到来，报告说你们的教会处于苦恼和骚动中时，我深感忧伤、内心痛苦，不能缄默。我恳求你们抛弃彼此间的一切争端，缔造和平，免得使你们的对手喜悦，并摧毁那如今传遍天下的教会的夸耀，即你们所有人在同一心灵和同一心志的治理下，如同一个身体生活。藉着你们的尊敬，我问候天主的全体子民，包括有职分和权位者，以及其余的圣职人员。我劝勉你们保持旧有的品德。我无法要求比这更多，因为你们善工的彰显已预先做到并使得任何改进成为不可能。\par

\SourceRecord{Translated by Blomfield Jackson.；Nicene and Post-Nicene Fathers, Second Series；Vol. 8.；Edited by Philip Schaff and Henry Wace.；Buffalo, NY: Christian Literature Publishing Co.,；1895.；<http://www.newadvent.org/fathers/3202219.htm>.}

% structure: p0001 source=h1 target=section reason=page-title

\section{\WorkTitle{第二二〇书}}

\Emph{\Person{凯撒勒雅的圣巴西略}}\par

致\Person{贝里亚人}。\par

凡被剥夺了亲自交往之乐的人，主赐予他们通过书信交流的巨大安慰。借由这种方式，我们固然无法得知身体的形像，却能了解灵魂本身的气质。因此，当下我收到诸位尊者的书信时，同时就认出了你们，并把你们对我的爱存于心中，无需长久时间便能与你们亲近。你们书信中表露的气质足以在我心中燃起对你们灵魂之美的爱慕。而除了你们那极佳的书信外，从担任我们之间沟通媒介的众位弟兄的和蔼可亲中，我也获得了关于你们现况的更明确证据。我所敬爱且可敬的司铎\Person{阿卡丘斯}，除你们所写的以外，还向我讲述了许多，并将你们日复一日所持守的争战以及你们为\OriginalTerm{真宗教}{true religion}所做的坚定抵抗呈现在我眼前。他如此打动我的钦佩，并在我心中激起如此热切的愿望，渴望享受你们的美善，以致我恳求主，愿有一天我能亲身认识你们和你们的一切。他向我讲述了你们中受托祭台职务者的严谨，此外还有全体民众的和睦一致，以及你们城中长官和首领人物的慷慨性格和对天主的真诚之爱。我因此祝贺教会拥有这样的成员，并祈求灵性的平安更丰盛地赐予你们，好使在较平静的时期，你们能在患难之日因自己的劳苦而喜乐。因为经历时痛苦的事，往往在回忆时带着愉悦。目前我恳求你们不要灰心。不要因接踵而至的苦难而绝望。你们的冠冕临近了，主的助佑临近了。不要让你们迄今所经受的一切付诸东流，不要使一场闻名于全世界的争战归于无效。人的生命何其短暂。\BibleQuote{凡有血肉的都像草，他的一切荣美都像野地的花……草必枯干，花必凋谢；但我们天主的言语必永远立定。}\BibleRef{依 40:6}让我们紧守那常存的诫命，轻视那消逝的虚幻。许多教会因你们的榜样而振作。你们在召唤新的战士进入战场时，为自己赢得了巨大的赏报，虽然你们自己并不知晓。赏报的赐予者是富足的，祂能不相称地酬报你们的英勇事迹。\par

\SourceRecord{Translated by Blomfield Jackson.；Nicene and Post-Nicene Fathers, Second Series；Vol. 8.；Edited by Philip Schaff and Henry Wace.；Buffalo, NY: Christian Literature Publishing Co.,；1895.；<http://www.newadvent.org/fathers/3202220.htm>.}

% structure: p0001 source=h1 target=section reason=page-title

\section{第二二一封信}

\Person{凯撒勒雅的圣巴西略}\par

致\Place{贝利亚人}。\par

亲爱的朋友们，此前我因你们远近闻名的虔诚，以及你们在基督内宣认所得的荣冠，已认识你们。或许有人会问：这些关于我们的消息是谁传得这么远？是主自己；因为他将敬拜他的人如同灯放在灯台上，使他们照亮全世界。难道竞技的胜利者不是因胜利的奖赏而闻名，工匠不是因作品的巧妙设计而著称吗？这些人和类似之人的记忆岂不永远不被遗忘？敬拜基督的人，主亲自说：\Quote{尊重我的，我必尊重他，}岂不因他而在众人面前成名并受光荣？他岂不将他们的荣耀光辉如同太阳的光线一样彰显出来？然而，由于你们惠然寄给我的信，我更加渴慕你们。在信中，你们为真理所做的努力，已超越以往，对真信仰的热情更加丰沛而强韧。我为此与你们一同喜乐，并同你们祈祷，愿那掌管宇宙、进行竞赛和赐予冠冕的天主，以热忱充满你们，使你们的灵魂强壮，并使你们的工作合乎他神圣的喜悦。\par

\SourceRecord{Translated by Blomfield Jackson.；Nicene and Post-Nicene Fathers, Second Series；Vol. 8.；Edited by Philip Schaff and Henry Wace.；Buffalo, NY: Christian Literature Publishing Co.,；1895.；<http://www.newadvent.org/fathers/3202221.htm>.}

% structure: p0001 source=h1 target=section reason=page-title

\section{第二二二封信}

\Emph{\Person{凯撒利亚的圣巴西略}}\par

\Emph{致\Place{卡尔基斯}居民。}\par

各位可敬者的信函，在我困苦之时来临，犹如在正午赛道中央，气喘吁吁、口吸灰尘的赛马被注入清水一般。虽遭接踵而至的试炼，我重得喘息，既因你们的话语而振奋，又因想到你们为以不屈勇气面对我眼前之事所作的奋斗而增添力量。因为那吞噬了东方大片地区的烈火，已逐渐蔓延至我们邻近之地，烧尽周围一切后，正试图触及\Place{卡帕多细亚}的教会，这些教会已因邻舍家园废墟升起的烟雾而泪流满面。火焰已几乎触及我。愿主以祂口中的气息使它们转向，并遏止这邪恶之火。谁是如此懦弱、如此怯懦、如此未经运动员般拼搏之人，竟不被你们的欢呼所激励，而不在你们身旁祈求被欢呼为胜利者呢？你们率先踏入了真正信仰的竞技场；在与异端者的较量中击退了多次攻击；你们承受了试炼的炽热强风，无论是你们这些教会的领袖，即祭台职务的受托人，还是每一位平信徒，包括那些身份较高者。因为这在你们中间尤其可赞叹、值得一切赞美，就是你们都在主内合而为一，你们中有些人是向善行进的领导者，其他人则心甘情愿地跟随。正是为此，你们对于攻击者的进攻过于强大，不给对手在你们任何肢体中留余地，为此我日夜祈求万代的君王，使百姓保持信仰的完整，并借他们保存圣职人员，如同在顶端完好无损的头颅，为身体每一部分行使它自己的警醒远虑。因为眼睛履行其职，双手方能各尽其工，双脚才能行走无绊，身体没有一部分被剥夺应有的照顾。因此我恳求你们，彼此紧密相连，正如你们所做、也必将做的那样。我恳求你们这些受托照顾灵魂的人，将每一个人都团结在一起，并像爱护心爱的子女一样爱护他们。我恳求百姓继续向你们显示应给予父亲般的尊敬和荣誉，好使你们在教会的美好秩序中，保持力量和在基督内的信仰根基；使天主的名字受光荣，爱德的美好恩赐增长而丰盛。愿我，当听到你们的消息时，为你们在天主内的进展而喜乐。若我仍蒙召在此世肉身寄居，愿我终有一天在平安中见到你们。若我现在就被召唤离世，愿我在圣人们的光辉荣耀中见到你们，与所有那些因忍耐和显出善工而被视为配得的人一起，你们的头上戴着冠冕。\par

\SourceRecord{Translated by Blomfield Jackson.；Nicene and Post-Nicene Fathers, Second Series；Vol. 8.；Edited by Philip Schaff and Henry Wace.；Buffalo, NY: Christian Literature Publishing Co.,；1895.；<http://www.newadvent.org/fathers/3202222.htm>.}

% structure: p0001 source=h1 target=section reason=page-title

\section{第二二三封信}

\Emph{\Person{圣巴西略·凯撒勒雅}}\par

\Emph{\Person{塞巴斯特的欧斯塔基}}。\par

1. 静默有时，言语有时，\BibleRef{训 3:7}是训道者的格言。我静默的时日已经足够，如今时候已到，要开口将迄今未为人知的真相公之于众。卓越的\Person{约伯}长久沉默地忍受着自己的灾难，并以忍受最难以承受的苦难来显示自己的勇气；但当他在沉默中挣扎了足够久，坚持将痛苦深藏在心底之后，他终于开口说出了那些著名的话。而我自己，现在已是沉默的第三年，我的夸耀变得如同圣咏作者那样：\BibleQuote{我像是一个不听闻的人，口中没有辩驳。}因此我将因针对我的诽谤所受的痛苦深藏于心底，因为诽谤使人卑微，诽谤使穷人眩晕。因此，如果诽谤的祸害如此之大，甚至能令完人从其高位坠落——因为圣经以\Quote{人}一词所指的正是如此，而\Quote{穷人}则指缺乏伟大教义者，这也正与先知所说的相符：\BibleQuote{这些是穷人，因此他们不会听；……我要去见那些伟人，}他所谓\Quote{穷人}是指那些缺乏理解的人；同样，这里他明确指那些在内里还未装备好、甚至未达到年龄的完满尺度之人；正是这些人，正如谚语所说，会眩晕并摇摆不定。不过，我曾认为自己应当沉默地忍受烦恼，等待从中出现一些迹象。我甚至不认为那些反对我的话是出自恶意；我以为那是由于对真理的无知。但现在我看到敌意与日俱增，我的诽谤者并不为当初所说的话后悔，也不努力弥补过去，反而继续下去并纠集在一起以实现其最初的目标：就是使我的生活痛苦，并设法在弟兄中玷污我的名誉。因此，我再也看不到沉默中的安全。我想起了\Person{依撒意亚}的话：\BibleQuote{我长久缄默，我是否总是静默并克制自己？我忍耐，如同临产的妇人。}愿天主使我既能获得沉默的赏报，又能获得驳倒对手的力量，如此，通过驳倒他们，我就能干涸那涌向我的谎言苦流。这样我可以说：\BibleQuote{我的灵魂跨越了洪流；}以及，\BibleQuote{若不是上主在我们这边，当人们起来攻击我们时，……他们早已活活吞灭我们，大水早已淹没我们了。}\par

2. 我曾在虚空中度过许多时间，几乎浪费了我整个青年时期在徒劳的劳作上，以获取那为天主所愚弄的智慧。后来有一次，如同从沉睡中醒来的人，我将目光转向福音真理的奇妙光芒，我领悟到\BibleQuote{这世界掌权者的智慧，终归无用的。} \BibleRef{格前 2:6}我为自己的悲惨生活流了许多眼泪，并祈祷天主赐我引导，使我得以进入真宗教的教义。首先，我定意要对我的生活方式作一些修补，它们因我与恶人的亲密往来而长期扭曲。然后我读了福音，我看到那里有一条达到成全的伟大途径：变卖财产，与穷人分享，放弃对此世的一切挂虑，不允许灵魂因任何同情而转向世俗事物。我祈祷能找到一位已经选择了这种生活方式的兄弟，与他一起渡过这短暂而充满风浪的人生海峡。我在\Place{亚历山大城}找到了许多这样的人，在\Place{埃及}其他地方也有许多，还有在\Place{巴勒斯坦}、在\Place{叙利亚}、在\Place{美索不达米亚}。我钦佩他们生活上的节制和劳作中的坚韧；我惊叹于他们祈祷的持久和战胜睡眠的胜利；不被任何自然需要所征服，始终保持灵魂的目标崇高而自由，在饥饿、口渴、寒冷、赤身露体时，\BibleRef{格后 11:27}他们从不向身体屈服；他们从不愿意浪费注意力在身体上；总是仿佛生活在不属于自己的肉身中，他们以实际行动显示了什么是在此生短暂寄居，什么是拥有天国的籍贯和家园。这一切令我赞叹。我称这些人的生活为有福的，因为他们确实以实际行动表明他们\BibleQuote{在身上带着耶稣的死。} \BibleRef{格后 4:10}我也祈祷，在我力所能及的范围内，效法他们。\par

3. 所以，当我看见本国有些人竭力效法他们的生活方式，我便觉得找到了有助于我自身救恩的帮助，并视有形之物为无形之事的证据。由于我们每个人心中的秘密是未知的，我便将衣着的卑微视作心灵谦卑的充分标志；粗布衣、腰带和生皮制成的鞋足以说服我。虽然许多人试图使我远离他们的团体，我却不允许，因为我看出他们将忍受的生活置于享乐之前；并且，由于他们生活的超凡卓越，我成了他们的热心支持者。于是，我听不得任何人指责他们的教义，尽管许多人坚持认为他们对天主的观念是错误的，而且他们已成为当今异端魁首的门徒，并秘密传播他的教导。但我从未亲耳听过这些事，便断定那些报告者是诽谤者。后来我被召管理教会。至于那些以帮助和友爱共融的假象赐予我的看守和窥探者，我闭口不言，以免因讲述难以置信的故事而损害自己的事业，或因被相信而给信徒提供憎恨同伴的机会。这几乎也成为我的情况，若非天主的仁慈阻止了我。因为几乎每个人都成了我怀疑的对象，我因被背信弃义地伤害而心中伤痛，似乎找不到任何值得信赖的人。尽管如此，我们之间仍保持着某种亲密关系。我们曾多次讨论教义问题，表面上似乎意见一致，保持合一。但他们开始发现，我所做的关于我对天主的信德的声明与以往他们从我这里听到的相同。因为，倘若我身上有别的令人叹息之事，我至少敢在\Person{主}内以此夸口：我从未有一刻对天主怀有错误观念，或持有后来改变的异端见解。我从孩提时代从我蒙福的母亲和外祖母\Person{玛克利纳}那里接受的关于天主的教导，我始终以日益坚定的信念持守。到了成熟的理性年龄，我没有改换观点，而是遵循父母传授给我的原则。正如种子生长时，起初微小，然后逐渐变大，但始终保持着同一性，种类不变，只是在成长中逐步完善；同样，我认为同样的教义在我身上通过逐步提升的阶段而成长。我现在所持守的并没有取代我起初所持守的。让他们省察自己的良心吧。让这些如今因伪教导的指控而使我成为众人谈论的话题，并用他们写来攻击我的诽谤信件震聋所有人耳朵的人回忆一下，他们是否曾听过我说过与现在不同的话，并让他们记住\Person{基督}的审判台。\par

4. 我被指控亵渎天主。然而，他们无法依据任何他们用来攻击我的关于信德的论文定我的罪，也无法依据我曾在天主的教会中不时口头宣讲，而未记录为文字的言论来指控我。找不到一个证人说他曾听我在私下谈论过任何违反真宗教的话。那么，如果我不是一个不正统的作者，如果我的宣讲找不出毛病，如果我不误导在我家与我交谈的人，我凭什么受审判？但有一个新发明！有人指控说，在\Place{叙利亚}有某人写了违反真宗教的东西；而二十年前或更久之前，你给他写了一封信：所以你是那家伙的同谋，对他的指控也对你的指控。哦，爱真理的先生，我回答，你曾被教导谎言是魔鬼的产物；是什么向你证明我写了那封信？你从未送交；你从未询问；你也从未从我这里得知，我本可能告诉你真相。但若那封信是我的，你怎么知道现在到你手中的文件与我的信是同一时期的？谁告诉你它已有二十年？你怎么知道那是我去信之人的作品？若他是作者，我写信给他，我的信和他的作品属于同一时期，有什么证据表明我接受它作为我的判断，并持有那些观点？\par

5. 请问你自己。你曾多少次探访我在\Place{伊里斯河}上的隐修院，当时我极其敬爱天主的兄弟\Person{额我略}与我同在，过着和我一样的生活方式？你曾听过这类事吗？有这种事的任何迹象，无论大小吗？我们在河对岸的村庄、在我母亲那里同住过多少天，像朋友一样生活，昼夜交谈？你发现我持有任何这类观点吗？当我们一起去看望蒙福的\Person{息耳瓦诺}时，在路上我们不是谈论了这些事吗？在\Place{尤西诺厄}，当你准备与其他主教一同出发去\Place{兰普萨古斯}时，我们的谈话不正是关于信德吗？在我口授对异端的反驳时，你的速记员不是一直在我身边吗？你最忠实的门徒不也在那里吗？当我探访弟兄们，与他们一同祈祷过夜，不断谈论并聆听关于天主的事而无争议时，我所给出的关于我情感的见证难道不是确切而明确的吗？那么，你怎么将这个腐朽而微弱的猜疑视为比如此长时间的亲身经历更重要呢？\par

你有什么理由宁愿相信自己的判断而不相信我心灵的证据呢？我在\Place{加采东}、又多次在\Place{赫拉克莱亚}、早先在\Place{凯撒勒雅}郊区关于信仰的言论，难道有一丝不和谐之处吗？它们不是彼此完全一致吗？只有一点例外，就是我刚才所说的随着进步而来的力量增强，这不应被视为从坏到好的转变，而是以知识的增加填补了先前的不足。你怎么能忘记：\BibleQuote{父亲不担负儿子的罪恶，儿子也不担负父亲的罪恶，各人因自己的罪恶死亡。}\BibleRef{则 18:20} 我没有父亲也没有儿子受你诽谤；我没有老师也没有门徒。但如果父母的罪恶必须归咎于子女，那么将\Person{亚略}的罪恶归咎于他的门徒就更加合理；无论谁生了异端者\Person{艾提乌}，儿子的罪名都应归于父亲。另一方面，如果说任何人因他们的缘故被指控是不公正的，那么我因那些与我毫无关系的人而承担责任就更不公正了，即使他们在各方面都是罪人，并且他们写下了值得谴责的东西。如果我不完全相信所有对他们提出的指控，请原谅我，因为我的亲身经历表明，控告者多么容易滑入诽谤。\par

6. 即使他们真的前来指控我，因为他们受了欺骗，以为我与他们四处散播的\Person{撒伯里乌}言论的作者有联系，那么他们在没有获得明确的证据之前，在我没有对他们做任何错事的情况下，就立刻攻击和伤害我，这是不可原谅的。我不愿说自己与他们有最亲密的关系；也不愿说他们显然没有受圣神引导，因为他们怀有虚假的猜疑。一个即将与兄弟断绝友谊的人，必须经过许多焦虑的思考，度过许多不眠之夜，流许多眼泪从天主那里寻求真理。就连今世的统治者，当他们即将判处某个作恶者死刑时，也会拉开帷幕，请来专家审查案情，花相当多的时间权衡法律的严苛与人性的共同软弱，并随着许多叹息和许多对严峻必要性的哀叹，向全体人民宣布他们是出于必要而服从法律，而不是为了满足自己的意愿而判刑。那么，一个即将与长期建立的兄弟友谊决裂的人，应该采取多少更大的谨慎和勤奋，多少更多的商议呢！在这种情况下，只有一封信，而且其真实性可疑。不可能辩称它通过签名而被知晓，因为他们没有原件，只有抄本。他们依赖一份单一的文件，而且是一份旧文件。至今已有二十年没有给那个人写过任何东西。至于中间的这段时间里我的观点和行为，我所能提出的最好见证人正是那些攻击和指控我的人。\par

7. 但分离的真正原因不是这封信。另有疏远的缘由。我羞于提及；若非近期的事件迫使我为了多数人的益处而公开他们的所有心思，我本会永远保持沉默。善良的人们认为与我共融是恢复他们权威的障碍。有些人受到我向他们提出的某一信经的签署的影响——我承认，并不是我不信任他们的情感，而是因为我希望消除那大部分与我意见相同的弟兄们对他们的猜疑。因此，为了避免因那份宣认而产生任何障碍，致使他们不被当今的当权者接纳，他们便断绝了与我的共融。这封信是后来想出来的，作为分离的借口。我说这话的一个非常明显的证据是：在谴责了我、并按自己的喜好编造了对我如此这般的控诉之后，他们就在与我沟通之前向四面八方发送了他们的信件。他们的信在传递过程中已被其他人持有，这些人正准备在信到达我手里七天之前就将其寄出。其想法是，信会从一个人传到另一个人，从而迅速传遍整个地区。当时那些清楚告知我他们所有行动的人向我报告了这一切。但我决定保持沉默，直到一切秘密的启示者通过明白无误的证明显明他们的行为。\par

\SourceRecord{Translated by Blomfield Jackson.；Nicene and Post-Nicene Fathers, Second Series；Vol. 8.；Edited by Philip Schaff and Henry Wace.；Buffalo, NY: Christian Literature Publishing Co.,；1895.；<http://www.newadvent.org/fathers/3202223.htm>.}

% structure: p0001 source=h1 target=section reason=page-title

\section{第二十四封信}

\Emph{\Person{圣巴西略·凯撒勒雅}}\par

致司铎格内提乌斯。\par

1. 我已收到阁下的来信，并为你恰当地称呼他们所撰写的文稿为\Quote{离婚书}而感到欣喜。\BibleRef{玛 19:7} 我完全无法想象，这些作者在基督的审判台前——那里任何借口都无济于事——能作出怎样的辩护。他们控告我，猛烈攻击我，并按照他们的愿望而非真相编造故事，却装出极大的谦卑，并指责我因拒绝接见他们的使者而傲慢。他们写的几乎全是谎言——我不想夸大其词——企图说服人而非天主，讨人喜欢而非天主，而天主眼中没有比真理更宝贵的了。此外，在写给我的信中，他们引入了异端言论，并隐藏了这种不敬之心的作者，以便让大多数较单纯的人被我遭受的诽谤所欺骗，并假定这些插文是我的。因为这些巧言诽谤者丝毫未提及这些邪恶学说的作者姓名，所以单纯的人容易怀疑这些编造——即便不是文字表述——是出于我。既然你们已知道这一切，我劝你们自己不要烦乱，并安抚那些激动者的情绪。我说这话，也知道我的辩护不易被接受，因为在我之前已有权势者对我发出的恶毒诽谤。\par

2. 至于那些传阅的文稿并非我所作，他们对我的愤怒情绪如此混淆理智，以致无法看到有益之处。然而，如果你们亲自向他们提出这个问题，我认为他们不会顽固到如此地步，竟敢亲口撒谎，指控那文件是我的。如果那不是我的，为何我要为别人的作品受审判？但他们可能会坚持说，我与阿波利纳尔共融，并从心底怀有这种乖谬的学说。那就请他们出示证据。如果他们能洞察人心，那就让他们说出来；而你们就接受他们对一切所说的话。另一方面，如果他们试图以明显公开的理由证明我与阿波利纳尔的共融，那就请他们拿出我写给他的或他写给我的共融信函。让他们指出我曾与他的神职人员进行过交往，或曾接待过他们中的任何一位进入祈祷共融。如果他们引用那封二十五年前写的信——那是平信徒写给平信徒的，甚至不是照我写的原样，而是被篡改过的（天主知道是谁所为）——那么请认清他们的不公。没有主教因为在身为平信徒时对无关紧要之事稍欠谨慎地写了些东西而受指控；那并非关乎信仰，而仅是友好的问候。甚至我的对手们也可能知道曾写信给犹太人和外邦人，却未受任何责备。迄今为止，从未有人因像我如今被这些吹毛求疵者所定罪的这类行为而受审判。深知人心的天主知道，我从未写过这些东西，也从未赞同过它们，而是谴责所有持守那关于位格混淆的邪恶意见的人——这混淆正是最不敬的撒伯流异端所复兴的。所有亲自认识我卑微之人的弟兄们也同样清楚。就让那些如今激烈控告我的人反躬自省，他们会承认，自幼年起我就远离任何此类学说。\par

3. 如果有人询问我的意见，他可以从那份附有他们亲笔签名的小文件中了解到。他们希望销毁这份文件，并急于掩盖自己通过诽谤我而改变立场的事实。因为他们不愿承认对自己签署我交给他们的手书感到后悔；相反，他们指控我不敬，以为无人察觉他们与我的分裂只是一个借口，而实际上他们已背离了那曾在众多见证人面前多次书面宣认、并最终从我手中领受并签署的信德。任何人都可以阅读那些签名并从文件本身了解真相。他们的意图将显而易见，如果任何人读了我交给他们的签名字样后，又读他们交给格拉修的信经，并观察到两个宣认之间巨大的差异。对于如此轻易改变立场的人，最好别审视别人眼中的木屑，而应拔除自己眼中的大梁。我在另一封信中对每一点作了更全面的辩护；这将使需要更多保证的读者满意。你们如今收到这封信后，请放下一切沮丧，并确认对我的爱——这爱使我热切渴望与你们合一。实在说，如果针对我的诽谤如此盛行，以致冷却你们的爱心，使我们彼此疏远，那对我将是巨大的忧伤和心中无法平息的痛苦。祝安。\par

\SourceRecord{Translated by Blomfield Jackson.；Nicene and Post-Nicene Fathers, Second Series；Vol. 8.；Edited by Philip Schaff and Henry Wace.；Buffalo, NY: Christian Literature Publishing Co.,；1895.；<http://www.newadvent.org/fathers/3202224.htm>.}

% structure: p0001 source=h1 target=section reason=page-title

\section{书信 225}

\Emph{圣巴西略（凯撒勒雅的）}\par

\Emph{致\Person{德摩斯梯尼}}，\Emph{由主教会议发出}\par

我常对天主及我们在其统治下的蒙天主钟爱的皇帝心怀感激，尤其在您来到我们中间时，我更有理由向天主和蒙天主钟爱的皇帝表达这份感激。我意识到有些和平的敌人正准备煽动您威严的法庭来反对我，我一直等待阁下召唤我，以便您从我这里了解真相——如果您的高见认为对教会事务的调查属于您的职权范围的话。法庭忽略了我，但阁下受\Person{斐洛卡勒斯}指控的影响，下令我的兄弟暨同僚\Person{额我略}拉至您的审判席前。他服从了您的召唤；他岂能不服从？但他突然遭受胁痛，同时因受寒引发旧患肾病。因此，尽管他被您的士兵强行拘押，却不得不被送到一个安静的地方，以便治疗他的疾病，并在难以忍受的痛苦中得到些许安慰。在此情况下，我们联合向阁下恳请，切莫因审判延期而恼怒。公共事务未曾因我们的延误而受损，教会亦未受伤害。若涉及浪费金钱的问题，教会资金的司库就在那里，随时准备向任何人交账，并揭露那些敢于在您面前进行详细听证的人所提出的指控之不公。因为从已故主教的真实著作中，他们不难向任何寻求真相的人阐明事实。如果有其他教会法规需要调查，且阁下屈尊俯听并审理，我们全体必须到场，因为如果法规有任何缺失，责任在于祝圣者，而非被强行逼迫履行职务者。因此我们恳请阁下将案件保留在我国境内审理，不要强迫我们出境，也不要迫使我们去与那些在教会问题上尚未达成一致的会晤。我也请求您怜悯我的年老和疾病。若蒙天主恩许，您将在实际调查中发现，在主教的任命中，无论大小法规均未被忽略。我祈求在您的治理下，能与我的弟兄们实现合一与平安；但只要这尚未实现，我们甚至连会面都困难，因为我们之间许多较为单纯的弟兄因彼此争论而受苦。\par

\SourceRecord{Translated by Blomfield Jackson.；Nicene and Post-Nicene Fathers, Second Series；Vol. 8.；Edited by Philip Schaff and Henry Wace.；Buffalo, NY: Christian Literature Publishing Co.,；1895.；<http://www.newadvent.org/fathers/3202225.htm>.}

% structure: p0001 source=h1 target=section reason=page-title

\section{第226封信}

\Emph{\Person{凯撒勒雅的圣巴西略}}\par

致他属下的隐修士。\par

神圣的天主或许会赐予我与你们相见的喜乐，因为我始终渴望见到你们，并听到你们的消息；因为除了你们在基督的诫命上长进与成全，我没有别的为我的灵魂寻得安息。但在这希望尚未实现之时，我觉得必须借着我们亲爱的、敬畏天主的弟兄们的媒介，以书信问候你们，我亲爱的朋友们。因此，我派遣了可敬而亲爱的弟兄、福音的同工、司铎默肋提乌斯前往你们那里。他将告诉你们我对你们的殷切疼爱，以及我灵魂的焦虑：我日日夜夜为你们恳求主，盼能在我们的主耶稣基督的日子，因着你们的得救而拥有胆量；当你们的工作因天主的公义审判而受考验时，你们能在圣人的光明中闪耀。与此同时，当前的困境使我深为焦虑：所有教会都动荡不安，所有人灵魂都受到筛簸。有些人甚至毫无保留地张口攻击他们的同仆。谎言被大胆说出，真理被隐藏。被告未经审判就被定罪，原告无凭无据就被相信。我听说许多人携带书信攻击我，用毒刺、羞辱和打击来对付我，为一些事——我在真理的审判台前已有所辩护；我本想保持沉默，事实上也是如此；因为三年来，我承受着诽谤的打击和控告的鞭笞，自认为拥有那知晓一切秘密的主作见证，证明这些是虚假的。但我现在看到，许多人把沉默当作这些诽谤的佐证，并以为我的沉默不是出于我的忍耐，而是因为我无法开口反对真理。为此，我尝试写信给你们，恳求你们在基督内的爱，不要接受这些片面的诽谤为真实，因为正如经上所记载：\BibleQuote{法律不审判任何人，除非它听过并知道他的行为。}\par

然而，在公正的法官面前，事实本身足以证明真理。因此，即使我沉默，你们也可以看看发生了什么事。那些如今控告我异端见解的人，曾被公开看见与异端派别同列。那些以他人的著作定我罪的控告者，明显违反了他们曾以书面形式给我的宣认。看看这些胆大妄为之徒的行径。他们惯常的做法是投靠当权派，践踏较弱的友人，巴结强者。那些写过著名信件反对欧多克修斯及其所有派别的人，曾将信件寄给所有弟兄，声明他们避之如同避祸，并且不接受罢免他们的投票，因为那些投票是异端者作出的——正如他们当时说服我的——这些人完全忘记了一切，加入了他们的派别。他们没有否认的余地。当他们在安卡拉私下与他们共融时——那时他们尚未被公开接纳——他们已暴露了心思。那么请问他们：如果曾经给厄克迪修斯共融的巴西里得斯现在是正统的，为什么他们从达达尼亚回来时，在冈格拉地区推翻了他的祭台，而摆设了自己的祭台？为什么他们最近攻击了阿马塞亚和泽拉的教会，在那里自行任命了司铎和执事？如果他们作为正统者与他们共融，为什么又像对待异端者一样攻击他们？如果他们认为他们是异端，为什么他们不回避与他们的共融？我的尊贵弟兄们，难道连孩童的智力也能明白：他们总是为了某种个人利益而竭力诽谤或支持？所以他们疏远我，不是因为我未回信（据称这是主要冒犯），也不是因为我未接待他们声称派来的\OriginalTerm{乡村主教}{chorepiscopi}。那些捏造故事的人将向主交账。有一个人——欧斯塔提乌斯——曾被派遣，向代理官的法庭递交了一封信，并在城里住了三天。当他准备回家时，据说他在深夜来到我家，当时我正在睡觉。听说我睡着了，他便离开了；次日他没来找我，就这样走完了对我尽义务的形式后离开了。这就是我犯下的罪过。这就是这些长久忍耐的人忽视了我曾在爱心中服侍他们的先前功劳，所要定我的罪。因这一失误，他们对我大发雷霆，以至于使我在全世界所有教会中——至少在他们所能到达的地方——受到谴责。\par

3. 但这当然不是我们分离的真正原因。当他们发现若与我疏远，便能博得欧佐伊乌斯的欢心时，便捏造了这些借口。其目的是为了在当局面前找到某种推荐理由，以便攻击我。如今，他们甚至开始诋毁《尼西亚信经》，并给我起了个绰号叫\Quote{\OriginalTerm{同体论者}{Homoousiast}}，因为在那信经中，独生子被称为与天主圣父\OriginalTerm{同体}{homoousios}。这不是说一个本体被分成了两个相似的部分；决非如此！那神圣而受天主钟爱的公会议并非此意；他们的意思是：圣父在本质上是什么，圣子也就是什么。而他们自己也用\Quote{\OriginalTerm{光明出自光明}{Light of Light}}这句话向我们解释了这一点。如今，正是他们自己从西方带来的《尼西亚信经》，被他们呈递给提雅纳公会议，并凭此信经得以被接纳。但他们对于此类变更有一套精巧的理论：他们使用信经的言辞，如同医生根据特定时刻使用药方，为适应不同的疾病而替换一种又一种说法。这种诡辩的缺陷理应由你们的思考来辨别，而非由我来证明。因为\Quote{主必赐给你们理解力}\BibleRef{弟后 2:7}，使你们知道什么是正确的教义，和什么是弯曲和乖谬的。如果我们今天签署一份信经，明天又签署另一份，随季节而变换，那么那说\Quote{同一主、同一信德、同一圣洗}\BibleRef{弗 4:5}的宣告便是虚假的。但若它是真实的，那么\Quote{不要让任何人用这些虚浮的话语欺骗你们}。他们妄称我在关于圣神的事上引进了新事。请问这新事是什么？我承认我所领受的：护慰者与圣父和圣子并列，不被列于受造之物中。我们宣认了信德于圣父、圣子及圣神，并因圣父、圣子及圣神之名受洗。因此，我们从不将圣神与圣父及圣子的结合分离。因为我们的心智，被圣神光照，便注视圣子，并在圣子内，如同在肖像中，瞻仰圣父。我并不自己杜撰名号，而是称圣神为护慰者；我也不允许损毁祂当得的荣耀。这些确实是我的教义。若有人愿意为此控告我，就让他控告罢；让我的迫害者迫害我罢。凡相信对我的毁谤之人，就当预备审判。\Quote{主临近了。}\Quote{我一无所挂虑。}\par

4. 若有人在叙利亚写作，这与我无关。因为经上说：\Quote{凭你的话，你将成义；凭你的话，你将受定罪。}\BibleRef{玛 12:37} 让我的话审判我自己吧。不要因别人的错误而定我的罪，也不要以二十年前写的信来证明我会与写那些东西的人共融。在那些事被写之前，在针对他们的任何这类嫌疑被激起之前，我确实曾以平信徒的身份写给平信徒。关于信仰，我从未写过任何像他们现在拿来诽谤我的话之类的东西。我所寄出的不过是一封单纯的问候，以回复友好的通信，因为我躲避并诅咒所有染有撒伯里乌斯谬论的人，以及所有坚持亚略意见的人，视他们为不虔敬者。若有人说圣父、圣子、圣神是同一的，并认为在同一位格下多个名称，以及由三个位格描述的一个\OriginalTerm{本体}{hypostasis}，我就视这样的人属于犹太人派别。同样，若有人说圣子在本质上不像圣父，或把圣神贬低为受造物，我便诅咒他，并说他正接近异教者的错误。然而，我的信无法约束控告我之人的口；反倒可能因我的辩护而激怒他们，并策划出新的、更猛烈的攻击。但你们的耳朵却容易得到保护。所以，尽你们所能，按我所吩咐的去做。保持你们的心灵洁净，不受他们的毁谤之偏见；并坚持要我针对指控作出答辩。如果你们发现真理在我这边，就不要向谎言屈服；反之，若你们感到我为自己辩护无力，那么就可以相信我的控告者，认为他们是可信的。他们彻夜不眠，意图加害于我。我并不向你们要求这一点。他们正走向经商生涯，把对我的诽谤变成谋利的手段。相反，我恳求你们安居家中，过着端正的生活，安静地做基督的工作。我劝你们避免与他们交往，因为那常会导致听者偏差。我说这些，是为使你们保持对清白之人的爱，完整地保存祖传的信德，并在主面前被证明为真理之友。\par

\SourceRecord{Translated by Blomfield Jackson.；Nicene and Post-Nicene Fathers, Second Series；Vol. 8.；Edited by Philip Schaff and Henry Wace.；Buffalo, NY: Christian Literature Publishing Co.,；1895.；<http://www.newadvent.org/fathers/3202226.htm>.}

% structure: p0001 source=h1 target=section reason=page-title

\section{第二十七书}

\Emph{\Person{凯撒勒雅的圣巴西略}}\par

\Emph{安慰信，致\Place{科洛尼亚}的圣职员}\par

在天主和人面前，有什么像完美的爱那样美好而可敬的呢？完美的爱，正如智者所教导的，就是法律的满全。见：\BibleRef{罗 13:10} 因此，我赞许你们对主教的热忱之爱，因为正如一个充满感情的儿子无法忍受好父亲的去世一样，基督的教会也无法忍受牧者和教师的离开。因此，你们对主教的极大爱慕，证明了你们善良高尚的品性。然而，你们对灵性之父的这种善意，只要是在理性和适度中表现出来，就值得赞许；一旦越界，就不再值得同样的称赞。至于你们极其敬爱天主的兄弟，我们的同工\Person{欧弗罗尼乌斯}，那些获得教会管理职责的人已经展现了良好的治理；他们按照时势所迫而行，既造福了他被调任去的教会，也造福了你们——他离开的你们。不要视此事为仅仅是人的规定，也不是出于那些关注世俗事物之人的算计。要相信那些对教会怀有焦虑关怀的人，正如他们所行的，是在圣神的帮助下行动的；要将这事的开端铭刻在你们心中，并尽力使之完善。要安静而感恩地接受所发生的事，坚信凡拒绝接受教会在天主的教会中所规定之事的人，就是在反抗天主的命令。不要与你们在\Place{尼科波利斯}的母教会发生争执。不要对那些承担了你们灵魂焦虑责任的人发怒。当\Place{尼科波利斯}的局面稳固时，你们在其中也能得以保存；但如果一些动荡影响那里，即使你们有无数保护者，头若被毁，心也会被毁灭。如同住在河边的人；当他们看到远处上游有人建造坚固的水坝阻挡水流时，他们知道，通过阻挡急流，这人在为他们提供安全。同样，那些如今承担了教会关怀重担的人，通过保护其他人，也在为你们自己的安全提供保障。你们将免受一切风暴侵袭，而其他人则必须承受攻击的冲击。但你们也应该考虑到这一点：他并没有抛弃你们；他只是将其他人也纳入他的照管之下。我并非如此嫉妒，以至于强迫那位能将自己的恩赐分施给他人的人，只将他的恩惠局限于你们，仅限于你们自己的城市。给泉源围上篱笆、破坏水流的人，难免染上嫉妒之病；同样，试图阻止丰富教导进一步涌流的人也是如此。让他也对\Place{尼科波利斯}有所关怀，并让他对那里的忧虑中加上你们的利益。他增加了劳苦，但为你们的勤勉并未减少。有一件事你们所说的话令我确实感到痛心，在我看来十分过分，那就是：\Quote{如果你们不能达到目的，你们将诉诸法庭，将此事交给那些以推翻教会为主要祈祷对象的人手中。}要当心，免得有人被不明智的激情冲昏头脑，说服你们在法庭上提出任何诉状，从而给你们带来伤害，以致灾难发生，后果的重担落在那些引发了这一切的人头上。接受我的劝告。我以慈父般的精神向你们提出。同意与那些极其敬爱天主的主教们已达成的安排，这安排是符合天主旨意的。等待我的到来。当我与你们同在时，在天主的助佑下，我将亲自向你们传达所有我在这封信中无法表达的劝勉，并尽我所能给予你们一切可能的安慰，不是用言语，而是用行动。\par

\SourceRecord{Translated by Blomfield Jackson.；Nicene and Post-Nicene Fathers, Second Series；Vol. 8.；Edited by Philip Schaff and Henry Wace.；Buffalo, NY: Christian Literature Publishing Co.,；1895.；<http://www.newadvent.org/fathers/3202227.htm>.}

% structure: p0001 source=h1 target=section reason=page-title

\section{第二十八封信}

\Emph{\Place{凯撒勒雅}的\Person{圣巴西略}}\par

致\Place{科洛尼亚}的行政官们。\par

我已收到诸位大人的信函，并感谢至圣的天主，因你们虽身负国政，仍不忘将教会事务置于次要位置。我感激你们每人所表现出的焦虑，仿佛在为自己私利、甚至为保卫自己的生命而行，且在我因你们那深受天主所爱的主教\Person{厄乌弗若尼乌斯}被调离而苦恼时，给我写信。\Place{尼科波利斯}并未真的将他从你们那里夺走；若她能在法官面前陈词，或许会说她正在收回属于自己的东西。若受到礼遇，她将像一位慈母般告诉你们，她愿与你们共享这位父亲，他将分赐恩宠给你们每人：他既不让你们因敌人的入侵而受到丝毫损害，同时也不使你们失去所习惯的关怀。因此，请你们思量当前的紧急情况；运用你们的睿智去理解善治何以要求某种行动；然后宽恕那些为建立我们主耶稣基督的教会而采取此行动的主教们。请你们为自己设想什么是合宜的。你们的智慧无需教导。你们知道如何采纳爱你们之人的建议。你们未必知晓许多正在激荡的问题，这也是自然的，因为我们远在\Place{亚美尼亚}；然而我们身处事务之中，每日从四面八方听到教会倾覆的消息，充满深切的忧虑，恐怕那共同的敌人，嫉妒我们生活的长久和平，也能在你们的田地里播下莠子，使\Place{亚美尼亚}如同其他地方一样，被交给我们的敌人吞噬。目前请保持安静，不要拒绝让邻居与你们分享一艘好船的用途。不久之后，若主允许我来到你们中间，如果你们认为必要，你们将从所发生的事中得到更大的安慰。\par

\SourceRecord{Translated by Blomfield Jackson.；Nicene and Post-Nicene Fathers, Second Series；Vol. 8.；Edited by Philip Schaff and Henry Wace.；Buffalo, NY: Christian Literature Publishing Co.,；1895.；<http://www.newadvent.org/fathers/3202228.htm>.}

% structure: p0001 source=h1 target=section reason=page-title

\section{第二二九封信}

\Emph{\Person{圣巴西略}\Place{凯撒勒雅}}\par

致\Place{尼科波利}的圣职人士。\par

我确信，由一位或两位虔诚之人所成就的工作，并非没有圣神的合作。因为当惟有属天之事摆在我们面前，当圣人们被激动行事，丝毫不顾及个人私利，而只以悦乐天主为唯一目标时，显然是主在引导他们的心。当属神之人在会议中领导，主的子民以同心合意之心跟随他们时，毫无疑问，他们的决定是在我们的主耶稣基督的参与下达成的，祂曾为教会倾流了自己的血。因此，你们正确地认为，我们那极其蒙天主所爱的弟兄和同工\Person{波哀默尼乌斯}，他适时抵达你们中间，并发现了这种安慰你们的方法，乃是出于天主的感动。我不仅赞扬他发现了正确的行动方针；我更钦佩他的坚定，他不允许任何拖延介入，以免减弱请求者的热忱，或给反对派以采取预防措施的机会，并启动暗中敌人的诡计，他立刻以成功的结局为他幸福的事业加冕。愿主以其特殊恩宠保全他和他的家人，好使教会合乎体统地延续在一个绝不堕落的传承中，不给那恶者留地步，这恶者如今——若有这样的时候——正因教会的稳固建立而烦恼。\par

2. 我也已详细写信劝勉我们在\Place{科洛尼亚}的弟兄们。此外，你们宁可以忍耐对待他们的心态，也不可增加他们的恼怒，仿佛你们因他们的卑微而轻视他们，或因你们的轻蔑而激起争执。争论者行事不经充分商议，并以惹恼对手为目的而糟糕地处理自己的事务，这是很自然的。没有一个人渺小到不能给那些想要借口的人提供重大麻烦的机会。我不是随意说的。我是根据自己患难的经验说的。愿天主因你们的祈祷而保护你们免受这些患难。也请为我祈祷，使我旅途顺利，抵达后能与你们一同因你们现任的牧者而喜乐，并与你们一同为我们共同之父的离世得到安慰。\par

\SourceRecord{Translated by Blomfield Jackson.；Nicene and Post-Nicene Fathers, Second Series；Vol. 8.；Edited by Philip Schaff and Henry Wace.；Buffalo, NY: Christian Literature Publishing Co.,；1895.；<http://www.newadvent.org/fathers/3202229.htm>.}

% structure: p0001 source=h1 target=section reason=page-title

\section{书信第二三〇号}

\Emph{\Person{凯撒勒雅的圣巴西略}}\par

致\Place{尼科波利斯}的长官们。\par

教会的治理由受托担任教会首领职务的人执行，但他们的力量因平信徒而加强。那些属于天主所爱的主教们已采取的措施，已经完成。其余的事关乎你们，如果你们愿意欣然接纳已赐给你们的主教，并有力抵制外来的攻击。因为没有什么比一致爱戴被任命的主教，并坚定持守你们的位置，更可能使所有（无论是统治者，还是嫉妒你们和平境遇的其余人）灰心丧气。若他们看到自己的计谋既不被圣职人员也不被平信徒接受，他们很可能对任何邪恶的企图感到绝望。那么，请促成你们自己的正确情感为全城所共享，并以此向市民和全区所有居民讲话，以巩固他们的良好情感，好使你们对天主的挚爱真诚无伪，遍传各处。我相信有一天必蒙允许我探访并巡视这教会，她是真宗教的乳母，我尊之为正统信仰的母城，因她自古以来一直受极其可敬、天主所选的人治理，他们持守\Quote{正如我们所受教的忠信之言}。你们已认可新近被任命者堪当先贤的继承者，我也赞同。愿你们在天主的恩宠中蒙保守。愿祂驱散我们仇敌的恶谋，并在你们灵魂中坚定力量与恒心，以持守已正确决定的事。\par

\SourceRecord{Translated by Blomfield Jackson.；Nicene and Post-Nicene Fathers, Second Series；Vol. 8.；Edited by Philip Schaff and Henry Wace.；Buffalo, NY: Christian Literature Publishing Co.,；1895.；<http://www.newadvent.org/fathers/3202230.htm>.}

% structure: p0001 source=h1 target=section reason=page-title

\section{书信二三一}

\Person{圣巴西略}\par

致\Place{依科尼雍}主教\Person{安菲洛基乌斯}。\par

我很少有机会写信给尊座，这使我颇感困扰。这就像我本可以时常与您相见并享受您的陪伴，却很少这样做一样。然而，我无法写信给您，因为从此地前往您那里的人太少；否则，我的信本该成为我生活的日记，将我每日所遇之事都告诉您，我亲爱的朋友。向您诉说我的事是我的安慰，我也知道您最关心的莫过于与我有关的事。如今，埃尔丕狄乌斯要回到他自己的主人那里，去驳斥某些敌人对他不实捏造的诽谤，他向我索要一封信。因此，我藉他问候尊座，并向您推荐这位值得您保护的人，既为正义之故，也为我本人之故。虽然我对他别无赞词，但因他十分重视为我送信一事，请将他列为我们朋友之中，并为我祈祷，为教会祈祷。\par

您要知道，我那位深受天主所爱的兄弟正在流亡中，因为他无法忍受无耻之人对他的骚扰。多阿拉正处于动荡之中，因为那肥胖的海怪把一切都搅乱了。据知情者告知，我的敌人在宫廷中阴谋反对我。但迄今为止，主的手一直护佑着我。只求您祈祷，使我最终不致被抛弃。我的兄弟安之若素。多阿拉收留了那个老骡夫。她已无能为力。主必会挫败我敌人的计谋。我目前和将来所有麻烦的唯一良药就是见到您。如果可能，趁我还在世，务请来看我。关于圣神的书我已写成，如您所知，已经完成。我这里的弟兄们阻止我将它写在纸上送给您，并告诉我他们已接到您的命令，要用羊皮纸誊写。因此，为了不违背您的吩咐，我暂缓发送，但稍后我会送出，只要找到合适的人携带。愿您藉圣者的仁慈，得以赐给我和天主的教会，健康幸福，并为我向主祈祷。\par

\SourceRecord{Translated by Blomfield Jackson.；Nicene and Post-Nicene Fathers, Second Series；Vol. 8.；Edited by Philip Schaff and Henry Wace.；Buffalo, NY: Christian Literature Publishing Co.,；1895.；<http://www.newadvent.org/fathers/3202231.htm>.}

% structure: p0001 source=h1 target=section reason=page-title

\section{书信第232号}

\Person{凯撒勒雅的圣巴西略}\par

致\Place{依科尼雍}的主教\Person{安菲洛基乌斯}\par

每一天，只要带来您的一封信，对我而言就是庆节，最大的庆节。当庆节的象征被带来时，我还能称它什么呢？只能称它为庆节中的庆节，正如旧约所常说的安息日中的安息日。我感谢主，您身体康健，并在和平的教会中庆祝了救恩工程的纪念日。我曾因一些困扰而心烦意乱，也因我那位蒙主所爱的兄弟被流放而深感痛苦。请为他祈祷，愿天主有朝一日恩赐他看见教会从异端者的咬伤创伤中得到痊愈。趁我尚在此世时，请来看望我。请按您自己的意愿和我最恳切的祈祷行事。请允许我对您所赐福的礼物感到惊奇，因为您神秘地祝愿我享有矍铄的晚年。您以灯火激励我夜间劳作；您以甜美的食物似乎稳妥地担保我全身健康。但以我这把年纪，已无法大嚼，因为我的牙齿早已被岁月和病痛磨损。至于您所询问的事，我在寄给您的文件中有一些答复，尽我所能力求妥善，并趁便写好。\par

\SourceRecord{Translated by Blomfield Jackson.；Nicene and Post-Nicene Fathers, Second Series；Vol. 8.；Edited by Philip Schaff and Henry Wace.；Buffalo, NY: Christian Literature Publishing Co.,；1895.；<http://www.newadvent.org/fathers/3202232.htm>.}

% structure: p0001 source=h1 target=section reason=page-title

\section{第233封书信}

\Emph{\Person{圣巴西略·凯撒勒雅}}\par

致\Person{安菲洛希乌}，答复某些问题。\par

一、我知道我自己曾听说过这些，也明白人类的构造。我该说什么呢？理智是一件奇妙之物，在其中我们拥有那按造物主肖像的要素。理智的运作也是奇妙的；由于它在不断运动中，常就非存在之物如同存在之物般形成想象，又常直趋真理。但在理智中有两种能力：按照我们信仰者的观点，一种是魔鬼的恶力，牵引我们趋向它们自己的背教；另一种是神圣而良善的，带领我们趋向天主的相似。因此，当理智独自而无所援助时，它只默观与自身相称的小事。当它屈服于欺骗者时，它便废弃了正确的判断，沉溺于怪诞的幻想。于是它认为木头不再是木头，而是一个神；它视黄金不再是金钱，而是崇拜的对象。反之，若它赞同其神圣部分，并接受圣神的恩赐，那么，就其本性所允许，它便成为觉察神圣的。可以说，有三种生活状态和三种理智运作。我们的道路可能是邪恶的，我们理智的活动也是邪恶的；诸如奸淫、偷盗、崇拜偶像、毁谤、争斗、情欲、叛乱、虚荣，以及宗徒\Person{保禄}所列的一切本性的作为。或者灵魂的运作可说是处于中间，既无可指责也无可称赞，如同我们通常视为无所谓的机械工艺的感知，其本身既不偏向德性也不偏向恶习。因为舵手或医生的技艺有什么恶呢？这些运作本身也不是德性，而是根据使用者的意愿偏向一方或另一方。但那受圣神之天主性所孕育的理智，立刻能观照伟大之物；它注视神圣的美，尽管只限于恩宠所赐及其本性所能接受的程度。\par

二、因此，让他们放弃这些辩证的问题，以敬畏而非有害的精确来审查真理。我们理智的判断是为理解真理而赐予我们的。如今我们的天主就是那真理本身。因此，我们理智的首要功能是认识独一的天主，但认识祂只限于无限伟大者能被极其微小者所认识的程度。当我们的眼睛初次被引向可见之物的知觉时，并非所有可见之物都立即映入眼帘。天穹并非一眼被看见，我们被某一种表象所包围，然而实际上其中许多事物，甚至可以说所有事物，都是未被察觉的——众星的本性、它们的宏大、距离、运动、会合、间隔、其他状况，穹苍的实际本质，从凹圆周到凸表面的深度距离。尽管如此，没有人会因未知而声称天穹是不可见的；因我们对其有限的知觉，天穹应被称为可见的。天主的情况也是如此。如果理智曾受魔鬼伤害，它将犯崇拜偶像之罪，或堕入其他形式的不敬。但如果它顺服于圣神的援助，它将理解真理，并认识天主。但它认识祂，如宗徒所说，是局部的；在来世则更圆满。因为\Quote{当那圆满的来到，那局部的就要消逝。}\BibleRef{格前 13:10} 因此，理智的判断是好的，是为着好的目的——认识天主——而赐予我们的；但它只在其所能的范围内运作。\par

\SourceRecord{Translated by Blomfield Jackson.；Nicene and Post-Nicene Fathers, Second Series；Vol. 8.；Edited by Philip Schaff and Henry Wace.；Buffalo, NY: Christian Literature Publishing Co.,；1895.；<http://www.newadvent.org/fathers/3202233.htm>.}

% structure: p0001 source=h1 target=section reason=page-title

\section{\WorkTitle{第234封书信}}

\Person{圣巴西略·凯撒勒雅}\par

致同一人，答复另一问题。\par

你敬拜你所知道的，还是你所不知道的？如果我回答\Quote{我敬拜我所知道的}，他们立即反问：\Quote{敬拜对象的\OriginalTerm{本质}{essence}是什么？}如果我承认我不认识本质，他们又反过来质问我：\Quote{那么你所敬拜的，你并不认识。}我回答说，\Quote{认识}一词有多种含义。我们说我们知道天主的伟大、祂的能力、祂的智慧、祂的良善、祂对我们的眷顾，以及祂审判的公正；但并非祂的终极本质。因此，这个问题只是为了争辩而提出的。因为否认自己认识本质的人，并不承认自己不认识天主，因为我们关于天主的观念是从我所列举的一切属性中汇集而来的。但祂说，天主是单纯的，你所认为可认识的任何属性都是祂的本质。然而，这种诡辩所包含的荒谬之处不可胜数。当所有这些崇高的属性被列举出来时，它们都是一个本质的名称吗？祂的威严与祂的慈爱、祂的公义与祂的创造权能、祂的眷顾与祂的预知、祂的赏罚报应、祂的威严与祂的眷顾，这些都有同等的效力吗？当我们提及其中任何一个时，我们是在宣示祂的本质吗？如果他们说是，那么不要问我们是否认识天主的本质，而要问我们是否知道天主是威严的，或公义的，或仁慈的。我们承认我们知道这些。如果他们说是本质是某种不同的东西，那么不要在单纯性上责备我们。因为他们自己也承认本质与所列举的每一个属性之间是有区别的。\OriginalTerm{行动}{operations}是多样的，本质是单纯的，但我们说，我们从天主的行动认识我们的天主，但不敢接近祂的本质。祂的行动降临到我们，但祂的本质仍遥不可及。\par

2. 但有人反驳说：\Quote{如果你不认识本质，你就不认识祂本身。}你可以反问：\Quote{如果你说你知道祂的本质，你就不认识祂本身。}一个人被疯狗咬了，看见盘子里有一只狗，他并不比健康的人看到更多；他值得可怜，因为他以为他看到了他实际没有看到的东西。因此，不要因他的宣告而钦佩他，而要因他的疯狂而怜悯他。要知道，当他们说\Quote{如果你不认识天主的本质，你就敬拜你所不知道的}时，那是嘲笑者的声音。我知道祂存在；至于祂的本质是什么，我视之为超越理智。那么我如何得救呢？藉着信德。信德足以知道天主存在，而不必知道祂是什么；并且\BibleQuote{祂是寻求祂者的报答。}\BibleRef{希 11:6} 因此，认识神圣本质包括对其\OriginalTerm{不可理解性}{incomprehensibility}的觉察，我们敬拜的对象不是我们领悟其本质的那一位，而是我们领悟其本质存在的那一位。\par

3. 还可以向他们提出以下反问。\BibleQuote{从来没有人见过天主，只有在圣父怀里的\OriginalTerm{独生子}{Only-begotten}，祂彰显了祂。}\BibleRef{若 1:18} 独生子彰显了圣父的什么？是祂的本质还是祂的权能？如果是祂的权能，我们知道祂所彰显给我们的那么多。如果是祂的本质，请告诉我，祂在什么地方说祂的本质是\Quote{\OriginalTerm{非受生}{unbegotten}}？亚巴郎何时敬拜？岂不是在他相信的时候？他何时相信？岂不是在他蒙召的时候？此处圣经哪里有什么证据表明亚巴郎领悟了？门徒们何时敬拜祂？岂不是在他们看见受造物服从祂的时候？是从海洋和风服从祂，他们才认出祂的\OriginalTerm{天主性}{Godhead}。因此，认识来自行动，敬拜来自认识。\BibleQuote{你信我能做这事吗？}\BibleQuote{主，我信；}于是他就敬拜了祂。所以，敬拜随信心而来，信心藉权能得以坚固。但如果你说，相信的人也知道，他因他所信的而知道；反之亦然，他因他所知的而相信。我们从天主的权能认识天主。因此，我们信那已被认识的天主，并敬拜那已被信的天主。\par

\SourceRecord{Translated by Blomfield Jackson.；Nicene and Post-Nicene Fathers, Second Series；Vol. 8.；Edited by Philip Schaff and Henry Wace.；Buffalo, NY: Christian Literature Publishing Co.,；1895.；<http://www.newadvent.org/fathers/3202234.htm>.}

% structure: p0001 source=h1 target=section reason=page-title

\section{\WorkTitle{第二百三十五封书信}}

\Person{凯撒勒雅的圣巴西略}\par

致同一人，答复另一问题。\par

1. 在秩序上，认识与信德哪一个在先？我的回答是：通常对于门徒而言，信德先于认识。但在我们的教导中，若有人认为认识先于信德，我也不反对；此处的\Quote{认识}仅限于人类理解所能达到的范围。在我们学习中，必须先相信字母\Emph{a}是对我们说的，然后学习符号及其发音，最后才获得对字母力量的确切概念。但在我们关于天主的信仰中，首先是\Quote{天主存在}这一概念。我们从祂的工程中领悟这一点。因为正如我们藉着所造之物，从世界的创造中领悟祂的智慧、祂的良善及祂一切不可见的事，同样我们也如此认识祂。同样，我们也接受祂为我们的主。因为既然天主是宇宙的创造者，而我们是宇宙的一部分，那么天主就是我们的创造者。这一认识之后便是信德，而信德之后便是敬拜。\par

2. 但\Quote{认识}一词有多种含义，因此那些嘲弄较单纯心思的人，喜欢用惊人言论使自己显得非凡（就像变戏法者当着众人的面让球消失），便匆忙将一切纳入他们笼统的询问中。我说，认识的应用范围极广，认识可以来自事物的本质、数目、体积、力量、存在方式、产生的时期、本质。因此，当我们的对手将整体纳入他们的问题中时，若是他们抓住我们承认自己认识，便立刻要求我们认识本质；反之，若他们见我们对作任何断言持谨慎态度，就将不虔敬的污名强加于我们。但我承认，我知道天主可知的事，也知道什么是超越我理解的。因此若你问我是否知道沙是什么，我回答知道，你若立刻问我沙的数目，你显然是在诽谤我；因为你的第一个询问只关于沙的形状，而你的第二个不当反对则涉及它的数目。这诡辩就好比有人说：你认识弟茂德吗？啊，你若认识弟茂德，你就认识他的本性。既然你承认认识弟茂德，就给我讲讲弟茂德的本性吧。不错，但我是同时既认识又不认识弟茂德，尽管并非以相同方式和相同程度。我不是以我知道的方式不知道，而是以一种方式知道，以一种方式不知道。我根据他的形态和其他属性认识他，但我不认识他的本质。的确，同样地，我也既认识又不认识我自己。我确实知道我是谁，但就我不认识我的本质而言，我又不认识我自己。\par

3. 让他们告诉我，保禄说：\BibleQuote{我们现在所知道的，只是局部的。}（见：\BibleRef{格前13:9}）我们是否局部地认识祂的本质，即知道祂本质的部分？不，这很荒谬，因为天主是没有部分的。那么，我们是否认识全部的本质？既然如此，\BibleQuote{及至圆满的来到，局部的就必归于无用。}（见：\BibleRef{格前13:10}）拜偶像者为何受责备？岂不是因为他们认识天主，却没有当作天主来荣耀祂吗？为何\BibleQuote{糊涂的迦拉达人}（见：\BibleRef{迦3:1}）受保禄责备，他说：\BibleQuote{你们既然认识了天主，或者更好说被天主所认识，怎么又再转向那懦弱无能、鄙陋的要素呢？}（见：\BibleRef{迦4:9}）天主在犹太地区是怎样被认识的？是否因为在犹太地区，人们知道祂的本质是什么？经上说：\BibleQuote{牛认识自己的主人。}（见：\BibleRef{依1:3}）\BibleQuote{驴认识主人的槽。}（见：\BibleRef{依1:3}）这样，驴认识槽的本质，但\Quote{以色列却不认识我。}所以，照你们看来，以色列受责备是因为不知道天主的本质是什么。\BibleQuote{愿你将你的愤怒倾注在不认识你的外邦人身上。}但，我再重复一次：认识是多样的——它包括感知我们的创造者、辨认祂奇妙的作为、遵守祂的诫命以及与祂密切的共融。这一切他们都推到一边，强迫认识落入单一含义，即默观天主的本质。\newline{}经上说：\BibleQuote{你要将这些物件放在约证前，我将在那里使你认识我。}\newline{}\Quote{我将被你所认识}这个措辞，难道不是\Quote{我要启示我的本质}的另一种说法吗？\newline{}\BibleQuote{主认识属于祂的人。}（见：\BibleRef{弟后2:19}）祂是否认识属于祂的人的本质，而不认识那些不顺从祂的人的本质？\newline{}\BibleQuote{亚当认识了自己的妻子。}（见：\BibleRef{创4:1}）他认识她的本质吗？\newline{}论到黎贝加说：\BibleQuote{那少女容貌很美丽，还是个处女，从未有人认识过她。}（见：\BibleRef{创24:16}）又：\BibleQuote{这事怎能成就？因为我不认识男人。}（见：\BibleRef{路1:34}）难道没有人认识黎贝加的本质？玛利亚的意思是\Quote{我不认识任何男人的本质}？\newline{}圣经岂不惯常以\Quote{认识}一词指称婚姻结合？\newline{}\Quote{天主将从赎罪盖上被认识}这句话的意思是，祂将被祂的敬拜者所认识。\newline{}而\BibleQuote{主认识属于祂的人}的意思是，因着他们的善行，祂接纳他们与祂密切共融。\par

\SourceRecord{Translated by Blomfield Jackson.；Nicene and Post-Nicene Fathers, Second Series；Vol. 8.；Edited by Philip Schaff and Henry Wace.；Buffalo, NY: Christian Literature Publishing Co.,；1895.；<http://www.newadvent.org/fathers/3202235.htm>.}

% structure: p0001 source=h1 target=section reason=page-title

\section{\WorkTitle{书信二三六}}

\Signature{凯撒勒雅的圣巴西略}\par

致同一位\Person{盎腓罗基}。\par

有关吾主耶稣基督不知末日的那一日和那一时辰的福音言论，早已屡次受到探究；\BibleRef{谷 13:32} 这是\OriginalTerm{无类似者派}{Anomoeans}为破坏独生子的荣耀而经常提出的反对，旨在证明祂的本质不同、尊严从属；因为若祂不全知，就不能与那位凭自己的预知和预测未来的能力而知识遍及宇宙者拥有相同的本质，也不能被视为与祂同形。这一问题如今被你明智地当作新问题提出给我。我可以给出答复，即我幼年时从我们教父们那里听到的、并因我喜爱善美而无疑接受的答案。我不期望它能挫败那些反对基督之人的无耻，因为有什么推理足以抵挡他们的攻击呢？然而，它足以说服所有爱主之人，他们心中由信德所赋予的确信比任何理性的论证更为坚强。\par

现在，\Quote{没有人}一词似乎是一个泛指，以至于没有一个人被排除在外，但这不是它在圣经中的用法，正如我在经文中观察到的那样：\Quote{除了一个，就是天主，没有良善的}。因为即使在这里，圣子也没有将自己排除在良善的本性之外。而是，因为圣父是最初的善，我们相信\Quote{没有人}一词是在隐含\Quote{最初}的情况下说出的。同样，经文\Quote{除了圣父外，没有人认识圣子}；\BibleRef{玛 11:27} 即使这里也没有指控圣神无知，而只是证明对其本性的认识首先自然属于圣父。同样，我们理解\Quote{没有人知道}，\BibleRef{玛 24:36} 是指圣父首先知道现在和将来之事，并且一般是为了向人展示第一因。否则，这段经文如何能符合圣经的其他证据，或与我们这些相信独生子是不可见天主的肖像的人们的共同观念相一致呢？这肖像不是身体形象，而是天主性本身和归于天主本质的大能属性的肖像，是能力的肖像，是智慧的肖像，因为基督被称为\Quote{天主的德能和天主的智慧}。\BibleRef{格前 1:24} 现在，智慧显然包含知识；如果祂在任何部分有欠缺，祂就不是整体的肖像；我们怎能理解圣父没有将那日和那时——万代的最小部分——显示给祂借着祂创造了万代的那一位呢？宇宙的创造者怎能缺乏祂所创造的极小部分的知识呢？当祂说末日临近时，天上地下将有如此这般的征兆，祂怎能不知末日本身呢？当祂说\Quote{但还不是结局}时，\BibleRef{玛 24:6} 祂是在明确陈述，似乎是有知识而不是不确定的。此外，对公允的探究者来说，显然我们的主以人的身份对人说许多话；例如，\Quote{请给我水喝} \BibleRef{若 4:7} 是我主因身体需要而说的话；然而，求水者并不是没有灵魂的肉身，而是使用有灵魂的肉身的天主性。同样，在目前的情况下，没有人会越出真正宗教解释的界限，如果他将那位根据救恩计划（\OriginalTerm{救恩计划}{œconomy}）接受了一切，并在恩宠和智慧上同天主和人一同进步者的无知理解为出于其人性的话。\par

2. 将福音的句子并排对照，比较玛窦和马尔谷的句子，是值得你勤勉的，因为只有这两部福音在此处吻合。玛窦的措辞是：\Quote{至于那日子和那时辰，没有人知道，连天上的天使也不知道，唯独我的父知道。} 马尔谷的措辞是：\Quote{至于那日子和那时辰，没有人知道，连天上的天使也不知道，圣子也不知道，唯独父知道。} \BibleRef{谷 13:32} 这些经文中值得注意的是，玛窦没有提到圣子的无知，并且似乎在说\Quote{唯独我的父}时与马尔谷意思一致。现在，我理解\Quote{唯独}一词的使用是与天使相对，而圣子并不与祂的仆人们一同被包括在无知之中。\par

祂不能说谎，祂曾说：\Quote{凡父所有的，都是我的}；\BibleRef{若 16:15} 但父所有之一就是那日子和那时辰的知识。因此，在玛窦的经文中，主没有提到祂自己的位格，这是没有争议的，祂说天使不知道，唯独祂的父知道，含蓄地宣称祂父的知识也是祂自己的知识，因为祂在别处说过：\Quote{如同父认识我，我也认识父}；\BibleRef{若 10:15} 如果父完全认识圣子，毫无例外，以致祂知道一切知识都居于圣子内，那么圣子也必同样被认识为拥有祂固有的一切智慧和一切对未来之事的知识的充满者。我想，可以这样解释玛窦的话：\Quote{唯独我的父。} 至于马尔谷的话，他似乎明确地排除圣子于知识之外，我的看法如下：没有人知道，连天主的天使也不知道；然而，圣子若不是因为父知道，也不会知道；也就是说，圣子知道的原因来自父。对公允的听者而言，这种解释并不牵强，因为玛窦中没有加上\Quote{唯独}一词。因此，马尔谷的意思如下：至于那日子和那时辰，没有人知道，连天主的天使也不知道；但圣子若父不知道，也不会知道，因为自然属于祂的知识是由父赐予的。这种说法很得体，并符合圣子的天主性，因为祂的知识和祂的存在，都是由与祂\OriginalTerm{同体}{consubstantial}的那一位所赐，并在一切适宜于祂天主性的智慧和光荣中被瞻仰。\par

3. 关于\Person{耶苛尼雅}，\Person{耶肋米亚}先知用以下的话宣告他被排斥于\Place{犹大}地之外：\Quote{耶苛尼雅如同一件无用的器皿，受人轻贱；因为他被逐出，他和他的后裔；他的后裔中再没有一人坐在\Person{达味}的宝座上统治\Place{犹大}，}。这事是清楚明白的。\Person{拿步高}摧毁\Place{耶路撒冷}时，王国已被消灭，不再有从前那样的世袭继承。然而，在那时，被废黜的\Person{达味}后裔仍生活在充军之中。\Person{沙耳提耳}和\Person{则鲁巴贝耳}归来后，最高统治权更多地归于人民，后来由于司祭支派和君王支派的混合，主权转到了司祭职；因此，主在有关天主的事上，既是君王，也是大司祭。此外，君王支派直到基督来临之前并未断绝；然而，\Person{耶苛尼雅}的后裔不再坐在\Person{达味}的宝座上。显然，\Quote{宝座}一词描述的是君王尊位。你记得历史：整个\Place{犹太}、\Place{厄东}、\Place{摩阿布}、\Place{叙利亚}邻近地区直至\Place{美索不达米亚}的遥远地区，以及另一边直至\Place{埃及}河的\Place{埃及}地区，都曾向\Person{达味}纳贡。如果他的后裔没有一人具有如此广大的主权，那么先知的话——\Person{耶苛尼雅}的后裔再无人坐在\Person{达味}的宝座上——怎会不真实呢？因为他的后裔中似乎没有一人达到这一尊位。然而，\Place{犹大}支派并未断绝，直到那预定的一位来到。但即使是他，也没有坐在物质的宝座上。\Place{犹太}的王国转给了\Place{阿市刻隆}人\Person{安提帕特}的儿子\Person{黑落德}及其儿子们，他们将\Place{犹太}分为四个公国，当时\Person{比拉多}任总督，\Person{提庇留}是罗马帝国皇帝。主所坐的\Person{达味}宝座，他指的是不可毁灭的王国。他是外邦人所期待的\BibleRef{创 49:10}，并非世上最小的一份，因为经上记载：\Quote{到那日，\Person{叶瑟}的根必要站立，作为万民的旗帜；外邦人将寻求他。}\Quote{我召叫你……作人民的盟约，外邦人的光明。}因此，\Person{天主}虽然没有接受\Place{犹大}的权杖，却仍然是司祭，是普世的君王；这样，\Person{雅各伯}的祝福就应验了，在他内\BibleRef{创 22:18}\Quote{万民都要蒙受祝福，}万国都要称基督为有福。\par

4. 至于诙谐的\OriginalTerm{恩克拉提派}{Encratites}提出的那个可怕的问题，为什么我们不吃一切东西？可以这样回答，我们厌恶自己的排泄物。就尊严而言，对我们来说，肉如同草；但就什么有用什么无用而言，正如在蔬菜中，我们区分有害和有益，同样在肉中，我们区分适于食物和有害于食物的。毒芹是蔬菜，正如兀鹰的肉是肉；然而，一个理智的人不会吃天仙子或狗肉，除非他处于极度困境。但如果他吃了，他并不会犯罪。\par

5. 至于那些认为人类事务受命运支配的人，不要向我请教，而是用他们自己的修辞之箭刺穿他们。这个问题在我目前虚弱的状态下太长了。关于在圣洗中浮现——我不知道你怎么会想到问这样的问题，如果你确实明白浸入是为了符合三天的图像。没有人可以被浸入三次而不浮现三次。我们写\OriginalTerm{贪食者}{\GreekText{φάγος}}一词是次末重音的。\par

6. \OriginalTerm{本质}{\GreekText{οὐσία}}和\OriginalTerm{位格}{\GreekText{ὑπόστασις}}之间的区别，就像一般和特殊之间的区别一样，例如，动物与个别的人之间的区别。因此，关于天主性，我们宣认一个本质或实体，以免对存在给出不同的定义，但我们宣认个别的位格，以便我们对圣父、圣子和圣神的概念清晰而不混淆。如果我们对各自的特性，即父性、子性和圣化，没有清晰的认识，而是从存在的普遍概念形成对天主的认识，我们就不可能对信仰给出正确的说明。因此，我们必须将特殊添加到共同之中来宣认信仰。天主性是共通的，父性是特殊的。因此，我们必须将两者结合并说：\Quote{我信天主圣父。}同样，在宣认圣子时，也必须遵循同样的步骤：我们必须将特殊与共同结合并说：\Quote{我信天主圣子，}同样，在圣神的情况下，我们必须使我们的表述符合称呼，并说：\Quote{在天主圣神。}由此，通过宣认一个天主性，令人满意地保存了合一，同时在各自中考虑的个别属性的区分中，宣认了位格的特殊属性。另一方面，那些将本质或实体与位格等同的人，被迫只宣认三位，并且由于他们犹豫不提三个位格，而被证明未能避免\Person{撒伯里乌}的谬误；因为甚至\Person{撒伯里乌}本人，尽管在许多地方混淆了概念，却通过主张同一位格为了当时需要改变形式，确实努力区分了位格。\par

7. 最后，关于你询问的中性及无关紧要的事物是如何为我们安排的，是出于某种偶然的自发作用，还是出于天主的公义眷顾，我的回答如下：健康与疾病、富足与贫穷、声望与耻辱，既然它们并不使拥有者成为善人，就不属于本性为善的范畴；但就它们能在某种程度上使生命之流更顺畅而言，在每种情况下，前者比其反面更可取，并具有某种价值。这些事物中，有些是天主为着管家职务的缘故赐予某些人的，例如赐予\Person{亚巴郎}、\Person{约伯}以及类似的人。对于品格较低的人，它们则是促其改善的挑战。因为那在领受了天主如此美好的爱之标记后仍坚持不义的人，使自己陷入无可辩驳的定罪。善人既不因拥有财富而倾心，也不在没有财富时去追求。他视所赐之物为赐予他，并非为了自私的享乐，而是为了明智的管理。除非有人试图博取世人的称赞——世人钦佩并嫉妒任何掌权者——否则头脑清醒的人不会自找麻烦去分配他人的财产。\par

善人视疾病如同运动员的比赛，等待着酬报他们坚忍的冠冕。若将这些事物的安排归诸其他任何一位，既不合乎真宗教，也不合乎常识。\par

\SourceRecord{Translated by Blomfield Jackson.；Nicene and Post-Nicene Fathers, Second Series；Vol. 8.；Edited by Philip Schaff and Henry Wace.；Buffalo, NY: Christian Literature Publishing Co.,；1895.；<http://www.newadvent.org/fathers/3202236.htm>.}

% structure: p0001 source=h1 target=section reason=page-title

\section{书信237}

\Emph{\Person{圣巴西略·凯撒勒雅}}\par

致\Place{萨摩撒塔}主教\Person{欧瑟比}。\par

1. 我通过\Place{色雷斯}的代理官给尊座写过信，又通过一位从我国前往\Place{色雷斯}的\Place{腓立波波利}财务官员送了其他信，并恳请他出发时带上这些信。但代理官从未收到我的信，因为当我巡视教区时，他在傍晚进城，清晨就出发了，以致教会职员不知道他来过，信便留在了我家。财务官也因意外且紧急的事务，没见到我就出发了，也没带上我的信。找不到其他人了；所以我留了下来，因无法给尊座写信，也收不到尊座的信而感到遗憾。然而，我仍希望，若可能的话，将每日发生的所有事都告诉尊座。发生了如此多令人惊奇的事，以至于需要每日记述；尊座可放心，若非事务的压力使我心意转移，我早就写下来了。\par

2. 我烦恼中首要且最大的事是代理官的来访。他是否真有异端思想，我不知道；因我认为他对于教义全然外行，对此毫无兴趣和经历，因为我日日夜夜见他身心灵忙于其他事。但他确实是异端者的朋友；他对他们的友善并不亚于他对我的敌意。他在隆冬时节于\Place{加拉达}召集了一群恶人的主教会议。他罢黜了\Person{希普西努斯}，立了\Person{厄克迪西乌斯}在其位。他仅凭一个人的控诉就下令罢免我的兄弟，而那人微不足道。然后，他在军中忙了一些时日，又来到我们这里，怒气冲冲，杀气腾腾，一句话便将\Place{凯撒勒雅}整个教会交给了议会。他在\Place{塞巴斯特}停留了数日，将朋友与敌人分开，将与我共融的人称为议员，判处他们从事公务，而提升\Person{尤斯塔修斯}的追随者。他下令在\Place{尼撒}召开第二次\Place{加拉达}和\Place{本都}主教的主教会议。他们顺从了，聚集了，并向各教会派了一个人——我不愿谈论他的品性；但尊座足可理解，一个甘愿受如此计谋支配的人会是何种人。此刻，我正在写这些事时，同一伙人纷纷赶到\Place{塞巴斯特}与\Person{尤斯塔修斯}联合，与他一同扰乱\Place{尼科波利斯}教会。因为可敬的\Person{德奥多图斯}已安眠。至今，尼科波利斯人勇敢而坚定地抵抗了代理官的第一波攻击；因为他试图说服他们接纳\Person{尤斯塔修斯}，并接受他任命的主教。但见他们不愿屈服，他正试图通过更激烈的行动，促成那曾试图给予他们的主教的设立。此外，据说有传言期待一次主教会议，借此他们打算召我接纳他们共融，或与他们友好。这便是各教会的处境。至于我个人的健康，我最好不说。我不忍不说实话，而说实话只会使尊座忧伤。\par

\SourceRecord{Translated by Blomfield Jackson.；Nicene and Post-Nicene Fathers, Second Series；Vol. 8.；Edited by Philip Schaff and Henry Wace.；Buffalo, NY: Christian Literature Publishing Co.,；1895.；<http://www.newadvent.org/fathers/3202237.htm>.}

% structure: p0001 source=h1 target=section reason=page-title

\section{\WorkTitle{238号书信}}

\Emph{圣巴西略（凯撒勒雅）}\par

\Emph{致\Place{尼科波利}的司铎们。}\par

我已收到你们的来信，我可敬的弟兄们，但信中并未告诉我任何我不知道的事，因为整个周边地区早已充满传言，宣告你们中间那位跌倒者的耻辱——他因贪图虚荣而自招极其羞耻的不光彩，并因自爱而丧失了信德所应许的赏报。不仅如此，由于敬畏主者公义的憎恨，他甚至失去了那点微末的虚名，为了贪求虚名，他竟将自己卖给了不虔敬。通过他如今所显露的品性，他已非常清楚地证明，在他整个一生中，从未曾活在对主为我们所存留之应许的望德中，反而在他所有的人间事务中，使用信德的话和嘲弄虔诚，以此欺骗他所遇到的每一个人。然而，你们受了什么损害？你们比从前更糟吗？你们中一人跌倒了，就算有一两个人随他而去，他们的跌倒值得怜悯，但靠着天主的恩宠，你们的身体是完整的。无用的部分已经离去，剩下的并未受损。你们或许因被驱逐出城墙而痛苦，但你们必住在天上天主的护佑下，而那守护教会的天使也与你们一同出去了。他们日复一日地躺卧在空旷之处，因民众的分散而使自己招致沉重的审判。如果在这一切中有忧苦需要承受，我在主内相信，这对你们必不会没有益处。因此，你们的考验越多，就越要从公义的审判者那里领受更完美的赏报。不要因当前的磨难而烦恼。不要失去望德。再过不久，你们的帮助者就要来到你们那里，必不迟延。\BibleRef{哈 2:3}\par

\SourceRecord{Translated by Blomfield Jackson.；Nicene and Post-Nicene Fathers, Second Series；Vol. 8.；Edited by Philip Schaff and Henry Wace.；Buffalo, NY: Christian Literature Publishing Co.,；1895.；<http://www.newadvent.org/fathers/3202238.htm>.}

% structure: p0001 source=h1 target=section reason=page-title

\section{信函239}

\Emph{圣巴西略·凯撒勒雅}\par

致\Place{撒摩撒塔}主教\Person{欧瑟比}。\par

1. 主赐予我现在藉我们敬爱的、极可敬的兄弟，司铎\Person{安提约基}向你的圣德致敬的恩惠，劝勉你如常为我祈祷，并在我们书信的交流中为我们长久的分离提供一些安慰。当你祈祷时，我请求你向主恳求这第一且最大的恩惠，即我能脱离卑鄙邪恶之人，他们对人民拥有如此大的权势，以至于现在我似乎确实看到了\Place{耶路撒冷}陷落事件的重演。因为教会越软弱，人们的权力欲就越增强。如今主教的名衔竟被授予了悲惨的奴隶，因为没有天主的仆人愿意站出来反对要求那个教座——除了像\Person{欧依比乌}的党羽\Person{阿尼西乌}和\Person{帕纳索斯的厄克迪奇乌}的使者这样可怜的家伙：无论谁任命了他，就是把一种可怜帮助自己进入来生的手段送进了教会。\par

他们把我的兄弟从\Place{尼撒}赶走了。取代他，他们引入了一个几乎算不上人的人——一个只值一两块钱的无赖，但就败坏信德而言，他与那些把他放在这个位置上的人不相上下。\par

在\Place{多阿拉}镇，他们玷污了主教这可怜的名号，并送去一个可怜虫，孤儿院的仆役，从他主人那里逃出来的，去谄媚一个不信天主的女人，她以前随意利用\Person{乔治}，现在又让这个家伙接替他。\par

谁能适当地哀叹在\Place{尼科波利}发生的事呢？那个不幸的\Person{弗龙托}，的确有一段时间假装站在真理一边，但现在他可耻地背叛了信德和自身，并以背叛的代价换来了一个可耻的名声。他想象自己从这些人那里得到了主教的品级；实际上，因天主的恩宠，他成了全\Place{亚美尼亚}的厌恶对象。但他们没有不敢做的事；他们也从不缺少合适的同谋。但\Place{叙利亚}的其他消息，我的兄弟比我更清楚，也更能告诉你。\par

2. 西方的消息你已经从兄弟\Person{多罗泰}的叙述中知道了。他离开时该给他什么信函？也许他会与非常热心的卓越的\Person{桑克蒂西穆斯}同行，游历东方，从所有知名人士那里收集信函和签名。他们该写什么，或者我如何与写信者达成一致，我不知道。如果你听说有人很快就走我这条路，请通知我。我不禁要说，如同\Person{狄奥墨得}所说，\par

\Quote{但愿，\Person{阿特瑞德斯}，你的请求尚未承担；\newline{}……他够骄傲了。}\par

真正高尚的灵魂，当被奉承时，变得比以往更傲慢。若主眷顾我们，我们还缺什么？若主的愤怒持续，西方的皱眉又能给我们什么帮助？那些不认识真理、也不愿学习真理、却被虚假猜疑所偏见的人，现在正如同他们在\Person{玛尔凯鲁}事件中所做的一样，当时他们与告诉他们真理的人争吵，并以自己的行动助长了异端的事业。除了共同的文件外，我还想写信给他们的领袖——但确实不涉及教会事务，只是温和地暗示他们对这里发生的事一无所知，并且不接受唯一能让他们了解此事的方法。我一般要说，他们不应该苛待被试炼压垮的人。他们不可把尊严当作傲慢。罪只能产生对天主的敌意。\par

\SourceRecord{Translated by Blomfield Jackson.；Nicene and Post-Nicene Fathers, Second Series；Vol. 8.；Edited by Philip Schaff and Henry Wace.；Buffalo, NY: Christian Literature Publishing Co.,；1895.；<http://www.newadvent.org/fathers/3202239.htm>.}

% structure: p0001 source=h1 target=section reason=page-title

\section{\WorkTitle{书信第二百四十}}

\Emph{\Person{凯撒勒雅的圣巴西略}}\par

致\Place{尼科波利斯}的司铎们。\par

1. 你们给我来信，并且通过一位即使你们没有写信也完全能够在我所有的焦虑中给我莫大安慰、并能如实报告事态的人送来，这是做得非常正确的。不断有许多不确定的传闻传到这里，因此我一直希望从某个能够凭借准确知识提供信息的人那里了解许多情况。关于这一切，我从我们可敬爱而尊贵的弟兄、司铎\Person{德奥多西}那里得到了令人满意而清晰的叙述。现在我把我给自己的劝告写给你们诸位尊者，因为我们的处境在许多方面是相同的；这不仅是在目前，而且在过去的时代也是如此，许多事例可以证明。其中有些我们拥有文字记录，另一些则通过没有文字记录的回忆从了解事实的人那里得来。我们知道，为了主的名，考验曾降临于个人和那些信赖祂的城市。然而，这些考验一个接一个地过去了，黑暗日子带来的痛苦并不是永远的。因为正如冰雹和洪水以及所有自然灾害，既能伤害并摧毁软弱之物，而它们本身只有在撞击坚固之物时才会受到影响，同样，那针对教会发动的可怕考验，已被证明比我们在基督内的信德的坚实基础更为软弱。冰雹过去了；洪流冲过了它的河床；晴朗的天空取代了前者，而后者留下了它所流过的无水干涸的河床，并消失在深渊中。同样，再过片刻，那现在向我们袭来的风暴也将止息。但这将取决于我们是否愿意不看眼前，而是怀着希望注视稍远一点的未来。\par

2. 我的弟兄们，考验是沉重的吗？让我们忍受这劳苦。没有一个人能回避战斗的打击和尘土而赢得冠冕。魔鬼的嘲弄和那被派来攻击我们的仇敌是微不足道的吗？他们之所以麻烦，是因为他们是他的差役，但他们是可鄙的，因为天主在他们身上将邪恶与软弱结合在一起。让我们谨慎，免得因一点小痛苦就大声喊叫而受谴责。只有一件事值得痛苦，那就是失去自我，当一个人为了片刻的声望——如果真能把自己当众出丑称为声望——而剥夺了自己义人永久的赏报时。你们是宣信者的子女，你们是殉道者的子女；你们抵抗罪恶直到流血。你们每个人都要运用身边亲人的榜样，为真正的宗教而勇敢。我们中没有一个人被鞭打撕裂，没有一个人遭受房屋被没收；我们没有被迫流亡，没有遭受监禁。我们经历了什么大苦呢？除非或许我们什么也没有经历，没有被认为配得上基督的苦难，这倒是可悲的。但如果你们悲伤，因为某位我无需提其名的人占据了祈祷之所，而你们在露天敬拜天地的主，那么请记住：当那些钉死主的人在犹太人闻名的圣殿里敬拜时，十一位门徒被关在楼上的房间里。或许，\Person{犹达斯}选择了吊死而非在耻辱中活下去，他证明了自己比那些现在面对普遍的谴责而毫不脸红的人更好。\par

3. 只是不要被他们的谎言所欺骗，当他们声称拥有正统信德时。他们不是基督徒，而是基督的贩子，总是偏爱今生的利益，胜过按真理生活。当他们以为能得到这空虚的尊荣时，他们就加入了基督的敌人；现在他们看到民众的愤慨，便又一次假装正统。我不承认他是主教——我决不把他算作基督的圣职人员——一个通过被玷污的手被提升到首位，以致破坏信德的人。这是我的决定。如果你们与我有什么共融，你们无疑会与我想法一致。如果你们自己负责商议，那么每个人都是自己思想的主人，而我在这个血案上是清白的。我这样写，不是因为我不信任你们，而是为了通过表明我的想法来坚定某些人的犹豫，并防止任何人在和平实现后过早地被接纳进入共融，或者在接受了我们敌人的覆手之后，后来试图强迫我把他列入神圣职务的行列。通过你们，我问候城市和教区的圣职人员，以及所有敬畏主的平信徒。\par

\SourceRecord{Translated by Blomfield Jackson.；Nicene and Post-Nicene Fathers, Second Series；Vol. 8.；Edited by Philip Schaff and Henry Wace.；Buffalo, NY: Christian Literature Publishing Co.,；1895.；<http://www.newadvent.org/fathers/3202240.htm>.}

% structure: p0001 source=h1 target=section reason=page-title

\section{书信二四一}

\Emph{\Person{圣巴西略}·\Place{凯撒勒雅}}\par

\Emph{致\Place{萨莫萨塔}主教\Person{厄色佳}}\par

我并非要加重您的痛苦，才在写给阁下的信中大量触及令人伤心的话题。我的目的是在哀叹中为自己寻求一些安慰——哀叹是一种天然的宣泄，能驱散深埋的痛苦——同时也是为了激励您，我宽厚的朋友，为教会作更恳切的祈祷。我们知道\Person{梅瑟}曾不停地为百姓祈祷；然而，当他与\Person{阿玛肋克}的战斗开始时，他从早到晚都没有垂下双手；圣者的举手只随着战斗的结束而结束。\par

\SourceRecord{Translated by Blomfield Jackson.；Nicene and Post-Nicene Fathers, Second Series；Vol. 8.；Edited by Philip Schaff and Henry Wace.；Buffalo, NY: Christian Literature Publishing Co.,；1895.；<http://www.newadvent.org/fathers/3202241.htm>.}

% structure: p0001 source=h1 target=section reason=page-title

\section{书信二四二}

\Emph{ \OriginalTerm{圣巴西略·凯撒勒雅}{ST. BASIL OF CAESAREA} }\par

\Emph{致西方人。}\par

圣洁的天主曾应许，凡信赖祂的人，必能从一切困苦中获得幸福的结局。因此，我们虽然被弃于苦难的汪洋之中，虽然被邪恶之灵掀起的巨浪所颠簸，但在那加强我们力量的基督内，我们仍然坚持。我们未曾减弱对教会热忱的力量，也不曾因绝望于自己的救恩，当暴风雨中的波涛高过我们的头顶时，就认为自己将被毁灭。相反，我们仍以最大的热诚坚持，记念那被海怪吞下的人，因他不绝望于自己的性命，而向主呼求，便得了救恩。同样，我们虽然已到达危险的极点，也不放弃对天主的希望。我们处处看见祂的援助环绕我们。因此，我们现在将目光转向你们，可敬的弟兄们。在我们受苦的许多时刻，我们曾期望你们会与我们同在；但希望落空，我们便自言自语说：\Quote{我指望有人怜悯，却没有一个；指望有人安慰，却寻不着一个。}我们的苦难已达到帝国的边境；既然一个肢体受苦，所有肢体都一同受苦，\BibleRef{格前12:26}，那么你们理应对我们这些长久处于困境的人表示怜悯。因为你们因爱德而对我们所怀的同情，与其说是由地域的接近，不如说是由精神的合一而来的。\par

2. 然而，为何我们按爱的律法当得之物，一无所获？没有安慰的书信，没有弟兄的探访。如今已是异端之战攻击我们的第十三个年头。在此期间，教会所受的磨难，远超过自基督福音最初宣讲以来所有记载的磨难。我不愿一一细述，唯恐我叙述的无力使灾祸的证据显得不够确凿。况且，我原无需向你们讲述，因为你们早已从传至你们那里的报告中知道了所发生的事。我们困境的总结与实质是：百姓离开了祈祷之所，在旷野中聚集。那是一幅悲惨的景象。妇女、孩童、老人，以及其他体弱多病者，留在露天之中，在暴雨、大雪、寒风与冬霜中，也在夏日灼热的阳光下。他们忍受这一切，是因为他们拒绝与亚略那邪恶的酵有任何关联。\par

3. 仅仅言语如何能让你清楚明了这一切，而不需要亲身经历与目击者的证词来激起你的同情？因此，我们恳求你们，向那些已倒地的人伸出援手，并派遣使者来提醒我们那为所有为基督忍耐受苦者所预备的奖赏。我们习惯的声音自然不如远方传来的声音那样能安慰我们，而那声音来自那些天主的恩宠使他们在普世皆知为最尊贵的人；因为到处都有传言说你们持守坚定，在信德上未受创伤，并保持了宗徒的纯正信仰传承。但我们并非如此。我们中间有些人，由于贪求荣耀和那种尤其常毁灭基督徒灵魂的自高自大，竟胆敢说出一些新奇的措辞，结果使教会变得像破裂的锅碗，让异端的不洁涌入。但你们，我们所爱所渴望的，请为我们作伤者的外科医生，为健全者的教练，医治有病的肢体，并为真宗教的职事傅油那健全的肢体。\par

\SourceRecord{Translated by Blomfield Jackson.；Nicene and Post-Nicene Fathers, Second Series；Vol. 8.；Edited by Philip Schaff and Henry Wace.；Buffalo, NY: Christian Literature Publishing Co.,；1895.；<http://www.newadvent.org/fathers/3202242.htm>.}

% structure: p0001 source=h1 target=section reason=page-title

\section{书信第二四三号}

\Emph{\Person{凯撒勒雅的圣巴西略}}\par

关于教会状况与混乱，致意大利和高卢的主教们。\par

1. 致我们真正属天主所爱、至亲的弟兄，以及心意相同的同工——高卢和意大利的主教们：凯撒勒雅（在卡帕多细亚）的主教巴西略。我们的主耶稣基督，祂曾屈尊称普世天主的教会为祂的身体，并使我们彼此互为肢体，更恩赐我们众人得以彼此亲密结合，相称于肢体的和谐。因此，虽然我们彼此相距遥远，但在紧密的联系上，我们却近在咫尺。既然头不能对脚说：\BibleQuote{我不需要你们}\BibleRef{格前 12:21}，我相信你们必不会忍心拒绝我们；反之，你们必会同情我们因犯罪所受的苦难，正如我们因主所赐予你们的平安而一同喜乐。以前我们也曾另有一次恳请你们的爱德，派遣援助和同情；但我们的惩罚还未满全，你们也未被允许起来援助我们。我们的一大愿望是，藉着你们，我们这混乱的状况能传达给你们那地区的皇帝。如果这有困难，我们恳求你们派遣使者来探访并安慰我们的痛苦，好使你们拥有目击者的证据，亲眼目睹东方无法以言语述说的苦楚，因为言语不足以清楚报告我们的处境。\par

2. 最有尊荣的弟兄们，迫害已临到我们身上，而且是极其严酷的迫害。牧者受迫害，好使羊群四散。而最糟的是，受虐待的人不能将他们的痛苦当作殉道的证据，人民也不能尊敬这些斗士如同殉道者的军队，因为基督徒的名号竟被加在迫害者身上。现今唯一必定招致重罚的罪名，就是谨慎持守祖传传统。为此，虔诚者被逐出家园，流放到遥远地区。不义的审判官毫不尊重白发苍苍的老人、实际的虔诚、以及自幼至老按照福音生活的人。没有一个罪犯在无证据下被定罪，但主教们仅凭毁谤就被定罪，并基于全无证据的指控而受到刑罚。有些人甚至不知道谁控告了他们，也没有被带到任何法庭面前，甚至从未被诬告。他们在深夜被暴力逮捕，被流放到遥远之地，并因这些偏远荒野的艰辛而丧命。其余的事尽管我不提，也昭然若揭——司铎的逃亡、执事的逃亡、全体圣职人员的劫掠。要么崇拜偶像，要么我们就被交付给鞭笞的恶火。平信徒呻吟；公共和私下里眼泪不绝；人人相互哀叹自己的苦难。没有人的心肠如此坚硬，以至于失去父亲而能温顺地忍受丧亲之痛。城里有哀哭的声音——在田野、在道路、在旷野也有声音。但所有发出悲伤可怜话语的人只发出一个声音。喜乐和属灵的欢愉已被夺去。我们的庆节已变为悲哀\BibleRef{亚 8:10}。我们的祈祷之所已关闭。属灵事奉的祭台闲置。基督徒不再聚集；教师不再主持。救恩的道理不再被传授。我们不再有庄严的集会，不再有晚祷的圣歌，不再有那因领受圣体及分享属灵恩惠而在所有信主的人灵魂中产生的蒙福喜乐。我们真可以说：\Quote{此时没有君王，没有先知，没有读经者，没有祭品，没有乳香，没有在祢面前献祭并找到仁慈的地方。}\par

3. 我们正在写信给了解这些事的人，因为世上没有一个地区不知道我们的灾祸。不要以为我们是用这些话来提供信息，或使我们重新被你们记起。我们知道，你们不会忘记我们，如同母亲不会忘记她腹中的儿子一样\BibleRef{依 49:15}。但所有被痛苦重压的人，在发出痛苦呻吟时，都会找到某些自然的减轻痛苦的方法；我们这样做正是为此。我们向你们诉说我们多重的灾祸，以此卸下我们悲伤的重担，并表达希望你们或许会更加感动而为我们祈祷，并能求主与我们和好。如果这些苦难仅限于我们自己，我们甚至可能决定保持沉默，并为基督的缘故在苦难中喜乐，因为\Quote{现时的苦楚，与将来在我们身上要显示的光荣，是不能较量的}\BibleRef{罗 8:18}。但目前我们担忧，恐怕这祸患日益增长，像火焰在燃烧的木材中蔓延，烧尽近处之物后，也会波及远处之物。异端的瘟疫正在蔓延，我们有理由担心，当它吞食了我们的教会后，甚至会偷偷蔓延到你们地区中健全的部分。或许是因为我们这里罪恶满盈，所以我们首先被交付给天主仇敌的残酷牙齿吞噬。但天国的福音始于我们的地区，然后传遍全世界。因此，或许——这是最可能的——我们灵魂的共同仇敌正竭力使背教的种子从同一源头散布到全世界。因为不虔敬的黑暗图谋临到那曾蒙受基督\Quote{认识的光照}的心上。\par

4. 那么，你们作为主的真门徒，要视我们的苦难为你们的苦难。我们受攻击不是为了财富、荣耀或任何暂时利益。我们站在竞技场上，为我们的共同遗产，为从祖先传承的纯正信德的宝藏而战。所有爱弟兄的人，请与我们一同哀伤，因为我们真正虔敬之人的口被封住，而那些对天主说不义之言的狂妄亵渎之口却大开。真理的柱石和根基被驱散。我们这些因卑微而得以被忽视的人，被剥夺了言论自由。请你们为人民的缘故投入这场斗争。不要只顾自己安稳地停泊在安全港湾，那里天主的恩宠保护你们免受邪恶风暴的袭击。向那些在风浪中颠簸的教会伸出援手，以免它们被抛弃后，遭遇信德的彻底沉船。为我们哀哭，因为独生子被亵渎，却无人反驳。圣神被藐视，能驳斥谬误的人被流放。多神论占了上风。我们的对手尊一位大神和一位小神。\Quote{圣子}不再是本性的名称，而被视为某种尊荣的头衔。圣神不被视为天主圣三的圆满，也不参与神圣可敬的本性，而是被看作受造物中的一员，随意地依附于圣父和圣子。\BibleQuote{但愿我的头是水，我的眼是泪的泉源} \BibleRef{耶 9:1}。我要为那些被这些邪恶教义引向毁灭的人多日哭泣。简单人的耳朵被误导，现在已习惯于异端的不敬。教会的乳儿在邪恶的教义中成长。他们该怎么办？我们的对手掌管圣洗，为临终者送行，探望病人，安慰悲伤者，帮助困苦者，提供各种援助，分施奥迹。只要这一切由他们执行，就都是系紧人民于他们观点的纽带。结果就是，即使稍后我们获得一些自由，那些长期受错误影响的人也几乎没有希望被召回认识真理。\par

5. 在这种情况下，我们许多人本应前往你们尊前，各自报告自己的处境。现在你们可以从我们甚至不能自由外出旅行这一事实，看出我们所处的困境。因为如果有人离开他的教会，哪怕很短时间，就会把自己的百姓留给那些图谋毁灭他们的人。赖天主的仁慈，我们代替众人派遣了一位，就是我们可敬可爱的弟兄司铎多罗岱。他完全能够通过个人报告补充我们信中所遗漏的任何事情，因为他仔细跟进了所发生的一切，并热心于纯正信德。请平安地接待他，并尽快送他回来，带给我们你们乐意援助弟兄们的好消息。\par

\SourceRecord{Translated by Blomfield Jackson.；Nicene and Post-Nicene Fathers, Second Series；Vol. 8.；Edited by Philip Schaff and Henry Wace.；Buffalo, NY: Christian Literature Publishing Co.,；1895.；<http://www.newadvent.org/fathers/3202243.htm>.}

% structure: p0001 source=h1 target=section reason=page-title

\section{Letter 244}

\Emph{\Person{圣巴西略} \Place{凯撒勒雅}}\par

\Emph{致\Place{埃盖}主教\Person{帕特罗斐禄}}\par

1. 我愉快地，非常愉快地阅读了你由司铎\Person{斯特拉特吉乌斯}带来的信。我怎能不这样阅读呢？这封信出自一位智者之手，并由一颗学会了遵守主诫命所教导的普世的爱的心所口述。或许我并非不知道您迄今保持沉默的原因。您仿佛是惊愕和震惊，一想起著名的\Person{巴西略}的改变。为什么？从他小时候起，他就如此这般地为某人服务；在某某时候他做了某某事；他为了效忠于一人而与无数敌人作战；现在他完全变成了另一个人；他以战争取代了爱；他就是您所写的一切；因此您对如此出乎意料的事态转变表现出相当的惊讶也是自然的。如果您发现了什么过错，我也不会介意。我并非不可纠正，以至于对弟兄们的亲切责备感到惊讶。事实上，我不仅没有因您的信而烦恼，反而几乎让我发笑，想到当我认为已经有很多强有力的原因存在以巩固我们的友谊时，您却对这些报告给您的琐事表现出如此巨大的惊讶。您确实遭了所有那些不去探究情况本质，而只关注被讨论的人之命运；所有那些不考察真理，而以人的身份判断，忘记了\Quote{你们审判时不可徇私}的劝诫。\BibleRef{申 1:17}\par

2. 然而，既然天主审断人不看情面，我也不拒绝向您说明我为那大审判所准备的辩护。从我这一方来说，从一开始就没有任何争吵的原因，无论大小；但是那些恨我的人——原因只有他们自己最清楚（我必须对他们不置一词）——不停地诽谤我。我一次又一次地澄清了这些诽谤。事情似乎没有尽头，我不断的辩护也无济于事，因为我远离现场，而捏造谎言的人近在咫尺，能够用他们的诽谤伤害一颗易感的心，一颗从未学会为缺席者保留一只耳朵的心。当\Person{尼科波利人}——您自己也有部分了解——要求某种信德的证据时，我决定诉诸书面文件。我想我可以同时实现两个目的；我期望既说服\Person{尼科波利人}不要对这人有恶评，又堵住我诽谤者的口，因为信德的一致将排除双方的诽谤。事实上，信经已经拟定，由我带来并签署。签署后，又指定了第二次会面的地点，并确定了另一个日期，以便我教区的弟兄们可以聚集在一起，彼此合一，我们今后的共融是真诚而恳切的。我这边按时到达，与我同工的弟兄中，有些已在现场，另一些正赶往那里，全都喜乐而热切，仿佛走在通往和平的大道上。信使和我自己的信宣布了我的到达；因为为聚集之人所指定的地点是我的。但对方没有一个人出现；没有一个人提前到来；没有一个人宣布期待的主教们即将到来。所以我派去的人回来报告说，聚集的人深感沮丧和抱怨，好像我颁布了新的信经。据说他们还决定，他们肯定不允许他们的主教投向我。接着，一个信使给我带来了一封仓促写就的信，信中完全没有提及原先商定的事。我的\Person{德奥斐洛}弟兄——一个我应予以一切尊重和敬意的人——派了他的一个追随者，做出了一些声明，他认为说这些话并非不当，我听这些话也无妨。他不屑于写信，与其说是因为他害怕被书面证据定罪，不如说是因为他急于避免被迫称呼我的主教头衔。的确，他的言辞激烈，出自一颗剧烈激动的心。在这种情况下，我羞愧而沮丧地离开了，不知道如何回答质问我的人。然后，没过多久，就有了\Place{基里基雅}之行，从那里返回后，立即收到一封拒绝与我共融的信。\par

3. 分裂的原因是指控我写信给\Person{阿波利纳里}，并与司铎\Person{狄奥多若}共融。我从未视\Person{阿波利纳里}为敌人，甚至出于某些原因尊敬他。但我从未与他如此联合，以至于承担对他的指控；实际上，在阅读了他的一些著作后，我自己也有一些对他的指控。我不知道我曾向他索要关于圣神的著作，或收到他寄来的；我听说他成了多产作家，但我的著作读得很少。我甚至没有时间研究这类问题。事实上，我回避承认任何较新著作，因为我的健康甚至不允许我勤奋地、按应有的方式阅读圣经。那么，如果某人写了令他人不悦的东西，与我有何干系？然而，如果一个人要代他人负责，那么因\Person{阿波利纳里}而指控我的人，就应该因他自己的主人\Person{亚略}和他自己的门徒\Person{埃提乌}而向我辩护。我从未从这位归咎于我的人那里学到什么，也从未教过他什么。至于\Person{狄奥多若}，作为真福\Person{西尔瓦诺}的哺育者，我从一开始就接纳了他；我现在爱他，尊敬他，因他言语的恩宠，许多遇见他的人因此变得更好。\par

4. 这封信对我的影响正如所料，如此突然而可喜的变化使我惊愕。我几乎无法作答。心脏几乎无法跳动；舌头结结巴巴，手也麻木了。我感觉自己像个可怜虫（因为要说实话；但这是可以原谅的）；我几乎陷入厌世的状态；对每个人都投以怀疑的目光，认为人类之中找不到爱德。爱德似乎只是动听的言辞，用作使用之人的装饰，而人的心中并不真正存在这种情感。难道一个似乎从小到老都自我训练的人，竟会因如此微不足道的理由轻易变得粗野，丝毫不顾及我，丝毫没有想过他多年的经验应当比这卑劣的诽谤更有分量？他难道真的像一匹尚未学会正确驮载骑手的未经驯服的小马，因一点小小的猜疑就人立而起，摔下骑手，把曾经引以为傲的丢在地上？如果是这样，那么其他与我并无如此深厚友谊、且未受过如此良好生活训练的人又该如何看待？这一切在我灵魂中翻腾，在我心中不断萦绕，或者不如说，是这些东西在回忆中争斗、刺痛我，从而使我的心被翻腾？我没有写回信；并非出于蔑视而沉默；我的兄弟，不要这样想我，因为我不是在人前为自己辩护，而是在基督内于天主前发言。我沉默是因为完全说不出与我的悲痛相称的话。\newline{}\par

5. 在我处于这种境况时，另一封信寄到了，是写给某位\Person{达齐扎斯}的，但实际上是写给全世界的。这一点从其迅速的传播中显而易见，因为几天之内就在整个\Place{本都}传开，并在\Place{迦拉达}流传；据说这些好消息的传递者还穿过了\Place{彼提尼雅}，甚至到达了\Place{赫勒斯滂}本身。写给\Person{达齐扎斯}反对我的内容你十分清楚，因为他们并不认为你超出他们的友谊范围，以至于不让你得到这一荣誉。然而，如果信没有到你手上，我会寄给你。你在其中会看到我被指控为诡诈和背叛，指控我败坏教会、毁灭灵魂。他们认为最真实的一项指控是，我发表那篇信仰宣言是出于隐秘和不诚实的原因，不是为\Person{尼科波利斯人}服务，而是以不诚实的方式从他们那里骗取信德的宣认。这一切全由主来判断。对心中的思想能有什么清楚的证据？我对他们有一点感到惊奇，就是他们签署了我提交的文件之后，却表现出如此大的分歧，以致混淆真理与谬误以满足那些控告他们的人，完全忘记了他们在罗马保存的尼西亚信经的书面宣认，以及他们亲手交给\Place{提雅纳}会议的、从\Place{罗马}带来的文件（那文件在我手中，包含同一信经）。他们忘记了自己曾经的发言——他们站出来哀叹被骗而同意了\Person{欧多克修}派系起草的文件，于是想到了补救那个错误的办法：他们应该去\Place{罗马}，在那里接受教父们的信经，以便通过引入更好的东西来弥补他们因赞同邪恶而对教会造成的损害。如今，正是那些为信德远行并发表所有这些漂亮言辞的人，却在辱骂我，说我行走诡诈，在爱的外衣下玩弄阴谋。从正在流传的信中可以清楚看出，他们已经谴责了\Place{尼西亚}信经。他们看到了\Place{基齐库斯}，却带着另一信经回家了。\newline{}\par

6. 但何必谈论单纯言辞上的矛盾？他们的行为所提供的立场转变的实践证据要强得多。他们拒绝服从五十位主教针对他们作出的判决。他们拒绝放弃对各自教会的管理，尽管赞同罢免他们决议的主教人数如此之多，理由是那些主教不分享圣神，不是靠天主的恩宠治理教会，而是借助人的权力、出于虚荣的贪欲攫取了他们的尊位。如今他们却准备接受由同一批人所祝圣的人为主教。我希望你替我问问他们（尽管他们藐视所有人，视为没有眼睛、耳朵和常识），他们是否意识到自己行为的矛盾，以及他们内心真正怀有怎样的情感。怎么会有两位主教：一位由\Person{尤伊皮乌斯}罢免，另一位由他祝圣？两者都是同一人的行为。倘若他没有领受赐予\Person{耶肋米亚}的恩宠，即拔除和建造、根除和栽种，他肯定不会根除一个而栽种另一个。承认了前者，就必须承认后者。看来他们的唯一目标就是处处寻求自己的利益，把凡是按他们意愿行事的人都当作朋友，而把任何反对他们的人都当作敌人，并且不遗余力地以诽谤诋毁他。\newline{}\par

7. 如今他们正在采取什么措施反对教会？其始作俑者的诡诈令人震惊，所有受其影响之人的冷漠令人痛惜。欧依丕乌斯的子孙和孙辈经一个可敬的委员会从偏远地区被召到塞巴斯提亚，民众被托付给了他们。他们占据了祭台。他们成了那教会的酵母。我被他们作为同体论者而迫害。欧斯塔修虽然未能获准加入他们那令人羡慕的共融——或因惧怕，或因尊重曾一致谴责他的众多人士的权威——但他从罗马将 \OriginalTerm{同体}{Homoousion} 带到提亚纳，如今却与他们紧密结盟。我只希望我永远不会有足够的时间来讲述他们的所有行径——谁被聚集，每人如何被祝圣，每人从何种先前生活达到现今的地位。我受教祈祷：\Quote{愿我的口不致说出这些人的行为。}你若询问，你自己会了解这些事；若它们对你隐藏，它们必然不会长久对审判者隐藏。\par

8. 然而，我亲爱的朋友，我无法省略不告诉你我的状况。去年我患上猛烈高烧，濒临死亡之门。因着天主的慈悲，我得以复原，但我因想起面前的所有烦恼而对自己重获生命感到痛苦。我心中思忖，在天主智慧深处隐藏的理由中，为何还允许我度过更多肉身的时日。但当我听说这些事时，我断定主愿意我看到教会在此前因那些人的疏远——因他们虚假的庄重品格，人人曾对他们完全信赖——而遭受风暴之后，如今得以安息。或者，主旨在增强我的灵魂，使之未来更加警醒，以便我不再听从人，而是藉着那些不随人类季节和处境变迁而变化，永远如一、由不能说谎的福唇所宣示的福音训诫而得以成全。\par

9. 人如同云彩，随风在空中飘移不定。在我所遇到的所有人中，我此刻所说的这些人最不稳定。至于人生的其他事务，与共处者或可作证；但就我自身所知，他们在信仰上的反复无常，我从未亲眼见过或从他人那里听说过类似的情形。起初他们是亚略的追随者；然后他们转向赫尔莫真，此人明确反对亚略的错误，正如他在尼西亚最初诵读的信经所显明。赫尔莫真安眠后，他们又转向欧瑟伯——我们亲身所知，亚略集团的首领。离开后者，无论出于何种原因，他们回到家乡，再次隐藏其亚略派情绪。在达到主教职位后，且忽略期间所发生的事，他们颁布了多少信经？一次在安齐拉；另一次在塞琉西亚；另一次在君士坦丁堡，那著名的信经；另一次在兰普萨库斯；然后是色雷斯的尼凯那一次；现在又是基齐库斯的信经。关于后者我一无所知，只听说他们已压制了\OriginalTerm{同体}{homoousion}，转而支持\Emph{类似本质}，并与欧诺弥乌一同签署亵渎圣神的言论。尽管我所列举的所有信经可能并非彼此对立，但它们同样显示了这些人思想的不一致，因为他们从未坚持相同的言辞。我对无数其他方面闭口不谈，但我所说的这一点属实。如今他们已转向你们，我请求你同一位使者——我们同是司铎的斯特拉特吉乌——回信告诉我：你对我是否保持同一心意，还是因与他们会面而疏远？因为他们不可能保持沉默，而你本人，既然曾如此写信给我，也不会不对他们直言。你若保持与我的共融，那就很好；这是我极其恳切祈求的。若他们把你拉拢过去，那令人悲伤。与如此一位兄弟分离怎会不悲伤？即使在其他事上，在承受这样的损失时，我们已深受他们的折磨。\par

\SourceRecord{Translated by Blomfield Jackson.；Nicene and Post-Nicene Fathers, Second Series；Vol. 8.；Edited by Philip Schaff and Henry Wace.；Buffalo, NY: Christian Literature Publishing Co.,；1895.；<http://www.newadvent.org/fathers/3202244.htm>.}

% structure: p0001 source=h1 target=section reason=page-title

\section{第245封信}

\Signature{圣巴西略·凯撒勒雅}\par

\Emph{致狄奥斐卢斯主教}。\par

我收到您的信已有时日，但我一直等待能通过合适的人回复，以便我的回信的信使能补充其中可能不足之处。如今我们深爱且极其可敬的弟兄斯特拉特吉乌斯到了，我认为借助他的服务是好的，因为他既了解我的心思，又能得体而恭敬地传达我的消息。因此，我亲爱的、受尊敬的朋友，当知道，我非常珍视我对您的感情，就我内心的情感而言，我并不觉得在任何时候有负于此，尽管我有很多严重且合理的抱怨理由。但我决定如同在天平上衡量好坏，并在善的一面倾斜时添加自己的心意。如今，那些最不应该允许这类事情发生的人却做出了改变。所以，请宽恕我，因为我没有改变心意，如果说我转向了任何一方，或者更确切地说，我仍然站在同一方，但有些人在不断改变立场，如今公开投靠敌人。您自己知道，只要他们属于正统一方，我是多么重视与他们的共融。如果现在我拒绝跟随这些人，并回避所有与他们想法相同的人，我应当得到公平的宽恕。我把真理和自身的救恩放在一切之上。\par

\SourceRecord{Translated by Blomfield Jackson.；Nicene and Post-Nicene Fathers, Second Series；Vol. 8.；Edited by Philip Schaff and Henry Wace.；Buffalo, NY: Christian Literature Publishing Co.,；1895.；<http://www.newadvent.org/fathers/3202245.htm>.}

% structure: p0001 source=h1 target=section reason=page-title

\section{第二百四十六封书信}

\Emph{圣巴西略}\par

致尼科波利坦人。\par

我看见恶势力在成功的大道上一往无前，而你们，我尊敬的各位朋友，却在连续不断的灾祸中疲弱衰竭，内心充满悲痛。但当我再次想到天主的全能之手，想到祂知道如何扶起跌倒的人，如何爱慕正义者、粉碎骄傲者、并从高位上推下有权势者时，我的心又因希望而轻松起来，并且我知道，藉着你们的祈祷，主将要赐予我们的平安很快就会来临。只要你们在祈祷中不厌倦，并在当前的紧急情况下，努力以实际行动为你们口中所教导的一切立下明显的榜样。\par

\SourceRecord{Translated by Blomfield Jackson.；Nicene and Post-Nicene Fathers, Second Series；Vol. 8.；Edited by Philip Schaff and Henry Wace.；Buffalo, NY: Christian Literature Publishing Co.,；1895.；<http://www.newadvent.org/fathers/3202246.htm>.}

% structure: p0001 source=h1 target=section reason=page-title

\section{\WorkTitle{书信二四七}}

\Signature{凯撒勒雅的圣巴西略}\par

致尼科波利人。\par

我读了诸位圣座的信函，如何能不叹息哀恸？我听闻了这些新的苦难：你们遭受拳打脚踢和侮辱，家园被毁，城市遭蹂躏，整个地区沦为废墟，教会受迫害，流放，司铎们被驱逐，豺狼入侵，羊群四散。然而，我仰望天上的主，便不再叹息哀哭，因为我完全确信（希望你们也知道），救援必将迅速来临，你们绝不会永远被遗弃。我们所受的一切，都是因我们的罪而受的。然而，仁爱的主为了祂对众教会的慈爱与怜悯，必将向我们显示祂的援助。尽管如此，我并未忽略亲自恳求当权者。我已写信给宫廷中关爱我们的人，但愿我们那贪婪之敌的怒气得以平息。此外，我认为许多人都会谴责他，除非这动荡的时局使得我们的官吏无暇顾及这些事。\par

\SourceRecord{Translated by Blomfield Jackson.；Nicene and Post-Nicene Fathers, Second Series；Vol. 8.；Edited by Philip Schaff and Henry Wace.；Buffalo, NY: Christian Literature Publishing Co.,；1895.；<http://www.newadvent.org/fathers/3202247.htm>.}

% structure: p0001 source=h1 target=section reason=page-title

\section{书信二四八}

\Person{圣巴西略·凯撒勒雅}\par

致\Person{安菲洛奇乌斯}，\Place{伊科尼雍}主教。\par

就我个人的愿望而言，我为远离尊驾而居感到悲伤。但就您生活的平安而言，我感谢\Person{主}，祂使您免受这场特别蹂躏我教区的火灾。因为公义的\Person{审判者}按我的作为，给我派遣了一位\OriginalTerm{撒殚的使者}{messenger of Satan}，他重重地击打我，并奋力维护\OriginalTerm{异端}{heresy}。事实上，他向我们的战争已到了如此地步，甚至不惜流出信靠\Person{天主}之人的血。您一定听说了，一个名叫\Person{阿斯克勒庇乌斯}的人，因为不肯与\Person{多厄格}共融，被他们殴打致死——或者说，因他们的击打而进入了生命。其余的行径，您可以想见与此一致：\OriginalTerm{司铎}{presbyters}和教师遭受迫害，以及那些滥用\OriginalTerm{皇权}{imperial authority}为所欲为的人可能做出的一切。但因您的祈祷，\Person{主}将赐我们摆脱这些事，并赐我们忍耐，以堪当承受试炼，相称于祂内的\OriginalTerm{望德}{hope}。请常常写信告诉我您的一切。如果您找到可靠的人可以将我写完的那本书带给您，恳请派人来取，这样，当我因您的认可而欢欣时，也可将它传给他人。愿藉\OriginalTerm{圣者}{the Holy One}的恩宠，您得以安然无恙，在主内欢欣，并为我祈祷，归于我和\Person{主}的教会。\par

\SourceRecord{Translated by Blomfield Jackson.；Nicene and Post-Nicene Fathers, Second Series；Vol. 8.；Edited by Philip Schaff and Henry Wace.；Buffalo, NY: Christian Literature Publishing Co.,；1895.；<http://www.newadvent.org/fathers/3202248.htm>.}

% structure: p0001 source=h1 target=section reason=page-title

\section{书信二四九}

\Emph{\Person{圣巴西略·凯撒勒雅}}\par

无收信人。推荐信。\par

我祝贺这位兄弟，因他得以脱离我们这里的困扰，并趋近您的尊威。他选择与敬畏主的人一同度善生，就是为来世作了美好的预备。我将此兄弟推荐于您的卓越，并藉他恳求您为我卑微的生命祈祷，使我得以脱离这些试探，开始按福音事奉主。\par

\SourceRecord{Translated by Blomfield Jackson.；Nicene and Post-Nicene Fathers, Second Series；Vol. 8.；Edited by Philip Schaff and Henry Wace.；Buffalo, NY: Christian Literature Publishing Co.,；1895.；<http://www.newadvent.org/fathers/3202249.htm>.}

% structure: p0001 source=h1 target=section reason=page-title

\section{书信二五〇}

\Emph{\Person{凯撒勒雅的圣巴西略}}\par

致\Person{埃加的帕特罗斐禄斯}主教。\par

我收到您对我上一封信的回复有些延迟；但经由文雅的\Person{斯特拉特吉乌斯}之手递到我手中，我已因您对我持续的情谊而感谢了主。您现在乐意就同一主题所写的，证明了您的善意，因为您所思所想合乎情理，且劝勉我获益。但若我必须回复您阁下所写的一切，我看我的话将过于冗长。因此，我只说这一点：如果和平的恩赐仅止于和平的名号，那么逐个挑出一些人、独让他们分享此福祉，而将无数其他的人排斥在外，便是荒唐可笑的。但如果与恶人同谋，表面上是在和平之中，实则对一切同意者造成如同敌人般的伤害，那么请想想，那些被接纳与之爲伴、对我怀有不义仇恨的人是谁？难道不是那些与我未共融的派系之人吗？现在无需指名道姓。他们被邀请到\Place{塞巴斯特亚}，承担了教会的管理职责，在祭台上服务，将他们自己的饼分给全体百姓，被那里的神职人员宣布为主教，并以圣徒和共融者的身份被护送经过整个地区。若必须采纳这些人的派系，那么从末梢开始、而不与他们的首领交往便是荒谬的。因此，如果我们要计算异端者而不回避任何人，那么告诉我，您为何要与某些人断绝共融？但若有人应被回避，就让那些在一切事上如此合乎逻辑的人告诉我，他们从\Place{加拉提亚}邀请来与他们同伙的人属于哪一派？若这类事情在您看来是可悲的，就请将分裂归咎于那些负责的人。若您认为他们无关紧要，那么就请原谅我拒绝成为谬误教义导师之酵的同类。因此，如果您愿意，不要再理会那些似是而非的论据，而要完全公开地驳斥那些不按福音真理正直行走的人。\par

\SourceRecord{Translated by Blomfield Jackson.；Nicene and Post-Nicene Fathers, Second Series；Vol. 8.；Edited by Philip Schaff and Henry Wace.；Buffalo, NY: Christian Literature Publishing Co.,；1895.；<http://www.newadvent.org/fathers/3202250.htm>.}

% structure: p0001 source=h1 target=section reason=page-title

\section{书信 251}

\Emph{\Person{凯撒勒雅的圣巴西略}}\par

\Emph{致\Place{厄瓦厄撒}的民众。}\par

1. 我的工作非常繁多，我的头脑充满许多焦虑的挂虑，但我从未忘记你们，我亲爱的朋友们，我一直为你们祈求我的天主，使你们在信仰上坚定不移，你们站在其中，并在天主光荣的希望中夸耀。诚然，现今很难找到并罕见看到一座纯洁、不受时代困扰所伤害、并完整无缺地保存着宗徒的教义的教会。你们的教会如今由那在每个世代中使配得祂呼召的人彰显出来的那一位所显明。愿主赏赐你们天上耶路撒冷的祝福，以回报你们将那些诽谤我的人对他们的诬蔑扔回头上，并拒绝让他们进入你们的心中。我知道，并在主内深信，\BibleQuote{你们的赏报在天上是丰厚的}\BibleRef{玛 5:12}，甚至仅因这种行为。因为你们明智地在你们中间得出结论，正如事实所是的，那些\Quote{以恶报善，以恨报爱我的人}如今指责我的正是他们自己所承认并签押的那些事。\par

2. 他们向你们呈交他们的签名作为对我的控告，这并不是他们陷入的唯一矛盾。他们被在君士坦丁堡聚集的主教们一致罢免。他们拒绝接受这一罢免，并上诉到一个不虔敬之人的会议，拒绝承认他们的审判者的主教职权，以免接受加于他们的判决。\par

他们不承认的理由是他们是不虔敬异端的首领。这一切发生在大约十七年前。罢免他们的主要人物是\Person{厄乌铎克西乌斯}、\Person{厄乌伊丕乌斯}、\Person{乔治}、\Person{阿卡基乌斯}，以及你们不认识的其他人。\par

现今教会的暴君们是他们的继承者，有些被祝圣来填补他们的位置，另一些实际上由他们提升。\par

3. 现在让那些指责我教义不纯的人告诉我，那些他们拒绝接受其罢免的人是在什么方式上是异端的。让他们告诉我，那些由他们提升并持有与其父辈相同观点的人是如何正统的。如果\Person{厄乌伊丕乌斯}是正统的，那么被他罢免的\Person{厄乌斯塔提乌斯}如何能不是平信徒？如果\Person{厄乌伊丕乌斯}是异端，那么由他祝圣的人现在如何能与\Person{厄乌斯塔提乌斯}共融？但这一切行为，这种试图控告人又把他们重新扶起的做法，是儿戏，是针对天主的教会为他们自己的利益而搞出来的。\par

当\Person{厄乌斯塔提乌斯}经过\Place{帕夫拉戈尼亚}时，他推翻了\Person{帕夫拉戈尼亚的巴西利德}的祭台，并常在自己的桌子上举行礼仪。现在他却乞求\Person{巴西利德}接纳他共融。他拒绝与我们可敬的弟兄\Person{厄尔丕迪乌斯}共融，因为他与\Place{阿马色内人}结盟；而现在他却作为恳求者来到\Place{阿马色内人}那里，请求与他们结盟。连你们自己也知道他对\Person{厄乌伊丕乌斯}的公开言论是多么令人震惊：现在他称赞那些持有\Person{厄乌伊丕乌斯}观点的人为正统，只要他们合作促进他的复职。而我却一直受到诽谤，不是因为我做错了什么，而是因为他们想象这样会得到\Place{安提约基雅}派系的推荐。去年他们从\Place{加拉达}请来的那些人，指望通过他们的手段恢复他们主教职权的自由运用，其品格凡是与他们同住过短时间的人都太清楚了。我祈求主永远不会允许我有闲暇来叙述他们所有的行径。我只说，他们走遍全国，有主教的尊荣和随从，由他们最可敬的护卫和同情者护送；并隆重入城，以全权召开集会。人民被交给了他们。祭台被交给了他们。他们如何去了\Place{尼科颇里}，在那里未能兑现他们所承诺的一切，以及他们如何回来，归来时显出的样子，是现场的人都知道的。他们显然每一步都是为了自己的利益和好处。如果他们说他们悔改了，让他们以书面形式表明悔改；让他们诅咒\WorkTitle{君士坦丁堡信经}；让他们与异端者分离；让他们不再欺骗头脑简单的人。关于他们和他们的，就说这么多。\par

4. 然而，我，亲爱的弟兄们，虽然渺小卑微，但赖天主的恩宠始终如一，从未随着世间的变化而改变。我的信经在\Place{塞琉西亚}、在\Place{君士坦丁堡}、在\Place{泽拉}、在\Place{兰普萨科斯}和在\Place{罗马}，没有变化。我现在的信经与从前并无不同；它始终如一。我们如何从主领受，就如何受洗；我们如何受洗，就如何宣认信仰；我们如何宣认信仰，就如何献上我们的\OriginalTerm{光荣颂}{doxology}，不将圣神与圣父和圣子分离，也不在尊荣上让圣神高于圣父，或断言祂先于圣子，如同亵渎者的舌头所捏造的。谁能如此鲁莽，拒绝主的命令，大胆地为这些名号发明自己的顺序？但我不会将那与圣父和圣子同列的神称为受造物。我不敢称那为王者的为奴仆。我恳请你们记住主在以下话语中发出的警告：\BibleQuote{凡人一切的罪和亵渎都可得赦免；惟独亵渎圣神，总不得赦免，今世来世也不得赦免}\BibleRef{玛 12:31}。你们要防备反对圣神的危险教导。\BibleQuote{你们要站稳在信仰上}\BibleRef{格前 16:13}。观看全世界，看看有多少部分是不纯正的。教会其余部分从世界的一端到另一端，都领受了福音，坚守这纯正无误的教义。我祈求我永不离开他们的共融，并祈求在你们中间，在我们主耶稣基督公义的日子，当祂来照各人的行为施行报应时，我得以有份和同得基业。\par

\SourceRecord{Translated by Blomfield Jackson.；Nicene and Post-Nicene Fathers, Second Series；Vol. 8.；Edited by Philip Schaff and Henry Wace.；Buffalo, NY: Christian Literature Publishing Co.,；1895.；<http://www.newadvent.org/fathers/3202251.htm>.}

% structure: p0001 source=h1 target=section reason=page-title

\section{书信252}

\Emph{\Place{凯撒勒雅}的\Person{圣巴西略}}\par

\Emph{致\Place{本都教区}的主教们。}\par

凡仰望主的人，都应热切渴求殉道者的荣耀，尤其是你们这些追求德行的人。你们对待同仆中伟大善良者的态度，正显明你们对我们共同的主的爱。此外，还有一个特别的原因，就是那如同血统的联系，将严谨纪律的生活与那些因忍耐而达于成全的人结合在一起。既然\Person{欧普绪基}、\Person{达玛斯}及其同伴是殉道者中最杰出的，他们的纪念每年在我们的城中和所有邻近地区举行，教会藉我的声音召叫你们，提醒你们保持前来朝拜的古老习惯。在你们面前，有一项伟大而美好的工作，就是在百姓中间，他们渴望得到你们的建树，并切望因尊敬殉道者而获得的奖赏。因此，请接受我的恳求，并仁慈地应允，以你们小小的辛劳，赐予我莫大的恩惠。\par

\SourceRecord{Translated by Blomfield Jackson.；Nicene and Post-Nicene Fathers, Second Series；Vol. 8.；Edited by Philip Schaff and Henry Wace.；Buffalo, NY: Christian Literature Publishing Co.,；1895.；<http://www.newadvent.org/fathers/3202252.htm>.}

% structure: p0001 source=h1 target=section reason=page-title

\section{\WorkTitle{书信第253号}}

\Emph{圣巴西略·凯撒勒雅}\par

\Emph{致安提约基雅的长老们。}\par

你们对天主的教会所怀的焦虑关切，将因我们极为敬爱的弟兄、司铎\Person{桑克提西穆斯}告诉你们整个西方对我们所怀的慈爱和善意，而得到一定程度的舒缓。但另一方面，当他亲自告诉你们当前的局势要求何等热忱时，这份关切又将被重新激起，变得更加敏锐。其他权威人士在向我们讲述西方人的心意和那里的状况时，几乎都只说了一半。他极善于理解人的心意，并精确地探询事态的真相；他将告诉你们一切，并引导你们的善意贯穿整个事务。摆在你们面前的材料，正适合你们在为天主的教会焦虑中所一贯表现出的卓越善意。\par

\SourceRecord{Translated by Blomfield Jackson.；Nicene and Post-Nicene Fathers, Second Series；Vol. 8.；Edited by Philip Schaff and Henry Wace.；Buffalo, NY: Christian Literature Publishing Co.,；1895.；<http://www.newadvent.org/fathers/3202253.htm>.}

% structure: p0001 source=h1 target=section reason=page-title

\section{书信254}

\Emph{\Place{凯撒勒雅}的\Person{圣巴西略}}\par

\Emph{致\Person{佩拉纠}}，\Emph{叙利亚\Place{劳狄刻雅}的主教。}\par

愿主再赏我亲自得见你的真正虔诚，并在实际交往中补充我信中所欠缺的一切。我迟迟动笔，必须诸多推诿。但我们有亲爱的可敬的弟兄\Person{桑克提西默斯}司铎在。他将告诉你一切，包括我们的消息和西方的消息。你听到后会感到振奋；但当他告诉你我们所陷入的困境时，或许会增添一些痛苦和忧虑，困扰你仁慈的灵魂。然而，你受此考验并非没有意义，因为你能感动主。你的忧虑将转化为我们的益处，我知道，只要有你祈祷的援助，我们必从天主获得扶助。也请与我一同祈祷，使我摆脱焦虑，并增加我体力的些许力量；然后主必使我旅途顺利，实现愿望，并见到阁下。\par

\SourceRecord{Translated by Blomfield Jackson.；Nicene and Post-Nicene Fathers, Second Series；Vol. 8.；Edited by Philip Schaff and Henry Wace.；Buffalo, NY: Christian Literature Publishing Co.,；1895.；<http://www.newadvent.org/fathers/3202254.htm>.}

% structure: p0001 source=h1 target=section reason=page-title

\section{第二五五封书信}

\Emph{该撒利亚的圣巴西略}\par

\Emph{致\Place{卡雷}主教\Person{维图斯}。}\par

但愿我能每天都给您写信！因为自从我体验到您的情谊以来，我就非常渴望与您交谈，若不能如此，至少也借书信与您往来，好向您告知我的近况，并了解您的状况。然而，我们所得的不是我们所愿的，而是主所赐予的，我们应当怀着感恩之心领受。因此，当我们的至爱且可敬的弟兄、司铎\Person{桑克蒂西穆斯}到达时，我感谢了圣洁的天主，因祂赐我机会给您写信。他在旅途中经历了相当多的辛劳，将会准确地向您讲述他在\Place{西方}所获悉的一切。为这一切，我们应当感谢主，并恳求祂也赐予我们同样的平安，使我们能自由地彼此接待。请以我的名义接待所有在主内的弟兄。\par

\SourceRecord{Translated by Blomfield Jackson.；Nicene and Post-Nicene Fathers, Second Series；Vol. 8.；Edited by Philip Schaff and Henry Wace.；Buffalo, NY: Christian Literature Publishing Co.,；1895.；<http://www.newadvent.org/fathers/3202255.htm>.}

% structure: p0001 source=h1 target=section reason=page-title

\section{第256封书信}

\Emph{凯撒勒雅的圣巴西略}\par

\Emph{致至爱可敬的弟兄：司铎阿卡基、埃提乌斯、保禄、西尔瓦努斯；执事西尔维努斯、路济乌斯，以及其余隐修士弟兄——巴西略主教。}\par

我已得知你们遭受了严厉迫害，复活节后立即有那些\BibleQuote{为争论和斗殴而禁食}\BibleRef{依 58:4}的人袭击你们的住所，将你们的劳作付之一炬。他们确实为你们预备了天上非人手所造的房屋\BibleRef{格后 5:1}，却为自己积蓄了那曾用来伤害你们的火。我一听到这事，便为所发生的事而叹息；并非怜悯你们，我的弟兄（愿天主宽恕！），而是怜悯那些如此深陷邪恶、竟将恶行推行至此的人。我曾期望你们全都立刻逃向我卑微自身为你们预备的避难所；并盼望主在我不断的烦恼中，借拥抱你们、以我这怠惰的身体接纳你们为真理而洒下的崇高汗水，使我得到慰藉，从而在真理审判者为你们所预备的奖赏中有一份。但你们并无此意，甚至未曾期望从我这里得到任何慰藉。因此，我至少渴望时常找机会写信给你们，为要像那些在竞技场上为斗士鼓劲的人一样，我自己能借信函在你们美好的战斗中给予鼓励。然而，有两个原因使此事不易。首先，我不知道你们住在何处。其次，我们的人中很少前往你们的方向。如今，主带来了至爱可敬的弟兄桑克提西穆斯司铎。通过他，我能向你们致敬，并恳请你们为我祈祷，欢喜踊跃，因为你们在天上的赏报是大的\BibleRef{玛 5:12}，而且你们在主面前有自由，昼夜不停地呼求祂止息教会的这场风暴，将牧者归还给羊群，使教会恢复其应有的尊严。我深信，若有人声感动我们良善的天主，祂必不远远施行仁慈，而会\BibleQuote{在受试探的时候，给你们开一条出路，叫你们能忍受得住}\BibleRef{格前 10:13}。请代我问候所有在基督内的弟兄，不拘何名。\par

\SourceRecord{Translated by Blomfield Jackson.；Nicene and Post-Nicene Fathers, Second Series；Vol. 8.；Edited by Philip Schaff and Henry Wace.；Buffalo, NY: Christian Literature Publishing Co.,；1895.；<http://www.newadvent.org/fathers/3202256.htm>.}

% structure: p0001 source=h1 target=section reason=page-title

\section{第二五七书}

\Emph{\Person{凯撒利亚的圣巴西略}}\par

致受亚略派扰乱的隐修士们。\par

我认为，当我听闻天主之敌加于你们的考验时，我曾对自己说：在一个被视为和平之时的时期，你们为自己赢得了为基督之名的缘故而受迫害的人所应许的福分，这是正当的，理应用信向你们宣告。依我判断，同国之人对我们发动的战争最难忍受，因为对付公开宣示的敌人我们容易自卫，而面对与我们联合的人，我们必然受其摆布，从而处于持续的危险之中。你们的情形正是如此。我们的祖先曾受迫害，但那是拜偶像者所为，他们掠夺财产，拆毁房屋，将他们放逐——这是我们的公开敌人为了基督之名做的。近来出现的迫害者，他们对我们的仇恨丝毫不逊于那些人，但为了欺骗许多人，他们标榜基督之名，使得受迫害者因宣认基督之名而失去一切安慰，因为大多数淳朴之人虽承认我们受冤屈，却不愿将我们为真理而死视为殉道。我因此确信，公义的审判者为你们存留的赏报，比赐给先前那些殉道者的更大。他们既得了人的公开赞美，又得了天主的赏报；而你们，虽然善行不亚于他们，却未从人民获得任何荣誉。故将来世所应得的报偿理应大得多。\par

因此，我劝勉你们，不要在苦难中沮丧，而要因天主的爱而振奋，日日增进你们的热情，深知在你们身上应当保存那真正的虔诚遗民，就是主在降临时所要寻得的。即使主教被逐出他们的教会，也不要惊慌。如果背叛者从神职人员本身兴起，也不要让这动摇你们对天主的信赖。我们得救不是凭名字，而是凭心思、意向和对造物主的真诚之爱。请回想，当攻击我主时，大司祭、经师和长老们如何设谋，而民众中有多少人真正接纳了圣言。要记住，得救的不是群众，而是天主所拣选的人。因此，不要被随风漂荡如海浪的一大群人吓倒。即使只有一个得救，像索多玛中的罗特，他也应持守正确的判断，坚定不移地持守对基督的希望，因为主必不遗弃祂的圣徒。请代我问候在基督里的诸位弟兄。请为我可怜的灵魂恳切祈祷。\par

\SourceRecord{Translated by Blomfield Jackson.；Nicene and Post-Nicene Fathers, Second Series；Vol. 8.；Edited by Philip Schaff and Henry Wace.；Buffalo, NY: Christian Literature Publishing Co.,；1895.；<http://www.newadvent.org/fathers/3202257.htm>.}

% structure: p0001 source=h1 target=section reason=page-title

\section{\WorkTitle{书信 258}}

\Emph{\Person{凯撒勒雅的圣巴西略}}\par

\Emph{致主教\Person{厄丕法尼}}\par

长期以来，按照我们主的预言，由于罪恶的增多，多数人的爱德已经冷却，这是意料之中的事。如今经验证实了这一点。虽然这种情形在我们这里已经出现，但由您的圣德带来的信件似乎反驳了这一点。因为您记得像我这样不配和微不足道的人，这确实是爱的非凡证明；其次，您派遣了适合管理和平通信的弟兄来看望我。如今，当每个人都以怀疑的眼光看待他人时，没有什么景象比您所呈现的更罕见了。处处不见怜悯，不见同情，不见为困苦的弟兄流下的兄弟之泪。为真理所受的迫害，全教会及其子民的泪水，围绕我们的这一大堆麻烦，都不足以激励我们彼此关心对方的福祉。我们践踏跌倒的人，我们抓挠和撕扯伤口；我们这些被认为彼此一致的人，竟然发出异端者所发出的诅咒；在最重要的事情上一致的人，却在某一点上完全分裂。那么，我怎能不钦佩他呢？他在这样的情况下，对邻人的爱纯真无伪，虽然与我有如此遥远的海洋和陆地相隔，却竭尽全力照顾我的灵魂。\par

我特别钦佩的是，您甚至为橄榄山上隐修士们的争论感到忧伤，并表示希望找到某种方法使他们彼此和好。我还高兴地听说，您并非不知道某些人采取的不幸步骤，这些步骤在弟兄们中引起了骚动，您甚至对这些事深感兴趣。但我认为，您将如此重要之事的纠正委托给我，这不太符合您的智慧：因为我没有得到天主的恩宠，因为我生活在罪中；我没有口才，因为我愉快地退出了虚浮的研究；我对真理的教义还不够精通。因此，我已经写信给我在橄榄山上的可爱弟兄们，我们的帕拉迪乌斯和意大利人依诺增爵，回复他们给我的信，说我不可能对尼西亚信经作丝毫补充，除了将荣耀归于圣神，因为我们的教父们对此点只是粗略地处理，当时没有关于圣神的问题。至于建议对该信经所作的关于我们主降生成人的补充，我既没有检验，也没有接受，因为它们超出我的理解。我深知，一旦我们开始干涉信经的单纯性，我们将陷入无休止的讨论，矛盾会使我们不断地前进，而我们只会通过引入新词语来搅扰更简单之人的灵魂。\par

至于安提约基雅的教会（我指的是在相同教义上一致的教会），愿主准许我们有一天能看到它合一。它特别容易遭受敌人的攻击，敌人对它发怒，因为在那里基督徒的名字最初得以确立。在那里，异端与正统分裂，正统自身也分裂。然而，我的立场是这样的：可敬的梅利特主教是第一个为真理勇敢发言的人，并在君士坦丁时代进行了那场美好的战斗。因此，我的教会因他那勇敢而坚定的立场而对他怀有深厚的感情，并与他共融。因此，因着天主的恩宠，我一直与他共融直到如今；如果天主愿意，我将继续如此。此外，至福的教父阿塔纳削从亚历山大里亚来，非常希望梅利特和他之间建立共融；但由于谋士们的恶意，他们的联合被推迟到另一个时机。但愿不是如此！我从未接受任何后来被引入该教座的人的共融，不是因为我认为他们不配，而是因为我看不到谴责梅利特的理由。尽管如此，我听到了许多关于弟兄们的事，但没有留意，因为被指控者没有被带到指控者面前，正如经上所记载：\BibleQuote{如果不先听取他，查明他所做的事，我们的法律难道就定他的罪吗？}\BibleRef{若 7:51}因此，我目前不能给他们写信，可敬的弟兄，我也不应该被强迫这样做。对于您平和的性情，合适的不是一方面造成合一，另一方面造成分裂，而是将分离的肢体恢复到原始的合一。那么，首先要祈祷；其次，尽您所能劝告，使野心从他们心中被驱逐，使他们之间实现和解，以恢复教会的力量，毁灭敌人的愤怒。令我的灵魂深感安慰的是，除了您在神学中其他正确而准确的陈述之外，您承认有必要声明\OriginalTerm{位格}{hypostases}是三个。请您以此方式教导安提约基雅的弟兄们。事实上，我相信他们已受过这样的教导；因为我确信，若不是您仔细确认了他们在这一点上的立场，您绝不会接受与他们的共融。\par

4. 正如您在其他信函中善意指出的那样，这里的\OriginalTerm{马古赛人}{Magusæans}人数众多，散布在整个地区，因为很久以前就有移民从巴比伦被引入这些地方。他们的习俗独特，不与其他民族交往。与他们交谈几乎是不可能的，因为他们已成了魔鬼的猎物，遵行魔鬼的旨意。他们没有书籍，也没有教义导师。他们在毫无意义的制度中长大，虔诚由父亲传给儿子。除了普遍可见的特征外，他们反对宰杀动物，认为这是玷污，并让自己所需使用的动物由他人宰杀。他们狂热追求非法婚姻；他们认为火是神圣的，等等。迄今为止，没有谁告诉过我关于马古赛人出自亚巴郎的传说：他们称某个\OriginalTerm{匝尔努阿斯}{Zarnuas}为他们的种族始祖。关于他们，我没有更多可以向阁下禀报的了。\par

\SourceRecord{Translated by Blomfield Jackson.；Nicene and Post-Nicene Fathers, Second Series；Vol. 8.；Edited by Philip Schaff and Henry Wace.；Buffalo, NY: Christian Literature Publishing Co.,；1895.；<http://www.newadvent.org/fathers/3202258.htm>.}

% structure: p0001 source=h1 target=section reason=page-title

\section{第二五九封信}

\Emph{\Person{凯撒勒雅的圣巴西略}}\par

致隐修士\Person{帕拉迪乌斯}与\Person{依诺增爵}。\par

从你们对我的情感，你们应能推测我对你们的情感。我一直渴望成为和平的传报者，当我未能达成目标时，我便忧伤。岂能不然？我无法因此对任何人动怒，因为我知道，和平的祝福早已从我们当中被收回。如果分裂的责任在他人身上，愿主赐予那些制造纷争的人停止如此行。我甚至不能要求你们常来看我。因此你们没有理由为此辩解。我很清楚，那些选择了劳动生活、并常亲手提供生活所需的人，不能长期离家；但无论你们在哪里，要记得我，为我祈祷，使我心中没有不安的缘由，并能与自己及天主和平相处。\par

\SourceRecord{Translated by Blomfield Jackson.；Nicene and Post-Nicene Fathers, Second Series；Vol. 8.；Edited by Philip Schaff and Henry Wace.；Buffalo, NY: Christian Literature Publishing Co.,；1895.；<http://www.newadvent.org/fathers/3202259.htm>.}

% structure: p0001 source=h1 target=section reason=page-title

\section{书信 260}

\Emph{\Person{圣凯撒勒雅的巴西略}}\par

\Emph{致主教奥普提摩}。\par

1. 在任何情况下，我都乐意看到那些好青年，既因他们超越年龄的稳重性格，也因他们与阁下您的亲密关系，这使我对他们有所期待。而当我看到他们带着您的来信来到我面前时，我对他们的爱加倍了。但现在我读了信，看到了信中所有对教会的焦虑关怀，以及它证明您对阅读圣经的热忱，我感谢主。我祝福那些带来如此信的人，甚至在他们之前，祝福写信者本人。\par

2. 您曾要求解答那著名段落，它被各处不同地诠释：\Quote{凡杀死加音的，将因七罪而受惩罚。}您的问题表明您自己已仔细注意了保禄对弟茂德的嘱咐，因为您显然用心于阅读。您还唤醒了我这个老人，因年龄和身体虚弱，以及身边许多搅扰、压在我生命上的苦难而迟钝。您自己充满热忱，现在正唤醒我，如同冬眠在洞穴中的野兽，使我有些许清醒和活力。所问的段落可以简单地解释，也可以加以详尽说明。简单而任何人都能立即想到的解释是：加音应为他的罪受七倍的惩罚。\par

因为公义的审判者不是以同等回报的原则来界定报应，而是邪恶的肇始者必须加增地偿还债务，如果他通过惩罚得到改善，并因他的榜样使其他人更明智。因此，由于规定加音须为他的罪付七倍的惩罚，据说杀死他的人将解除由天主的审判所宣告的对他的判决。这是我们在第一次阅读这段经文时自然产生的含义。\par

3. 但拥有更大好奇心的读者，自然倾向于进一步探究这个问题。他们问，正义怎能满足七次？惩罚是什么？是为所犯的七种罪吗？还是只犯了一次罪，但为一罪而有七种惩罚？圣经持续将七定为赦罪的数目。有人问：\Quote{多少次我的兄弟得罪我，我该宽恕他？}（这是伯多禄对主说。）\Quote{直到七次吗？}然后主回答说：\Quote{我不对你说直到七次，而是直到七十个七次。}\BibleRef{玛 18:21}我们的主没有改变数目，而是乘以七，从而固定了宽恕的限度。七年后，希伯来人从奴隶制中被释放。\BibleRef{申 5:12}七个七年在古时构成著名的禧年，\BibleRef{肋 25:10}在那时土地休息，债务豁免，奴隶得释放，可以说，新生命重新开始，旧生命世世代代在某种意义上在七这个数字上完成。这些事是现世生命的预像，它绕行七天而流逝，其中较轻的罪受到惩罚，按照我们仁慈的主的爱护，以免我们被交付于永无穷尽的世代中的惩罚。因此，\Emph{七次}这个说法被引入，因为它与现世有关，对于喜爱现世的人，应当特别在使他们选择过邪恶生活的事物上受惩罚。如果你把\Emph{惩罚}理解为加音所犯的罪，你会发现那些罪有七种。或者你若把它们理解为审判者对他的判决，你也不会大错。试列举加音的罪行：第一罪是\Emph{嫉妒}亚伯尔被偏爱；第二是\Emph{诡诈}，借此他对兄弟说：\Quote{我们到田间去：}\BibleRef{创 4:8}第三是\Emph{谋杀}，进一步邪恶；第四是\Emph{杀兄弟}，更大的不义；第五是他犯了\Emph{第一件谋杀}，为人类立了坏榜样；第六是\Emph{错}，因他使父母悲伤；第七是\Emph{谎言}对天主；因为当被问：\Quote{你的兄弟亚伯尔在哪里？}他回答：\Quote{我不知道。}\BibleRef{创 4:9}因此，七罪在加音的毁灭中受到惩罚。因为主说：\Quote{地开了口，从你手中接受你兄弟的血，你是应受诅咒的，}又说：\Quote{你在地上摇摆不定，}加音说：\Quote{你今天将我由地面上驱逐，我就要逃避你的面，在地上摇摆不定，凡遇见我的，必要杀我。}主回答他说：\Quote{凡杀死加音的，将偿清七项惩罚。}加音以为他容易成为每个人的猎物，因为在地里没有他的安全（因为地因他受了诅咒），并且他被天主除去援助，天主因他杀人而对他发怒，因此从地或天上都没有帮助。所以他说：\Quote{凡遇见我的，必要杀我。}圣经证明他的错误，用这句话：\Quote{不然；}\Emph{即}你不会被杀。因为对于受惩罚的人，死亡是一种收益，因为它带来痛苦解脱。但你的生命要延长，使你的惩罚与你的罪行相称。既然\OriginalTerm{被惩罚者}{\GreekText{ἐκδικούμενον}}一词可以从两方面理解：既指受惩罚的罪，也指惩罚的方式，现在让我们检查罪犯是否受了七倍的折磨。\par

4. \Person{加音}的七项罪在已经说过的话中已经列举了。现在我要问，施加于他的惩罚是否也是七项，我陈述如下。主询问说：\BibleQuote{你的弟弟\Person{亚伯尔}在哪里？}并不是因为他想知道，而是为了给\Person{加音}一个悔改的机会，正如这些话本身所证明的，因为当\Person{加音}否认时，主立即谴责他说：\BibleQuote{你弟弟的血从地里向我喊冤。}因此，\BibleQuote{你的弟弟\Person{亚伯尔}在哪里？}这个询问不是为了天主的信息，而是为了让\Person{加音}有机会认识他的罪。因为若不是天主访问他，他可能会辩解说自己被独自留下，没有给他悔改的机会。如今医生出现了，好让病人能向他求助。然而，\Person{加音}不但没有隐藏他的创伤，反而在谋杀上又加上了谎言，制造了另一个创伤。\BibleQuote{我不知道。难道我是看守我弟弟的吗？}现在就从这个点开始计算惩罚。\BibleQuote{地因你而受诅咒}，第一个惩罚。\BibleQuote{你要耕田}，这是第二个惩罚。某种隐秘的必然性被强加给他，迫使他耕种土地，以致他永远不被允许在愿意时休息，反而要一直与他的仇敌——土地——一同受苦，这土地被他用弟弟的血玷污而成了可咒骂的。\BibleQuote{你要耕田}。可怕的惩罚：与仇恨自己的人同住，与敌人、不共戴天的仇敌为伴。\BibleQuote{你要耕种田地}，也就是说，你要在田间劳苦，永不休息，日日夜夜不得解脱于你的工作，被隐秘的必然性束缚，这比任何野蛮的主人都更严厉，并被不断地催逼去劳动。\BibleQuote{她不再为你出产力量}。虽然无休止的劳作有一些果实，但对于一个被迫永不松懈的人来说，劳作本身就是不小的折磨。但劳苦是无休止的，土地的劳作是无果的（因为\BibleQuote{她不再出产力量}），而这种劳作的无效是第三个惩罚。\BibleQuote{你在世上要呻吟和颤抖}。这里在这三个之外又加上了两个：不断的呻吟和身体的颤抖，四肢被剥夺了源于力量的稳定。\Person{加音}滥用了自己身体的力量，所以它的活力被摧毁了，它蹒跚摇晃，他很难把肉和饮料送到嘴里，因为在他不敬的行为之后，他邪恶的手不再被允许服侍他身体的需要。另一个惩罚是\Person{加音}所说的：\BibleQuote{你把我从地面上驱逐出去，我要躲开你的面。}这个从地面上被驱逐是什么意思？意思是剥夺从土地而来的利益。他不是被转移到另一个地方，而是成为地上一切美好事物的陌路人。\BibleQuote{我要躲开你的面}。对于心地善良的人来说，最重的惩罚是被天主疏远。\BibleQuote{凡是遇见我的人，都要杀我}。他从前面的话推出这一点。如果我被从地上驱逐，而且躲开你的面，那么我就剩下被每个人杀死。主说什么？决不这样。但他给他一个记号。这是第七个惩罚：惩罚不应被隐藏，而应通过一个明显的记号向所有人宣告，这是第一个行不义之事的人。对所有正确推理的人来说，最重的惩罚是羞耻。我们也在审判的情况下学到了这一点，当\BibleQuote{有些人要复活进入永生，有些人要蒙受羞耻和永远的轻蔑}。\BibleRef{达 12:2}\par

5. 你下一个问题是关于\Person{拉默客}对他妻子所说的话，属于类似的性质：\BibleQuote{我杀了一个人，使我受伤；杀了一个青年，使我受创}。\BibleRef{创 4:23}有些人认为\Person{加音}是被\Person{拉默客}杀死的，并且他活到了这一代，以便遭受更长时间的惩罚。但事实并非如此。\Person{拉默客}显然犯了两宗谋杀，从他自己的话可以看出：\BibleQuote{我杀了一个人和一个青年}，这个人使他受伤，那个青年使他受创。\BibleQuote{如果\Person{加音}受七倍罚，那么\Person{拉默客}要受七十七倍罚}。\BibleRef{创 4:23}如果我应当承受四百九十项惩罚，这是公义的，因为天主对\Person{加音}的判断是公义的，他的惩罚应当是七项。\Person{加音}没有从别人那里学到杀人，也从未看见过杀人犯受惩罚。但我，眼前有\Person{加音}呻吟和颤抖，以及天主愤怒的浩大，却没有被眼前的榜样教训变得更明智。因此我应当承受四百九十项惩罚。然而，也有人做出了以下的解释，这与教会的教义并不冲突：他们说，从\Person{加音}到洪水，过去了七代，惩罚降临到整个大地，因为罪恶到处蔓延。但\Person{拉默客}的罪所需要的医治不是洪水，而是那亲自除去世人之罪的祂。\BibleRef{若 1:29}数算从亚当到基督来临的世代，你会根据路加的家谱发现，主是在第七十七代出生的。\par

这样，我已经尽我所能考察了这个问题，尽管我略过了其中一些可探究的事项，以免我的观察超过书信的篇幅。但对于你的智慧，小小的种子就足够了。\BibleQuote{教导明智的人，他会更加明智}。\BibleRef{箴 9:9}\BibleQuote{如果一个有技巧的人听到智慧的话语，他会称赞它，并增加它}。\BibleRef{德 20:18}\par

6. 关于\Person{西默盎}对\Person{玛利亚}所说的话，并无晦涩或多种解释。\BibleQuote{西默盎祝福了他们，并对他的母亲玛利亚说：看，这孩子已被立定，为使以色列中许多人跌倒和复起，并成为反对的记号；至于你，一把利剑将刺透你的心灵——为叫许多人的思念显露出来。} \BibleRef{路 2:34} 这里我感到惊奇，你忽略前文无需解释的话，却询问\Quote{至于你，一把利剑将刺透你的心灵}这句话。对我来说，同一个孩子如何既是跌倒又是复起，以及那反对的记号是什么，这些问题似乎并不比利剑如何刺透\Person{玛利亚}的心更少困惑。\par

7. 我的看法是：主是为跌倒和复起，并非因一些人跌倒而另一些人复起，乃因在我们内，那更坏的跌倒，那更好的被建立。主的来临摧毁我们身体的情欲，并激发灵魂的固有品质。如\Person{保禄}所说：\BibleQuote{我软弱时，正是我刚强时。} \BibleRef{格后 12:10} 同一个人既是软弱又是刚强，但他软弱在肉身，刚强在圣神内。因此，主并非给一些人跌倒的机会，给另一些人复起的机会。那些跌倒的人，从他们曾站立的位置跌倒；但显然，不信的人从未站立，总是被他所追随的蛇拖着在地上爬行。那么他就没有可跌倒之处，因为他已因自己的不信而被摔下。因此首要的恩惠是：那站立在罪中的人应跌倒并死去，然后应在正义中生活并复起，这两样恩宠都因我们对基督的信德赐予我们。让那更坏的跌倒，好使那更好的有机会复起。若奸淫不跌倒，贞洁就不复起。除非我们的无理性被摧毁，我们的理性便无法达到完美。在此意义上，他是许多人的跌倒和复起。\par

8. \Emph{成为反对的记号}。在圣经中，我们通常理解\Quote{记号}是指十字架。据说\Person{梅瑟}把蛇放在\BibleQuote{竿子上} \BibleRef{户 21:8}，那就是放在十字架上。或者，记号是表明某种奇特而隐晦的事物，为简单人所见，却为聪明人所理解。关于主的降生成人，争论从未止息：有些人断言祂取了身体，有些人说祂的存有是无形的；有些人说祂有可受苦的身体，有些人说祂藉一种显现完成了身体的\OriginalTerm{经济}{oikonomia}；有些人说祂的身体是属地的，有些人说是属天的；有些人说祂在万世以前已存在，有些人说祂从\Person{玛利亚}取得开始。正是为此，祂成了反对的记号。\par

9. 所谓利剑，是指那试探并判断我们思念的圣言，它刺入甚至剖开灵魂与精神、关节与骨髓，且能辨明我们的思念。如今，在受难之时，每个灵魂都仿佛经历一种搜查。按照主的话，经文说：\BibleQuote{你们都要为我的缘故而跌倒。} \BibleRef{玛 26:3} 因此，\Person{西默盎}预言\Person{玛利亚}自己，当她站在十字架旁，目睹所发生的事，并听见那些声音时——在\Person{加俾额尔}的见证之后，在她对天主受孕的隐秘知识之后，在伟大的奇迹展示之后——她的灵魂中将感到一股巨大的风暴。主必须为每人尝到死亡，成为世界的赎罪祭，并藉自己的血使所有人成义。连你，那从高处受教关于主的事的人，也将被某种疑惑触及。这就是那利剑。\Quote{为叫许多人的思念显露出来。} 祂指明，在基督十字架的绊跌之后，必有一种迅速的医治从主而来，临到门徒和\Person{玛利亚}自己，坚固他们的心对祂的信德。同样，我们看到\Person{伯多禄}，在绊跌之后，对基督的信德反而更加坚定。他人性中的软弱被证明，好显示主的大能。\par

\SourceRecord{Translated by Blomfield Jackson.；Nicene and Post-Nicene Fathers, Second Series；Vol. 8.；Edited by Philip Schaff and Henry Wace.；Buffalo, NY: Christian Literature Publishing Co.,；1895.；<http://www.newadvent.org/fathers/3202260.htm>.}

% structure: p0001 source=h1 target=section reason=page-title

\section{\WorkTitle{书信第二六一}}

\Person{凯撒勒雅的圣巴西略}\par

致\Person{索佐波利坦人}。\par

我收到了你们这些尊贵的弟兄关于你们处境的信。我感谢主，你们让我分享了你们对必要和值得认真关注之事的焦虑。但我忧伤地听说，除了亚略派给教会带来的骚动以及他们在信德定义上造成的混乱之外，你们中间又出现了一种创新，使弟兄们极为沮丧，因为据你们告知，有些人在信徒面前宣讲新奇陌生的教义，声称是从圣经教导中推导出来的。你们写道，你们中间有人试图摧毁我们主耶稣基督的救恩降生成人，并尽其所能推翻那从永远隐藏、但在祂自己时期显现的伟大奥秘的恩宠，当时主遍历了所有有关医治人类的事，将祂自己与我们同在的恩赐赐给了我们。祂首先通过圣祖们帮助祂的受造物，他们的生活被立为榜样和规则，给所有愿意跟随圣人脚步并以同样热忱达到善工成全的人。其次，祂赐下法律作为救助，由天使经\Person{梅瑟}之手颁布；然后有先知，预言将要来临的救恩；有民长、君王和义人，以大能的手行大事。在这些之后，在末世，祂自己在肉身内显现，\Quote{生于女人，生于法律之下，为救赎在法律之下的人，使我们获得义子的地位。}\BibleRef{迦 4:4}\par

2. 如果主的肉身降临从未发生，救赎者就没有替我们向死亡支付赎金，也没有通过自己摧毁死亡的统治。因为如果死亡所统治的不是主所取的那个，死亡就不会停止运作它自己的目的，天主之肉的苦难也不会成为我们的益处；祂不会在肉身内杀死罪恶；我们这些在亚当内死了的，就不会在基督内复活；破碎的不会被重新塑造；毁坏的不会被重新建立；因蛇的诡计而与天主疏远的，不会重新成为祂的。所有这一切恩惠都被那些断言主以属天之体降临的人所否定。如果天主之肉不是注定要从亚当的团块中取用，那圣童贞又有何必要？但谁有胆量现在再次借助诡辩的论证以及当然还有圣经证据，来复兴那早已沉寂的\Person{瓦伦提诺}的古老教义？因为这虚假的显现教义并非新奇。它很久以前由心智薄弱的\Person{瓦伦提诺}发起，他在撕扯了几句宗徒的论述后，为自己构造了这个邪恶的捏造，断言主取了\Quote{仆人的形状}\BibleRef{斐 2:7}，而没有取仆人本身，并且祂取了\Quote{相似}，但实际的人性并未被祂取。类似的情感由这些可怜的人表达，他们给你们带来新的麻烦。\par

3. 至于说人类的情感受到达了实际的天主性本身，这是那些在思想中不保持秩序、且不知道肉身的情感、拥有灵魂的肉身的情感、以及使用肉身的灵魂的情感之间存在区别的人所说的话。分裂、减少、溶解是肉身的特点；感到疲劳、痛苦、饥饿、口渴、被睡眠压倒，是拥有灵魂的肉身的特点；感到忧伤、沉重、焦虑等，是使用肉身的灵魂的特点。其中有些是自然的，对每个生物都是必要的；另一些来自邪恶的意志，是因生活缺乏适当的纪律和德行的训练而附加的。因此，显然我们主取了自然的情感以确定祂的真实降生成人，而不是以幻象或魔法的方式，而所有来自邪恶、玷污我们生命纯洁的情感，祂都拒绝了，认为它们不配祂无瑕的天主性。正是因此，祂被称为\Quote{取了有罪肉身的相似}，而不是像这些人所认为的，是取了肉身的相似，而是有罪肉身的相似。因此，祂取了我们的肉身及其自然的痛苦，但\Quote{没有犯罪}\BibleRef{伯前 2:22}。正如通过亚当传给我们的肉身中的死亡被天主性吞没，同样，罪被耶稣基督内的义除去，以致在复活中我们重新得到的肉身既不再死亡，也不再犯罪。\par

弟兄们，这就是教会的奥迹；这就是圣传。我嘱咐每一个敬畏主、等候天主审判的人，不要被各种教义带走。如果有人教导不同的教义，拒绝接受信德的健康言语，弃绝圣神的神谕，并把自己的教导看得比福音的教训更有权威，要提防这样的人。愿主准许我们有一天相见，以便我在面对面时可以补充我论证中遗漏的一切！我有许多话要说，却只写了少许，因为我不愿超出信函的界限。同时，我不怀疑所有敬畏主的人，简短的提醒就足够了。\par

\SourceRecord{Translated by Blomfield Jackson.；Nicene and Post-Nicene Fathers, Second Series；Vol. 8.；Edited by Philip Schaff and Henry Wace.；Buffalo, NY: Christian Literature Publishing Co.,；1895.；<http://www.newadvent.org/fathers/3202261.htm>.}

% structure: p0001 source=h1 target=section reason=page-title

\section{书信 262}

\Signature{凯撒勒雅的圣巴西略}\par

致隐修士\Person{乌尔比丘}\par

1. 你写信给我，做得很好。你已显示出爱德的果实何其伟大。请继续这样做。不要以为，当你写信给我时，你需要找借口。我认识自己的地位，并且我知道按本性，每个人与其余的人都享有同等的尊荣。我身上任何卓越之处，并非来自家族，也非来自过多的财富，亦非来自身体状况；它只来自在天主敬畏上的优长。那么，有什么阻止你更加敬畏主，并在这方面比我更大呢？常写信给我，告知我你那里弟兄们的状况。告诉我你们那里有哪些教会成员是健全的，使我知道我该写信给谁，可以信任谁。我听说有些人正试图以乖僻的意见败坏主降生成人的正确教义，因此我通过你们呼吁他们远离那些不合理的观点，这些观点据报有人持守。我的意思是，天主本身变成了肉身；祂并未通过圣玛利亚取用了亚当的本性，而是以祂自己的天主性，被改变成一种物质的本性。\par

2. 这个荒谬的立场很容易被驳倒。亵渎本身就是定罪，因此我认为，对于敬畏主的人来说，只需提醒一下就够了。因为如果祂被转变了，那么祂就改变了。但我不敢说或想这样的事，因为天主曾宣告：\BibleQuote{我是上主，决不改变}\BibleRef{拉 3:6}此外，降生成人的益处如何能传达给我们，除非我们的身体与天主性结合后，变得超越死亡的管辖？如果祂改变了，祂就不再构成一个真正的身体，如同在结合了神性的身体之后所存留的那样。但是，如果独生子的全部本性都被改变了，那不可理解的天主性怎么能被局限于一个渺小身体的体积之内呢？我确信，凡是头脑清醒且敬畏天主的人，都不会患有这种不健全。但我收到报告说，你们中有些人正受这种精神疾病之苦，因此我认为有必要，不是仅仅向你们致以礼节性的问候，而是在我的信中包括一些甚至能建立敬畏主之人的灵魂的内容。因此我敦促这些谬误应受到教会的纠正，并且你们应避免与异端者共融。我知道我们在基督内的自由正因对这些问题的漠不关心而被剥夺。\par

\SourceRecord{Translated by Blomfield Jackson.；Nicene and Post-Nicene Fathers, Second Series；Vol. 8.；Edited by Philip Schaff and Henry Wace.；Buffalo, NY: Christian Literature Publishing Co.,；1895.；<http://www.newadvent.org/fathers/3202262.htm>.}

% structure: p0001 source=h1 target=section reason=page-title

\section{\WorkTitle{书信二六三}}

\Person{圣巴西略}（\Place{凯撒勒雅}）\par

致西方教会。\par

1. 愿我们所信靠的主天主，赐予你们每一位足够的恩宠，使你们能实现你们的希望，如同你们透过那可敬爱的同僚司铎们送来的信函，以及你们对我患难中所感同身受的关怀所带给我的喜乐一样。你们已穿上了怜悯的心肠，\BibleRef{哥 3:12}，正如上述司铎们向我描述你们的那样。虽然我的创伤依旧，但身边有良医，若能得着时机，迅速为我的伤痛敷上药膏，这仍然使我得到缓解。因此，我藉着我们所爱的朋友们回礼问候你们，并劝告你们，若主赐予你们能力到我这里来，不要迟疑来访。因为看望病人是最大诫命的一部分。但若善良的天主、我们生命的明智管家将这恩惠保留到另一时节，那么至少请写信给我，写下你们当写的、足以安慰受压迫者和扶起被压碎者的言语。教会已遭受许多沉重打击，我为此极其痛苦。如今到处都无指望得到救助，除非主藉着你们这些祂真正的仆人将良药赐予我们。\par

2. 亚略异端的狂妄无耻，虽已被公开从教会身体中割除，却仍居于其错误之中，并不给我们造成太大伤害，因为其不虔敬已是众所周知。然而，那些披着羊皮、外表温和可亲、内里却毫不留情地蹂躏基督羊群的人，容易伤害较单纯的人，因为他们从我们中间出来。这些人正是难防又难对付的。我们恳请你们勤勉，向东方所有教会公开谴责他们；好使他们或转向正道，与我们真诚联合；或若执迷不悟，则只将祸害留给自己，不能因不谨慎的共融而将自身的瘟疫传给邻人。我不得不提名道姓，为使你们亲自认识那些在此挑起纷争的人，并将他们告知我们的教会。我自己的话遭大多数人猜疑，仿佛我因某些私怨对他们怀有恶意。然而，你们在民众中却更有信誉，因为你们的家乡离他们更远，况且天主赐予你们恩宠以帮助那些困苦者。若你们更多人在发表相同意见上一致，那么显然，发表这些意见的人数之众将使反对接受它们成为不可能。\par

3. 给我带来巨大悲伤的人之一，是小亚细亚\Place{塞巴斯特}的\Person{欧斯塔丢}；他昔日是\Person{亚略}的门徒，当亚略在\Place{亚历山大}得势并炮制出那些反对独生子的邪恶亵渎时，他是亚略最忠实的追随者之一。他返回本国后，向\Place{凯撒勒雅}极其有福的主教\Person{赫尔莫革内斯}提交了一份正确信仰的宣认，因为赫尔莫革内斯正要因他错误教义而谴责他。在如此情况下，他被赫尔莫革内斯祝圣为主教；那位主教去世后，他急忙投奔\Place{君士坦丁堡}的\Person{欧瑟比}，后者在积极支持亚略不虔敬教义方面不亚于任何人。他从君士坦丁堡因某种原因被驱逐，回到本国，第二次为自己辩护，试图隐藏其不虔敬的思想，并以某种言语上的正统掩饰它们。他刚一获得主教品秩，就在\Place{安卡拉}的亚略派主教会议上，立即公开写作谴责\OriginalTerm{同质}{Homoousion}。从那里他去了\Place{塞琉西亚}，参与了其异端同伙的臭名昭著的行径。在\Place{君士坦丁堡}，他第二次同意异端者的主张。因先前在\Place{梅利特内}被罢免，他失去了主教职位，于是想到前往你们那里旅行，作为恢复自身地位的手段。有福的主教\Person{利贝流}向他提出什么条件，他又同意了什么，我不得而知。我只知道他带回一封恢复职务的信函，在\Place{提雅纳}主教会议上出示，从而恢复了主教职位。如今他正诋毁那曾使他被接纳的信经；他与那些谴责同质的人为伍，并成为\OriginalTerm{神格反对者}{pneumatomachi}异端的首要领袖。既然他是在西方获得伤害教会的力量，并利用你们所赋予的权威来颠覆众人，那么他的纠正也应来自同一地方，且应发送一封信函给诸教会，说明他是在什么条件下被接纳，以及他如何改变了行为，使当时教父们赐予他的恩惠归于无效。\par

4. 接下来是\Person{阿波里纳利}，他给教会带来的悲伤毫不逊色。他写作便捷，口舌锋利，能就任何主题展开辩论，以致他的著作充满天下，无视那说\Quote{慎防著作繁多}的劝言。著作多了，错误定然不少。多言多语，岂能无罪？阿波里纳利的神学著作以圣经论证为基础，却以人的起源为根基。他关于复活的著作，是从神话的、甚至犹太人的观点出发；主张我们将回到法律崇拜，受割损，守安息日，戒避某些肉类，向天主献祭，在\Place{耶路撒冷}圣殿朝拜，完全从基督徒变成犹太人。还有什么比这更可笑的？或者更确切地说，还有什么比这更违背福音教义的呢？此外，他关于降生成人又在弟兄们中造成了如此混乱，以致他的读者中很少有人保留真正宗教的古旧印记；而大部分人在追求新奇的热情驱使下，转向对他那些无益论著的探究和争辩。\par

5. 至于保利努斯的言论是否有可指摘之处，你们自己可以判断。使我痛苦的是，他竟对玛尔塞鲁的学说表现出倾向，并毫无保留地接纳他的追随者进入共融。最可敬的弟兄们，你们知道，玛尔塞鲁的学说牵涉我们一切希望的颠覆，因为他不承认圣子以其本身的\OriginalTerm{位格}{\GreekText{ὑπόστασις}}存在，却将祂描述为被派遣出去，然后又回到派遣祂者那里；他也不承认\OriginalTerm{护慰者}{\GreekText{Παράκλητος}}拥有自己的自立体。因此，任何人若宣称这异端完全与基督教相悖，并称之为腐化的犹太主义，都不会有错。我恳请你们对此予以适当关注。如果你们愿意写信给东方各教会，告知那些曲解这些教义的人，若他们改过，则与你们共融；但若他们固执坚持自己的创新，则你们与他们分离，这样就能实现上述关注。我自己深知，本该与你们共同在会议中商议这些事，但时间不允许。拖延是危险的，因为他们造成的祸害已经扎根。因此，我不得不派遣这些弟兄，以便你们能从他们那里得知我信中所遗漏的一切，并且他们能激励你们向东方各教会提供我们所祈求的援助。\par

\SourceRecord{Translated by Blomfield Jackson.；Nicene and Post-Nicene Fathers, Second Series；Vol. 8.；Edited by Philip Schaff and Henry Wace.；Buffalo, NY: Christian Literature Publishing Co.,；1895.；<http://www.newadvent.org/fathers/3202263.htm>.}

% structure: p0001 source=h1 target=section reason=page-title

\section{Letter 264}

\Emph{凯撒勒雅的圣巴西略}\par

\Emph{致流亡中的厄德萨主教巴尔塞斯。}\par

致主教巴尔塞斯，真正为天主所爱、堪受一切敬重与尊荣者，巴西略在主内问安。\newline{}我亲爱的弟兄多姆尼努斯启程前往你处，我欣然抓住写信的机会，藉他向你致候，并祈求圣洁的天主，使我们能在此世存活足够长久，得以见到你，并享受你所拥有的美好恩赐。\newline{}但请你，我恳求你，祈祷主不要永远将我们交付给基督十字架的仇敌，而是愿祂保守祂的教会，直到公义的审判者亲自知道何时赐予的那一和平的时刻。因为祂必会赐予。祂不会永远抛弃我们。正如祂为惩罚以色列人的罪，将他们的被掳期限定为七十年（\BibleRef{耶 25:12}），同样，或许全能者会在某个指定的时期交出我们之后，再次召叫我们，并将我们恢复到起初的平安——除非那背道之事现已临近，而最近发生的事件正是假基督来临的开始。\newline{}若是如此，请祈祷善良的主要么除去我们的苦难，要么使我们在苦难中不败地保持自己。\newline{}藉着你，我问候所有堪当与你共处的人。凡与我在一起的人，都向你的尊荣致意。\newline{}愿藉圣者的恩宠，你身体健康，为天主的教会所保存，信赖于主，并为我祈祷。\par

\SourceRecord{Translated by Blomfield Jackson.；Nicene and Post-Nicene Fathers, Second Series；Vol. 8.；Edited by Philip Schaff and Henry Wace.；Buffalo, NY: Christian Literature Publishing Co.,；1895.；<http://www.newadvent.org/fathers/3202264.htm>.}

% structure: p0001 source=h1 target=section reason=page-title

\section{第二六五封书信}

\Emph{\Person{圣巴西略·凯撒勒雅}}\par

致在流徙中的\Place{埃及}主教们\Person{欧罗基乌斯}、\Person{亚历山大}、\Person{哈波克拉提翁}。\par

在一切事上，我们发现我们良善的\Person{天主}对其教会施行的眷顾是强有力的，因此，那些看似阴暗、不按我们所愿发展的事，凭着\Person{天主}隐藏的智慧，以及祂正义的不可测度的判断，也被安排为大多数人的益处。如今主已将你们从\Place{埃及}地区迁出，将你们安置在\Place{巴勒斯坦}中间，如同古时的\Place{以色列}，祂曾藉掳掠将他们带到\Place{亚述}之地，并藉祂的圣者的寄居在那里消灭了偶像崇拜。现在我们看到同样的事，当我们观察到主使你们为真宗教的奋斗为人所知，透过你们的流徙为你们开启了蒙福竞赛的竞技场，并为所有看见你们高贵坚毅的人，赐予你们善表的恩惠，引导他们获得救恩。藉\Person{天主}的恩宠，我听说你们信仰的正确，以及对弟兄们的热忱，并且你们并非以粗心大意或敷衍了事的态度，提供有益及必要的救恩之事，并支持一切有助于建立教会的事。因此我认为自己应当与你们的良善共融，并通过书信与你们的尊威结合。为此，我派遣了我至亲爱的弟兄执事\Person{厄尔彼狄乌斯}，他不仅带去我的信，而且完全有资格向你们宣布我信中所可能省略的一切。\par

2. 我尤其因听闻各位尊长在正统信仰事业上的热忱报告，而渴望与你们合一。你们心灵的坚毅既没有被诸多书籍的纷繁，也没有被各种巧妙论点的多样所动摇。相反地，你们认出了那些试图引入与宗徒教义相对立的革新的人，并且拒绝对他们正在造成的危害保持沉默。我确实在所有持守主平安的人中，因\Place{劳狄刻雅}的\Person{阿波利纳里}的种种革新而大为苦恼。他尤其使我苦恼的是，起初他似乎站在我们一边。受苦者在某种意义上可以忍受来自公开敌人的一切，即使极其痛苦，正如经上所载：\BibleQuote{因为不是仇敌辱骂我，那时我还能忍受。}但是，被一位亲密又同情的朋友冤屈，则是无法忍受、超出安慰能力的。如今，那个我曾期望站在我右边捍卫真理的人，我却发现他以多种方式阻碍那些正在得救的人，诱惑他们的心思，将他们从正直的教义中拉走。他有什么鲁莽仓促的行为没有做？有什么欠考虑的危险论点他没有冒险？自从他派遣人到由正统主教管理的教会，撕裂他们并设立某种奇特且非法的礼仪以来，整个教会岂非自相分裂？当主教们没有人民和圣职人员，只有空名和头衔，对促进和平与救恩的福音毫无作为时，真宗教的伟大奥迹岂非受到嘲笑？他关于\Person{天主}的言论岂非充满不敬的教义，那疯狂的\Person{撒伯里乌}的古老不敬如今由他笔端重现？因为如果在\Person{塞巴斯特人}中流传的作品不是敌人的伪造，而真是他的著作，那么他达到了无法超越的不敬高度，声称圣父、圣子、圣神是同一的，还有其他我不愿听闻的黑暗不虔之词，我祈祷自己与说出这些话的人毫无关系。他岂非混淆了降生成人的教义？救恩的经纶岂非因他关于此事的黑暗含混的推测而令许多人变得可疑？要收集所有这些并加以驳斥，需要长时间和许多讨论。但是福音的应许在哪里像被他的虚构所钝化和摧毁？他竟然如此卑劣贫乏地解释为所有按基督福音生活之人所预备的蒙福希望，将其降低为纯粹的老妇寓言和犹太人的教义。他宣告圣殿的更新、法律敬礼的遵守、在真正大司祭之后又一个预像性的大司祭，以及在除免世罪的天主羔羊之后的赎罪祭。\BibleRef{若 1:29}他在一次洗礼之后传讲部分洗礼，以及用牛犊灰洒在教会——这教会藉着对基督的信仰，没有斑点、皱纹或任何这类东西；\BibleRef{弗 5:27}在复活的无痛苦状态之后癞病的洁净；当他们不娶不嫁时的疑忌祭；在从天降下的食粮之后的供饼；在真光之后点燃的灯。总之，如果诫命的规条被教义所废除，那么显然在这些情况下，基督的教义也会被法律的规条所废止。面对这些事，羞耻和耻辱遮蔽了我的脸，沉重的痛苦充满我的心。因此，我恳求你们，作为熟练的医生，并受教如何以温和惩戒反对者，尽力使他回到教会的正当秩序，并劝他轻视自己著作的冗长；因为他证实了这句谚语的真理：\BibleQuote{多言难免有罪。}\BibleRef{箴 10:19}请大胆地将正统教义摆在他面前，以便他的改过能被宣扬，他的悔改能为弟兄们所知。\par

3. 我也应当提醒尊驾关于\Person{玛尔塞鲁斯}的追随者，使您不致对他们的案件轻率或欠考虑地作出决定。因他不敬神的教义，他已被逐出教会。因此，他的追随者只应在弃绝那异端的条件下才被接纳共融，好使那些通过您与我联合的人能被所有弟兄接纳。如今，许多人听说您竟在这些人前来投奔阁下时接纳了他们并允许他们进入教会共融，都感到非常悲伤。然而，您应当知道，靠着天主的恩宠，您在东方并不孤立，而是有许多与您共融的人，他们维护教父的正统，并宣扬在\Place{尼西亚}所确立的信仰的正统教义。西方教会的弟兄们也完全赞同您和我，我已接受并保存他们所阐述的信仰，同意他们纯正的教义。因此，您本应使所有赞同您的人都满意，使正在采取的行动得到普遍认同，而不致因接纳一些人、排斥另一些人而破坏和平。如此，您本当同时严肃而温和地商议对普世各教会都重要的事务。赞美不属于仓促作决定的人，而属于坚定而不可更改地掌管每项细节的人，这样，他的判断即使日后被审查，也能更受尊重。这样的人，作为谨慎措辞之士，为天主和人悦纳。这样，我已在信中尽我所能向尊驾进言。愿主恩赐我们有一日能相见，好使我们在共同为治理教会安排一切后，能与您一同领受那正义的审判者为忠信而明智的管家所预备的赏报。同时，请惠告您基于何种意向接纳了\Person{玛尔塞鲁斯}的追随者，须知即便您本人各方面都稳妥，也不应独自处理如此重要的事务。此外，还需要西方教会的人以及与他们在东方共融的人，共同同意接纳这些人。\par

\SourceRecord{Translated by Blomfield Jackson.；Nicene and Post-Nicene Fathers, Second Series；Vol. 8.；Edited by Philip Schaff and Henry Wace.；Buffalo, NY: Christian Literature Publishing Co.,；1895.；<http://www.newadvent.org/fathers/3202265.htm>.}

% structure: p0001 source=h1 target=section reason=page-title

\section{书信二六六}

\Emph{圣巴西略（凯撒勒雅）}\par

\Emph{致亚历山大主教佩特鲁斯。}\par

1. 你非常合宜地责备了我，且是以一位被主教导过真实爱德的属灵弟兄相称的方式，因为我未能将这里发生的一切精确而详细地告知你；因你理当关心我的事，而我也理当将关于我的一切都告诉你。但我必须告诉你，极可敬且亲爱的弟兄，我们不断的苦难，以及如今摇撼各教会的巨大动荡，使我将所发生的一切视为理所当然。正如在打铁铺里，耳朵被震聋的人习惯了那声音，同样，由于不断有奇异的消息传来，我已习惯于对非常事件保持平静而不惊恐。因此，亚略派长期损害教会所推行的策略，尽管其成就众多且巨大，并传遍全世界，对我而言仍是可忍受的，因为这是真理之言的公开仇敌所作的事。只有当这些人做出不寻常之事时，我才感到惊讶，而非当他们试图对真宗教做出大胆凶暴之举时。但使我悲伤和困扰的，是那些与我同感同思之人的所作所为。然而，他们的行为如此频繁且不断传到我耳中，以致连这些也不再显得惊人。于是，我对最近那些混乱的行为并不激动，部分原因是由于我深知，无需我开口，传闻便会将之带到你那里；部分原因是由于我宁愿等别人来告诉你这些不愉快的消息。再者，我认为自己对这些行为显出愤慨是不合理的，仿佛我因遭受轻慢而恼怒。我已以相称的言辞写信给此事的具体经办人，劝诫他们，因那里有些弟兄中起的纷争，不要背离爱德，而要等待那些有权按正当教会秩序纠正混乱的人来解决问题。你被尊贵合宜的动机所激励，如此行事，值得我称赞，并使我对主心怀感恩，因你身上仍保存着古代纪律的遗风，且教会在我所受的迫害中没有丧失她的能力。教规并未像我一样遭受迫害。虽再次被迦拉达人纠缠，我始终未能给他们答复，因我等候你的决定。如今，若主旨意如此，且他们肯听从我，我希望我能将人民领回教会。这样，就不能责备我投靠了玛尔西安派，相反，他们将变成基督教会身体的肢体。如此，由异端导致的耻辱将因我所采用的方法而消失，我也将避免投靠他们的责难。\par

2. 我也因我们的弟兄多洛泰而悲伤，因为据他自己所写，他并未以温和的方式将一切禀报给你的阁下。我将此归咎于时世艰难。若我最好的弟兄因未能按我意愿履行职责而被证明心怀恶意、不胜任，似乎我因自己的罪而被剥夺了一切事业的成功。多洛泰回来后，向我报告了他在极可敬的达玛苏斯主教面前与你的阁下交谈的内容；他提到我们蒙天主所爱的弟兄和同工梅利休斯与欧瑟比被列为亚略派狂人，这使我痛苦。若他们的正统性别无其他证明，那么亚略派对他们发起的攻击，在所有思想正确的人看来，正是他们正直的不小证据。即使你因基督之名与他们一同受苦，也理当使你的尊敬在爱德中与他们联合。请确信，极可敬的先生，没有一句正统的话语不是这些人以全部勇气宣讲过的。天主是我的见证，我亲耳听过他们。若我曾发现他们在信仰上失足，我定然不会如今仍接纳他们进入共融。但若你认为合适，让我们把过去放下。让我们为未来和平开篇。因为我们必须在肢体的共融中彼此需要，尤其此刻，当东方各教会正注视着我们，并将你们的合一看作力量与巩固的保证。反之，若他们觉察你们相互猜疑，便会垂手，放松对信仰仇敌的抵抗。\par

\SourceRecord{Translated by Blomfield Jackson.；Nicene and Post-Nicene Fathers, Second Series；Vol. 8.；Edited by Philip Schaff and Henry Wace.；Buffalo, NY: Christian Literature Publishing Co.,；1895.；<http://www.newadvent.org/fathers/3202266.htm>.}

% structure: p0001 source=h1 target=section reason=page-title

\section{书信 267}

\OriginalTerm{圣巴西略·凯撒勒雅}{ST. BASIL OF CAESAREA}\par

致在流放中的厄德萨主教巴尔塞斯。\par

我因对您怀有的情感，渴望与您相会，亲爱之友，亲爱地拥抱您，并光荣那在您身上受显扬的主，祂使您可敬的晚年得以在天下敬畏祂的人中闻名。然而，沉重的病痛折磨着我，我因对教会的挂虑而更加沉重，甚至无法用言语表达。我不能自主，随心所欲地前往或探访任何我想见的人。因此，我试图通过书信来满足我对您内里美善的渴望，并恳请尊座为我并为教会祈祷，求主恩赐我在余下的时日或寄居此世的时刻中无过犯生活。愿祂使我得见祂教会的平安。至于您的同工与同道，愿我得知我所祈求的一切，也得知您得蒙属地人民日夜向正义之主祈求的福分。我写信不多，甚至不及当写的次数，但我已给尊座写信。或许，我所托付弟兄们保存的问候未能得以保存。如今，我得知有些弟兄正前往尊座之处，便欣然将书信托付给他们，并附上一些嘱托，恳请尊座不弃我卑微，并请依圣祖依撒格的榜样为我祝福。\BibleRef{创 27:27}我事务繁多，心思多虑，或许有所遗漏而未曾言及。若有，求您不算为过，亦请不要因此忧伤。愿您行事处处合乎您的高尚品格，使我能与他人同享您美德的果实。愿您安康，喜悦于主，为我祈祷，与我及教会同在。\par

\SourceRecord{Translated by Blomfield Jackson.；Nicene and Post-Nicene Fathers, Second Series；Vol. 8.；Edited by Philip Schaff and Henry Wace.；Buffalo, NY: Christian Literature Publishing Co.,；1895.；<http://www.newadvent.org/fathers/3202267.htm>.}

% structure: p0001 source=h1 target=section reason=page-title

\section{\WorkTitle{第268封信}}

\Signature{凯撒勒雅的圣巴西略}\par

致流亡中的\Person{欧瑟伯}。\par

即使在我们这个时代，主也以其伟大而有力的手保护了您圣德的生命，以此教导我们：祂不会遗弃祂的圣者。我认为您的情况几乎就像那位圣人在深海巨兽腹中毫发无损，或像那些敬畏主的人在烈火中安然无恙一样。因为尽管战争四面环绕您，但据我所闻，祂仍然保全了您安然无恙。愿全能的天主保守您，如果我再活久一点，好让我能实现迫切祈求见到您的祷告！若不为我，愿祂为其余的人保守您，他们等待您的归来如同等待自己的救恩。我深信主以祂的仁慈，必会垂顾众教会的眼泪和众人为您发出的叹息，并保全您的生命，直到祂允准日夜向祂祈祷者的祈求。关于对您所采取的一切措施，直至我们可爱的弟兄\Person{利巴纽斯}执事到来，我在他途中已从他那里得到充分了解。我急切想知道后来的事情。我听说同时您那里发生了更大的困难；关于这一切，如果可能，尽早，若不能，至少在我们可敬的弟兄\Person{保禄}长老返回时，但愿我能得知，如我所祈求的，您的生命安然无恙。但因听说所有道路都充满了盗贼和逃兵，我不敢委托弟兄携带任何东西，以免导致他的死亡。如果主赐予一点平静（因我听说军队即将到来），我将设法派遣我自己的一个人去拜访您，以便带给我关于您一切的消息。\par

\SourceRecord{Translated by Blomfield Jackson.；Nicene and Post-Nicene Fathers, Second Series；Vol. 8.；Edited by Philip Schaff and Henry Wace.；Buffalo, NY: Christian Literature Publishing Co.,；1895.；<http://www.newadvent.org/fathers/3202268.htm>.}

% structure: p0001 source=h1 target=section reason=page-title

\section{书信第二百六十九号}

\Emph{\Person{凯撒勒雅的圣巴西略}}\par

致将军\Person{阿林泰乌斯}之妻。慰唁书。\par

1. 按理说，我本应在场，并亲历目前所发生的事，这样我既能减轻自己的悲痛，也能给阁下一些安慰。但我的身体已无法忍受长途跋涉，因此只好以书信的方式接近您，以免显得对发生的事全然漠不关心。谁不为那人哀悼？谁的心硬如石，不为他流下热泪？尤其是我，想起从他那里受到的一切尊敬，以及他对天主的教会所给予的普遍保护，心中充满了哀伤。然而，我想到他是人，在世完成了当行之事，如今按既定时间，被安排我们命运的天主收回。我恳求您，以您的智慧，将这些事放在心上，以温柔面对这一事件，并尽可能以节制承受您的损失。时间或许能安抚您的心，让理性得以接近。同时，您对丈夫的深爱以及对众人的善良，使我担心，因您性格的纯朴，悲伤的创伤可能刺入极深，使您完全沉溺于情感。圣经的教导总是有益的，尤其是在这样的时刻。请记住我们的造物主所宣判的判决。借此，我们所有由尘土而来的人都要归回尘土。\BibleRef{创 3:19} 没有人伟大到能超越腐朽。\par

2. 您令人钦佩的丈夫是一位善良而伟大的人，他的体力与灵魂的德行相匹敌。我必须承认，他在两方面都无人能及。但他是人，他死了，像亚当、像亚伯、像诺厄、像亚巴郎、像梅瑟，或任何您能想到的同性之人。那么，我们不要因为他被带走而抱怨。倒要感谢那使我们与他结合的那一位，感谢我们从起初就与他同住。失去丈夫是您与其他妇女共有的命运，但与这样一位丈夫结合，却是任何其他妇女都无法夸耀的。事实上，我们的造物主将那人塑造为我们的人类本性应当成为的模范。所有人的目光都被他吸引，所有人的口舌都在讲述他的事迹。画家和雕塑家都不及他的卓越，当历史学家讲述他在战争中的成就时，似乎陷入了神话和不可信的区域。因此，大多数人甚至无法相信那传达噩耗的报告，或接受阿林泰乌斯已死这一事实的真相。然而，阿林泰乌斯经历了将要发生在天、太阳和地上的事。他光荣地死去，没有因年老而弯腰，没有失去一丝荣誉；在世时伟大，在未来生命中伟大；在希望的光荣面前，他现世的辉煌没有任何减损，因为他在离开此世的时刻，在重生的洗濯中，洗净了灵魂的一切污秽。你安排并参与了这一礼仪，这是莫大的安慰。现在，将你的思想从现世转向未来，使你能够通过善工，获得与他一样的安息之所。怜悯年迈的母亲，怜悯柔弱的女儿，你现在是她唯一的安慰。作其他妇女的坚毅榜样，这样节制你的悲伤，既不要将其从心中驱除，也不要被痛苦压垮。始终着眼那极大的忍耐赏报，那是我们的主耶稣基督应许作为此世善行的回报。\par

\SourceRecord{Translated by Blomfield Jackson.；Nicene and Post-Nicene Fathers, Second Series；Vol. 8.；Edited by Philip Schaff and Henry Wace.；Buffalo, NY: Christian Literature Publishing Co.,；1895.；<http://www.newadvent.org/fathers/3202269.htm>.}

% structure: p0001 source=h1 target=section reason=page-title

\section{书信二七〇}

\Signature{凯撒勒雅的圣巴西略}\par

\Emph{无地址。论掠婚。}\par

我感到痛心，发现你们对那些被禁止的罪毫无义愤，似乎不能明白，这种已经发生的\OriginalTerm{掠婚}{raptus}行为是对社会和人性的非法与暴虐，是对自由人的凌辱。我相信，如果你们在这件事上同心一致，这种恶习早就可以从你们的地方驱除。现在，你要彰显基督徒的热忱，按罪行当受的惩罚而行动。无论你在哪里找到那女子，务必把她带走，归还她的父母；将那男子逐出祈祷，并使他受绝罚。按照我已经颁布的教规，他的同伙，连同其全家，都断绝共祷。那在掠婚之后收留女子并藏匿她，甚至抗拒交还的村庄，连同其全体居民，都不得参与祈祷；为的是使众人知道，我们将那掠婚者视为公敌，如同毒蛇或任何野兽一般，要追捕他，并帮助那些被他伤害的人。\par

\SourceRecord{Translated by Blomfield Jackson.；Nicene and Post-Nicene Fathers, Second Series；Vol. 8.；Edited by Philip Schaff and Henry Wace.；Buffalo, NY: Christian Literature Publishing Co.,；1895.；<http://www.newadvent.org/fathers/3202270.htm>.}

% structure: p0001 source=h1 target=section reason=page-title

\section{第二七一书}

\Emph{\Place{凯撒勒雅}的\Person{圣巴西略}}\par

\Emph{致我的同伴\Person{尤西比乌斯}，推荐司铎\Person{西里亚库斯}}\par

你离开后，我立即匆忙赶回城里。何必告诉一个无需被告知、因经验而知的人，我因未找到你而多么苦恼？若能再次见到卓越的尤西比乌斯，拥抱他，在记忆中重游我们的青年时代，并追忆旧日时光——那时我们两人有同一个家，同一个炉灶，同一位老师，同样的闲暇，同样的工作，同样的款待，同样的艰辛，一切共享——那该是多么令人愉快啊！你以为我为了通过实际会面而追忆这一切，为了摆脱我年迈的重负，并仿佛从一个老人变回少年，会有什么不愿付出的呢？但我失去了这一乐趣。至少，通过书信与阁下您相会的特权，以及用我力所能及的最好方式安慰自己的特权，我并未被剥夺。我很幸运遇见了极为可敬的司铎西里亚库斯。我羞于将他推荐给您，并使他通过我成为您的人，以免我似乎在向您献上您已经拥有并珍视的自己的东西，从而在做多余之举。但为真理作证，并将我拥有的最好恩赐给予那些在精神上与我结合的人，这是我的责任。我认为此人在其神圣职位上的无可指责是您所熟知的；但我加以确认，因为我不知道那些不敬畏主、并向所有人下手的人曾对他提出任何指控。即使他们曾做过这种事，此人也不会不配，因为主的敌人反而为他们所攻击之人的品秩辩护，而不是剥夺他们由圣神所赐的任何恩宠。然而，如我所说，甚至没有人考虑过任何不利于此人的事。那么，请仁慈地视他为与我联合、无可指责、堪受一切尊敬的司铎。这样，您将有益于您自己，并使我在喜悦中获得满足。\par

\SourceRecord{Translated by Blomfield Jackson.；Nicene and Post-Nicene Fathers, Second Series；Vol. 8.；Edited by Philip Schaff and Henry Wace.；Buffalo, NY: Christian Literature Publishing Co.,；1895.；<http://www.newadvent.org/fathers/3202271.htm>.}

% structure: p0001 source=h1 target=section reason=page-title

\section{第二七二封书信}

\Person{圣巴西略（凯撒勒雅）}\par

\Emph{致\Person{索夫罗尼乌斯}\OriginalTerm{官署大臣}{magister officiorum}。}\par

执事\Person{阿克提亚库斯}向我报告，有人激怒您反对我，谎称我对阁下心怀恶意。对于像阁下这样地位之人被某些谄媚者追随，我毫不惊讶。高位似乎在某种程度上天然伴随着这类可怜的依附者。他们没有丝毫自身可被人称道的优点，却试图借着他人的不幸来推荐自己。或许，正如霉病是谷粒中滋生的病害，谄媚悄悄潜入友谊，亦是友谊的病害。所以，如我所说，这些人嗡嗡围绕您那明亮而卓越的炉边，如同雄蜂围绕蜂巢，我丝毫不惊讶。但令我惊奇、且全然震惊的是，像您这样以品性庄重特别出众的人，竟被诱导去双耳倾听这些人，并接受他们对我的诽谤。从我的青年直到如今老年，我对许多人怀有感情，但我不知道曾对任何人怀有比对你阁下更深的感情。即使理性没有诱导我尊重如此品性之人，我们自幼的亲密也足以使我的灵魂依附于您。您自己知道习惯与友谊何干。请原谅我的不足，如果我无法展示任何配得上这份偏爱的事。您不会要求我以某些行为来证明我的善意；您会满意于一种心境，它确实为您祈祷，愿您获得一切最好的。愿您的命运永不会低落到需要像我这样微不足道之人的帮助！\par

2. 那么我怎可能说过任何反对您的话，或在\Person{墨摩尼乌斯}的事情上采取任何行动？这些事是执事向我报告的。我怎能把\Person{许梅提乌斯}的财富置于像您这样慷慨大方之人的友谊之前？这些全无真实。我未曾说过或做过任何反对您的事。或许我对一些制造骚动的人说过一些话，给那些谎言提供了某些依据，我说：\Quote{如果那人已决心完成他心中所想，那么，无论你们是否制造骚动，他要做的事必定会做成。你们说话或缄口，没有区别。若他改变主意，你们要当心，不要毁谤我朋友可敬的名声。你们不要以热心于保护人的事为借口，企图借着恐吓与威胁谋取私利。}至于那人立遗嘱之事，我从未对此说过一个字，无论大小，直接或间接。\par

3. 您不可拒绝相信我的话，除非您视我为完全绝望的人，对说谎的大罪毫不在意。请抛开关于此事对我的任何猜疑，此后请视我对您的爱超越一切诽谤的触及。请效法\Person{亚历山大}，他正欲饮药时收到一封信，说他的医生正图谋害死他；他丝毫不信诽谤者，同时读了信并饮了药。我拒绝承认自己在任何方面逊于那些以友谊著称的人，因为我从未被发现有任何背信弃义之处；此外，我由我的天主领受了爱的诫命，不仅因您是整个人类中的一员而爱您，更因我特别认出您是我国家和我个人的恩主。\par

\SourceRecord{Translated by Blomfield Jackson.；Nicene and Post-Nicene Fathers, Second Series；Vol. 8.；Edited by Philip Schaff and Henry Wace.；Buffalo, NY: Christian Literature Publishing Co.,；1895.；<http://www.newadvent.org/fathers/3202272.htm>.}

% structure: p0001 source=h1 target=section reason=page-title

\section{第二七三封信}

\Person{凯撒勒雅的圣巴西略}\par

无收信人。关于赫拉。\par

我确信阁下如此爱我，以致凡关乎我的事，皆视为关乎您自己。因此，我将我可敬的兄弟\Person{赫拉}——我不仅以惯用语称他为兄弟，更因他无限的情谊——推荐予您的极大仁厚与高度重视。我恳求您视他如同与您关系密切一般，并在您可能的情况下，在他需要您慷慨且周详援助的事上给予保护。如此，我将在您已赐予我的众多恩惠之外，再增添这一项恩惠。\par

\SourceRecord{Translated by Blomfield Jackson.；Nicene and Post-Nicene Fathers, Second Series；Vol. 8.；Edited by Philip Schaff and Henry Wace.；Buffalo, NY: Christian Literature Publishing Co.,；1895.；<http://www.newadvent.org/fathers/3202273.htm>.}

% structure: p0001 source=h1 target=section reason=page-title

\section{\WorkTitle{第274封信}}

\Emph{\Person{圣巴西略}}\par

致\Person{希梅里乌斯}大人。\par

我与可敬的兄弟\Person{赫拉}之间的友谊和情谊始于我年少之时，并赖天主的恩宠持续至我的晚年，这件事没有谁比您更清楚。因为主大约在让我结识\Person{赫拉}的同时，也赐予了我对您阁下的爱戴。他现在需要您的庇护，因此我恳求并哀求您，看在我们旧日情谊的份上，施予恩惠，并留意我们现今所处的急需。我请求您将他的事当作您自己的事，使他无需其他保护，而能在他所祈求的一切事上顺利返回我身边。这样，我就能在您赐予我的众多恩惠之上再添这一件。我无法为自己要求比这更重要的恩惠，或更关乎我个人利益的恩惠。\par

\SourceRecord{Translated by Blomfield Jackson.；Nicene and Post-Nicene Fathers, Second Series；Vol. 8.；Edited by Philip Schaff and Henry Wace.；Buffalo, NY: Christian Literature Publishing Co.,；1895.；<http://www.newadvent.org/fathers/3202274.htm>.}

% structure: p0001 source=h1 target=section reason=page-title

\section{书信二七五}

\Emph{\Person{圣巴西略·凯撒勒雅}}\par

无收信人。论\Person{赫拉}。\par

你以对你的至为可敬的弟兄\Person{赫拉}的关爱，先于我的恳求，并以你赐予他的丰厚荣耀和你对他每次的庇护，待他比我所祈求的更好。但我不能一言不发而忽略他的事务，我恳求阁下，为我的缘故，请你加上你对他已表现出的关怀，并使他凯旋胜过敌人的诽谤，返回本国。现在许多人试图羞辱他生活的平安，他也未能免于嫉妒之箭的伤害。针对他的敌人，我们将找到一个可靠的安全保障，只要你愿意将你的庇护延伸到他身上。\par

\SourceRecord{Translated by Blomfield Jackson.；Nicene and Post-Nicene Fathers, Second Series；Vol. 8.；Edited by Philip Schaff and Henry Wace.；Buffalo, NY: Christian Literature Publishing Co.,；1895.；<http://www.newadvent.org/fathers/3202275.htm>.}

% structure: p0001 source=h1 target=section reason=page-title

\section{书信二七六}

\Person{圣巴西略}\par

致伟大的\Person{哈马提乌斯}。\par

人性的普遍法律使长辈成为年轻人的父亲，而我们基督徒特殊的独特法律把我们老年人置于父母的地位对待年轻人。那么，不要以为我无礼，或显示自己不可原谅地多管闲事，如果我为你的儿子向你恳求。在其他方面，我认为你要求他服从是全然正当的；因为，就他的身体而言，他服从于你，既因自然律，也因我们所生活的民法。然而，他的灵魂源自一个更神圣的来源，并且可以合理地认为服从另一种权威。它欠天主的债比任何其他债具有更高的要求。既然他选择了我们基督徒的天主、真天主，而不是你们藉着物质象征所敬拜的众多神祇，不要对他生气。反而要钦佩他灵魂的高贵坚定，牺牲对父亲应有的畏惧和尊敬，以通过真知识和德性生活与天主密切结合。本性自身会感动你，以及你一贯的温和与善良性情，不允许你对他有丝毫怒气。我相信你不会轻视我的调解——或者更确切地说，我是作为你的同乡们的代言人，他们都如此爱你，如此恳切地为你祈求一切祝福，以致他们认为在你身上也接纳了一位基督徒。他们对突然传到城里的报告欣喜若狂。\par

\SourceRecord{Translated by Blomfield Jackson.；Nicene and Post-Nicene Fathers, Second Series；Vol. 8.；Edited by Philip Schaff and Henry Wace.；Buffalo, NY: Christian Literature Publishing Co.,；1895.；<http://www.newadvent.org/fathers/3202276.htm>.}

% structure: p0001 source=h1 target=section reason=page-title

\section{书信277}

\Signature{\Emph{凯撒勒雅的圣巴西略}}\par

致博学之士\Person{马克西姆}。\par

杰出的\Person{德奥特克诺}向我讲述了您的事迹，由此激起了我对您相识的渴望，因为他用言语如此清晰地描绘了您心智的品格。他点燃了我对您炽热的爱慕，以至于若不是我年事已高，先天疾病缠身，且被教会无数的挂虑所束缚，没有任何事能阻止我前来见您。因为的确，一个名门望族的成员、出身显赫的人，若采纳福音的生活方式，以理性克制青春的冲动，使肉体的情感服从理智；展示与他基督徒身份相符的谦卑，按本分反思自己从何而来、往何处去——这并非小小的收获。正是对我们本性的这种反思，抑制了心灵的傲慢，驱散了所有自夸和骄傲。总之，它使人成为我们主的门徒，主曾说：\BibleQuote{你们向我学习，因为我是良善心谦的。}\BibleRef{玛 11:29}实在的，极亲爱的儿子，唯一值得我们的努力和赞美的是我们永恒的福祉；这是来自天主的尊荣。\par

人间万事比影子更虚幻，比梦更欺骗。青春凋谢比春天的花朵更快；容貌因年老或疾病而憔悴。财富不确定；荣华易逝。对技艺和学问的追求局限于现世生命；人人都向往的雄辩魅力仅止于耳；而修德的实践对其拥有者是宝贵的财富，对所有目睹者是赏心悦目的景象。请以此为您的学业；这样您就配得上主所应许的美好事物。\par

但是，列举获得并享有这些福乐的方法，超出了本信的意图。然而，在听了我的兄弟\Person{德奥特克诺}的讲述后，我想到了给您写这些。我祈祷他永远说实话，尤其是在他对您的描述中；愿主在您内更受光荣，因为您虽出自异域之根，却充盈着最珍贵的虔敬之果。\par

\SourceRecord{Translated by Blomfield Jackson.；Nicene and Post-Nicene Fathers, Second Series；Vol. 8.；Edited by Philip Schaff and Henry Wace.；Buffalo, NY: Christian Literature Publishing Co.,；1895.；<http://www.newadvent.org/fathers/3202277.htm>.}

% structure: p0001 source=h1 target=section reason=page-title

\section{第278封书信}

\Emph{\Person{凯撒勒雅的圣巴西略}}\par

致\Person{瓦勒里亚努斯}。\par

我曾在\Place{奥尔法嫩}时渴望见到阁下；也曾希望，当您居住在\Place{科萨盖纳}时，不会有任何事阻止您前来参加我预期在\Place{阿塔盖纳}举行的教务会议。然而，我未能举行该会议，于是我的愿望是到山区见您；因为在那里的\Place{埃维苏斯}，由于靠近那一带，曾给我的会面带来希望。但既然两次都失望了，我便决定写信，恳请您屈尊来访我；因为我认为年轻人前来探望老年人是理所当然的。此外，在我们会面时，我将向您提出我的忠告，涉及您与\Place{凯撒勒雅}的某些人的商谈：此事要得到正确解决需要我的介入。如果您同意的话，就请不要迟疑地到我这里来吧。\par

\SourceRecord{Translated by Blomfield Jackson.；Nicene and Post-Nicene Fathers, Second Series；Vol. 8.；Edited by Philip Schaff and Henry Wace.；Buffalo, NY: Christian Literature Publishing Co.,；1895.；<http://www.newadvent.org/fathers/3202278.htm>.}

% structure: p0001 source=h1 target=section reason=page-title

\section{书信二七九}

\Emph{\Person{凯撒勒雅的圣巴西略}}\par

致总督\Person{莫德斯图}。\par

尽管我经由同样众多的信差呈递给阁下的信件如此之多，但鉴于您对我所显示的特别尊荣，我无法认为它们的数量之多会令您感到任何烦扰。\par

因此，我毫不犹豫地将随信托付给这位弟兄：我知道他将如愿以偿，而您也会视我为恩人，因为我提供了机会满足您仁慈的意愿。他恳求您的庇护。他的事将由他亲自陈明，只要您屈尊以善意的眼光看待他，并鼓励他在如此威严的权威面前自由陈辞。请接受我的保证：凡向他所施的仁慈，我将视同对我本人。他离开\Place{提雅纳}前来见我的特殊原因，是他极为重视呈上一封由我亲笔所写、支持他请求的信函。为使他不会失望，为使我继续蒙受您的眷顾，为使您对一切美善的热忱能在当前之事中充分施展——基于这些理由，我恳求您仁慈地接待他，并赐他一个在您近侍中的位置。\par

\SourceRecord{Translated by Blomfield Jackson.；Nicene and Post-Nicene Fathers, Second Series；Vol. 8.；Edited by Philip Schaff and Henry Wace.；Buffalo, NY: Christian Literature Publishing Co.,；1895.；<http://www.newadvent.org/fathers/3202279.htm>.}

% structure: p0001 source=h1 target=section reason=page-title

\section{书信 280}

\Emph{\Place{凯撒勒雅}的\Person{圣巴西略}}\par

致长官\Person{莫德斯图}。\par

我深知，以书信向阁下这般地位的人恳求实属冒昧；然而您过去赐予我的尊荣已打消我的一切顾虑。因此我满怀信心地写信。\par

我恳求的对象是我的一位亲戚，他为人正直，值得尊敬。他是这封信的送信人，于我而言犹如儿子一般。他实现愿望只需您的恩惠。因此，伏请屈尊通过上述送信人接受我这封助其恳求的信。我恳求您给他机会，与能助他一臂之力的人面谈，以便他陈述自己的事务。如此，在您的指示下，他将迅速得偿所愿；而我则将有机会夸耀，因天主的恩宠，我找到了一位保护者，他将我友人的恳求视作自己理应庇护之事。\par

\SourceRecord{Translated by Blomfield Jackson.；Nicene and Post-Nicene Fathers, Second Series；Vol. 8.；Edited by Philip Schaff and Henry Wace.；Buffalo, NY: Christian Literature Publishing Co.,；1895.；<http://www.newadvent.org/fathers/3202280.htm>.}

% structure: p0001 source=h1 target=section reason=page-title

\section{\newline{}\WorkTitle{第二八一封书信}\newline{}}

\newline{}\Emph{\Place{凯撒勒雅}的\Person{圣巴西略}}\newline{}\par

\newline{}致总督\Person{莫德斯多}。\newline{}\par

\newline{}我铭记着您与其他人一同赐予我的鼓励，使我得以致函阁下，这份殊荣令我深感荣幸。我珍视这份特权并享受您的仁恩。\newline{}\par

\newline{}我为自己能有您这样的通信者而庆幸，也为阁下有机会通过回信赐予我荣誉而欣慰。\newline{}\par

\newline{}我为我的挚友\Person{赫拉迪乌斯}恳求您的宽仁。我祈求他得以免除税务评估官的烦扰，从而能为我们祖国的利益效力。\newline{}\par

\newline{}您已如此恩允在前，故我再次恳求，请您致函行省总督，下令解除\Person{赫拉迪乌斯}所受的这一苦役。\newline{}\par

\SourceRecord{Translated by Blomfield Jackson.；Nicene and Post-Nicene Fathers, Second Series；Vol. 8.；Edited by Philip Schaff and Henry Wace.；Buffalo, NY: Christian Literature Publishing Co.,；1895.；<http://www.newadvent.org/fathers/3202281.htm>.}

% structure: p0001 source=h1 target=section reason=page-title

\section{\WorkTitle{第282封信}}

\Emph{\Person{圣巴西略\Place{凯撒勒雅}}}\par

致一位主教。\par

你责怪我没有邀请你；而当你受到邀请时，你又不来。你先前那个借口是空洞的，从你第二次的行为可以清楚看出。因为如果你先前被邀请过，很可能你根本就不会来。\par

不要再鲁莽行事，要服从这次邀请；因为你知道，邀请的重复加强了控诉，第二次使先前的指控变得可信。\par

我劝你常常容忍我；或者即使你不能，至少你有责任不要忽视殉道者，你被邀请参加他们的纪念。因此，请你同时为我们两者提供服务；如果你不愿意这样，至少为更值得的一方。\par

\SourceRecord{Translated by Blomfield Jackson.；Nicene and Post-Nicene Fathers, Second Series；Vol. 8.；Edited by Philip Schaff and Henry Wace.；Buffalo, NY: Christian Literature Publishing Co.,；1895.；<http://www.newadvent.org/fathers/3202282.htm>.}

% structure: p0001 source=h1 target=section reason=page-title

\section{\WorkTitle{书信二八三}}

\Signature{圣巴西略·凯撒勒雅}\par

致一位寡妇。\par

我希望在为山地定下日期之后，能找到适合会谈的一天。我看不出我们有会面的机会（除非主超出我的期望另有安排），除非在一次公开的会谈中。\par

你可以从自己的经验想象我的处境。若你在管理一个家庭时都被如此多的忧虑所困扰，想想看，每天要给我带来多少分心的事？\par

我认为，你的梦更充分地揭示了必须为属神的默观做准备，并培养那种通常用来看见天主的内心视觉。既然你享有圣经的安慰，你既不需要我的帮助，也不需要任何其他人的帮助来理解你的职责。你有圣神全备的劝告和指引，引领你走向正义。\par

\SourceRecord{Translated by Blomfield Jackson.；Nicene and Post-Nicene Fathers, Second Series；Vol. 8.；Edited by Philip Schaff and Henry Wace.；Buffalo, NY: Christian Literature Publishing Co.,；1895.；<http://www.newadvent.org/fathers/3202283.htm>.}

% structure: p0001 source=h1 target=section reason=page-title

\section{\WorkTitle{第284封信}}

\Emph{\Person{凯撒勒雅的圣巴西略}}\par

致隐修士案件的评税官。\par

关于隐修士，阁下已有现行规则，我相信，因此我无需为他们请求特别恩惠。\par

他们早已辞别了此世，克制了自己的身体，既无钱财可花，也无体力可献，为公共福祉效力，因此我认为，应将他们免除赋税。因为若他们的生活与誓愿相符，则他们既无钱财也无身体：钱财已分施给穷人，身体因祈祷和守斋而憔悴。\par

我知道，您必会特别尊敬度此生活的人；甚至渴望得到他们的转求，因为他们按福音生活，能在天主面前有效力。\par

\SourceRecord{Translated by Blomfield Jackson.；Nicene and Post-Nicene Fathers, Second Series；Vol. 8.；Edited by Philip Schaff and Henry Wace.；Buffalo, NY: Christian Literature Publishing Co.,；1895.；<http://www.newadvent.org/fathers/3202284.htm>.}

% structure: p0001 source=h1 target=section reason=page-title

\section{第285号书信}

\Person{圣巴西略凯撒勒雅}\par

无收信人称呼。\par

这封信的收信人，就是那位肩负我们教会的照料及其财产管理之责的人——我亲爱的儿子。\par

请恩准他在向尊者您呈报的事项上享有言论自由，并聆听他表述自己的见解；这样，我们的教会终将恢复元气，从此摆脱这多头海蛇的纠缠。\par

我们的财产就是我们的贫穷；以至于我们一直在寻找一个能把我们从这重担中解脱出来的人；因为教会财产的开支超过了它所带来的任何收益。\par

\SourceRecord{Translated by Blomfield Jackson.；Nicene and Post-Nicene Fathers, Second Series；Vol. 8.；Edited by Philip Schaff and Henry Wace.；Buffalo, NY: Christian Literature Publishing Co.,；1895.；<http://www.newadvent.org/fathers/3202285.htm>.}

% structure: p0001 source=h1 target=section reason=page-title

\section{书信二八六}

\Emph{\Person{圣巴西略·凯撒勒雅}}\par

\Emph{致主理司}。\par

鉴于某些流浪者在教会内被逮捕，因为他们违背天主的诫命，偷窃了一些穷人的衣物，这些衣物本身价值不高，但却是他们宁愿穿在身上也不愿脱下的；而您认为，根据您的职务，您本人应负责看管这些犯法者：——我在此声明，我愿你知道，在教会内所犯的罪行应由我们负责施行惩罚，世俗政权的介入在此类情况下是多余的。因此，对于偷盗的财物，正如您手中的文件以及在目击证人面前制作的副本所记载的，我责令您予以扣留，一部分留作将来可能的索赔，其余的分发给目前的申请人。\par

至于犯法者——希望他们在主的管教和劝戒中得到纠正。我期望通过这种方式使他们逐步改过自新。因为当人间法庭的鞭笞无效时，我常常见到天主可畏的审判却能奏效。然而，如果您希望也将此事提交给伯爵，鉴于我对他的公正和正直深具信心，我任由您自行决定。\par

\SourceRecord{Translated by Blomfield Jackson.；Nicene and Post-Nicene Fathers, Second Series；Vol. 8.；Edited by Philip Schaff and Henry Wace.；Buffalo, NY: Christian Literature Publishing Co.,；1895.；<http://www.newadvent.org/fathers/3202286.htm>.}

% structure: p0001 source=h1 target=section reason=page-title

\section{书信第287号}

\Emph{\Person{圣巴西略}（\Place{凯撒勒雅}）}\par

无收信人地址。\par

和此人打交道很困难。我几乎不知道该如何对付这样一个诡计多端、而且从证据来看如此绝望的人。当被传唤到法庭时，他不出庭；如果他真的出庭，他又如此能言善辩、满口誓言，以至于我认为自己能尽快摆脱他是幸运的。我常见他把对他的指控转嫁到指控者头上。总之，世上没有一个活物像他这样狡猾多端、善于作恶。略加了解就足以证明这一点。那么你们为何还要向我申诉？为何不立即让自己忍受他的恶待，如同对待天主义怒的惩罚呢？\par

同时，你们不可因与恶人接触而被玷污。\par

因此我命令：禁止他及其全家参与教会的礼仪，并与教会圣职人员的一切来往。这样，他既成了儆戒，或许能认识到自己罪孽深重。\par

\SourceRecord{Translated by Blomfield Jackson.；Nicene and Post-Nicene Fathers, Second Series；Vol. 8.；Edited by Philip Schaff and Henry Wace.；Buffalo, NY: Christian Literature Publishing Co.,；1895.；<http://www.newadvent.org/fathers/3202287.htm>.}

% structure: p0001 source=h1 target=section reason=page-title

\section{书信 288}

\Person{圣巴西略（该撒利亚的）}\par

无收信人。绝罚性质。\par

当公开惩罚不能使一个人觉悟，或排除于教会祈祷之外不能驱使他悔改时，只剩下按照我们主的指示对待他——如经上所记载：\Quote{如果你的兄弟得罪了你……你要单独指出他的错处……} \Quote{如果他}然后\Quote{不听，你要告诉教会；如果他连教会也不听，你就把他看作外邦人和税吏。}\BibleRef{玛 18:15} 这一切我们已经在那个人的事上做了。首先，他被指认了他的过错；然后，他在一两个人面前被定罪；第三，在教会面前。这样我们做了庄严的见证，他却没有听从。今后，让他被绝罚。\par

此外，要在全区发布公告，禁止他参与任何日常交往；这样，由于我们断绝与他的一切来往，他完全成为魔鬼的食物。\par

\SourceRecord{Translated by Blomfield Jackson.；Nicene and Post-Nicene Fathers, Second Series；Vol. 8.；Edited by Philip Schaff and Henry Wace.；Buffalo, NY: Christian Literature Publishing Co.,；1895.；<http://www.newadvent.org/fathers/3202288.htm>.}

% structure: p0001 source=h1 target=section reason=page-title

\section{\WorkTitle{书信 289}}

\Emph{\Person{圣巴西略（\Place{凯撒勒雅}）}}\par

无收信人。关于一位受苦的妇人。\par

我认为让有罪者不受惩罚，以及超过适当的刑罚限度，是同样的错误。因此，我判定这个人应受的处罚——逐出教会——这是我认为必须施加的。我劝勉那位受苦者不要为自己复仇，而要将冤屈交给天主来纠正。这样一来，如果我的劝告有任何分量，她就会听从，因为我所用的言辞远比任何要求服从的信件更易获得信任。\par

因此，即使在听完她的陈述——其中包含了足够严重的内容——之后，我仍然保持缄默；甚至现在，我也不确定是否适合再次处理这个相同的问题。\par

\Quote{因为，她说，我放弃了丈夫、孩子以及生活中的一切享乐，以求获得这一单一目标：天主的恩宠和在人间的好名声。然而有一天，那个从少年时代就以败坏家庭为能事的冒犯者，以他惯常的无耻，强行闯入我的家；就这样，仅仅在一次会面的限度内，就结下了相识。只是因为我对这个人的无知，以及因缺乏经验而来的胆怯，我犹豫了，没有公开将他赶出家门。然而，他的不敬和傲慢达到了如此地步，以致他用诽谤充满了全城，并公开张贴诽谤性的布告于教会门上攻击我。诚然，这种行为使他触犯了法律；但他仍然继续对我进行诽谤攻击。市场再次充满他的辱骂，还有体育馆、戏院以及那些因习气相投而接纳他的房屋。他的极端行为并未使人们认识到我所显著的美德，因为我的名声已被普遍描绘为一种放荡不羁的性格。在这些诽谤中，她继续说道，有些人感到快乐——因为人们自然乐于贬低他人；有些人自称痛苦，却不同情；另一些人相信这些诽谤的真实性；还有一些人因他誓言的一贯性而犹疑不定。但我没有同情。现在我真的开始意识到我的孤独，并为自己哀叹。我没有兄弟、朋友、亲戚，没有奴仆，无论奴仆或自由人，总之，没有任何人与我分担痛苦。尽管如此，我认为，在这邪恶的憎恨者如此之少的城市里，我比任何人都更值得怜悯。他们彼此施加暴力；但暴力，尽管他们看不到，却在循环运动，最终会追上他们每一个人。}\par

她用如此甚至更恳切的言辞讲述她的故事，泪流无数，然后离去。她并未完全免除我的责任；认为当我应该像父亲一样同情她时，却对她的烦恼漠不关心，并且过于冷静地看待他人的痛苦。\par

因为她坚持说：\Quote{你所劝我忽视的并非金钱的损失，也非身体痛苦的忍受，而是受损的名誉，一种涉及整个教会损失的伤害。}\par

这是她的呼吁；现在我得请求您，最尊贵的先生，考虑您希望我给她什么答复。我心中已作出的决定是：不将冒犯者交给官府；但也不解救那些已在官府拘禁中的人，因为宗徒早已宣布，官长对于作恶的人应当是可畏的；因为经上说：\BibleQuote{他并不是无端地佩剑。} \BibleRef{罗 13:4} 因此，把他交出去有违于我的人道；而释放他又会鼓励他的暴力。\par

然而，也许您会推迟行动，直到我到达。届时我将向您证明，因为无人听从我，我什么都无法实现。\par

\SourceRecord{Translated by Blomfield Jackson.；Nicene and Post-Nicene Fathers, Second Series；Vol. 8.；Edited by Philip Schaff and Henry Wace.；Buffalo, NY: Christian Literature Publishing Co.,；1895.；<http://www.newadvent.org/fathers/3202289.htm>.}

% structure: p0001 source=h1 target=section reason=page-title

\section{\WorkTitle{二九〇書}}

\Emph{\Person{凯撒勒雅的圣巴西略}}\par

致\Person{纳克塔留斯}。\par

愿丰厚的祝福归于那些鼓励阁下与我保持经常通信的人！不要把这祝愿仅仅视为客套，而要看作是我对您话语价值之真诚信念的表达。有谁能比\Person{纳克塔留斯}更值得我尊敬呢？——我从他早年最美好的童年时期就认识他，如今他借着实践每一种美德，已达到了极高的地位。因此，在我所有的朋友中，最亲爱的就是你的信使。\par

关于地区主管的选举，天主禁止我为了取悦人而做任何事，无论是听从恳求还是屈服于恐惧。那样我就不是管家，而是小贩，用天主的恩赐换取人的欢心。但既然选票是由凡人投出的，他们只能根据外表作证；而合适人选的选择，则应谦卑地交托给那洞悉人心秘密的天主。或许对每个人来说，在提供投票证据后，最好避免一切激动和争论，好像证明中涉及某些私利，并祈祷天主，使有利的事不致被隐藏。这样结果就不再归于人，而是感恩于天主。因为这些事情，如果是出于人，就不能说它们是真实的；它们只是虚假，毫无实质。\par

此外，当一个人竭尽全力去达到目的时，他极有可能把罪人也拉到自己一边；而罪过很多，这是人性的软弱，甚至在我们最意想不到的地方也是如此。\par

再者，在私下商议中，我们常给朋友忠告，即使他们不采纳，我们也不生气。那么，当不是人在出主意，而是天主在决定时，我们若因未被选上而反对天主的决定，又岂能愤愤不平呢？\par

如果这些事是由人赐予人的，我们又何必向别人求取呢？每个人自己取来岂不更好？但如果是天主的恩赐，我们就该祈祷，而不该忧伤。在祈祷中，我们不应寻求自己的意愿，而应交给那安排一切的最好的天主。\par

愿圣洁的天主使你的家远离一切悲伤的滋味；并赐予你和你的家人免于伤害和疾病的生活。\par

\SourceRecord{Translated by Blomfield Jackson.；Nicene and Post-Nicene Fathers, Second Series；Vol. 8.；Edited by Philip Schaff and Henry Wace.；Buffalo, NY: Christian Literature Publishing Co.,；1895.；<http://www.newadvent.org/fathers/3202290.htm>.}

% structure: p0001 source=h1 target=section reason=page-title

\section{书信二九一}

\Emph{圣\Person{巴西略}·\Place{凯撒勒雅}}\par

\Emph{致乡间主教\Person{弟茂德}}.\par

书信的适当篇幅以及我对您的称呼方式，使我无法写出我所有想法；但同时，当我心中因对您义怒燃烧时，缄默不言又几乎不可能。我将采取折中之道：写一些，省略一些。因为我希望，若可以的话，以坦诚而友好的言辞责备您。\par

是的！我所认识的那个自幼年起便如此专注于正直与苦修生活，甚至被认为过分的\Person{弟茂德}，如今却放弃了探究如何与天主结合的方法；如今他首先想到的是别人可能怎么看他，从而过着依赖他人意见的生活；主要焦虑如何服侍朋友而不招致敌人嘲笑；惧怕世俗的羞辱，认为那是大不幸。他难道不知道，当他忙于这些琐事时，正不知不觉地忽视了最高的利益吗？因为我们不能同时从事于两者——此世之事与天上之事——圣经充满了对我们的教导。不仅如此，大自然本身也充满这样的实例。在心灵官能的运用中，同时思考两个想法是完全不可能的。在感官知觉中，同时接受两个声音进入耳朵并分辨它们，即使我们有两条通畅的通道，也是不可能的。同样，我们的眼睛，除非它们都注视我们所见之物，否则无法准确执行其功能。\par

关于大自然就说到这里；但向你背诵圣经的证词，正如俗语所说，\Quote{把猫头鹰送到\Place{雅典}}，一样荒谬。那么，为什么将不相容的事物——世俗生活的喧嚣与宗教实践——结合在一起呢？\par

远离喧嚣；不再成为烦扰的起因或目标；让我们独处。我们早已将宗教作为目标；让我们以达成它为实践，并向那些想要侮辱我们的人表明，他们不能随意烦扰我们。但这只有在我们清楚地向他们表明我们并未提供可攻击的把柄时才会实现。\par

目前就到此为止吧！但愿有一天我们能见面，更圆满地考虑那些关乎我们灵魂福祉的事情；这样我们就不会过于专注于虚妄的思绪，因为死亡终有一日会追上我们。\par

我很高兴收到您好意送来的礼物。它们本身就很受欢迎；想到是谁送来的，更使它们受欢迎数倍。来自\Place{本都}的礼物，即书写板与药品，当我送出时请善意接纳。目前这些东西不在我身边。\par

注：本笃会版编者将编号292至366的书信归入\Quote{\OriginalTerm{第三类}{Classis Tertia}}，没有时间标记。有些存疑，有些明显是伪作。其中我收录了那些看似最重要的。\par

\SourceRecord{Translated by Blomfield Jackson.；Nicene and Post-Nicene Fathers, Second Series；Vol. 8.；Edited by Philip Schaff and Henry Wace.；Buffalo, NY: Christian Literature Publishing Co.,；1895.；<http://www.newadvent.org/fathers/3202291.htm>.}

% structure: p0001 source=h1 target=section reason=page-title

\section{第二九二封书信}

\Emph{\Place{凯撒勒雅}的\Person{圣巴西略}}\par

致\Person{帕拉迪乌斯}。\par

天主已实现了我一半的渴望，即祂恩准我与我们那位贤淑的姊妹、你的妻子会面。另一半祂也能完成；因此，一旦得见尊驾，我必将向天主献上我全然的感恩。\par

我更加渴望见到你，因为我听说你已佩上了那伟大的装饰——不朽的衣袍，它遮盖我们的必死性，并使肉身的死亡隐没不见；藉此，那可朽坏的被不朽所吞没。\par

如此，仁慈的天主现已使你与罪分离，使你与祂自己结合，打开了天堂之门，并指明了通往天上福乐的道路。因此，我凭那使你超越所有人的智慧恳求你，谨慎地接受天主的恩宠，作这宝藏的信实守护者，作这君王恩赐的保管者，以万般小心看守它。持守这正义的印章毫无玷污，好使你能站在天主面前，在圣人的光辉中闪耀。不要让任何污点或皱纹玷污这不朽的纯洁衣袍；却要在你所有肢体中保持圣洁，因为你们已穿上基督。因为，正如经上所说：\BibleQuote{你们凡受洗归于基督的，都已穿上基督了。}\BibleRef{迦 3:27}因此，你们所有的肢体都当圣洁，以配得上那圣洁与光明的衣袍。\par

\SourceRecord{Translated by Blomfield Jackson.；Nicene and Post-Nicene Fathers, Second Series；Vol. 8.；Edited by Philip Schaff and Henry Wace.；Buffalo, NY: Christian Literature Publishing Co.,；1895.；<http://www.newadvent.org/fathers/3202292.htm>.}

% structure: p0001 source=h1 target=section reason=page-title

\section{Letter 293}

\Emph{\Person{凯撒勒雅的圣巴西略}}\par

致\Person{犹利安}。\par

你长久以来可好？你的手是否完全恢复功能？其他事情进展如何？是如你所愿、如我所祈，还是按你的计划？\par

人若轻易善变，其生活自然难以有序；但若心思坚定、稳固而不移，则其生活必然与志向相符。\par

的确，舵手无法随心所欲地使海面平静；但对我们而言，平息内心澎湃的情欲风暴，超越外在的侵袭，使生活宁静安稳，却是轻而易举的。正直之人不为损失、疾病或人生的其他灾祸所动；因为他内心与天主同行，目光定睛于未来，轻松且从容地度过人间兴起的风暴。\par

不要为世上的烦恼所困扰。这样的人如同肥鸟，徒有飞翔之能，却如野兽般在地上爬行。然而你——因为我曾见你在困境中——恰似在海上竞渡的泳者。\par

一爪可见全狮；从浅交之中，我想我已完全认识你。我视为大事的是，你对我有所珍视，我并未从你思念中消失，而是常在你的记忆里。\par

如今，写信便是记忆的证明；你写得越勤，我便越是喜悦。\par

\SourceRecord{Translated by Blomfield Jackson.；Nicene and Post-Nicene Fathers, Second Series；Vol. 8.；Edited by Philip Schaff and Henry Wace.；Buffalo, NY: Christian Literature Publishing Co.,；1895.；<http://www.newadvent.org/fathers/3202293.htm>.}

% structure: p0001 source=h1 target=section reason=page-title

\section{书信二九四}

\Emph{\Person{圣巴西略}（\Place{凯撒勒雅}）}\par

致\Person{斐斯托斯}与\Person{玛格努斯}。\par

父亲之为子女筹谋，农夫之于植物与庄稼，师长之于学生，尤其是当日禀之善显示出前途之兆时，均属分内之事，此无疑义。\par

农夫见谷穗成熟或作物生长，便以劳作为乐；师长见学生知识增进，父亲见儿子身材增长，均皆欢喜。然而我对尔之牵挂更大，所怀希望更高；因为虔诚超越一切技艺，超越一切动物与果实之总和。\par

我已将虔诚种植于尔心，当时它尚纯洁柔嫩；我培育它，期望它达到成熟并按季结果。同时，我的祈祷因尔之好学而得以推进。尔深知我祝愿尔，且天主的恩宠临于尔之努力；因为若方向正确，天主无论是否被呼求，都近在咫尺，予以助成。\par

凡爱天主者皆倾向于教导；非但如此，凡有能力教授有益之事者，其热忱几乎无法抑制；但首先必须清除听众心中一切障碍。\par

身体的分离并非教导的障碍。造物主以其圆满的爱与智慧，并未将我们的心智局限于身内，也未将言谈之能囿于舌上。获益的能力甚至能因时间间隔而得益；因此我们不仅能将教导传给远处的人，甚至能传给后世的人。经验证实了此言：多年前生活的人，通过保留在著作中的教导教诲后人；而我们虽身体相隔遥远，却在思想中常相近，能轻松交谈。\par

教导既不为海所限，也不为陆所囿，只要我们关心自己灵魂的益处。\par

\SourceRecord{Translated by Blomfield Jackson.；Nicene and Post-Nicene Fathers, Second Series；Vol. 8.；Edited by Philip Schaff and Henry Wace.；Buffalo, NY: Christian Literature Publishing Co.,；1895.；<http://www.newadvent.org/fathers/3202294.htm>.}

% structure: p0001 source=h1 target=section reason=page-title

\section{\WorkTitle{Letter 295}}

\Emph{\Place{凯撒勒雅}的\Person{圣巴西略}}\par

致隐修士们。\par

我不认为在亲自对你们说过话之后，还需要进一步将你们交托于\Person{天主}的恩宠。那时我曾命令你们效法\Person{宗徒们}的生活方式，度共同生活。你们接受了这有益的教诲，并为此感谢\Person{天主}。\par

因此，你们的行为并非主要归因于我所说的话，而是归因于我将这些话付诸实践的训导；这同时有利于你们这些接受者的裨益，也有利于我这位提出劝告者的安慰，并归光荣与赞颂于\Person{基督}，我们是因\Person{祂}的名而被称呼的。\par

为此，我曾派遣我们所爱的弟兄到你们那里，以便他了解你们的热情，激发你们的懈怠，向我报告反对的情况。因为我极切希望看到你们全体合为一体，并听说你们不满足于过一种没有见证的生活，而是担当起彼此监督勤勉、互为成功见证的责任。\par

这样，你们每个人都将获得圆满的赏报，不仅为了自己，也为了弟兄的进步。而且，正如所应得的，你们将因不断的交往和劝勉，在言语和行为上成为彼此获益的来源。但我尤其劝告你们要牢记\Person{教父们}的信德，不要被那些在你们退隐中试图使你们背离的人所动摇。因为你们知道，除非因着对\Person{天主}的信德而受到光照，生活的严格毫无益处；同样，缺乏善工的正统信德宣认，也不能使你们站在\Person{主}面前。\par

信德与善工必须结合：这样，\Person{天主}的人才能成全，他的生命不致因任何缺陷而跛行。\par

因为那拯救我们的信德，正如\Person{宗徒}所说，是藉着爱德运作的。\par

\SourceRecord{Translated by Blomfield Jackson.；Nicene and Post-Nicene Fathers, Second Series；Vol. 8.；Edited by Philip Schaff and Henry Wace.；Buffalo, NY: Christian Literature Publishing Co.,；1895.；<http://www.newadvent.org/fathers/3202295.htm>.}

% structure: p0001 source=h1 target=section reason=page-title

\section{第299封信}

\Person{圣巴西略（凯撒勒雅）}\par

\Emph{致一位监察官}。\par

在您告诉我之前，我已知道您不喜欢担任公职。有一句古老的谚语：那些渴望度虔敬生活的人，不会欣然投身官职。在我看来，官吏的情形与医生相似：他们目睹可怕的景象，遭遇恶臭，从他人的灾难中为自己招来麻烦——至少那些真正的官吏是如此。所有从商之人也追求利润，并为这种荣耀所激动，认为获取一些权柄和影响力是最大的好处，借此得以有益于朋友、惩罚敌人、为自己谋取所需。您不是这种人。您怎么会是呢？您甚至自愿从国家的高位上引退。您本可像管理一栋房屋那样治理城市，却选择了无忧无虑的生活。您将不为己添烦、不扰他人置于比他人招人厌烦更高的价值。然而，主认为伊波拉地区不应受小贩的辖制，也不应沦为纯粹的奴隶市场。祂的旨意是让其中的每个人都按正当方式登记。那么，请接受这一职责吧？我知道它令人烦恼，但它可能为您带来天主的赞许。不要谄媚权贵，也不要轻视贫苦之人。向您治下的所有人显示一种比任何天平还要精确的公正之心。这样，在那些托付您这些职责的人眼中，您对正义的热忱将显而易见，他们会以无比的钦佩看待您。即使他们未留意您，我们的天主也不会忽视您。祂为我们善工所预备的赏报是伟大的。\par

\SourceRecord{Translated by Blomfield Jackson.；Nicene and Post-Nicene Fathers, Second Series；Vol. 8.；Edited by Philip Schaff and Henry Wace.；Buffalo, NY: Christian Literature Publishing Co.,；1895.；<http://www.newadvent.org/fathers/3202299.htm>.}

% structure: p0001 source=h1 target=section reason=page-title

\section{书信三〇三}

\Emph{凯撒勒雅的圣巴西略}\par

\OriginalTerm{私产伯爵}{Comes Privatarum}。\par

我想，是当地居民提供的不实信息，导致您向这些人征收了母马税。此事极不公平，必然令阁下不悦，也因我与受害者的密切关系而感到痛心。因此，我立即恳求阁下，切勿让这些不义之徒的恶谋得逞。\par

\SourceRecord{Translated by Blomfield Jackson.；Nicene and Post-Nicene Fathers, Second Series；Vol. 8.；Edited by Philip Schaff and Henry Wace.；Buffalo, NY: Christian Literature Publishing Co.,；1895.；<http://www.newadvent.org/fathers/3202303.htm>.}

% structure: p0001 source=h1 target=section reason=page-title

\section{\WorkTitle{书信 306}}

\Emph{\Person{圣巴西略}（\Place{凯撒勒雅}）}\par

\Emph{致\Place{塞巴斯特亚}总督.}\par

我知道阁下欣然接受我的信件，也明白原因。您热爱一切善事，乐于行善。因此，每当我给您机会展示您的宽宏大量时，您都急切地期待我的信件，因为您知道它们为善行提供了机会。现在，再一次，请看一个机会，让您展现一切正直的标记，同时公开彰显您的德行！有几个人从\Place{亚历山大}前来，为了履行众人对死者应尽的必要义务。他们请求阁下下令，允许他们在官方许可下，将一位在\Place{塞巴斯特亚}去世的亲属的遗体运走，当时军队驻扎在那里。他们进一步恳求，尽可能给予他们公费旅行的帮助，以便因您的慷慨，他们在长途旅行中得到一些帮助和安慰。这个消息将远传至伟大的\Place{亚历山大}，并在那里传播阁下惊人仁慈的报道。这一点您很清楚，无需我赘言。我将为这一额外的恩惠，对您为我所做的一切表示感激。\par

\SourceRecord{Translated by Blomfield Jackson.；Nicene and Post-Nicene Fathers, Second Series；Vol. 8.；Edited by Philip Schaff and Henry Wace.；Buffalo, NY: Christian Literature Publishing Co.,；1895.；<http://www.newadvent.org/fathers/3202306.htm>.}

% structure: p0001 source=h1 target=section reason=page-title

\section{书信 334}

\Person{圣巴西略（凯撒勒雅的）}\par

致一位抄写员。\par

写直了，并使线条保持平直。不要让手写得太高或太低。避免迫使笔斜着走，就像\Person{伊索}的螃蟹那样。要径直向前，仿佛循着木匠的直尺线前进，这尺子始终精确，防止任何歪斜。斜线是不雅观的。直线才悦目，不会让读者的眼睛像摆杆一样上下摇晃。我在读你的字时就遭遇了这种命运。因为字母像梯子一样排列，当我不得不从一行移到下一行时，我得爬到前一行末尾；发现没有连接时，又得退回去重新寻找正确的顺序，就像故事中的\Person{忒修斯}循着\Person{阿里阿德涅}的线团一样后退着沿着犁沟走。写直了，不要用你歪斜不整的字迹使我们困惑。\par

\SourceRecord{Translated by Blomfield Jackson.；Nicene and Post-Nicene Fathers, Second Series；Vol. 8.；Edited by Philip Schaff and Henry Wace.；Buffalo, NY: Christian Literature Publishing Co.,；1895.；<http://www.newadvent.org/fathers/3202334.htm>.}

% structure: p0001 source=h1 target=section reason=page-title

\section{第三三五函}

\Emph{\Place{凯撒勒雅}的\Person{圣巴西略}}\par

\Emph{\Person{巴西略}致\Person{里巴尼乌}}。\par

我实在羞于将\Place{卡帕多细亚}人一个个地送到你那里去。我更愿意促使我们所有的青年都献身于文学与学问，并在他们的训练中受益于你的教导。但一下子把他们全都召集起来是不现实的，因为他们各自选择适合自己的道路。因此，我将那些不时被争取过来的人送到你那里；这样做时，我确信我赐予他们的恩惠，如同那些将口渴之人带到泉源的人所赐予的一样大。我现在送去的这个少年，一旦与你相处，必将因其本身而受到高度珍视。他已经因其父亲而闻名，其父在我们中间因生活正直和在我们团体中的权威而享有声誉。此外，他是我自己的密友。为回报他对我的友谊，我将他的儿子引荐给你——这是一份所有能够评判一个人高尚品格的人都应当热切祈求的恩惠。\par

\SourceRecord{Translated by Blomfield Jackson.；Nicene and Post-Nicene Fathers, Second Series；Vol. 8.；Edited by Philip Schaff and Henry Wace.；Buffalo, NY: Christian Literature Publishing Co.,；1895.；<http://www.newadvent.org/fathers/3202335.htm>.}

% structure: p0001 source=h1 target=section reason=page-title

\section{第336封信}

\Emph{\Person{圣巴西略（凯撒勒雅）}}\par

\Person{利巴尼乌斯}致\Person{巴西略}。\par

1. 过了一些时候，一位年轻的卡帕多细亚人到了我这里。一个收获是，他是卡帕多细亚人。但这位卡帕多细亚人乃是出身名门。这是另一收获。此外，他带来了可敬的\Person{巴西略}的一封信。这是最大的收获。你以为我已忘了你。我曾在你的青年时代非常敬重你。我看见你以克己自律与长者争胜，而这是在充满享乐的城市中。我看见你已拥有相当的学问。那时你认为也应去看看雅典，于是你说服了\Person{凯尔苏斯}陪伴你。有福的\Person{凯尔苏斯}，竟能蒙你眷爱！然后你返回，住在家里，我对自己说：不知\Person{巴西略}现在在做什么？他投身于什么工作？是追随古代的演说家，在法庭中操练吗？还是把有福之人的子弟培养成演说家？后来有人来报告我，说你在采纳一种比这一切更好的生活方式，你更多思虑如何赢得天主的友谊，而非成堆的黄金。我祝福了你和卡帕多细亚人：祝福你，因你以此为目标；祝福他们，因他们能以如此高尚的同乡为荣。\par

2. 我知道你所提到的\Person{菲尔穆斯}，到处都始终获胜；因此他是一位大有能力的演说家。但在他所受的一切颂扬中，我不知道曾有像我在你信中所听到的那样高的赞美。因为由你宣称他的名声不亚于任何人，这是何等的光荣啊！\par

看来你是在见到\Person{菲尔米努斯}之前，就把这位年轻人打发到我这里来了；你若是先见了他，你的信不会不提到他。\Person{菲尔米努斯}现在做什么，或打算做什么？他仍急于结婚吗？还是那一切现在都已过去？元老院的职责对他沉重吗？他是否必须留在原地？他还有希望重新投入学业吗？让他给我一个答复，我希望是满意的答复。如果答复令人沮丧，至少可以免去他看见我登门拜访。如果\Person{菲尔米努斯}现在正在雅典，你们的元老们会做什么？他们会派出萨拉米尼亚号去追他吗？你看，只有你的同乡们亏待了我。然而我绝不会停止爱护和称赞卡帕多细亚人。我希望他们对我更好一些，但如果他们继续如此行事，我也将忍受。\Person{菲尔米努斯}和我在一起四个月，没有一日懈怠。你将知道他学到了多少，或许不会抱怨。至于他能否再来这里，我能求助于什么盟友？如果你们的元老们如受过教育的人所应有的那样明理，他们会在第二次情形中尊敬我，因为他们在第一次情形中使我伤心。\par

\SourceRecord{Translated by Blomfield Jackson.；Nicene and Post-Nicene Fathers, Second Series；Vol. 8.；Edited by Philip Schaff and Henry Wace.；Buffalo, NY: Christian Literature Publishing Co.,；1895.；<http://www.newadvent.org/fathers/3202336.htm>.}

% structure: p0001 source=h1 target=section reason=page-title

\section{书信第337号}

\Emph{\Person{圣巴西略（凯撒勒雅）}}\par

\Person{巴西略}致\Person{利巴尼乌斯}。\par

你瞧，又有一位卡帕多细亚人到你那里来了；这是我的一个儿子！然而我现在的地位使所有人都成为我的儿子。因此，他可被视为先前那位的兄弟，值得从我——他的父亲——以及从你——他的导师——得到同样的关注——若果真有可能的话，这些从我这里来的年轻人还能获得更多恩惠。我这样说，并非指阁下不可能再给予你的老同伴们什么，而是因为你的帮助对所有人都如此慷慨。对于这位尚未有经验的青年，若能成为你亲密相知者中的一员，便已足够。我相信你会将他送回，使他配得上我的祈祷以及你在学识与口才上的伟大声誉。与他同行的是一位同龄青年，同样热衷于求学；此子家世良好，与我关系密切。我确信他将处处得到同等善待，尽管他的家境不如其他人优渥。\par

\SourceRecord{Translated by Blomfield Jackson.；Nicene and Post-Nicene Fathers, Second Series；Vol. 8.；Edited by Philip Schaff and Henry Wace.；Buffalo, NY: Christian Literature Publishing Co.,；1895.；<http://www.newadvent.org/fathers/3202337.htm>.}

% structure: p0001 source=h1 target=section reason=page-title

\section{\WorkTitle{第三三八封信}}

\Emph{\Place{凯撒勒雅}的\Person{圣巴西略}}\par

\Person{利巴尼乌斯}致\Person{巴西略}。\par

我知道您会经常写道：\Quote{又是一个卡帕多细亚人！}我预料您会给我送来许多。我深信您到处用您对我的所有美言向父亲们和儿子们施加压力。但我若不提及您那封好信所引发的事，就是不仁慈。当时与我在一起的有不少我们的显贵人物，其中包括非常卓越的\Person{阿利比乌}，\Person{希厄罗克勒斯}的表弟。信使们递上了信。我一声不吭地从头到尾读了一遍；然后带着微笑，显然很高兴地喊道：\Quote{我输了！}\Quote{怎么？何时？何处？}他们问道。\Quote{您为何不为失败而沮丧？}\Quote{我被打败了，}我回答说，\Quote{在优美的书信写作上。巴西略赢了。但我爱他；所以我感到高兴。}听到这话，他们都想从信中亲自听听这场胜利。\Person{阿利比乌}朗读了信，大家倾听。大家一致认为我说得非常正确。然后朗读的人拿着信走了出去，我想是去给别人看。我费了些劲才把信要回来。继续写其他像这样的信；继续赢。这对我来说就是赢。您认为我的服务不是用金钱来衡量，这完全正确。对于一无所有的人来说，渴望接受就足够了。如果我看出任何穷人热爱学问，他就比富人优先。诚然，我从未遇到过这样的教师；但没有什么能阻止我至少在那方面比我的教师有所改进。那么，任何人都不应该因为贫穷而犹豫前来，只要他拥有唯一的资格，知道如何工作。\par

\SourceRecord{Translated by Blomfield Jackson.；Nicene and Post-Nicene Fathers, Second Series；Vol. 8.；Edited by Philip Schaff and Henry Wace.；Buffalo, NY: Christian Literature Publishing Co.,；1895.；<http://www.newadvent.org/fathers/3202338.htm>.}

% structure: p0001 source=h1 target=section reason=page-title

\section{\WorkTitle{书信三三九}}

\Emph{\Person{圣巴西略（凯撒勒雅的）}}\par

\Person{巴西略}致\Person{利巴尼乌斯}。\par

智者有何不能言？更何况是这样一位智者！他的独特技艺是，只要他愿意，就能使大事变小，使小事变大。这就是你在我身上所展示的。我那封又脏又短的信——或许，你这饱享辞藻丰盛的人会如此称呼它——这封信绝不比你现在手中的这封更可容忍，而你却如此盛赞它，说是要被它吃掉，并且把写作的奖赏让给了我！你的所作所为，颇像父亲们参与孩子的游戏，让小家伙们为那些他们让出的胜利而自豪，这对自己毫无损失，对孩子的竞争心却大有裨益。当真，当你拿我开玩笑时，你的话必定带来了无法形容的乐趣！就好像某个\Person{波利达玛斯}或\Person{米隆}拒绝与我进行搏击或摔跤一般！仔细审视之后，我没有发现任何软弱的迹象。因此，那些寻找夸张的人，更惊讶于你能够屈尊降到我的水平来戏耍，胜过你曾带领野蛮人全帆驶过\Place{阿索斯山}。然而，我亲爱的先生，我现在正与\Person{梅瑟}、\Person{厄里亚}以及像他们那样的圣人们相处，他们用蛮族的语言向我讲述他们的故事，我讲出我从他们那里学到的——的确，意思真实，但言辞粗鲁，正如我现在所写的这封信所证明的。如果我曾从你那里学到过什么，如今已随时日遗忘。但请你继续写信给我，从而为通信提出其他话题。你的信将展现你自己，而不会定我的罪。我已经将\Person{安尼修斯}的儿子介绍给你，如同我自己的儿子。如果他是我的儿子，他就是他父亲的儿子，贫穷，且是穷人的儿子。我所说的，这位既是智者又是明哲的人十分明白。\par

\SourceRecord{Translated by Blomfield Jackson.；Nicene and Post-Nicene Fathers, Second Series；Vol. 8.；Edited by Philip Schaff and Henry Wace.；Buffalo, NY: Christian Literature Publishing Co.,；1895.；<http://www.newadvent.org/fathers/3202339.htm>.}

% structure: p0001 source=h1 target=section reason=page-title

\section{书信第三四〇号}

\Emph{\Person{圣巴西略}}\par

\Person{利巴尼乌斯}致\Person{巴西略}。\par

若您长久思忖如何最佳回应我关于您信件的来信，依我判断，您不可能比现在所写的这封信表现得更好。您称我为诡辩家，并断言诡辩家之能事在于化小为大、化大为小。您又坚称，我那封信的目的是要证明您的那封信是佳作，而它本非佳作，且它并不比您最后寄来的那封更好——总而言之，您毫无表达能力，因为您手头钻研的书籍并未产生这种效果，您曾拥有的口才也已尽失。然而，为了证明这一点，您却使这封您正在贬斥的书信如此出色，以至于我的访客在听人诵读时，不禁跃起赞叹。我感到惊讶的是，您试图通过这封信来贬低前信，声称前信与它相似，实际上您却借此恭维了前信。要实现您的目的，您本该把这封信写得更糟，这样才能诋毁前信。但我认为，您不会违背真理。若您故意写得拙劣，而不发挥您的能力，那才是违背真理。您的作风应当是不去挑剔值得称赞的事物，以免在试图化大为小时，您的行为使自己沦入诡辩家的行列。请继续阅读您那些风格虽逊、意义却佳的著作。无人阻止您。至于那些永远属于我、也曾属于您的原则，其根基既已存留，也将继续存留，只要您存在。即使您浇灌它们甚微，任何漫长岁月也无法将它们彻底摧毁。\par

\SourceRecord{Translated by Blomfield Jackson.；Nicene and Post-Nicene Fathers, Second Series；Vol. 8.；Edited by Philip Schaff and Henry Wace.；Buffalo, NY: Christian Literature Publishing Co.,；1895.；<http://www.newadvent.org/fathers/3202340.htm>.}

% structure: p0001 source=h1 target=section reason=page-title

\section{第341封书信}

\Emph{凯撒勒雅的圣巴西略}\par

李巴尼乌斯致巴西略。\par

你仍未停止对我生气，因此我写信时战栗不已。你若在意，我亲爱的先生，为何不写信？你若仍在生气——这实在与任何理性灵魂，尤其与你的灵魂格格不入——为何当你向他人宣讲不可含怒到日落时，自己却含怒多日？或许你意在剥夺我聆听你甜蜜声音的权利，以此惩罚我？不，卓越的先生，请温和些，让我享受你的金口吧。\par

\SourceRecord{Translated by Blomfield Jackson.；Nicene and Post-Nicene Fathers, Second Series；Vol. 8.；Edited by Philip Schaff and Henry Wace.；Buffalo, NY: Christian Literature Publishing Co.,；1895.；<http://www.newadvent.org/fathers/3202341.htm>.}

% structure: p0001 source=h1 target=section reason=page-title

\section{书信三四二}

\Emph{\Person{圣巴西略}·\Place{凯撒勒雅}}\par

\Person{圣巴西略}致\Person{里巴尼乌斯}。\par

凡爱慕玫瑰的人，正如爱好美者所期望的那样，甚至连那花朵从中绽放的刺也不讨厌。我甚至听人谈论玫瑰时说过，或许出于戏谑，或许甚至是认真的，说大自然给这花朵配上了那些纤细的刺，如同爱者之间的爱之刺，为的是通过这些巧妙设计的刺，激发采摘者更强烈的渴望。但我将这玫瑰引入书信，是何用意？你毋需被告知，因为你记得你自己的信。那信确实有玫瑰的芬芳，并以优美的言辞向我展开了整个春天；但它也带有某些挑剔和指责我的刺。不过，就连你话语中的刺也令我愉悦，因为它在我心中燃起更强烈的对你友谊的渴望。\par

\SourceRecord{Translated by Blomfield Jackson.；Nicene and Post-Nicene Fathers, Second Series；Vol. 8.；Edited by Philip Schaff and Henry Wace.；Buffalo, NY: Christian Literature Publishing Co.,；1895.；<http://www.newadvent.org/fathers/3202342.htm>.}

% structure: p0001 source=h1 target=section reason=page-title

\section{\WorkTitle{第343封书信}}

\Emph{\Person{凯撒勒雅的圣巴西略}}\par

\Person{利巴尼乌斯}致\Person{巴西略}。\par

如果这些话语出自未经训练的唇舌，那么你若将它们雕琢，又会成为何等样人？你的唇上涌流着比泉水更美的言辞泉源。而我，与此相反，若得不到每日的浇灌，便静默无声。\par

\SourceRecord{Translated by Blomfield Jackson.；Nicene and Post-Nicene Fathers, Second Series；Vol. 8.；Edited by Philip Schaff and Henry Wace.；Buffalo, NY: Christian Literature Publishing Co.,；1895.；<http://www.newadvent.org/fathers/3202343.htm>.}

% structure: p0001 source=h1 target=section reason=page-title

\section{书信 344}

\Person{圣巴西略}（凯撒勒雅的）\par

巴西略致利巴尼亚。\par

我因怯懦和无知，不敢经常写信给像你这样有学问的人。但你的持久沉默则是另一回事。这有什么借口可言呢？如果有人考虑到你生活在文字之中，却很少写信给我，他将责备你忘记了我。一个善于说话的人，并不是不擅长写作。如果一个这样有天赋的人保持沉默，显然他的行为出于忘记或轻视。然而，我将以问候来回报你的沉默。再见，尊贵的阁下。如果你喜欢，就写信。如果你更喜欢，就不要写。\par

\SourceRecord{Translated by Blomfield Jackson.；Nicene and Post-Nicene Fathers, Second Series；Vol. 8.；Edited by Philip Schaff and Henry Wace.；Buffalo, NY: Christian Literature Publishing Co.,；1895.；<http://www.newadvent.org/fathers/3202344.htm>.}

% structure: p0001 source=h1 target=section reason=page-title

\section{书信第三四五}

\Emph{\Person{凯撒勒雅的圣巴西略}}\par

\Person{利巴尼乌斯}致\Person{巴西略}。\par

我认为，为自己辩护为何这么久才开始给你写信，比为自己现在开始写信找借口更为必要。我还是那个你每次露面时总会跑上去迎接的人，最喜悦地聆听你雄辩的洪流；高兴地听你讲话；很难拽开自己；对我的朋友们说：\Quote{这个人远比\Person{阿刻罗俄斯}的女儿们优越，因为像她们一样，他抚慰人，却不像她们那样伤害人。}确实，不伤害并不是什么大事；但这个人的歌声对听者来说是一种实在的收获。我处在这样的心境中，自以为被关爱，又似乎能言善辩，却不敢写信，这是一个极端懒惰的人的标志，同时也是在惩罚自己。因为很明显，你会用一封漂亮的长信来回复我这封蹩脚的小信，并会注意不再委屈我。听到这句话，我想，许多人会喊叫起来，并会围拢过来大喊：\Quote{什么！\Person{巴西略}做了错事——哪怕是一点点错事？那\Person{埃阿科斯}、\Person{米诺斯}和他的兄弟也一样。}在其他方面我承认你赢了。谁看见你而不嫉妒你呢？但在一件事上你得罪了我；如果我提醒你这件事，使那些为此愤慨的人不要公开喧嚷。从来没有人到你那里去求一个容易给予的恩惠而空手而归的。但我却是那些渴求恩惠却没有得到的人之一。那么我求了什么呢？当我和你待在营地时，我常常渴望借着你的智慧进入\Person{荷马}灵感的深处。\Quote{如果全部可能是不可能的，}我说，\Quote{请你带我到我所求的一部分。}我渴望那一部分，其中希腊人处境不顺时，\Person{阿伽门农}用礼物向被他侮辱的人献殷勤。当我这样说时，你笑了，因为你不能否认，你如果愿意就能做到，但你不愿意给予。在你的朋友们因我说你做了错事而愤慨时，我真的看起来像是被冤枉了吗？\par

\SourceRecord{Translated by Blomfield Jackson.；Nicene and Post-Nicene Fathers, Second Series；Vol. 8.；Edited by Philip Schaff and Henry Wace.；Buffalo, NY: Christian Literature Publishing Co.,；1895.；<http://www.newadvent.org/fathers/3202345.htm>.}

% structure: p0001 source=h1 target=section reason=page-title

\section{第346封信}

\Emph{\Person{圣巴西略·凯撒勒雅}}\par

\Person{利巴尼乌斯}致\Person{巴西略}。\par

你自己来判断，我是否在以学问教导你所派来的青年方面添加了什么。我希望这添加，无论多么微小，因着你对我的友谊，而得以算为伟大。但是，既然你对学问的赞美不如对节制以及拒绝将我们的灵魂交付给可耻的享乐的赞美之多，他们便将主要注意力放在后者上，并且正如他们应当做的，怀着对遣送他们到此的朋友的应有记念而生活。\par

因此，请接纳属于你自己的一切，并赞美那些以其生活方式为你我都增光的人。但是，请求你对他们有帮助，如同请求一位父亲对他的孩子们有帮助一样。\par

\SourceRecord{Translated by Blomfield Jackson.；Nicene and Post-Nicene Fathers, Second Series；Vol. 8.；Edited by Philip Schaff and Henry Wace.；Buffalo, NY: Christian Literature Publishing Co.,；1895.；<http://www.newadvent.org/fathers/3202346.htm>.}

% structure: p0001 source=h1 target=section reason=page-title

\section{\WorkTitle{第三四七封信}}

\Emph{\Person{圣巴西略·凯撒勒雅}}\par

\Person{利巴尼乌斯}致\Person{巴西略}。\par

每一位主教都是极难从中得到任何东西的人。你在学识上越超越他人，你越使我害怕你会拒绝我的请求。我需要一些椽木。任何其他智者都会称它们为桩子或杆子，倒不是因为他需要桩子或杆子，而是为了炫耀他的小字眼，而非出于任何实际需要。你若不予供应，我将不得不在露天过冬。\par

\SourceRecord{Translated by Blomfield Jackson.；Nicene and Post-Nicene Fathers, Second Series；Vol. 8.；Edited by Philip Schaff and Henry Wace.；Buffalo, NY: Christian Literature Publishing Co.,；1895.；<http://www.newadvent.org/fathers/3202347.htm>.}

% structure: p0001 source=h1 target=section reason=page-title

\section{\WorkTitle{书信三四八}}

\Emph{凯撒勒雅的圣巴西略}\par

\Person{巴西略}致\Person{利巴尼乌斯}。\par

如果\Quote{\OriginalTerm{获取}{\GreekText{γριπίζειν}}}与获利同义，而这正是您的诡辩巧思从\Person{柏拉图}的深处所获得之短语的含义，那么，亲爱的先生，请考虑谁更难以被说服：是我，被您的书信技巧如此嘲弄之人，还是那些诡辩士一族，他们的技艺是从言辞中赚钱。有哪位主教曾通过言语征收贡赋？有哪位主教曾使他的门徒纳税？是您使您的言辞成为可买卖之物，就像糖果商制作蜜饼一般。看您如何使这位老人跳跃不已！然而，对于您如此夸耀您的演说，我已下令供应与温泉关战士数量相等的椽子，全部优良长度，正如\Person{荷马}所说，\Quote{长影，}神圣的\Place{阿尔法约}承诺将其归还。\par

\SourceRecord{Translated by Blomfield Jackson.；Nicene and Post-Nicene Fathers, Second Series；Vol. 8.；Edited by Philip Schaff and Henry Wace.；Buffalo, NY: Christian Literature Publishing Co.,；1895.；<http://www.newadvent.org/fathers/3202348.htm>.}

% structure: p0001 source=h1 target=section reason=page-title

\section{书信第三四九封}

\Emph{\Person{圣巴西略} \Place{凯撒勒雅}}\par

\Person{利巴尼乌斯}致\Person{巴西略}。\par

\Person{巴西略}，你还不肯罢手，把这\Person{缪斯}的神圣居所塞满卡帕多细亚人，这些人身上带着霜雪和\Place{卡帕多细亚}所有好东西的气息？他们几乎也把我变成了卡帕多细亚人，总是唱着\Quote{敬礼}。\par

我必须忍耐，因为是\Person{巴西略}下的命令。但你要知道，我正在仔细研究那地方的风俗习惯，我打算把这些人都变成我的\Person{卡利俄佩}的高贵与和谐，使他们在你看来像是从野鸽变成了家鸽。\par

\SourceRecord{Translated by Blomfield Jackson.；Nicene and Post-Nicene Fathers, Second Series；Vol. 8.；Edited by Philip Schaff and Henry Wace.；Buffalo, NY: Christian Literature Publishing Co.,；1895.；<http://www.newadvent.org/fathers/3202349.htm>.}

% structure: p0001 source=h1 target=section reason=page-title

\section{书信350}

\Emph{\Person{圣巴西略（凯撒勒雅）}}\par

\Person{巴西略}致\Person{利巴尼乌斯}。\par

你的烦恼已经结束。就让这成为我信的开端吧。继续嘲笑和辱骂我和我的，无论是笑还是认真。为什么要提到霜或雪，当你可以在嘲笑中享乐？至于我，\Person{利巴尼乌斯}，为了让你开怀大笑，我写这封信时裹在雪白的面纱中。当你拿起信时，你会感到它多么冰冷，以及它如何象征寄信人的状况——待在家里，不能把头伸出门外。因为我的房子像坟墓，直到春天来临，将我们从死亡带回生命，再次赐给我们，如同植物，生存的恩赐。\par

\SourceRecord{Translated by Blomfield Jackson.；Nicene and Post-Nicene Fathers, Second Series；Vol. 8.；Edited by Philip Schaff and Henry Wace.；Buffalo, NY: Christian Literature Publishing Co.,；1895.；<http://www.newadvent.org/fathers/3202350.htm>.}

% structure: p0001 source=h1 target=section reason=page-title

\section{\WorkTitle{书信351}}

\Emph{圣巴西略·凯撒勒雅}\par

\Person{巴西略}致\Person{利巴尼乌斯}。\par

许多从你那里来的人，都赞赏你的雄辩能力。他们都说，这是这方面一个非常出色的范例，也是一场非常伟大的比赛，据他们称，结果所有人都挤在一起，全城没有人出现，只有\Person{利巴尼乌斯}一人在竞技场上，所有人，无论老少，都在听。因为没有人愿意缺席——没有一位达官贵人，没有一位杰出士兵，没有一位工匠。连妇女也急忙赶来观看这场较量。那是什么呢？是哪篇演说聚集了如此庞大的集会？我听说它描写了一个暴躁脾气的人。请速将此篇备受赞赏的演说寄来，好让我也能加入赞美你口才的行列。如果我在未见你作品的情况下就已是\Person{利巴尼乌斯}的赞美者，那么在得到赞美的理由之后，我又会成为怎样的人呢？\par

\SourceRecord{Translated by Blomfield Jackson.；Nicene and Post-Nicene Fathers, Second Series；Vol. 8.；Edited by Philip Schaff and Henry Wace.；Buffalo, NY: Christian Literature Publishing Co.,；1895.；<http://www.newadvent.org/fathers/3202351.htm>.}

% structure: p0001 source=h1 target=section reason=page-title

\section{信件第三五二}

\Emph{\Person{圣巴西略}（\Place{凯撒勒雅}）}\par

\Person{利巴尼乌斯}致\Person{巴西略}。\par

看啊！我把我那篇演讲词送给你了，我真是汗流浃背！我怎能不如此呢？当我把这篇演讲词送给一位能以修辞技巧证明\Person{柏拉图}的智慧和\Person{德摩斯梯尼}的才能是徒受赞誉的人时，我真是如同蚊子比大象。我一想起你将审阅我的作品的那一天，就战战兢兢，几乎失去理智！\par

\SourceRecord{Translated by Blomfield Jackson.；Nicene and Post-Nicene Fathers, Second Series；Vol. 8.；Edited by Philip Schaff and Henry Wace.；Buffalo, NY: Christian Literature Publishing Co.,；1895.；<http://www.newadvent.org/fathers/3202352.htm>.}

% structure: p0001 source=h1 target=section reason=page-title

\section{书信353}

\Emph{凯撒勒雅的圣巴西略}\par

\Person{巴西略}致\Person{利巴尼乌斯}。\par

我已阅读了你的演说，并极为欣赏。哦，缪斯！哦，学问！哦，\Place{雅典}！你们不把什么赐给爱慕你们的人呢！那些与你们共度短暂时光的人，会收获何等果实啊！哦，你们那丰沛涌流的泉源！所有畅饮它的人，都展现出何等的人物啊！我仿佛在你的演说中看到了那人本人，与他那喋喋不休的小妇人为伴。\Person{利巴尼乌斯}在地上写下了一个活生生的故事，唯有他赐予了他的言辞生命的恩赐。\par

\SourceRecord{Translated by Blomfield Jackson.；Nicene and Post-Nicene Fathers, Second Series；Vol. 8.；Edited by Philip Schaff and Henry Wace.；Buffalo, NY: Christian Literature Publishing Co.,；1895.；<http://www.newadvent.org/fathers/3202353.htm>.}

% structure: p0001 source=h1 target=section reason=page-title

\section{\WorkTitle{书信 354}}

\Emph{\Place{凯撒勒雅}的\Person{圣巴西略}}\par

\Person{利巴尼乌斯}致\Person{巴西略}。\par

现在，我认清了人们对我描述！巴西略称赞了我，我便被欢呼为万物的征服者！既然我得到了你的投票，我便有资格以一个人傲慢地俯视世界的那种骄傲步伐行走。你写了一篇反对酗酒的演说辞。我想读一读。但我不愿尝试说任何机巧的话。当我看到你的演说时，它将教给我表达的艺术。\par

\SourceRecord{Translated by Blomfield Jackson.；Nicene and Post-Nicene Fathers, Second Series；Vol. 8.；Edited by Philip Schaff and Henry Wace.；Buffalo, NY: Christian Literature Publishing Co.,；1895.；<http://www.newadvent.org/fathers/3202354.htm>.}

% structure: p0001 source=h1 target=section reason=page-title

\section{\WorkTitle{书信三五五}}

\Emph{\Place{凯撒勒雅}的圣\Person{巴西略}}\par

\Person{利巴尼乌斯}致\Person{巴西略}。\par

\Person{巴西略}，你如今住在\Place{雅典}吗？你忘记自己了吗？\Place{凯撒勒雅}人的子弟无法忍受听到这些事情。我的舌头对这些还不习惯。就好像我走在危险的地面上，被这些新奇的声响所震惊，它就对我——它的父亲说：\Quote{我的父亲，你从未教过这个！这个人就是\Person{荷马}，或\Person{柏拉图}，或\Person{亚里士多德}，或\Person{苏萨里昂}。他无所不知。}我的舌头到此为止。我只希望，\Person{巴西略}，你能以同样的方式赞美我！\par

\SourceRecord{Translated by Blomfield Jackson.；Nicene and Post-Nicene Fathers, Second Series；Vol. 8.；Edited by Philip Schaff and Henry Wace.；Buffalo, NY: Christian Literature Publishing Co.,；1895.；<http://www.newadvent.org/fathers/3202355.htm>.}

% structure: p0001 source=h1 target=section reason=page-title

\section{\WorkTitle{书信 356}}

\Emph{\Person{凯撒勒雅的圣巴西略}}\par

\Person{巴西略}致\Person{里巴尼乌斯}。\par

我很高興收到你的來信，但當你要求我回信時，我感到為難。面對如此雅致的阿提卡方言，除非我承認——並且欣喜地承認——我是漁夫們的門徒，我還能說什麼呢？\par

\SourceRecord{Translated by Blomfield Jackson.；Nicene and Post-Nicene Fathers, Second Series；Vol. 8.；Edited by Philip Schaff and Henry Wace.；Buffalo, NY: Christian Literature Publishing Co.,；1895.；<http://www.newadvent.org/fathers/3202356.htm>.}

% structure: p0001 source=h1 target=section reason=page-title

\section{\WorkTitle{书信第357号}}

\Emph{\Person{圣巴西略} \Place{凯撒勒雅}}\par

\Person{利巴尼乌斯}致\Person{巴西略}。\par

是什么使\Person{巴西略}反对这封信，即哲学的证明？我学会了从你那里取笑，但尽管如此，你的取笑是可敬的，可以说是老成持重。但是，凭着我们真正的友谊，凭着我们共同的消遣，我恳求你，消除你的信所带来的苦恼……毫无二致。\par

\SourceRecord{Translated by Blomfield Jackson.；Nicene and Post-Nicene Fathers, Second Series；Vol. 8.；Edited by Philip Schaff and Henry Wace.；Buffalo, NY: Christian Literature Publishing Co.,；1895.；<http://www.newadvent.org/fathers/3202357.htm>.}

% structure: p0001 source=h1 target=section reason=page-title

\section{信函358}

\Emph{\Person{圣巴西略}·\Place{凯撒勒雅}}\par

\Person{利巴尼乌斯}致\Person{巴西略}。\par

啊，但愿能回到从前，那时我们彼此亲密无间！如今我们不幸地分离了！你们彼此相伴，我却无一人能替代你。我听说\Person{阿尔西莫}在年迈之际竟冒险做出年轻人的举动，匆忙赶往\Place{罗马}，却将逗留于此照料这些青年的劳苦强加于你。你一向如此仁慈，必不会为此生气。你甚至未曾因我不得不先写信而恼怒。\par

\SourceRecord{Translated by Blomfield Jackson.；Nicene and Post-Nicene Fathers, Second Series；Vol. 8.；Edited by Philip Schaff and Henry Wace.；Buffalo, NY: Christian Literature Publishing Co.,；1895.；<http://www.newadvent.org/fathers/3202358.htm>.}

% structure: p0001 source=h1 target=section reason=page-title

\section{\WorkTitle{第三五九封书信}}

\Emph{圣凯撒勒雅的巴西略}\par

\Person{巴西略}致\Person{利巴尼乌斯}。\par

你将古人的所有技艺都包含在了自己的心智中，却如此沉默，甚至不让我在信函中有所得。我，如果\Person{代达洛斯}的技艺仍然安全，我本会造出\Person{伊卡洛斯}的翅膀来到你身边。但蜡不可托付给太阳，所以，我不送你\Person{伊卡洛斯}的翅膀，而是送你话语以证明我的情谊。话语的本性是指明内心的爱。目前，只有话语。你随意处置它们，虽然你拥有你所拥有的一切权能，却保持沉默。但请将涌自你口中的话语之泉赐予我。\par

\SourceRecord{Translated by Blomfield Jackson.；Nicene and Post-Nicene Fathers, Second Series；Vol. 8.；Edited by Philip Schaff and Henry Wace.；Buffalo, NY: Christian Literature Publishing Co.,；1895.；<http://www.newadvent.org/fathers/3202359.htm>.}

% structure: p0001 source=h1 target=section reason=page-title

\section{书信360}

\Emph{\Person{凯撒勒雅的圣巴西略}}\par

论天主圣三、降生成人、圣人转求及其圣像。\par

按照我们从天主所获得的基督徒无瑕信仰，我宣认并同意，我相信惟一的天主，全能的圣父；天主圣父、天主圣子、天主圣神；我钦崇并朝拜一个天主，即那三位。我宣认圣子降生成人的\OriginalTerm{经济}{\GreekText{οἰκονομία}}，并且，按照肉体生下祂的圣玛利亚是天主之母。我也承认圣宗徒、先知和殉道者；我呼求他们向天主祈求，好藉着他们——即藉着他们的转求——仁慈的天主对我开恩，并为我的罪赐予赎价。因此，我也尊敬并亲吻他们圣像的容貌，因为这些圣像由圣宗徒传授下来，并未被禁止，反而存在于我们所有教堂中。\par

\SourceRecord{Translated by Blomfield Jackson.；Nicene and Post-Nicene Fathers, Second Series；Vol. 8.；Edited by Philip Schaff and Henry Wace.；Buffalo, NY: Christian Literature Publishing Co.,；1895.；<http://www.newadvent.org/fathers/3202360.htm>.}

% structure: p0001 source=h1 target=section reason=page-title

\section{第366封书信}

\Emph{\Person{圣巴西略} \Place{凯撒勒雅}}\par

\Emph{\Person{巴西略} 致隐修士 \Person{乌尔比奇}，论节欲}\par

你为我们作出精确的定义，好让我们不仅认识节欲，也认识它的果实，这做得很好。它的果实是与天主的共融。因为不被败坏，就是与天主有分；正如被败坏，就是与世俗为伴。节欲是弃绝肉身，并向天主宣认。它远离一切属死的事物，如同拥有天主圣神的身体。它没有竞争与嫉妒，并使我们与天主结合。凡爱肉身的人，必嫉妒他人。凡未将败坏的疾病引入内心的人，将来必能胜任任何劳苦，即使他死于肉身，却活在不朽中。确实，若我正确领悟此事，天主在我看来便是节欲，因为祂一无所欲，却拥有万有在自身内。祂不追求什么，也没有眼目的感觉；一无所缺，祂在各方面都是圆满与完满的。贪欲是灵魂的疾病，而节欲是它的健康。并且，节欲不应仅被视为某一种类，例如在感官之爱方面。它应被视为灵魂以邪恶方式贪求的一切事上，不满足于所需之物。嫉妒因黄金而起，无数不义因其他贪欲而起。不醉酒是节欲。不过饱是节欲。克制肉身是节欲，并使邪恶思想服从，每当灵魂被任何虚假恶劣的幻象扰乱，内心被虚妄的挂虑分心时。节欲使人自由，它既是良药又是力量，因为祂不教导节制，而是赐予节制。节欲是天主的恩宠。耶稣仿佛是节欲本身，当祂成了陆地和海洋的光时；因为祂既未被大地也未被海洋承载，正如祂在海面上行走，祂也没有压沉大地。因为若死亡来自败坏，而不死来自没有败坏，那么耶稣所成就的并非有死性，而是神性。祂以独特方式饮食，而不转化祂的食物。在祂内节欲的力量如此强大，以至于祂的食物在祂内没有腐败，因为祂没有败坏。只要我们内有一点节欲，我们就比一切都高。我们被告知，天使因贪欲被逐出天堂，变得不节欲。他们被击败，并非主动降下。如果我想到的那种眼睛在那里，那灾祸又能造成什么影响呢？因此我说，如果我们有一点忍耐，不爱世界，而爱天上的生命，我们就将出现在我们心智所指向的地方。因为显然，心智是看见未见之事的眼睛。因为我们说：\Quote{心智看见；心智听见。}我写得冗长，尽管在你看来可能不多。但我所说的一切都有深意，你读到时自会明白。\par

\SourceRecord{Translated by Blomfield Jackson.；Nicene and Post-Nicene Fathers, Second Series；Vol. 8.；Edited by Philip Schaff and Henry Wace.；Buffalo, NY: Christian Literature Publishing Co.,；1895.；<http://www.newadvent.org/fathers/3202366.htm>.}
